% 本文件由 LaTeX 排版 Agent 根据有效 translation.json 自动生成。
% author=John Chrysostom; work_id=2102; title=罗马书讲道集

\chapter{\WorkTitle{罗马书讲道集}}

% structure: p0001 source=h1 target=omitted reason=page-title-already-rendered-by-parent

讲道1\newline{}讲道2\newline{}讲道3\newline{}讲道4\newline{}讲道5\newline{}讲道6\newline{}讲道7\newline{}讲道8\newline{}讲道9\newline{}讲道10\newline{}讲道11\newline{}讲道12\newline{}讲道13\newline{}讲道14\newline{}讲道15\newline{}讲道16\newline{}讲道17\newline{}讲道18\newline{}讲道19\newline{}讲道20\newline{}讲道21\newline{}讲道22\newline{}讲道23\newline{}讲道24\newline{}讲道25\newline{}讲道26\newline{}讲道27\newline{}讲道28\newline{}讲道29\newline{}讲道30\newline{}讲道31\newline{}讲道32\par

\SourceRecord{Translated by J. Walker, J. Sheppard and H. Browne, and revised by George B. Stevens.；Nicene and Post-Nicene Fathers, First Series；Vol. 11.；Edited by Philip Schaff.；Buffalo, NY: Christian Literature Publishing Co.,；1889.；<http://www.newadvent.org/fathers/2102.htm>.}

% structure: p0001 source=h1 target=section reason=page-title

\section{罗马人书讲道 一}

\BibleRef{罗 1:1-2}\par

\BibleQuote{保禄，耶稣基督的仆人，蒙召为宗徒，被选拔去传天主的福音——这福音是天主先前藉自己的先知在圣经上所预许的。}\par

\Person{梅瑟}写了五本书，却没有在任何地方署上自己的名字；后来编写历史的人也都没有署自己的名字；\Person{玛窦}、\Person{若望}、\Person{马尔谷}、\Person{路加}也没有。但是真福\Person{保禄}在每封书信中都署上了自己的名字。这是为什么呢？因为前几位是写给在场的人，既然他们在场，表明自己就是多余的。但\Person{保禄}是从远方以书信形式寄出他的著作，因此署上名字是必要的。不过，在\WorkTitle{希伯来书}中他并没有这样做，这同样出于他明智的判断。因为希伯来人对他怀有偏见，唯恐一听到名字，他们就会堵塞一切听他讲论的通道，所以他巧妙地隐去名字以赢得他们的注意。至于有些先知和\Person{撒罗满}署上了自己的名字，这一点我留给你们日后进一步探究：为什么他们中有些人愿意如此，有些人却不然。因为你们不应凡事都从我学习，你们自己也要努力并进一步探究，以免变得迟钝。\par

\Quote{保禄，耶稣基督的仆人。} 为什么天主更改了他的名字，称扫禄为保禄？这是为了使他即使在这一点上也不逊于宗徒们，而且让门徒之长所拥有的那种卓越，他也能获得\BibleRef{谷 3:16}，并得以建立与他们的更紧密的联系。他称自己为基督的仆人，但不仅止于此；因为仆役有多种。一种是由于受造，经上说：\Quote{因为一切都是祢的仆人}\BibleRef{咏 118:91}；又按另一种意义，经上说：\Quote{我的仆人拿步高}\BibleRef{耶 25:9}——因为作品是创造它的那位仆役。另一种来自信仰，关于此经上说：\Quote{感谢天主，虽然你们曾作罪恶的奴隶，却从心里听从了那传授给你们的教理，从罪恶中解放出来后，就成了正义的奴隶。}\BibleRef{罗 6:17} 另一种来自公民隶属关系\OriginalTerm{公民隶属}{\GreekText{πολιτείας}}，按此经上说：\Quote{我的仆人梅瑟死了}\BibleRef{苏 1:2}；事实上，所有犹太人都曾是仆人，但\Person{梅瑟}因在团体中特别出众而更显光明。既然\Person{保禄}在所有奇妙仆役的形式上都是仆人，他就把这当作最大尊荣的称号，说：\Quote{耶稣基督的仆人。} 接着他从最低的开始往上，展示了属于降生奥迹的名号。因为当祂由童贞女怀孕时，天使从天带来了耶稣的名字；祂又被称作基督，因为被傅油，这也属于肉身。那么，可以问，祂是用什么油被傅的？祂不是用油被傅，而是用圣神。圣经有称这样人为\Quote{受傅者}的例子；因为傅油的关键在于圣神，油只为傅油而用。圣经在哪里称那些没有用油傅过的人为\Quote{受傅者}？经上说：\Quote{不可触犯我的受傅者，不可伤害我的先知}\BibleRef{咏 104:15}，但那时用油傅油的规定甚至尚未存在。\par

\Quote{蒙召的宗徒。} 他自称\Quote{蒙召的}，以此显示他的坦诚\OriginalTerm{坦诚}{\GreekText{εὐγνωμοσύνην}}，并表明不是他自己寻求而找到的，而是蒙召后前来遵从的。至于信友们，他称他们为\Quote{蒙召为圣徒的}，但他们在信德上蒙召，而他则受托了另一项不同的事，即宗徒职务，这职务充满无数祝福，同时比所有恩赐更大并包含一切恩赐。\par

对此还有什么可说的呢？无非是基督在世时所行的一切，祂在离去时都托付给了他们。保禄也高声宣告，论及此事并彰显宗徒职务的尊贵：\Quote{我们为基督作大使，好像天主藉着我们劝勉；}即代替基督。\Quote{被分别出来归于天主的福音。}\BibleRef{格后 5:20}因为如同在一座房屋里，各人被分派做不同的工作；同样，在教会中也有各种职务的分配。他在这里似乎暗示，他并非仅凭抽签被选，而是从起初就已被预定担任此职；这也正如耶肋米亚所说，天主论及他：\Quote{在你出母胎以前，我已祝圣了你，立你为万民的先知。}\BibleRef{耶 1:5}因为他写信给一座好虚荣、事事自大的城市，所以用尽各种方式表明他的蒙选出自天主。因为是天主亲自召叫了他，亲自分别了他。他这样做，是为了使书信值得信赖，并易于被人接受。\Quote{归于天主的福音。}那么，不独玛窦是福音作者，不独马尔谷是福音作者，同样，此人也不独是宗徒，他们也都是；即使他被称为这身份的特出者，他们也一样。他称之为福音，不仅是因为已经成就的善事，也是因为将来的事。他怎么说福音\Quote{归于天主的福音}由他自己宣讲呢？因为他说：\Quote{被分别出来归于天主的福音}——因为圣父在福音之前已经显现。然而即使祂已显现，也只是向犹太人，而且并非向所有犹太人，如所应当的。因为他们既不认识祂为父，对祂也怀有许多不配的想法。因此基督也说：\Quote{真正的朝拜者}将要来到，而且\Quote{圣父就是寻找这样朝拜祂的人。}\BibleRef{若 4:23}但后来，祂与圣子一同向全世界启示了自己，这也是基督预先说过的：\Quote{使他们认识你，唯一的真天主，和你所派遣的耶稣基督。}\BibleRef{若 17:3}但他称之为天主的\Quote{福音}，是要在开始时使听者感到喜乐。因为他带来的消息并非使人面容忧愁，如先知们带着指控、责难和责备而来，而是带来喜讯，即天主的\Quote{福音}，无数的永恒不变之福的宝藏。\par

第2节。\BibleQuote{这福音是天主先前藉着祂的先知在圣经中所预许的。}\BibleRef{罗 1:2}\par

因为主，他说，\Quote{必赐给传报佳音者大有能力的话语}\BibleRef{咏 67:12}；又说，\Quote{那传报和平福音的脚是多么美丽啊}\BibleRef{依 52:7}。请看，在旧约中福音的名称和性质都已明确表明。因为他说，我们不仅用言语传报，也用行动传报；因为这事不是人所能及的，而是神圣的、不可言喻的、超乎一切本性的。既然他们也以新奇为罪名攻击福音，他便表明福音比希腊人的事更古老，并且早已在先知书中预述了。如果天主因那些不愿接受的人而没有从起初就赐下福音，但愿意的人确实听到了它。祂说：\Quote{你们的父亲亚巴郎曾欢欣希望看到我的日子，他看见了，也高兴了。}\BibleRef{若 8:56}那么祂怎么说：\Quote{许多先知和义人渴望看到你们所看见的，却没有看到呢？}\BibleRef{玛 13:17}祂的意思并非如你们所见所闻，即看到肉身本身和眼前的奇迹。但请容许你们观察，这些事是在多么久远以前被预言的。因为当天主将要公开行某些大事时，祂会预先在很久以前宣布，以训练人的耳朵，使他们届时能接受。\par

\Quote{在圣经中。}因为先知们不仅说话，还写下他们所说的话；他们不仅写下，还用行动预示，如亚巴郎带领依撒格，梅瑟举起铜蛇，他向阿玛肋克伸开双手，以及他奉献逾越节羔羊。\par

第3节。\BibleQuote{论及祂的儿子，按肉身说，是从达味的后裔生的。}\BibleRef{罗 1:3}\par

你做什么，噢，保禄，你先把我们的灵魂高举、提升，使伟大而不可言喻的事在我们眼前展现，谈论福音，而且是天主的福音，引入先知们的合唱，显示他们多年以前就预告了那些要来的事；然后为什么又将我们带回到达味？你在谈论一个人，给我说说，并称祂为叶瑟的儿子之父？这些事怎能与你刚才所说的相称呢？是的，完全相称。因为他说，我们的谈论并非关于一个纯粹的人。正是为此我才加上了\Quote{按肉身说}，以暗示同一位也有按圣神的出生。那么他为什么从较低的出生开始，而不是从较高的呢？因为玛窦、路加和马尔谷也从这里开始。因为凡要将人一步步引向天堂的人，必须从低处向上领他们。实际的救恩计划也是如此。首先，他们看见祂在地上是一个人，然后他们才领悟祂是天主。因此，正如祂自己建构祂的教导，祂的门徒也铺设了通向那里的道路。所以，按肉身的出生在他的语言中排在首位，并不是因为这是最先的，而是因为他要将听者从这一点引向更高之处。\par

第4节。\BibleQuote{按圣德的神性，藉着从死者中复活，被立为具有大能的天主之子，就是耶稣基督。}\BibleRef{罗 1:4}\par

所说的因措辞紧凑而显得隐晦，故有必要加以分解。那么，他所说的是什么呢？他说，我们宣讲那一位，祂是由达味所出的。但这是显而易见的。那么，这降生成人的\Quote{位格}也是天主子，这从何而见呢？首先，从先知而来；因此他说：\Quote{这福音是天主先前藉先知在圣经上所预许的}\BibleRef{罗 1:2}。这种论证方式并不薄弱。其次，从祂的生成方式本身：他藉着说\Quote{按肉身是由达味的后裔}而表明这一点；因为祂打破了自然的规律。第三，从祂所行的奇迹，显示出极大的德能，因为\Quote{按德能}即指此。第四，从祂赐给相信祂之人的圣神，并藉此使他们全都成圣，因此他说\Quote{按圣德的圣神}。因为唯有天主才能赐予这样的恩赐。第五，从复活而来；因为祂是首位且唯独祂自己复活了：这事本身被视为超越一切的奇迹，足以堵住那些甚至无耻之徒的口。因为祂说：\Quote{拆毁这座圣殿，三天内我要把它重建起来}\BibleRef{若 2:19}；又说：\Quote{当你们举起人子以后，你们便知道我就是那一位}\BibleRef{若 8:28}；又说：\Quote{这世代寻求征兆；除了约纳的征兆外，必不给它任何征兆}\BibleRef{玛 12:39}。那么，\Quote{被立定}是什么意思呢？即被显示、被彰显、被判定、被众人以情感和共识所承认：藉先知，藉肉身中奇妙的诞生，藉奇迹中的德能，藉圣神（祂藉此赐下圣化），藉复活（祂藉此终结了死亡的统治）。\par

第5节：\Quote{藉着祂，我们领受了恩宠和宗徒职务，为使万民信服。}\BibleRef{罗 1:5}\par

请看这位仆人的坦率。他不愿有任何事物归于自己，而全归于他的主人。实际上，正是圣神赐予了这事。因此祂说：\Quote{我还有许多事要告诉你们，但你们现在不能承受。当那一位真理之神来到时，祂要引导你们进入一切真理}\BibleRef{若 16:12}；又说：\Quote{你们给我选出保禄和巴尔纳伯来}\BibleRef{宗 13:2}。在致格林多人书里，他说：\Quote{这人从圣神领受了智慧的言辞，那人从同一圣神领受了知识的言辞}\BibleRef{格前 12:8}；并且圣神按自己的意愿分施一切。在向米肋托人致辞时，他说：\Quote{圣神立你们为监督，牧养天主的教会}\BibleRef{宗 20:28}。你看，他把圣神的事称为圣子的，把圣子的事称为圣神的。\Quote{恩宠和宗徒职务}：也就是说，并非我们靠自身成就了什么而成为宗徒。因为我们并非藉着许多劳苦和努力才得到这尊荣，而是我们领受了恩宠，成功的结果属于天上的恩赐。\Quote{为使万民信服}：所以不是宗徒们成就了这事，而是恩宠为他们铺平了道路。因为他们的职责是四处宣讲，但说服人是天主在他们内运行的工作。正如路加所说：\Quote{上主开明了她的心}\BibleRef{宗 16:14}；又说：\Quote{他们得以听见天主的话}。\Quote{信服}，他说，不是要质问和展示\OriginalTerm{}{\GreekText{κατασκευὴν}}论证，而是要\Quote{信服}。因为我们的使命不是去辩论，而是把我们受托付的事物交给他人。因为当主人宣布某事时，听众不应吹毛求疵、好奇地处理所告诉之事，而只应接受；因为宗徒们正是为此被派遣，去讲说他们所听见的，而不是添加任何自己的东西，而我们则应当相信——相信什么呢？——\Quote{因祂的名}。不是要好奇于本质，而是要相信这名字；因为正是这名字行出了奇迹。因为经上说：\Quote{因耶稣基督之名，起来行走罢}\BibleRef{宗 3:6}。这也需要信德，人无法凭推理\OriginalTerm{}{\GreekText{λογισμῷ} \GreekText{καταλαβεῖν}}把握其中任何事物。\Quote{在万民中，其中你们也是蒙耶稣基督所召的人。}什么？保禄那时已向所有民族宣讲了吗？从他写给罗马人的信，可知他跑遍了从耶路撒冷直到依里黎苛的全部空间，然后又从那里出发直到地极；但即使他没有到达所有人，他的话仍不虚假，因为他不仅指自己，也指十二宗徒和所有后来传扬这道的人。从另一种意义上看，如果指他自己，当我们想到他热忱的心志，以及他在死后仍不停地在世界各地宣讲，这话也无可指摘。请注意他如何颂扬这恩赐，表明它伟大且远超过前者，因为旧事只限于一个民族，而这恩赐却吸引了海洋和陆地。也请注意，保禄的心多么远离一切谄媚；因为他与当时仿佛处于世界顶峰的罗马人交谈时，并不把他们置于其他民族之上，也不因他们当时掌权统治，就说他们在属灵事物上也有任何优势。而如同（他意指）我们向万民宣讲，同样我们也向你们宣讲，把他们与西古提人和色雷斯人同列：因为他若不想表明这一点，说\Quote{其中你们也是}便是多余的。他这样做是为了抑制他们的高傲精神\OriginalTerm{}{\GreekText{κενῶν} \GreekText{τὸ} \GreekText{φύσημα}}，压伏他们心头的虚荣膨胀，并教导他们以同样的方式尊重他人：于是接着便论述这一点。\par

第6节：\Quote{其中你们也是蒙耶稣基督所召的人。}\BibleRef{罗 1:6}\par

也就是说，你们也一同在其中；他并没有说，祂与你们一同召叫了别人，而是说祂与别人一同召叫了你们。因为如果在基督耶稣内既没有为奴的也没有自主的，更何况君王与平民呢？因为连你们也是蒙召的，并不是凭自己而来的。\par

第七节。\BibleQuote{给所有在罗马、为天主所爱、蒙召为圣人的人：愿恩宠与平安由天主我们的父和主耶稣基督赐予你们。}见：\BibleRef{罗 1:7}\par

请看，他是多么不停地说\Quote{蒙召}，他说：\Quote{蒙召为宗徒；你们也是其中蒙召的；给所有在罗马的、蒙召的人。}他这样做并非词藻堆砌，而是为了提醒他们所得的恩惠。因为在那些信友中，很可能有一些是执政官（\OriginalTerm{执政官}{\GreekText{ὑ} \GreekText{πάτων}}；本。\Emph{执政官}）和统治者，也有穷人和普通人，他摒弃了等级的差异，用一个称谓写给所有人。如果在更必要、更属灵的事上，一切都为奴隶和自由人共同拥有，例如：来自天主的爱、蒙召、福音、义子名分、恩宠、平安、圣化，以及其他一切，那么，天主所结合、并在更大的事上同享尊荣的人，若因世间的事而分裂，岂非极端愚蠢？因此，我想，这位有福的宗徒从一开始就驱逐了这种有害的疾病，引导他们走向万福之母——谦卑。这使仆人变得更好，因为他们明白，只要拥有真正的自由，奴役本身不会伤害他们；这也使主人变得温和，因为他们受教导：除非信德之善居首位，自由本身并无优势。为使你明白他这样做并不是要搅乱一切、混淆一切，而是仍然知道最好的区别，他写信的对象并非简单地是\Quote{所有在罗马的人}，而是加上了一个界定：\Quote{为天主所爱。}因为这正是最好的区分，并指出了圣化的根源。那么，圣化从何而来？来自爱。因为说了\Quote{所爱}之后，他接着就说\Quote{蒙召为圣人}，表明万福之泉源于此。他称所有信友为圣人。\Quote{愿恩宠与平安赐予你们。}\par

啊，这为我们带来无量祝福的称呼！这也是基督曾命令宗徒们进入房屋时首先使用的话。\BibleRef{路 10:5}因此，保禄也在所有地方以此为起点，从恩宠与平安开始；因为基督平息的不是一场小战争，而是一场确实变化多端、形形色色、历时久远的战争（\OriginalTerm{形形色色}{\GreekText{τοικίλον} \GreekText{καὶ} \GreekText{ταντοδαπὸν}}）；这并非由于我们的劳苦，而是凭借祂的恩宠。既然爱为我们带来了恩宠，恩宠带来了平安，他就按照问候的合宜顺序写下这些，并祈求它们能永久不变地存留，以免另一场战争再度燃起，他恳求那位赐予者稳固地保守这些事物，说道：\Quote{愿恩宠与平安由天主我们的父和主耶稣基督赐予你们。}请注意，在此处\Quote{由}对圣子和圣父是共同的，相当于\Quote{出自于。}因为他并没有说：愿恩宠与平安由天主父\Quote{藉着}我们的主耶稣基督赐予你们；而是说：\Quote{由天主父和主耶稣基督。}多么奇妙！天主的爱何其伟大！我们本是仇敌和羞辱之人，竟一下子成了圣人和儿子。因为当祂称祂为父时，就表明他们是儿子；而当他说儿子时，他就揭示了所有恩惠的宝藏。\par

那麼，讓我們不斷展示與這恩賜相稱的品行，持守於平安與聖德。因為其他的尊榮只是暫時的，隨此現世生命而終結，且可用金錢購得（由此可說它們根本不是尊榮，僅是尊榮之名，其力量在於華服與僕役的奉承），但這恩賜既由\OriginalTerm{天主}{God}所賜，即聖化與義子的恩賜，甚至死亡也不能將其打破，反而在此世使人顯耀，並且伴隨我們走向來世的生命。因為那持守於義子身份，並精確守護其聖德的人，比那頭戴冠冕、身披紫袍的人更為光輝和幸福；他在現世享有豐盛平安的喜樂，被美好的希望所滋養，沒有憂慮和困擾的緣由，而是享受恆常的喜樂；因為論到喜樂與歡愉，並非權勢的偉大、財富的豐裕、權威的炫耀、身體的強健、筵席的奢華、衣飾的點綴，或任何其他人力可及的事物所通常產生的，而是屬靈的成功和一顆純潔的良心所獨自產生的。那擁有這純潔良心的人，即使衣衫襤褸、忍受饑荒，也比那些生活安逸的人更為喜樂。同樣，那自知惡行的人，即使聚斂天下財物，也是最悲慘的人。為此緣故，\Person{保祿}雖處於持續的饑餓與赤身，每日受鞭打，卻仍然喜樂，比當時的皇帝們更加安詳。但\Person{阿哈布}儘管是國王，縱情奢華，在犯了那一次罪之後，便呻吟哀嘆、垂頭喪氣，無論犯罪前後都面帶愁容。因此，如果我們希望享受喜樂，就當首先避開邪惡，追隨美德；因為除非如此，否則即使登上王位，也無法分享喜樂。為此，\Person{保祿}也說：\BibleQuote{聖神的果實是愛、喜樂、平安。}\BibleRef{迦 5:22}那麼，讓我們使這果實在我們內不斷成長，好使我們在此世享受喜樂，並獲得將來的天國，藉著\OriginalTerm{我們的主耶穌基督}{our Lord Jesus Christ}的恩寵和對人的愛，藉著祂並與祂一起，光榮歸於\OriginalTerm{聖父}{Father}及\OriginalTerm{聖神}{Holy Spirit}，從現在直到永遠，萬世無窮。阿們。\par

\SourceRecord{Translated by J. Walker, J. Sheppard and H. Browne, and revised by George B. Stevens.；Nicene and Post-Nicene Fathers, First Series；Vol. 11.；Edited by Philip Schaff.；Buffalo, NY: Christian Literature Publishing Co.,；1889.；<http://www.newadvent.org/fathers/210201.htm>.}

% structure: p0001 source=h1 target=section reason=page-title

\section{罗马书第二篇讲道}

\BibleRef{罗 1:8}\par

\BibleQuote{第一，我藉着耶稣基督为你们众人感谢我的天主，因你们的信德传遍了全世界。}\par

这样的开场白与这位有福之灵相称，足以教导众人将善行和言语的初果奉献给天主，不仅为自己的善行，也为别人的善行而感恩；这亦使灵魂远离嫉妒与怨恨，并使天主更接近那些如此感恩者的爱心。因此他在别处也说：\BibleQuote{愿我们的主耶稣基督的天主和父受赞美，祂以各种属神的祝福祝福了我们。}\BibleRef{弗 1:3} 理当感恩，不仅富足时，贫穷时亦然，不仅健康时，患病时亦然，不仅顺利时，逆境时亦然。因为当事情顺风顺水时感恩并不希奇；但当大风暴袭来，船只摇摆颠簸、危在旦夕时，那才是显示忍耐与心善的大好时机。为此，\Person{约伯}也因此获得冠冕，堵住了魔鬼的无耻之口，清楚地表明：即使在他过好日子时，感恩并非出于财富，而是出于他对天主的深厚爱慕。再看他为何事感恩：并非为属地易朽之物，如权势、权威和荣耀（这些算不了什么），而是为真正的福分——信德与宣讲的胆量。他以多么深的情感感恩啊，因为他不说\Quote{给天主}，而说\Quote{给我的天主}，先知们也是如此，把众人公有之物变为自己的。先知们这样行有什么奇特呢？因为天主自己也常这样对祂的仆人，称自己是亚巴郎的天主、依撒格的天主、雅各伯的天主，好似这是他们特有的。\BibleQuote{你们的信德传遍了全世界。}那么，全世界都听说了罗马人的信德吗？是的，据他说，全世界都听说了（\OriginalTerm{自从那时起}{\GreekText{πἅσα} \GreekText{ἐξ} \GreekText{ἐκείνου}}）。这并非不可能，因为那城市并非无名小城，而是位于某种高处，各方面都引人注目。但请思考这宣讲的力量：在短时间内，它藉着税吏和渔夫之手，竟然占据了所有城邦之首，而叙利亚人成了罗马人的导师和领路人。于是他为他们作证了两项卓越之处：他们不仅相信了，而且相信得有胆量，如此之大以致他们的名声传遍天下。他说：\BibleQuote{因为你们的信德，}他说，\BibleQuote{传遍了全世界。}他说，你们的信德，而非你们的言语辩论、质疑或三段论。然而那里有许多阻碍教训的障碍。因为他们刚获得世界霸权，趾高气扬，生活奢华，渔夫们把宣讲带到那里，而他们又是犹太人，出自犹太人，一个被万民憎恶的民族；他们被要求崇拜那位被钉十字架者，祂在犹太长大。同时，教师们在教义之外还向那些惯于安逸、为现世事物焦虑的人宣讲一种严苛的生活。而那些宣讲者是贫穷平凡的、出身低微、生于低微之家。但这一切都没有阻碍福音的进程。被钉十字架者的力量如此巨大，以致把福音传遍四方。他说：\BibleQuote{因为它是被传说的，}他说，\BibleQuote{在全世界。}他没有说\Quote{被彰显}，而是说\Quote{被传说}，仿佛所有人都把这放在嘴边。事实上，当他在得撒洛尼人那里为此作证时，他还加上了另一件事。因为他说了\BibleQuote{天主的圣言由你们发出后，}他又说，\BibleQuote{以致我们不需要说什么。}\BibleRef{得前 1:8} 因为门徒们取代了教师的位置，以他们宣讲的勇气教导众人，吸引众人归向他们。因为宣讲从未在任何地方停滞，而是比火更快地传遍全球。但在此处只有这么多——\BibleQuote{被传说}。他正确地说\BibleQuote{被传说}，表明无需对所说的话增添或删减。因为信使的任务就是只把所告诉他的传达给别人。为此，司铎也被称为\BibleQuote{使者}\BibleRef{拉 2:7}，因为他不说自己的话，而说那派遣他者的话。然而\Person{伯多禄}曾在那里宣讲过。但他把别人的看作自己的，正如我之前所说，他如此超越度量地清除了所有嫉妒之心！\par

第9节：\BibleQuote{因为天主是我的见证，我以心神在祂儿子的福音中事奉祂。}\BibleRef{罗 1:9}\par

\Person{保禄}的这些话语是宗徒慈爱肺腑的话语，这是父亲关怀的彰显！他说的是什么？为何呼求天主作证？他必须向他们表达自己的情感。既然他尚未见过他们，便不呼求任何人作证，而是呼求那洞察人心的一位。因为他说了\Quote{我爱你们}，并以此作为证据，提到自己不断为他们祈祷，且渴望前来看望他们；而这一点并非显而易见，他便转向可靠的见证。那么，你们中有人能夸口说，当在家（\OriginalTerm{在家}{\GreekText{ἐπὶ} \GreekText{τἥς} \GreekText{οἰκίας}}）祈祷时，他记念整个教会的团体吗？我想没有。但\Person{保禄}不仅代表一座城，而是代表全世界亲近天主，且不是一次、两次或三次，而是不停地。但如果不断地将任何人记在心中，不能没有极大的爱；那么，在祈祷中有他们，且不停地有他们，试想这表示何等大的情爱与友谊。但当他引用\Quote{我以心神在他圣子的福音中所事奉的那位}时，他同时向我们展示了天主的恩宠和他自己的谦卑之心：天主的恩宠是因为祂将如此重大的事交托给他；他自己的谦卑是因为他将这一切不归功于自己的热忱，而归于圣神的援助。而加上\Quote{福音}一词，表明了服务的方式。因为有许多种不同的服务方式。正如在君王之下，所有人都隶属于一位掌王权者，并非所有人都服事（\OriginalTerm{服事}{\GreekText{διακονοὕνται}}）同一件事，而是有人负有统帅军队的职务，有人治理城市，还有人保管仓库中的财宝；同样，在属神的事上，一人事奉天主并劳力（\OriginalTerm{敬拜与服务}{\GreekText{λατρεύει} \GreekText{καὶδουλεύει}}）于在信德中并妥善安排自己的生活，另一人承担照顾旅客，另一人着手保护有需要的人。正如在宗徒时代，\Person{斯德望}的同道们在照料寡妇的事上事奉天主，其他人（\OriginalTerm{其他人}{\GreekText{ἄλλοι}}，两份抄本均为\OriginalTerm{异文}{\GreekText{ὡν}}）在教导圣言上事奉；\Person{保禄}也属于后者，在传福音上服务。这便是他服务的样式：因为他正是为此被委任。因此，他不仅呼求天主作证，而且还说出自己所受托的事，为表明既然将如此大事交托在他手中，他不会呼求那托付给他的那一位为虚假之事作证。此外，他还想指出另一件事，即他不得不对这城市怀有爱与关怀。为了不让他们说\Quote{你是谁？从哪里来？你说你为如此伟大且最帝制的城市挂虑}，他表明自己必须有此关怀，至少就他所受托的服务是宣告福音而言：因为凡受托此事者，必须不断挂念那些将要领受圣言的人。他还通过说\Quote{以我的心神}来表明另一件事：这服务远高于外邦人或犹太人的服务。因为外邦人的服务既属肉体且处于谬误，犹太人的服务诚然真实，但仍是属肉体的。但教会的服务正与外邦人相反，且远高于犹太人的服务。因为我们服务的方式不是用羊、牛、烟和脂肪，而是借着一个属神的灵魂。这一点基督也表明了，祂说\Quote{天主是神，凡朝拜祂的，应当以心神以真理朝拜祂}。\BibleRef{若 4:24}\par

\Quote{在他圣子的福音中}。上面他说是圣父的福音，这里他说是圣子的福音。说圣父的或圣子的，竟如此无关紧要！因为他从那蒙福的声音学得，圣父的一切是圣子的，圣子的一切是圣父的。因为\Quote{我的一切都是你的，你的一切都是我的}。\BibleRef{若 17:10}\par

\Quote{我不断在祈祷中常常提到你们}。这是真诚之爱的表现。他看似只说一件事，但在此却陈述了四件事：他记念他们，他不断地记念，是在祈祷中记念，并且是为他们祈求大事。\par

第10、11节。\Quote{恳求，或许现在终能因天主的旨意，顺利来到你们那里。}\par

你见他痛苦地渴望见他们，然而若不合天主看为善的事，他就不忍见他们，却将他的渴望与对天主的敬畏相融合。因为他爱他们，并急切地想去他们那里。然而他没有因爱他们而渴望见他们，若那不合天主看为善的事。这才是真爱，不像我们这些在爱之律法两方都犯错的人：我们或是不爱任何人，或是若我们曾爱，我们便爱得与天主看为善的事相反，在两方面都违犯天主的律法。若这些话在谈论时就显得沉重（\OriginalTerm{沉重的}{\GreekText{φορτικά}}），那么行出来就更沉重了。而我们是怎样相反天主看为善的事而爱呢？（你将说）当我们忽略了因饥饿而消瘦的基督，却为我们的子女、朋友和亲戚提供超出他们所需之物。或更不必说要将这话题推进了。因为若有人愿省察自己的良心，他会发现这在许多事上发生。但那位有福之人并非如此，他却知道如何去爱，并适当地去爱\EditorialNote{三部抄本省略\InnerQuote{as he ought}}，且合宜地爱，尽管他在爱上超越众人，却没有逾越爱的尺度。请看在他内两样事极其兴盛：对天主的敬畏，以及对罗马人的渴望。因为不断祈祷，且未得允时不停止，显出极度的爱。然而在爱时，仍继续顺服天主的旨意，显出强烈的敬畏。但在另一处，他曾\Quote{三次祈求主}\BibleRef{格后 12:8}，他不仅没有得到，反而当未得时，他因未被垂听而非常感恩。所以，他在一切事上都注目于天主。但在这里他得到了，尽管不是在他求的时候，而是在延迟之后，而他也未因此不满。我提这些事，是要我们不要因未蒙垂听或蒙垂听缓慢而抱怨。因为我们并不比保禄好，他承认对两者都感恩，且有理。因为当他一旦将自己交付给那统御万有的手，并将自己如泥土对陶匠般顺服于它，他便随天主所引。在说了他渴望见他们之后，他提起了他渴望的原因；那是什么呢？\par

第11节。\BibleRef{罗 1:11} \Quote{为要将一些属灵的恩赐分给你们，使你们得以坚固。}\par

因为他并非如现今许多人无益无目的地旅行那样前去，而是出于必要且极其紧迫的目的。他没有公开告诉他们他的意思，而是通过暗示，因为他并没有说\Quote{我可以教导你们}、\Quote{我可以教训你们}、\Quote{我可以补足那缺乏的}；而是说\Quote{我可以分给}；表明他给他们的不是他自己的东西，而是将自己所受的分施给他们。在这里他再次谦逊，说\Quote{一些}，意思是一点，且与我的能力相称。而这将分施的一点是什么呢？他说，就是\Quote{为要使你们得以坚固}。这同样来自恩宠，即不摇动并站稳。但当你听到恩宠时，不要以为我们方面决心的报酬便被摒弃了；因为他讲论恩宠，并非要贬低我们方面决心的辛劳，而是要抑制（\OriginalTerm{抑制}{\GreekText{ὑποτεμνόμενος}}，如刺破膨胀之物）狂妄精神的骄傲（\OriginalTerm{狂妄}{\GreekText{ἀπονοίας}}）。所以，不要因为保禄称此为恩宠的恩赐，就变得懈怠。因为他知道，在他极大的坦诚中，他甚至称善行为恩宠；因为即使在善行中，我们也需要许多从上而来的影响。但他说\Quote{为要使你们得以坚固}，暗含他们需要许多纠正：因为他想说：我早已渴望并祈祷见你们，\Quote{长久以来我既渴望}见你们，并无别因，只为我可以彻底地\Quote{坚固、加强、固定}你们在天主圣言中，使你们不再不断摇摆。但他没有如此直说（因为这会震惊他们），而是以另一种方式向他们暗示同一件事，尽管语气缓和。因为当他说\Quote{为要使你们得以坚固}，他就表明了这一点。既然这也令人不悦，请看他是如何通过下文柔和了它。为避免他们说：我们岂是摇摆不定的，飘来飘去？岂需要你发言才能站稳？他预见到并消除了这类反驳，这样说：\par

第12节。\BibleRef{罗 1:12} \Quote{这就是说，我在你们中间，因你我彼此的信德，同得安慰。}\par

彷彿他曾說：不要懷疑我說這話是為了指責你們。我並非懷著這種心情說我所說的。但我想說的是什麼呢？你們正經歷許多磨難，從四面八方被淹沒（由於那些迫害你們的人\OriginalTerm{被淹没}{\GreekText{περιαντλούμενοι}}。三個手抄本作\OriginalTerm{被骚扰}{\GreekText{παρενοχλούμενοι}}，被騷擾）。那時我渴望見到你們，好能安慰你們；或者更確切地說，不僅是安慰你們，也是為了我自己得到安慰。請看這位師傅的智慧。他說：\Quote{使你們堅固}；他知道他所說的會使門徒感到沉重和不悅。他又說：\Quote{使你們得到安慰}。但這也同樣沉重，雖然不如前者，但仍然沉重。於是他進一步削減其中刺痛的部分，從各方面使他的話語平滑，易於接受。因為他並不只是說\Quote{得到安慰}，而是說\Quote{與你們一同得到安慰}；他並不滿足於此，又加上一個緩和，說：\Quote{藉著你我彼此的信德}。哦，他的謙卑何其偉大！他表明自己也需他們，而不只是他們需要他。他將門徒置於師傅的地位，不讓任何優越感留在自己一方，而是指出他們的完全平等。因為他說，這益處是相互的，我需要你們的安慰，你們也需要我的安慰。這如何可能呢？\Quote{藉著你我彼此的信德}。如同火，若有人聚集許多燈火，便燃起明亮火焰；同樣，在信友中自然也是如此。因為當我們獨處，與他人分離時，我們總是有點沮喪。但當我們彼此相見，與我們自身的肢體交織時，我們便得到極大的安慰。你們不要只看現今的時代——因天主的恩寵，無論在城市還是在曠野中，都有許多信友的群體，一切不敬神之事已被驅逐——但要思想在那個時代，門徒見到師傅，以及從另一城市來的弟兄被弟兄看見，是何等大的益處。為使我說的話更清楚，讓我舉例說明。假如竟發生了這樣的事（但願永不如此！），我們被擄到波斯人或西徐亞人或其他蠻族的土地，並被分散（七個手抄本作\Quote{撕裂}）兩三人在他們的城中，然後突然看到任何一位從這裡來的人來到我們中間，請想我們會收穫何等大的安慰！你們不見那些在監獄中的人，若見到任何熟人，他們會如何復甦，因快樂而雀躍不已嗎？但我將那些日子比作被擄和監禁，請不要奇怪。因為那些人所受的痛苦遠比這些人更為艱難：他們四散流離，被驅逐，生活在饑荒和戰爭中，每天戰戰兢兢地等待死亡，懷疑朋友、親屬和戚族，活在世上如處異地，甚至比生活在異鄉的人處境更為艱苦。為此他說：\Quote{使你們與我們藉著彼此的信德，得以堅固並一同得到安慰。}他說這話，並非因為他自己需要他們的任何幫助（絕非如此；因為教會的柱石，比鐵和岩石更堅固，屬靈的金剛石，足以擔當無數城市的責任），而是為了使他的言語不過於激烈，責備不過於嚴厲，才說他自己也需要他們的安慰。但若有人說，這安慰是他因他們信德的增長而喜樂，而且 \Person{保祿} 需要這喜樂，那麼這樣理解他的意思也並不錯。那麼，若人說：你既然渴望、祈禱，並要從中獲得安慰和給予安慰，又有什麼妨礙你來呢？為了消除這種懷疑，他就繼續說。\par

第十三節。\Quote{\BibleQuote{弟兄們，我不願意你們不知道，我屢次定意往你們那裡去，只是到如今仍有阻隔。}\BibleRef{罗 1:13}}\par

這裡有奴僕般的順從，以及他優秀品格\OriginalTerm{優秀的品格}{\GreekText{εὐγνωμοσύνης}}的明顯表現！因為他說有阻隔，但卻沒有繼續說明原因。他並不追究主人的命令，只是服從。然而，有人可能會提出問題：為什麼天主阻止一個如此顯著而偉大的城市——全世界都注視著它——享受這樣的師傅如此之久呢？因為那征服了首府城市的人，本可以輕易前往其屬城。但那忽略更為王者的城市，而伺機於附屬地的人，卻忽略了要點。然而，他卻不為這些事煩惱，而是順服於天主眷顧的不可測度，藉此既顯出他靈魂的正直，又教導我們眾人絕不要要求天主為所發生的事負責，即使所發生的事似乎困擾許多人的心思。因為主人只有發令的份，僕人則要服從。這就是為什麼他說有阻隔，卻不說原因；意思是：連我也不知道，那麼就不要問我天主的計劃或心意。因為\Quote{\BibleQuote{受造之物豈能對造它的說：你為什麼這樣造我呢？}}告訴我，你為什麼甚至想要知道它？難道你不知道萬事都在祂的照顧下，祂是智慧的，祂不憑偶然行事，祂愛你超過生養你的人，祂對你的慈愛遠遠超過父親的疼愛和母親的憂慮嗎？那麼就不要再尋求，也不要再進一步；因為這對你已是足夠的安慰：因為即使在那時，對於羅馬人而言，事情也是妥善安排的。即使你不知道方式，也不要介意；因為這正是信德的主要特點：即使不明白安排的方式，也要接受那告訴我們關於祂眷顧的事。\par

保禄既然在他所热切关注的事上获得了成功（什么事呢？就是表明他之所以没有到他们那里去，并非由于轻视他们，而是因为他虽然极度渴望，却被阻止了），并且消除了自己懈怠的指责，又说服他们相信他渴望见到他们的心情并不亚于他们自己，于是便通过其他事进一步显示了他对他们的爱。他说：\Quote{即使我被阻止，} 我也没有放弃尝试，而是不断尝试，却总是被阻止，但我从未因此放弃，不与天主旨意相违，仍然保持我的爱。因为他自己定意并不放弃，表明了他的情感；而由于被阻止却没有与之抗争，表明了他对天主的全部爱。\BibleQuote{为叫我在你们中也能获得一些果实。} 他之前已经告诉了他们他渴望的原因，并表明那是相称的；但在这里他仍然陈述它，消除他们的一切猜疑。因为既然那座城市引人注目，在整个海陆范围内无与伦比，以至于许多人仅因渴望认识它就成了前往那里的理由（\OriginalTerm{理由}{\GreekText{πρόφασις}}），免得他们对保禄产生任何类似的想法，或怀疑他只是为了荣耀而把他们归为己有才渴望去那里，他便反复陈明他渴望的缘由，之前他说：\BibleQuote{为将一些神恩分施给你们，} 我渴望见到你们；但这里说得更清楚：\BibleQuote{为叫我在你们中也能获得一些果实，如同在其他外邦人中一样。} 他将统治者与臣民并列，在无数次凯旋、胜利和执政官的荣耀之后，他将他们与野蛮人并列，这也是有充分理由的。因为在哪里有信德的高贵，那里就没有野蛮人、没有希腊人、没有外邦人、没有公民，而所有人都上升到同一尊严的高度。请看他也在这里谦逊，因为他没有说：为叫我能教导和训诲，而是说什么？\BibleQuote{为叫我能获得一些果实。} 而且不是简单的果实，而是\BibleQuote{一些果实。} 再次贬低他自己在其中所分得的部分，正如他上面所说的：\BibleQuote{为将一些恩赐分施。} 然后，为了也抑制他们，正如我之前也说的，他说：\BibleQuote{如同在其他外邦人中一样。} 因为，我不因为你们是富有的、比别人优越，就较少关心其他人。因为我们寻找的不是富人，而是有信德的人。如今那些希腊的智者在何处？他们留着长须，穿着敞开的外袍，高谈阔论（\OriginalTerm{高谈阔论}{\GreekText{τὰ} \GreekText{μεγάλα} \GreekText{φυσὥντες}}）？整个希腊和所有蛮族之地都被那位制帐篷的人归化了。但那个在他们当中备受赞誉、被推崇的柏拉图，第三次带着他那些言辞的浮夸、带着他辉煌的名声（\OriginalTerm{名声}{\GreekText{ὐ} \GreekText{πσλήψεως}}）来到西西里，却连一个国王都没有说服，反而如此悲惨地离开，甚至失去了自由。但这制帐篷的人不仅跑遍了西西里或意大利，而是整个天下；并且在讲道的同时他也没有放弃他的手艺，甚至还在缝制皮革，管理作坊。甚至连这一点也没有冒犯那些出生自执政官的人，而且很有理由，因为通常使教师成为轻蔑对象的不是他们的行业和职业，而是虚假和伪造的教义。正因如此，雅典人至今仍在嘲笑前者。但这个人连野蛮人也听从，甚至愚笨无知的人也听从。因为他的宣讲平等地摆在全众人面前，它不区分等级，不讲究民族的优越，也不讲究其他任何这一类的事；因为它只要求信德，而不是推理。因此它最值得钦佩，不仅因为它有益且救灵，而且因为它容易接受、\Quote{可爱的}，并为所有人所能理解：这正是天主眷顾的一个主要目的，祂将祂的恩惠普遍地赐给众人。\par

因为祂对太阳、月亮、大地、海洋和其他事物所做的，并不是让富人和智者享受这些恩惠的更大一份，而让穷人享受较少，而是让众人同等地享受它们，祂对宣讲也这样做了，甚至更甚，因为这宣讲比它们更为不可或缺。因此保禄多次说：\Quote{在所有的外邦人，} 以表明他丝毫不偏袒他们，而是在履行他主的命令，并引导他们归向对万有之天主的感恩。他说：\par

第 14 节。\BibleQuote{我对希腊人和野蛮人、对智者和愚者都是一个欠债者。} \BibleRef{罗 1:14}\par

这也是他在写给格林多人的书信中所说的。他说这话是为了将一切归于天主。\BibleRef{格前 9:16}\par

第 15 节。\BibleQuote{所以，尽我所能，我也愿意向你们在罗马的人宣讲福音。} \BibleRef{罗 1:15}\par

哦，高贵的灵魂！ 他承担了一项充满如此巨大危险的任务——远渡重洋，遭遇诱惑、阴谋、叛乱（因为一个向如此伟大的城市发言、而该城市正受不虔敬的暴政统治的人，很可能要承受多如雪花的诱惑；也正是用这种方式他在这座城市丧了命，被该城的暴君所杀）——然而，尽管预料要承受如此巨大的艰难，他却没有因此变得懈怠，反而心急如焚、热切等候、决心已定。因此他说：\Quote{所以，尽我所能，我也愿意向你们在罗马的人宣讲福音。}\par

第 16 节。\BibleQuote{因为我并不以福音为耻。} \BibleRef{罗 1:16}\par

\Quote{保禄，你说什么？当可以说：我夸耀、骄傲、以此自得时，你却不这么说，反而说了比这更轻的话：你\InnerQuote{不以它为耻}，这通常不是我们用来形容极荣耀之事的话语。} 那么他说的这话是什么意思？他为何这样说呢？然而他为此所感到的喜乐，甚至超过对天堂的喜乐。至少，在写给迦拉达人的信里，他说：\BibleQuote{但愿我绝不夸耀，除非夸耀我们的主耶稣基督的十字架。} \BibleRef{迦 6:14} 那么他在这里怎么既不说他夸耀，反而说\Quote{我不以它为耻}呢？罗马人因他们的财富、帝国和胜利，对世俗之事极其热衷，他们甚至把他们的君王视为神明，并如此称呼他们。为此，他们也用圣殿、祭台和祭献来崇拜他们。既然他们如此骄傲，而保禄将要宣讲的耶稣，被认为是个木匠的儿子，在犹太长大，且在一位卑微妇人的家中，他没有卫兵，没有环绕的财富，甚至像罪犯一样与强盗同死，还忍受了许多其他可耻的事；而他们很可能因为尚未知道那些不可言传的伟大之事而羞于承认。所以他说：\Quote{我不以它为耻}，仍然要教导他们不要羞愧。因为他知道，如果他们在这方面成功，就会迅速前进，达到夸耀的地步。所以，如果你听到有人说：你敬拜被钉十字架的那位吗？不要羞愧，不要低头，反要以此自得，面带荣光，以自由人的眼神和昂首的姿态，宣认你的信仰。如果他再问：你敬拜被钉十字架的那位吗？你要回答他：是的！而且不是那奸夫，不是那辱骂父亲的人，不是那杀害自己儿女的人（因为他们所拜的神都是这样的），而是那通过十字架堵住魔鬼之口、消除他们无数诡计的那位。因为十字架是为了我们，是不可言喻的爱人之情的工程，是祂对我们极大关怀的标志。此外，由于他们因浮夸的言辞和外在智慧的伪装而自高自大，我（意思是要说）完全告别这些推理，前来宣讲十字架，并不因此感到羞愧：\BibleQuote{因为它是天主的德能，为救恩。} \BibleRef{罗 1:16} 因为也有天主的德能用于惩罚（因为当祂惩罚埃及人时，祂说：\BibleQuote{这是我的大能} \BibleRef{岳 2:25}），也有德能用于毁灭（因为祂说：\BibleQuote{你们要害怕那能在地狱里毁灭身体和灵魂的那位} \BibleRef{玛 10:28}），为此他说：我来带来的不是这些，不是惩罚和刑罚的德能，而是救恩的德能。那么怎样？难道福音没有告诉我们这些事吗？即关于地狱、外面的黑暗和毒虫的叙述？然而我们除了从福音，别无来源知道这些。那么他在什么意义上说：\BibleQuote{天主的德能，为救恩}？只看下文：\BibleQuote{对一切信的人，先是犹太人，后是希腊人。} \BibleRef{罗 1:16}\par

因为他不是对所有人，而是对接受它的人。即使你是一个希腊人（即外邦人），甚至是一个陷入各种恶习的人，即使是西古提人，即使是野蛮人，即使是一头畜生，充满一切非理性，背负着无数罪恶的重担，只消你接受了关于十字架的话，并领受了圣洗，你就将这一切都涂抹了。他为什么在这里说：\BibleQuote{先是犹太人，后是希腊人}？这区别是什么意思？然而他常说：\BibleQuote{割损不算什么，不割损也不算什么} \BibleRef{迦 5:6}；那么他为什么在这里区分，把犹太人放在希腊人之前？这是为什么呢？因为虽然先到的人并没有因此得到更多的恩宠（因为同等的恩赐赐给了这人和那人），但\Quote{先}只是在时间顺序上的荣誉。因为他并没有像得到更大义德那样的优势，只是在首先接受上得到荣誉。因为在那些受光照的人中（你们这些受过洗的知道是什么意思），所有人都跑向圣洗，但并非都在同一时刻，而是一个先，另一个后。然而先到的人并不比后到的人得到更多，也不比之后的人多，所有人都享受同等的恩赐。这里的\Quote{先}只是在言语上的荣誉，而不是在恩宠上的优越。然后说了\BibleQuote{为救恩}之后，他进一步彰显这恩赐，表明它不止于目前，而是继续前进。因为这就是他接下来所阐述的：\par

节17. \BibleQuote{因为天主的义德正在这福音上显示出来。} \BibleRef{罗 1:17}\par

但成了正义的人必要生活，不只为现世的生命，也为那未来的生命。他不仅暗示这个，还暗示另一件事，即这种生命的光荣与荣耀。因为，既然有可能得救，却不免羞辱（如许多因君主的仁慈而免于刑罚的人得救那样），为免人一听到安全就猜疑这一点，他就加上了正义；而且不是你们自己的正义，而是天主的正义；同时也暗示其丰富和容易。因为你们不是靠劳苦和努力获得它，而是藉从上而来的恩赐领受它，只从你们自己的库存中贡献一件事，就是\Quote{相信}。然后，既然他的陈述似乎不可信，如果奸淫者、柔弱者、盗墓者、行邪术者不仅要突然免于刑罚，而且要成为正义的，并且是以最高的正义成为正义的；他就从旧约证实他的断言。首先用一句简短的话，他为有见识的人打开了一片历史的广阔海洋。因为说了\Quote{由信德到信德}之后，他将听者引回到天主在旧约中如此实行的安排，当给希伯来人写信时，他用他惯常的大智慧解释了这一点，表明甚至那时义人和罪人都以那种方式成义，因此他也提到了妓女和亚巴郎。但在这里，仅仅暗示之后（因为他正赶往另一个紧迫的主题），他又从先知们证实他所说的，将哈巴谷带到他们面前，大声呼喊说，要生活的人除非藉信德，否则无法生活；因为他说：\BibleQuote{义人必因信德而生活} \BibleRef{哈 2:4}，讲的是未来的生命。因为既然天主所赐完全超越理性，我们有理由需要信德。但那轻视它、傲慢虚荣的人，将一事无成。愿异端者听从圣神的声音，因为推理的本性就是如此。它们像某种迷宫或谜题，无处不在尽头，不让人把理性立在磐石上，其根源在于虚荣。因为他们羞于接受信德，并显得对天上的事无知，所以将自己卷入无数推理的尘埃云中。那么，啊，可怜痛苦的人，配得无尽的眼泪，若有人问你天是如何造的、地是如何造的——我为何说天和地？你自己是如何生的，如何滋养的，如何长大的——你难道不以无知为耻吗？但若谈到独生子，你却因羞耻将自己推入毁灭的坑中，认为不知道一切有损你的尊严？然而好争论是可耻的，不合时宜的好奇心也是如此。我为何谈教义？甚至我们摆脱今生的败坏，也无非藉信德。同样，古时众人如此发光，亚巴郎如此，依撒格如此，雅各伯如此，妓女也得救，旧约中的一位，新约中的一位也是如此。因为他说：\BibleQuote{因信德，妓女辣哈布没有与那些不相信的人一同灭亡，因为她接待了侦探。} \BibleRef{希 11:31} 因为她若对自己说：\InnerQuote{那些被掳、流亡、逃难、过游牧生活的人，怎能胜过我们有城、有墙、有塔的我们？}她就会毁灭自己和她们。当时得救之人的祖先也遭受了这痛苦。因为当看见高大的人时，他们追问胜利的方式，就未经战斗和阵列全部灭亡。你看到不信的坑有多深吗！信德的墙有多坚固！前者带下无数成千的人，后者不仅拯救了妓女，还使她成为如此众多民族的保护者！\par

既然我们知道这些和更多，绝不要叫天主为所做的事负责，而是无论祂加给我们什么，我们都接受，不要陷入琐碎和好奇的追问，即使对人理性而言所命令的事似乎不妥。因为，请问，有什么比父亲亲手杀死他唯一合法的儿子更不妥的呢？\BibleRef{创 22:3} 但那时当命令这位义人时，他没有提出琐碎的疑虑，而是由于命令者的尊威，他仅接受了命令。另一位受天主命令击打先知的人，当他因命令表面的不合理而提出琐碎疑虑，没有简单地服从时，受到了极重的惩罚。\BibleRef{列上 20:35} 但击打的人却得了称赞。撒乌耳也因违背天主的谕令而饶恕人，从王国坠落，并受到不可挽回的惩罚。此外还可以找到其他例子：所有这些教导我们，永不要为天主的命令要求理由，只要服从。如果对祂所吩咐的任何事提出琐碎疑虑是危险的，好奇追问者被指定受极重的惩罚，那么那些对更隐秘可怕的事好奇追问的人，比如祂如何生了圣子，以何种方式，祂的本质是什么，还有什么借口呢？既然我们知道这一点，就让我们以善意接受万福之母——信德；好让我们仿佛航行在平静的港湾，既保持教义纯正，又安全地引导生活直行，得以获得那些永恒的福乐，藉着我们的主耶稣基督的恩宠和慈爱，祂与圣父及圣神，同享光荣，直到永远。阿们。\par

\SourceRecord{Translated by J. Walker, J. Sheppard and H. Browne, and revised by George B. Stevens.；Nicene and Post-Nicene Fathers, First Series；Vol. 11.；Edited by Philip Schaff.；Buffalo, NY: Christian Literature Publishing Co.,；1889.；<http://www.newadvent.org/fathers/210202.htm>.}

% structure: p0001 source=h1 target=section reason=page-title

\section{罗马书讲道集 第三篇}

\BibleRef{罗 1:18}\par

\BibleQuote{因为天主的忿怒，从天上显示在一切不虔敬和不义的人身上，就是那些以不义压制真理的人。}\par

请观察保禄的智慧：他用温和之事鼓励后，如何转向更可畏的话题。因他说福音是救恩和生命的根源，是天主的大能，产生救恩和义德后，便提及那可能使忽略福音者惧怕的事。因为一般人多半不为美善的应许所动，反为痛苦的恐惧所驱使，他便从两方面吸引他们。为此天主不但应许天国，也恐吓地狱。先知们对犹太人说话也是如此，常将灾祸与福乐交织。为此保禄也这样变化他的言论，但并非任意，而是先置福乐，后置灾祸，以显明前者出于天主的引导旨意，后者则出于悖逆者的邪恶。先知也是这样先提好事，说：\BibleQuote{如果你们愿意服从我，就必吃地上的美物；如果你们不愿意不服从，刀剑必吞灭你们。}\BibleRef{依 1:19}保禄在此也如此行。但请看：他意谓基督来是要带来赦免、义德和生命，然而并非任意，而是借着十字架——那是至极且奇妙的事，祂不仅赐予这些，也亲自承受了这些。因此，如果你们傲慢地轻视恩赐，刑罚就会等待你们。再看他是如何提升言辞：他说：\BibleQuote{因为天主的忿怒，}他说，\BibleQuote{从天上显示出来。}这从何显明？如果说是信徒说的，我们会告诉他基督的宣告；但如果是不信者或希腊人，保禄就以他所说的关于天主的审判来使他们无言，从他们所行的事上带来无可辩驳的证据。这也是他最显著的一点：他如何使那些反对真理的人，借着他们日常言行，自己为真理的道作证。但这些留待后话；目前我们专注于眼前的话。\BibleQuote{因为天主的忿怒从天上显示出来。}的确，即使在今世，这事常借着饥荒、瘟疫和战争发生：个人和全体都受惩罚。那么新的情况是什么？刑罚将更大，且普遍临到众人，不再按同样的规律。因为现在发生的是为了管教，那时却是为了惩罚。这也正是圣保禄所显示的，他说：\BibleQuote{我们现在受管教，免得和世人一同被定罪。}\BibleRef{格前 11:32}如今对许多人来说，这些事似乎不是来自上天的忿怒，而是出于人的恶意。但那时，天主的惩罚将是显明的，当审判者坐在可畏的审判席上，命令一些人被拖到火炉，另一些到外面的黑暗，还有一些到其他无情且难忍的刑罚。为何他不直白地说天主子带着千万天使来临，并要各人交账，反倒说\BibleQuote{天主的忿怒显示出来}？因为听众尚为初学者，所以他先用他们完全认可的事吸引他们。此外，除了这里提到的，在我看来他似乎也针对希腊人。这就是为什么他从这开始，但之后便引入基督审判的主题。\par

\BibleQuote{针对一切不虔敬和不义的人，就是那些以不义压制真理的人。}这里他显示不虔敬的路径很多，而真理的路径只有一条。因为错误是多样、多形和复合的，但真理是唯一的。在谈论教义之后，他谈论生活，提到人的不义。因为不义也有多种。一种涉及金钱，如有人在这方面对邻人行不义；另一种关于妇女，如有人离开自己的妻子，侵扰他人的婚姻。因为圣保禄也称此为欺骗，说：\BibleQuote{谁也不可在这件事上越分欺压他的弟兄。}\BibleRef{得前 4:6}还有的人不损害妻子或财产，而是损害邻人的名誉，这也是不义。因为\BibleQuote{美名胜过大财。}\BibleRef{箴 22:1}但有些人认为这也是保禄谈论教义的话。即便如此，也不妨碍它同时针对两者。至于\BibleQuote{以不义压制真理}是什么意思，可从下文得知。\par

\BibleRef{罗 1:19} \BibleQuote{因为天主的事情，人所能知道的，在他们里面原是明显的，因为天主已经给他们显明了。}\par

但这份荣耀，他们却赋予了木石。正如那受托管理君王财物的人，奉命将其用于君王的荣耀，若他将这些财物浪费在强盗、娼妓和巫婆身上，并用君王的库藏使他们光彩，他便会因对王国犯下极大的恶行而受罚。同样，那些在领受了天主及其荣耀的知识后，却将其赋予偶像的人，\BibleQuote{以不义压制了真理}\BibleRef{罗 1:18}，并且，至少就他们能力所及，不义地对待了这知识，未将其用于合适对象。现在，我所说的你明白了吗？或者还需要更清楚些？或许有必要再多说一点。那么，这里所说的是什么呢？天主从起初就将认识祂的知识放在人里面。但这些人却将此知识赋予木石，从而不义地对待了真理，至少在他们力所能及的范围内如此。因为真理保持不变，拥有自己不变的荣耀。\Quote{那么，祂将这知识放在他们里面，这从何明显呢，哦，\Person{保禄}？} \BibleQuote{因为}\BibleRef{罗 1:19}，他说，\BibleQuote{那关于祂可知的事，在他们当中是显而易见的。}\BibleRef{罗 1:19}然而，这是一个断言，而非证明。但请你证明它，并向我表明天主的认识对他们而言是明显的，而他们却情愿偏离。那么，它从何明显呢？祂是否从天上向他们发出声音？绝不。但什么比声音更能吸引他们呢？祂所做的，就是将受造物摆在面前，使得无论智慧人、无知者、西古提人、还是蛮夷，都能通过眼见学习所见之物的美，从而上升到天主。因此他说：\par

第20节：\BibleQuote{因为从创世以来，天主的不可见之事，藉着所造之物，就被清楚看见了，从而被理解。}\BibleRef{罗 1:20}\par

先知也曾说：\BibleQuote{诸天宣示天主的荣耀。}\BibleRef{咏 18:1} 因为那些希腊人（即外邦人）在那一天将说什么呢？\Quote{我们不知道你吗？} 难道你们没有听见天藉着景象发出声音，而万有并然有序的和谐比号角更响亮地说话吗？难道你们没有看见昼夜时刻永恒不动地持续，冬、春和其他季节的美好秩序既稳固又不动摇，以及海在狂涛巨浪中的\OriginalTerm{明理}{\GreekText{εὐγνωμοσύνην}}吗？万有按秩序存在，并以其美和宏大高声宣讲造物主。因为保禄在总结这一切和更多事情时说：\BibleQuote{天主的不可见之事，从创世以来，藉着所造之物就被清楚看见了，从而被理解，即祂永远的大能和祂为神的本性，以致他们无可推诿。}\BibleRef{罗 1:20}然而，天主造这些事并非为此，即使这一结果随之而来。因为祂并非要剥夺他们一切推诿的借口，才在他们面前设置如此伟大的教导体系，而是使他们能认识祂。但因他们没有认识祂，便使自己丧失了所有借口；然后，为了显示他们是如何被剥夺借口的，他说：\par

第21节：\BibleQuote{因为他们虽然认识天主，却没有以祂为天主而光荣祂。}\BibleRef{罗 1:21}\par

这是第一个最大的控诉；第二个则是他们也崇拜偶像，正如\Person{耶肋米亚}在责备他们时所说：\BibleQuote{这人民犯了双重的罪恶：他们离弃了我这活水的泉源，却为自己掘了破裂的水池。}\BibleRef{耶 2:13} 然后，作为他们虽认识天主却未将知识用于合适对象的标志，保禄引用了这件事本身，即他们知道神。因此他补充说：\BibleQuote{因为他们虽然认识天主，却没有以祂为天主而光荣祂。}\BibleRef{罗 1:21}并且他指出了他们陷入如此愚昧的原因。那是什么呢？他们事事依赖自己的推理。然而，他并没有这样措辞，而是用更尖锐的语言说：\BibleQuote{反而在推理中变为虚妄，他们愚昧的心就昏暗了。}\BibleRef{罗 1:21}因为如在无月之夜，有人试图在陌生的道路上行走，或在陌生的海上航行，他非但不能很快到达目的地，反而会迅速迷失。同样，他们试图走上通往天堂的道路，却毁灭了来自自身的光明，反而把自己交付给自己推理的黑暗，并在有形之物中寻找无形的那一位，在形状中寻找无形状的那一位，经历了一场最悲惨的沉船之祸。但除却所说的以外，他还指出了他们错误的另一个原因，当他说：\par

第22节：\BibleQuote{他们自称聪明，反而成了愚拙。}\BibleRef{罗 1:22}\par

因为他们对自己有极大的自负，且不愿走天主所命令的路，就陷入了无感的推理（1抄本作\OriginalTerm{思想}{\GreekText{διανοίας}}）。然后，为了展示并概述这是多么悲惨的浪潮，以及多么缺乏借口，他继续说：\par

第23节：\BibleQuote{并且将不可朽坏的天主的荣耀，改变为如同可朽坏的人、飞鸟、四足走兽和爬虫的形像。}\BibleRef{罗 1:23}\par

第一个指控是，他们没有找到天主；第二个是，当时他们有伟大而清晰（\EditorialNote{边注：智慧}）的方法去做这事；第三个是，他们竟仍自称有智慧；第四个是，他们不仅没有找到那位可敬畏的存在，甚至将祂贬低为魔鬼和石头木头。现在，他在致\WorkTitle{格林多书}中也打击了他们的傲慢，但那里与这里不同。因为在那里，他是从十字架给他们打击，说：\Quote{天主的愚妄比人更有智慧。}\BibleRef{格前1:25}但这里，他毫无比较地单独将他们的智慧拿来嘲笑，表明那是愚妄，纯粹是虚妄夸耀的表现。然后，为让你知道，当他们有天主的真知识时，却如此背信弃义地放弃了，他说\Quote{他们换了}。现在，换的人有东西可换。因为他们想要发现更多，不肯忍受所给他们的限制，于是连这些也被剥夺了。因为他们是新奇的贪求者，希腊人的一切便是如此。这也是他们彼此对立、亚里士多德起来反对柏拉图、斯多葛派对他大吹大擂（希腊文作\OriginalTerm{大吹大擂}{\GreekText{ἐφρυάξαντο}}，六份抄本作\InnerQuote{围护自己}\OriginalTerm{围护}{\GreekText{ἐφράξαντο}}，菲尔德倾向于后者）的原因；一个人成了这个的敌人，另一个则成了那个的敌人。所以，人们不应那么惊讶于他们的智慧，而该转身对他们愤怒仇恨，因为正是通过这智慧，他们成了愚昧人。因为若不是他们把自己的东西托付给推理、三段论和诡辩，他们就不会遭受所遭受的。然后，为加重对他们的指控，他嘲笑了他们全部的偶像崇拜。因为首先，改变本身已是极好的嘲笑对象。但改变成这些东西，更是无可推诿。那么，他们把天主的荣耀变成了什么？又拿什么来取代祂的荣耀？他们本应对祂有某些概念，例如，祂是天主，祂是万有的主，祂创造了原本不存在的东西，祂施行眷顾，祂关心他们。因为这些都是\Quote{天主的荣耀}。他们将它归于谁？甚至不归于人，而是\Quote{归于像必朽坏的人一样的偶像}。他们没有停在这里，甚至降到了畜生，或者说这些畜生的偶像。但请注意，我祈祷，保禄的智慧，他如何取了两个极端：至高无上的天主和最低下的爬虫；或者不如说，不是爬虫本身，而是它们的偶像；为能清楚表明他们明显的疯狂。因为他们本该对那无比卓越于一切者所具有的知识，却加给了无比毫无价值于一切者之上。但这与哲学家有何相干？有人会说。我刚才说的与他们关系最大。因为他们有发明这些东西的埃及人为师。而柏拉图，被认为比其余人更可敬，却以这些师傅为荣。他的老师（其遗言出自\WorkTitle{斐多}）对偶像怀有愚蠢的敬畏，因为他命人把公鸡祭给阿斯克勒庇俄斯，那里（即在他的神庙里）有这些野兽和爬虫的偶像。可以看到阿波罗和巴克斯与这些爬虫一同受崇拜。有些哲学家甚至把公牛、蝎子、龙和所有其余的虚荣举到天上。因为魔鬼在各方面都狂热地努力，要把人带到爬虫的偶像面前，把天主愿意高举到天上的那人排在一切最无感觉的东西之下。不仅从这一点，也从其他理由，你会看到他们的首要人物落入了刚才的议论。因为他收集了诗人的话，说我们应该在这些关乎天主的事上相信他们，如同他们有准确的知识，结果除了这些荒唐话的\Quote{连锁甜蜜}外，他没有拿出任何别的东西，然后说，这完全可笑的琐事要被当作真理。\par

第24节：\BibleQuote{因此，天主任凭他们随从心里的情欲，陷于不洁，以致彼此侮辱自己的身体。}\BibleRef{罗1:24}\par

由此他表明，即使律法的歪曲也是不虔敬的原因，但祂\Quote{任凭他们}，意思是祂放任不管。正如那在军队中发号施令的人，如果作战沉重压向他，他后退并离开，就等于把士兵交给敌人，不是因为他亲自推送他们，而是因为他撤去了自己的援助；同样，天主也离开了那些不愿接受来自祂之物、反而首先从祂那里跑开的人，尽管祂已完全尽了自己的本分。但请思量：祂为他们设立世界作为教义的形式；祂赐给他们理智和理解力，足以感知必要之事。这些人没有一样用来得救，反而将所领受的扭曲为相反。那么该怎么办？用强迫和暴力拖他们？但这不会使他们成为有德的人。此后，祂只好让他们自己去，祂也确实这样做了，为叫他们在无其他办法时，通过经验知道他们所贪恋的事物，从而逃避如此可耻的事（三份抄本加\OriginalTerm{合情合理}{\GreekText{εἰκότως}}）。因为假如有国王的儿子，羞辱父亲，宁愿与强盗、杀人犯和掘墓人为伍，并喜好他们的所作所为胜过父家，父亲会离开他，好让他通过实际经历，认识自己疯狂的极端。但他为何不提别的罪，例如凶杀、贪婪或其它此类的，而只提不洁？在我看来，他暗指当时的听众和将要接收这封信的人。\BibleQuote{陷于不洁，以致彼此侮辱自己的身体。}\par

请注意这里的强调语气，它极其严厉。因为意思是说，他们不需要别人对他们施以强暴的侮辱，而是敌人本该对他们做的，他们却对自己做了。然后，他再次提起指控说：\par

第25节。\BibleQuote{他们将天主的真理变成了虚谎，崇拜事奉受造物，而不崇拜事奉造物主。}\BibleRef{罗 1:25}\par

他把那些极可鄙的事特别列出，而那些看起来比其余更严重的则笼统论之；藉此表明崇拜受造物是希腊人的行为。且看他如何有力地陈述他的断言，因为他并不只是说\Quote{他们崇拜受造物}，而是说\Quote{超过崇拜造物主}：这样处处给指控增添新的力度，并通过比较剥夺他们一切减轻责任的余地。\Quote{祂是永远受赞美的。阿们。}他以此表示，天主未曾受到丝毫损害。因为祂自己\Quote{永远受赞美}。这里他证明，天主并不是为了自卫而任凭他们——祂自己并未受到任何损失。因为即使这些人侮辱祂，祂却不曾被侮辱，祂荣耀的特性也未受丝毫损伤，而是始终受赞美。因为如果人常常因哲学而不感受到别人对他的侮辱，那么不朽不变的天主，那不变不动的荣耀，更当如何呢？\par

因为人在这件事上相似天主：当他们不为那些企图对他们施以羞辱之人的行为所动，既不被别人的侮辱所激怒，也不因别人的殴打而烦恼，也不因别人的嘲笑而恼怒时。有人会说：这种事在本质上怎么可能呢？这是可能的，是的，确实可能，当你对所发生的事不感到烦恼时。又有人说：怎么可能不烦恼呢？不如说：怎么可能烦恼呢？请问你，如果你的小孩子侮辱你，你会把这种侮辱当作侮辱吗？你会烦恼吗？肯定不会。但如果你烦恼了，岂不显得可笑吗？同样，让我们对邻人也抱有这样的态度，这样我们就不会感到不快了。因为侮辱我们的人比孩子更无知。让我们甚至不要寻求免于侮辱，而是在被侮辱时忍受它们。因为这是唯一可靠的光荣。为何如此？因为这是你自己能掌控的，而那却是别人掌控的事。难道你没看见金刚石在受到打击时反弹回去吗？但你会说，是本性赋予了它这种特性。然而，你也有能力通过自由选择成为像它那样天生如此的状态。如何实现？难道你不知道那在火窑中的孩童没有被烧伤吗？以及达尼尔在狮穴中没有受到伤害吗？这种事如今也能发生。我们身边也有狮子：忿怒和贪欲，它们用可怕的牙齿撕裂掉进它们中间的人。（柏拉图\WorkTitle{理想国}卷八）你要成为像\OriginalTerm{那位}{\GreekText{ἔκεινον}}（三份抄本作\Quote{达尼尔}）一样，不要让这些情欲的牙齿咬住你的灵魂。但你会说，那完全是恩宠的作用。是的，因为自由意志的行为预先铺平了道路。因此，如果我们愿意把自己塑造成类似的品格，现在恩宠也近在咫尺。即使这些野兽饥饿，它们也不会触碰你的肋旁。因为如果它们看见一个仆人的身体就退缩，那么当它们看见基督的肢体（这就是我们信友）时，怎能不安静呢？但如果它们不安静，那是由于掉进它们中间的人的过失。的确，许多人纵容这些狮子：养娼妓、破坏婚姻、向敌人报复。因此，甚至还没有到狮穴底部，他们就被撕碎了。\BibleRef{达 6:24}但在达尼尔身上并未发生这样的事，如果我们如此决心，我们身上也不会发生，反而会发生比那时更伟大的事。因为狮子没有伤害他；如果我们保持清醒，那么那些伤害我们的人甚至会使我们得益。正如保禄从反对他和谋害他的人中得以显耀，约伯从许多鞭笞中，耶肋米亚从泥坑中，诺厄从洪水中，亚伯尔从背信弃义中，梅瑟从嗜血的犹太人中，厄里叟，以及古时的每一位圣贤，不是从松懈和安逸中，而是从磨难和试探中，才得以戴上灿烂的荣冠。因此，基督深知这是美名的根基，便对祂的门徒说：\BibleQuote{在世界上你们要受苦，但放心，我已战胜了世界。}\BibleRef{若 16:33}他们会说：那么，岂非有许多人被这些恐怖吓跑了吗？是的，但那不是试探的本性，而是他们自己的懈怠。然而，祂\BibleQuote{在试探中也会开辟一条出路，使你们能承受得住}\BibleRef{格前 10:13}，愿祂扶持我们众人，伸出祂的手，使我们光荣地获胜，获得永久的荣冠，藉着我们的主耶稣基督的恩宠和爱人之心\EditorialNote{五份抄本加上了其余部分，Field 多处如此}，愿光荣归于圣父，偕同圣神，世世无穷。阿们。\par

\SourceRecord{Translated by J. Walker, J. Sheppard and H. Browne, and revised by George B. Stevens.；Nicene and Post-Nicene Fathers, First Series；Vol. 11.；Edited by Philip Schaff.；Buffalo, NY: Christian Literature Publishing Co.,；1889.；<http://www.newadvent.org/fathers/210203.htm>.}

% structure: p0001 source=h1 target=section reason=page-title

\section{罗马书讲道 第四讲}

\BibleRef{罗 1:26-27}\par

\BibleQuote{因此，天主任凭他们陷于可耻的情欲：甚至他们的女人把顺性之用变为逆性之用；同样，男人也放弃了与女人的顺性之用，彼此欲火中烧。}\par

所有这些情感都是卑鄙的，尤其是对男性的疯狂淫欲；因为灵魂在犯罪中受苦更深，遭受的羞辱比身体在疾病中更甚。但请看，正如在教义方面一样，他在这里也剥夺了他们的借口，论到女人时，他说：\Quote{她们改变了本性的用途}。他的意思是，没有人可以说，她们是因为被阻止了合法的结合才落到这般地步，或是由于无法满足欲望才被驱入这骇人听闻的疯狂中。因为改变暗示着拥有。同样，在论述教义时，他也说过：\Quote{他们以谎言交换了天主的真理}。关于男人，他也同样用这话表明：\Quote{离开了女人的自然用途}。与此相似，他也将这些男人置于无可辩驳之地，指责他们不仅拥有满足的手段，却放弃已有的，去追求另一种，而且，在羞辱了那合乎本性的事物之后，竟然追逐那违反本性的事物。但违反本性的事物本身带有厌烦和不悦，所以他们甚至不能合理地声称其中有快乐。因为真正的快乐是合乎本性的。但当人被天主遗弃时，一切就都颠倒了。因此，不仅是他们的教义是撒殚式的，他们的生活也是魔鬼般的。当他论述他们的教义时，他将世界和人的理智摆在他们面前，告诉他们，藉着天主赐予他们的判断力，他们本可以通过可见之物被引导至创造主那里，但由于他们不愿意这样做，他们便无可推诿。在此，他用合乎本性的快乐代替世界，这种快乐本可以让他们更安全、更欣慰地享受，因而远离可耻的行为。但他们不愿意；因此他们完全超出了宽恕的界限，并且侮辱了本性本身。甚至比这些更可耻的是，当妇女寻求这些交合时，她们本应比男人更有羞耻感。在此，保禄的判断也值得钦佩，他如何面对两个相反的目标，却精确地完成了两者。因为他既想说话纯洁，又想刺痛听者。这两件事并非他所能同时做到，而是彼此妨碍。因为如果你说话纯洁，就无法严厉责备听者；但如果你想触及他的痛处，你就不得不以明确的言辞将赤裸的事实摆在他面前。但他那谨慎而神圣的灵魂能够精确地做到两者：通过提及本性，既加强了指控，又以此作为某种帷幕，保持其描述的纯洁。接着，他先责备了女人，然后继续论到男人，说：\Quote{同样，男人也离开了女人的自然用途}。这是极端败坏的一个明显证据：当两性都被抛弃，他（男人）被立为女人的导师，她（女人）被命令成为男人的助手，却彼此行敌对之事。也要思考他如何使用词语。他没有说他们彼此迷恋、彼此贪恋，而是说：\Quote{他们彼此在情欲中焚烧}。你看，整个欲望源自一种无法安于适当界限的越轨。因为凡违反天主所定律法的事物，都会贪恋异常之事，而非寻常之事。正如许多人时常放弃对食物的欲望，转而吃土和小石子；另一些人因过度干渴，甚至渴望泥浆；同样，这些人也陷入这种无法无天之爱的沸腾。但如果你问，这种情欲的强烈从何而来？它来自对天主的离弃；而对天主的离弃又来自那些离弃祂的人的不法行为：\Quote{男人与男人行可耻之事}。他的意思是，不要因你听到他们焚烧，就以为邪恶仅在于欲望。其实，大部分来自他们的纵欲，这又使他们的情欲燃成火焰。这就是为什么他没有说\Quote{被冲走}或\Quote{被胜过}——他在别处使用的表达——而是说什么？\Quote{行}。他们将这罪当作事业，不仅是事业，甚至是热切追求的事业。他称之为不体面的事，而不是情欲，并且恰如其分。因为他们既侮辱了本性，又践踏了法律。看看双方产生的巨大混乱。不仅头向下，脚也向上，他们成为自己和彼此之间的敌人，带来一种有害的争斗，一种比任何内战更无法无天的争斗，一种充满分裂且形式多样的争斗。因为他们将此分为四种新的、无法无天的种类。因为（据三份手稿）这场战争不是二重或三重，而是四重。那么想一想。二者本应成为一体，我指的是女人和男人。因为经上说：\BibleQuote{二人}，\BibleQuote{成为一体。} \BibleRef{创 2:24}但交合的欲望实现了这一点，使两性彼此结合。魔鬼夺走了这种欲望，并将其流向转为另一种方式，从而分裂了两性，使一体成为两部分，违反天主的法律。因为经上说：\BibleQuote{二人成为一体}；但他将一体分为二：因此这就是第一场战争。同样，这两部分他挑动彼此作战，也各自与自己作战。因为甚至女人也彼此凌辱，而不仅仅是男人。男人彼此对抗，也对抗女性，如同夜间的战斗。你看到第二、第三场战争，第四、第五场；还有另一场，因为除了已经提到的，他们还违反本性本身行事。因为魔鬼看到，主要是这种欲望将两性吸引在一起，他便决心切断这种联系，从而毁灭人类，不仅通过使他们无法合法交合，还通过挑动他们彼此战争、彼此分裂。\par

\Quote{他们在自己身上受到了他们错谬所应得的报应。}请看，他如何再次回到邪恶的根源，即由他们的教义产生的不敬，并称这是那不法的报应。因为当论及地狱和惩罚时，似乎目前对那些不虔敬和刻意选择这种生活的人而言并不可信，甚至被轻蔑，他便表明惩罚就在这快乐本身之中。（\Person{柏拉图}\WorkTitle{泰阿泰德篇}第176页7行。）但若他们不察觉，反而仍感快乐，不必惊讶。因为即便是那些疯狂和患了癫狂的人（参\Person{索福克勒斯}\WorkTitle{埃阿斯}265-277行），虽严重伤害自己，成为令人怜悯的对象，他人为他们哭泣，他们自己却微笑并庆幸所发生的事。然而，我们不仅不因此说他们免于惩罚，反而正因此他们处于更严重的报应之下，因为他们对自己所处的境况毫无察觉。因为需要争取的是清醒之人的赞成，而非混乱之人的。古时，这事甚至被视为法律，他们中的一位立法者禁止家奴在干燥时（即除沐浴外）使用香膏，也不可豢养少年，将这一荣誉——或更确切地说，这一可耻——赐予自由人。然而，他们并不认为此事可耻，反而视之为一种特权，且对奴隶而言过于高贵，雅典人——最智慧的民族，以及他们中如此伟大的\Person{梭伦}——只允许自由人享有。人们可以看到许多哲学家的书籍充满这种病症。但我们并不因此说这事是合法的，而是说接受这法律的人是可怜的，应流许多眼泪。因为他们所受的对待与卖淫的妇女相同。或者更确切地说，他们的处境更悲惨。因为就妇女而言，交合虽不合法，却合乎自然；但这既违背法律又违背自然。因为即使没有地狱，没有惩罚的威胁，这也比任何惩罚更糟。然而，若你说\Quote{他们从中找到快乐}，你告诉我的正是加重报复的事。因为假设我看见一个人赤身奔跑，全身沾满泥浆，却不遮盖自己，反而为此欢欣，我不会与他一同快乐，反而应为他不察觉自己在做可耻之事而哀哭。但为更清楚地显示这行为的凶残，请容忍我再举一例。假如有人判一个处女住在密室里（\OriginalTerm{被关在密室中}{\GreekText{θαλομευομένην}}），并与无理性的野兽交合，而她对此交合感到快乐，难道她不应尤其值得哀哭，因为她甚至不察觉自己的不幸而无法摆脱这痛苦？这显然人人都明白。但若那事是悲惨的，这事也不亚于它。因为被自己的亲属侮辱比被陌生人侮辱更可怜：这些（五份手稿作\Quote{我认为}）甚至比谋杀者更恶，因为死亡倒比在这种侮辱下活着更好。谋杀者使灵魂与身体分离，而这人却连灵魂带身体一同毁灭。你无论说出什么罪恶，都不会有与这不法相等的。若那些受害者察觉了，他们会宁愿死一万次也不愿受这恶。因为这种侮辱没有、确实没有比这更严重的恶。因为当论及奸淫时，保禄说：\Quote{凡人所犯的每一罪恶都在身体以外，但那犯奸淫的，却是得罪自己的身体。}\BibleRef{格前 6:18}；那么对这疯狂之事，我们该说什么？它比奸淫更甚，无法言表。因为我不仅要说你成了女人，而且你失去了男儿气，既未变成那本性，也未保持你原有的，而是同时背叛了二者，应被男人和女人赶出并用石头砸死，因为你伤害了两种性别。而且为让你明白这事的真正力量，若有人来向你保证他要把你从人变成狗，你岂不避之如瘟疫？但看！你把自己从人变成了比狗更可耻的动物，因为狗尚有用途，而这样自甘堕落的人毫无用处。或者，若有人威胁让男人怀孕生子，我们岂不愤慨？但看！现在陷入这狂热的人对自己做了更恶的事。因为变成女人本性是一回事，继续做男人却成了女人——或者二者皆非——则是另一回事。但若你从其他方面了解这恶的严重性，请问立法者为何惩罚阉割他人者？你必会发现，除了他们损毁本性之外，绝无其他原因。然而他们所做的不义于此无法相比。因为有些被阉割的人阉割后仍能在许多方面有用。但自甘堕落者则毫无价值。因为不仅灵魂，连如此受对待的身体也受羞辱，应被处处驱逐。多少地狱才够这样的人？但若你嘲笑听到地狱，不相信那火，请记住\Place{索多玛}。因为我们确实看到，就在今世，我们看到了地狱的雏形。因为许多人完全不相信复活后的事，如今听到不灭的火，天主便借现世之事使他们醒悟。因为\Place{索多玛}的焚烧和那大火正是如此！去过那地方、亲眼看见那神圣灾祸和来自上天的闪电效果的人都知道。\BibleRef{犹 1:7}请想那罪有多大，竟迫使地狱在时间之前显现！因为许多人轻视祂的话，天主便以行为用一种新奇的方式向他们显示了地狱的图像。因为那雨是反常的，因为那交合是违反本性的，那雨淹没了大地，正如情欲淹没了他们的灵魂。因此那雨与常雨相反。它不仅没有使大地孕育果实，反而使之对接收种子无用。因为男人的交合也是如此，使这样的身体变得比\Place{索多玛}的土地更无用。还有什么比自甘堕落的人更可憎、更可咒骂？哦，何等疯狂！哦，何等迷乱！这淫荡狂欢的情欲从何而来，使人的本性成为敌人所能做的一切？或者甚至更糟，因为灵魂比身体更尊贵。哦，你比无理性的动物更无感觉，比狗更无耻！因为动物之间从不发生这样的交合，本性自知其界限。但你却以彼此施加的侮辱，使我们的种族沦为低于无理性之物。那么这些邪恶从何而生？从奢侈，从不认识天主。因为一旦他们弃绝了对祂的敬畏，一切美好便立即崩溃。\par

如今，为使此事不致发生，让我们常将天主的敬畏摆在眼前。因为确实，再没有什么比脱离这锚更摧毁一个人，也没有什么比不断仰望它更拯救人。因为若因将人放在眼前，我们在犯罪时便更退缩，且常常甚至因对更好品级的仆役感到羞愧，而不做任何错事；那么试想，将天主放在眼前，我们将享有何等安全！因为当处于此种状态时，魔鬼绝不会攻击我们，因为他劳而无获。但若他看见我们四处游荡，无缰而行，他便能从我们自身得一个开端，随后将我们驱往任何地方。正如市场上那些轻率的仆人，放下主人交托的必要工作，随便聚焦于偶然遇到的人，并在那里消磨闲暇；当我们偏离天主的诫命时，也遭受同样的事。因为我们随即驻足，羡慕财富、美貌及其他与我们无关的事物，正如那些仆人注视表演戏法的乞丐，然后迟到，在家中被痛打。许多人因跟随行为同样不当的他人，而偏离定准的道路。但我们切不可如此。因为我们被派遣去处理许多急务。若我们放下这些，却张口呆望无用之物，所有时间都将徒然浪费，毫无益处，并要受最重惩罚。因为若你愿让自己忙碌，你便有应当惊叹之物，终日志在它上；这些事不是笑料，而是惊叹与多般赞颂的对象。正如羡慕可笑事物的人，自己常成为那样，甚至比惹人发笑者更甚。为使你不陷于此，立刻逃离它。因为，请问，你为何见财富便张口呆望、心绪不定？你看见什么如此奇妙，能让你注目不移？这些金鞍马匹、这些仆役，部分野蛮人，部分阉人，还有华服，以及在其中变得完全软弱的灵魂，傲慢的额头，忙碌与喧闹？这些事何值得惊叹？它们比市场上跳舞吹笛的乞丐好在哪里？因为这些人同样染上严重的德行饥荒，跳着比那些乞丐更可笑的舞，一时被带往奢华宴席，一时被引至娼妓居所，一时又被引向一群奉承者和众多依附者。但若他们真的穿戴黄金，正因此他们最可怜，因为那些与他们无关之物，却成为他们热切欲望的对象。现在，我恳求你，不要看他们的衣服，而要敞开他们的灵魂，看它是否充满无数伤口，身披破衣，穷困无援！那么这疯狂炫耀有何用？因为贫而德行远胜于富而邪恶；因为穷人自身享受灵魂的一切喜乐，因内在富足甚至不觉外在贫穷。但国王，纵情于那些完全不属于他的事物，却在他最切身的事上受罚，即灵魂、思想与良心，这些将随他前往另一个世界。因此，既然我们知道这些，让我们放下镀金的衣裳，拿起德行及其所带来的快乐。因为如此，在今世及来世，我们将享受极大的欢愉，藉着我们的主耶稣基督的恩宠与对人的爱情。愿光荣归于圣父及圣神，藉着祂并与祂，至于无穷之世。阿们。\par

\SourceRecord{Translated by J. Walker, J. Sheppard and H. Browne, and revised by George B. Stevens.；Nicene and Post-Nicene Fathers, First Series；Vol. 11.；Edited by Philip Schaff.；Buffalo, NY: Christian Literature Publishing Co.,；1889.；<http://www.newadvent.org/fathers/210204.htm>.}

% structure: p0001 source=h1 target=section reason=page-title

\section{罗马书讲道词 第五讲}

\BibleRef{罗 1:28}\par

\BibleQuote{他们既然不肯保留对天主的认识，天主就任凭他们存有败坏的心思，去作那些不适宜的事。}\par

怕人以为他长久停留在论逆性之罪上是在暗示他们，他便接着也谈及别的罪，并为此将整篇讲论作为对他人而言。正如他常对信徒论罪时所为，欲表明这些罪应避免时，他引外邦人为例，说：\Quote{不要随着情欲的贪恋，像那不认识天主的外邦人一样。}\BibleRef{得前 4:5}又说：\Quote{不要忧伤，像那些没有望德的人一样。}\BibleRef{得前 4:13}同样在此，他表明这些罪属于他们，并夺去他们一切托辞。因为他说，他们的妄行不是出于无知，而是出于惯行。正因如此，他没有说\Quote{他们不认识天主}，而是说\Quote{他们不肯保留对天主的认识}；意思是，这罪更多出于顽固执拗的败坏决心，而非突然的冲动，并表明这罪不属于肉身（如某些异端者所说），而属于心思，邪恶的贪欲属于它，邪恶的源头也从此流出。因为心思一旦失去辨别力，其他一切便被拖离正轨而倾覆——当执缰者败坏时！（\Person{柏拉图}\WorkTitle{费德罗篇}246A、B）\par

第29节。\BibleQuote{充满了各种不义、邪恶、贪婪、恶毒。}\BibleRef{罗 1:29}\par

看，这里一切都加重了。因为他说\Quote{充满了}，又说\Quote{各种}，并先总提恶毒，后又逐条细述，而且这些也越了分，说\Quote{满有嫉妒、凶杀}，因为后者出自前者，如亚伯尔和若瑟的例子，然后说了\Quote{争竞、诡诈、毒恨}；\par

第30节。\Quote{谗毁的、背后说人的、憎恨天主的、侮慢人的，}并将那些在许多人看来无足轻重的事列在指控中，进一步加重了控诉，直逼他们恶行的堡垒，称他们为\Quote{自夸的}。因为即使在善行中自高自大也会丧失一切，若有人在罪中自高，那还有什么惩罚是他不配的呢？因为这样的人再也不能悔改了。接着，他说\Quote{捏造恶事的}，表明他们不满足于已有的，甚至还发明别的。这又像是那些意志坚决、认真的人，而非被匆忙拖迫的人；在列举了几种恶意，并表明在此他们也违背了本性（因为他说\Quote{违背父母的}）之后，他便转而指摘那大瘟疫的根源，称他们是\BibleRef{罗 1:30}\BibleRef{格前 5:2}，\par

第31节。\BibleQuote{无亲情、不解怨。}\BibleRef{罗 1:31}\par

因为基督自己也宣称这是邪恶的原因，说：\Quote{由于罪恶的增多，许多人的爱情必会冷淡。}\BibleRef{玛 24:12}圣保禄在此也这样说，称他们为\Quote{背约的、无亲情的、不解怨的、不怜悯人的}，表明他们甚至背弃了本性的恩赐。因为我们彼此之间有一种天然的同感，就连野兽之间也有。因为经上说：\Quote{一切走兽都爱自己的同类，同样，每个人也爱自己的邻人。}\BibleRef{德 13:15}但这些人变得比它们更凶残。于是，他以这些见证向我们证明邪说给世界带来的混乱，并清楚表明两种病症都来自患病者的疏忽。此外，他还表明，正如在教义上一样，在此他们也毫无托辞；因此他说：\par

第32节。\BibleQuote{他们虽明知天主的判决，行这样事的人是当死的，然而他们不但自己去行，还喜欢别人去行。}\BibleRef{罗 1:32}\par

在这里他预先设了两个反驳，首先予以排除。因为他说：你有什么理由说你不知道该做的事？即使你不知道，你也有罪，因为你离弃了教导你的天主。但我们已经用许多论据证明你知道，却故意犯法。难道是受情欲所牵引？那么你为什么既与之合作，又赞美它？因为他说：他们\Quote{不但自己去行，还喜欢别人去行}。既然首先提出了这更严重、更不可饶恕的罪，好以此了结它（或\Quote{定罪于你}，\OriginalTerm{（或\Quote{定罪于你}）}{\GreekText{ἵνα} \GreekText{ἑλῃ}}）；因为赞美罪的人远不如犯罪者的可恶；他以此法首先提出这一点，随后便更有力地对付他，这样说：\par

第二章第1节。\BibleQuote{所以，你是无可推诿的，人啊！无论你是谁，你判断人；因为你判断别人，就是定你自己的罪。}\BibleRef{罗 2:1}\par

这些话是针对统治者说的，因为当时那城被赋予了统治世界的权柄。因此他预先说：你是在剥夺自己的辩护，无论你是谁；因为当你判断奸夫，自己却犯奸淫时，虽然没有别人判断你，但你在对有罪者的判决中已经定了自己的罪。\par

第2节。\BibleQuote{我们知道，天主的裁判是按真理，对行这样事的人。}\BibleRef{罗 2:2}\par

因为惟恐有人说：\Quote{至今我逃脱了}，使他害怕，他就说：在天地间并非如此。因为在此世（柏拉图在《\WorkTitle{Theæt. et Phædon.}》中）一个人受罚，另一个人做同样的事却逃脱。但来世并非如此。他既已说了：那审判者那时必知道正义，但他没有补充从何知道；因为那是多余的。因为论及不敬，他既表明了不敬者即使认识天主仍如此，也表明他从何处得到那认识，即从受造物。既然并非人人明白，他也给出了原因；但此处他视之为公认而略过。但当他说：\Quote{凡你审判人的}，他不是只对统治者说话，也是对私人和臣民。因为所有人，即使没有公义之座、没有刑吏、没有枷锁，他们也在交谈和公共集会（希腊文：\OriginalTerm{公共集会}{\GreekText{κοινοἵς} \GreekText{συλλόγοις}}）中，凭良心的投票审判犯罪者。没有人敢说奸夫不应受罚。但他却说，他们谴责别人，却不谴责自己。为此，他强烈地反对他们，说道：\par

第3节。\Quote{你以为这事吗？}（3份抄本省略\Quote{这事}）\Quote{人哪，凡审判那些行这样事的人，自己却行同样的事，你以为你能逃脱天主的审判吗？}\BibleRef{罗 2:3}\par

因为他既已表明世界的罪是大的，无论是从教义还是行为，并且他们虽是智者，虽有受造物引领，却仍然犯罪，不仅离弃天主，还选择爬虫的偶像，而且不尊重德行，不顾本性之抗拒，甚至违反本性而侍奉恶习：他接着要表明，行这样事的人也要受罚。他确实立刻指出一种刑罚，即提及他们的行为本身。因为他说：\Quote{他们在自己身上受了那错误应得的报应。}但因他们不觉察，他又另提一种他们最惧怕的刑罚。事实上，他主要指向的就是这个。因为当他说：\Quote{天主的审判是按真理}，他说的正是这个。但他又用其他进一步的理由确立同一论点，说：\Quote{人哪，凡审判那些行这样事的人，自己却行同样的事，你以为你能逃脱天主的审判吗？}你没有被你自己的审判宣告无罪，难道你会藉着天主的审判逃脱吗？谁会这样说呢？其实你已定了自己的罪（3份抄本：\Quote{且没有被宣告无罪}）。既然审判庭如此严厉，你连自己都不能宽恕，那无过错的、至公义的天主岂不更肯定做同样的事？但你自己定了罪，难道天主会赞同你、赞美你吗？这岂有道理？而且你比你审判的人更该受更重的刑罚。因为仅仅犯罪，与犯下你曾谴责别人所犯的同样的罪，是不同的。看，他如何加重了控诉！因为你若惩罚一个犯较小罪的人，尽管这将使你自己羞惭，何况天主，祂永不羞惭，岂不更要定你的罪，更严厉地谴责你，因你犯了更大的过犯，且你已由你自己的清算定了罪？但如果你说：我知道我该受罚，但因祂的容忍而轻看它，并因没有立即受罚而自信，这恰恰是你应当害怕战栗的原因。因为你还没有受罚，并不会导致你不再受罚，反而会使你若不悔改，受更重的刑罚。所以他接着说：\par

第4节。\Quote{还是你藐视祂丰盛的恩慈、宽容和忍耐，不晓得天主的恩慈是领你悔改吗？}\BibleRef{罗 2:4}\par

因为在赞美天主的忍耐之后，他表明这对留意的人有极大的益处（即引领罪人悔改），他又增加了恐惧。因为对那些正确利用它的人，它是安全的根基；对那些轻视它的人，它导致更大的惩罚。因为每当你说出这个共同观念，即天主因为祂是良善、忍耐的，就不施行公义，他说，你只是提到那会使惩罚更严厉的事。因为天主显示祂的恩慈，是为叫你脱离罪，而不是添加罪。所以你若不善用，审判将更可怕。因此，天主的忍耐是禁戒犯罪的主要理由，而不是利用恩惠作为顽固的借口。因为如果祂忍耐，祂必定会惩罚。这从哪里显现？从接下来所说的。因为若是罪大恶极而恶人未受报应，则他们必定要受报应。因为如果人不忽视这些事，天主岂会忽视呢？所以从这一点他引入了审判的主题。因为指出许多若不悔改就有罪的人，却仍在此世不受罚，这必然引入审判，并且加剧。因此他说：\par

第5节。\Quote{但你竟任着你刚硬和不悔改的心，为自己积蓄忿怒。}\BibleRef{罗 2:5}\par

因为若一个人既不能被良善软化，也不能被恐惧回转，还有什么比这样的人更顽固呢？因为在他展示了天主对人的良善之后，他就显示了祂的审判，对于甚至这样仍不归向悔改的人，那审判是无法承受的。并且请注意，他是多么恰当地使用这些词语！他说：\Quote{你为自己积蓄忿怒}，以此清楚表明什么是确实积蓄的，并指出不是祂审判，而是被定罪的人自己制造了这忿怒。他说：\Quote{你为自己积蓄}，而非天主为你积蓄。因为祂做了一切合宜的事，创造了你，赋予你辨别善恶的能力，对你显示忍耐，召叫你悔改，并警告了那可怕的日子，用各种方式吸引你悔改。但如果你继续顽梗，\Quote{你就在忿怒和启示（据所有手稿，除两份外）与天主公义审判的日子为自己积蓄忿怒}。因为恐怕你听到忿怒时就想到某种激情，他加上了\Quote{天主的公义审判}。他很有理由说到\Quote{启示}，因为当各人按自己的功过领受时，这事才被启示。因为在此世，许多人常常不公地烦扰和伤害他人。但在来世则不然。\par

第六至第七节。\BibleQuote{祂将按各人的行为报应各人：凡恒心行善，}等等。\par

既然他因谈论审判和将来的惩罚而显得令人敬畏和严厉，他并没有如人所料立即进入报复的话题，而是将话题转向更甜蜜的事，即善行的报酬，他接着说：\par

第七节。\BibleRef{罗 2:7} \BibleQuote{凡恒心行善，寻求荣耀、尊贵和永生的人，就有永生。}\par

在这里，他也唤醒了那些在试炼中退缩的人，表明只信赖信德是不对的。因为那个审判也要审问行为。但是请注意，当他谈论将来的事时，他无法清楚说出福乐，只说荣耀和尊贵。因为这些福乐超越人所拥有的一切，他没有从世间取任何形象来显示，而是借着那些在我们中间带有光明表象的事物，尽他所能地将它们摆在我们面前，即借着荣耀、尊贵和生命。因为这些都是人努力追求的，但那些事不是这些，而是远比这些更好，因为那是不能朽坏、不能死的。请看，他如何借着说不朽坏，为我们打开了通往身体复活的门。因为不朽坏属于可朽坏的身体。然后，既然这还不够，他又加上了荣耀和尊贵。因为我们众人将复活而不朽，但并非都得荣耀，而是有些得惩罚，有些得生命。\par

第八节。\BibleRef{罗 2:8} \BibleQuote{惟有结党不顺从真理，}他说。再次，他以这种方式剥夺了那些活在邪恶中的人的借口，并表明他们是因为好争论和疏忽而坠入不义。\par

\BibleQuote{反顺从不义。}看，这又是一项控告。因为那逃避光明、选择黑暗的人能提出什么辩护呢？他没有说\Quote{被强迫}、\Quote{被支配}，而是说\Quote{顺从不义}，好让人知道堕落是自由选择的结果，罪行并非出于必须。\par

第九节。\BibleRef{罗 2:9} \BibleQuote{将患难、困苦加给一切作恶的人。}\par

也就是说，一个人若是富有，若是执政官，若是至高无上的君主（所以菲尔德：几份抄本和版本说\Quote{皇帝本人}），没有一个人的身份能改变审判的账目。因为在审判中，地位没有位置。既然他已指出疾病极其严重，又加上了原因，即源自混乱之人的疏忽，并最终指出毁灭等着他们，而改正并不难，在惩罚上他也再次给了犹太人更重的分量。因为享有更大教导的人，如果违法，也应当承受更大的报应。因此，我们越有智慧或越有权势，犯罪时受的惩罚也越重。因为如果你富有，你将被要求比穷人交出更多的钱财；如果你比别人更有智慧，更严格的服从；如果你被授予了权威，更光耀的善行；所以对于其他一切事物，你都必须按你的能力交出相应的行动。\par

第十节。\BibleRef{罗 2:10} \BibleQuote{却将荣耀、尊贵与平安加给一切行善的人，先是犹太人，后是外邦人。}\par

他这里所说的犹太人是指谁？他谈论的又是哪些外邦人？是指基督来临之前的人。因为他的论述尚未进入恩宠的时代，而是仍逗留在较早的时期，先是从远处拆毁并清除希腊人与犹太人之间的隔阂，以便当他在恩宠之事上这样做时，不再显得是在构思某种新颖且贬低人的观点。因为在早期，当这恩宠尚未如此彰显，犹太人的地位在众人面前是庄严、著名和荣耀的时候，并无分别；那么，在如此伟大的恩宠显示之后，他们还能为自己说什么呢\OriginalTerm{他们还有什么可说的？}{\GreekText{τίνα} \GreekText{ἂν} \GreekText{ἔχοιεν} \GreekText{λόγον} \GreekText{εἰπεἵν}}；这就是为什么他以如此大的热诚确立这一点。因为当听者得知这在早期就已成立，他在信德之后会更加接受这一点。但他这里所说的\Quote{希腊人}，并不是指那些崇拜偶像的人，而是指那些敬拜天主、服从自然法律、严格持守一切（除了有助于虔敬的犹太礼仪之外）的人，如默基瑟德和\OriginalTerm{他们等}{\GreekText{οἱ} \GreekText{περί}}、约伯、尼尼微人、科尔乃略。在这里，他首先突破了割礼与未受割礼之间的隔阂；并预先从远处消除了这一区别，以便不被怀疑地做到这一点，并像是被另一场合所迫而切入其中，这始终是他宗徒智慧的一个特点。因为如果他是在恩宠时期显示这一点，他所说的话就会显得非常可疑。但在描述了占据世界的恶习以及邪恶之路的终结之后，从那里连续过渡到对这些点的处理，使他的教导不可疑。而且他正是这个意思，并为此目的如此编排，从此处可以明显看出：因为如果他无意达成此目的，他只需说\BibleQuote{照你的硬心和不悔改的心，为自己积蓄忿怒，直到忿怒的日子}，然后便不再谈论这个话题，因为那已经完整了。但是，因为他所着眼的不只是谈论将来的审判，也要表明犹太人在这样的希腊人面前毫无优势，因此不可心高气傲，他便更进一步，依次谈论他们。但请注意！他使听者感到恐惧，向那人提出那可畏的日子，告诉他生活在邪恶中是何等邪恶，向他表明没有人因无知而犯罪，也没有人不受惩罚，即便现在不受惩罚，将来也一定受惩罚；然后他接下来想要证实法律的教导并不是一件多么重要的事。因为惩罚和赏报都取决于行为，而不在于割礼与未受割礼。既然他已经说过外邦人绝不可能不受惩罚，并以此为当然，又在此基础上证实了他也将得到赏报，他接下来就表明法律和割礼是多余的。因为他在这里主要反对的是犹太人。因为他们有点吹毛求疵：首先出于骄傲，不屑于与外邦人同列；其次认为如果信德能消除所有罪恶，那是可笑的；因此他先指控外邦人——他正是在为外邦人说话——为的是可以毫无嫌疑、大胆地攻击犹太人。然后，在谈到惩罚的问题时，他表明犹太人非但没有从法律中得到任何好处，反而被法律所重压。这就是他之前的意图。因为如果外邦人在这方面无可推诿——当受造物和他自己的理性引导他时，他仍不改正——那犹太人更是如此，因为除了这些之外，犹太人还有法律的教导。在说服了他在别人犯罪的情况下轻易接受这些推论之后，现在他迫使他在自己的罪行上也如此行，即便违背他的意愿。而且为了让他所说的话更容易被接受，他也以更美好的事物引导他向前，这样说道：\BibleQuote{但是荣耀、尊贵与平安归于一切行善的人，先是犹太人，后是外邦人。}因为在这里，无论一个人拥有什么好东西，他都是带着争斗得到的，即使他是富人、王子或国王。即使他与别人没有不和，也常常与自己不和，心中充满战争。但在那里却并非如此，一切都平静、无忧，并拥有真正的平安。然后，他从以上所说的证实了，那些没有法律的人也将享受同样的福分，他在下面的话中补充了他的理由：\par

第11节。\BibleQuote{因为天主不看人的情面。}\BibleRef{罗 2:11}\par

因为当他说，犹太人也好，外邦人也好，若犯罪都受罚时，他无需推理；但当他想证明外邦人也受尊敬时，他便需要为此建立根基；因为一个既未听过法律，也未听过先知的人，若因行善而受尊敬，似乎奇妙而过分。这也是为什么（正如我以前也说过的）他在恩宠时代之前训练他们的耳朵，以便日后能更容易地连同信德一同带入对这些事的认同。因为在这里，他全然不被怀疑，似乎没有在论证自己的观点。因此，在他说了\BibleQuote{凡行善的，得光荣、尊贵与平安，先是犹太人，后是外邦人}之后，他又加上\BibleQuote{因为天主不看情面}。奇妙啊！他取得了何等超越胜利的成就！因为他通过归谬法表明，在天主面前情况相反是不合适的。因为那将是偏袒，但天主并非如此。他没有说\Quote{因为若非如此，天主就成了偏袒的}，而是更庄严地说\BibleQuote{因为天主不看情面}。这表明它不是人的品质，而是行为的差异，天主正是按此进行审判。通过这样说，他表明犹太人与外邦人的区别不在于行为，而只在于身份。由此得出的结论是：并不是因为是犹太人或是外邦人，就有人受尊敬、有人受羞辱，而是出于行为才得到不同的对待。但他没有这样说，因为那会激起犹太人的愤怒，而是进一步陈述了某种更深的道理，从而降低他们的傲慢精神，并压制它以接纳另一真理。但这是什么？就是下一个论点。\par

第12节。\BibleQuote{因为凡没有法律而犯罪的人，也必没有法律而丧亡；凡在法律之下犯罪的人，将依据法律受审判。}\BibleRef{罗 2:12}\par

因为在这里，正如我以前说过的，他不仅表明犹太人与外邦人的平等，而且表明犹太人反而因有法律的恩赐而负担更重。因为外邦人没有法律而受审判。但这里的\Quote{没有法律}\OriginalTerm{没有法律}{\GreekText{ἀνόμως}}所指的并非更糟的处境，而是更容易的，即他没有法律来控告他。因为\Quote{没有法律}（即没有由此产生的谴责），他仅凭自然的推理就被定罪；而犹太人\Quote{在法律之下}，即既有自然又有法律来控告他。因为他所享有的关注越大，所受的惩罚也越大。你看，他多么强调犹太人急需快速投奔恩宠！因为他们说他们不需要恩宠，因法律而成义，但他表明他们比外邦人更需要恩宠，因为他们要受更重的惩罚。然后他又添加另一个理由，进一步为所说的话辩护。\par

第13节。\BibleQuote{因为听法律的，在天主面前不是义人。}\BibleRef{罗 2:13}\par

他恰当地加上了\Quote{在天主面前}；因为或许在人面前，他们能显得高贵并夸耀大事，但在天主面前却完全相反——只有法律的遵行者才能成义。你看他以何等优势进行争辩，把他们的话扭转到相反的意思。因为如果你们声称靠法律得救，在这方面，他说，外邦人将站在你们面前，当他们被看到是法律所写之事的实行者时。一个人从未听过怎能实行？他说，这不仅可能，而且有比这更多的事。因为不仅可能没有听过而实行，甚至可能听过却不实行。最后这一点他讲得更清楚，并且以更大的优势胜过他们，当他说：\BibleQuote{你这教导别人的，难道不教导自己吗？}\BibleRef{罗 2:21}但在这里，他仍在证明前一点。\par

第14节。\BibleQuote{因为当外邦人，没有法律时，却顺着本性去做法律所要求的事，这些人虽然没有法律，对自己就是法律。}\BibleRef{罗 2:14}\par

他的意思是，我并非排斥法律，而是就这一点也使外邦人成义。你看，当他削弱犹太人的自负时，并没有给自己留下诋毁法律的口实，反而通过赞扬法律、彰显其伟大来巩固自己的全部立场。但每当他说\Quote{顺着本性}时，他指的是自然的推理。他表明其他人比他们更好，而且更优异的是，这些人没有接受法律，没有犹太人似乎占优势的东西。基于此，他认为他们应当受赞美，因为他们不需要法律，却显出了法律的一切行为，将行为而非文字铭刻在心中。这就是他所说的。\par

第15节。\BibleQuote{他们显示法律的功用刻在他们心上，良心也为此作证，他们的思念互相控告或辩白。}\BibleRef{罗 2:15}\par

第16节。\BibleQuote{在天主藉耶稣基督，照我的福音审判人隐秘事的日子。}\BibleRef{罗 2:16}\par

请看，他如何再次将那一日摆在他们面前，使之临近他们，打垮他们的狂妄，并指出：那些没有法律却努力履行法律之事的人，反而更应受到尊崇。但是在宗徒的明智中，最令人惊叹之处，现在值得提及。因为他从所给的依据中，已经表明外邦人比犹太人更伟大；但在推理和结论中，他却未明说，以免激怒犹太人。但为使我说得更清楚，我要引用宗徒的原话。因为他说过，不是听法律的人，而是行法律的人，才将成义；接着他本可以说：\\Quote{外邦人没有法律，若顺着本性行法律之事，}他们比受法律教导的人更优越。但他没有这样说，而是止于对外邦人的称赞，暂时不进行比较论述，以便至少能让犹太人接受所说的话。所以他没有像我所做的那样措辞，而是怎样呢？\\BibleQuote{外邦人没有法律，若顺着本性行法律之事，他们虽然没有法律，自身就是法律；他们显示法律的功用刻在他们心上，他们的良心也一同作证。}\\BibleRef{罗 2:14-15} 因为良心和理智足以代替法律。藉此，他首先表明天主创造人使其自主，能选择德行，避免邪恶。不要惊讶他多次证明这一点，因为由于那些说\\Quote{基督为何现在才来？先前（多数手稿作：这强大的）眷顾的安排在哪里？}的人，这个主题对他极为必要证明。他现在正是要驳斥这些人，因为他表明即使在先前、在法律赐下之前，人类（\\OriginalTerm{人类}{nature}）就充分享受了天主眷顾的关怀。因为\\BibleQuote{关于天主所能知道的事，在他们身上是明显的}\\BibleRef{罗 1:19}，他们知道什么是善、什么是恶；藉此他们判断别人，他为此责备他们，当他说：\\BibleQuote{你判断别人，正是定你自己的罪}\\BibleRef{罗 2:1}。至于犹太人，除了已提到的，还有法律，不仅是理智和良心。那么，他为什么说\\BibleQuote{控告或也辩护？}\\BibleRef{罗 2:15}？因为，如果他们刻有法律，并在他们身上显示出法律的功效，理智怎么还能控告他们呢？但他不再只讲那些行善的人，而是普遍论及人类（\\OriginalTerm{人类}{nature}）。因为那时，我们的理性站立起来，有的控告，有的辩护。在那审判台前，人不需要别的控告者。然后为增加他们的恐惧，他没有说人的罪行，而是说人的隐秘事。因为他说过：\\BibleQuote{你判断那些行这样事的人，自己却行同样的事，你以为你能逃脱天主的审判吗？}\\BibleRef{罗 2:3}，为使你们不要指望得到像你们自己判决那样的判决，而要知道天主的审判远比你们自己的更精确，他提到了\\BibleQuote{人的隐秘事}\\BibleRef{罗 2:16}，并加上\\BibleQuote{藉耶稣基督，按照我的福音}\\BibleRef{罗 2:16}。因为世人只根据外表行为审判。前面他也只提到了圣父，但一旦他以恐惧压垮了他们，就也提到了基督。他不仅这样做，而且即使在提及圣父之后，他也这样引入祂。并通过同样的事，提高了他宣讲的尊严。因为，他意思是，这宣讲公开说出本性预先教导的事。你看见他以何等的智慧，将他们既与福音联系，又与基督联系，并表明我们的事并非在此止步，而是继续前行。这一点他以前也已证明，当他说：\\BibleQuote{你为自己积蓄天主的忿怒，直到忿怒的日子}\\BibleRef{罗 2:5}；而在此又再说：\\BibleQuote{天主将审判人的隐秘事}\\BibleRef{罗 2:16}。\par

现在让每个人进入自己的良心，清算自己的过犯，严格地审问自己，好使我们那时不随同世界被定罪。\BibleRef{格前 11:32}因为那法庭是可怕的，审判席是恐怖的，账目令人战栗，一条火河翻滚\OriginalTerm{被拖}{\GreekText{ἕλκεται}}。\BibleQuote{兄弟不能赎回，人岂能赎回？}\BibleRef{咏 49:8} 那么要记起福音中所说的，天使们来回奔跑，洞房关闭，灯盏熄灭，权能将人拖向火炉。请考虑这一点，如果我们当中任何一个人的隐秘行为今天被公开在教会面前，他除了求死、求大地裂开口吞他，以免有这么多见证者看到他的邪恶之外，还能做什么呢？那么，当在全世界面前，一切都被公开在如此明亮广阔的剧场中，认识和不认识我们的人都看穿一切时，我们将有何感受？但哀哉！我不得不以什么来恐吓你们！以人的评价！当我本当使用对天主的敬畏和祂的定罪时。请问，那时我们被捆绑、咬牙切齿、被带到外面的黑暗中，将会怎样？或者，更确切地说——这是最可怕的想法——当我们冒犯\OriginalTerm{冒犯}{\GreekText{προσκρούσωμεν}}天主时，我们将做什么？因为如果有人有理智，当看不见天主时，他已经忍受了一个地狱。但既然这不痛苦，因此火被用来恐吓。因为我们不当在受罚时感到刺痛，而应在犯罪时。所以要听保禄为那些他本不会受罚的罪恶哀哭痛悔，因为他说：\BibleQuote{我不配，}他说，\BibleQuote{称为宗徒，因为我迫害了教会。}\BibleRef{格前 15:9} 也要听达味，当他从惩罚中被释放后，仍认为他得罪了天主，便呼求报复落到自己身上，说：\BibleQuote{请祢的手加在我和我父的家上。}\BibleRef{撒下 24:17} 因为得罪天主比受罚更令人痛苦。但如今我们如此可悲地倾向，如果不怕地狱，我们甚至不愿乐意做任何善事。因此即使不为别的，至少为此我们该当地狱，因为我们怕地狱超过怕基督（部分手抄本作：天主）。但蒙福的保禄却不是这样，恰恰相反。但是既然我们感觉不同，因此我们被判下地狱：因为如果我们爱基督当如何爱祂，我们就该知道，得罪我们所爱的祂比地狱更痛苦。但既然我们不爱祂，我们就不知道祂惩罚的严重。这正是我最哀哭痛惜的！然而天主为了被我们爱，还有什么没有做？还有什么没有设计？还有什么忽略？我们侮辱了祂，而祂在任何事上都没有亏待我们，反而以无数不可言喻的恩惠恩待了我们。我们背离了祂，当祂用各种方式召唤、吸引我们时，但即使这样祂也没有惩罚我们，反而亲自跑向我们，在我们逃跑时拉住我们，而我们却摆脱祂，跳向魔鬼。即使如此祂也没有远离，而是派遣了无数使者再次召唤我们，先知、天使、圣祖：我们不仅不接受使命，反而侮辱了来的人。但即使为此，祂也没有将我们从口中吐出，而是像那些被轻视却极其热切的情人一样，四处恳求众人，天地、耶肋米亚、米该亚，并不是为了压制我们，而是要为祂自己的行为辩护\BibleRef{米 6:1}：并且祂自己也与先知一同到那些背离祂的人那里，准备接受审查，屈尊俯就商议，把那些对一切呼吁充耳不闻的人吸引到与祂辩论之中。因为祂说：\BibleQuote{我的百姓，我做了什么对你们？我在什么事上使你们厌倦？答覆我。}\BibleRef{米 6:3} 在这一切之后，我们杀了先知，用石头砸死他们，对他们行了无数残酷的虐待。然后呢？祂不再派遣先知、天使、圣祖，而是派遣了圣子自己。祂来后也被杀了，但即使如此祂也没有熄灭祂的爱，反而更加炽烈，即使在祂的独生子被杀害之后，祂仍然继续恳求我们、劝勉我们，做一切事来使我们归向祂。保禄大声呼喊说：\BibleQuote{所以我们代基督作大使，好像天主藉着我们劝你们：你们要与天主和好。}\BibleRef{格后 5:20} 然而这一切都没有使我们和好。但即使那时祂也没有离开我们，而是继续既以地狱恐吓，又应许王国，以便这样吸引我们归向祂。但我们仍然麻木不仁。有什么比这兽性更糟糕？因为如果一个人做了这些事，我们岂不是多次让自己成为他的奴隶？但当天主这样做时，我们却转离祂！哦，何等麻木！何等无情！我们这些不断活在罪恶和邪恶中的人，如果我们偶然做了些许善事，就像无情的家仆，我们何等吝啬地索取报酬，如果所做的有任何将来的报答，我们多么计较报答。然而如果你做善事毫无得赏的希望，报答反而更大。说这一切并仔细计算，是更适合雇工的语言，而非甘心的家仆。因为我们一切当作在基督的缘故，不为赏报，而为祂。因为这也是祂以地狱恐吓、应许王国的原因，好使我们爱祂。那么让我们按所当爱的来爱祂。因为这就是伟大的赏报，这就是王权与喜乐，这就是享受、光荣与尊荣，这就是光明，就是伟大的幸福，是言语（或理智）无法陈明、人心无法想象的。然而我确实不知道我为何在说话中这样远行，竟来劝勉那些甚至不肯为基督的缘故轻视此世权势和光荣的人，要轻视天国。然而那些伟大高尚的人甚至达到了这种爱的程度。例如，听听伯多禄如何热烈爱祂，将祂置于灵魂、生命和万有之上。而且当他否认祂时，他悲伤的不是惩罚，而是他否认了他所渴望的那一位，这对他比任何惩罚更痛苦。这一切都是在圣神恩宠赐下之前显出的。他坚持追问：\BibleQuote{祢往哪里去？}\BibleRef{若 13:36} 在此前又说：\BibleQuote{我们去投奔谁？}\BibleRef{若 6:67} 又说：\BibleQuote{我愿跟随祢不论祢往哪里去。}\BibleRef{路 22:33} 因此祂是他们的万有，他们既不看重天堂，也不看重天国，与所渴望的祂相比较。因为祢就是我的一切，他意指。你为何惊异伯多禄如此心志？现在听先知说什么：\BibleQuote{在天上除祢以外，为我还能有谁？在地上除祢以外，为我一无所喜。}\BibleRef{咏 72:25} 他所说的差不多是：既不求上面的事，也不求下面的事，只求祢自己。这是热忱，这是爱。如果我们能如此爱，则不但现世的事物，就连将来的事物，我们也会视为与那爱慕的魅力无法相比，并且在此世我们就能享受天国，以祂的爱为乐。这如何可能？有人会说。让我们反省我们多少次在无数的善事之后侮辱祂，而祂仍站着召唤我们，多少次我们从祂旁边跑过，但祂仍不忽视我们，反而跑向我们，吸引我们，抓住我们归于祂自己。因为如果我们思考这些和类似的事，我们就能燃起这种渴望。因为如果是一个普通人如此爱，而是一个君王如此被爱，他岂不会因爱的伟大而升起敬意吗？必定会。但当情况颠倒，祂的美（圣\Quote{那美}）是不可言喻的，爱我们的那位的荣耀和富有也是如此，而我们的卑贱如此之大，我们当然应受最大的惩罚，我们如此卑贱和被遗弃，竟被如此伟大奇妙的一位以如此出奇的爱对待，却竟然对祂的爱放肆妄为？祂不需要我们的任何东西，却至今仍不停爱我们。我们非常需要属于祂的东西，却仍不依恋祂的爱，反而将金钱、人的友谊、身体的安逸、权势、名声看得比祂更重，而祂却珍视我们超过一切。因为祂有一个儿子，真是（字面意为\Quote{真生}）独生子，为了我们祂连自己的儿子都没有保留。但我们却把许多事物看得比祂更重。那么，地狱和折磨岂不是有充分的理由，即使是两倍、三倍或多倍于现在的？因为如果连撒殚的命令我们都看得比基督的法律更重，不顾自己的救恩，选择邪恶的行为，而不选择那为我们承受了一切的祂，我们还有什么可为自己辩护的？这些事配得什么宽恕？有什么借口？一个也没有（五个手抄本作：\OriginalTerm{连一个也没有}{\GreekText{οὐδὲ} \GreekText{μιἵς}}）。那么让我们在此之后停止狂奔，重新清醒；数算这一切，以我们的行为（因为仅言语不足以）向祂献上光荣，好使我们也能享受由祂而来的光荣，愿我们众人都因我们的主耶稣基督的恩宠和对人的爱而获得这光荣，借着祂，同祂，并偕同圣神，愿光荣归于圣父，至于无穷世之世。阿们。\par

\SourceRecord{Translated by J. Walker, J. Sheppard and H. Browne, and revised by George B. Stevens.；Nicene and Post-Nicene Fathers, First Series；Vol. 11.；Edited by Philip Schaff.；Buffalo, NY: Christian Literature Publishing Co.,；1889.；<http://www.newadvent.org/fathers/210205.htm>.}

% structure: p0001 source=h1 target=section reason=page-title

\section{罗马书讲道集第六篇}

\BibleRef{罗 2:17,18}\par

\Quote{\Emph{看}，\Emph{你称为犹太人，又依靠法律，且以天主夸耀，并知道祂的旨意，且鉴别更美的事，从法律受了教训。}}\par

说完外邦人若遵行法律，并不缺少任何与救恩有关之事，并作了那奇妙的比较之后，他接着陈述犹太人的荣耀，他们因这些荣耀而轻视外邦人：首先就是这名称本身，它曾大有威严，正如现今基督徒之名一样。因为即在那时，这个称呼所带来的区别是巨大的。因此他从这一点开始，且看他如何贬低它。因为他并不说：\Quote{看，你是犹太人}，而是说\Quote{称为}如此，且\Quote{以天主夸耀}；即，作为被祂所爱，并受尊荣超乎众人。这里他似乎温和地讥讽他们的无理，以及对荣誉的极大狂热，因为他们滥用这个恩赐，不是为了自己的救恩，而是去高抬自己，轻视他人。\Quote{知道祂的旨意，并鉴别更美的事。}这确实是一种不利，若没有实行；然而它似乎仍是一种有利，所以他准确地陈述。因为他不是说：你实行，而是知道；且鉴别，而不是追随和实践。\par

第19节。\Quote{且自信你自己。}\BibleRef{罗 2:19}\par

在此他又不说你是盲目的向导，而是说你自信，所以你夸耀，他说。犹太人的荒谬如此之大。因此他也几乎重复了他们夸口所用的原话。例如看他们在福音中所说的。\Quote{你完全是（\OriginalTerm{完全}{\GreekText{ὅλος}}，四份手稿作\OriginalTerm{全然}{\GreekText{ὅλως}}）生在罪中，你还教训我们吗？}\BibleRef{若 9:34}他们对所有人极其鄙视，为了说服他们，保禄不断颂扬他们并贬低其他人，以便能更抓住他们，并使他的指控更有分量。因此他接着添加类似的事，并通过不同的叙述方式使他们更为显著。因为他说：\Quote{你自信你自己是盲目的向导。}\par

第20节。\Quote{愚昧者的训导者，婴孩的教师，拥有知识和真理的模型，就是在法律内的。}\BibleRef{罗 2:20}\par

在此他又不说在良心、行为与善行中，而是说\Quote{在法律内}；说了这些之后，他在这里也做了他对付外邦人所做的事。因为正如在那里他说：\Quote{因为你在何事上判断别人，就定了你自己的罪}，他在这里也这样说。\par

第21节。\Quote{你既教导别人，难道不教导自己吗？}\BibleRef{罗 2:21}\par

但那里他以更尖锐的方式构造他的言辞，这里则更温和。因为他不是说：因此你应受更大惩罚，因为受托如此大事却未善用任何一件，而是以提问的方式将话语引向他们，使他们自省（\OriginalTerm{使他们自省}{\GreekText{ἐντρέπων}}），说：\Quote{你既教导别人，难道不教导自己吗？}在这里我要你们看\Person{保禄}在另一情形中的智慧。因为他列出了犹太人的好处，这些好处并非由于他们的热忱，而是来自上主的恩赐，并且他显示这些好处若不实行不仅无益，反而带来更大的惩罚。因为被称为犹太人并非他们的善行，领受法律也不是，其他他刚才列举的也不是，而是来自上主的恩宠。在开头他曾说，听法律若无行为便无价值（他说：\Quote{不是听法律的，就在天主面前为义人}），但现在他进一步显示，不仅听，连法律的教导本身也不能庇护教导者，除非他实行所说的话；并且不仅不能庇护，反而将更严厉地惩罚他。他也恰当地使用了措辞，因为他不是说：你领受了法律，而是说\Quote{你依靠法律}。因为犹太人不必辛苦寻求当做之事，而是轻易得到指引德行的法律。因为如果连外邦人也有自然理性（正因此他们优于犹太人，因他们未听却行法律），但其他人仍有更大的便利。但若你说：我不但是听，而且是教导者，这本身正是你惩罚的加重。因为他们以此自豪，所以他从这一点上显明他们的可笑。但当他说：\Quote{盲目的向导，愚昧者的训导者，婴孩的教师}，他是在说他们自己夸夸其谈的话。因为他们极其恶劣地对待归化者，而这些都是他们称呼归化者的名字。这就是为什么他详细讲述那些被认为是他们赞美的话，深知所说的话为更大的控诉提供了基础；\Quote{拥有知识和真理的模型，就是在法律内的。}好像一个人拥有国王的画像，却什么都不描画，而那些未被托付画像的人却凭记忆完美模仿了。然后，在提到天主给他们的好处之后，他告诉他们他们的失败，引用了先知对他们的指控。\Quote{你既教导别人，难道不教导自己吗？你宣讲不可偷盗，你偷盗吗？你说不可奸淫，你奸淫吗？你憎恶偶像，你犯盗窃圣物之罪吗？}因为他们被严格禁止接触偶像上的任何珍宝（根据手稿：武加大译本作\Quote{在偶像庙中}），因着玷污。但是，他说，贪婪的暴政已经说服你（4份手稿和边注作\Quote{我们}）也将这条法律践踏在脚下。然后他随后带来了更为严重的指控，说：\par

第23节：\BibleQuote{你既然以法律夸耀，自己却违犯法律而侮辱天主吗？}\BibleRef{罗 2:23}\par

他提出了两项指控，或者更恰当地说，三项。第一，他们侮辱了天主；第二，他们侮辱了那使他们受尊荣的；第三，他们侮辱了那位尊荣他们的天主——这已是麻木不仁的极致。然后，为免显得是他自己任意指控他们，他引用了先知作为他们的指控者，此处简要而概括地像是一个总结，但后来更加详细；这里引用了依撒意亚，之后是达味，当他已指出不止一个可以责备的理由时。因为他的意思是：要说这些话并非出于我自己的责备，请听依撒意亚所说的。\par

第24节：\BibleQuote{因为天主的名在你们中间因着你们被外邦人亵渎。}\BibleRef{罗 2:24} \BibleRef{依 52:5}\par

看，又是一项双重的指控。因为他们不仅自己放肆，甚至还引诱别人也这样作。那么，你们教导别人，却不教导自己，这有什么用处呢？然而，以上他只是这样说，但在这里他甚至将这事颠倒过来。因为你们不仅自己不教导别人应作的事，甚至更糟——你们不仅不教导法律的事，反而教导相反的事，即亵渎天主，这与法律是相反的。但是，有人会说：割礼是一件大事。是的，我也承认，但什么时候呢？是当它具备内在的割礼时（所有手抄本都作：\Quote{那时，当}）。请注意他的判断：他是如此适时地引出了关于割礼的讨论。因为他并没有立即以割礼开始，因为人们对它的自负心很重。但在显示出他们已在更重大的事上犯罪，并应对亵渎天主负责之后，于是他便取得了读者对他们的不利判断，并剥夺了他们的优越感，然后才引出关于割礼的讨论，确信没有人会再为它辩护，于是他说：\par

第25节：\BibleQuote{割礼确有益处，如果你遵行法律。}\BibleRef{罗 2:25}\par

然而，若不是这样，人可能会拒绝它并说：割礼是什么？难道割礼本身是善行吗？难道它是正确选择的体现吗？因为它是在不成熟的年龄施行的，而且在旷野中的那些人长期未受割礼。从许多其他角度来看，人们也可以视它为不必要。然而，他并非从这个角度拒绝它，而是从最恰当的理由出发，即从亚巴郎的实例。因为这是最卓越的胜利——即从他们尊崇它的原因出发，来证明它的微不足道。他本可以说，连先知也称犹太人为未受割礼的。但这并非对割礼的贬低，而是对那些不善持守它的人。因为他的目的是要表明，即使在最完美的生活中，它毫无力量。这正是他接下来要证明的。在这里，他并没有抬出圣祖，而是先在其他基础上推翻了它，留待后来再说关于信德的事，当他说到：\Quote{那么，它是如何算为的}给亚巴郎？\Quote{是在受割礼的时候，还是在未受割礼的时候？}因为只要它是在与外邦人和未受割礼的人争辩时，他不愿提及此事，以免给他们过于难堪。但当它反对信德时，他便更完全地摆脱了它，投入战斗。因此，到目前为止，争辩的对象是未受割礼，这就是为什么他说话时语气克制，并说：\par

\Quote{割礼确有益处，如果你遵行法律；但如果你违犯法律，你的割礼就成了未受割礼。}因为这里他谈到两种未受割礼和两种割礼，正如两种法律一样。有一种自然的法律和一种成文的法律。但还有一种是介于两者之间的，即由行为所体现的法律。请看他是如何指明这三种，并将它们呈现在你面前。\par

\Quote{因为当外邦人，}他说，\Quote{没有法律。}什么法律？成文的法律。\Quote{顺着本性去遵行法律的事。}什么法律？由行为所体现的法律。\Quote{他们没有法律。}什么法律？成文的法律。\Quote{他们自己就是自己的法律。}怎么这么说呢？是通过运用自然的法律。\Quote{他们显示法律的功用。}什么法律？由行为所体现的法律。因为成文的法律是外在的，而这内在的是自然的法律，另一种则是行为上的。成文的法律由文字宣告；自然的法律由本性宣告；行为上的法律由行为宣告。这第三种是必要的，前面两种——自然的法律和成文的法律——都是为了它而存在。如果这第三种不存在，前两者就毫无益处，甚至带来极大的害处。为了在自然法律的例子中表明这一点，他说：\Quote{因为你判断别人，就定了你自己的罪。}至于成文的法律，则是这样：\Quote{你教导人不可偷盗，自己却偷盗。}同样，有两种未受割礼：一种是本性的未受割礼，另一种是由行为而来的未受割礼；还有一种是肉身上的割礼，另一种是由意志而来的割礼。我的意思是，一个人在第8天受了割礼，这是肉身的割礼；一个人遵行了法律所命令的一切，这是心灵的割礼，正是圣保禄所首要要求的，甚至法律也是如此要求的。请看，他是如何一面在言辞上承认它，一面在实际上废除了它。因为他没有说割礼是多余的，割礼毫无益处，毫无用处。而是说：\Quote{割礼确有益处，如果你遵行法律。}\BibleRef{申 10:16}他如此赞成它，说：我承认并不否认割礼是光荣的。但什么时候呢？当有法律的遵守与之相伴时。\par

\Quote{但是你若是犯法律的，你的割损反倒成为未割损。} 他并没有说它不再有益，免得他好像侮辱了割损。但他在剥去犹太人后，继续打击他。这不再是侮辱割损，而是侮辱那因懈怠而失去其好处的人。正如那些身居尊位、后来被判定犯有重大罪行的人，法官剥夺他们职位的荣誉，然后再惩罚他们；同样，保禄也这样做了。因为说了：若是犯法律的，你的\Quote{割损就成为未割损}，并表明他是未受割损的，然后他就毫无顾忌地谴责他。\par

26节。\Quote{因此，如果未割损遵守了法律的义德，他的未割损岂不成为割损吗？}\BibleRef{罗 2:26}\par

请看他是如何行事的。他没有说未割损胜过割损（因为这在当时听的人听来非常刺耳），而是说未割损成了割损。接着他探究什么是割损，什么是未割损，他说割损是善行，未割损是恶行。他先将行善的未割损者纳入割损，而将生活腐败的割损者驱逐到未割损中，然后这样把未割损者置于优先。他并没有说对未割损者，而是针对事情本身，如此说道：\Quote{他的未割损岂不成为割损吗？} 他没有说\Quote{算为}，而说\Quote{转为}，这更富表现力。正如上文他也没有说你的割损算为未割损，而是成了未割损。\par

27节。\Quote{并且那本性的未割损，岂不审判吗？}\BibleRef{罗 2:27}\par

你看，他区分了两种未割损：一种来自本性，一种来自意志。然而，这里他指的是来自本性的，但并没有停留于此，而是继续说：\Quote{如果它完成了法律，就要审判你，你这凭着文字和割损违犯法律的。} 请看他的精妙判断。他没有说那来自本性的未割损会审判割损，而是在得胜之处，他引入未割损；但在失败之处，他并沒有揭露割损是失败的，而是揭露拥有割损的犹太人自己；这样措辞是为了避免冒犯听者。他也没有说\Quote{你这有法律和割损的}，而是更温和地说\Quote{你这凭着文字和割损违犯法律的}。也就是说，这样的未割损甚至替割损辩护，因为它是受屈的，前来帮助法律，因为它受到了侮辱，并取得了显著的胜利。因为当时犹太人不是被犹太人审判，而是被未割损者审判，这决定了胜利；正如他所说：\Quote{尼尼微人要起来审判这一世代，并要定它的罪。}\BibleRef{玛 12:41} 因此，他不是侮辱法律（因为他非常尊重法律），而是侮辱那使法律蒙羞的人。接着，在清楚确立了这些基础之后，他自信地继续定义真正的犹太人是什么；并表明不是犹太人，也不是割损，而是他所排斥的既非犹太人也非未割损的人。他表面上似乎为割损辩护，但实际上却废除了关于割损的意见，通过得出的结论争取人们的同意。因为他不仅表明犹太人与未割损者之间没有差别，而且表明如果未割损者谨慎自守，甚至还有优势，并且他才是真正的犹太人；于是他说道：\par

12节。\Quote{因为外表上的犹太人，并不是犹太人。}\BibleRef{罗 2:12}\par

在这里，他攻击他们为了一切都出于表现。\par

29节。\Quote{唯有内在的犹太人才是犹太人；割损是内心的割损，在于圣神，不在于文字。}\BibleRef{罗 2:29}\par

他说这话，便排除了所有属肉体的事物。因为割损是外在的，安息日、祭献和洁净也是如此；他用一句话暗示了这一切，当他说：\Quote{因为外表上的犹太人，并不是犹太人。} 但既然割损受到重视，甚至安息日也让步于它\BibleRef{若 7:22}，他特别针对割损是有充分理由的。当他说了\Quote{在于圣神}之后，他便为教会的转变铺平了道路，并引入了信德。因为信德也在于心与圣神，并且从天主那里得赞许。他怎么没有证明行事正直的外邦人不逊于守法的犹太人，反而证明行事正直的外邦人比违法的犹太人更好呢？这是为了使胜利不容置疑。因为一旦这一点达成一致，肉身的割损就必然被废弃，善行的需要就在各处得到彰显。因为希腊人没有这些却得救，而犹太人有了这些却仍受惩罚，犹太教便毫无功效。他在这里所说的\Quote{希腊人}，不是指拜偶像的希腊人，而是指虔敬而有德、不受任何法律规条束缚的人。\par

第三章第1节。\Quote{那么，犹太人有什么长处呢？}\BibleRef{罗 3:1}\par

既然他藉着说\Quote{他不是外在的犹太人，而是内在的犹太人}将一切撇下：听闻、教导、犹太人的名号、割礼以及所有其他细节，接下来他看到有一个异议产生，并对此进行辩护。这个异议是什么呢？他的意思是，如果这些东西没有用处，那么那民族被拣选以及割礼被赐予的理由是什么？那么他如何解决这问题？用他先前同样的方式：因为正如在那里，他谈的不是他们的赞美，而是天主的恩惠；也不是他们的善行（因为被称为犹太人、认识祂的旨意、并赞同更卓越的事，并不是他们自己的善行，而是来自天主的恩宠：先知也这样斥责他们说：\BibleQuote{祂没有这样对待任何其他民族，也没有将祂的法律显示给他们}\BibleRef{咏 147:20}；梅瑟又说：\BibleQuote{请你们现在问问，是否有这样的事？}\BibleRef{申 4:32}他说，\BibleQuote{有哪一个民族听见天主从火中说话的声音，还能活着？}\BibleRef{申 4:32}），这里他也这样做。因为正如他在论及割礼时，并没有说割礼若无善行便无价值，而是说割礼与善行并存时才有价值，指出了同样的事，但语气更缓和。再者，他并没有说：如果你是个犯法的人，你虽是受割礼的，也毫无益处，而是说\BibleQuote{你的割礼反成了未受割礼}；之后他又说，\BibleQuote{未受割礼的}将要\BibleQuote{审判}，不是审判割礼，而是\BibleQuote{审判你这犯法律的人}，所以宽容法律的事，而打击人本身。这里他也这样做。因为他在提出这个异议并说\Quote{那么犹太人有什么长处？}之后，他并没有说\Quote{没有}，而是赞同这个说法，又通过后续驳斥它，并表明他们甚至因这优越而受到惩罚。我将在陈述异议后告诉你们他是如何做的。他说：\BibleQuote{那么犹太人有什么长处？割礼有什么益处？}\par

第二节：\BibleQuote{各方面都有很多。首先，因为他们被信托了天主的圣言。}\BibleRef{罗 3:2}\par

你看到吗，正如我上面所说的，他所数算的处处不是他们的善行，而是天主的恩惠。而\Quote{被信托}这个希腊词（\OriginalTerm{被信托}{\GreekText{ἐπιστεύθησαν}}）是什么意思？它意味着法律被交在他们手中，因为祂如此看重他们，以致将从上而来的圣言托付给他们。我知道有些人将\Quote{被信托}不是理解为关于犹太人，而是关于圣言，意思是法律被相信了。但上下文不允许这样理解。因为首先，他这样说是为了控告他们，并表明尽管他们享受了许多从上面来的祝福，却表现出极大的忘恩负义。然后，上下文也清楚表明这一点。因为他接着说：\BibleQuote{即使有些人不相信，那又如何？}\par

第三节：\BibleQuote{即使有些人不相信，那又如何？}\BibleRef{罗 3:3}\par

而接下来的部分也表明了同样的一点。因为他再次添加并接着说：\BibleQuote{难道他们的不信会使天主的信实失效吗？}\BibleRef{罗 3:3}\par

第四节：\BibleQuote{绝对不可。}\BibleRef{罗 3:4} 那么，\OriginalTerm{被信托}{\GreekText{ἐπιστεύθησαν}}这个词宣告了天主的恩赐。\par

我希望你也在此注意他的判断。因为他并不亲自控告他们，而是如同通过反驳的方式，好像说：\Quote{但也许你会说：\InnerQuote{那么这割礼有什么用，既然他们没有按应有的方式使用它，既然他们受托于法律却对托付不忠？}} 到目前为止他并非严厉的控告者，而是仿佛要替天主洗脱对祂的抱怨，通过这种方式将全部控告转回到他们自己身上。因为，他会说，为什么你们抱怨他们不信？这与天主何干？因为对于祂的恩惠，受惠者的忘恩负义能推翻它吗？或者使这尊荣不再是尊荣？因为这就是这些话的意思：\Quote{他们的不忠信会使天主的忠信失效吗？} \Quote{断乎不是！} 就好像有人说：我如此尊敬某人。如果他未曾接受这尊敬，这并不给我任何控告的理由，也不损害我的善意，而是显示了他的无情。但\Person{保禄}不仅仅说这些，而是说了更重要的。那就是，他们的不信不仅没有给天主留下抱怨的余地，反而更显示祂的尊荣和爱人之心，因为祂被看见将尊荣赐给了一个会羞辱祂的人。看，他如何通过他们所夸耀的使他们显出过失；因为天主对待他们的尊荣是如此之大，甚至当祂预见到后果时，也没有收回对他们的善意！然而他们却将所受的尊荣当作侮辱那尊荣者的手段！接着，既然他说了\Quote{即便有些人不信}（而显然他们全都不信），为了避免在这里像历史所允许的那样说话，以免显得像敌人一样严厉控告他们，他把实际发生的事放在推理和三个段论的方法中，如下所说：\Quote{但愿天主是真实的，众人都是虚谎的。} 他所说的是这样意思：我不是说有些人不信，但如果你愿意，假设所有人都不信，这样略过实际发生的事，以迎合反对者，免得他显得傲慢或可疑。好了，他说，这样天主更显为义。'义'这个词是什么意思？意思是，如果可能有一场审判和审查，审查天主为犹太人所作的事，以及他们为天主所作的事，那么胜利将属于天主，一切正义在祂那边。在从前面所说的清楚显示这一点后，他接着也引入先知赞同这些事，说：\BibleQuote{叫祢在祢的言语上显为义，在祢被审判时得胜。} \BibleRef{咏 50:4} 祂那边已尽了一切，但他们对这也没有好转。然后他接着又提出另一个反对意见，说：\par

第5节。\BibleQuote{但若我们的不义显明了天主的义，我们该说什么呢？难道降怒的天主不义吗？——我按人的说法说。} \BibleRef{罗 3:5}\par

第6节。\BibleQuote{断乎不是！} \BibleRef{罗 3:6}\par

他用另一个难题来解决一个难题。然而因为这话不清晰，我们必须更清楚地说明。那么他是什么意思呢？天主尊荣了犹太人：他们却侮辱了祂。这给了祂胜利，并显示了祂的爱人之心是何等伟大，因为祂尊荣了他们，尽管他们是如此。那么，他的意思是，我们既侮辱了祂、亏欠了祂，天主恰恰因这件事得胜，祂的义就显明了。那么（有人可以说）为什么我要受惩罚，我因侮辱祂而成了祂胜利的原因？现在他如何应对？正如我所说，又是另一个荒谬。因为如果你是胜利的原因，而又受惩罚，那便是不义。但如果祂并非不义，而你却受惩罚，那么你就不再是胜利的原因。且注意他宗徒的虔敬（或谨慎：\OriginalTerm{虔诚}{\GreekText{εὐλάβεια}}）；因为说了\Quote{难道降怒的天主不义吗？}之后，他又加上\Quote{我按人的说法说。} 仿佛他是说，任何人都可能以人的方式来推理。因为在我们看来是正义的事，天主公正的审判远超过它，并有某些其他不可言说的理由。接着，因为还不清晰，他又重复同样的话：\par

第7节。\BibleQuote{因为若天主的真理因我的谎言越发显耀，而归荣于祂，为什么我还要被判定为罪人呢？} \BibleRef{罗 3:7}\par

因为，他的意思是，若天主因你们的悖逆行为被显为爱人、正义和良善，你们不仅不应受罚，甚至应当受善待。但若如此，就会产生那个流传甚广的荒谬说法：善出于恶，恶是善的原因；二者必居其一：要么祂在惩罚时显为不义，要么若祂不惩罚，那么祂的胜利来自我们的恶行。这两者都极其荒谬。他自己也意在表明这点，就引进了外邦人（即异教徒）作为这些意见的创始者，认为只要指出说这些话的人的身份就足以反驳他所提及的。因为那时他们常讥笑我们说：\Quote{让我们行恶以成善吧。} 这就是为什么他在以下话语中清楚陈述了它。\par

第8节。\BibleQuote{为什么不？——正如我们被人毁谤，有人肯定我们说过的——让我们行恶以成善？对这样的人，他们的定罪是公道的。} \BibleRef{罗 3:8}\par

因为保禄曾说：\Quote{罪恶在那里增多，恩宠在那里也格外丰富}（\BibleRef{罗 5:20}），他们却嘲笑他，曲解他的话，说：我们必须坚持恶习，好能得着善。但保禄并非如此说；然而为纠正这种观念，他才说：\Quote{那么，我们该怎么办？要留在罪恶中，好叫恩宠增多吗？绝对不可！}（同上\BibleRef{罗 6:1-2}）。他说，我指的是过去的时候，并非要我们把此当作常行。为使他们远离这种疑虑，他说此后甚至是不可能的。因为他说：\Quote{我们这些死于罪恶的人，怎能再活在其中呢？}然后他毫无困难地\OriginalTerm{谴责}{\GreekText{κατέδραμεν}}了希腊人。因为他们的生活极为放荡。但关于犹太人，即使他们的生活似乎漫不经心，他们仍有很大借口在律法和割礼中掩盖这些事，以及天主曾与他们交谈的事实，并且他们是众人的导师。这就是为什么他甚至剥去这些，表明因这些他们更受惩罚，这也是他在这里得出讨论的结论。因为他说，如果他们不因这样做而受惩罚，那么那句亵渎的话——\Quote{我们作恶好叫善来到}——必然流行。但如果这是不虔敬的，而说这话的人将受惩罚（因为他藉着说：\Quote{他们的惩罚是公义的}来表明），显然他们受惩罚。因为如果说话的人应受报应，那么行动的人更应受；但如果应受报应，那就是因为他们犯了罪。因为惩罚他们的不是人，免得有人怀疑判决，而是天主，祂行事皆公义。但如果他们公义地受惩罚，那么嘲笑我们的人所说的话就是不正当的。因为天主曾做并一直做一切，使我们的行为在各方面都光明正直。\par

那么，我们不可懈怠；因为这样我们也能将希腊人从错误中挽回。但当我们口称爱智，行为却不端时，我们有何面目见他们？用何嘴唇谈论教义？因为他会对我们每个人说：你在小事上失败，怎能妄图教导我更大的事？你尚未学会贪财是恶，怎能通晓天上的事？但你明知它是恶？那么，罪责更大，因为你明知故犯。我何必提希腊人？因为连我们的法律也不允许我们在生活堕落时如此大胆发言。因为经上说：\Quote{天主对罪人说：你有什么权利宣读我的律例？}（\BibleRef{咏 49:16}）曾有一次犹太人被掳，波斯人催促他们，叫他们唱那些神圣的歌曲，他们说：\Quote{我们怎能在外邦的土地上唱上主的歌曲？}（\BibleRef{咏 136:4}）如果在异乡唱天主的圣言是不合法的，那么疏远的灵魂更不能唱。因为无怜悯的灵魂是疏远的。如果律法使那些被掳、在异乡成为奴隶的人静坐不语，那么那些被罪恶奴役、在\OriginalTerm{外邦社群}{\GreekText{πολιτεί}\& 139}中的人更应静坐不语。\par

\SourceRecord{Translated by J. Walker, J. Sheppard and H. Browne, and revised by George B. Stevens.；Nicene and Post-Nicene Fathers, First Series；Vol. 11.；Edited by Philip Schaff.；Buffalo, NY: Christian Literature Publishing Co.,；1889.；<http://www.newadvent.org/fathers/210206.htm>.}

% structure: p0001 source=h1 target=section reason=page-title

\section{罗马书讲道第七篇}

\BibleRef{罗 3:9-18}\par

\BibleQuote{\Emph{那么，我们比他们更有何长处呢？}\Emph{比他们？}\Emph{因为我们已指控了犹太人和外邦人，他们都处于罪恶之下。正如经上所载：没有义人，连一个也没有；没有明智的，没有寻求天主的；他们都离弃了正道，一同变得无用；没有行善的，连一个也没有。他们的喉咙是敞开的坟墓；他们用舌头行欺骗；虺蛇的毒汁在他们的嘴唇之下；他们的口充满咒骂和苦毒；他们的脚急于流血；毁灭和苦难在他们的行径上；和平的道路他们不认识；他们眼中没有对天主的敬畏。}}\par

他曾指控了外邦人，又指控了犹太人；按顺序接下来应提及因信德而来的正义。因为如果自然律无效，成文法律也无益，反而压制了那些不善用它们的人，并更加证明他们应受更大的惩罚，那么之后因恩宠而来的救恩就是必需的。那么，保禄啊，请讲论它，并显扬它吧。但他还不敢，因为他顾及犹太人的激烈，于是又转向对他们的指控；他首先引用了达味作指控者，详细谈论了依撒意亚简略提及的相同的事，从而为他们提供了一个强有力的约束，使他们不致暴跳，也不致当信德的事向他们敞开时，他们此后会逃避；因为他们被先知们的指控预先牢牢地压制住了。因为先知指出了三种过分：他说他们全都作恶，并且不是善恶不分，而是只追随邪恶，而且热切地追随。为了使犹太人不能说：\Quote{那么，如果这些话是对别人说的呢？}他接着说道：\par

第19节。\BibleQuote{如今我们知道，凡法律所说的，都是对在法律之下的人说的。}\BibleRef{罗 3:19}\par

这就是为什么在依撒意亚之后，他引用了达味，因为依撒意亚明确针对他们，好表明这些话也属于同一个主题。他意思是，一个为纠正你们而派遣的先知，有什么必要去指控别的民族呢？因为法律也不是赐给别人的，而是赐给你们的。他为什么不说\Quote{我们知道，凡先知所说的}，而说\Quote{凡法律所说的}呢？因为保禄习惯将整个旧约称为法律。他在另一处也说：\BibleQuote{你们没有听见法律说，亚伯拉罕有两个儿子吗？}\BibleRef{迦 4:21} 这里他把圣咏也称为法律，当他说：\BibleQuote{我们知道，凡法律所说的，都是对在法律之下的人说的}时。接着他表明，他说这些话也不仅仅是为了指控，而是为了再次为信德铺路。旧约与新约的关系如此密切，因为连指控和责备也完全是为了这个目的：使信德的门为听的人明亮地敞开。因为犹太人最大的祸害是他们自高自大（他曾在后面提到：\BibleQuote{他们怎样因为不认识天主的正义，而企图建立自己的正义，便没有服从天主的正义}\BibleRef{罗 10:3}），因此法律和先知预先驳斥了他们，打倒了他们的高傲，压制了他们的自大，好使他们因思考自己的罪过，倒空所有的愚妄，看见自己处于最后的危险之中，从而热切地奔向那位赐予他们罪赦的主，并藉信德接受恩宠。这就是圣保禄在这里暗示的，当他说：\par

\BibleQuote{如今我们知道，凡法律所说的，都是对在法律之下的人说的，使众口闭住，使普世都成为有罪的于天主前。}\par

在这里，他表明他们因行为而缺乏宣讲的胆量，只是空谈和厚颜无耻。这就是他用如此字面的表达，说\BibleQuote{使众口闭住}，指出他们言语的厚颜和几乎无法控制的浮夸，以及他们的舌头现在被严格勒住。因为如同不可阻挡的洪流，它就这样被冲走。但先知阻止了它。当保禄说\BibleQuote{使众口闭住}时，他的意思不是说他们犯罪的原因是为了堵住他们的口，而是他们受责备的原因是为了使他们不致无知地犯下这个罪。\BibleQuote{普世都成为有罪的于天主前。}他说的不是犹太人，而是全人类。因为\BibleQuote{使众口闭住}这句话是暗示他们，尽管他没有明确说明，以免言语过于严厉。但\BibleQuote{普世都成为有罪的于天主前}这句话，同时指的是犹太人和希腊人。这对于破除他们的愚妄并非小事。因为即使在这里，他们也不比外邦人有什么优势，就救恩而言，都被同样地放弃。因为一个正当地被称为有罪的人，是不能借助任何借口自救，而需要别人帮助的人：我们所有的人当时都处于这种状态，因为我们失去了有关救恩的事物。\par

第20节。\BibleQuote{因为法律只是叫人认识罪。}\BibleRef{罗 3:20}\par

他再次对法律发难，然而仍带着宽容（因为他所说的并非指摘法律，而是指摘犹太人的懈怠）。不过，他在这里显然怀着热忱，旨在（当他即将引入关于信德的论述时）揭示法律的极度无力。他的意思是：倘若你以法律自夸，它反倒使你更蒙羞；它庄重地将你的诸罪宣告在你面前。只是他没有用如此严厉的措辞，而是再次以平和的语调说：\BibleQuote{因为法律只叫人认识罪。}因此刑罚更重，但这是因为犹太人的缘故。因为法律为你达成了罪的揭露，而你本该因此逃避罪。既然你没有逃避，你便更沉重地将刑罚引到自己身上，法律的美善反倒成了你更大报复的根源。如今，既然他已增添了他们的恐惧，他便接着引入恩宠之事，使他们强烈渴望罪的赦免，于是说道：\par

第21节。\BibleQuote{但如今，天主的正义在法律之外已显示出来。}\BibleRef{罗 3:21}\par

在这里他讲出一件重大的事，且需要许多证明。因为如果那些生活在法律之下的人不仅未能逃脱惩罚，反而因此更加被压垮，那么，没有法律，怎么可能不仅逃脱惩罚，甚至成义呢？因为他在这里提出了两个高超之点：成义且无须法律便获得这些祝福。正因如此，他没有简单地说\Quote{正义}，而是说\Quote{天主的正义}，以便以位格的尊贵彰显恩宠的更大程度，以及应许的可能性。因为在天主，一切都是可能的。他没有说\Quote{被赐予}，而是说\Quote{显示出来}，从而取消了\Quote{新奇}的指责。因为那被显示出来的，乃是古老的，只是被隐藏了。不仅如此，下文也表明这并非新事。因为说过\Quote{显示出来}之后，他接着说：\par

\BibleQuote{有法律和先知作证。}\par

他意思是：不要因它刚刚被赐下而烦恼，也不要因它像新奇陌生的事而惊恐。因为古时法律和先知都预言了它。他在论述过程中已指出一些经文，不久还将指出另一些；他在前面引用了\Person{哈巴谷}的话：\BibleQuote{义人将因信德而生活}\BibleRef{罗 1:17}，而在后面则引用了\Person{亚巴郎}和\Person{达味}，因为他们本人也与我们谈论这些事。他们对这些人的尊敬是大的，因为一位是族长兼先知，另一位是君王兼先知；而且关于这些事的应许都已赐予他们二人。正因如此，\Person{玛窦}在他的福音开头首先提到这两人，然后按顺序引出先祖。因为他说了\BibleQuote{耶稣基督的族谱}\BibleRef{玛 1:1}之后，他没有在提到\Person{亚巴郎}后立即也提\Person{依撒格}和\Person{雅各伯}，而是将\Person{达味}与\Person{亚巴郎}一同提及。令人惊叹的是，他甚至将\Person{达味}置于\Person{亚巴郎}之前，如此说：\BibleQuote{达味之子，亚巴郎之子}，然后才开始列出\Person{依撒格}、\Person{雅各伯}以及其余所有按顺序排列的人。正因如此，\Person{宗徒}在此交替呈现他们，并谈到天主的正义有法律和先知作证。然后，为免有人说：我们怎么能什么都不贡献就得到拯救呢？他表明我们也为此贡献了不小的事，即我们的信德。因此，在说了\Quote{天主的正义}之后，他立即加上：\BibleQuote{因信德，为所有相信的人。}\par

在这里，犹太人再次因自己并不比其余人更好，且被列于全世界之中而惊惶。为免他感到这一点，他又以恐惧降伏他，加上：\BibleQuote{因为并没有分别，众人都犯了罪。}因为不要对我说：这是某某希腊人，这是某某西古提人，这是某某色雷斯人；因为众人都处于同样境地。纵然你领受了法律，你从法律所学的只有一件事——认识罪，而非逃避罪。其次，为免他们说：纵然我们犯了罪，仍未与他们一样，他加上：\BibleQuote{并欠缺了天主的光荣。}所以，即使你没有犯与他人同样的罪，你仍同样被剥夺了光荣，因为你属于那些冒犯者；冒犯者不属于得光荣之人，而属于蒙羞之人。然而，不要害怕：我说这话，并非要叫你陷入绝望，而是要显示主（\OriginalTerm{主}{\GreekText{Δεσπότου}}）对人的爱；于是他接着说：\par

第24、25节。\BibleQuote{但因他的恩宠，藉着在基督耶稣内的救赎，白白地成义；这耶稣天主立为赎罪祭，藉着信德，靠他的血，为显示他的正义。}\par

请看，他是藉着多少证据来证实他所说的话。首先，是从这一位本身的尊荣：因为行这些事的不是一个人，不应因之软弱，而是全能的天主。因为他说道，义德是属于天主的。其次，是从法律和先知来的。因为你听到\Quote{没有法律}时，不必害怕，因为法律本身也赞成这一点。第三，是从旧约下的祭献来的。正因如此，他才说\Quote{因祂的血}，以唤起他们对那些绵羊和牛犊的回忆。因为如果无理性之物的祭献，他说，能洁净罪恶，那么这血岂不更能？而且他说的不仅仅是\OriginalTerm{赎价}{\GreekText{λυτρώσεως}}，而是\OriginalTerm{完全赎价}{\GreekText{ἀπολυτρώσεως}}，指完全的赎救，以表明我们不应再陷入那样的奴役中。基于同样的理由，他称之为赎罪祭，以表明如果预像有那样的效力，那么实体更会显明同样的效力。但为了再次表明这并非新奇或最近的事，他说\Quote{预定了}\EditorialNote{参考钦定本边注}；且藉着说天主\Quote{预定了}，并表明这善工是父的，他也表明它也是子的。因为父\Quote{预定了}，但基督藉着自己的血正确地完成了全部。\par

\Quote{为显明祂的义德。}显明义德是什么意思？如同显明祂的丰富，不仅是为了祂自己富有，也是为了叫别人富有；或如显明生命，不仅是为了祂自己活着，也是为了叫死人活过来；又如显明祂的能力，不仅是为了祂自己有能力，也是为了叫软弱的人有能力。同样，显明祂的义德也不仅是为了祂自己正义，也是为了叫那些充满罪之\OriginalTerm{腐烂}{\GreekText{κατασαπέντας}}的人忽然成为正义的。而为了解释这一点，即什么是\Quote{显明}，他接着加了这句话：\Quote{为使他成为正义的，并使他成为使那信耶稣基督的人成义的那一位。}所以不要疑惑：因为不是出于行为，而是出于信德；也不要逃避天主的义德，因为它有两方面的益处：既容易，又对所有人开放。不要羞愧难当。因为如果祂自己公开\OriginalTerm{表明}{\GreekText{ἐνδείκνυται}}祂这样做，而且，可以这样说，祂以此为乐、以此为荣，那么，当你的主以此夸耀时，你怎能沮丧掩面呢？现在，在他藉着说所发生的事是显明天主的义德而提高了听者的期望之后，他接着用恐惧来催促那迟延和懈怠、不愿前来的人，这样说：\par

\Quote{为宽恕从前所犯的罪。}你看到他多么频繁地提醒他们想起自己的过犯？之前，他藉着说\Quote{藉着法律便是对罪的认识}；之后又藉着说\Quote{众人都犯了罪}，但在这里用了更强烈的语气。因为他不是说\Quote{为罪}，而是说\Quote{为宽恕}，即死寂。因为再也没有恢复健康的希望了，但如同瘫痪的身体需要从上而来的手，那已经死去的灵魂也是如此。而更糟的是，有一件事他作为指控提出来，并指出这是更大的控告。那是什么呢？就是最后的状态是在天主的容忍中发生的。因为，他的意思是，你不能辩解说你未曾享受过大量的容忍和良善。而\Quote{在此时期}这句话，是一个指出那大能（或容忍）和对人的爱之伟大的人说的。因为，他可以说，在我们完全放弃之时，在应该判我们罪之时，在邪恶已大、罪恶已满之时，祂彰显了祂自己的能力，好叫你明白祂的义德是何等丰盈。因为这件事，若在起初发生，就不会像现在这样显得如此奇妙和不寻常，那时各种疗治都被发现无效了。\par

27.'\Quote{那么，哪里还有可夸耀之处呢？它被排除了，}他说。\Quote{藉着什么法律？藉着行为的法律吗？不，而是藉着信德的法律。}\newline{}\newline{}\BibleRef{罗 3:27}\par

保禄极费苦心，要证明信德之强大远超法律所能臆想。因为在他说明天主藉信德使人成义之后，又再次与法律辩论。他并未说：犹太人的善行何在？他们的正义行为何在？而是说：\BibleQuote{那么，夸耀在哪里呢？}借此抓住每个机会指出，他们只是空言自夸，仿佛比别人多些什么，却拿不出任何业绩。在说了\BibleQuote{那么，夸耀在哪里呢？}之后，他没有说：它已消失、终结，而是说：\BibleQuote{它被排除在外。}这个词更表达不合时宜，因为其理由已不复存在。正如审判来临后，想要悔改的人不再有时间；同样，如今判决已然下达，众人濒临灭亡，那要藉恩宠消除恐怖的一位已在眼前，他们便不再有藉法律夸耀自身改善的时机。因为若要以此自恃，本应在他来临之前。但现在藉信德施行拯救的那一位已经来临，那些努力的时机已然被夺走。既然所有人都被定罪，他便藉恩宠拯救。正因如此，他才现在来临，免得他们像当初他若立即来临那样说：原是可能藉法律和我们自己的劳作与善行得救的。为遏止他们的厚颜，他等待了很长时间，好让他们在一切论证中都清楚证明自己无力自救时，才藉恩宠拯救他们。也正因如此，他在上文说了\BibleQuote{为彰显他的正义}之后，又加上\BibleQuote{在现今的时候}。若有人反驳，便如同一个人犯了大罪，却无法在法庭上辩护，已被定罪将要受罚，然后蒙王赦免，却在赦免之后厚颜夸口说自己从未犯罪。因为在赦免来临之前，本应证明自己无罪；赦免之后，他再没有夸口的时机。犹太人的情形正是如此。因为他们背叛了自己，所以他来临，藉着他的来临本身消除了他们的夸耀。因为那自称\BibleQuote{愚蒙人的教师}，\BibleQuote{在法律中自夸}，\BibleQuote{愚昧人的导师}的人，若与他们一样需要一位导师和救主，便再无夸口的借口。因为即使此前，割礼已算为未受割礼，如今更是如此，因它已被两个时代一同弃绝。在说了\BibleQuote{它被排除}之后，他也表明如何排除。那么他是如何说被排除的呢？\BibleQuote{藉着什么法律？藉着行为的法律吗？不，是藉着信德的法律。}看，他也称信德为法律，乐于沿用这些名称，以缓解表面的新奇。但\Quote{信德的法律}是什么？就是靠恩宠得救。此处他显示天主的权能：他不仅拯救，还使人成义，并引向夸耀，且无需行为，只期待信德。他说这话，意在使已信主的犹太人行于节制，并安抚未信者，从而引向他的观点。因为那已得救的，若因坚守法律而骄傲，则会被告知：他自己封住了自己的口，控告了自己，放弃了对自身得救的诉求，并排除了夸耀。而那未信的，同样因这些手段而自卑，便能被引向信德。你看到信德何等卓越？它如何将我们从旧事中移除，甚至不容我们对它们夸耀？\par

第28节。\BibleQuote{因此我们认为：人成义是因着信德，而不在于法律的行为。}\BibleRef{罗 3:28}\par

当他表明藉信德他们优于犹太人之后，便满怀信心地继续论及信德，并将其中似乎令人困扰之处再次化解。因为有两件事困扰犹太人：一是若法律行为不能拯救人，人是否可以不经行为而得救；二是未受割礼者是否可与那些长期受法律培育者共享同样福分，后者比前者更令他们困扰。为此，他先证明了前者，接着处理后者，这令犹太人极其困惑，以至他们在信主后还因这立场指责伯多禄。那么他说了什么呢？\Quote{因此我们认为：人成义是因着信德。}他没有说犹太人，或法律之下的人，而是将论题引入广阔空间，为世界打开信德之门，他说\Quote{人}，这是我们族类的通称。然后借此机会，他又回应了一个未明言的异议。因为犹太人听到信德使每个人成义，很可能会不悦而反感，于是他继续：\par

第29节。\BibleQuote{难道天主只是犹太人的天主吗？}\BibleRef{罗 3:29}\par

仿佛在说：如果每个人都应得救，这怎能使你们觉得不妥？难道天主偏心吗？由此他表明，他们想藐视外邦人，实际是对天主光荣的侮辱——如果，他们不让天主作众人的天主。但若他是众人的，他就眷顾众人；若他眷顾众人，他便同样藉信德拯救众人。正因如此，他说：\BibleQuote{难道天主只是犹太人的天主吗？不也是外邦人的天主吗？是的，也是外邦人的。}因为他不像外邦神话（参考奥维德\WorkTitle{哀歌} I. ii. 5 等）那样偏爱谁，而是众人共同的、唯一的天主。正因如此，他接着说：\par

第30节。\BibleQuote{既然天主是唯一的。}\BibleRef{罗 3:30}\par

这就是说，同一位就是这些人和那些人的主。但若你对我提及古代的情形，那么天主眷顾的作为也由双方共享，尽管方式不同。因为正如你们领受了成文的法律，他们领受了自然律；他们在任何事上都不欠缺，只要他们愿意，甚至还能超过你们。因此他接着就暗指这一点说：\Quote{谁将凭信德使受割礼者成义，并藉信德使未受割礼者成义？}这样提醒他们先前关于未受割礼和割礼所说的话，借此表明没有差别。但若那时没有差别，如今就更没有了。因此他现在基于更清楚的根据建立这一点，并如此证明，双方同样需要信德。\par

第31节。他说：\Quote{那么我们藉信德废除了法律吗？绝对不可！反而巩固了法律。}见：\BibleRef{罗 3:31}\par

你看到他多样而不可言喻的判断吗？因为仅仅使用\Quote{巩固}这个词，就表明它当时并未站立，而是已被废弃了（\OriginalTerm{废弃的}{\GreekText{καταλελυμένον}}）。也要注意保禄的卓越能力，以及他如何超额地坚持他所愿望的。因为他在这里表明，信德不但没有对\Quote{法律}有任何贬损，反而协助了它，正如法律反过来为信德铺了路。因为正如法律本身先前为它作证（因为他说：\Quote{由法律和先知所证明}），同样，这里这也证实了，既然法律已失去力量。它如何巩固？他会说。法律的目标是什么，它所有规定的内容是什么？就是使人成义。但法律没有能力做到这一点。它说：\Quote{因为众人都犯了罪，}但信德来到时却成就了此事。因为一旦一个人成为信徒，他就立刻成义。因此法律的本意，它确实巩固了，而所有规定所追求的目标，它将其带到了完成。所以它没有废除法律，而是成全了它。这里他证明了三点：第一，没有法律也可能成义；第二，法律不能成就此事；第三，信德与法律并不对立。因为使犹太人困惑的主要原因就是信德似乎与法律对立，他比犹太人更清楚地表明，信德非但不相反，反而与法律紧密结盟且合作，这正是他们特别渴望听得到证明的。\par

但既然在我们成义之后，还需要与此恩宠相称的生活，就让我们显露出配得上这恩赐的热忱吧。我们若以热忱持守那善行之母——爱德，便能如此。爱德不在于空言或仅仅对人说话\OriginalTerm{言谈}{\GreekText{προσρήσεις}}，而在于实际关怀\OriginalTerm{照顾}{\GreekText{προστασία}}他们，并以行动表现出来，例如救济贫穷、扶助病患、解救危难、支持困苦者、\BibleQuote{与哭泣的人一同哭泣，与喜乐的人一同喜乐。}\BibleRef{罗 12:15}因为连后者也是爱德的一部分。然而，这似乎是一件小事——与喜乐的人同乐，但它极其伟大，需要真正智慧的精神。我们发现许多人能履行那更艰难的部分\OriginalTerm{更苦的事}{\GreekText{πεικρότερον}}，却在此事上缺乏力量。因为许多人能与哭泣的人同哭，却不能与喜乐的人同乐，反在别人喜乐时流泪；这出于嫉妒和猜忌。那么，与兄弟同乐的善行并非小事，甚至比前者更大；也许不仅比与哭泣者同哭更大，甚至比支持身处危险中的人更大。无论如何，有许多人曾与危险中的人共患难，但当那些人获得尊荣时，他们却心如刀割。嫉妒之心是多么专横！然而，前者是辛劳和工作，后者只涉及选择和心态。然而，许多忍受了更艰难任务的人，却没有完成这更容易的事，反而在看见别人得荣耀、整个教会因教义或其他方式受益时，憔悴消耗。还有什么比这更糟的呢？因为这样的人不再是与兄弟争斗，而是与天主的旨意争斗。请思考这一点，并摆脱这疾病：即使你不愿释放你的邻人，至少使自己摆脱这无数的邪恶。为何将战争带入你自己的思想？为何使你的灵魂充满烦恼？为何掀起风暴？为何倾覆一切？在这种心态下，你怎能祈求罪过的赦免？因为如果那些不容忍他人对自己的冒犯的人，天主也不宽恕他们，那么对于那些试图伤害那些没有伤害他们的人，天主会赐予什么样的宽恕呢？因为这是极端邪恶的证明。这样的人是与魔鬼一同作战，反对教会，甚至可能比魔鬼更坏。因为我们可以防备魔鬼，而这些人却伪装在友善的面具下，秘密点燃火堆，先把自己投入火炉，患上的疾病不仅不值得怜悯，反而成为极大的笑柄。因为请告诉我，你为何脸色苍白、颤抖、恐惧？发生了什么灾祸？是因为你的兄弟得了荣耀、受尊敬、被重视吗？你本该制作花冠、欢喜并光荣天主，因为你自己的肢体得了荣耀和尊敬！但你却因天主受光荣而痛苦吗？你看出这场战争导致什么结局吗？但有人会说：我并非因天主受光荣而痛苦，而是因我的兄弟。然而，荣耀通过他上达于天主；同样，你的战争也上达于天主。但那人会说：使我忧伤的不是这个，因为我愿天主通过我得到光荣。那么好吧！当你兄弟得荣耀时，你也喜乐，这样天主也通过你得到光荣；众人都会说：\Quote{愿天主受赞美，祂的家庭如此同心，完全摆脱嫉妒，为彼此的好处同乐！}我为何只提你的兄弟？即使他是你的仇敌或敌人，而天主通过他得了光荣，你也应因此把他当作朋友。但你却因天主因他的荣耀而受光荣，把朋友当作敌人。若有人治好你病痛的躯体，即使他是敌人，你从此也会把他列为最亲密的朋友；而那个使基督的身体（即教会）容光焕发的人，又是你的朋友，你却视他为敌人？你还能怎样显明你对基督的战争呢？因此，即使一个人行奇迹、守独身、禁食、睡在地上，凭借这样的美德甚至达到天使的境界，但他若有这缺陷，仍将是所有人中最受咒诅的，比奸淫者、娼妓、强盗和盗墓者更是违犯律法的人。为了使无人指责这是夸张之词，我乐于向你提出这个问题：如果有人带着火和镐来，毁坏焚烧这所圣殿，挖倒这祭台，在座的每一位难道不会像对待被咒诅的违法者一样用石头砸死他吗？那么，如果一个人带来比那火更炽烈的火焰——我指的是嫉妒——它不毁坏石造建筑，也不挖倒金祭台，而是推翻并轻蔑地毁坏比墙壁和祭台更宝贵的东西——教师的建筑，他该受怎样的惩罚呢？因为没有人可以告诉我：他多次尝试但未能成功；因为行为是由心态来判断的。\BibleQuote{因为撒乌耳确实想杀达味，即使他没有击中他。}\BibleRef{撒上 19:10}告诉我，当你与牧人作战时，你没有意识到你在谋害基督的羊群吗？那些羊群，基督也曾为它们倾流鲜血，并命令我们为它们做一切、受一切苦？你不提醒自己，你的主寻求你的荣耀而非自己的，而你却寻求自己的而非主的荣耀吗？然而，如果你寻求祂的，你也会得到你自己的。但因你优先寻求自己的而非祂的，你连自己的也永远得不到。\par

那么，补救办法是什么呢？让我们一同祈祷，齐声为他们代求，如同为附魔者代求一样，因为他们比附魔者更可怜，因为他们的疯狂是出于自愿。这种痛苦需要祈祷和恳求。\newline{}因为，若一个人不爱他的弟兄，即使他散尽钱财，甚至获得殉道的荣冠，也毫无益处；那么，想一想，那与未曾得罪过他的人争战的人，该受怎样的惩罚呢？他甚至比外邦人更糟。因为，如果爱那些爱我们的人，并不能使我们比他们更有优势，那么，嫉妒那些爱他的人，他又该列入哪一等呢？告诉我，因为嫉妒甚至比争战更糟；因为争战的人，当战争的原因结束时，也停止了仇恨；但嫉妒者永远不会成为朋友。一个表现出公开的战争，另一个则是隐蔽的；一个常常能给战争提出合理的理由，另一个则除了疯狂和撒殚的精神外，一无所有。\newline{}那么，将这样的灵魂比作什么呢？比作什么样的毒蛇？比作什么样的蝮蛇？比作什么样的蛀虫？比作什么样的蝎子？因为没有比这样的灵魂更可咒诅、更有害的了。\newline{}正是这个，正是这个，颠覆了教会，生出了异端，武装了兄弟的手，使他的右手浸在义人的血中，拔去了自然的律法，为死亡敞开了大门，使那个诅咒付诸行动，不让那个可怜的人想起生产的痛苦，或他的父母，或任何其他事情，反而使他如此疯狂，引他到这样的癫狂地步，以至于天主劝告他说：\BibleQuote{它必迷恋你，你要制服它}\BibleRef{创 4:7}；他甚至仍不屈服。然而，天主既赦免了他的过错，又使他的兄弟顺服于他；但他的怨言是如此不可救药，即使施用千万种药，它仍因自身的腐坏而溃烂。\newline{}因为你为什么如此恼怒，你这最可怜的人？是因为天主得了荣耀吗？不，这倒显出了撒殚的精神。那么，是因为你的弟兄在名誉上超越了你吗？说到这点，你也可以反过来超越他。所以，若你想成为得胜者，不要杀害，不要毁灭，而要让他存活，好使斗争的素材得以保存，并活活地征服他。因为这样，你的冠冕将是光荣的；但如此毁灭，你等于对自己宣告了更严厉的失败。然而，嫉妒的心对此毫无感觉。\newline{}你在这如此孤独的环境中有什么可贪求荣誉的呢？因为那时地上只有他们二人居住。然而，即便如此，也未能阻止他，他反而抛却一切，站到了撒殚的行列中；因为当时\Person{撒殚}正是站在\Person{加音}一边作战。因为既然人已经必须死于死亡，这对撒殚还不够，他还试图通过死亡的方式使悲剧更加惨烈，并说服他成为杀弟者。因为他急于并切望看到判决的执行，从不满足于我们的苦难。就像一个人把敌人关在监狱里，见他已被判刑，却在他出城之前就急着要在城内看他被屠杀，连适当的时间也不肯等；那时\Person{撒殚}也是这样，虽然他已听到人必须归于尘土，却仍切望看到更糟的事，即儿子先于父亲而死，兄弟杀死兄弟，以及过早而暴烈的杀戮。\newline{}你看，嫉妒为他成就了多大的功业？它怎样满足了\Person{撒殚}贪得无厌的精神，并为他预备了如他所愿的大宴？\par

那么，让我们从这疾病中逃离吧；因为要逃脱为魔鬼预备的火是不可能的，确实不可能，除非我们摆脱这病患。但我们若能牢记基督如何爱了我们，以及祂如何命令我们彼此相爱，我们就能获得自由。祂向我们显示了何等的爱？当我们还是仇敌，并对祂犯下极大不义时，祂为我们倾流了祂的\OriginalTerm{宝血}{Blood}。你也要对你的兄弟这样做（因为祂这话的结局是：\Quote{我给你们一条新诫命：你们该彼此相爱，如同我爱了你们一样}\BibleRef{若 13:34}；或者更确切地说，即使这样，尺度也没有停止）。因为祂这样做是为了祂的仇敌。你难道不愿意为你的兄弟流血吗？那么你为什么反而流他的血，违背诫命甚至颠倒它？然而祂所做的并非理所当然；但如果你做了，你不过是尽一份债务。因为那个领了一万塔冷通而后索要一百德纳的人，受罚不仅因为他索要，而且因为他受了恩惠却未变得更好，甚至没有跟随他主人所开始的，或豁免债务。因为对仆人来说，如果做了，也终究只是债务。因为我们所做的一切，都是为了还债。正因如此，祂亲自说：\Quote{你们做完了一切，要说：我们是无用的仆人，我们只做了我们应做的事。}\BibleRef{路 17:10}因此，如果我们显示\OriginalTerm{爱德}{charity}，将我们的财物给需要的人，我们就是在尽一项债务；这不仅因为祂首先开始了善行，更因为如果我们给予，我们分配的其实是祂的财物。那么，你为什么要剥夺祂愿意你拥有的权利呢？因为祂命令你给别人，正是为了让你自己拥有它们。因为只要你为自己拥有，你甚至不拥有它们；但当你给了别人，你才真正得到了它们。还有什么魅力能与此相比？祂自己为仇敌倾流了祂的宝血，而我们甚至不为恩人花钱。祂流的是自己的血，我们连不属于自己的钱都不肯花。祂在我们之前做了，我们甚至不跟随祂的榜样。祂为我们的\OriginalTerm{救恩}{salvation}做了，我们甚至不为自己的益处去做。因为祂并不从我们对人的爱上得到任何好处，而是全部益处归于我们。这正是我们被命令舍弃财物的原因，免得我们被剥夺它们。就像一个人给小孩钱，命令他紧紧握住，或交给仆人保管，以免被随意抢夺；天主也是如此。因为祂说：给需要的人，免得有人从你那里抢夺，比如告密者、诽谤者、盗贼，或者避免这一切之后，死亡。因为只要你亲自握着它，你就没有安全地持有它。但如果你通过穷人把它给我，我就为你完全保存它，并在适当的时候加倍归还。因为我接受它并不是要拿走，而是要使它增多并更精确地保存，以便为你保留到那个无人可借无人可怜的时候。那么，还有什么比在如此应许之后，我们仍不能下定决心借给他更冷酷的呢？是的，正因此我们站到祂面前时，是赤贫、赤裸、匮乏的，没有拥有交托给我们的东西，因为我们没有把它们存放在那位保管得比任何人都精确的祂那里。为此我们将受最严厉的惩罚。因为当我们被质问时，关于它们的损失我们能说什么？能提出什么借口？什么辩护？你有什么理由不给？难道你不相信会再得回来吗？这怎么合理？因为那曾给过没有给的人，祂在收到之后岂不更愿给予？看到它们使你高兴吗？那么，正因为此，你要更慷慨地给予，以便在那无人能夺走的地方，你会更喜悦。因为现在如果你保留它们，你甚至会遭受无数祸患。因为如同狗扑向富人一样，魔鬼也这样扑向他们，想要从他们那里抢夺，如同从一个拿着饼干或蛋糕的孩子那里抢夺。那么，让我们把它们交给我们的天父；如果魔鬼看到这事，他一定会退去；当他退去后，天父就会安全地把它们全部交给你，在那来世，他不能再骚扰你。因为如今富人确实无异于被狗骚扰的小孩子，众狗都在他们周围吠叫、撕咬、拉扯；不仅是人，还有卑贱的情欲，如贪饕、醉酒、谄媚、各种不洁。当我们必须借贷时，我们非常担心那些给得多的人，并特别寻找那些坦率的借贷者。但在此我们却相反。因为对于那位坦率行事、不按百分之一而按百倍给予的天主，我们抛弃了祂，而对于连本金都不归还的人，我们却追求他们。因为我们的肚腹消耗我们财物的大部分，它能回报我们什么？粪土和腐朽。虚荣能回报什么？嫉妒和怨恨。亲近能回报什么？忧虑和焦虑。不洁能回报什么？地狱和毒虫！因为这些就是富人的债主，他们以这些为资本付利息：眼前的邪恶和将来的可怕之事。那么，让我们投靠这些以如此惩罚为利息的债主吗？而我们却不信赖那向我们展示天堂、永生和说不尽恩福的基督\EditorialNote{四份抄本省略了\GreekText{τἥ}}？我们还有什么借口？你怎么能不给那必定回报、且更丰盛回报的祂呢？也许是因为祂偿还的时间太久。然而祂确实在此世也偿还。因为祂是信实的，祂说：\Quote{你们先寻求天国，这一切自会加给你们。}\BibleRef{玛 6:33}你看到这极致的慷慨吗？祂说，那些财物已为你贮存，并不减少；而这里的这些，我以加增和盈余的方式赐给你。但除此之外，你将在很久以后才收到这个事实，反而使你的财富更大，因为利息更多。\par

在借钱给人这件事上，我们看到放贷者正是这样做的：他们更愿意借给那些很久之后才还钱的人。因为立刻还清全部债务的人，切断了利息的增长；而拖延更久的人，则使收益更大。那么，既然在人身上我们并不因延迟而生气，反而设法使之更大，难道在天主身上，我们竟如此心胸狭隘，以至于以此为由退缩和反悔吗？然而，如我所说，祂在此世也给予，并且除了已经提及的理由，祂还为我们计划其他更大的益处，所以将一切存留于彼处。因为所赐之物的丰盛及其卓越，超越了今世短暂的生命。在这朽坏与注定死亡的身体中，甚至不可能接受那些永不衰败的冠冕；在我们这动荡不安、充满烦恼、易受诸多变化的状态中，也不可能达到那不变不动之境。现在，如果有人欠你金子，而你在异乡停留，既无仆人，也无办法将金子运回住处，他承诺偿还贷款，你岂不会以无数方式恳求他不要在外地支付，而是在家乡支付吗？但你却认为应当在此世接受那些属灵的、不可言喻的事物？这是何等疯狂！因为如果你在此世接受，它们必定是可朽坏的；但如果你等待那时，祂将以不朽坏、纯全的方式偿还你。如果你在此世接受，你得到的是铅；但若在那里，你得到的是精炼的金子。然而，祂甚至没有剥夺你今生的美物。因为在那应许之外，祂还附上了另一项，大意是：凡爱慕来世事物的人，将在此生获得\BibleQuote{百倍，并承受永生}（\BibleRef{玛 19:29}）。如果我们没有得到百倍，那是我们自己的错，因为我们没有借给那能赐予如此之多者；因为凡是给予的，即使只给了少许，都得到了许多。告诉我，伯多禄给了什么大事？不就是一个破了的网（\BibleRef{路 5:6}），只有一根竿子和一个钩子吗？然而天主向他敞开了世界的房屋，将海陆铺展在他面前，并邀请所有人将自己的财产带到他脚前。或者更确切地说，他们变卖自己的财物，放在宗徒脚前，甚至没有放入他们手中，因为他们不敢；他对他们的尊敬和慷慨是如此之大。但你会说：他是伯多禄！那又怎样？人啊！因为祂不仅向伯多禄作了这应许，也没有说：伯多禄，唯独你得到百倍，而是说\BibleQuote{凡离开房屋或兄弟的，必要得到百倍}。因为祂不看人的身份，而是看所行的正义。有人说：\Quote{我身边有一群小孩子，我想给他们留下好运气}。那么我们为什么使他们成为穷人？如果你把所有东西都留给他们，你仍是将你的财产交托给一个可能背叛的信托。但如果你让天主成为他们的共同继承人和守护者，你就给他们留下了无数的宝藏。因为正如当我们为自己复仇时，天主不帮助；但当我们把这事交给祂时，事情会超出我们的期望；同样，在财产上，如果我们自己操心，祂就会撤回眷顾；但如果我们把一切交给祂，祂将安全地安置财产和我们的孩子。你为什么对天主这样的事感到惊讶？因为在人身上也能看到这事。如果你在临终时没有邀请任何亲戚照顾你的孩子，常常会发生这样的情况：一个非常愿意的人会感到犹豫，不好意思主动承担。但如果你把责任托付给他，他会认为你给了他极大的荣誉，从而以极大的回报作为报答。因此，如果你想为你的孩子留下巨大财富，就给他们留下天主的照顾。因为祂没有你任何作为，就赐予你灵魂，塑造你的身体，赐予你生命，当祂看到你如此慷慨，将自己的财物与祂一起分施时，祂必定会向他们开启一切财富。如果厄里亚在吃了少许面饼之后，看到那妇人尊敬他超过自己的儿子，就在寡妇的小屋里使谷仓和油榨出现，那么厄里亚的主将显示何等慈爱！因此，我们不要考虑如何使孩子富有，而是如何使他们有德。因为如果他们拥有财富的自信，就不会在意其他任何事，因为他们有手段在丰裕的财富中掩盖其恶行。但如果他们发现自己缺乏来自财富的安慰，他们就会尽力以德行来弥补贫穷。因此，不要留给他们财富，以便留给他们德行。因为我们活着时不让他们成为我们所有财物的主人，而在我们死后却给青春的放纵以完全免于恐惧的自由，这是极不合理的。而且，我们活着时还能好好规劝他们，如果他们滥用财物，我们还能约束他们；但如果我们死后，在他们失去我们以及他们青春年少之时，给予他们财富带来的能力，那我们将把那不幸可怜的人推入无尽的深渊，如同在火焰上添柴，在烈火中浇油。因此，如果你希望他们既富有又安全，就让天主成为他们的债务人，并将遗产交付在祂手中。因为如果他们自己收到钱，他们甚至不知道给谁，会遇到许多狡猾无情的人。但如果你预先将钱借给天主，这宝藏从此就不可侵犯，偿还也极为容易。因为天主乐于偿还祂所欠的债，祂更青睐那些借给祂的人，胜过那些没有借给祂的人；祂最爱那些祂欠得最多的人。因此，如果你希望祂持续作你的朋友，就要让祂成为你大量的债务人。因为没有放贷者像基督（有的抄本作\Quote{天主}）那样，因有欠债者而喜悦；祂甚至奔向那些欠祂债的人。那么，让我们用尽一切方法使祂成为我们的债务人；因为现在是借贷的季节，祂如今有需要。如果你现在不给予祂，祂在你离世后不会向你索取。因为祂在这里口渴，在这里饥饿。祂口渴，是因为祂渴望你的得救；为此祂甚至乞讨；为此祂赤身露体，为你谈判永生。不要漠视祂；因为祂不是要得滋养，而是要滋养你；祂不是要穿衣，而是要给你穿上金色衣裳、君王的外袍。你没有看到那些更有爱心的医生，在给病人洗澡时，自己也洗了，尽管他们不需要？同样，祂所做的一切都是为了生病的你。为此，祂在要求时不用强迫，以便给你巨大的回报：使你明白祂向你要求不是因为祂有需要，而是为了补救你的需要。为此，祂以卑微的姿态来到你面前，伸出右手。如果你给祂一个铜钱，祂不会转脸；即使你拒绝祂，祂也不离开，而是再次来到你面前。因为祂极其渴望我们的得救。那么，让我们轻视钱财，以免被基督轻视。让我们轻视钱财，甚至是为了获得钱财本身。因为如果我们在此世保留它，我们将在此世和来世完全失去它。但如果我们以大量的花费分施它，我们将在今生和来世都享受丰裕的财富。因此，谁想富有，就让他贫穷，以便他富有；让他花费，以便积累；让他分散，以便聚集。但如果这听起来新奇而矛盾，请看撒种的人：他唯有通过分散手中已有的才能收集得更多。现在，让我们播种并耕耘天国，以便获得大丰收，并通过恩宠和祂对人的爱，等等。\par

\SourceRecord{Translated by J. Walker, J. Sheppard and H. Browne, and revised by George B. Stevens.；Nicene and Post-Nicene Fathers, First Series；Vol. 11.；Edited by Philip Schaff.；Buffalo, NY: Christian Literature Publishing Co.,；1889.；<http://www.newadvent.org/fathers/210207.htm>.}

% structure: p0001 source=h1 target=section reason=page-title

\section{罗马书讲道 第八篇}

罗马书 4:1-2\par

\BibleQuote{那么，我们按肉身的祖先亚巴郎找到了什么呢？因为如果亚巴郎是因行为而成义，他就有可夸耀之处，但在天主前却不能。}\par

他曾说过（五个抄本作\OriginalTerm{他说}{\GreekText{εἶπεν}}），世界在天主前已变为有罪，众人都犯了罪，夸耀已被排除，除非藉着信德，否则无法得救。他现在致力于表明，这救恩非但不是可耻的事，反而是辉煌荣耀的缘由，且大于透过行为所得的荣耀。因为既然得救而带着羞耻，会有些沮丧，他接下来也消除了这种疑虑。而且他早已暗示了这一点，因他称它不只是救恩，更是\Quote{义德}。他说：\BibleQuote{天主的义德在其中显示出来。}\BibleRef{罗 1:17} 他称这不只是\Quote{义德}，也是天主义德的彰显。但天主是在光荣、灿烂、伟大的事物中彰显的。然而，他仍从当前所讨论的内容中为此寻求支持，并通过提问的方式推进他的论述。他总是这样做，既为清晰，也为使所说的话令人确信。例如，他先前就这样做了，他说：\BibleQuote{那么，犹太人有什么长处呢？}\BibleRef{罗 3:1}，以及\BibleQuote{那么，我们比他们强在哪方面呢？}\BibleRef{罗 3:9}，又说：\BibleQuote{那么，夸耀在哪里？已被排除}\BibleRef{罗 3:27}；而在这里，他说：\BibleQuote{那么，我们该说什么呢？我们的祖先亚巴郎……}等等。既然犹太人不断反复提及那圣祖、天主的友人第一个接受了割损礼的事实，他就想表明，他也是因信德而成义的。而这是一个非常有利的立足点（\OriginalTerm{丰盛的得胜}{\GreekText{περιουσία} \GreekText{νίκης} \GreekText{πολλἥς}}）。因为对于一个没有行为的人，因信德而成义，并不令人意外。但对于一个富有多项善行的人，不是由此而成义，而是因信德，这才会引起惊叹，并将信德的能力鲜明地显扬出来。这就是为什么他越过所有其他人，而将话题引回此人。他称他为\Quote{按肉身的父亲}，目的是将他们从与他的真实关系（\OriginalTerm{真实亲缘关系}{\GreekText{συγγενείας} \GreekText{γνησίας}}）中排除出来，并为外邦人铺平与他成为亲属的道路。然后他说：\BibleQuote{因为如果亚巴郎是因行为而成义，他就有可夸耀之处，但在天主前却不能。}在说过天主\Quote{因信德使受割损者成义，也藉信德使未受割损者成义}之后，并且已在先前的话中充分确证了这一点，他现在藉亚巴郎更清楚地证明这一点，超过了他的承诺，并为信德而战反对行为，并使这位义人成为整个争论的焦点；这并非没有特殊含义。因此，他也通过称他为\Quote{祖先}来高举他，并约束他们在一切事上与他一致。因为他会说：不要向我提犹太人，也不要提这个人或那个人。因为我要上溯到众人的元首，以及割损礼的起源。因为他说：\BibleQuote{如果亚巴郎是因行为而成义，他就有可夸耀之处，但在天主前却不能。}这里所说的并不清楚，因此必须使其更加清晰。因为有两种夸耀，一种是行为的，一种是信德的。在说了\BibleQuote{如果他是因行为而成义，他就有可夸耀之处，但在天主前却不能}之后，他指出他也可以因信德而有可夸耀之处，而且理由大得多。因为保禄的大能尤其表现在此：他将对方的反对论点转向另一边，并表明那似乎更属于行为救恩的一面，即夸耀或大胆声称（\OriginalTerm{坦然无惧}{\GreekText{παρρησιάζεσθαι}}），实则更真实地属于信德。因为以行为夸耀的人，可提出自己的劳苦；但那以信德在天主内找到荣耀的人，有更大的夸耀理由，因为他光荣并赞美的是天主。因为那些可见世界的本性所不能告诉他的事，他因信德接受这些，就立即显示出对他真诚的爱，并清楚宣告他的大能。这正是最高贵灵魂、哲学精神和高超心智的特征。因为不偷盗、不杀人，是微不足道的成就，但相信天主能做不可能的事，则需要一位不平凡的、热切爱慕他的灵魂；这是真诚之爱的标记。因为遵守诫命的人的确光荣天主，但藉着信德这样追求智慧（\OriginalTerm{爱智}{\GreekText{φιλοσοφὥν}}）的人，光荣的幅度更大。前者服从他，后者则接受关于他适当的意见，光荣他，并对他感到惊叹，超过行为所表达的。因为那种夸耀属于行善的人，但这光荣是天主，完全在于他。因为他因思想关于他的伟大事物而夸耀，这些事归光荣于他。这就是为什么他谈到在天主前有可夸耀之处。不仅为此，还有另一个原因：因为信者再度夸耀，不仅因为他真诚爱天主，也因他享受了来自天主的大爱与尊荣。因为正如他藉着对天主有伟大的思想来表达对他的爱（这是爱的证明），同样天主也爱他，尽管他该为无数的罪受苦，不仅使他免受惩罚，甚至使他成为义人。那么，他就有可夸耀之处，因为他被算为配得上那大爱。\par

第4节。\BibleQuote{因为工作的人，他的工资不算是恩宠，而是债务。}\BibleRef{罗 4:4}\par

那么，这最后的岂不是最大的吗？他意指。断乎不是：因为它是算给相信者的。但若非他自己有所贡献，它就不会被算。所以，他也使天主成为他的债务人，且债务人不是为平常的事，而是为伟大崇高的事。为显示他的高尚精神和属灵悟性，他不仅说\Quote{给相信者}，而是\par

\BibleRef{罗 4:5} \BibleQuote{相信那使不虔敬者成义的天主的人。}\par

请思考，这是多么伟大的事：能完全确信天主不仅能突然释放一个生活在不虔敬中的人免受惩罚，而且还能使他成为义人，并让他配得上那些不朽的光荣。因此，不要以为此人在恩宠尚未算给前者方面被贬低了。因为正是这件事使相信者蒙受光荣：他享受如此伟大的恩宠，并彰显如此伟大的信德。也要注意，赏报更大。因为给前者的是一份酬报，而给后者的却是义德。如今义德远大于酬报。因为义德是一种最完全包含多项赏报的酬报。因此，在从亚巴郎证明了这一点之后，他又引到达味，作为对所说之事的赞同。那么，达味说了什么？他称谁为有福？是那因行事夸胜的人，还是那享受恩宠的人？是那获得赦免和恩赐的人？当我论及真福时，我是指一切美善中最主要者；因为正如义德大于酬报，真福也大于义德。既已证明义德更好，不仅因为亚巴郎领受了它，而且从理性来看（因为他有可夸耀之处，他说，在天主面前），他再次用另一种方式显示其更尊贵，即带进达味以这种方式作证。因为达味也说，他称那如此成义的人为有福，说：\par

\BibleRef{罗 4:7} \BibleQuote{罪恶蒙赦免的人，是有福的。}\par

他似乎带来了一个偏离其目的的见证。因为经文并没有说：\InnerQuote{他们的信德算为义德的人是有福的。}但他故意这样做，并非由于疏忽，而是为显示更大的优越性。因为如果那因恩宠获得赦免的人是有福的，那么那成义且显示信德的人更有福。因为哪里有真福，那里一切羞愧就被除去，并且有许多光荣，因为真福是赏报和光荣的更大程度。为此，他陈述另一方的优势为不成文的：\InnerQuote{行事的，赏报不算为恩宠}；但为了证明信者的优势，他引进圣经的见证说，正如达味所说：\InnerQuote{罪恶蒙赦免，过犯得掩盖的人，是有福的。}他说，你说什么？难道\InnerQuote{他不是出于债务，而是出于恩宠获得赦免}吗？但请看，正是此人被宣称为有福。因为他若没有看见他在享受大光荣，就不会这样称他。而且他并没有说这\InnerQuote{赦免}是临到割损上；而是说什么呢？\par

\BibleRef{罗 4:9} \BibleQuote{那么，这真福}（这是更大的事）\BibleQuote{是落在割损上，还是落在未受割损上呢？}\par

因为现在所探讨的问题是：这美好伟大的事物与谁同在？是与割损，还是与未受割损？并注意其优越性！因为他表明它非但不逃避未受割损，反而在割损之前就已欣喜地与它同住。因为既然称他为有福的是达味，而达味自己也处于割损状态，并且他是对处于那种状态的人讲话，请看保禄如何热切地争辩，将达味所说的话应用到未受割损的人身上。因为在他将真福的宣告与义德联系起来，并表明它们是一回事之后，他询问亚巴郎是如何成为义人的。如果真福属于义人，而亚巴郎成了义人，那么让我们看看他是如何成为义人的：是在未受割损时，还是在受割损时？他说，是在未受割损时。\par

\Quote{因为我们说，信德算为亚巴郎的义德。}\par

在以上引用了圣经之后（因为他说：\Quote{圣经说什么？亚巴郎相信了天主，这就算为他的义德}），他在这里继续确保说话者的判断，并表明成义是在未受割损时发生的。然后从这些根据他又解决了另一个出现的异议。因为如果他在未受割损时就成义了，有人会说，那么引进割损有什么目的呢？\par

\BibleRef{罗 4:11} 他说：\BibleQuote{他领受了}，\BibleQuote{割损的标记和印号，作为那基于信德的义德的印证，这义德是他未受割损时已拥有的。}\par

你们看他如何指出犹太人几乎像是客人的身份（即客人），而不是那些未受割损者，而且这些人是后来加入的？因为如果他在未受割损时已被成义并获得冠冕，犹太人随后才来，那么\Person{亚巴郎}首先是未受割损者之父，这些未受割损者因信德而属于他，然后才是受割损者之父。因为他是两族的先祖。你们看信德大放光明？因为在信德到来之前，圣祖并未成义。你们看未割损并不构成阻碍？因为他未受割损，却未被阻碍成义。因此割损次于信德。何以奇怪它次于信德，当它甚至次于未割损呢？它不仅次于信德，而且远不如信德，正如标记与所标记之实物相比；例如，如同印章之于士兵。（参见\WorkTitle{格林多后书讲道集}第三篇末尾。）他又问：他那时为何需要印章？他自己并不需要。那么他为何接受印章呢？为使他成为相信的人——不论在未割损或割损中的——共同的父亲。但并非绝对地是受割损者的父亲；因此他接着说：\BibleQuote{给那些不只是受割损的人。}因为如果对于未受割损者，他之所以是他们之父，不是因他未受割损（虽然他在未割损中成义），而是因他们效法了他的信德；那么对于在割损状态下的人，他成为他们的先祖更不是由于割损，除非信德也加上了。因为他接受割损的理由是：让我们两派中的人都能以他为先祖，并且让未受割损者不排斥受割损者。看，前者首先要以他为先祖。如果割损因宣讲义德而尊贵，那么未割损甚至在割损之前就已领受了义德，这也不是小的优越。那时，当你行走在那信德的足迹中，不争辩，也不因引入法律而成为分裂的缘由，你才能以他为先祖。什么信德？请告诉我。\par

第十二节. \BibleQuote{这是他在未受割损时所有的。}\BibleRef{罗 4:12}\par

此处他再次贬抑犹太人高傲的精神，提醒他们成义的时间。他适切地说：\Quote{足迹}，为使你与\Person{亚巴郎}一样相信死者身体的复活。因为他在这方面也显示了他的信德。所以，如果你拒绝未割损，你当确知，割损对你再无用处。因为如果你不跟随他信德的足迹，即使你千万次在割损状态中，你也不会是\Person{亚巴郎}的后裔。因为他接受割损正是为此目的：使未受割损的人不会抛弃你。那么，你不要向他要求这个。因为这要帮助的是你，而不是他。但他（保禄）称之为义德的标记。这也是为了你，因为现在连这个也不是了：因为你那时需要身体的标记，但现在不需要了。有人会说：\Quote{难道不可能}，\Quote{从他的信德学习他灵魂的良善？}是的，可能，但你也需要这个附加物。因为你没有效法他灵魂的良善，也不能看见它，一个可感知的割损被赐予你，使你在习惯了这个身体上的标记后，能逐渐被引导到灵魂的真正智慧之爱，并且当你极其认真地把它当作一个极大的特恩接受时，你被教导去效法和尊敬你的先祖。天主有这个目的，不仅在于割损，也在所有其他礼仪中，即祭献、安息日和庆节。现在，你要从后续知道，为了你的缘故他接受了割损。因为他说他接受了一个标记和印鉴后，也给出了理由如下：使他成为割损之父——也是那些接受了属灵割损的人，因为如果你只有这个（即肉体的），再不会有好结果。因为只有当你发现它所标记的现实——即信德——时，它才是标记；如果你没有信德，标记对你不再有标记的力量，因为它是何物的标记？或者，当没有什么可封印时，什么是印鉴？好比有人向你展示一个带封印的钱袋，里面却空无一物。因此，如果里面没有信德，割损就是可笑的。因为如果它是义德的标记，而你却没有义德，那么你也没有标记。因为你接受标记的原因，是要你勤勉地寻求那标记所指向的现实；所以如果你没有标记就能确信随后勤勉地寻求，那么你就不需要它。但割损所宣告的不仅是义德，而且是未割损之人的义德。因此割损只是宣告：割损不需要。\par

第十四节. \BibleQuote{因为如果属于法律的人成为继承者，信德就被废弃，恩许就无效了。}\BibleRef{罗 4:14}\par

他已显示信德是必要的，它比割损更古老，比法律更有力，它建立了法律。因为如果众人都犯了罪，它是必要的；如果一人未受割损而成义，它是更古老的；如果罪的认知来自法律，而它却无需法律而显明，它是更有力的；如果它获得法律的证明，并建立法律，它并不反对法律，而是友好且与法律结盟。此外，他还从其他根据显示，甚至通过法律也不可能获得产业，并且在将信德与割损比较并使它获胜之后，他又用这些话将信德与法律对比：\Quote{因为那些属于法律的若是继承人，信德就落空了。}为了防止有人可能说人可以既有信德又遵守法律，他表明这是行不通的。因为那紧抓法律、视其为拯救力量的人，就贬低了信德的能力；所以他说：\Quote{信德就落空了，}也就是说，不需要因恩宠而得救。因为那时它不能显示其本身的能力；\Quote{那恩许也就被废弃了。}这是因为犹太人可能会说：我需要信德做什么？如果这样，那所恩许的就会与信德一起被废除。请看他是如何从一开始、从圣祖起就处处与他们争辩的。因为从那里显示了义德和信德在产业中并行之后，他现在显示恩许也是如此。为了防止犹太人说：如果亚巴郎是因信德成义，与我何干？保禄说：你所关心的产业恩许，没有信德也无法生效：这是最令他们害怕的事。但他所说的是哪一个恩许？即作\Quote{世界的继承人，}以及万民将因他而受祝福。他如何说这个恩许被废弃了呢？\par

第15节。\Quote{因为法律激起义怒：哪里没有法律，哪里就没有违犯。}\BibleRef{罗 4:15}\par

如果它激起义怒，并使他们成为违犯者，显然也使他们受到诅咒。但那些受诅咒、惩罚和违犯的人，不配承继，只配受罚和被弃。那么会发生什么？信德来到，带来了恩宠，从而使恩许生效。因为哪里有恩宠，哪里就有赦免；哪里有赦免，哪里就没有惩罚。惩罚既被除去，义德从信德而来，我们就毫无阻碍地成为恩许的继承人。\par

第16节。他说：\Quote{因此，它是出于信德，}\Quote{为的是恩宠，使恩许对一切后裔都是确定的。}\BibleRef{罗 4:16}\par

你看，信德不仅建立了法律，也建立了天主的恩许，不容许它落空。反之，法律若不合时宜地遵守，甚至会使信德落空，并阻碍恩许。由此他显示信德非但不是多余的，反而是如此必要，以致没有它就不能得救。因为法律激起义怒，因为众人都违犯了它。但信德根本不让义怒产生：因为他说：\Quote{哪里没有法律，哪里就没有违犯。}你看到他是如何不仅消除了已经存在的罪，甚至不允许它产生吗？这就是为什么他说\Quote{出于恩宠。}为了什么目的？不是为了使他们蒙羞，而是为了使恩许对一切后裔都是确定的。在这里他列出了两个福分：恩赐是确定的，并且是给一切后裔的，从而聚集外邦人，并显示犹太人若反对信德，就被排除在外。因为信德比法律更确定。信德不会伤害你（不要争辩），甚至现在你处于法律的危险之中，它保护你。接着他说了\Quote{对一切后裔}之后，他定义了他指的是什么后裔。他说：那出于信德的，从而将他们的血缘与外邦人混合，显示那些不像亚巴郎那样相信的人不应以他为傲。看信德还成就了第三件事：它使与这位义人的关系更加明确（\OriginalTerm{更明确}{\GreekText{ἀκριβεστέραν}}），并使他成为更众多子嗣的祖先。这就是为什么他不只说亚巴郎，而是说\Quote{我们的父亲，}我们信者之父。然后他又用证词印证了他所说的话——\par

第17节。他说：\Quote{正如经上所载：}\Quote{我已立你为万民之父。}\BibleRef{罗 4:17}\par

你注意到这是自远古由上智安排的吗？那么，他说这话是因依市玛耳人、阿玛肋克人、或哈加尔人吗？然而，在继续中，他更明确地证明并非指这些人。但此时他仍推进到另一点，借此他通过定义关系的方式，并以宽广的心智确立了这一点。那么他说什么？\par

\Quote{在祂（或：面对，\OriginalTerm{在……面前}{\GreekText{κατέναντι}}）所信的天主面前。}\par

但他的意思大致如此：正如天主不是一部分人的天主，而是众人的父，他也是如此。再者，正如天主不是藉本性的关系而为父，而是藉信德的结合而为父，他也是如此，因为正是服从使他成为我们众人的父。既然犹太人轻视这种关系，如同依附那粗俗的关系一样，他便藉将话题提升到天主而指明这种关系更为真实。与此并行，他清楚表明这就是他所获得的信德的赏报。因此，如果并非如此，而他是地上所有居民的父，那么\Quote{之前（或回应）}一词就不合适，而天主的恩赐也会被削减。因为\Quote{之前}相当于\Quote{与……一样}。请问，一个人成为自己亲生的子女之父，这有什么可惊奇的？这是每个人都会遇到的。但不可思议的是，那些按本性他并不拥有的人，却藉天主的恩赐获得了。因此，你若相信圣祖受到尊荣，你也当相信他是众人的父。但说了\Quote{在他所相信的天主面前}之后，他并没有停在这里，而是继续如此说：\Quote{祂叫死者复活，又使那不存在的成为存在的}，这样预先为讨论复活奠定基础。这也适用于他当前的意图。因为如果祂能\Quote{叫死者复活}并使\Quote{那不存在的成为存在的}，那么祂也能使那些非由他所生的人成为他的子女。这也是为什么他没有说\Quote{使那不存在的成为存在的}，而是说\Quote{召唤}，以显示其更容易。因为正如我们很容易称呼那些存在的事物，对祂来说也很容易，甚至更容易，使不存在的事物得以存立。但在说了天主的恩赐伟大而不可言喻，并论及祂的能力之后，他又进一步表明亚巴郎的信德配得这恩赐，免得你以为他无缘无故受到尊荣。在引起听众的注意力以防止犹太人喧嚷和提出疑问，并说\Quote{那些不是子女的人如何能成为子女？}之后，他转而谈论圣祖，说道：\par

\BibleRef{罗 4:18}\BibleQuote{他在绝望中仍怀着希望而相信了，为使他成为万民之父，正如所吩咐的：\InnerQuote{你的后裔也要这样。}}\par

他怎么\Quote{在绝望中仍怀着希望而相信了}？这是在人性的绝望中，怀着属于天主的希望。（因为他在显示行为的崇高，并为所讲的话不容置疑留有余地。）相互矛盾的事，信德却将它们调和。但如果他说的是关于从依市玛耳所出的那些，这话就是多余的：因为那些不是藉信德，而是藉本性生的。但他也将依撒格呈现在我们面前。因为他所相信的不是关于那些外邦人，而是关于那要从他不生育的妻子所生的儿子。因此，如果作万民之父是赏报，那显然是针对他如此相信的那些外邦人。为叫你知道他是在说他们，请听下文。\par

\BibleRef{罗 4:19}\BibleQuote{他在信德上并不软弱，反而想到自己的身体已死。}\par

你看见他如何既指出障碍，又指出义人超越一切的高迈精神吗？\Quote{在绝望中}，他说，那所预许的就是如此：这是第一个障碍。因为亚巴郎没有其他以这种方式得了儿子的人可以效法。他以后的人效法他，但他却无人可效法，只唯独仰望天主。正因为此，他说\Quote{在绝望中}。然后，\Quote{他的身体已死}。这是第二个。\Quote{撒辣的子宫已死}。这是第三个，也是第四个障碍。\par

\BibleRef{罗 4:20}\BibleQuote{但面对天主的恩许，他没有因不信而怀疑。}\par

因为天主既没有给出任何证据，也没有显示任何征兆，只有赤裸裸的话语许诺本性不存任何希望的事。然而他仍然说：\Quote{他没有怀疑}。他不是说\Quote{他没有不信}，而是说\Quote{他没有怀疑}，即他既没有怀疑也没有犹豫，尽管障碍如此巨大。由此我们学到：如果天主应许无数不可能的事，而听者却不接受，那不是事物本性的过错，而是不接受者的无理。\Quote{反而在信德上坚强。}请看保禄的坚持。因为既然这篇讲论是关于行事的人和相信的人，他便表明相信的人比另一人更勤劳，需要更大的力量和巨大的毅力，并且承受非常人所能承受的辛劳。因为他们认为信德无价值，好像其中没有辛劳。因此，他坚持这一点，表明不仅是在节制或其他此类德行上成功的人，而且表现信德的人也需要更大的力量。因为正如前者需要力量以击退不节制的思虑，信德的人也需要刚强的灵魂，以推开不信的暗示。那么他如何变得\Quote{坚强}？他答道：通过将事情托付给信德而不是思虑：否则他就跌倒了。但他如何在信德本身上茁壮成长？他说：通过将光荣归于天主。\par

\BibleRef{罗 4:21}\BibleQuote{且深信天主所应许的，也必能成就。}\par

因此，戒绝好奇的质询就是光荣天主，而沉溺其中则是干犯。然而，若我们因涉足好奇的问题、探究下界之事而不能光荣祂，那么，若我们在主的受生之事上过于好奇，我们更将因自己的狂妄而受极大的惩罚。因为，倘若复活的预像不可深究，那么那些不可言喻、令人敬畏的主题就更不必说了。而他不仅用了\Quote{相信}一词，还用了\Quote{完全确信}。因为信仰正是如此，它比理由的论证更清晰，也更令人信服。因为随后产生的另一个理由不可能动摇它。那被言辞说服的人，也可能被言辞改变他的信服。但那依靠信仰的人，从此便坚固了自己的听觉，免受可能伤害信仰的言辞影响。既然他说了他是因信德而成义，他便显示了他是因那信德而光荣了天主；这特别属于善生之事。因为，\BibleQuote{照样，你们的光也当在人前照耀，好使他们看见你们的善行，光荣你们在天之父。}\BibleRef{玛 5:16}但是看哪！这也被证明属于信德！再者，正如善行需要能力，信德也需要能力。因为在善行中，身体常分担劳苦；但在信德中，善行只属于灵魂。因此，劳苦更大，因为它无人分担挣扎。你看到他是如何显示原属于善行的一切，以远大的程度依附于信德吗？因为信德在天主面前有可夸耀之处——需要能力和劳苦——并且再次光荣天主？在他说了\Quote{祂所应许的，祂也能成就}之后，在我看来，他预言了未来的事。因为祂所应许的不仅是现在的事，也是未来的事。因为现在是未来的预像。因此，不相信是软弱、渺小和可怜心智的标志。因此，当有人以信德指控我们时，让我们反以不信指控他们，认为他们可怜、心胸狭隘、愚蠢、软弱，品性不比驴好。因为正如相信属于高尚和高贵的灵魂，不信则属于最无理性、最无价值的灵魂，并且是昏沉（\OriginalTerm{昏沉}{\GreekText{κατενηνεγμένης}}）地沉入畜类的无感觉。因此，让我们离弃这些（行为），效法圣祖，正如他光荣了天主那样光荣天主。那么，\Quote{光荣了祂}是什么意思？他心中思念祂的威严、祂的无限能力。而且，在形成了对祂的正确认识后，他也\Quote{完全确信}祂的应许。\par

那么，我们也要因信德并因善工光荣祂，好使我们也能获得被祂光荣的赏报。祂说：\BibleQuote{光荣我的，我必光荣他}\BibleRef{撒上 2:30}；而且，即便没有赏报，单单光荣天主这件事本身就是一种光荣。因为，如果人们仅仅因为能称颂君王便引以为傲，即使没有其他成果；试想，我们的主被我们光荣，这是何等荣耀的事；同样，因我们的缘故导致祂被亵渎，又是何等大的惩罚。然而，祂愿意这种光荣为我们而发生，因为祂自己并不需要它。因为你认为天主与人之间的距离有多大？像人与蛆虫之间的距离吗？或者像天使与蛆虫之间的距离？但即使我提到如此巨大的距离，我也完全没有表达出来：因为要表达其伟大是不可能的。那么，你爱慕光荣，会愿意在蛆虫中拥有显赫的声誉吗？当然不会。既然你爱慕光荣尚且不愿如此，那远离这种情感、比我们高出无限多的祂，又怎么会需要你的光荣呢？然而，尽管祂没有这种需要，祂仍说祂渴望它，是为了你们。因为，如果祂为了你们而甘愿成为奴隶，那么祂同样为了这个缘故而要求其他细节，又有什么奇怪呢？因为任何有益于我们得救的事，祂都不会认为有损于祂的尊贵。既然我们知道了这一点，就让我们完全避开罪恶，因为罪恶使祂受亵渎。因为经上说：\BibleQuote{当逃避罪恶，如同逃避蛇的面；你若走近它，它必咬你}\BibleRef{德 21:2}：因为不是它来就我们，而是我们投靠了它。天主如此安排，使魔鬼不能以强制（希腊文：暴政）胜过我们：否则，无人能抵挡它的权势。因此，祂把它安置在遥远的居所，如同强盗和暴君。除非它发现一个人毫无防备且孤身一人，它才敢袭击他。除非它看见我们行经旷野，它才敢靠近我们。但魔鬼的旷野和居所无非是罪恶。那么，我们需要信德的盾牌、救恩的头盔、圣神的剑，不仅为了不受虐待，而且一旦它想扑向我们，我们就能砍下它的头。我们需要不断的祈祷，使它被踏在我们脚下，因为它无耻而鲁莽，即使它从下面作战。但即便如此，它却获胜了：原因是我们没有认真立志要高于它的攻击。因为它甚至没有能力将自己举得很高，而是沿着地面爬行。蛇就是它的预象。既然天主起初已将它置于那样的地位，如今更会如此。但如果你不知道\BibleQuote{从下面}\BibleRef{若 8:23}是什么意思，我也将试着向你解释这种战争的方式。那么，这\BibleQuote{从下面}\BibleRef{若 8:23}的作战是什么呢？就是站在世界的低处来攻击我们，诸如快乐、财富和今世的一切美物。因此，凡它看见飞向天堂的人，首先，它甚至无力跳那么远；其次，即使它尝试，也会迅速坠落。因为它没有脚，不要怕；它没有翅膀，不要怕。它在地上爬行，属于地上的事。那么，你不可与地上的事有份，你就甚至无需劳力。因为它不擅长公开作战：而是像蛇一样藏在荆棘中，常常潜伏在\BibleQuote{财富的迷惑}\BibleRef{玛 13:22}里。如果你砍去荆棘，它就会被发现而轻易被赶走；如果你知道如何用天主所感发的咒语来迷惑它，它就会立刻被击倒。因为我们是拥有，确实拥有属神的咒语，即我们的主耶稣基督的圣名和十字架的德能。这咒语不仅能把蛇从藏身处引出，抛入火中\BibleRef{宗 28:5}，而且还能医治伤口。但若有人念了这圣名却没有痊愈，那是由于他们自己的信德小，并非他们说的话语有什么无能。因为有些人拥挤耶稣、紧靠祂\BibleRef{路 8:44}，却没有得到好处。但那患血漏的妇人，甚至没有触摸祂的身体，只触摸祂衣服的繸头，就止住了长久以来的血漏。（参考圣奥思定\WorkTitle{讲道集}LXII.iii.4，p.124 O.T.）这圣名使魔鬼、情欲和疾病都恐惧。因此，我们要以此为乐，以此武装自己。保禄就是这样变得伟大，而他与我们一样具有相同的本性，众门徒的整个歌咏团也是如此。但信德使他成为完全不同的人，而且信德在他们身上如此丰盈，甚至连他们的衣服都有极大的能力\BibleRef{宗 19:12}。那么，如果那些人的影子和衣服都能驱走死亡\BibleRef{宗 5:15}，而我们的祈祷却连压制情欲都做不到，我们该当何罪？原因何在？我们的气质相差甚远。因为本性所赐予的，我们与他们一样拥有。因为他出生、成长都和我们一样，住在世上、呼吸空气，也和我们一样。但在其他方面，他比我们伟大得多、好得多：在热忱、信德和爱德上。那么，让我们效法他。让我们让基督藉我们说话。祂比我们更渴望如此：为此，祂预备了这乐器，不愿它闲置无用，而希望常常使用它。那么，你为何不让它适合造物主的手，反而使它松驰，因纵欲而松懈，使整个竖琴不适合祂使用？你本应使它的弦紧紧绷直，用属神的盐来坚固。因为如果基督看到我们的灵魂如此调准，祂也会藉它发出祂的乐声。当这事发生时，你将会看到天使们跳跃（\OriginalTerm{跳跃}{\GreekText{σκιρτῶντας}}），总领天使们亦然，以及革鲁宾们。\par

那么，让我们成为配得上祂无玷双手的人。让我们邀请祂甚至触动我们的心。因为祂并不需要任何邀请。只要使心灵配得上那触动，祂就会率先奔向你。因为若考虑到他人尚未达到的成就，祂就奔向他们（因为当\Person{保禄}尚未如此进步时，祂已为他编织了那赞美），当祂看见一人完全预备好，有什么事祂不会做呢？但如果基督发言，圣神确实降临在我们身上，我们将比天堂更好，不是把日月安置在我们身体内，而是让日月和天使的主居住并行走在我们内。我说这话，不是为了我们能复活死人、洁净癞病人，而是为了我们能显出一种比这一切更大的奇迹——爱德。因为无论这光荣之物在哪里，圣子就与圣父一同居住，圣神的恩宠也常临。因为经上说：\BibleQuote{哪里有两个或三个人因我的名聚在一起，我就在他们中间。} \BibleRef{玛 18:20}现在这是出于深切的情感，也为那些非常亲密的朋友，使他们所爱的人在两边。那么，他意思是，谁这样可怜，不希望基督在他们中间呢？就是那些彼此不和的我们！也许有人会嘲笑我问：你这是什么意思？难道你没有看见我们都在同一墙内，在教会的同一围栏下，在同一羊栈内一致站立；没有人争斗，我们在同一牧者之下，共同高呼，共同聆听所说的话，共同献上祈祷——而你却提到争斗和不和？我确实提到争斗，我不是疯了，也不是失去理智。因为我看见我所看见的，我知道我们在同一羊栈、同一牧者之下。然而正因如此，我更加哀叹，因为尽管有这么多因素吸引我们在一起，我们却不和。有人会说：你在这里看见什么叛乱？我确实没有看见。但当我们散会后，某人指责另一个人，另一个人公开侮辱，又一个人心怀嫉妒，又一个人欺诈、贪婪、暴力，又一个人沉溺于非法之爱，又一个人编织无数诡计。如果可能打开你们的灵魂，那么你们就会清楚地看见一切，并知道我没有疯。你们没有看见在军营里，当和平时，人们放下武器，毫无防备地穿过敌营；但当他们被武器、卫兵和哨所保护时，夜晚都在警戒，火不断燃烧，这种状态不再是和平而是战争吗？现在这就是我们中间所见的。因为我们彼此提防、彼此害怕，每个人都在邻人耳边说话。如果我们看见别人在场，就闭口不言，吞下所有要说的话。这不像有信心的人，而像严格戒备的人。\Quote{但我们做这些事（有人会说）不是为了作恶，而是为了逃避受恶待。} 是的，我为此悲伤，因为我们生活在弟兄中，却需要提防被恶待；我们点燃这么多火，设置卫兵和哨所！原因是虚妄盛行、狡猾盛行、爱德普遍分离，以及无休止的战争。由此可以找到一些人，对外邦人（希腊人）比对基督徒更有信心。然而，我们对此应当何等羞愧；我们应当何等哭泣哀悼！\Quote{那么，有人会说，我该怎么办呢？某某人脾气古怪，令人烦恼。} 那么，你的\OriginalTerm{宗教}{religion}（\OriginalTerm{哲学}{\GreekText{φιλοσοφία}}）在哪里？宗徒的法律在哪里？那法律命令我们彼此担代负担。\BibleRef{迦 6:2} 因为如果你不知道善待你的弟兄，你何时能善待陌生人？如果你没有学会如何对待你自身的肢体，你何时能吸引外人到你身边，并与他结合？但我该如何感觉？我极其烦恼，几乎流泪，因为我本可使我的眼流泪成泉（\BibleRef{耶 9:1}），正如那位先知所说，因为我看见平原上无数敌人，比他看见的更令人痛心。因为他看见外敌来攻击他们时说：\BibleQuote{我的内脏！我的内脏疼痛！} \BibleRef{耶 4:19} 但当我看见人们在一首领下排列，却彼此对立，撕咬撕裂自己的肢体，有些为金钱，有些为荣耀，还有些随便地嘲笑、戏弄、用无数方式伤害彼此，而且尸体比战争中的更糟，只剩下弟兄的空名，我自己感到无力设计任何适合如此灾难的哀歌！该敬畏，啊，敬畏这我们众人所共食的祭台！\BibleRef{格前 10:16} 基督，那位为我们而被杀的牺牲，被置于其上！\BibleRef{希 13:10} 强盗一旦与人同食盐，就不再是对同食者行抢；那餐桌改变了他们的性情，使比野兽更凶猛的人变得比羔羊更温顺。但我们虽分享这样的祭台、同享那样的食物，却武装起来彼此敌对，而我们应该武装起来对抗那向我们所有人作战的——魔鬼。然而这就是为什么我们日渐软弱，他日渐强壮。因为我们没有联合起来防御他，反而与他一同彼此对立，并利用他作为如此敌对阵列的指挥官，而我们应该只与他战斗。但现在放过他，我们只向弟兄拉弓。什么样的弓？你会说。就是舌头和嘴巴的弓。因为不仅标枪和箭矢，言语也更锋利得多，造成伤口。我们怎样才能结束这场战争？有人会问。如果你意识到，当你诽谤你的弟兄时，你是在从嘴里吐出污泥；如果你意识到你是在毁谤基督的一个肢体，是在吃你自己的肉（\BibleRef{咏 26:2}），是在使你定下的审判更加苦涩（它本是可怕而不可腐坏的），那箭矢杀死的不是被击中的他，而是发射它的你自己。\par

但他可能得罪了你，或许伤害了你？为此叹息，不要谩骂。哭泣，不是为你所受的委屈，而是为他的丧亡，正如你们的主也曾为\Person{犹达斯}哭泣，不是因为自己将要被钉十字架，而是因为他是一个负卖者。他侮辱了你，咒骂了你？为他祈求天主，好使天主快快与他修和。他是你的弟兄，他是你身上的肢体，与你承受同样的产痛，他已被邀请到同一祭台。但有人会说：他只对我发动新的攻击。那么，你的赏报就更大。正是为此，我们有最好的理由平息愤怒，因为他受到了致命的创伤，魔鬼击伤了他。你不要再给他一击，也不要和他一同跌倒。因为只要你站着，你就有办法拯救他。但如果你以侮辱相报，自己跌倒了，那么谁能扶起你们两人呢？受伤的人吗？不，他既已倒下，不能。与你一同跌倒的你吗？你连自己都站立不住，又怎能伸出手去帮助别人呢？所以，你现在要刚强地站立，拿起盾牌保护自己，趁他死了，以你的忍耐把他从战场上拖开。忿怒击伤了他，你不要再击伤他，反而要把那第一支箭拔出来。因为如果我们这样彼此相处，我们所有人很快都会获得健康。但如果我们互相武装起来，那么就连魔鬼也不必费心毁灭我们了。因为一切战争都是邪恶的，尤其是内战。但这场战争比内战更严重，因为我们之间的权利大于公民权，甚至大于亲族权。古时，\Person{亚伯尔}的兄弟杀了他，流了亲族的血。但这次谋杀比那次更无法无天，因为亲族权利更大，死亡是更严重的邪恶。因为他伤害了身体，而你对准灵魂磨快了剑刃。\Quote{但你首先受了苦。} 是的，但受苦不是真正的受苦，造孽才是。现在想一想：\Person{加音}是杀人者，\Person{亚伯尔}是被杀者。那么死者是谁？是那死后喊叫的（因为祂说：\BibleQuote{你兄弟的血从地里向我喊叫}\BibleRef{创 4:10}），还是那活着时却战栗恐惧的人？他确实比任何死人都更可怜。你看见了吗？受冤枉更好，即使人甚至被杀害。要知道：行不义更糟糕，即使人强大到可以杀人。他击倒并杀死了他的兄弟，然而后者得了冠冕，前者受罚。\Person{亚伯尔}被谋害，无辜被杀，但他甚至死后还控告\BibleRef{若 5:45}，定罪并得胜；另一个虽活着，却哑口无言，羞惭，被定罪，适得其反。因为他看见他被人爱，就除灭他，以为能把他从爱中驱逐出去。但他反而使爱更加炽热，天主在他死后更加寻找他，说：\BibleQuote{你的兄弟\Person{亚伯尔}在哪里？}\BibleRef{创 4:9}因为你的嫉妒没有熄灭对他的渴望，反而燃烧得更旺。你杀死他没有减少他的尊荣，反而使之更丰富。在这之前，天主甚至使他服从于你，而自从你杀了他，他即使死了，也要向你复仇。我对他的爱是如此伟大。那么谁是被定罪的人？是惩罚者还是受罚者？是那个从天主那里享受如此大尊荣的，还是那个被交付给某种前所未有、意想不到的惩罚的？你活着时不畏惧他，（他会说）所以你死后要畏惧他。你拔剑时没有战栗，现在血流了，你要被不断的战栗抓住。他活着时是你的仆人，你对他毫不宽容。为此，他死后成了你惧怕的主人。亲爱的，默想这些事，让我们逃离嫉妒，熄灭恶意，以爱德彼此相待，好使我们收获由此而来的祝福，无论是今生还是来世，借着恩宠和爱人，等等。阿们。\par

\SourceRecord{Translated by J. Walker, J. Sheppard and H. Browne, and revised by George B. Stevens.；Nicene and Post-Nicene Fathers, First Series；Vol. 11.；Edited by Philip Schaff.；Buffalo, NY: Christian Literature Publishing Co.,；1889.；<http://www.newadvent.org/fathers/210208.htm>.}

% structure: p0001 source=h1 target=section reason=page-title

\section{罗马书讲道辞第九篇}

\BibleRef{罗 4:23}\par

\Quote{\InnerQuote{算为他的成义}这句话，不是单为他写的，也是为我们这些将来要算为成义的人写的，就是我们这些信那使我们的主耶稣基督从死者中复活的人。}\par

在论及亚巴郎的许多伟大事迹、他的信德、成义和在天主面前的荣耀之后，恐怕听众会说：\Quote{这于我们何干？因为成义的是他。}于是又使我们与圣祖亲近。属神言语的能力何等伟大！因为对于一个外邦人、一个最近才来归依、并未做过任何善工的人，他不仅说此人丝毫不逊于相信的犹太人（即作为犹太人），甚至也不逊于圣祖，反而——若必须说出这奇妙真理——远超过他。因为我们的出身如此高贵，以至于他的信德只是我们信德的预像。他并没有说：\Quote{若算为他，大概也会算为我们}以免使之成为推理问题，而是以神圣律法确凿的话语发言，使整件事成为圣经的宣告。因为他说，圣经之所以写下这事，岂不正是要让我们看见，我们也是以这方式成义的？因为我们所信的是同一位天主，所信之事也相同，即便对象不是同一个人。论述了我们的信德之后，他又提到天主对人不可言喻的爱，他在各方面都始终彰显这爱，将十字架置于我们眼前。现在他以下面的话使之显明：\par

\BibleRef{罗 4:25} \Quote{祂为了我们的过犯被交付，又为了我们的成义而复活。}\par

请看，在提及祂死亡的原因之后，他又使同一原因也成为复活的确证。他的意思是：祂为何被钉？并非因为祂自己的任何罪。这从复活可以清楚看出。因为祂若是个罪人，怎能复活呢？但若祂复活了，显然祂不是罪人。然而祂若并非罪人，又怎会被钉呢？——是为了别人。若是为了别人，那么祂当然复活了。为防你问：\Quote{我们虽负如此大罪，怎能成义呢？}他便指出那涂抹一切罪过的一位，既从亚巴郎因信成义，又从救主的苦难（我们由此得以脱免罪过）来证实他所说的。提及祂的死亡后，他也论及祂的复活。因为祂死的目的是要使我们不致负罪受罚，而是要向我们施恩。因此祂死而复活，为使我们成义。\par

\BibleRef{罗 5:1} \Quote{所以我们既因信德成义，就借着我们的主耶稣基督与天主和平相处。}\par

\Quote{和平相处}是什么意思？有人说：\Quote{我们不要再因固执引入法律而彼此不合。}但在我看来，他现在似乎是在论及我们的行为。因为在信德上说了许多之后，他曾将信德置于因行为而成的义德之前，以免有人以为他所说的是懈怠的理由，于是他说：\Quote{让我们和平相处，}也就是说，不要再犯罪，也不要回到我们先前的情况。因为那是在与天主作战。\Quote{怎能不再犯罪呢？}有人说。\Quote{从前的事怎能做到呢？}因为若我们负如此多罪却由基督全得解放，那么借着祂，我们就更能常保我们所处的状态了。因为接受那本来不存在的和平，与保守那已赐予的和平，并非同一回事。因为获得确实比保守更难。然而，更难的却已变得容易，并已实现。那么，较易的事我们也会轻易成功，只要我们紧靠那位也为我们成就了更难之事的主。但我认为，他在这里不仅暗示易行，也暗示合理。因为若我们在公开作战时，祂就使我们与祂和好，那么按理我们应停留于和好状态，并将这报偿献给祂，免得祂似乎既已使不顺从、麻木不仁的受造物与圣父和好，却又被辜负了。\par

\BibleRef{罗 5:2} \Quote{借着祂，我们也因信德得以进入这恩宠中。}\EditorialNote{七份手抄本加有\Quote{进入这恩宠中}等字。}\par

若当时我们远离时，祂已使我们亲近祂，如今我们既已亲近，祂必更保守我们。并请你们思考，他如何处处设定这两点：祂的部分和我们的部分。然而在祂的部分，事物是多样、繁多且不同的。因为祂为我们死了，又进一步使我们和好，带领我们归向祂，并赐予我们说不尽的恩宠。但我们只献上了信德作为我们的贡献。因此他说：\Quote{因信德，进入这恩宠中。}这恩宠是什么？请告诉我。这是堪当被称为认识天主、从错误中被强行带出、认识真理、获得藉圣洗而来的一切福佑。因为祂使我们亲近我们的结局，正是要我们领受这些恩赐。因为我们和好，不只是为使我们获得简单的罪赦，而是为使我们领受无法计量的益处。祂甚至在这一切上并未止步，还许诺了其他，即是那些超越悟性和言语的不可言喻的福佑。这就是为什么他也将这两者都记下来。因为藉着提到恩宠，他清楚地指出我们目前所领受的；但藉着说：\Quote{我们在天主光荣的望德中也夸耀，}他揭示了整个未来之事。他也很好地说：\Quote{我们也站在这恩宠中。}因为这就是天主恩宠的本性：它没有终结，不知界限，反而总是向更伟大的事物前进；在人间之事中却非如此。举个我所说的例子：一个人获得了统治、光荣和权柄，却不能持续站立在其中，而是迅速被逐出。即使人不从他手中夺走，死亡也会来临，必然从他手中夺走。但天主的恩赐并非如此：因为人、时机、事态的危急、甚至魔鬼或死亡，都不能来将我们从其中赶出。相反，当我们死后，严格来说我们才拥有它们，并继续享受越来越多。因此，如果你对未来之事怀疑，从现在之事和你已领受的，也相信其他之事。因为这就是他所说的原因：\Quote{我们在天主光荣的望德中也夸耀（\OriginalTerm{夸耀}{\GreekText{καυχώμεθα}}）}，使你学习，信友应当拥有何种灵魂。因为不仅为已赐予的，也为将赐予的，我们应当充满信心，如同已经赐予一样。因为一个人在已赐的事上\Quote{夸耀}。既然未来之事的望德与已赐之事同样确定明白，他说在那事上我们同样\Quote{夸耀}。为此原因，他也称它们为光荣。因为如果它有助于天主的光荣，必然会发生，即使不是为我们，也为祂而成就。我为何要说（他意指）未来的福佑堪当被夸耀（\OriginalTerm{夸耀}{\GreekText{καυχήσεως}}）？因为甚至当前时代的恶事也能照亮我们的面容，并使我们从中找到安息。因此他也加上，\par

\BibleQuote{不但如此，我们在患难中也夸耀。}\BibleRef{罗 5:3}\par

试想，未来的事物何等伟大，就连在看似痛苦的事上，我们也能振奋；天主的恩赐如此伟大，任何其中的不快都微不足道！因为在外在福乐的情况下，为它们奋斗带来麻烦、痛苦和烦恼；而花冠和奖赏才带来快乐。但在此却非如此，因为摔跤对我们来说与奖赏一样有滋味。因为在那些日子里，有各样的诱惑，王国也存在望德中，恐怖迫在眉睫，而好事只在期望中；这使较为软弱的人气馁，甚至在他们未得花冠之前，他现在就给了他们奖品，藉着说我们应当在\Quote{患难中也要夸耀}。他所说的不是\Quote{你们应当夸耀}，而是我们夸耀，以他自身给予他们鼓励。接着，既然他所说的显得奇怪而矛盾——如果一个人在饥荒中挣扎，被锁链、折磨、侮辱、虐待，却应当夸耀，于是他继续确认这一点。而且（更重要的是），他说患难本身堪当被夸耀，不仅为了未来之事的缘故，也为了自身之故。因为患难本身是一件好事。如何如此？因为它们膏抹我们，使我们能忍耐。因此，在他说我们在患难中夸耀之后，他加了理由，用这些话：\BibleQuote{知道患难造就忍耐。}\BibleRef{罗 5:3}再注意保禄论辩的精神，他如何把他们的论点转向相反的方向。因为正是患难最使他们放弃对未来之事的望德，并使他们沉入沮丧，但他说这些正是产生信心的原因，也是不沮丧未来之事的原因，因为\Quote{患难}，他说，\Quote{造就忍耐。}\par

\BibleQuote{忍耐生老练，老练生望德，而望德不令人蒙羞。}\BibleRef{罗 5:4-5}\par

磨难非但不反驳这些望德，反而证实了它们。因为在将来的事物实现之前，磨难带来一个非常大的果实——忍耐，并使那受试炼的人成为有经验的。它也在某种程度上为将来的事物贡献，因为它在我们内注入望德的活力，因为没有什么像良好的良心那样使人渴望祝福。没有一个度正直生活的人对将来的事缺乏自信，而那些疏忽的人中，许多因不良良心的重负，希望既没有审判也没有报应。那么如何？我们的福乐在于望德吗？是的，在于望德——但不是单纯人的希望，这种希望常常溜走，使那寄望者蒙羞；当某个预期帮助他的人死了，或虽然活着却改变了时。我们没有这样的命运：我们的望德是确实而稳固的。因为那作出应许者永远活着，而我们这些将享有它的人，即使死了，也要复活，绝对没有任何东西能使我们因随便得意、毫无目的地寄望于不稳固的希望而蒙羞。他既以这些话充分消除了他们所有的疑虑，就不让自己的论述停留在现时，而是再次催促将来，知道有性格较弱的人，他们也寻求现世的利益，并对上述的不满足。于是他提供了一个从已赐祝福而来的证明。免得有人说：但假使天主不愿赐给我们呢？因为祂能，且祂永存活着，我们都知道：但我们如何知道祂也愿意如此做呢？从已经完成的事。\Quote{什么事完成了？}是祂为我们显示的爱。做什么呢？有人会说。在赐予圣神。因此，在说了\BibleQuote{望德不会令人蒙羞}\BibleRef{罗 5:5}之后，他接着这样证明：\par

\BibleQuote{因为天主的爱，}他并没有说\Quote{赐予}，而是说\BibleQuote{倾注在我们心中}，以此显示它的丰富。然后，那最大的恩赐，祂已赐下；不是天、地、海，而是比其中任何更宝贵的，并使我们从人成为天使，甚至成为天主的子女、基督的弟兄。但这恩赐是什么？是圣神。如今，若祂不愿意在我们劳作后以大冠冕赏报我们，祂绝不会在劳作前赐予我们如此伟大的恩赐。但现在，祂爱的温暖由此显明，祂不是逐渐地、一点一点地尊荣我们，而是倾注了祂全部祝福的源泉，且是在我们的争战之前。所以，若你不够相称，不要沮丧，因为你拥有那作为你审判者的爱，作为你强有力的辩护者。为此，他本人通过说\BibleQuote{望德不会令人蒙羞}\BibleRef{罗 5:5}，将一切归因于天主的爱，而非我们的善行。但在提及圣神的恩赐后，他又转向十字架，这样说：\par

6-8节。\BibleQuote{当我们还在软弱的时候，基督在预定的时间为不虔敬的人死了。因为为义人死，是很少有的；为善人或许有人敢死。但是天主向我们显示祂的爱。}\BibleRef{罗 5:6-8}\par

现在，他所说的类似这样：因为若为有德之人，没有人会轻易选择死亡，那么当你的主的爱，显明祂不是为有德之人，而是为罪人和仇敌被钉在十字架上时，请思量——他在此之后也说了：\BibleQuote{当我们是罪人时，基督为我们死了}\BibleRef{罗 5:8}\par

9-10节。\BibleQuote{那么，现在我们既因祂的血成义，更要藉着祂免于天主的义怒。因为，如果我们还在为仇敌的时候，藉着祂圣子的死，与天主和好了，那么，在和好之后，我们更要藉着祂的生命得救了。}\BibleRef{罗 5:9-10}\par

他所说的话看似重复，但仔细留意的人会发现并非如此。请思考。他希望就未来之事给与他们信心。首先，他通过义人的决断使他们感到羞愧，他说，他也\Quote{深信天主所应许的，祂也必能成就}；其次，通过所赐的恩宠；然后通过患难，足以引导我们进入希望；再次，通过我们领受的圣神。接着，通过死亡和我们从前的邪恶，他证实了这一点。正如我所说，他提及的似乎是一件事，但发现它变成了两件、三件，甚至更多。第一，\Quote{祂死了}；第二，那是\Quote{为不虔敬的人}；第三，祂\Quote{使我们和好、拯救、成义}，使我们不死，使我们成为子女和继承人。他说，我们不仅从祂的死亡本身获得强有力的保证，也通过祂的死亡赐予我们的恩赐。如果祂仅为像我们这样的受造物而死，那么祂所做的已是最大爱情的证明！但当祂同时死去，又赐予我们恩赐，而且是这样的恩赐，赐给这样的受造物，那么祂的作为便超越了我们最高的理解，并引导最迟钝的人进入信德。因为除了那在我们还是罪人时如此爱我们，甚至为我们舍弃自己的那一位，没有别人能拯救我们。你看到这个主题为希望提供了何等基础吗？因为在此之前，我们得救有两大障碍：我们是罪人，以及我们的得救需要主的死亡——这在发生前是难以置信的，并且需要极大的爱才能实现。但如今这已成为事实，其他条件就更容易了。因为我们成了朋友，不再需要死亡。那么，那位曾对敌人如此宽恕，甚至不舍弃自己儿子，如今他们成为朋友，祂还会不去保护他们吗？既然祂不再需要交出祂的儿子？因为一个人不拯救，要么是因为他不愿意，要么虽然他可能愿意，却无能为力。这些都不能用来谈论天主。祂愿意是显然的，因为祂交出了自己的儿子。祂也能，这也正是祂所证明的，因为祂使罪人成义。那么还有什么能阻止我们获得未来之事？没有！另外，当你听到罪人、敌人、软弱和不虔敬的人时，不要感到羞愧脸红，请听他说的话。\par

第11节。\BibleQuote{不但如此，我们现今既藉着我们的主耶稣基督获得了和好，也藉着祂而欢跃于天主。}\BibleRef{罗 5:11}\par

\Quote{不但如此}是什么意思？他意思是，我们不仅得救了，而且甚至因此荣耀，而有些人以为我们应当羞愧。因为像我们这样生活在极大邪恶中的人得以得救，正是我们被拯救者深深爱着的最好证明。因为拯救我们的不是天使或总领天使，而是祂的独生子本身。因此，祂拯救我们，且在我们处于如此困境时拯救我们，并藉着祂的独生子，不仅藉着独生子，还藉着祂的血，这为我们编织了无尽的荣耀冠冕。因为没有什么比被天主当作朋友对待（\OriginalTerm{被爱}{\GreekText{φιλεἵσθαι}}）、并且在爱我们的那一位（\OriginalTerm{爱}{\GreekText{ἀγαπώντα}}）中找到一位朋友（\OriginalTerm{爱人}{\GreekText{φιλεἵν}}），更算为荣耀和信心的事了。这使天使荣耀，使率领者和掌权者荣耀。这比王国更大，所以保禄将它置于王国之上。因为我也因此认为无形体的是有福的，因为他们爱祂，并在一切事上服从祂。先知也曾为此称赞他们：\BibleQuote{你们这些大有能力、执行祂话的。}\BibleRef{咏 102:20}因此依撒意亚也颂扬色辣芬，从他们接近那荣耀而彰显他们的卓越，这是最大爱情的表征。\par

让我们效法天上的权能，不仅渴望站立在宝座近旁，更渴望那坐在宝座上的那一位住在我们内。祂在我们恨祂时爱了我们，并且继续爱我们。\BibleQuote{因为祂使太阳上升，光照恶人，也光照善人；降雨给义人，也给不义的人。}\BibleRef{玛 5:45}\newline{}正如祂爱我们，你也当爱祂。因为祂是我们的朋友（\OriginalTerm{他是爱}{\GreekText{φιλεἵ} \GreekText{γὰρ}}）。有人会说：既然祂是我们的朋友，为什么还要以地狱、惩罚和报复相威胁呢？这完全是出于祂对我们的爱。因为祂所做的一切，祂所忙碌的一切，都是为了消除你的邪恶，用恐惧如同缰绳一般约束你倾向于恶的倾向，藉着祝福和痛苦将你从堕落中挽回，引领你归向祂，并使你远离一切比地狱更坏的恶行。但若你嘲笑这些话，宁愿终生生活在痛苦中，也不愿受一日惩罚，这并不奇怪。因为这只是你判断不成熟（\OriginalTerm{不成熟的判断}{\GreekText{ἀ} \GreekText{τελοὕς} \GreekText{γνώμης}}）、沉醉和无可救药的混乱的标记。正如小孩子看到医生要用火或刀，就尖叫着痉挛着逃跑跳开，宁愿让持续的溃疡侵蚀身体，也不愿忍受暂时的疼痛，之后得享健康。但那些有见识的人知道，生病比受刀更糟，同样，邪恶比受惩罚更糟。因为前者是得医治而健康，后者是毁坏体质而持续衰弱。健康优于衰弱，这对每个人都是显而易见的。那么，盗贼应当哭泣的时刻不是当他们肋旁被刺穿时，而是当他们凿穿墙壁、杀人时。因为若灵魂优于身体（事实也确实如此），当灵魂被毁时，更应当哀叹悲恸；但若人感觉不到，就更有理由哀哭。因为那些以不纯洁之爱去爱的人，比患激烈热病的人更值得怜悯，那些沉醉的人比受酷刑的人更值得怜悯。但有人会说：如果这些更痛苦，我们怎么会偏爱它们呢？因为人类中有许多人，如谚语所说，喜欢较坏的东西，他们选择这些，而错过了更好的。这在食物、政体、人生竞逐的目标、享乐、妻子、房屋、奴隶、土地以及所有其他事物上都可以看到。因为请问，与女人同居还是与男人同居，哪个更愉悦？与女人还是与牲畜？然而，我们仍会发现许多人撇开女人，与无理性的牲畜同居，并滥用男性的身体。然而自然的快乐大于不自然的快乐。但仍有许多人追求可笑、无趣且伴有惩罚的事，如同它们是快乐的一样。但有人可以说，这些事在他们看来就是如此。仅此一点就足以使他们可怜，因为他们以为那些不是快乐的东西是快乐的。因此，他们以为惩罚比罪恶更坏，而其实恰恰相反。然而，如果惩罚对罪人是坏事，天主就不会给恶上加恶；因为祂所做的一切都是为了消除恶，不会增加它。因此，受惩罚对犯罪的人并不是坏事，而是处在那种状态下不受惩罚才是坏事，正如病人不就医一样。（\Person{柏拉图}, \WorkTitle{Gorg.} 第 478 页以下）没有什么比放纵的欲望更坏。我说放纵，是指奢侈、不当的荣誉、权力，以及一般而言所有超出必要的事物。因为一个过著放荡生活的人，似乎是最幸福的，但实际上是最悲惨的，因为他给灵魂加上了严酷而暴虐的统治者。为此，天主使今生成为劳动的生活，为使我们摆脱那种奴役，进入真正的自由。为此，祂以惩罚相威胁，使劳动成为我们人生的一部分，从而约束我们的傲慢精神。同样，犹太人被束缚在粘土和制砖中时，立刻成为自律的，并不断呼求天主。但当他们享受自由时，却抱怨、激怒上主，并以无数邪恶刺透自己。那么，有人会说，对那些因患难而变坏的人，你作何解释？我要说，这不是患难的结果，而是他们自己软弱的结果。因为如果一个人胃弱，不能服用有泻药作用的苦药，反而因此更糟，我们不会指责药物，而是指责那部分的软弱；同样，这里也要指责性格的软弱。因为一个因患难而改变的人，更可能因松弛而受影响。如果在被夹板（或捆绑）时（这就是患难）他都失败了，那么当绷带除去时，他更会失败。如果在被束缚时他都改变了，那么在膨胀的状态（\OriginalTerm{膨胀的状态}{\GreekText{χαυνούμενος}}）下更会改变。有人会问：我怎样才能不被患难这样改变呢？如果你反省：无论你愿意与否，你将不得不承受那施加的事；但若你怀着感恩之心去做，你将因此大得益处；但若你对此愤怒、暴躁、亵渎，你并不会减轻灾祸，反而会使它的波浪更加汹涌。因此，让我们用这种心态，将必然之事变为自愿选择。我的意思是：假设有人失去了他的儿子，另一个人失去了全部财产。如果你反省：已发生的事无法挽回，但即使无法补救，你仍能从不幸中获得果实，即高尚地承受环境；如果你不说亵渎之言，反而向上主献上感恩之词，那么那些加在你身上的恶事，就会成为你自由选择的好行为。你看见一个儿子过早被带走吗？说：\BibleQuote{上主赐的，上主收回。}\BibleRef{约 1:21}你看见财产耗尽吗？说：\BibleQuote{我赤身脱离母胎，也要赤身归去。}\BibleRef{约 1:21}你看见恶人亨通，义人受苦，遭受无数患难，却不知原因何在？说：\BibleQuote{我像畜牲在你面前，但我是常与你同在。}\BibleRef{咏 72:22}但若你要探求原因，反省祂已定下审判世界之日，这样你就能摆脱困惑，因为那时每人将得到应得的报应，正如拉匝禄和财主一样。要想起宗徒们，因为他们也为受鞭打、被驱赶、遭受无数苦难而喜乐，因为他们\BibleQuote{算是配为这名字受辱。}\BibleRef{宗 5:41}那么，你若有病，就当高尚地承担，并承认自己欠天主这笔债，你将得到与他们同样的赏报。但如何在软弱和痛苦中能对上主心怀感激呢？你若真诚爱祂，就能。因为那三个被投入火窑的青年，以及其他在监狱和无数患难中的人，都不停地感恩，更何况在疾病中的人呢？因为，确实，没有什么能胜过强烈的爱。但当这爱是对天主的爱时，它便高于一切，无论是火、刀剑、贫穷、疾病、死亡，还是其他任何类似的东西，对于拥有这爱的人都不再可怕；他藐视这一切，飞向天堂，他的情感丝毫不逊于天上的居民，只看见那荣耀的唯一美善，看不见别的，无论天地海洋。今生的烦恼不会使他沮丧，美好而快乐的事也不会使他得意或自满。让我们以这种爱去爱（因为没有什么能与它相比），既为现在，也为将来。或者更确切地说，超过这些，为了爱本身的性质。因为我们将从今生和来世的惩罚中得释放，并享受天国。然而，逃脱地狱和享有天国，与将要说的相比，也算不了什么。因为超过这一切的是拥有基督，祂既是我们的爱人，又是我们的爱者。因为若这在人间发生时已超越一切快乐，那么当两者都来自天主时，什么言语或思想能描绘这灵魂的幸福呢？除了亲身经验，没有别的。那么，为使我们因经验而认识这属灵的喜乐、幸福的生活和不可言喻的宝藏，让我们放下一切，紧抱那爱，为我们的喜乐，也为天主的荣耀。因为光荣和权能归于祂，与祂的独生子，及圣神，从今直到永远，世世无尽。阿们。\par

\SourceRecord{Translated by J. Walker, J. Sheppard and H. Browne, and revised by George B. Stevens.；Nicene and Post-Nicene Fathers, First Series；Vol. 11.；Edited by Philip Schaff.；Buffalo, NY: Christian Literature Publishing Co.,；1889.；<http://www.newadvent.org/fathers/210209.htm>.}

% structure: p0001 source=h1 target=section reason=page-title

\section{罗马书讲道词第十篇}

\BibleRef{罗 5:12}\par

\BibleQuote{因此，正如罪藉着一人进入了世界，死亡藉着罪也进入了，于是死亡便临到（\OriginalTerm{传世读法}{\GreekText{διἥλθεν}}，六份抄本作\OriginalTerm{另一读法}{\GreekText{εἴς}}……）所有人，因为众人都犯了罪。}\par

最好的医生总是煞费苦心去发现疾病的根源，直捣祸患的源头，有福的保禄也是如此。因此，他在说过我们已成义，并从族长、圣神和基督的死（因为祂若不打算使我们成义，就不会死）加以证明之后，又从其他来源来确认他如此长篇论证的结论。他从相反的事物，即死亡和罪，来确认他的命题。怎么确认的呢？以什么方式呢？他探究死亡从何而来，又如何得势。那么死亡是怎样进来并得势的？\BibleQuote{藉着一人的罪。}那么，\BibleQuote{因为众人都犯了罪}是什么意思呢？这是指：他既然一度堕落了，就连那些没有吃那棵树果子的人，也都从他而成为必死的。\par

\BibleRef{罗 5:13}　\BibleQuote{因为没有法律以前，罪已经在世上，但没有法律，罪就不被追究。}\par

\BibleQuote{没有法律以前}这句话，有些人认为是指颁布法律以前的时代——例如亚伯尔、诺厄、亚巴郎的时代——直到梅瑟出生。那么照这样，那时有什么罪呢？有人说他指的是在乐园里的罪。因为直到那时，这罪（他要说）还没有被消除，但其果实仍然有效。因为它产生了那众人都分沾的死亡，这死亡得势并辖制了我们。那么，他为什么接着说\BibleQuote{但若没有法律，罪就不被追究}呢？那些站在我们这边的人说，他是以犹太人提出的异议为前提而说这句话的：也就是说，若没有法律就没有罪，那么法律以前的人是如何被死亡吞灭的呢？但在我看来，即将给出的含义更值得接受，也更符合宗徒的意思。这是什么含义呢？他说\BibleQuote{没有法律以前，罪已经在世上}，在我看来，他的意思是，法律颁布之后，因违反法律而生的罪得势了，并且只要法律存在，它就一直得势。因为他说，如果没有法律，罪就不能存在。那么，如果因违反法律而生的这个罪是产生死亡的原因，那么法律以前的人是如何死的呢？因为如果死亡源于罪，但若没有法律，罪就不被追究，那么死亡如何得势呢？由此可见，不是违反法律的这个罪，而是亚当的悖逆败坏了万物。有什么证据呢？就是法律以前的人都死了，因为他说：\BibleQuote{从亚当到梅瑟，死亡做了王，连那些没有像亚当一样犯罪的人也管辖了。}\par

它是怎样做王的？\BibleQuote{仿照亚当的悖逆，他是那要来的那一位的预像。}这就是亚当何以是基督的预像。怎样是预像呢？将会有人说。因为正如前者成为他后裔的原因——即使他们没有吃那棵树上的果子——他的吃导致了死亡的引入；同样，基督也成为他后裔的原因——即使他们没有行义——祂藉着十字架恩赐给所有人的那义，供应给他们。为此，他处处紧扣\BibleQuote{一人}，并不断将它摆在我们面前，当他说：\BibleQuote{正如罪藉着一人进入了世界}——以及\BibleQuote{若因一人的过犯，众人都死了}；\BibleQuote{但恩惠和过犯不同}；\BibleQuote{审判是因一人的过犯而成了判定}；又再说：\BibleQuote{若因一人的过犯，死亡藉着一人做了王}；以及\BibleQuote{因此，正如因一人的过犯}。还有\BibleQuote{正如因一人的悖逆，众人都成了罪人。}所以他总是不放鬆\BibleQuote{一人}，以便当犹太人对你说：怎么会由于这一人基督的善行，世界就被拯救了呢？你就能对他说：怎么会由于这一人亚当的悖逆，世界就被判罪了呢？然而罪与恩宠是不相等的，死亡与生命是不相等的，魔鬼与天主是不相等的，其间有无限的距离。那么，既然从事物的本质、行事者的权能以及其本身的适宜性来看（因为拯救比惩罚更符合天主的本性），这优越性和胜利是在这一边的，告诉我，你还有什么话来说不信呢？然而，他在以下的话语中表明所做的事是合理的。\par

\BibleRef{罗 5:15}　\BibleQuote{但恩惠和过犯不同：若因一人的过犯，众人都死了，那么，天主的恩惠和那因耶稣基督一人的恩宠所赐与的恩惠，更加丰富地倾注到众人身上。}\par

他所说的意思大致是这样的：如果罪的效力如此广泛，而且仅是一人的罪，那么恩宠，而且是天主的恩宠——不仅圣父，还有圣子的恩宠——怎么能不更加丰盛呢？因为后者是远更合理的假设。因另一人而惩罚一人，似乎不太合理。但因另一人而拯救一人，则更适宜、更合理。既然前者发生了，后者就更可能了。因此，他根据这些理由显明了它的可能性和合理性。然而，他接着从以下内容又证明了它甚至是必然的。那么他如何证明呢？\BibleQuote{因一人的过犯，众人都死了}\par

\BibleRef{罗 5:16}　\BibleQuote{而且恩惠也和由于一人的犯罪而导致的后果不同：因为审判是由一人而被判定，但恩惠却是从许多过犯而成义。}\par

他所说的是什么呢？就是罪恶有能力带来死亡和定罪；但恩宠不但消除了那一个罪，也消除了随后而来的众罪。为避免\Quote{如同}和\Quote{同样}这些词可能使福乐与灾祸的尺度显得相等，并且为使你不致于听到亚当时就以为他所带来的那一个罪已经被消除，他说，因许多过犯而获得赔偿。这如何显明呢？因为在乐园中那件事之后所犯的无数罪恶，最终竟导致成义。但是，哪里有了义，就必然随之有生命和数不尽的福乐，如同哪里有了罪就带来死亡一样。因为义胜于生命，它甚至是生命的根源。于是他指出，当时带来的善果不止一个，被消除的也不止那一个罪，而是连同其余一切都被除去了。他说：\Quote{这恩赐因许多过犯而至成义。}其中必然包含一个证明：死亡也被连根拔除了。但是，由于他曾说第二者大于第一者，他就必须再为此事提供进一步的依据。因为之前他曾说，如果一个人的罪杀死了所有人，那么一个人的恩宠就更能拯救众人。此后，他表明恩宠不仅消除了那一个罪，也消除了其余一切罪，而且不仅消除了罪，还赐下了义。基督所行的善不仅与亚当所行的恶相等，而且远远更大、更善。既然他作了这样的声明，他就需要在此再次进一步证实这些。他如何证实呢？他说：\par

第17节。\BibleQuote{若因一人的过犯，死亡就藉着那一人而为主，那么，那些领受丰盛的恩宠和恩赐\EditorialNote{Field 依多数抄本作\Quote{和恩赐}}以及成义的人，将更要藉着一个人耶稣基督在生命中为王了。} \BibleRef{罗 5:17}\par

他所说的，大致如此。什么武装了死亡来对抗世界？只是一个人吃树上的果子。如果死亡一次过犯就获得了如此大的权柄，那么，当发现某些人领受了与那过犯完全不成比例的恩宠和义时，他们如何还会受死亡辖制呢？为此，他此处不说\Quote{恩宠}，而说\Quote{恩宠的丰满}。因为我们所领受的恩宠不仅仅是足以消除那罪的分量，甚至远远更多。因为我们不仅立刻从惩罚中得释，脱下了一切不义，且从上而生 \BibleRef{若 3:3}，与旧人一同埋葬而复活，蒙救赎、成义、被收纳为义子、被圣化、成为独生者的兄弟、同嗣业者和与祂同一身体，被算作祂的肢体，如同身体与头一样，我们竟与祂结合！这一切，保禄称之为恩宠的\Quote{丰满}，显示我们所领受的不仅是一剂足以抵消伤口的药物，更是健康、美丽、荣誉、荣耀和远超我们本性的尊严。所有这些，每一件本身都足以消除死亡，但当它们全部显然汇聚为一成时，死亡就丝毫痕迹也不存留，连影子也看不见，完全被消除了。正如有人将一个欠十枚小钱（\OriginalTerm{十枚小钱}{\GreekText{ὀ} \GreekText{βόλους}}）的人投进监狱，不仅如此，连他的妻子、儿女和仆人也因他而受累；另一个人来了，不仅付清那十枚小钱，还赐予一万金塔兰特，并将那囚犯带进君王的宫廷，到最高权柄的宝座前，使他分享最高的荣耀和一切华贵，那么债主就再也记不起那十枚小钱了。我们的情况正是如此。基督所偿付的远超我们所欠的，是的，就像无边的海洋多于一小滴水。所以，人啊，当你看到如此浩大的恩宠宝库时，不要犹疑，也不必追问那小小的死亡和罪恶的火星是如何被消除的，因为有如此恩赐的海洋涌来掩盖它。这就是保禄所说的\Quote{那领受了恩宠和成义的丰满者，将要在生命中为王}。现在他清楚证明了这一点，就再次运用之前的论点，重新拿起同一个词来加以强调，说如果因那一过犯所有人都受了惩罚，那么他们也同样可以藉这些手段而成义。于是他说：\par

第18节。\BibleQuote{因此，就如因一人的过犯，审判临到众人，以致定罪；同样，因一人的成义，恩赐临到众人，以致生命的成义。} \BibleRef{罗 5:18}\par

他再次坚持说：\par

第19节。\BibleQuote{因为正如因一人的悖逆，众人成了罪人；同样，因一人的服从，众人也将成为义人。} \BibleRef{罗 5:19}\par

他所说的似乎确实引起一个大问题；但只要有人仔细注意，这也容易解决。那么问题是什么？就是那句话：因一人的过犯，众人都成了罪人。因为当他犯罪并成为可朽的之后，出于他的人也成为可朽的，这并不奇怪。但怎么能从他的不顺服导致另一个人成为罪人呢？若如此，这样的人甚至不应受惩罚，因为如果他是从自己成为罪人的话。那么\Quote{罪人}这个词在这里是什么意思？在我看来，它指的是应受惩罚和被判处死刑。他此前已经清楚而详细地说明，因\Person{亚当}的死亡我们都成为可朽的。但现在的问题是，这样做有何目的？但他没有继续补充这一点，因为这对当前的主题没有贡献。因为争辩的对象是犹太人，他怀疑并嘲笑由一人而来的正义。因此，在表明惩罚也是由一人临到众人之后，他没有补充这样做的原因。因为他不为多余之事，只保持必要。因为辩论的原则并不要求他对犹太人多说；所以他留下它未解决。但如果你们中有人为了学习而询问，我们应该给出这样的回答：我们远没有受这死亡和惩罚的伤害，如果我们警醒，我们反而因成为可朽的而获益：第一，因为我们不在一个不朽的身体里犯罪；第二，因为我们获得了无数从事\OriginalTerm{灵修}{\GreekText{φιλοσοφίας}}的理由。因为节制、克己、谦卑、远离一切邪恶，是死亡因其临在和被期待而说服我们去做的事。但伴随着这些，甚至在此之上，它还引入了其他更大的祝福。因为殉道者的荣冠和宗徒的赏报都由此而来。\Person{亚伯尔}就这样成了义人，\Person{亚巴郎}也如此，因为他献上了自己的儿子；\Person{若翰}也如此，他因基督而被杀；三个青年也如此，\Person{达尼尔}也如此。因为如果我们这样警醒，不仅死亡，连魔鬼自己也不能伤害我们。此外，还有这一点可说：永生正等待着我们，在受到少许管束之后，我们将无忧无惧地享受未来的祝福，仿佛在今生的一所学校里，借着疾病、磨难、诱惑、贫穷和其他明显的恶事受训，以便我们配得承受来世的祝福。\par

20节。\BibleQuote{此外，法律插进来，为要使过犯增多。}\BibleRef{罗 5:20}\par

由于他已表明世界从\Person{亚当}被定罪，但从基督得救并脱离定罪，他现在适时地开始讨论法律，再次削弱对法律的高傲见解。他的意思是，法律远非带来好处或任何帮助，反而因它的介入使混乱加剧。但\Quote{使}一词再次不表示原因，而表示结果。因为赐下法律的目的并不是\Quote{为了}使过犯增多，因为赐下法律是为了减少和消灭过犯。但结果却相反，这不是由于法律的本质，而是由于接受法律之人的懈怠。但他为什么不说\Quote{法律被赐下}，而说\Quote{法律插进来}？这是为了表明法律的需要是暂时的，而非绝对或必须的。他也对\Place{迦拉达}人说了同样的话，以另一种方式表明同一件事。他说：\Quote{在信仰来临以前，我们被看守在法律之下，被圈入那将要启示的信仰。}因此，法律不是为自己，而是为别人看守羊群。因为犹太人有些粗鲁、软弱、对恩赐本身漠不关心，所以法律被赐下，为要更多地定罪他们，清楚地让他们知道自己的状况，并通过增加控诉来更抑制他们。但不要害怕，因为这样做不是为了加重惩罚，而是为了显明恩宠更为伟大。这就是他继续说的原因：\par

\Quote{但是，罪在那里增多，恩宠在那里也格外丰盛。}\par

他没有说\Quote{增多}，而是说\Quote{格外丰盛}。因为他赐给我们的不只是脱离惩罚，还有脱离罪和生命。就好像一个人不仅要使发烧的人痊愈，还要给他美貌、力量和地位；或者，不是只给饥饿的人食物，而是让他拥有巨大的财富，并安置他于最高的权位。罪如何增多呢？有人会说。法律给出了无数命令。由于他们违犯了所有命令，违犯便更加增多。你看恩宠和法律之间有多大的区别？因为法律增加了定罪，而恩宠则更加丰盛地赐与。在提到这不可言喻的慷慨之后，他再次讨论死亡和生命的根源。那么死亡的根源是什么？是罪。因此他又说：\par

21节。\BibleQuote{就如罪恶藉死亡为王，恩宠也藉着正义为王，使人藉着我们的主耶稣基督获得永生。}\BibleRef{罗 5:21}\par

他说这话，为表明后者（罪恶）居于君王之位，前者（死亡）如同士兵，受后者管辖，且被后者武装。既然如此，后者（即罪恶）武装了死亡，那么很明显的，毁灭死亡的正义——由恩宠引入的正义——不仅解除了死亡的武装，甚至摧毁了死亡，彻底消除了它的统治，因为正义是两者中更大的，它不是由人和魔鬼引入的，而是由天主和恩宠引入的，并引领我们的生命趋向更美好的状态和无限的祝福。因为正义永无止境（也由此让你看到它的优越性）。因为罪恶将我们逐出今生，而恩宠来到时，赐给我们的不是今生，而是不朽的永生。但这一切，基督是我们的担保。因此，若你拥有正义，就不要为你的生命担忧，因为正义大于生命，是生命之母。\par

第六章第一节。\Quote{这样，我们该说什么呢？我们要常留在罪恶中，好叫恩宠增多吗？断乎不可。}\BibleRef{罗 6:1}\par

他再度转向劝诫，却又不是直接引入，免得被许多人视为惹厌和烦扰，而是好像从教义中自然生发出来。因为，即使他如此变换表达方式，仍担心他们因他所说的话而感到冒犯，因而说：\Quote{我有些冒昧地给你们写了信，}\BibleRef{罗 15:15} 更何况他若不这样做，他们就更会觉得他过于严厉。既然他表明了恩宠的伟大是因它所医治的罪恶的伟大，正因如此，在无思想的人看来，这似乎成了对罪恶的鼓励（因为他们会说：如果显示更大恩宠的原因，是我们行了大罪，那么我们就不要停止犯罪，好让恩宠更显多）。为避免他们这样说或这样怀疑，请看他是如何再次反驳这个异议的。首先，他通过谴责来做。\Quote{断乎不可。}他惯于对众所公认荒谬之事这样做。然后，他提出一个无可辩驳的论证。是什么呢？\par

第二节。\Quote{我们这些已死于罪恶的人，如何还能在罪恶中生活呢？}\BibleRef{罗 6:2}\par

\Quote{我们已死}是什么意思？是指就罪恶而言，我们都已接受了死刑判决吗？还是指我们因相信而成为光明之子，从而向罪恶死了？我们更应说后者，因为下文清楚表明这是正确的。那么，向罪恶死是什么意思？就是不再在任何事上服从它。因为圣洗一次而永远地成就了这事，它使我们向罪恶死了。但此后，我们必须靠自己的努力持续维护这种状态，以致尽管罪恶向我们发出无数命令，我们永不再服从它，而像死人一样坚定不移。其实，他在别处说罪恶本身已经死了。在那里，他如此说是想表明德行是容易的（\BibleRef{罗 7:8}）。但在这里，他急切地想激发听者，就将死归到人这一边。接下来，由于所说的话有些隐晦，他再次解释，并用已经说过的话来责备他们。\par

第三、四节。\Quote{难道你们不知道，我的弟兄们，我们这受过圣洗归于基督耶稣的人，就是受圣洗归于他的死亡吗？所以我们借着圣洗已归于死亡与他同葬了。}\par

\Quote{受圣洗归于他的死亡}是什么意思？就是指以我们像他一样死去为目的。因为圣洗就是十字架。那么，十字架与埋葬之于基督，就是圣洗之于我们，尽管并非完全等同。因为基督自己在肉身中死了并被埋葬，但我们却向罪恶死了并被埋葬。所以他没有说：与他的死亡同扎根，而是说：与他死亡的相似。因为二者都是死亡，但主体不同：基督的死亡是肉身的，而我们的死亡是向罪恶，这是我们自己的。既然基督的死亡是真实的，那么我们的死亡也是真实的。既是真实的，那么我们自己这方面也必须有所贡献。因此他接着说：\par

\Quote{正如基督借着父的光荣从死者中复活了，同样，我们也该在生命的新度中行走。}\par

这里，他借着谨慎行走的义务，也暗示了复活的主题。怎样暗示呢？他的意思是：你相信基督死了，并复活了？那么也相信你自己也是如此。因为这与基督相似，因为十字架和埋葬也是你的。如果你分享了死亡和埋葬，那么更会分享复活和生命。因为如今那更大的——我指的是罪恶——已被除去，就不应再怀疑那较小的——死亡的除去。\par

但此事他暂且留给听众的良心去推理，而自己在向我们展示未来的复活之后，又向我们要求另一种复活，即新的生活，这生活是在现世通过习性的改变而实现的。当淫乱者变得贞洁、贪婪者变得仁慈、粗暴者变得温顺时，甚至在这里就已发生了复活，那是另一种复活的序曲。这如何是复活呢？因为罪恶被治死，正义复生，旧生活消失，而这新而天使般的生活正在活出来。但当你听到新生活时，要期待巨大的转变、宽广的变化。然而我眼含泪水，深深叹息，想到保禄要求我们多么大的宗教修持（\OriginalTerm{宗教修持}{\GreekText{φιλοσοφίαν}}），而我们却让自己陷入何等懈怠，在领洗后又回到从前的旧态，返回埃及，在玛纳之后仍记挂葱蒜。\BibleRef{户 11:5}因为在我们受光照的时刻，我们只改变十天或二十天，然后又重拾从前的行为。但保禄要求我们的不是固定天数，而是终生这样的生活。但我们却回到从前的呕吐，在恩宠的青年期之后，又建立起罪过的老年期。因为贪爱金钱，或不当的欲望奴役，或其他任何罪恶，都惯常使人衰老。\BibleQuote{那衰败和衰老的，即将消失。}\BibleRef{希 8:13}因为没有身体，确实没有身体，会像灵魂因许多罪恶而衰败摇晃那样因时间长久而瘫痪。这样的人被带到昏聩的极致，发出模糊的声音，如同极老而疯癫的人，被痰湿压得沉重，心智极度扭曲，健忘，眼上有鳞片，令人厌恶，易被魔鬼捕获。这就是罪人的灵魂；义人则不然，他们年轻、秀美，始终处于生命的盛期，随时准备任何战斗或拼搏。但罪人的灵魂即使受到微小震动，也会立刻跌倒而毁灭。先知曾表明这一点，他说：\BibleQuote{像糠秕被风吹散}\BibleRef{咏 1:4}，这样，生活在罪恶中的人也如此被飘摇，暴露于各种伤害。因为他们既不像健康人那样看，也不单纯地听，说话不清晰，反而被极度的气短所压迫。他们口中满是唾沫。但愿仅是唾沫，而非污秽之物！但如今他们发出的言辞比任何泥沼更臭，更糟的是，他们甚至没有能力将这唾沫般的言辞吐掉，反而用极大的猥亵把它接在手中，再涂抹回去，以至凝结而难以发散。也许你们听了这番描述感到恶心。你们岂不更应面对现实吗？因为如果这些发生在身体上是令人厌恶的，那么在灵魂中就更甚了。那儿子就是如此，他挥霍了全部家产，落到最悲惨的境地，比任何痴傻或混乱的人更虚弱。但当他愿意时，仅凭他的决意和转变，就突然变得年轻。因为他一说：\Quote{我要回到我父亲那里去}\，这一句话就带给他所有福分；或者更确切地说，不是空话，而是他在话中附加的行动。因为他没有说\Quote{让我回去}\，然后留在那里；而是说：让我回去，就回去了，走完了那整段路。让我们也这样做；即使我们被带到了边界之外，也要上到我们父亲的家里，不要因路途遥远而犹豫。因为如果我们愿意，回程是容易且非常迅速的。只要我们离开那陌生和异邦之地；因为罪恶就是这样，将我们远远带离父亲的房子；那么让我们离开她，好能迅速回到我们父亲的家中。因为我们的父亲对我们怀有自然的渴望，如果我们改变，他会像对待那些无玷者一样尊重我们，甚至更多，正如那父亲给予那儿子更大的尊荣。因为他自己在接纳儿子回来时得到了更大的喜悦。\Quote{我该怎样回去呢？}\有人会说。只要开始行动，一切就完成了。远离恶行，不再深入其中，你就已经抓住了全部。因为在病人身上，病情不再恶化就是好转的开始，在罪恶方面也是如此。不再深入，你的恶行就会终结。如果你这样做两天，第三天就会更容易保持；三天之后，你会增加十天，然后二十天，然后一百天，然后整个一生。\EditorialNote{参见：玛窦福音讲道集第十七篇，第267页，旧译版}因为你走得越远，就会看到道路越容易，并且你将站在顶峰本身，立刻享有许多好处。因为当荡子回来时，有笛声、竖琴、舞蹈、宴席和聚会；那个本可以责备他不合时宜的挥霍和逃到如此遥远地方的父亲，什么也没有做，反而视他如同无玷者，甚至不能使用责备的话语，或甚至向他提起从前的事，反而扑向他，亲吻他，宰了牛犊，给他披上袍子，赐予他丰盛的尊荣。让我们既然有这样榜样在前，就应振作精神，避免绝望。因为天主喜欢被称为父亲胜过主人，喜欢有儿子胜过有奴隶。这正是祂所喜爱的。这就是祂所做一切的原因；祂\BibleQuote{连自己的独生子也没有怜惜}\BibleRef{罗 8:32}，好使我们获得义子的名分，好使我们爱祂，不仅当作主人，也当作父亲。如果我们这样对祂，祂就以此为荣耀，并向所有人宣告，尽管祂不需要我们任何东西。例如，在亚巴郎的事上，祂处处这样做，用这些话：\Quote{我是亚巴郎、依撒格和雅各伯的天主。}\然而，应该是祂家中的人以此为荣；但现在显然是主这样做；因此祂对伯多禄说：\BibleQuote{你爱我比这些更深吗？}\BibleRef{若 21:17}，为了表明祂从我们这里寻求的莫过于此。为此祂也命亚巴郎献上自己的儿子，为要向所有人表明祂深受这位圣祖的爱。渴望被深深敬爱，源于深深去爱。为此祂也对宗徒们说：\BibleQuote{爱父亲或母亲超过我的，不配属于我。}\BibleRef{玛 10:37}为此祂命令我们甚至将与我们最密切的东西，即我们的灵魂\BibleRef{若 12:25}，置于爱祂之后，因为祂希望我们以完全的至诚爱祂。因为如果我们对一个人没有强烈感情，我们也不会强烈渴望他的友谊，即使他伟大尊贵；而当我们真诚热烈地爱一个人时，即使被爱的人地位卑微，我们仍将他的爱视为极大的荣耀。因此，祂自己也称被我们爱为荣耀，甚至为我们经受那些可耻的事\BibleRef{若 12:23}。然而，那些事仅是出于爱而成为荣耀。但无论我们为祂受什么苦，不仅是出于爱，也因我们所渴慕者的伟大和尊贵，那既有理由被称为荣耀，也确实是荣耀。那么，让我们为祂冒险，如同为最大的冠冕奔跑，不要认为贫穷、疾病、侮辱、诽谤，甚至死亡本身，在为我们受这些苦时是沉重和负担的。因为如果我们心态正确，我们由此获得最大的利益，正如从相反的事，如果我们心态不正，也一无所获。但想想：有人侮辱你、攻击你吗？他岂不因此使你警惕，给你机会变得像天主吗？因为如果你爱那谋害你的人，你就像那\BibleQuote{让太阳上升，光照恶人和善人}\BibleRef{玛 5:45}的天主。另一个人夺走你的钱财吗？如果你高尚地承受，你将获得与那些把所有财产施舍给穷人的人同样的赏报。因为经上说：\BibleQuote{你们欣然接受了家产的被抢，因为知道自己在天上有更美且常存的产业。}\BibleRef{希 10:34}有人辱骂你、毁谤你，不论是真是假，如果你温和地承受他的侮辱，他就为你编织了极大的冠冕；因为诽谤者也为我们提供了丰厚的赏报。因为经上说：\BibleQuote{欢喜踊跃吧！当人为了我而捏造一切坏话毁谤你们时，你们在天上的赏报是丰厚的。}\BibleRef{玛 5:12}而那说实话攻击我们的人也大有裨益，只要我们温和地承受所说的话。因为法利塞人毁谤了税吏，且说的是真话，但反而使税吏成了义人\BibleRef{路 18:11}。何须一一列举呢？因为任何人愿意去查看约伯的奋斗，就能准确了解这一切。这正是保禄所说：\BibleQuote{天主若为我们，谁能反对我们呢？}\BibleRef{罗 8:31}因此，正如通过勤奋，我们甚至从烦恼的事中获益，同样，通过懈怠，我们甚至从有利的事中也没有长进。因为犹达斯与基督在一起，告诉我，有什么益处呢？法律对犹太人有益吗？乐园对亚当有益吗？梅瑟对那些在旷野中的人有益吗？所以，我们应撇开一切，只看一点：如何妥善管理自己的资源。如果我们这样做，连魔鬼本身也绝不能胜过我们，反而因使我们警惕而使我们更加进步。保禄正是这样唤醒厄弗所人，描述他的凶猛。但我们却在睡觉、打鼾，尽管我们要与如此狡猾的敌人争战。如果我们意识到一条蛇卧在我们的床边，我们会竭力杀死它。但当魔鬼卧在我们的灵魂中时，我们却以为没有受害，安然躺卧；原因是，我们不用身体的眼睛看他。然而，正因为如此，我们更应该警觉和清醒。因为对于一个能感知的敌人，我们容易防范；但对于一个看不见的敌人，如果我们不时刻武装，就难以逃脱。而且，因为他并不打算公开作战（因为那样他很快就被击败），而是常常装作朋友的模样，注入他残忍恶意的毒液。他就是这样收买了约伯的妻子，披着自然亲情的外衣，给出那邪恶的建议。同样，与亚当交谈时，他装出一副关切和为他着想的样子，说：\BibleQuote{你们吃那果子的日子，你们的眼睛就开了。}\BibleRef{创 3:5}他也以此说服了依弗大，假借宗教的名义，要杀死自己的女儿，献上律法禁止的祭献。你看到他的诡计是什么，他的作战如何多变吗？那么你要警惕，用圣神的武器全面武装自己，准确了解他的计划，这样你既能避免被捉，也能轻易捉住他。因为保禄正是这样胜过了他，准确了解了这些。所以他说：\BibleQuote{因为我们不是不知道他的阴谋。}\BibleRef{格后 2:11}那么，让我们也勤奋学习并避开他的诡计，以便在战胜他之后，无论在此生或在来世，都被宣告为胜利者，并获得那纯净的福分，赖天主的恩宠和爱人，等等。\par

\SourceRecord{Translated by J. Walker, J. Sheppard and H. Browne, and revised by George B. Stevens.；Nicene and Post-Nicene Fathers, First Series；Vol. 11.；Edited by Philip Schaff.；Buffalo, NY: Christian Literature Publishing Co.,；1889.；<http://www.newadvent.org/fathers/210210.htm>.}

% structure: p0001 source=h1 target=section reason=page-title

\section{罗马书讲道集第十一篇}

\BibleRef{罗 6:5}\par

\BibleQuote{\Emph{因为我们若同栽于祂死亡的形像，也必同栽于祂复活的形像。}}\par

我先前曾有机会指出的一点，这里也要重提：他不断转向劝勉，不像在其它书信中那样作两分法，将前一部分用于教义，后一部分用于道德训诲的料理。在此他并非如此，而是将后者与主题通篇融合，使之易于接受。因此他说有两种致死、两种死亡，一种由基督在圣洗中完成，另一种则需我们此后藉勤勉努力来实现。因为我们的先前之罪被埋葬，是出于祂的恩赐。但圣洗后继续死于罪恶，则必须是我们自身勤勉的工夫，尽管我们也发现天主在此大大帮助我们。因为圣洗所能做到的不仅是抹去我们先前的过犯，它也能确保我们免于后续的过犯。正如在先前的过犯中，你的贡献是信德，使它们被抹去，同样，在此后的过犯中，你也要在目标上显出改变，以免再次玷污自己。因为他劝告你的正是这一点以及类似的事，当他说：\BibleQuote{因为我们若同栽于祂死亡的形像，也必同栽于祂复活的形像。}你是否注意到，他如何藉着立刻把听众领到他们的主那里，并竭力强调这相似之处来激发听众？这就是为什么他不说\Quote{在死亡中}，免得你反驳，而是说\Quote{在祂死亡的形像中}。因为我们的本质本身并未死亡，而是那罪恶之人，即邪恶。他也没有说\Quote{因为我们已有分于}祂死亡的\Quote{形像}，而是说\Quote{我们若同栽}，藉着''栽种''这个词，暗示由此为我们结出的果实。因为正如祂的身体被埋在地里，结出的果实是世界的救恩；同样，我们的身体被埋在圣洗中，结出的果实是正义、圣化、义子身份和无数祝福。并且它将来还会结出复活的恩赐。既然我们被埋在水中，祂被埋在土里，我们相对于罪恶，祂相对于身体，这就是为什么他没有说\Quote{我们同栽于祂的死亡}，而是\Quote{同栽于祂死亡的形像}。因为两者都是死亡，但主体不同。因此如果他说\Quote{我们同栽于祂的死亡，也必同栽于祂的复活}，这里指的是那将来的复活。因为当他先前谈到死亡时说：\BibleQuote{弟兄们，岂不知我们这受洗归于基督的人，是受洗归于祂的死亡吗？}他并没有清楚说明复活，只谈了圣洗后的生活方式，吩咐人在新生活中行走；因此他在这里重拾同一主题，并清楚预言了那复活。并且为叫你知道他不是指由圣洗产生的复活，而是指另一个，他说\BibleQuote{因为我们若同栽于祂死亡的形像}之后，并没有说我们必在祂复活的形像里，而是说我们必属于那复活。以免你问：如果我们不像祂那样死，怎能像祂那样复活？当他提到死亡时，他没有说\Quote{同栽于死亡}，而是\Quote{同栽于祂死亡的形像}。但当他提到复活时，他没有说\Quote{在复活的形像里}，而是说我们必属于复活本身。他也没有说我们已经成了，而是说我们必要是，藉这个词再次明确指那尚未发生、将来才有的复活。然后为使他的话可信，他指出了在此之前在此世发生的另一个复活，好叫你因这现世的复活也相信那将来的。因为他说\Quote{我们必同栽于复活}之后，又加上：\par

第6节。\BibleQuote{知道这一点：我们的旧人与祂同钉十字架，使罪恶之身被毁灭。}\BibleRef{罗 6:6}\par

这样他同时给出了将来复活的原因和证明。他没有说''被钉十字架''，而是说''与祂同钉十字架''，如此将圣洗与十字架拉近。为此他在上面也说过：\Quote{我们同栽于祂死亡的形像，使罪恶之身被毁灭}，他并不是指我们这身体，而是指一切邪恶。因为正如他把邪恶的总和称为旧人，同样他把由不同邪恶部分组成的邪恶称为那个人的身体。为证明我所说的不是猜测，请听保禄自己在下一节对这件事的解释。因为他说\BibleQuote{使罪恶之身被毁灭}之后，又加上\BibleQuote{今后我们不再作罪恶的奴仆}。随着他继续，他把这点说得更清楚。\par

第7节。他说：\BibleQuote{因为死了的人，已脱离（原文：成义）了罪。}\BibleRef{罗 6:7}\par

这是对每个人说的：正如死了的人因身体已死而不再犯罪，同样从圣洗中上来的人，既然在那里一次死了，就应永远对罪保持死亡。因此如果你在圣洗中死了，就保持死亡，因为任何死了的人不能再犯罪；但如果你犯罪，你就破坏了天主的恩赐。在要求我们如此程度的\OriginalTerm{英勇}{哲学}之后，他随即又带来了冠冕，说：\par

第8节。\BibleQuote{我们若与基督同死。}\BibleRef{罗 6:8}\par

诚然，甚至在加冕之前，这本身就是更大的冠冕——与我们的主有分。但祂说，我还要赐另一种赏报。是怎样的呢？就是永生。因为祂说：\Quote{我们相信，我们也要与祂一同生活。} 这如何显明呢？\par

第9节。\BibleQuote{基督既从死者中复活，就不再死。}\BibleRef{罗 6:9}\par

请再次留意他的无畏，以及他如何从相反的立场来证实这事。既然很可能有些人会对十字架和死亡感到困惑，他便指出，正是此事从此成为使人充满信心的理由。\par

他说：不要以为祂既然死过一次，便是可朽的；因为这正是祂不朽的原因。祂的死是死之死，正因祂死了，所以祂不再死。因为那死亡本身\par

第10节。\BibleQuote{祂向罪恶死了。}\BibleRef{罗 6:10}\par

\BibleQuote{向罪恶} 是什么意思？意思是祂连那一次也并非受其辖制，而是为了我们的罪，为要消灭罪，斩断其筋腱和一切能力，祂才死了。你看出他如何惊吓他们吗？因为如果祂不再死，就没有第二次\OriginalTerm{圣洗}{laver}；那么你当远离一切易犯罪的倾向。他说这一切，都是为了驳斥\Quote{我们作恶以成善，让我们留在罪中使恩宠增多} 的说法。因此，为了根除这种观念，他才写下这一切。但祂活着，\BibleQuote{祂活于天主}，他说——即是说，不再改变，以至于死亡再不能管辖祂。因为如果祂先前死，不是因为祂自身有任何必死的缘故，只是为了别人的罪，那么如今祂既已除灭了那罪，更不会再次死去。他在\BibleBook{希伯来}{书}中也这样说：\BibleQuote{但如今，在末世一次出现，以祂自己的祭献除掉了罪。就如规定人只死一次，以后便是审判；同样，基督也只一次奉献了自己，为承担大众的罪；对那些期待祂的人，祂要第二次显现，与罪无关，为施救恩。}\BibleRef{希 9:26} 他既指出合乎天主的生活的能力，也指出罪的力量。关于合乎天主的生活，他表明基督不再死；关于罪，他说明既然它甚至导致了无罪的那位的死亡，那么它怎能不使那些受它辖制的人灭亡呢？然后，既然他论及祂的生活，恐怕有人说：你所说的话与我们何干？他便加上，\par

第11节。\BibleQuote{你们也要照样算定自己向罪恶确实是死的，但在耶稣基督我们的主内，却活于天主。}\BibleRef{罗 6:11}\par

他很好地用了\Quote{算定}，因为他所谈论的目前尚未显于眼前。我们要算什么？有人会问。就是我们\BibleQuote{向罪恶确是死的，但在耶稣基督我们的主内，却活于天主}。因为如此生活的人必能抓住一切美德，有耶稣自己作他的盟友。这就是\BibleQuote{在基督内} 的意思；因为如果祂使死人复活，那么当人还活着时，祂更能保守他们如此。\par

第12节。\BibleQuote{所以，不要让罪恶在你们必死的身体上为王，以致顺从它的情欲。}\BibleRef{罗 6:12}\par

他并没有说：不要让肉身生活或行动，而是说：\BibleQuote{不要让罪恶在你们必死的身体上为王}；因为祂来不是要毁灭我们的本性，而是要纠正我们的自由选择。然后，为表明我们被不义所束缚并非出于武力或必然，而是出于自愿，他不用\Quote{不让它暴虐}（这个词会暗示必然），而用\Quote{不让它作王}。因为那些被引导进入天国的人竟让罪恶作他们的女王，那些蒙召与基督一同为王的人竟选择作罪恶的囚徒，这是荒谬的；如同有人将王冠从头上掷下，宁肯作一个前来乞讨、衣衫褴褛的疯妇的奴隶。接着，既然战胜罪恶是一项艰巨的任务，请看他如何表明它甚至很容易，以及他如何减轻劳苦，说\BibleQuote{在你们必死的身体上}。因为这表明斗争只是暂时的，很快就会结束。同时，他提醒我们过去的苦境和死亡的根源，因为正是从死亡开始，身体变成了可朽的。然而，一个必朽的身体仍可能不犯罪。你看到基督恩宠的丰盛吗？因为亚当，虽然那时还没有必朽的身体，却堕落了；而你这承受了甚至可朽身体的人，却能获得冠冕。那么，如何\Quote{罪恶作王}呢？他说：这不是出于它自身的能力，而是出于你的懈怠。因此，在说了\Quote{不让它作王}之后，他也指出了这种作王的方式，接着说道：\BibleQuote{以致顺从它的情欲}。因为凡事随它（即身体）所欲，不是荣耀，反而是极端的奴役和最大的耻辱；因为当它为所欲为时，它便丧失一切自由；但当它受约束时，它才最好地保持自己应有的地位。\par

第13节。\BibleQuote{也不要把你们的肢体献给罪，作不义的武器…却要献给天主，作正义的武器。}\BibleRef{罗 6:13}\par

身体对恶与德是无关的，正如武器（或兵器）也是一样。但无论是哪种效果，都来自使用它的人。好比一个士兵为自己的国家作战，而一个强盗武装起来对抗居民，却使用同样的武器进行防御。因为过错不在于盔甲的装备，而在于那些将它用于邪恶目的的人。关于肉身也可以这样说，它成为这样或那样，是由于心灵的决断，而非由于它自身的本性。因为如果它贪恋他人的美貌，眼睛就成为了不义的武器，并非出于眼睛自身的作为（因为眼睛的本分只是看，并非看错），而是由于支配它的思想的过错。但如果你约束它，它就成为义德的武器。舌头如此，手如此，所有其他肢体也是如此。他恰当地称罪为不义。因为人犯罪，要么对自己不义，要么对邻人不义，或者更准确地说，对自己比对邻人更加不义。那么，在带领我们远离了邪恶之后，他以这些话引导我们走向德行：\par

\Quote{但要将你们自己奉献给天主，如同从死者中复活的人。}\par

请看，他仅以简单的话语鼓励他们，在那边提到\Quote{罪}，在这边提到\Quote{天主}。因为他通过表明统治者之间有多么大的差异，使离开天主而渴望在罪的统治下服役的士兵没有任何借口。但他不仅借此方式，还通过后续的话确立了这一点，说道：\Quote{如同从死者中复活的人。}因为通过这些，他显示了前者的不幸和天主恩赐的伟大。因为，他说，请想想你们从前是什么，以及你们被造成了什么。那么从前你们是什么呢？是死的，被一种从任何方面都无法修复的毁灭所摧毁。因为没有人有能力帮助你们。而你们从这些死者中被造成什么呢？活着的，得享不朽的生命。凭着谁？凭着全能的天主。因此，你们应当以如同那些从死者中复活的人那样的欣然乐意，将自己列于祂的麾下。\par

\Quote{并将你们的肢体作为义德的武器。}\par

由此可知，身体并非邪恶，因为它可以成为义德的兵器。但通过称它为兵器，他清楚地表明我们面临着艰苦的战争。为此，我们需要坚固的盔甲，也需要高贵的灵魂，并且要通晓这战争的方式；最重要的是我们需要一位统帅。然而，统帅正站立一旁，随时准备帮助我们，且常胜不败，祂也赐予我们坚固的兵器。此外，我们需要有正确运用这些兵器的心志，以便我们既能服从统帅，又能为国家出战。那么，在给予我们如此有力的劝勉、并提醒我们兵器、战斗和战争之后，请看，他如何再次鼓励士兵，滋养他欣然的精神。\par

第14节. \BibleQuote{因为罪不再统治你们，因为你们不在法律之下，而在恩宠之下。}\BibleRef{罗 6:14}\par

既然如此，如果罪不再统治我们，为什么他在\Quote{不要让罪在你们必死的身体中为王}和\Quote{不要将你们的肢体献给罪作不义的武器}这些话中，对他们提出如此重大的要求呢？那么这里所说的话是什么意思呢？他是在这话中播种一种种子，以后要发展它，并用强有力的论证来培养它。那么这话是什么呢？就是说：在基督来临之前，我们的身体容易成为罪的攻击目标。因为在死亡之后，一大群情欲也进入了。为此，在德行竞赛中并不轻松。因为没有圣神前来帮助，也没有具有德能的圣洗来致死。\BibleRef{若 7:39} 但正如一匹马（柏拉图\WorkTitle{斐德罗篇}第74节）不听从缰绳，它确实奔跑，但经常失足，法律同时宣告什么是该做的、什么是不该做的，却没有给赛跑者带来任何超出言语劝诫之外的东西。但当基督来临后，努力变得更容易，因此为我们设立了\OriginalTerm{更远的目标}{\GreekText{μείζονα} \GreekText{τὰ} \GreekText{σκάμματα}}，因为我们得到的帮助更大。因此基督也说：\BibleQuote{除非你们的义德超过经师和法利塞人的义德，你们决不能进入天国。}\BibleRef{玛 5:20} 但他在后面更清楚地说出了这一点。然而现在，他在这里简要地提及，以表明除非我们非常屈就于罪，罪就不能胜过我们。因为不仅法律劝诫我们，而且恩宠也已赦免了我们从前的罪，并保护我们免受将来的罪。因为法律应许劳苦后的冠冕，而恩宠（即恩宠）先加冕了他们，然后带领他们进入竞赛。现在，在我看来，他在这里并非指信徒的整个生命，而是建立圣洗与法律之间的比较。他在另一处也这样说：\BibleQuote{文字带来死亡，但圣神带来生命。}\BibleRef{格后 3:6} 因为法律证明过犯，而恩宠消除过犯。因此，正如前者通过证明而坚固罪，后者通过宽恕使我们不受罪的辖制。这样，你们便在两方面摆脱了这轭：既不在法律之下，又享有恩宠。那么，在以此话让听者稍作喘息之后，他再次为他们提供保护，通过引入一个针对异议的劝勉，并如下所说。\par

第15节. \BibleQuote{那么，怎么样呢？我们要犯罪，因为我们不在法律之下，而在恩宠之下吗？绝对不可！}\BibleRef{罗 6:15}\par

所以他首先采用了一种誓愿的形式，因为他所提及的是一个荒谬之事。然后他使讲论转入劝勉，并以以下话语表明斗争的极大容易。\par

第16节. \BibleQuote{难道你们不知道：你们将自己献给谁作奴仆服从谁，就是谁的奴仆，或是犯罪的奴仆以至于死，或是顺服的奴仆以至于成义吗？}\BibleRef{罗 6:16}\par

他说：我还不提地狱，也非那极大的（抄本博德利长稿）惩罚，而是现世的羞耻，当你成为奴隶，而且是自愿的奴隶，罪恶的奴隶，并且报酬是第二次死亡。因为如果圣洗之前，它带来了肉体死亡，伤口需要如此大的照料，以至万有的主降来受死，从而遏止了邪恶；如果在如此大的恩赐、如此大的自由之后，它再次抓住你，而你甘心屈服其下，它还有什么做不出来？那么，不要跳入这样的坑，或甘心交出自己。因为在战争中，士兵常被迫交出。但在这种情况，除非你自行叛离，没有人能胜过你。于是，他试图以责任感羞愧他们，也以赏报惊吓他们，并将两者的报酬摆在他们面前：义德和死亡，且这死亡不像从前，而是更糟。因为如果基督不再死，谁能消除死亡？没有人！那么我们必须永远受罚，受报复。因为一种感官可觉察的死亡在此不会来临，如同前者那样，给身体安息并与灵魂分离。\BibleQuote{因为最后被毁灭的仇敌，便是死亡}\BibleRef{格前 15:26}，因此惩罚将是不死的。但是，对于服从的人，义德及其所产生的祝福将是他们的报酬。\par

第17节。\BibleQuote{但感谢天主，你们曾作罪恶的奴隶，却从心里听从了那传给你们的教导模式。}（直译：\InnerQuote{你们被交付了的}）\BibleRef{罗 6:17}\par

在借此奴隶制羞愧他们，借此赏报惊吓他们，并如此劝诫他们之后，他又通过回想恩惠来扶正他们。因为借此他显示，他们从极大的邪恶中被释放，且不是靠自己的任何辛劳，并且以后的事将更易管理。正如有人从暴君手中救出一名俘虏，并劝他不要逃回去，提醒他过去的痛苦奴役；同样，\Person{保禄}通过感谢天主，最有力地摆在我们面前已过去的邪恶。因为不是人的能力能使我们摆脱这一切邪恶，而是\Quote{感谢天主}，祂愿意且能够成就如此伟大的事。他又恰当地说：\Quote{你们从心里听从了}。你们既非被迫，也非受压，而是自愿以甘心之心过来的。这就像一个人既赞美又指责。因为既然自愿过来，且未经历任何必要，若你逃回从前的情形，你还能要求什么宽恕，或作什么借口？其次，为使你们知晓这并非仅出于你们自愿的气质，而全是天主的恩宠，在说了\Quote{你们从心里听从了}之后，他又加上\Quote{那传给你们的教导模式}。因为从心里的听从显示自由意志，而被交付则暗示天主的助佑。但教导模式是什么？是生活正直，与最佳行为相符。\par

第18节。\BibleQuote{如此，你们既然从罪恶中获得了自由，就成了义德的奴仆。}\BibleRef{罗 6:18}\par

他在这里指出天主的两种恩赐：\Quote{从罪恶中解放}，以及\Quote{使他们成为义德的奴仆}，后者胜过任何自由。因为天主所做的，如同有人夺取一个被野人掳到他们国家的孤儿，不仅使他脱离囚禁，而且为他安排一位慈父，并使他到达极大的尊荣。这同样发生在我们身上。因为祂不仅将我们从旧日的邪恶中解放，甚至引导我们到天使的生命，为我们铺平通向最佳行为的道路，将我们交给义德的安全守护，杀死我们旧日的邪恶，使旧人致死，并引导我们到达不朽的生命。\par

那么，我们就继续活在这生命中吧；因为许多看似呼吸行走的人，其处境比死者更为悲惨。因为死有多种：有一种是肉身的死，按此说亚巴郎曾死过，却仍未死。因为\BibleQuote{天主}，祂说，\BibleQuote{不是死人的天主，而是活人的天主。}\BibleRef{玛 22:32}另一种是灵魂的死，基督提及此，当祂说：\BibleQuote{让死人埋葬他们的死人。}\BibleRef{玛 8:22}另一种甚至是应受称赞的，是由\OriginalTerm{哲学}{\GreekText{φιλοσοφίας}}所带来的，关于此保禄说：\BibleQuote{你们要致死那属于地上的肢体。}\BibleRef{哥 3:5}另一种甚至是这（死）的原因，即发生在圣洗中的死。他说：\BibleQuote{因为我们的旧人已被钉在十字架上，}\BibleRef{罗 6:6}也就是说，已被处死。既然我们知道这些，就让我们逃避那种虽生犹死的死，不要害怕那带来普通死亡的死。至于另外两种，一种是有福的，由天主所赐；另一种是可称赞的\EditorialNote{参亚里士多德《尼各马可伦理学》I.12}，由我们与天主一同完成，让我们选择并效法它们吧。在这两种中，达味称一种有福，他说：\BibleQuote{罪恶蒙赦免的人是有福的。}\BibleRef{咏 31:1}；另一种，保禄赞赏，在致迦拉达人的书信中说：\BibleQuote{凡属于基督的人，已将肉身钉在十字架上了。}\BibleRef{迦 5:24}至于另一对，基督宣布其中一种容易轻蔑，当祂说：\BibleQuote{不要害怕那些杀害身体却不能杀害灵魂的；}另一种则是可怕的，因为祂说：\BibleQuote{要害怕那把身体和灵魂都能毁灭在地狱里的。}\BibleRef{玛 10:28}因此，让我们逃避这种死，选择那被称为有福和可赞叹的死；这样，我们就能逃避另一种死，不惧怕那一种：因为看见太阳、吃喝，对我们并非最小的益处，除非我们拥有善言的生活。请问，一个身着紫袍、拥有武器却没有一个臣民、任由一切有意攻击侮辱他的国王，有什么益处呢？同样，一个基督徒有了信德和圣洗的恩赐，却向一切情欲敞开，也毫无益处。这样，耻辱更大，羞愧更甚。因为正如这样一个戴着王冠、穿着紫袍的人，非但没有因此获得任何尊荣，反而因自己的羞耻玷污了这装束；同样，一个信徒若过着败坏的生活，非但不能因此受人尊敬，反而只会更加被人鄙视。经上说：\BibleQuote{因为凡没有律法而犯罪的，也必没有律法而灭亡；凡在律法以下犯罪的，必按律法受审判。}\BibleRef{罗 2:12}。在\BibleBook{希伯来}{书}中，他说：\BibleQuote{凡废弃梅瑟律法的，凭两三个见证人，尚且不得怜悯而被处死；那践踏了天主子的人，你们想，该受怎样更重的刑罚呢？}\BibleRef{希 10:28}这有道理。因为（祂可能会说）藉着圣洗，我已使所有情欲受制于你。那么，你何以玷污了如此伟大的恩赐，变得面目全非？我已像杀死并埋葬了蠕虫一样，除掉了你以往的过犯——你怎能又生出别的呢？因为罪比蠕虫更坏，蠕虫伤害身体，罪伤害灵魂，且发出更恶臭的气味。然而我们却察觉不到，因此不刻意清除它们。因此，酒鬼不知道陈酒的恶臭，而不醉酒的人却能清楚感知。同样，对于罪，度清醒生活的人彻底了解那污秽和斑点。但那纵情邪恶性的人，像醉酒昏睡的人，甚至不知道自己生了病。这就是恶习最可悲之处：它不让陷于其中的人看见自己祸患的严重性，反而当他们躺在泥淖中时，自以为在享受芬芳。因此他们甚至无力挣脱；反而当满身蠕虫时，像人以宝石为荣那样，以这些为乐。于是他们连杀死它们的意愿都没有，反而滋养它们，使它们在自己内繁衍，直到把它们送往来世的蛆虫那里。因为它们是那些（蛆虫）的供应者，不仅是供应者，甚至是不死之虫的父；正如经上说：\BibleQuote{他们的虫不死。}\BibleRef{谷 9:44}这些点燃了永不熄灭的地狱之火。为避免此发生，让我们消除这邪恶之源，熄灭这火炉，从根部拔出邪恶之根，因为若根留在下面，再从上面砍断树木也无济于事，它仍会长出新芽。那么，什么是万恶之根？从那位好农夫学习吧\BibleRef{格前 3:6}，他精确知道这些事，照料属灵的葡萄树，耕耘整个世界。他说什么是万恶之根？贪爱钱财。因为\BibleQuote{贪爱钱财是万恶之根。}\BibleRef{弟前 6:10}。由此产生争斗、敌意和战争；由此产生忌妒、诽谤、猜疑、侮辱；由此产生谋杀、偷盗和盗墓。\par

藉着这，不仅是城市和国家，而且道路、可居住与不可居住之地、山岭、树林、丘陵，总而言之，所有地方都充满了流血与凶杀。甚至海洋也没有免除这邪恶，因为海盗从四面袭击，发明了一种新的抢劫方式，同样以极大的狂暴肆虐其中。藉着这，自然法则被颠覆，亲属关系被置诸脑后，虔敬的法则本身也被破坏。因为金钱的奴役武装了这些人的右手，不仅对付活人，甚至也对付死人。即使在死亡之时，他们也不与之和解，而是掘开坟墓，将不敬的手伸向尸体，连那已放弃生命的人，他们也不让他脱离他们的阴谋。无论你在屋里、市场、法庭、元老院、王宫或任何其他地方能找到的一切邪恶，你会发现全都由此而生。因为这邪恶，正是这邪恶，使所有地方充满流血与凶杀，点燃地狱的火焰，使城市如同荒野一般悲惨，甚至更加悲惨。因为那些拦路抢劫的人，人们可以轻易提防，因为他们不是随时攻击。但在城市中仿效他们的人，比前者更坏，因为更难防范，并且公然敢做那些暗中做的事。那些为遏止他们不义而制定的法律，他们甚至拉来作为盟友，使城市充满这类凶杀和污秽。请问，把穷人交给饥荒，投入监牢，不仅使之遭受饥荒，还有酷刑和无数侮辱，这不是凶杀，且比凶杀更甚吗？即使你没有亲自对他做这些事，但你却是这些事的起因，你比执行这些事的差役做得更多。凶手一刀刺入人体，使他痛苦短暂时间，就不再延长折磨。但你，藉着你的诽谤、骚扰、阴谋，使他的光明变为黑暗，使他千万次希望死去，你想一想，你实行了多少次死亡，而不是仅仅一次？最糟的是，你掠夺和贪得无厌，并非贫困所迫，也没有饥饿的驱使，而是为了使你的马嚼子镀上足够的金，或是你家的天花板，或是你的柱头。哪里的地狱会不配这种行为？当他是你的弟兄，与你分享了不可言喻的恩赐，并为主所如此尊荣，而你为了装饰石头、地板和那些没有理性、没有感知这些装饰的动物的身体，却使他陷入无数灾祸中？你的狗也备受关照，而人，或者说基督，却为了猎犬和所有这些我提到的东西，因极端饥饿而受困。有什么比这混乱更糟？有什么比这无法无天更可悲？什么样的火流才足够炼净这样的灵魂？那按照天主的肖像所造的人，因你的无情而处境不堪；而你妻子所拉骡子的脸，以及构成那顶篷的皮和木头，却闪耀着大量的金。如果要制作一个座位或脚凳，它们全是金和银的。但基督的肢体，基督为此从天堂降临，倾流祂的宝血，却因你的贪婪连必需的食物也得不到。而床榻到处蒙着银，圣人们的肢体却连必要的衣服都没有。对你而言，基督不如任何东西珍贵，仆役、骡子、床榻、椅子、脚凳都不如；因为我略过比这些更低微的家具，留给你自己去想。如果你听到这些感到震惊，那就停止这样做，那么所说的话就不会伤害你。停止吧，不要再疯狂。因为这确实是疯狂，对这些东西如此热衷。因此，让我们放下这些，迟了也好，仰望天上，记念那将要来的日子，思想那可怕的审判台，精确的账目，和永不腐败的判决。让我们想想，看见这一切的天主，并没有从天上降下闪电；然而所做的事不仅该当被雷劈。但祂并没有这样做，也没有让海洋泛滥我们，也没有把地裂开，没有熄灭太阳，也没有把天和星辰砸在我们身上。祂不动摇任何东西，却让它们继续运行，让整个受造物为我们服务。深思这一切吧，让我们敬畏祂对人的爱是何等伟大，让我们回归那属于我们的高贵根源，因为目前我们确实不比无理性的受造物更好，甚至更糟。因为它们爱自己的同类，只需自然的共性就能对彼此产生感情。但你，除了自然之外，还有无数因素将你拉近并附着于你自己的肢体：被圣言尊荣，参与同一宗教，共享无数恩赐，却变得比它们更野蛮，因为你如此关心无益之物，而让天主的殿宇在饥饿和赤身中消亡，常常还使它们遭受一千种邪恶。如果你做这些是出于爱光荣，那么你更应该关注你的弟兄，而不是你的马。因为接受善意的受造物越好，为这种关心编织的花冠就越亮。既然你现在陷入这一切的反面，你招来无数控告者，却还察觉不到。因为谁不会说你坏话？谁不会控告你犯有最可恶的暴行和厌世？当他看到你忽视人类，将无知的受造物置于人之上，此外还有你屋里的家具？你没有听到宗徒们说，那些首先接受圣言的人，变卖了\BibleQuote{房屋和田地}\BibleRef{宗 4:34}，以供养弟兄们？而你却掠夺房屋和田地，为了装饰马、木器、皮、墙壁或地板。更糟的是，不仅男人，女人也受这疯狂之苦，迫使她们的丈夫从事这种虚空的劳苦，强迫他们把金钱花费在任何东西上，而不是必需的事物。如果有人为此责备她们，她们便习于一种本身充满诸多可指责之处的辩护。因为有人说这两件事是同时做的。\par

你们说什么？难道你们不怕说出这样的事，并且把马、骡、床榻、脚凳看得与饥饿的基督同等吗？或者还不如说，根本没有把祂与这些相比，而是把更大的份给了这些东西，对祂却艰难地施以微薄的一份？你们不知道一切都属于祂，包括你们和你们所有的吗？你们不知道祂塑造了你们的身体，同样赐给你们灵魂，并将整个世界分配给你们吗？但你们却不愿给祂一点回报。如果你们出租一间小屋，你们会极其严格地要求租金，然而你们享受着祂全部的受造物，居住在如此广阔的世界中，却不敢支付一点租金，反而把你们自己和你们所有的一切都交给了虚荣。因为这一切都由此而来。马匹并不因为有了这装饰而超越其本性的优越，骑在它上面的人也是如此，有时反倒因此更不受尊重；因为许多人忽视骑手，却把目光转向马的装饰，以及前后的随从和持扇者。但他们却憎恨被这些簇拥的人，并转过头去，如同对待公敌。但当你装饰你的灵魂时，情况就不同了；那时人、天使，以及天使的主，都会为你编织冠冕。所以，如果你热爱荣耀，就远离你现在所做的这些事，把你的品味展示在你的灵魂中，而不是你的房屋中，那样你就会变得光彩夺目。因为现在没有比你们更廉价的了，你们的灵魂毫无装饰，仅以房屋的华美作为遮掩。但如果你们不耐烦听我这样说，那就听听那些外邦人中的一个做了什么，至少因他们的哲学而感到羞耻。据说他们中有一人进入一座宫殿，那里金光灿烂，大理石和柱子的华美闪闪发光，当他看到地面四处铺着地毯时，便吐唾沫在屋主的脸上；当有人为此责备他时，他说：\Quote{因为这屋里没有其他地方可以这样做，他只好当着面这样做。} 看，一个在外表展示品味的人是多么可笑，在所有明白人眼中是多么微不足道。这很有道理。因为如果有人让你们的妻子衣衫褴褛、无人照顾，却让你们的婢女穿着华丽的衣服，你们不会温和地忍受，反而会愤怒，说这是极大的侮辱。那么，关于你们的灵魂，也要这样想。当你们在墙壁、地板、家具以及其他类似的东西上展示品味，却不肯慷慨施舍，或实践修行生活的其他部分（\OriginalTerm{修行生活}{\GreekText{φιλοσοφίαν}}）时，你们所做的与此无异，甚至更糟。因为婢女与主母的差别不算什么，但灵魂与肉体之间却有巨大的差异。如果连肉体都是如此，那么房屋、床榻、脚凳就更不用说了。那么，你们值得什么原谅呢？你们将银子放在所有这些上，却毫不关心灵魂，尽管它覆盖着污秽的破布，肮脏、饥饿、满身伤痕，被无数恶狗撕咬（\BibleRef{路 16:20}）；而在这之后，竟幻想通过展示外在的装饰来获得荣耀？这正是疯狂之极，一边被人嘲笑、责备、羞辱、贬损，陷入最严厉的惩罚，仍然对这些事自夸！因此，我恳求你们，将这一切放在心上，让我们尽管为时已晚，仍要清醒过来，成为自己的主人，将这装饰从外在事物转移到我们的灵魂上。因为这样它才能免于毁坏，使我们与天使相等，并以永不改变的美善款待我们；愿我们所有人都能藉着祂的恩宠和爱人之爱获得这美善，等等。\par

\SourceRecord{Translated by J. Walker, J. Sheppard and H. Browne, and revised by George B. Stevens.；Nicene and Post-Nicene Fathers, First Series；Vol. 11.；Edited by Philip Schaff.；Buffalo, NY: Christian Literature Publishing Co.,；1889.；<http://www.newadvent.org/fathers/210211.htm>.}

% structure: p0001 source=h1 target=section reason=page-title

\section{罗马书第十二篇讲道}

\BibleRef{罗 6:19}\par

\BibleQuote{我因你们肉体的软弱，就照人的常话对你们说：你们从前怎样将肢体（所以四份抄本 Sav. 作\Quote{你们肉体的肢体}）献给不洁和不法作奴仆，以至于不法；现今也要照样将肢体献给义作奴仆，以至于成圣。}\par

既然他要求严格的生活，吩咐我们向世界而死，向邪恶而死，不与罪恶的作为有任何念头，似乎说出了一些伟大、沉重、超乎人性的话；为了表明他没有提出过分的要求，甚至没有像享受如此大恩的人所应期望的那样，而是非常适中且轻松的，他从反面证明，并说：\Quote{我照人的常话对你们说}，意思是：按照人的推理；按照通常遇到的方式。因为他要么指这个，要么指其适度，用\Quote{照人的常话}这个说法。因为他在其他地方也用了同样的词。\BibleQuote{你们所遇见的试探，无非是人所能受的} \BibleRef{格前 10:13}，即适度且小的。\Quote{你们从前怎样将肢体献给不洁和不法作奴仆，以至于不法；现今也要照样将肢体献给义作奴仆，以至于成圣。} 的确，主人非常不同，但我要求的奴役程度是相等的。因为人应当提供更大的，而且越大越好，因为这是一个更大更好的主人。然而我没有提出更大的要求，\Quote{因为肉体的软弱}，而且他不是说你们的自由意志或精神准备，而是\Quote{肉体的}，从而使他的话不那么严厉。然而一边是不洁，另一边是成圣；一边是不法，另一边是义。谁如此可怜、如此窘迫，以至于不为基督的服事付出像为罪恶和魔鬼的服事那样的热心呢？那么，听下面的话，你就会清楚地看到我们甚至没有付出这一点。因为当这样赤裸裸地陈述时，似乎不可信、不易接受，没有人愿意听到他服事基督不如他服事魔鬼，他通过下面的话证明这一点，并通过呈现那种奴役状态使之可信，并说他们如何服事罪。\par

第20节。\BibleQuote{因为你们作罪之奴仆的时候，就脱离了义。}\BibleRef{罗 6:20}\par

现在他所说的大意是这样：当你们生活在邪恶、不敬和最坏的事情中时，你们所处于的顺从状态是你们完全不行任何善事。因为这就是\Quote{脱离了义}。即你们不受它支配，完全远离了它。因为你们甚至没有在义和罪之间划分服事的方式，而是完全将自己交给了邪恶。因此，既然你们已转向义，现在就完全把自己交给美德，不行任何恶事，以便你们付出的量至少相等。然而，不仅是主人如此不同，而且服事本身也有巨大差异。这一点他也非常清晰地展开，并说明了他们当时服事的条件，以及现在的条件。而且目前他还没有提到事情带来的害处，只谈到了羞耻。\par

第21节。\BibleQuote{那么，你们现今所羞耻的事，当时有什么果子呢？}\BibleRef{罗 6:21}\par

奴役如此之大，以至于现在仅回忆就使你们羞耻；但如果回忆使人羞耻，现实就更甚。所以你们现在在两件事上获益：从羞耻中得释放，并且认识到你们所处的状态；正如当时你们在两方面受损：做了可耻的事，并且甚至不知道要羞耻。而后者比前者更糟。然而你们仍然处于奴役状态。他既然从羞耻方面极其充分地证明了当时所发生的事的危害，便谈到事情本身。那是什么？\BibleQuote{因为那些事的结局就是死亡。} 既然羞耻似乎不是什么严重的恶，他就谈到非常可怕的事，即死亡；尽管实际上他以前提到的已经足够了。试想那祸害何等巨大，以至于即使从应受的惩罚中得释放，仍不能摆脱羞耻。那么，当你们从回忆中，而且是在从惩罚中得释放之后，还掩面脸红，何况在如此恩宠之下时，你们从现实中会有什么报酬呢？但天主的一面远非如此。\par

第22节。\BibleQuote{现今你们既从罪中得了释放，作了天主的奴仆，就有成圣的果子，那结局就是永生。}\BibleRef{罗 6:22}\par

从前的果子是羞耻，即使在得释放之后。这些的果子是成圣，有成圣的地方就有全然的自信。但那些事的结局是死亡，这些的结局是永生。你看到他是如何指出有些东西已经赐下，有些存在于希望中，并从已赐下的东西——即生活的成圣——中证明其余的东西。为了阻止你（作为反对）说一切都存在于希望中，他指出你已收获了果子：第一，从邪恶和那些仅仅回忆就使人羞耻的坏事中得释放；第二，成为义的奴仆；第三，享有成圣；第四，获得生命，而且不是暂时的生命，而是永生的。然而，尽管如此，他说，你们只要像服事它那样服事就行了。因为尽管主人更优越，服事也有许多优点，还有你们服事的报酬，我仍不作进一步要求。然后，既然他提到了武器和国王，他继续用比喻说：\par

第23节。\BibleQuote{因为罪的工价乃是死；但天主的恩赐，在我们的主基督耶稣内，乃是永生。}\BibleRef{罗 6:23}\par

在谈论了罪恶的代价之后，对于祝福，他没有保持相同的秩序（\OriginalTerm{秩序}{\GreekText{τάξιν}}）：因为他没有说，善行的代价，\Quote{但是天主的恩赐}；以表明，他们获得自由并非出于自己，也不是应得的回报，也不是劳苦的报酬，而是靠恩宠这一切才发生。因此，也由于这个原因，存在着一种优越性，在于他不仅解放了他们，或者改变了他们的状况，使其变得更好，而且他这样做没有他们方面的任何劳动或麻烦：他不仅解放了他们，还给了他们比以前多得多的事物，并且是通过他的圣子。所有这一切他插入其中，因为已经讨论了恩宠的主题，并且接下来将要推翻法律。那么，这些事可能不会使他们更加懈怠，他插入了关于严格生活的一部分，利用每一个机会激励听众实践德行。因为当他称死亡为罪恶的代价时，他再次警告他们，保护他们免受未来的危险。因为他用来提醒他们以前状况的话语，他也使用以便使他们感恩，并更加安全地抵御任何诱惑的侵袭。然后，他在这里停止劝勉部分，并再次继续教义，这样说。\newline{}\par

\BibleQuote{你们不知道吗，弟兄们——我是向懂法律的人说话——}（\BibleRef{罗 7:1}）\newline{}\par

既然他曾说过，我们\Quote{死于罪恶}，他在这里表明，不仅罪恶，而且法律也不再统治他们。但如果法律没有统治权，那么罪恶更没有；为了使他的话语易于接受，他用一个人间的例子来说明。他似乎在陈述一点，但同时为他的论点提出了两个论证。一个是，当丈夫死了，妻子不再受丈夫管辖，没有什么能阻止她成为另一个男人的妻子；另一个是，在这种情况下，不仅丈夫死了，而且妻子也死了。所以一个人可以通过两种方式享受自由。如果丈夫死了，她就摆脱了他的权力，那么当妻子也被证明死了时，她就更加自由。因为如果一件事让她摆脱他的权力，那么两者同时发生更让她自由。当他准备证明这些要点时，他以赞扬听众开始，说道：\Quote{你们不知道吗，弟兄们——我是向懂法律的人说话——}，也就是说，我所说的是完全公认的、清楚的，而且是对那些准确知道这一切的人说的，\newline{}\par

\BibleQuote{法律怎样统治人，只要他活着？}\newline{}\par

他没有说丈夫或妻子，而是说\Quote{人}，这个名字对两种受造物都是共同的；\BibleQuote{因为死了的人从罪恶中得释放（\OriginalTerm{成义}{\GreekText{δικαιόω}}）}（\BibleRef{罗 6:7}）。那么法律是为活着的人赐下的，但对死人，它不再发出命令。你看到他如何提出双重自由吗？接下来，在开头暗示了这一点后，他通过女人这个案例继续提出他的证明，如下所述。\newline{}\par

\BibleQuote{因为已婚的妇女，在丈夫活着的时候，被法律约束；但如果丈夫死了，她就从丈夫的法律中解脱。所以，如果丈夫活着，她嫁给另一个男人，就被称为淫妇；但如果丈夫死了，她就从那个法律中解脱，所以如果她嫁给另一个男人，也不是淫妇。}（\BibleRef{罗 7:2-3}）\newline{}\par

他不断地坚持这一点，而且非常精确，因为他对此证据深感确定；他把法律放在丈夫的位置，而把所有的信徒放在妻子的位置。然后他以下列方式添加结论，使结论与前提不完全吻合；因为上下文要求的是：\Quote{所以，我的弟兄们，法律不再统治你们，因为它死了}。但他没有这样说，只是在前提出了暗示，在推论中，后来，为了避免他的话令人反感，他通过说：\newline{}\par

\BibleQuote{所以，我的弟兄们，你们也成了对法律死了的人}（\BibleRef{罗 7:4}）\newline{}\par

既然这两种情况都带来同样的自由，那么他对法律显示恩宠而丝毫无损于（福音）事工，又有什么妨碍呢？\Quote{因为有丈夫的女人，当丈夫活着的时候，被法律束缚。}那些诋毁法律的人如今成了什么？（三个抄本作\InnerQuote{那时}）让他们听听，即使被迫如此，他也没有剥夺法律的尊严，反而大谈其能力；如果当法律活着时，犹太人被束缚，而违反它、在它活着时离开它的人被称为奸淫者。但如果法律死了之后他们才放弃它，这并不奇怪。因为在人间事务中，没有人因这样做而受指责：\Quote{但如果丈夫死了，她就从丈夫的法律中得了解脱。}你看在这个例子中，他如何指出法律是死的，但在推论中却没有这样做。所以，如果丈夫活着，女人被称为奸妇。看他如何细数那些在法律还活着时违反它的人的罪状。但既然他已经终结了法律，随后他就让法律完全安稳，于此丝毫不伤害信仰。\Quote{因为如果丈夫活着，她嫁给了另一个男人，她就被称为奸妇。}这样，接下来自然要说：弟兄们，现在法律既然死了，你们若嫁给另一个丈夫，就不会被判定犯奸了。然而他没有用这些话，而是用了什么？\Quote{你们对法律已经死了；}如果你们已经死了，你们就不再受法律的约束了。因为如果丈夫死了，女人不再受约束，那么当她自己也死了时，她就更从前者中得了释放。你注意到保禄的智慧吗？他如何指出法律自身也愿意我们与它离婚，而嫁给另一个？因为他的意思是，既然前夫死了，你们与另一个丈夫生活没有什么不对；因为怎么会有呢？既然当丈夫活着时，法律允许有离婚书的女人这样做？但他没有写下这点，因为这是对女人的一种指控；尽管这一点被允许，仍然不免有责备。\BibleRef{玛 19:7}因为在那些他已经通过必要而公认的证据获胜的案例中，他不会深入到不相干的问题中去；他并不吹毛求疵。那么令人惊奇的是，正是法律本身宣告我们这些与它离婚的人无罪，而它的本意是要我们属于基督。因为法律本身死了，我们也死了；它在我们身上的权力以双重方式被移除。但他不满足于此，还增加了理由。因为他不是毫无目的地提及死亡，而是再次引入造成这些事情的十字架，并以此方式使我们负有义务。因为他意思是，你们不单得了释放，而且是借着主的死亡。因为他说：\par

\Quote{你们借着基督的身体，对法律已经死了。}\par

他不仅以此为基础进行劝勉，还因为第二个丈夫的优越性。于是他接着说：\Quote{使你们归于另一位，就是那位从死者中复活的。}\par

为了防止他们说：\Quote{如果我们不愿与另一个丈夫生活，那又怎样？因为法律确实不使再嫁的寡妇成为淫妇，但尽管如此，它并不强迫她过这种生活。}为使不这样说，他指出，从已经领受的恩惠来看，我们必须选择它；他在其他段落表述得更清楚，他说：\Quote{你们不是自己的人；}和\Quote{你们是重价买来的；}以及\Quote{不要作人的奴仆}\BibleRef{格前 6:19}；又说：\Quote{一人替众人死了，是叫那些活着的人不再为自己活，乃为替他们死而复活的主活。}\BibleRef{格后 5:15}这就是他在这里用\Quote{借着身体}这句话暗示的事。接着他鼓励更好的希望，说：\Quote{为叫我们结出果实归于天主。}因为那时，他的意思是，你们结出果实归于死亡，但现在却归于天主。\par

第5节。\Quote{因为我们还在肉体中的时候，那些借着法律而生的罪恶情欲，在我们的肢体中运行，为叫我们结出死亡的果实。}\BibleRef{罗 7:5}\par

你看，从前的丈夫带来的利益！他没有说\Quote{当我们还在法律中时}，他在每一处都避免给异端者留下把柄；而是说\Quote{当我们还在肉体中时}，即处于恶行中、属肉体的生活中。所以他说的不是他们以前在肉体中，而如今四处游荡没有身体；而是通过这样说，他既不说法律是罪的原因，也不免除法律的可憎。因为法律起了严厉控告者的作用，揭露了他们的罪过：因为那向根本不愿服从的人提出更多要求，使过犯更大。这就是为什么他没有说\Quote{借着法律所生的情欲}，而是说\Quote{借着法律}\BibleRef{罗 2:27}，没有加上\Quote{产生}，只是\Quote{借着法律}，即借着法律被显明、被知道。接着，为了他不控告肉体，他没有说\Quote{肢体做成的}，而是说\Quote{在我们的肢体中运行（或被运行）}，以表明祸害的根源在别处，来自在我们内运行的思想，而不是来自让它们运行的肢体。因为灵魂像是演奏者，肉体的结构像是竖琴，随着演奏者的意愿发声。所以不和谐的音调不应归于后者，而应归于前者，而非后者。\par

第6节。他说：\Quote{但现在，我们从法律中被释放了。}（\OriginalTerm{失效}{\GreekText{κατηργήθημεν}}，\Quote{归于无效}）\BibleRef{罗 7:6}\par

请看，他在这里再次为肉身和法律开脱。因为他并没有说法律无效，或肉身无效，而是说我们被无效化（即被释放）。我们是如何被释放的呢？是因着那被罪压制的老旧之人死去并被埋葬。这就是他在\Quote{我们因那曾束缚我们的事而死}这句话中所阐述的。仿佛他说：捆绑我们的锁链被削弱并打破了，以至于那捆绑我们的罪，不再捆绑了。但是，不要退却或懈怠。因为你被释放是为了再次成为仆人，尽管方式不同，而是\Quote{以心神的新生，而不以文字的老旧}。现在他在这里是什么意思？因为有必要在此揭示，以免我们遇到这段经文时感到困惑。当亚当犯罪（他是说）时，他的身体变得易于死亡和受苦，也承受了许多身体的损伤，这匹马变得不那么活跃和不那么顺从。但基督来了，通过圣洗圣事使它为我们更加敏捷，用圣神的翅膀唤醒它。因此，我们所要跑的赛程标记与古人不同。既然当时的赛程不像现在这么容易，所以他要求他们不仅通过不杀人来保持纯洁——正如他昔时要求他们的那样——而且连愤怒也要避免；不仅命令他们远离奸淫，甚至连不洁的眼神也要避免；不仅免除虚假的起誓，甚至连真实的起誓也要免除。\BibleRef{玛 5:21} 他对他们的朋友也命令他们爱他们的仇敌。在所有其他责任上，他给我们更长的赛程，如果我们只是服从，他就用地狱来威胁我们，以此表明这些事并非战斗者自愿献出的（如独身和贫穷），而是我们必须绝对履行的。因为它们属于必要且紧急的要求，不做的人将受到极重的惩罚。这就是为什么他说：\Quote{除非你们的义德超过经师和法利塞人的义德，你们决进不了天国。}\BibleRef{玛 5:20} 但看不见天国的，必然落入地狱。为此，保禄也说：\Quote{罪不再统治你们，因为你们不在法律之下，而在恩宠之下。}这里又说：\Quote{你们应以心神的新生，而不以文字的老旧来服事。}因为不是文字（即旧的法律）定罪，而是圣神帮助。因此，在古人中，若有人实行守贞，那十分令人惊讶。但现在，这事普及世界各地。死亡在过去只有少数人艰难地蔑视它，但现在在村庄和城市中，有无数的殉道者，不仅有男人，还有女人。接下来，在说完这些后，他再次面对一个浮现的异议，并在面对它时为自己的论点提供确认。他并不把解答作为主要论点，而是通过反驳这一点来引出；这样，通过应对的紧迫性，他获得一个机会说出他想说的话，使他的责备不那么令人不悦。于是，他说了\Quote{以心神的新生，而不以文字的老旧}之后，他继续道。\par

第7节。\Quote{那么，法律是罪吗？绝对不是。}\BibleRef{罗 7:7}\par

在此之前，他一直在说：\Quote{罪的冲动，藉着法律，在我们的肢体中活动}\BibleRef{罗 7:5}；以及\Quote{罪不再统治你们，因为你们不在法律之下}\BibleRef{罗 6:14}；以及\Quote{哪里没有法律，哪里就没有过犯}\BibleRef{罗 4:15}；以及\Quote{法律进来，为叫过犯增多}\BibleRef{罗 5:20}；以及\Quote{法律激起忿怒}\BibleRef{罗 4:15}。既然所有这些话似乎都使法律声名狼藉，为了纠正由此产生的怀疑，他也设想了一个异议，并说：\Quote{那么，法律是罪吗？绝对不是。}在证明之前，他用这个誓言来安抚听众，并安抚那些因此感到困扰的人。这样，当他听到这些，并确信说话者的立场时，他就会与他一起探究表面的困惑，而不会对他产生怀疑。因此，他提出了异议，将对方与自己联合起来。所以他不是说我该说什么？而是\Quote{那么我们该说什么？}好像有一个商议和判断在他们面前，一个全体会议被召集，异议不是来自他自己，而是来自讨论的过程和实际的情况。因为\Quote{文字杀死}（他是说），没有人会否认，或者\Quote{圣神赐予生命}\BibleRef{格后 3:6}，这也是明显的，没有人会争辩。如果这些是公认的真理，那么我们关于法律该说什么？说\Quote{它是罪？绝对不是。}那么解释这个难题。你看见他如何假设对方在场，并承担教师的尊严，开始解释它。那么这是什么？他说，法律不是罪。\Quote{然而，除非藉着法律，我不知道什么是罪。}注意他智慧的深度！法律不是什么，他已经以异议的形式放下，以便通过移除它，从而取悦犹太人，说服他接受较小的选择。这是什么？就是\Quote{除非藉着法律，我不知道什么是罪。因为我不会知道贪欲，除非法律说：\InnerQuote{不可贪}}。\par

你注意到，他如何逐渐表明法律不仅是罪的控告者，在某种程度上还是它的生产者？然而，他证明这发生不是由于法律本身的过错，而是由于倔强的犹太人的过错。因为他很谨慎地堵住指责法律的摩尼教人的口；所以在说\Quote{然而，除非藉着法律，我不知道什么是罪；}和\Quote{我不会知道贪欲，除非法律说：\InnerQuote{不可贪}}之后，他加上：\par

第8节。\Quote{但罪，藉着诫命，机会，在我内产生各种贪婪。}\BibleRef{罗 7:8}\par

你可看到他如何为法律洗脱一切责难？因为他（保禄）说：\Quote{罪}，藉着诫命\Quote{找到机会}，增加了私欲偏情，并产生了与法律意图相反的结果。这来自软弱，并非来自任何恶意。因为当我们渴望某物，且又受阻时，欲望的火焰反而增加。现在这并非来自法律；因为法律（三份手稿作\Quote{尽力}）阻止我们，意图使我们远离它；但罪，即你自己的疏忽和恶劣的秉性，却将善事用于相反的目的。但这并非医生的过错，而是病人误用药物的过错。因为颁布法律的目的，并非激起私欲偏情，而是熄灭它，虽然结果恰恰相反。然而责备不在于它，而在于我们。假如有人发烧，想在对他不利时喝冷饮，而有人不让他尽情饮用，因而增加了他对这毁灭性快感的贪恋，这个人并不该受责备。因为医生的事仅仅是禁止，而自我克制则是病人的事。即使罪从中找到了机会，那又怎样？确实有许多恶人因善规而在自己的邪恶中增长。这正是魔鬼毁灭犹达斯的方式：使他陷入贪婪，偷窃属于穷人的东西。然而，托付钱囊并非导致此事的原因，而是他自身精神的邪恶。厄娃通过让亚当吃树上的果子，把他逐出乐园。但那个例子中，树也不是原因，即使机会是通过它发生的。但如果他对关于法律的讨论有些激烈，不要惊讶。因为保禄是在对抗当前的急务，不允许他的言词给那些怀疑的人留下把柄，而是竭力使当前的陈述正确。那么，不要孤立地审视他接下来所说的（四份手稿作\Quote{此处所说}），而要结合引导他说这些话的目的，并考虑到犹太人的疯狂和他们强烈的争辩精神；他渴望消除这种精神，因而似乎对法律\OriginalTerm{猛烈抨击}{\GreekText{πολὺς} \GreekText{πνεἵν}}，并非要指责它，而是要削弱他们的锐气。因为如果说法律的一个耻辱是罪从中找到机会，那么在\WorkTitle{新约}中也会发现同样的情况。因为在\WorkTitle{新约}中有成千上万的诫命，涉及更多（Field本作\Quote{远为更多}）重要的事。人们可以看到同样的事情也在那里发生，不仅关于贪财\BibleRef{罗 7:7}，而且普遍关于一切邪恶。因为主说：\BibleQuote{我若没有来向他们讲话，他们就没有罪。}\BibleRef{若 15:22}这里罪因此找到了立足点，并导致更大的惩罚。再者，当保禄谈到恩宠时，他说：\BibleQuote{你们想想，那践踏了天主子的人，该当受怎样更严厉的惩罚呢？}\BibleRef{希 10:29}那么，更重的惩罚难道不正是来源于更大的恩惠吗？他说希腊人无可推诿的原因，是因为他们蒙受了理性的恩赐，认识到了受造界的美好，并有一条公平的道路可以被引向造物主，却没有按他们的本分善用天主的智慧。你看到吗？对于恶人来说，在所有情况下，更大的惩罚机会都来自善事。但我们不会因此指责天主的恩惠，反而因此更加钦佩它们；我们要将责备归咎于那些将祝福滥用于相反目的之人的精神。那么，这也应成为我们对法律的态度。但前一点是容易且可行的——困难的是另一点。他为何说：\Quote{除非法律说\InnerQuote{不可贪恋}，我就不晓得贪恋？}如果人在接受法律之前不晓得贪恋，那洪水或索多玛的焚烧又是为何呢？那他是什么意思？他指的是强烈的贪恋：这就是为何他没有说\Quote{贪恋}，而说\Quote{各样私欲偏情}，暗示其强烈程度。有人会说，如果法律反而增加了混乱，它有何好处？毫无好处，甚至带来许多祸害。但责难不在法律，而在接受法律之人的疏忽。因为罪藉着法律运行了这事，尽管这并非法律的目的，恰恰相反。他说，罪因此变得更加强烈和凶猛。但这又是对法律的责难吗？不，是对他们顽梗的责难。\Quote{因为没有法律，罪是死的。}也就是说，不是那么显然。因为甚至在法律之前的人也知道自己犯了罪，但在颁布法律之后，他们对罪有了更确切的认识。因此他们该受更大的控告：因为仅凭本性的控告是一回事，而除了本性之外，还有清晰指出每一项罪过的法律又是另一回事。\par

第9节。\BibleQuote{因为没有法律，我从前是活着的。}\BibleRef{罗 7:9}\par

请问那是何时？在梅瑟之前。你看他如何致力表明：法律无论通过它所做的事还是它没做的事，都压制了人性。因为当他说\Quote{没有法律，我从前是活着的}时，他的意思是，我没有受到那么大的定罪。\par

\BibleQuote{但诫命来到，罪复活了，我反而死了。}\par

这确实像是控告法律。但若有人仔细观察，就会发现这甚至是赞美法律。因为它并非使先前不存在的罪产生，而只是指出那被忽视的罪。这至少是法律的称赞，因为在此之前，他们犯罪而不自知。但当法律来到时，即使他们除此之外一无所获，至少他们清楚地认识到了自己犯罪的事实。对于摆脱邪恶而言，这一点并非小事。如果他们并未摆脱，这与法律无关；法律为达到此目的而制定了一切，但指控完全针对他们那超乎想像的乖僻心灵。所发生的事并非自然——他们竟被有益的事物所伤害。这就是为什么他说：\Quote{而那本是引向生命的诫命，我发现它却引致死亡。} 他并不是说\Quote{它成了}或\Quote{它产生了}死亡，而是说\Quote{被发现}，以此解释这种新奇而异常的差异，并让一切责任落在他们自己头上。因为，他说，如果你要知道它的目的，它引向生命，并且是为这目的而赐下的。但如果死亡是结果，那是领受诫命者的过错，而不是这引向生命的诫命之过。这一点他在下文进一步阐明了。\par

\BibleRef{罗 7:11} \BibleQuote{因为罪趁着机会，借着诫命欺骗了我，并且借着它杀了我。}\par

你注意到他如何处处归咎于罪，完全为法律开脱罪责。于是他接着说道。\par

\BibleRef{罗 7:12} \BibleQuote{所以，法律是圣的，诫命也是圣的、公义的、良善的。}\par

但若你愿意，我们将把那些曲解这些声明之人的话摆在你面前。因为这会使我们自己的陈述更清楚。因为有些人说，他在这里所说的并非关于梅瑟法律，而是有些认为是指自然律，有些认为是指乐园中的诫命。然而，保禄的目的处处都是要废除这法律，但他与那些并无冲突。这很有道理；因为正是对这法律的恐惧和战栗，使犹太人顽固地反对恩宠。但他似乎从未称乐园中的诫命为\Quote{法律}，其他作者也从未如此。为了从他实际所说的使这一点更清楚，让我们稍加回溯，顺着他的话来探讨。在谈论了生活的严格之后，他接着说：\Quote{弟兄们，难道你们不知道，法律只在人活着的时候管束人吗？因此你们已死于法律。} 因此，如果这些话是关于自然律的，那么我们就被发现是没有自然律的。如果这是真的，我们比无理性的受造物更愚钝。然而，当然并非如此。至于乐园中的法律，没有必要争辩，以免我们陷入多余的麻烦，与人们已下定决心的事情作对。那么，他说\Quote{若不是借着法律，我就不知罪}，这是什么意思？他不是指完全的无知，而是指更准确的知识。因为如果这是针对自然律说的，后面的话如何对应？他说：\Quote{我从前是活着的}，\Quote{没有法律的时候。} 然而，无论是亚当还是其他任何人，都不能证明曾脱离自然律而生活。因为天主一形成他，就把那自然律放在他里面，使它作为整个人类的保障常驻其中（希腊文：自然，见第365页）。此外，他从未在任何地方称自然律为\Quote{诫命}。但这里的却称为诫命，并且是\Quote{公义而圣善}的，是\Quote{属神的法律}。但自然律并非由圣神赐给我们。因为野蛮人、希腊人以及其他人都拥有这法律。因此显然，他前后所说的都是梅瑟法律。为此他也称它为圣的，说：\Quote{所以，法律是圣的，诫命也是圣的、公义的、良善的。} 因为即使犹太人自从法律以来变得不洁、不义、贪婪，这并不能摧毁法律的美德，正如他们的不信并不能使天主的信实失效。因此，从这一切可以清楚地看出，他在这里说的是梅瑟法律。\par

\BibleRef{罗 7:13} \BibleQuote{那么，那善的竟成了我的死亡吗？绝对不可！乃是罪，为要显出是罪。}（四个抄本省略\OriginalTerm{}{\GreekText{ἡ}}。）\par

这就是说，要显出罪是何等大的恶，即：懈怠的意志、倾向于更坏一方的倾向、实际的行为（\EditorialNote{三份手抄本省略此句}），以及扭曲的判断。因为这是一切恶的原因；但\Person{保禄}藉着指出\Person{基督}那超乎寻常的恩宠，并教导他们人性从何种恶中被解救出来——这恶因用来医治它的药物而变得更糟，并因预防措施而增长——来放大这一点。因此他接着说：\BibleQuote{罪借着诫命更显出是极端的罪恶}\BibleRef{罗 7:13}。你看见这些事如何处处交织在一起吗？他用以控告罪的同一手段，又显示了法律的卓越。他藉着显出罪是何等大的恶，揭示其全部毒素并展现在眼前，而获得的这一点并非小事。因为他正是藉着说\BibleQuote{罪借着诫命更显出是极端的罪恶}来表明这点。也就是说，要使人清楚罪是何等大的恶，是何等毁灭性的事。而这就是由诫命所显明的。他也由此显示了恩宠超越法律——是超越，而不是与法律冲突。因为不要看这一点：那些接受法律的人反而变得更糟；而要思量另一点：法律不仅无意将邪恶推向更甚，反而郑重地试图砍倒既存的邪恶。但如果它没有力量，那么固然要为它的意向加冕，但更要钦崇\Person{基督}的能力，祂废除了、砍断了、拔除了如此多种多样、难以推翻的邪恶的根基。然而，当你听我说到罪时，不要把它想成一个实体性的力量，而是恶行，它来到人身上又不断离去，在发生前不存在，发生后便又消失。这就是法律被赐下的原因。没有任何法律是为终结自然的事情而赐下的，而是为纠正故意邪恶的行为。连外邦的立法者也明白这点，全人类一般也都知道。因为只有出于邪恶的恶行，他们才着手纠正，而不试图根除与本性一起赋予我们的那些恶；因为这是他们做不到的。因为自然的事情保持不变（\Emph{\Person{亚里士多德}\WorkTitle{伦理学}卷二第一章}），正如我们在其他讲道中常常告诉你们的。\par

因此，让我们放下这些争辩，重新在劝勉中操练自己。或者更确切地说，这最后的部分也属于那些争辩。因为如果我们驱除了邪恶，也该引入德行；藉此我们将清楚教导，邪恶并非天生的恶，并能轻易堵住那些探询恶之起源者的口，不仅藉言辞，也藉行为，因为我们与他们共享同一本性，却从他们的邪恶中得解脱。所以，不要只看德行的劳苦，也要看其成功的可能性。只要我们认真对待，它立时会变得轻省而甘甜。但若你向我提及恶习的快乐，也请说出它的结局：因为它通向死亡，正如德行引导我们走向生命。或者，你若愿意，让我们在结局之前就审视两者；我们会看到，恶习伴随着极大的痛苦，而德行则带着极大的喜乐。因为，请问有什么比\OriginalTerm{不良的良心}{bad conscience}更痛苦？有什么比\OriginalTerm{良好的希望}{good hope}更令人愉悦？确实，没有什么比恶事的预期更深刻刺痛、压迫我们；也没有什么比良好的良心更能支撑我们，几乎赐予我们翅膀。这一点，我们甚至可以从眼前所见之事得知。因为那些囚禁在狱中、等待判决的人，即使享受无数次奢华，其生活也比那些在路旁乞讨、却良心无愧的人更为痛苦。因为对可怕结局的预期不让他们感受到手中的快乐。我为何只提囚犯？看那些不在狱中、却享有好运但良心不安的人，那些为面包劳作、终日辛劳的工匠，他们的处境远胜于前者。为此我们也说，角斗士是何等可怜（尽管我们在酒馆中看到他们醉醺醺、奢侈、暴饮暴食），并称他们为最悲惨的人，因为他们必须期待的结局之灾难，远非那种快乐可比。如果这种生活对他们而言似乎令人愉悦，请记住我不断告诉你们的：一个生活在恶习中的人不逃离恶习的痛苦与苦恼，并不奇怪。因为你看，连如此可憎之事，对行恶之人看来却似乎甘美。但我们并不因此说他们多么幸福，恰恰相反，这正是我们认为他们可怜的原因，因为他们不晓得自己所处的邪恶。至于奸淫者，你会怎么说？他们为了片刻的快乐，忍受着可耻的奴役、金钱的损失和永久的恐惧（贺拉斯 \WorkTitle{讽刺诗集} II. vii. 58-67），实际上过着\OriginalTerm{加音}{Cain}般的生活，甚至更糟：充满对现在的恐惧，对未来战栗，猜疑朋友和敌人，猜疑知情者和不知情者；即便入睡，也不得摆脱这场斗争，不良的良心为他们编织出充满各种恐怖的梦境，以此惊吓他们。贞洁之人则截然不同，他度过今生，无拘无束，完全自由。那么，将那些恐惧的种种波动与片刻的快乐相权衡，将\OriginalTerm{节制}{continence}的短暂辛劳与整个生命的平静相权衡，你会发现后者比前者更多喜乐。至于那些一心掠夺、侵占他人财物的人，请告诉我，他岂不需要忍受无数的痛苦：奔走、谄媚奴隶、自由人、守门人；恐吓、威胁、无耻行事；警醒、颤抖、痛苦、怀疑一切？藐视财富的人则截然不同，他也享有丰富的快乐，生活无所畏惧，完全安全。如果有人逐一检视恶习的其他例子，会发现许多烦恼和众多暗礁。但更重要的是，在德行中，困难在先，快乐在后，因此劳苦甚至得以减轻。在恶习中则相反：快乐之后是痛苦和惩罚，这样，连快乐也被消除。因为正如等候冠冕的人感受不到当前的烦恼，同样，在快乐之后期待惩罚的人无法拥有纯粹的喜乐，因为恐惧搅乱了一切。或者更确切地说，如果有人仔细审视，即使在随之而来的惩罚之前，在恶习刚被大胆涉足的那一刻，也会感受到极大的痛苦。如果你愿意，就让我们以侵占他人财物的人为例来考察。或者那些以任何方式聚敛钱财的人，撇开恐惧、危险、战栗、痛苦、忧虑等一切，假设一个人毫无烦恼地致富，并确信能保持现状（尽管他无法做到，但暂且假设以作争论）。那么，拥有这么多财富，他能从中得到何种快乐？相反，恰恰是这财富本身不让他快乐。因为只要他还贪求其他东西，他就仍在受折磨。因为\OriginalTerm{欲望}{desire}在停止时才带来快乐。例如，当我们喝足所需时，感到舒畅；但若持续口渴，即使喝尽世上所有泉源，我们的痛苦只会增长；即使饮干万条河流，我们的惩罚状态更为痛苦。你也是如此，即使你得到全世界的财富，仍贪求不止，你品尝的越多，你的惩罚越大。因此，不要以为聚敛巨资会带来任何快乐，反而要避免致富。但若你贪求致富，你将永远受鞭笞。因为这是一种无法达到目标的欲望；你走得越远，离终点越远。\par

这岂不是一个悖论、一种错乱、一种极端的疯狂吗？那么让我们远离这第一个恶，或者更好，让我们连碰都不碰这贪念。然而，如果我们已经碰了它，就让我们从它的起首（\OriginalTerm{起首}{\GreekText{προοιμίων}}）立即跳开。因为这是\BibleBook{}{箴言}的作者在谈论妓女时给我们的劝告：他说：\BibleQuote{跳开，不要停留，不要走近她家的门口}（\BibleRef{箴 5:8}）。我对你们也要说同样的话，关于贪爱金钱。因为如果逐渐进入，你陷入这疯狂的大海，你将不能轻易脱身，就像在漩涡中，无论你如何挣扎，也难以摆脱；同样，落入这更坏的贪财深渊后，你将毁灭自己及其所有的一切。（\BibleRef{宗 8:20}）所以我的劝告是：我们要警醒于开端，避免小恶，因为大恶由小恶生出。因为一个人若习惯于对每个罪说：\Quote{这没什么！}就会渐渐完全毁灭自己。总之，正是这句话引入了恶习；给强盗（五份抄本作\Quote{魔鬼}）敞开了门；推倒了城墙；就是我们每犯一罪时说：\Quote{这没什么！}同样，在身体方面，最大的疾病也是在轻视小病时产生的。如果\Person{厄撒乌}起初没有背叛他的长子名分，他就不会变得不配承受祝福。如果他没有使自己不配承受祝福，他就不会产生犯下杀害兄弟之罪的欲望。如果\Person{加音}没有贪爱首位，而是将这事留给天主，他就不会沦为次位。再者，当他处于次位时，如果听从了劝告，他就不至于生出杀意。再者，如果在杀人之后，当天主呼唤他时，他悔改了，并且没有以不敬的方式回答，他就不必遭受随后的灾祸。但如果那些在法律之前的人，因这种懈怠而堕入苦难的深渊，那么想想我们这些蒙召参与更大竞赛的人，若不严格督责自己，并迅速在恶行的火花点燃整堆柴薪之前将其熄灭，将会有何等结局？以我所言为例。你习惯发虚誓吗？不要只止于此，而要完全弃绝一切誓言，这样你就不再需要烦恼。因为对一个人来说，发誓而不发虚誓，远比完全禁戒发誓更难。你是一个侮辱和骂人的人吗？还是一个打人的人？你为自己立下规矩：丝毫不动怒、不争吵，这样，连根拔除，果实也会被除掉。你好色放荡吗？你要为自己定下规矩：连看女人都不要（\BibleRef{约 31:1}），也不要上剧场，也不要为所见其他人的美貌而烦恼。因为不看身材好的女人，远比看了并摄入淫欲后再驱除由此而来的搅扰容易得多，因为在起首（\OriginalTerm{起首}{\GreekText{προοιμίοις}}）时挣扎更容易。或者不如说，如果我们不向敌人敞开门户，也不接纳邪恶（\OriginalTerm{邪恶}{\GreekText{κακίας}}）的种子，就根本不需要挣扎。因此基督斥责那以不洁目光看女人的男人（\BibleRef{玛 5:28}），为的是在敌人变强之前，使我们免于更大的劳苦，吩咐我们在还能轻易逐出他时，就把他赶出房屋。因为何必多添麻烦、与敌人纠缠呢？我们本可在无纠缠时竖立凯旋柱，并在角力之前夺取奖赏。因为不看美女的麻烦，远不如看着时克制自己那么大。或者不如说，前者根本不算麻烦，而看后却带来巨大的辛劳与劳苦。既然这麻烦较小（大多数抄本附加：\Quote{对不能自制的人}），或者不如说根本无劳苦、无麻烦，反而有更大的收获，我们为何还要费力跳进无数罪恶的海洋呢？此外，不看女人的人，不仅更轻易地克服这种情欲，而且以更高的纯洁度克服；而看的人，则更麻烦地解脱，并且不无某种污点，即如果他确实能解脱的话。因为不看美女身形的人，也从可能产生的情欲中保持纯洁。但贪恋观看的人，在首先使理性屈辱、以无数方式玷污之后，还必须驱逐那来自情欲的污点——如果他确实驱逐它的话。所以基督为了防止我们这样受苦，不仅禁止杀人，而且禁止发怒；不仅禁止奸淫，而且禁止不洁的注视；不仅禁止伪证，而且禁止一切起誓。祂使人德的标准不止于此，在颁布这些诫命之后，祂更进一步。因为祂使我们远避凶杀、命令我们免于愤怒之后，祂命令我们甚至准备受亏，并不要只准备忍受攻击者所加的一切，更要更进一步，以我们自己的基督徒精神（\OriginalTerm{自己的哲学（基督徒精神）}{\GreekText{τἥς} \GreekText{οἱκείας} \GreekText{φιγοσοφίας}}）的洋溢战胜他极端的疯狂。因为祂所说的并非：\Quote{如果有人打你的右脸，你要高贵地忍受，默不作声；}而是还要加上把另一面也给他。因为祂说：\BibleQuote{连另一面也转给他。}（\BibleRef{玛 5:39}）这就是辉煌的胜利：给他超乎他所愿的，并以自己坚忍的丰盈超越他邪恶欲望的界限。因为这样你就能制止他的疯狂，并且从第二次行动再次获得第一次行动的赏报，同时还制止对他的愤怒。你看，在一切情况中，我们有权不遭受恶害，而非他们施加恶害？其实，不仅是不遭受恶害，甚至接受益处（\EditorialNote{抄本异文 \OriginalTerm{\GreekText{παθεἵν} \GreekText{εὖ}}{\GreekText{παθεἵν} \GreekText{εὖ}}，两份抄本}）也在我们权下。这是最真实的奇迹：如果我们持守正见，我们不仅不致受损，反而得益，而且是藉着别人加给我们不公之事本身。试想：有人凌辱了你吗？你有能力使这凌辱变为你的荣耀。因为如果你以凌辱回报，你只是增加耻辱。但如果你祝福那凌辱你的人，你将看见众人把胜利归给你，宣扬对你的赞美。你看出我们被亏负之事如何使我们得益吗？只要我们存此心意。在钱财、殴打以及其他一切事上，都可看到这一点。因为如果我们以相反的方式回报，我们只是为自己编织了双重冠冕：一重为所受的苦，另一重为所行的善。因此，若有人来告诉你：\Quote{某人凌辱了你，并不断向众人说你坏话，}你就当着那告诉你的人称赞那人。因为这样，即使你想报仇，你也能施加惩罚。因为那些听你的人，即使再愚蠢，也会赞美你，并恨那像野兽般凶残的人，因为他毫无理由就使你痛苦，而你即使在受亏时，仍以相反的方式回报他。这样，你就有能力证明他所说的一切都毫无价值。因为感到毁谤毒牙的人，正是以其烦恼证明他自觉所言属实。而那以微笑待之的人，恰好向在场者洗清了自己的一切嫌疑。那么试想，你能从这件事中收获多少好处。首先，你摆脱一切烦恼和麻烦。其次（其实这应摆在首位），即使你有罪，你也脱去它们，正如税吏以谦和地承受法利塞人的谴责所行的那样。此外，藉此操练，你的灵魂会变得英勇（希腊文：哲学式的），并享有无尽的赞美，同时脱去因所言而起的任何嫌疑。但即使你想报复那人，这也将充分实现：一方面天主因他所说的话惩罚他，另一方面在此之前，你的英勇行为对他而言有如致命一击。因为没有什么比我们受辱者以微笑面对侮辱，更能刺伤那侮辱我们之人的心。因此，既然以基督徒的英勇行事会为我们积累如此多的荣耀，那么从气量狭小（\OriginalTerm{气量狭小}{\GreekText{μικροψυχίας}}）出发，我们在一切事上都会遭遇相反的结局。因为我们使自己蒙羞，并在在场者眼中显得对所提之事有罪，使灵魂充满烦恼，使敌人喜悦，激怒天主，并增加我们先前的罪过。既然如此，让我们逃离气量狭小（\OriginalTerm{气量狭小}{\GreekText{μικροψυχίας}}）的深渊，躲进坚忍（\OriginalTerm{坚忍}{\GreekText{μακροθυμίας}}）的港湾，好在此立刻\BibleQuote{使我们的灵魂得享安息}（\BibleRef{玛 11:29}），正如基督所指示的，并藉着天主的恩宠与对人的慈爱等。\par

\SourceRecord{Translated by J. Walker, J. Sheppard and H. Browne, and revised by George B. Stevens.；Nicene and Post-Nicene Fathers, First Series；Vol. 11.；Edited by Philip Schaff.；Buffalo, NY: Christian Literature Publishing Co.,；1889.；<http://www.newadvent.org/fathers/210212.htm>.}

% structure: p0001 source=h1 target=section reason=page-title

\section{罗马书第十三篇讲道}

\BibleRef{罗 7:14}\par

\Quote{因为我们知道法律是属神的，我却是属血肉的，被卖给了罪恶。}\par

在说了大恶发生后，罪借着诫命得到机会变得更加强大，与法律主要目的相反的结果出现，并使听者陷入极大的困惑之后，他接着便解释这些事的原因，首先为法律洗脱任何恶毒嫌疑。免得人听到罪是因诫命而得到机会，且法律来到时罪复活了，并借着它欺骗和杀害人，便以为法律是这些邪恶的根源，他首先有力地提出辩护，不仅洗脱法律的指控，更给它加上极大的赞美。他如此立论，并非出于自己的认可，而是宣告普遍的判断。他说：\Quote{因为我们知道，法律是属神的。} 仿佛在说：这是公认的、自明的事，法律是属神的，它远非罪的根源，也不应为所发生的邪恶受责备。请注意，他不仅为法律洗脱指控，还给予它极大的赞美。因为称它为属神的，显明它是德行的导师、与罪恶为敌；属神的本意就是如此，使人远离各种罪恶。法律也确实如此行，通过恐吓、劝诫、惩罚、纠正、推荐各种德行。那么，既然导师如此卓越，罪从何而生？是从门徒的懈怠而来。于是他接着说：\Quote{我却是属血肉的，} 为我们勾勒出人在法律之下和法律之前的样子。\Quote{被卖给了罪恶。} 因为（他指）死亡也带来了情欲的纷扰。当身体成为可朽的，从此就必然接纳私欲偏情、愤怒、痛苦以及其它一切情欲；这些需要极大的\OriginalTerm{智慧}{\GreekText{φιλοσοφίας}}来防止它们淹没我们，使理性沉溺于罪的深渊。因为它们本身不是罪，但当它们放纵无度时，便产生这种效果。举例而言（让我取其一审视之）：欲望不是罪，但若过度放纵，不愿遵守婚姻的法则，甚至染指他人之妻，这事就变成奸淫，然而并非由于欲望，而是由于过度。请观察保禄的智慧。他赞美法律之后，立即转向更早的时期，以显示我们人类在当初和接受法律时的状态，并显明恩宠的临在是多么必要，这是他每次竭力证明的事。因为当他说\Quote{被卖给了罪恶}时，不仅指那些在法律之下的人，也包括那些在法律之前的人，以及从起初就存在的人。接着他提到他们被出卖和转让的方式。\par

\BibleRef{罗 7:15} \Quote{因为我的行为，我不明了。}\par

\Quote{我不明了}是什么意思？——我不懂。这何时发生过？因为没有人因无知而犯罪。你们看到了吗？如果我们不谨慎地接受他的话，不始终注视宗徒的目的，就会产生无数矛盾。如果他们是因无知而犯罪，就不该受罚。正如他上面所说：\Quote{因为没有法律，罪是死的，} 其意并非他们不知道自己在犯罪，而是他们确实知道，但不那么清楚；因此他们受罚，却不那么严厉。又说：\Quote{我不知贪欲为何物，} 其意并非完全不知，而是指最清晰的认识；又说：\Quote{它在我内发动了各种私欲偏情，} 其意并非诫命制造了私欲偏情，而是罪借着诫命引入了强烈的私欲偏情。同样，这里他说：\Quote{因为我的行为，我不明了，} 其意并非绝对无知；否则他内心的那个人怎能以天主的法律为乐呢？那么，\Quote{我不明了}是什么意思？他指，我头晕目眩，感到被掳走，发现自己受到暴力，不知如何被绊倒。正如我们常说：某人来了，把我带走了，我不知道怎么搞的；此时我们并非以无知为借口，而是要显示一种欺骗、诡计和阴谋。\Quote{我所愿意的，我不去做；我所憎恨的，我反而去作。} 那么，你怎么能说不知道自己在做什么呢？因为你若愿意善，憎恨恶，这就需要完备的认识。由此可见，他说\Quote{我不愿意}，并非否认自由意志，也不是指出任何强制的必然。因为如果我们犯罪不是出于自愿，而是被迫，那么先前的惩罚就不合理。但正如他说\Quote{我不明了}时，不是向我们呈现无知，而是我们所说的；同样，加上\Quote{我不愿意}时，他并非表示必然，而是对所做之事的否定。因为如果他说\Quote{我所不愿意的，我反去做}不是这个意思，他本该接着说\Quote{但我做我被迫强迫去做的事}。因为这才是与意愿和能力相对立的。\OriginalTerm{能力}{\GreekText{ἐξουσία}}\BibleRef{罗 7:17-18} \Quote{那么，现在做这事的已不是我，而是住在我内的罪。我知道，在我内，即在我肉性内，没有善性存在。}\par

关于这段经文，那些指摘肉身、并主张肉身不是天主创造的一部分的人，攻击我们。那么我们该说些什么呢？正如我们从前讨论法律时所说的：正如在那里，他使罪对一切负责，同样在这里也是如此。因为他不是说肉身做了这事，恰恰相反，他说：\BibleQuote{不是我做的，而是住在我内的罪做的。}\BibleRef{罗 7:17}但如果他说\BibleQuote{在它内不住着善，}\BibleRef{罗 7:18}这仍然不是对肉身的指控。因为\BibleQuote{不住着善}这一事实，并不表明它本身就是恶。现在我们承认，肉身不如灵魂伟大，且次于它，但并非相反或敌对，也不是恶；而是它低于灵魂，就如竖琴低于竖琴师，又如船低于舵手。这些并不与引导并使用它们的东西相反，而是完全一致，只是与技艺者不同等尊荣。因此，若有人说，技艺不在于竖琴或船，而在于舵手或竖琴师，这并非指摘工具，而是指出它们与技艺者之间的巨大差异。同样，保禄说\BibleQuote{在我肉身内不住着善，}不是指摘身体，而是指出灵魂的优越。因为灵魂肩负着全部责任或掌舵之职，以及演奏之职。保禄在这里指出，将治理权交给了灵魂，并将人分为灵魂和身体这两部分后，他说，肉身缺乏理性，没有辨别力，属于被领导之物，而非领导之物。但灵魂更有智慧，能看清该做什么、不该做什么，却无法随心所欲地勒住这匹马。这不仅是针对肉身的指控，也是针对灵魂的，因为灵魂确实知道该做什么，却未能在实践中执行它认为最好的事。他说：\BibleQuote{因为立志行善由得我，但行出来由不得我。}\BibleRef{罗 7:18}在这里，在\BibleQuote{行出来由不得我}这句话中，他并非指任何无知或困惑，而是指罪的一种阻挠和狡猾的攻击，他在接下来的话语中更加清晰地指出了这一点。\par

\BibleQuote{因为我所愿意的善，我反而不做；我所不愿意的恶，我倒去做。若我做我所不愿意的，就不是我做的，而是住在我内的罪做的。}\BibleRef{罗 7:19-20}\par

你们看，他如何使灵魂的本质和肉身的本质都免于指控，而将一切归咎于罪恶行为？因为如果灵魂不愿作恶，它就得释放；如果他自己不去行恶，身体也获得自由，整个责任就可以归咎于邪恶的道德选择。现在灵魂、肉身的本质与这个选择并非同一回事，因为前两者是天主的工程，而后者是我们自己朝向我们所愿意引导的方向的运动。因为意志确实是\OriginalTerm{天生的}{\GreekText{ἔμφυτον}}，来自天主；但以这种方式愿意则是我们自己的，来自我们自己的心智。\par

\BibleQuote{这样，我发现一条规律：就是我愿意为善的时候，恶常在我面前。}\BibleRef{罗 7:21}\par

他所说的不是十分清楚。那么，所说的究竟是什么呢？他说：我在良心上赞美法律，并且发现只要我愿意行善，法律就站在我这一边，并且激励这种意愿。因为正如我以此为乐，法律也赞美我的决定。你们看见他如何显出善恶的知识是我们本性中原初和基础的部分，而梅瑟法律赞美它，并从它得到赞美吗？因为上文他没有说\Quote{我由法律得教训}，而是说\Quote{我赞同法律}；后来又没有说\Quote{我受它教导}，而是说\Quote{我喜爱它}。那么\Quote{我喜爱}是什么意思？就是说，我认为它是正确的，正如它在我愿意行善时也赞同我一样。因此，愿意行善和不愿意作恶从一开始就是我们本性的基础部分。但当法律到来时，它立即成为恶事上更强烈的控告者，以及善事上更伟大的赞美者。你们是否观察到，他处处作证法律具有某种加强和附加的优势，但仅此而已？因为虽然法律赞美，我也喜爱它，并愿意行善，但\Quote{恶}仍然\Quote{常在我面前}，它的作用并未被消除。因此，对于一个决心行善的人，法律只起到辅助的作用，因为它怀有与他相同的意愿。随后，由于他表达得不够清晰，在继续论述时，他给出了更清晰的解释，表明恶如何常在，以及法律又是如何只成为那些心思行善之人的法律。\par

\BibleQuote{因为我按着内心的人，喜爱天主的法律。}\BibleRef{罗 7:22}\par

他的意思是：我在此之前就已知道什么是善，但当我发现它被书写下来时，我就赞美它。\par

\BibleQuote{但我看另一个法律，在我的肢体内与我心智的法律交战。}\BibleRef{罗 7:23}\par

这里，他再次将罪称为一条法律，与另一条法律交战，这不是出于良好秩序，而是因为那些顺从它的人对它所表示的严格服从。正如他将\OriginalTerm{主人}{\GreekText{κύριον}}之名赐给玛门（\BibleRef{玛 6:24}），将\Quote{神}之名赐给肚腹（\BibleRef{斐 3:19}），不是由于它们本身配得，而是由于隶属者的极端谄媚；同样，他称罪为一条法律，是因那些对它如此卑躬屈膝、不敢离弃它的人，就像那些接受了法律的人害怕离开法律一样。他意指这条法律与自然的法律相对立；这就是所谓\BibleQuote{我心思的法律}的含义。接着他描述了一场列阵与战斗，并将整个奋斗归诸自然的法律。因为梅瑟的法律是后来加添的；然而这两者，一个在教导，一个在赞美正确之事，在这场战斗中却没有产生什么大效果；罪恶的奴役如此巨大，它克服并占了上风。保禄陈述这一点，并表明它那\OriginalTerm{彻底胜利}{\GreekText{κατὰ} \GreekText{κράτος}}的胜利，就说：\BibleQuote{我看见另一个法律，与我心思的法律交战，并把我掳去。}他不仅用了\Quote{战胜}这个词，还说\BibleQuote{把我掳去，使我隶属于罪的法律。}他没有说肉体的倾向或肉体的本性，而是说\BibleQuote{罪的法律}。那就是奴役、权势。那么，他说\BibleQuote{在我肢体中的}是什么意思？这是什么？这当然不是使肢体成为罪，而是尽可能使它们与罪区分。因为在某物中的东西与该物本身是不同的。因此，正如诫命也不是邪恶的，因为罪藉着它得了机会，同样，肉体的本性也不是邪恶的，即使罪藉着它征服我们。因为这样，灵魂也成了邪恶的，而且更是如此，因为它在行动上拥有权柄。但事实并非如此，绝非如此。因为如果一个暴君和强盗占领了一座华丽的府邸和一座王宫，这情况也不会给那房屋带来任何耻辱，因为全部罪责都在于那些谋划此等行为的人。但真理的敌人，除了他们的不敬，还不自觉地陷入极大的不合理。因为他们不仅控告肉体，而且也贬低法律。但若肉体是邪恶的，法律就是好的。因为肉体与法律交战，反对它。然而，如果法律不好，那么肉体就是好的。因为即使按照他们自己的说法，肉体也与法律争战、对抗。那么，他们怎能断言两者都属于魔鬼，而把彼此对立的东西摆在我们面前呢？你看到，除了他们的不敬，他们的不合理是何等之大吗？但这样的教义不是教会的教义，因为教会只谴责罪；并且她宣称天主所赐的两条法律，即自然的和梅瑟的法律，都是与罪为敌，而非与肉体为敌；因为她否认肉体是罪，因为肉体是天主的化工，并且如果我们节制地生活，它对于德行也是非常有用的。\par

第24节。\BibleQuote{我真不幸啊！谁能救我脱离这死亡的身体呢？}\BibleRef{罗 7:24}\par

你们注意到罪恶的奴役是多么大吗？它甚至战胜了一个喜爱法律的心智。因为他意谓没有人可以反驳说：我憎恨法律，厌恶它，所以罪胜过了我。但他说：\BibleQuote{我喜爱它，且同意它，}并归向它寻求庇护，然而它却没有能力拯救那归向它的人。但基督甚至拯救了那逃离祂的人。看，恩宠拥有何等大的优势！可是宗徒并没有这样陈述；而是仅带着叹息和极大的哀号，仿佛没有任何帮助者，他以自己的困惑指出基督的大能，说道：\BibleQuote{我真不幸啊！谁能救我脱离这死亡的身体呢？}法律无能为力：良心证明自己不足，虽然它赞美了善事，并且不仅是赞美，甚至还与善的对立面争战。因为藉着\Quote{争战}这个词，他表明他自己也列阵反对它。那么，从何处可以希望救恩呢？\par

第25节。\BibleQuote{我藉着我们的主耶稣基督，感谢天主。}\BibleRef{罗 7:25}\par

请注意他是如何表明恩宠与我们同在最必要的，以及此处的善行同时属于圣父和圣子。因为他感谢的是圣父，但圣子却是这感谢的原因。当你听见他说：\BibleQuote{谁能救我脱离这死的身体呢？}不要以为他在指控肉体。因为他没有说\Quote{罪的肉身}，而是说\Quote{死的身体}：即必朽的身体——那被死亡制伏了的，而不是那产生死亡的身体。这并非证明肉体的邪恶，而是表明它所受的损毁（\OriginalTerm{损毁}{\GreekText{ἐπηρείας}}）。如同一个人被野蛮人俘虏，我们说这人属于野蛮人，不是因为他自己是野蛮人，而是因为被他们扣留；同样，身体被称为\Quote{死}的，是因为被死所压制，而不是因为它产生死。因此，他自己所愿摆脱的并非身体，而是这必死的身体，暗示（如我常说的）正因身体变得易于受苦，也成了罪轻易猎取的对象。那么，有人会说，在恩宠时代之前，罪的奴役如此严重，人为何因犯罪而受惩罚？因为他们领受了即使在罪的统治下也能遵行的诫命。因为他并没有把他们引向最高等的交往方式，而是允许他们享有财富，不禁止娶多个妻子，在正义的原因下发泄怒气，以及有节制地享乐。\BibleRef{玛 5:38} 这种屈就已如此之大，以至于成文法比自然法要求更少。因为自然法规定一个男人终生只与一个女人结合。基督在以下话语中表明：\BibleQuote{那起初造人的，是造男造女。}\BibleRef{玛 19:4} 但梅瑟法律既不禁止休妻另娶，也不禁止同时有两个妻子！\BibleRef{玛 5:31} 此外，律法中还有许多其他诫命，人们可以看到律法之前的人完全依自然法而行。因此，生活在旧约之下的人，受如此温和的法律约束，并无困难。如果他们在这种条件下仍然不能得胜，那只能怪他们自己的懈怠。所以保禄感恩，因为基督对这些事毫无苛刻，不仅不追究这微小的要求，甚至使我们能够面对更大的赛程。因此他说：\BibleQuote{我藉着耶稣基督感谢我的天主。}他放下众所同意的救恩，从已经证明的论点转向另一个更深的论题，其中他指出我们不仅摆脱了从前的罪，而且对于未来也成为不可战胜的。因为他说：\BibleQuote{如今那些在基督耶稣内的，不随从肉身而生活的人，已无罪可定。}然而，他在说出这话之前，先回顾了我们从前的光景，说：\BibleQuote{这样看来，我自己在理智上服事天主的法律，在肉身内却服事罪的法律。}\par

第八章　第一节　\BibleQuote{这样，在基督耶稣内，已无罪可定了。}\BibleRef{罗 8:1}\par

接着，因为许多人即使在领洗后仍陷于罪的事实造成了困难（\OriginalTerm{面临困难}{\GreekText{ἀντέπιπτεν}}），所以他急忙予以应对，说并非仅仅\Quote{在基督耶稣内的}，而是加上\Quote{不随从肉身而生活的人}，以此表明一切后来的行为都出于我们的懈怠。因为如今我们有能力不随从肉身，而在从前则是一项艰巨的任务。随后，他通过下文给出另一个证明：\par

第二节　\BibleQuote{因为生命之神的法律在基督耶稣内释放了我。}\BibleRef{罗 8:2}\par

他这里所称的\Quote{神的法律}，是指圣神。如同他称罪为罪的法律，他称圣神为圣神的法律。然而他也称梅瑟法律为\Quote{属神的}，如他说：\BibleQuote{我们知道法律是属神的。}那么区别何在？巨大且无限。因为梅瑟法律是属神的，而这是圣神的法律。两者区别是什么？前者仅仅是藉着圣神赐下的，但后者还使领受者获得丰盛的圣神。因此他称它为生命之神，与罪的法律相对，而不是与梅瑟的法律。因为他说它\Quote{使我摆脱了罪与死的法律}，这里指的并非梅瑟的法律，因为他从未称它为罪的法律：他岂会称那常被他说成\Quote{正义而神圣}且能摧毁罪的法律为罪的法律呢？他指的是那与理智的法律交战的法律。圣神的恩宠终止了这场剧烈的战争，它杀死罪，使争战对我们变得轻松，从一开始就加冕我们，然后用丰富的帮助引领我们战斗。接着，如同他惯于从圣神转向圣子和圣父，将我们的一切归于天主圣三，他在这里也是如此。在说了\BibleQuote{谁能救我脱离这死的身体呢？}之后，他指出圣父藉着圣子行了这事，然后又指出圣神与圣子一同行动。因为他说：\BibleQuote{生命之神的法律在基督耶稣内释放了我。}接着又指向圣父与圣子：\par

第三节　\BibleQuote{凡法律不能成就的，因它为肉身所弱化，天主派遣了自己的儿子，带着罪恶肉身的形状，且为罪，在肉身中定了罪。}\BibleRef{罗 8:3}\par

他又似乎在贬抑法律。但若仔细留意，他甚至高度赞扬法律，因为显示法律与基督和谐一致，并优先注重相同的事。他并非谈论法律的恶劣，而是谈论\BibleQuote{它所不能行的}；所以又说\BibleQuote{因它软弱}，而不是\BibleQuote{因它有害或奸诈}。他也没有把软弱归咎于法律，而是归咎于肉体，正如他所说\BibleQuote{因通过肉体而软弱}，此处使用\Quote{肉体}一词并非指本质和实体的本身，而是将它用来指更属肉性的心思。这样，他既使身体也摆脱了责备。不仅这样，而且凭借下文也是如此。因为假设法律是敌对的一方，基督怎会来帮助它，满足它的要求，并通过在肉体上定罪罪来助它一臂之力？因为这是所缺少的，因为主早已在灵魂中定它罪了。那么，法律成就了更大的事，而独生子成就了较小的事吗？当然不是。因为天主也主要是那成就者，他赐给了我们自然法律，并将成文的法律加给它。再者，若较小的没有提供，更大的也无用。知道什么事该做，若不实行又有什么用？毫无用处，因为那只会是更大的定罪。所以，那拯救灵魂的，也是使肉体易于约束的。因为教导容易，但此外还显示一条轻松成就这些事的途径，这才令人惊奇。独生子正是为此而来，在他将我们从这困难中解救出来之前，他并未离开。但更伟大的是胜利的方法：他并没有取别的肉体，而是取了这正充满困扰的同一个肉体。这就像一个人在大街上看到一个下流卑贱的女人被殴打，就说自己是她的儿子，尽管他是国王的儿子，这样就能让那女人从施暴者手中得到自由。他真的这样做了，他承认自己是人子，并支持肉体，定罪的罪。但他不忍心再击打肉体；或者更确切地说，他以自己的死击打肉体，但在这行动中，被击打的肉体并没有被定罪而消亡，而是那一直击打人的罪被定罪了。这是最大的奇迹。因为如果胜利不在肉体中发生，就不会如此令人惊奇，因为法律也成就了这一点。但奇妙的是，他的凯旋是\OriginalTerm{带着肉体}{\GreekText{μετὰ} \GreekText{σαρκὸς}}建立的，那曾无数次被罪打败的，竟自身获得了对罪的辉煌胜利。看啊，发生了多么奇异的种种！一是罪没有征服肉体；二是罪被征服了，而且是被肉体征服的。因为不被征服与征服那不断摧毁我们的，是不同的。三是它不仅征服了罪，还惩罚了它。因为通过不犯罪，它避免了被征服；但通过死亡，他战胜并定罪了罪，使那从前轻易被罪嘲弄的肉体，成了罪明显畏惧的对象。这样，他立刻卸除了罪的权能，废除了罪所引入的死亡。因为只要它抓住罪人，它就正义地一直走向终点。但找到无罪的身体后，当它把身体交付死亡时，它就被定为不义。你看到胜利有多少证据吗？肉体没有被罪征服，甚至征服并定罪它，不单定罪它，而且定罪它为有罪。因为证实了它的不义之后，他进而定罪它，不仅是靠权力和力量，甚至是按正义的法则。这就是他所说\BibleQuote{因罪在肉体上定罪了罪}的意思。就好像他说过他已经证实它犯了大罪，然后定罪它。所以你看，被定罪的是罪，而不是肉体，因为肉体甚至被加冕，并且必须对另一方宣判。但如果他说他派遣圣子带着\BibleQuote{相似于肉身的形状}，不要因此以为他的肉身是别样的。因为他称之为\BibleQuote{有罪的}，这就是为什么他用了\Quote{相似}这个词。因为基督并没有有罪的肉体，而是确实相似于我们有罪的肉体，却是无罪的，并且本性与我们相同。因此，即使从这一点也可清楚看出，肉体在本性上并非邪恶。因为基督并不是取了另一个代替先前的，也不是改变了这同一个的本体而使肉体重新得胜：他让它留在自己的本性中，却使它戴上战胜罪的荣冠，然后在胜利后复活它，使它不朽。那么，可以说，这些事发生在这肉体上，与我何干？不，它与你关系重大。因此他接着说：\par

第4节（\BibleRef{罗 8:4}）\BibleQuote{为使法律的正义，成全在我们这些不随从肉体生活的人身上。}\par

\Quote{义}这个字是什么意思？它是指终向、目标和善行。因为它的目的和它要求的是什么？那就是无罪。这如今已为我们成全了（\OriginalTerm{已为我们成全}{\GreekText{κατόρθωται} \GreekText{ἡμἵν}}），借着基督。而抵抗它并胜过它，是来自\Person{祂}。但享受胜利却是我们的事。那么，从此以后我们就不再犯罪了吗？除非我们变得极其松懈和怠惰，否则我们永远不会。这就是为什么他补充说：\Quote{不随从肉性的人}。因为恐怕你听到基督已将你从罪的战争中释放，并借着罪已在肉身上被\Quote{定罪}，而使法律的义之要求（\OriginalTerm{义之要求}{\GreekText{δικαίωμα}}）在你们身上得以满足，你就拆毁所有防御；因此，在那处地方，在说了\Quote{因此，如今为那些在基督耶稣内的人，已无处可定}之后，他补充道\Quote{不随从肉性的人}；而在这里，在说了\Quote{为使法律的义之要求得以满足在我们身上}之后，他又继续同一件事；或更确切地说，不仅是同一件事，而且是一件更强烈的事。因为说了\Quote{为使法律的义得以满足在我们这些不随从肉性的人身上}之后，他继续道\Quote{却随从圣神}。\par

这显示，我们不仅有义务避免恶行，而且还要以善行来装饰（\OriginalTerm{装饰}{\GreekText{κομᾅν}}）自己。因为赐予你冠冕是\Person{祂}的事；但牢牢持守它却是你的事。因为法律的义，就是使人不致承担其诅咒，基督已为你成就了。所以不要对这如此伟大的恩赐变节，而要守护好这美好的宝藏。因为在这段经文中，他表明圣洗池（Font）不足以保证我们得救，除非我们从其中出来之后，显示配得上这份恩赐的生活。因此，他再次为法律辩护而说这些话。因为一旦我们顺从基督，就必须用一切方式和方法，使基督所成全的法律的义，常在我们内，而不致落空。\par

\BibleQuote{因为随从肉性的人，思念肉性的事。}\BibleRef{罗 8:5}\par

然而这也不是对肉性的贬损。因为只要它保持在自己的位置，就不会发生任何错乱。但当我们让它完全任意而行，它就越过其应有的界限，起来反抗灵魂，那时它就毁灭并败坏一切，但这并非由于它自己的本性，而是由于失度以及随之而来的混乱。\Quote{但随从圣神的人，思念圣神的事。}\par

\BibleQuote{因为思念肉性的事就是死亡。}\BibleRef{罗 8:6}他并不是在谈论肉性的本质，或身体的实质，而是指肉性的\Quote{思念}，这是可以改正和消除的。他这样说，并不是赋予肉性本身的推理能力。绝不是。而是要表达心智较低的运动，并以下等部分来命名，正如他常常习惯称呼整个的人，视为有灵魂的人为\Quote{肉性}一样。\BibleQuote{但思念圣神的事，}这里他又说到属神的思念，正如他稍后所说：\BibleQuote{但那洞察人心的，知道圣神的心思，}\BibleRef{罗 8:27}；他指出这带来的许多祝福，既在今生，也在来世。因为因思念肉性的事所带来的邪恶，远远少于属神的思念所提供的祝福。他在这句话中指出了这一点：生命与平安。两者形成对比：因为他说思念肉性的事导致死亡，而思念圣神的事带来生命与平安。在提到平安之后，他继续说，\par

第七节。\Quote{因为随从肉性的思念，是与天主为仇：}\BibleRef{罗 8:7}这比死亡更糟。然后为显示它既是死亡又是仇恨，他说：\Quote{因为它不服从天主的法律，}\Quote{也决不能服从。}但听到\Quote{也决不能服从}时不要烦恼。因为这个难题容易解决。因为他在这里称的\Quote{随从肉性的思念}就是属地的、粗俗的、对今生的事及其邪恶行为热心渴望的推理（或\Quote{思维方式}，\OriginalTerm{推理}{\GreekText{λογισμὸν}}）。他说的正是这个\Quote{也决不能}服从天主。那么，如果坏人不可能变为好人，还有何得救的希望呢？他不是这个意思。否则，保禄如何变成这样？悔改的强盗，或默纳舍，或尼尼微人，或达味跌倒后如何重振？伯多禄否认后如何起来？\BibleRef{格前 5:5}生活在淫乱中的人如何能被列在基督的羊栈中？\BibleRef{格后 2:6}迦拉达人如何能从\Quote{失去了恩宠}（\BibleRef{迦 5:4}）恢复到先前的尊严？那么他说的是，一个邪恶的人不可能变成好人，而是一个继续邪恶的人不可能服从天主。然而人转变成为好人并服从祂，是容易的。因为他不说人不能服从天主，而是说恶行不能成为善。如同他也在这福音中说：\Quote{坏树不能结好果子}（\BibleRef{玛 7:18}），这不是禁止从善到恶的转变，而是说持续行恶不能结出好果子。因为祂不是说坏树不能变成好树，而是说它继续坏时不能结好果子。因为树能改变，祂在此处和另一个比喻中显示，当祂将莠子比喻为变成麦子时，也因此禁止拔除它们：\Quote{免得你们把麦子也拔出来}（\BibleRef{玛 13:29}），也就是，那将从它们生出的（\OriginalTerm{产生}{\GreekText{γίνεσθαι}}，四份抄本作\OriginalTerm{生出}{\GreekText{τίκτεσθαι}}）。那么，他所说的\Quote{随从肉性的思念}指的就是恶，而\Quote{随从圣神的思念}指所赐的恩宠及它在正确的心意决断中可察觉的功效，在此段中任何部分他都没有讨论实体和本质，而是讨论德行和恶行。因为你在法律下无力做到的，现在他指，你将能够做到，正直地前行，中间没有跌倒，只要你抓住圣神的帮助。因为不仅是不随从肉性，我们还必须随从圣神，因为离恶并不保证我们的救恩，我们还必须行善。这将实现，如果我们将自己的灵魂交给圣神，并说服我们的肉身认识其适当位置，因为我们这样做也会使它成为属神的；同样，如果我们懈怠，我们也会使我们的灵魂成为属肉性的。因为不是本性的必然性把礼物放在我们里面，而是自由的选择把它放在我们手中，所以此后它取决于你成为这样或那样。因为祂那一方面已经完成了一切。因为罪不再与我们的理智法律交战，也不再像先前那样掳掠我们，因为那个状态已经终结并破裂，情欲在圣神的恩宠面前战栗恐惧。但如果你熄灭光，扔出执缰者，赶走舵手，那么此后就把颠簸归咎于你自己。因为现在德行已变得更容易（因此我们也有更严格的虔诚生活的义务），思考当法律盛行时人类的状态，以及现在恩宠照耀之后的状态。先前似乎对任何人都不可行的事——童贞、轻视死亡和其他更大的苦难——现在在全世界各个角落都充满活力，不仅在我们中间，而且在司提亚人、色雷斯人、印度人、波斯人和其他几个野蛮民族中，都有童贞女团体、殉道者群体、隐修士团体，这些现在甚至比已婚者还多，还有严格的斋戒和彻底的舍弃财产。这些事，除了一两个例外，生活在法律下的人从未在梦中想象过。既然你看到真实状态比任何号角更响亮地宣告，不要让自己变得软弱和辜负如此大的恩宠。因为即使在信德之后，懈怠的人也不可能得救！因为摔跤变得容易，是为了让你争战并得胜，而不是让你睡觉，或者滥用恩宠的伟大作为懈怠的理由，从而再次在以前的泥沼中打滚。因此他接着说：\par

第八节。\BibleQuote{所以，随从肉性的人，不能得天主的欢心。}\BibleRef{罗 8:8}\par

那么，有人会说，我们是否要将身体切碎以取悦天主，并逃离肉体？你难道要我们成为杀人者，以此引导我们修德？你看见按字面理解这些话会产生何等矛盾。因为在这段经文中，他所说的\Quote{肉体}不是指身体或身体的本质，而是指那属肉、属世俗的生活，它完全沉溺于纵情和奢侈，从而使人整个成为肉体。因为正如那些拥有圣神翅膀的人使身体也成为属神的，同样那些与此分离、成为肚腹和享乐奴隶的人，使灵魂也成为肉体的，并非他们改变了灵魂的本质，而是他们败坏其高贵的出身。这种说话方式在\WorkTitle{旧约}许多地方也能见到，用肉体来指代粗俗和属尘世的生活，它纠缠于不正当的享乐之中。因为他对\Person{诺厄}说：\BibleQuote{我的圣神不常留在这等人里面，因为他们是血肉。}\BibleRef{创 6:3}然而\Person{诺厄}自己也为肉体所包围。但这并非抱怨之处，即被肉体包围，因为这乃是出于本性，而是指选择了肉性的生活。因此\Person{保禄}也说：\Quote{凡在肉体里的人，不能使天主喜悦。}他接着说：\par

第9节。\BibleQuote{但你们不在肉体内，而在圣神内。}\BibleRef{罗 8:9}\par

这里他同样不是绝对地指肉体，而是指那一类肉体，即处于情欲的旋涡和奴役中的肉体。那么，有人会说，他为何不这样说，也不指出区别呢？这是为了激励听者，并表明正确生活的人甚至不在身体内。因为既然在某种程度上人人都清楚属神的人不在罪恶中，他便说出更大的真理：属神的人不仅不在罪恶中，而且从此以后甚至不在肉体中，因为他从那一刻起已成为天使，升入天堂，从此只是勉强带着身体。如果这就是你贬低肉体的理由，因为他用肉体的名字称呼肉性的生活，那么照此说来，你也要贬低世界，因为罪恶常以世界之名被称呼，正如\Person{基督}也对门徒说：\BibleQuote{你们不属于这世界}\BibleRef{若 15:19}；又对他的弟兄说：\BibleQuote{世界不能恨你们，却是恨我。}\BibleRef{若 15:19}同样，\Person{保禄}后来也要把灵魂称为与天主疏远的，因为他给生活在错误中的人起名为\OriginalTerm{属魂的人}{A.V. natural}\BibleRef{格前 2:14}。但事实并非如此，确实不是这样。因为我们不可只看字面，总要顾及说话者的思想，从而对所说的话获得完全清晰的理解。因为有些事物是善的，有些是恶的，有些是中性的。灵魂和肉体属于中性事物，因为二者都可以变成善或恶。但圣神属于善的事物，绝不会变成其他。再者，肉性的心思，即作恶，属于总是恶的。因为\Quote{它不服从天主的法律。}如果你将灵魂和身体献给那更好的，你就将成为它的一部分。如果你献给那更坏的，那么你就在那里的败坏中有份，这不是因为灵魂和肉体的本性，而是因为那能选择二者的判断。为表明这些事的确如此，并且这些话并非贬低肉体，让我们再次拿起这句话本身，更彻底地分析它。他说：\Quote{但你们不在肉体内，而在圣神内。}那么怎样？难道他们不在肉体中，没有身体而四处行走？这有什么意义？你看，他所暗示的是肉性的生活。那他为何不说\Quote{但你们不在罪恶中}？这是为了使你知道，\Person{基督}不仅消灭了罪恶的统治，甚至还使肉体更少地压垮我们，使之更属神，并非改变其本性，而是赋予它翅膀。因为正如当火与铁相遇时，铁也变成火，虽然仍保有自己的本性；同样，对于相信并拥有圣神的人，肉体从此转入那种运作方式，变得完全属神，各部分都被钉在十字架上，与灵魂以同样的翅膀飞翔，就如这里说话者的身体一样。因此，他蔑视一切放纵和享乐，而在饥饿、鞭打和监禁中找到自己的纵情，甚至在忍受它们时也不感觉痛苦。\EditorialNote{格林多后书第十一章}为表明这点，他说：\BibleQuote{我们这轻微短暂的苦难}\BibleRef{格后 4:17}等等。他如此训练肉体与圣神和谐一致。\BibleQuote{如果天主的圣神住在你们内}\OriginalTerm{若是}{\GreekText{εἴπερ}}。他常用这个\Quote{若是}，不是表达任何怀疑，而是即使他完全确信某事时，也用来代替\Quote{既然}，正如他说：\BibleQuote{如果这是公义的}\BibleRef{得后 1:6}，意思是\Quote{因为天主是公义的，必将患难报应那加患难给你们的人}。又说：\BibleQuote{你们受了那么多的苦，若是白白地受，难道真是白白地受吗？}\BibleRef{迦 3:4}\par

\Quote{现在，若有人没有基督的圣神。}他不说\Quote{若你们没有}，而是以指涉他人的方式说出令人沮丧的话。他说：\Quote{他不属于基督。}\par

第10节。\BibleQuote{若基督在你们内。}\BibleRef{罗 8:10}\par

再一次，他将美好的事归于他们，而令人痛苦的部分则是短暂且插入的。那渴望的对象在其两侧，并且被详细陈述，以便使另一部分黯然失色。他说这话，并非断言圣神就是基督，绝非如此，而是表明拥有圣神的人不仅被称为属于基督的，甚至拥有基督自己。因为圣神在哪里，基督也必在那里。因为无论天主圣三的一位在哪里，整个天主圣三都在那里。因为祂自身是不可分割的，具有最完全的合一。那么，有人会说，如果基督在我们内，会发生什么？\\Quote{身体因罪而死，但圣神因义而成为生命。}你看见没有圣神所带来的巨大灾祸：死亡、与天主为敌、无法满足祂的律法、不能如我们所应然属于基督、缺乏祂的内住。现在也思考拥有圣神所带来的巨大祝福：属于基督，拥有基督自己，与天使媲美（因为这就是致死肉体），并活出不朽的生命，从此持有复活的凭据，轻松奔跑德行的赛程。因为他说身体从此对罪不再活跃，甚至说它死了，以此夸大赛程的容易。这样的人无需劳苦就赢得花冠。此后，为此原因他又加上\\Quote{于罪}，好叫你看见祂所一举废除的是身体的邪恶，而非其本质。因为若后者被废除，许多甚至对灵魂有益的事物也会被废除。然而他所说的不是这样，而是主张当身体还活着且存留时，它却是死的。这就是我们拥有圣子、圣神在我们内的标记：我们的身体在罪的运作上与躺在棺材里的身体毫无区别\\EditorialNote{所以手稿Sav.写\\Quote{身体}。前文略有缺损。}但听到致死不要害怕。因为在其中你拥有真正的生命，之后不会再有死亡：圣神的生命正是如此。它不再屈服于死亡，而是磨损死亡、消耗死亡，并且它所接受的，它保持不朽。这就是为什么他说了\\Quote{身体死了}之后，并没有说\\Quote{但圣神活着}，而是说\\Quote{是生命}，以指出祂（圣神）有能力将此也赋予他人。然后，为鼓励听者，他告诉他生命的原因和证据。这就是义；因为哪里没有罪，死亡也不被看见；但哪里看不见死亡，生命就不可分解。\par

第11节。\\BibleQuote{如果那使耶稣从死者中复活者的圣神住在你们内，那使基督从死者中复活的，也必要藉那住在你们内的圣神，使你们有死的身体复活。}\\BibleRef{罗 8:11}\par

再次，他触及复活的关键，因为这对听者是最令人鼓舞的希望，并借由发生在基督身上的事给他保证。现在你不要害怕，因为你被死身体所包围。让身体有圣神，它必会复活。那么，没有圣神的身体就不会复活吗？那么\\BibleQuote{所有人都要站在基督的审判台前吗？}\\BibleRef{罗 14:10} 或者地狱的报应如何可靠？因为若没有圣神的人复活，地狱就根本不存在了。那么所说的是什么？所有人都要复活，但不是都到生命，而是\\BibleQuote{有的到惩罚，有的到生命}\\BibleRef{若 5:29} 这就是为什么他没有说\\Quote{叫起来}，而是说\\Quote{使人活}\\BibleRef{达 12:2} 这是比复活更大的事，且只赐给义人。他加上这尊荣的原因时说了\\Quote{借着住在你们内的圣神}。因此，如果你在这里赶走圣神的恩宠，没有带着它安然离去，那么即使你复活，也必定灭亡。因为正如祂那时若看见祂的圣神在你内闪耀，就不愿将你交给惩罚，同样，祂若看见圣神熄灭，也不允许你进入婚筵，就如祂没有接纳那些童女一样。\\BibleRef{玛 25:12}\par

不要让你的身体在此世活着，为的是它那时能活着。使它死去，好使它不死。因为如果它继续活着，它就不会活；但如果它死去，那么它就会活。这就是一般复活的情况。因为它必须先死去并埋葬，然后成为不朽的。但这事已在圣洗中完成了。因此它先经历了被钉十字架和埋葬，然后被复活。主身体的情形也是如此。因为那身体也被钉十字架并埋葬（七部抄本作\Quote{死去}）并复活了。因此让我们也这样做：让我们不断在它的作为中将它致死。我并非指在它的实体上——切莫如此——而是指它向恶行的倾向。因为这也是一种生命，或者更确切地说，这才是生命，不受人类常情的束缚，也不做情欲的奴隶。因为凡将自己置于这些统治之下的人，甚至无力活过由此产生的消沉、恐惧、危险以及无数祸患。因为如果死亡必然临到，他早在死亡之前就因恐惧而死了。如果他惧怕疾病、侮辱、贫穷或任何其他不可预料的灾祸，他就已经败坏且灭亡了。还有什么比这种生活更悲惨的呢？然而为圣神而活的人却大不相同，因为他立刻超越恐惧、悲伤、危险和一切变化：并非因为他不会经历这些事，而是（这更伟大）当他受这些事攻击时，他藐视它们。这怎么可能呢？如果圣神持续地住在我们内，这事就会成就。因为他说的不是暂时的停留，而是持续的居住。因此他没有说\Quote{那圣神}曾居住，而是说\Quote{那在我们内居住的}，以此指出持续的同在。那么，最真实地活着的人，就是对此世死去的人。因此他说\Quote{圣神是生命，因为正义}。为使这事更清楚，让我把两个人摆在你面前：一个人沉溺于奢侈享乐和此生的欺骗；另一个人对这一切已死去；让我们看看哪一个才是真正活着的人。假设其中一人非常富有且备受尊崇，豢养食客和谄媚者，又假设他整日沉溺于宴乐醉酒；而另一人过着贫穷、斋戒、艰苦饮食和严格纪律（\OriginalTerm{严格纪律}{\GreekText{φιλοσοφί}\& 139}）。\par

\SourceRecord{Translated by J. Walker, J. Sheppard and H. Browne, and revised by George B. Stevens.；Nicene and Post-Nicene Fathers, First Series；Vol. 11.；Edited by Philip Schaff.；Buffalo, NY: Christian Literature Publishing Co.,；1889.；<http://www.newadvent.org/fathers/210213.htm>.}

% structure: p0001 source=h1 target=section reason=page-title

\section{罗马书第十四篇讲道}

\BibleRef{罗 8:12-13}\par

\BibleQuote{因此，弟兄们，我们并不是欠肉体的债，去顺从肉体生活。如果你们顺从肉体生活，必要死亡；但如果你们藉着圣神，治死身体的恶行，必要生活。}\par

在展示了属灵生活的赏报何等伟大，它使基督住在我们内，使我们必死的身体复活，并振翅高飞，使德行之路变得轻省之后，他接着恰当地引入一项劝勉。\BibleQuote{因此}我们不该\BibleQuote{顺从肉体生活}。但这并不是他所说的，因为他以更惊人、更有力的方式表述，即：\Quote{我们是欠圣神的债}。因为说\Quote{我们并非欠肉体的债}，正表示这一点。这是他处处在证明的一点：天主为我们所做的一切，并非出于债务，而纯粹是恩宠。但在此之后，我们所做的就不再是自由奉献的事，而是债务了。因为当他写道：\BibleQuote{你们是用高价买来的，不要作人的奴隶}\BibleRef{格前 7:23}；当他写道：\BibleQuote{你们不是属于自己的}\BibleRef{格前 6:19}；在另一处他又提醒他们同样的事，说：若（多数手抄本省略\Quote{若}）\BibleQuote{一人替众人死了，那么众人就都死了，是为使他们不再为自己而活}\BibleRef{格后 5:15}。正是为了确立这一点，他在这里也说：\Quote{我们是债务人}；然后，既然他说我们\Quote{并非}欠\Quote{肉体的}债，免得你又以为他在反对肉体的本性，他没有停止说话，而是继续说：\BibleQuote{去顺从肉体生活}。因为我们确实欠肉体许多东西，例如给予它食物、温暖、休息、生病时的医药、衣物以及千百种其他照料。因此，为了防止你以为他所说的\Quote{我们并非欠肉体的债}是要废除这种照料，他便以\BibleQuote{去顺从肉体生活}加以解释。他的意思是：我所要废除的照料，是那导向罪恶的照料，而我却赞成肉体得到有益健康的照料。他在后面也表明了这一点。因为当他说：\BibleQuote{不要为肉体作准备}时，他并未停留于此，而是加上：\BibleQuote{去满足私欲}\BibleRef{罗 13:14}。他在这里也给了我们同样的教导，意思是：确实要关心它，因为我们欠它这个，但让我们不要顺从肉体生活，也就是说，不要让它成为我们生命的主宰。因为它必须追随，而非领导；不是它来调节我们的生活，而是它必须接受圣神的法律。明确了这一点，并证明了我们是欠圣神的债之后，为了进一步说明我们因哪些恩惠而成为债务人，他没有谈及过去的恩惠（这足以作为他判断力最有力的证明），而是谈及未来的恩惠；尽管过去的恩惠也足以达到目的。然而，在这种情况下，他并没有列出它们，甚至没有提及那些不可言喻的祝福，而是未来的事物。因为一次性的恩惠通常不如所期待的未来恩惠那样能吸引人。在加上这一点之后，他首先用顺从肉体生活所带来的痛苦和灾祸来吓唬他们，以下列话语：\BibleQuote{因为如果你们顺从肉体生活，必要死亡}，以此向我们暗示那不朽的死亡、惩罚和地狱中的报应。或者，如果有人要准确地审视这一点，这样的人即使在今世也是死的。这一点我们已经在上一篇讲道中向你们阐明。\BibleQuote{但如果你藉着圣神，治死身体的恶行，必要生活}。你看，我们所讨论的不是身体的本体，而是肉体的作为。因为他不是说\Quote{如果你们藉着圣神治死}身体的本体，而是说身体的\Quote{恶行}，而且不是所有的作为，而是那些邪恶的。这在接下来是显而易见的：因为如果你们这样做，他说：\BibleQuote{你们必要生活}。如果他的话语适用于所有作为，这在情理上怎么可能呢？因为看、听、说、走都是身体的作为；如果我们治死了这些，我们不但不能生活，反而要承担杀人的刑罚。那么他要我们治死的是哪些作为呢？是那些趋向邪恶的作为，那些追随恶习的作为，这些作为除了藉着圣神之外，别无他法可以治死。因为你可以通过自杀来结束其他作为，但你无权这样做。而对于这些（你只能）藉着圣神来结束。因为如果圣神临在，一切风浪都会平息，情欲都会俯伏，没有任何事物能高抬自己反对我们。所以你看，正如我之前所说，他是以未来的事物为基础来劝勉我们，并表明我们不仅是因已成就的事而成为债务人。因为圣神的恩惠不仅在于他使我们从过去的罪恶中获得自由，还在于他使我们对未来的罪恶坚不可摧，并使我们配得上不朽的生命。然后，为了提出另一个赏报，他继续说：\par

\BibleQuote{凡受天主圣神引导的，都是天主的子女。}\BibleRef{罗 8:14}\par

现在，这比第一个光荣又大得多。因此他不只说\Quote{凡随从天主圣神而生活的人}，却说\Quote{凡受天主圣神引导的}，为表明他愿意圣神在我们的生命中使用这样的权能，如同舵手对于船，或御者对于一对马。他不仅要将身体，也要将灵魂本身置于这类缰绳之下。因为他连灵魂也不愿它独立，而是要把它的权威也置于圣神的能力之下。免得因信赖圣洗的恩赐，他们之后对生活漠不关心，他说，即使你领受了圣洗，但若你后来没有心意要\Quote{随圣神的引导}，你就失去了赋予你的尊荣和义子的优越地位。因此他不说\Quote{凡领受了圣神的人}，却说\Quote{凡受圣神引导的}，就是那些终生这样生活的人，\Quote{他们是天主的儿女}。既然这尊荣也赐给了犹太人，因为经上说：\BibleQuote{我曾说你们是神，都是至高者的儿子}\BibleRef{咏 81:6}；又说：\BibleQuote{我养育了儿女}\BibleRef{依 1:2}；并且\BibleQuote{以色列是我的长子}\BibleRef{出 4:22}；保禄也说：\BibleQuote{他们拥有义子的名分}\BibleRef{罗 9:4}——接着他宣称后者与前者光荣之间的重大区别。因为尽管名称相同，他意指事物却不相同。关于这些点，他根据那些\OriginalTerm{行善者}{\GreekText{κατορθούντων}}的情形以及赐予他们的事物和将来的事物引入一个比较，给出了清楚的证明。首先他表明旧时赐予他们的是什么。这又是什么呢？\Quote{奴役的精神。}他于是这样继续：\par

第十五节。\BibleQuote{你们没有再受奴役的精神，以致恐惧。}\BibleRef{罗 8:15}\par

然后，他没有停留于提及与奴役相对立的那一端，即自由的精神，却提到了远为伟大的，即义子的名分，借此他同时也引入另一端，说道：\BibleQuote{但是你们领受了义子的圣神。}\par

但这是显而易见的。然而什么是束缚之灵，却不那么明显，需要使其更加清晰。他所说的话远非清楚，实际上令人非常困惑。因为犹太民族并没有领受圣神。那么他在这里的意思是什么呢？他把这名给了文字，因为文字是属神的，所以他也称法律为属神的，并称出自磐石的水和玛纳为属神的。\BibleQuote{他们吃了}，他说，\BibleQuote{同样的神粮，并且都饮了同样的神饮。}\BibleRef{格前10:3} 他也给磐石这名，当他说：\BibleQuote{他们饮了那随从他们的神磐石。} 现在，他之所以称它们为属神的，是因为当时所举行的一切礼仪都是超自然的，而并非因为领受它们的人领受了圣神。那些文字在什么意义上是束缚的文字呢？在你面前摆出整个救恩安排，然后你也会对此有清晰的看法。因为对他们来说，报应近在咫尺，赏报随即而来，立即相称，如同给仆人的每日口粮，众多的恐惧在他们眼前达于顶点，他们的洁净只涉及身体，他们的节制只延伸至行为。但对我们来说并非如此，因为就连想象和良心也被净化了。因为祂不是说，\BibleQuote{不可杀人，} 只是，甚至不可动怒；同样，不是\BibleQuote{不可奸淫，} 而是不可贪色正视。因此，我们的习惯性德行以及所有个别的善行，不应来自对现世惩罚的恐惧，而应来自对祂的渴望。祂也没有应许流奶流蜜之地，而是使我们与独生子同为继承人，以此千方百计地使我们远离现世之物，并应许赐予那些特别配得上成为天主儿子的人所接受的事物，也就是说，没有任何感官或物质的事物，而全是属神的。所以，即使他们有儿子的名号，也不过是奴隶；但我们既已获得自由，便领受了义子身份，并等候天堂。祂对他们是通过他人的中介讲话，对我们则是亲自讲话。他们所行的一切都是出于恐惧的冲动，而属神的行为则是出于渴望和强烈的欲望。这一点由他们曾跨越诫命的事实所证明。他们如同雇工和顽固之人，故从未停止抱怨；但这些人所做的一切都是为了悦乐圣父。同样，他们在蒙受恩惠时曾亵渎；但我们却在遭受危险时感恩。若因我们犯罪而需要惩罚我们，即使在这种情况下，差别也极大。因为我们并非像他们那样被司铎用石头砸、烙印、残害才被挽回。对我们来说，被逐出圣父的筵席、在若干天内不得露面就足够了。对犹太人来说，义子身份的荣誉只是名义上的，但在这里，现实也随之而来——圣洗圣事的洁净、圣神的赐予、其他恩惠的提供。此外还有几点，表明我们出身高贵、他们地位低下。在借着谈论圣神、恐惧和义子身份暗示了这一切之后，他再次提出拥有义子之神的新的证据。这是什么？就是\BibleQuote{我们呼喊：阿爸，父啊！} 这有多么伟大，已受启迪者知道\EditorialNote{圣济利禄·耶路撒冷，《教理讲授》23, §11, p. 276, O.T.}，他们有充分的理由奉命在入门者的祈祷中首先使用这个词。那么，有人会说，难道他们不也称天主为父吗？你难道没有听见梅瑟说：\BibleQuote{你抛弃了那生你的天主吗？}\BibleRef{申32:15} 你难道没有听见玛拉基亚责备他们，说：\BibleQuote{一位神造了你们，} 并且\BibleQuote{你们众人有一位父？}\BibleRef{拉2:10} 然而，如果这些话语和其他话语都被使用，我们找不到他们任何地方用这称呼呼号天主，或用这种言语祈祷。但我们所有人——司铎与平信徒、统治者和被统治者——都奉命以此祈祷。这就是我们在那些奇妙的阵痛和那奇异而不寻常的生育方式之后所发出的第一句话。如果在其他任何场合他们这样称呼祂，那只是出于他们自己的心意。但处于恩宠状态的人这样做，是因为被圣神的运作所推动。因为正如有一位智慧之神，借着它，无知的人变得有智慧，这在他们教导中显明；有一位德能之神，借着它，软弱的人复活死人、驱逐魔鬼；还有一位神医之神、一位先知之神、一位舌音之神，同样也有一位义子之神。正如我们认识先知之神，在于拥有它的人预言未来，不是凭自己说话，而是被恩宠所推动；同样，义子之神也是如此，拥有它的人被圣神推动，呼号天主为父。为了表达这最真实的血缘，他也用了希伯来语，因为他不仅说\BibleQuote{父啊，} 而且说\BibleQuote{阿爸，父啊！} 这名字是亲生儿女对父亲的一个特殊标记。在提及他们生活方式的差异、所赐恩宠的差异以及自由的差异之后，他提出了另一个与这义子身份相伴随的优越性的证明。现在这证明是什么呢？\par

第16节。\BibleQuote{圣神亲自和我们的心神一同作证：我们是天主的子女。}\BibleRef{罗8:16}\par

因为他所说的，不是仅仅从语言本身，而是从语言产生的原因来提出我的主张；因为我们这样说话是出于圣神的感动。在另一处，他用更明白的话表达了这一点，说：\Quote{天主派遣了自己圣子的圣神，到我们心中，呼喊说：\InnerQuote{阿爸，父啊！}}\BibleRef{迦 4:6} 那\Quote{圣神与我们的心神一同作证}是什么意思？他是指护慰者与赐给我们的恩赐一同作证。因为这不仅是恩赐的声音，也是赐予恩赐的护慰者的声音，因为祂自己通过恩赐教导我们这样说话。但当圣神\Quote{作证}时，还有什么怀疑的余地呢？因为如果是一个人、一位天使、一位总领天使，或任何其他这样的权能许下这事，那么怀疑是有理由的。但当至高的本体赐予这恩赐，并藉着祂吩咐我们在祈祷中使用的话语\Quote{作证}时，谁还会怀疑我们的尊位呢？因为即使是皇帝选拔某人，并在众人面前宣告对他的尊荣，也没有人敢反对。\par

17. \Quote{你们若是子女，便是承继者。} \BibleRef{罗 8:17}\par

请注意他是如何逐步提升这恩赐的。因为，子女虽是可能的，却不一定成为承继者（因为并非所有子女都是承继者），所以他另外加上了这一点——我们是承继者。但犹太人，除了没有我们这种义子名分外，还被剥夺了承继权。因为\Quote{他要凶恶地消灭那些恶人，把葡萄园另租给别的佃农}\BibleRef{玛 21:41}；在此之前，祂说：\Quote{有许多人要从东方和西方来，同亚巴郎一同坐席，但本国的子民要被驱逐出去。}\BibleRef{玛 8:11-12} 然而，他并没有停留于此，还定下了比这更大的事。这事是什么呢？就是我们成为天主的承继者；所以他加上\Quote{天主的承继者}。更有甚者，我们不单是承继者，还是\Quote{与基督一同为承继者}。请注意他是多么渴望使我们接近主。因为并非所有子女都是承继者，所以他表明我们既是子女也是承继者；其次，既然并非所有承继者都承继很大的产业，所以他表明我们也拥有这一点，因为我们是天主的承继者。再者，人可能是天主的承继者，却丝毫不是与独生子\Quote{一同为承继者}，所以他表明我们也有这一点。再请思考他的智慧。因为当他把令人不悦的部分压缩得十分简短，当他谈到那些\Quote{随从肉体生活}之人的结局时，例如他们将\Quote{死亡}，当他进入比较安慰人的部分时，他把言论引入广阔的天地，如此扩展关于赏报的叙述，并指出恩赐也是多样而伟大的。因为，如果做子女已是说不尽的恩宠，那么成为承继者该是多么伟大！如果这已是伟大，那么成为\Quote{一同为承继者}则更为伟大！然后，为了表明这恩赐不仅是恩宠，同时也给予他所说的话以可信度，他接着说：\Quote{只要我们与祂一同受苦，也必与祂一同受光荣。}他的意思是，如果我们在痛苦的事上与他一同有份，那么在好事上岂不更加如此？因为那将如此恩赐赐予没有行善之人的天主，当祂看见他们如此劳苦和受苦时，岂能不赐予更大的报答？既然他已表明这事是一种回报，为了使人相信所说的话，并防止任何人怀疑，他进一步表明它具有恩赐的效力。一方面他表明了这一点，使所说的话甚至在怀疑的人面前也获得信任，并使领受者不因常常白白领受救恩而感到羞愧；另一方面，是为了使你看到天主在报答上超过了劳苦。他在\Quote{只要我们与祂一同受苦，也必与祂一同受光荣}这句话中表明了一面；但另一面他在继续添加时表明：\par

18. \BibleQuote{现时的苦楚，与那将要显示于我们的光荣，是不能相比的（\OriginalTerm{在我们内}{\GreekText{εἰς}}）。} \BibleRef{罗 8:18}\par

在前文里，他要求属神的人修正自己的习惯（Mar.和六份手稿作\Quote{情欲}），他说：\Quote{你们并不是欠债的人，要按肉体去生活。}例如，这样的人应当超脱情欲、愤怒、金钱、虚荣、忌妒。但在这里，他提醒他们整个恩赐——既已赐予的，也将要赐予的——并把他高举到希望之中，安置在基督身边，表明他是独生子的共同继承者；现在他勇敢地引领他甚至面对危险。因为战胜我们内在的恶情欲，与承受那些试炼、鞭打、饥荒、掠夺、枷锁、镣铐、处决并不相同。因为后者需要更崇高的精神。请注意他如何同时安抚和激励战斗者的精神。因为他已表明赏报大于劳苦之后，既劝勉更大的努力，又不让他们骄傲，仿佛仍被作为报偿的冠冕所超越。在另一处他又说：\BibleQuote{我们这暂时轻微的苦难，要为我们成就极重无比、永远的荣耀}（\BibleRef{格后 4:17}）。那是他当时对更深层次的人说话。然而在这里，他并不承认那些苦难是轻微的；反而用将来美好事物带来的补偿来安慰他们，说：\Quote{我实在认为现今的苦楚，与将来的荣耀相比，不值得一提。}他没有说与将来的休息（\OriginalTerm{休息}{\GreekText{ἄνεσιν}}）相比，而是说更重大的：\Quote{与将要来临的荣耀相比。}因为休息之处并不一定有荣耀；但有荣耀之处必定有休息。然后，既然他说荣耀是将来要来的，他却表明它现在已经存在。因为他不是说\Quote{那将要存在的}，而是说\Quote{那要在我们身上被启示出来的}，好像已经存在但尚未启示。正如在另一处他用更清楚的话说：\BibleQuote{我们的生命与基督一同藏在天主内。}所以你们要对此放心。因为荣耀已经预备好了，等待着你们的劳苦。但如果它尚未到来使你烦恼，那么这件事本身倒应使你欢喜。因为正是由于它伟大、不可言喻、超越我们现今的状况，所以被保管在那里。因此他不只说\Quote{现今的苦楚}，而是这样说话，以表明不仅性质上，而且数量上，来生也占优势。因为这些苦难无论是什么，都是与现世生命相连的；但将来的祝福却延伸到无尽的世代。既然他无法详细描述这些祝福，也无法用语言向我们呈现，他就用一个似乎是我们特别渴望的对象来命名它们：\Quote{荣耀}。因为祝福的顶峰和总和似乎就是荣耀。为了以另一种方式激励听者，他通过提及受造界使他的言谈崇高，从而在接下来的话语中达到两个目的：轻视现世之物，渴望来世之物；此外还有第三个目的，或者更确切地说第一个目的是表明天主如何眷顾人类，以及祂如何尊崇我们的本性。除此之外，哲学家们为自己构想的所有关于这个世界的学说，如同蛛网或孩童的沙堆，都被这一教义推翻了。但为更清楚了解这些事，让我们听宗徒自己所说的话。\par

第19、20节。\BibleQuote{受造界热切地期待着}，他说，\BibleQuote{天主的子女的启示。因为受造界被屈服于虚妄之下，不是出于自愿，而是由于那使之屈服的那位，在希望中屈服了。}\par

意义大致如此：受造物本身正在经历产痛，期待并等待着我们刚才所说的这些美善。因为\Quote{热切期待}（\OriginalTerm{热切期待}{\GreekText{ἀποκαραδοκία}}，即翘首以盼）意指强烈的期待。如此，他的论述便更具强调，并且他像先知们那样将整个世物人格化——先知们曾使洪流鼓掌、丘陵跳跃、山岳欢腾，这并非要我们认为它们有生命，或赋予它们任何理性，而是为了让我们明白祝福的伟大，以至于连无感觉之物也被触及。同样，他们在痛苦之事上也常如此，使葡萄树哀叹、酒也哀叹、山岳和圣殿的木板哀号，这同样是为了让我们理解灾祸的极致。宗徒正是仿效这些，使受造物在此成为一个活生生的人，说它呻吟并忍受产痛：并非他听到了从地和天传来的呻吟，而是为了显示未来美善的无比伟大，以及摆脱如今遍及它们的灾祸的渴望。\Quote{因为受造物被屈服于虚妄，不是出于自愿，而是由于那使它屈服的那一位。}\Quote{受造物被屈服于虚妄}是什么意思？就是说它变成了可朽坏的。原因何在？为了什么？为了你，人啊。因为你领受了可朽坏、易受苦的身体，大地也受了诅咒，长出荆棘和蒺藜。至于天，当它同地一起衰老之后，将改变为更好的分位（\OriginalTerm{分位}{\GreekText{λῆξιν}}），请听先知的话：\Quote{主，你从起初奠定了大地，诸天是你手的化工；它们将要毁灭，你却永存；它们都要像衣服一样破旧，你要像卷起斗篷一样卷起它们，它们就被更换。}（\BibleRef{咏 102:25}）依撒意亚也宣告了同样的事，他说：\Quote{你们要仰观上天，俯视大地；因为诸天如烟雾消散，大地要像衣服一样破旧，地上的居民也要如此死亡。}（\BibleRef{依 51:6}）现在你明白受造物\Quote{屈服于虚妄}是什么意思，以及它将如何从败坏中得到释放。因为那一位说：\Quote{你要像卷起斗篷一样卷起它们，它们就被更换}；依撒意亚说：\Quote{地上的居民也要如此死亡}，当然不是指完全的灭绝。因为地上的居民，即人类，也不会经历那样的灭绝，而是一种暂时的，并藉此转变为不朽的状态（\BibleRef{格前 15:53}），因此受造物也将如此。这一切他都顺便指出了，藉他的说法\Quote{如此}（\BibleRef{伯后 3:13}），保禄在稍后也这样说。但现在他谈论的是奴役本身，并说明它为何成为这样，指出我们自己就是原因。那么，它因别人的缘故而受苦待了吗？绝不是，因为它是为我而被造的。那么，它为我而被造，却因我的纠正而遭受这些，这有什么不义呢？或者，在无灵魂、无感觉的事物上，根本无需提出正义与否的问题。但保禄，既然已将它人格化，便没有使用我提到的任何这些论题，而是用了另一种说法，似乎想以最大的益处安慰听者。这是怎样的呢？你意如何？它的意思是：它为了你的缘故而受虐待，变成了可朽坏；然而它并没有受不义。因为为了你的缘故，它又将变成不朽的。这就是\Quote{在希望中}的意思。但当他说它\Quote{不是出于自愿}而被屈服时，并非要表明它拥有判断力，而是为了让你明白，这一切由基督的眷顾所成就，并非它自己的功绩。那么，在什么希望中呢？\par

21节：\Quote{受造物自己也要从败坏奴役中获得释放。}（\BibleRef{罗 8:21}）\par

这受造物是什么？不只是你，还有那比你低等、不具理性或感觉的，也要分享你的祝福。因为他说：\Quote{它要从败坏奴役中获得释放}，即它不再可朽坏，而要随你身体所得的美丽而行；正如当此身体变成可朽坏时，它也可朽坏一样；同样，现在它变成不朽的，它也要随之。为了表明这一点，他继续说（\OriginalTerm{进入}{\GreekText{εἰς}}）：\Quote{进入天主儿女的光荣自由。}即因他们的自由。如同一位抚养王子孩子的乳母，当孩子到达父亲的权位时，她自己也与孩子一同享受美善；同样，受造物也是如此，他意指。你看，人在各方面都占首位，一切都是为了他而造。看他如何安慰那挣扎者，并显示天主对人不可言喻的爱。因为他会说，为什么你为考验而烦恼？你是为你自己受苦，受造物是为你。他不仅安慰，还显示他的话是可信的。因为如果那完全为你而造的受造物\Quote{在希望中}，那么你——受造物将藉你而享受那些美善——更应当如此。正如人们，当儿子要在他获得尊荣时出现，也会给仆人穿上更光亮的衣服，以荣耀儿子；同样，天主也要为了儿女的光荣自由，给受造物穿上不朽的衣服。\par

22节：\Quote{我们知道，直到现在，整个受造物都一同呻吟，一同忍受产痛。}（\BibleRef{罗 8:22}）\par

注意他如何使听者羞愧，几乎在说：不要比受造物更差，也不要在安于现状中寻乐。我们不仅不应执着于它们，甚至要为延迟离开此世而叹息。因为如果受造物尚且如此，你这蒙受理性尊荣的人更当如此。但既然这还不足以强求他们的注意，他便继续说道。\par

第 23 节. \BibleQuote{不但如此，连我们这些有圣神初果的，也是自己心中叹息。} \BibleRef{罗 8:23}\par

那就是说，尝到了将来之事的滋味。因为即使有人心硬如石，他所意指的意思是：已经赐下的已足以提升他，使他脱离现状，并使他在两件事上向往将来之事：既因所赐之事的伟大，也因它们虽伟大且众多，却只是初果。因为如果初果如此伟大，以至我们藉此得以脱离罪恶，获得正义与圣化，而且那时的人甚至藉他们的影子——\BibleRef{宗 5:15}——或衣物——\BibleRef{宗 19:12}——驱魔并复活死者，那么想一想整体该有多么伟大。如果受造物虽没有心智和理性，且对此毫无所知，尚且叹息，我们岂不更当如此？其次，为了不给异端者可乘之机，或显得在贬低我们现今的世界，他说我们叹息，不是因为指责现世体制，而是出于对更大之事的渴望。他在以下话语中表明了这一点：\BibleQuote{等候儿子的名分。} \BibleRef{罗 8:23} 你说什么？我且听一听。你一再坚持并大声呼喊，说我们早已成了儿子，如今你怎把这件好事放在希望中，写说我们必须等待它？为了纠正这一点，他接着说道：\BibleQuote{就是我们身体的救赎。} \BibleRef{罗 8:23} 那就是完美的荣耀。事实上，我们的命运直到最后一息都处于不确定中，因为我们中许多曾是儿子的人变成了犬和囚徒。但如果我们怀着善终的希望离世，那么恩赐就不可动摇、更清晰、更伟大，不再惧怕死亡和罪恶带来的任何改变。因此，当我们的身体从死亡及其无数疾苦中得释放时，恩宠才会稳固。因为这就是完全的赎价（\OriginalTerm{赎价}{\GreekText{ἀπολύτρωσις}}），不仅是赎价，而且是我们将永不再返回先前的囚禁。为了让你在听到这么多荣耀却不了解其具体内容时不致困惑，他部分地揭示了将来之事，将你身体的改变（即改变你的身体）以及整个受造物的改变放在你面前。他在另一处更清楚地说明了这一点，说道：\BibleQuote{祂要改变我们卑贱的身体，使之相似祂光荣的身体。} \BibleRef{斐 3:21} 他在另一处又写道：\BibleQuote{几时这可朽坏的穿上了不朽坏的，那时就要应验经上所记载的话：死亡被胜利吞灭了。} \BibleRef{格前 15:54} 但为表明随着身体的朽坏，今世事物的结构也要完结，他在别处写道：\BibleQuote{因为这世界的局面正在逝去。} \BibleRef{格前 7:31}\par

第 24 节. \BibleQuote{因为我们得救是在于望德，} 他说。\BibleRef{罗 8:24}\par

既然他详述了将来之事的应许，而这似乎使较软弱的听者感到痛苦——如果这些福分都是希望之事——那么，在证明它们比现在可见之事更确定、并详尽论述已赐的恩赐、显示我们已领受了那些美物的初果之后，以免我们在今世寻求一切，而背叛信仰所赐予的高贵，他便说：\BibleQuote{因为我们得救是在于望德（希腊文是过去式）。} \BibleRef{罗 8:24} 这大概是他的意思。我们不应在今生寻求一切，而应拥有望德。因为这是我们在天主面前所献的唯一礼物，相信祂所应许的将要来临，并且唯独藉此方式我们得救。如果我们失去这望德，就失去了我们自己贡献的一切。因为我要问你，他会说：你岂不是身负无数罪恶？岂不是陷于绝望？岂不是被定罪？难道不是所有人都对你的得救灰心？那么，是什么救了你？是你唯独盼望天主，信赖祂有关应许和恩赐的话，此外别无他物。如果正是这望德救了你，如今也应紧握它。因为那曾带给你如此大恩的，必不会在将来的事上欺骗你。既然它找到了你——已死、败坏、囚徒、仇敌——却使你成为朋友、儿子、自由人、义人、同承继者，且赐下无人曾期望的如此大事，那么，在如此慷慨和眷顾之后，它怎会在后续的事上背叛你？不要对我说：又是希望！又是期待！又是信德！因为从一开始你正是藉此得救，而这嫁妆是你带给新郎的唯一礼物。如今你要紧握它、持守它。因为如果你要求在今生拥有一切，你就失去了使你成为光明的善行。这正是他接着说的原因：\BibleQuote{但可见的望德不是望德；因为人既然看见，为什么还希望呢？} \BibleRef{罗 8:24}\par

第 25 节. \BibleQuote{如果我们希望那看不见的，就要耐心等待。} \BibleRef{罗 8:25}\par

这就是说，倘若你期望在这世上获得一切，那么望德有何必要呢？什么是望德？就是对于未来之事的信心。那么，天主向你们要求什么大事呢？祂自己从自己的宝库中完全赐予你们恩赐，只向你们要求一件，就是望德，使你们也能为你们的救恩贡献一些属于你们自己的东西。这一点祂在接下来的话中暗示说：\BibleQuote{如果我们盼望那未见之物，就会耐心等候它。}因此，正如天主加冕那忍受劳苦、艰难和无数辛劳的人，祂也加冕那心怀望德的人。因为忍耐之名属于辛苦和很大的坚韧。然而，祂甚至将这忍耐赐给了那心怀望德的人，以便安慰疲惫的灵魂。然后，为了表明我们因这轻微的任务而享受丰富的助佑，他接着说：\par

第26节。\BibleQuote{同样，圣神也扶助我们的软弱。}\BibleRef{罗 8:26}\par

因为一方面属于你自己，即忍耐；另一方面来自圣神的恩赐，祂也\OriginalTerm{傅油}{anoints}你，使你怀有这望德，并借此再次减轻你的劳苦。为使你知道，这恩宠不仅在你的劳苦和危险中扶持你，即使在那些看似最容易的事情上，它也与你同工，并在一切情况下尽联盟之责，他接着说：\par

\BibleQuote{因为我们不知道该如何祈求才合适。}\par

祂这样说，是为显示圣神对我们极大的关怀，并教导他们不要以为那些在人看来似乎好的事就是确定可取的。因为很可能，当他们遭受鞭打、被驱逐、忍受无数苦难时，会寻求喘息，向天主祈求这恩惠，并认为这对他们有利，但（他说）绝不要以为那些看似祝福的事真的就是祝福。因为连这一点我们也需要圣神的助佑。人是如此软弱，靠自己什么也不是。为此祂说：\BibleQuote{因为我们不知道该如何祈求才合适。}为了使学习者不为自己的无知感到羞愧，祂没有说\Quote{你们不知道}，而是说\Quote{我们不知道。}祂这样做并非仅仅为了显得谦和，祂从其他段落清楚表明了这一点。因为祂在祈祷中不断渴望见到罗马。但祂获得此事的时间并非立即在祂渴望之时。至于那加给祂的\BibleQuote{刺在肉身内}\BibleRef{格后 12:8}，即危险，祂曾多次恳求天主，却完全未获应允。同样，旧约中的\Person{梅瑟}祈求看到巴勒斯坦\BibleRef{申 3:26}，\Person{耶肋米亚}为犹太人恳求\BibleRef{耶 15:1}，\Person{亚巴辣罕}为索多玛的人民代祷。\BibleQuote{但是圣神亲自以无可言喻的叹息为我们转求。}这句话不清楚，因为当时通常发生的许多奇迹如今已停止。因此，我必须向你们说明当时的情况，这样其余的问题就会清晰。那么当时的情况如何呢？天主在那些日子赐给所有受洗的人某些卓越的恩赐，这些恩赐的名称就是神恩。因为\BibleQuote{先知的神恩是顺服先知的}\BibleRef{格前 14:32}。有人得了预言之恩，预言未来的事；有人得了智慧之恩，教导众人；有人得了治病的恩典，治愈病人；有人得了奇迹之恩，复活死人；有人得了语言之恩，说不同的语言。此外还有祈祷之恩，也称为神恩，拥有这恩赐的人为全体百姓祈祷。因为我们对许多有益的事无知，并祈求无益的事，所以那祈祷之恩临到当时某个人身上，凡对全体教会有益的事，他就被指派代表众人祈求，并教导其他人。因此，此处所说的神恩乃是指这种恩宠的特性，以及那领受恩宠、向天主叹息转求的灵魂。因为那被认为配得上这种恩宠的人，怀着极大的痛悔和许多内心的叹息，俯伏在天主面前，祈求对众人有益的事。今日的执事在代表百姓献上祈祷时，正是这事的预像。这就是保禄说\BibleQuote{圣神亲自以无可言喻的叹息为我们转求}的意思。\par

第27节。\BibleQuote{但那洞察人心的。}\BibleRef{罗 8:27}\par

你们看，祂所说的并非关于护慰者，而是关于属神的心。因为若非如此，祂本应说\Quote{那洞察}圣神。但为使你们知道这话是指一个属神的人，他拥有祈祷之恩，他接着说：\BibleQuote{那洞察人心的知道圣神的心意，}也就是说，知道属神的人的心意。\par

\BibleQuote{因为祂按天主的旨意代圣人转求。}\par

并非（他意思说）祂好像无知似的告知天主，而是这样做是为了使我们学习祈求合适的事，并向天主求祂所喜悦的事。因为这就是所谓\Quote{按天主的旨意}。这样做是为了安慰那些来到祂面前的人，并给他们卓越的教导。因为那供应恩赐，并赐下无数恩福的，是护慰者。因此经上说：\BibleQuote{这一切事都是同一且相同的圣神所运行}\BibleRef{格前 12:11}。这事的出现是为教导我们，并显示圣神的爱，祂甚至屈尊就卑。正因如此，祈祷的人蒙垂听，因为祈祷是\Quote{按天主的旨意}而作的。\par

你看到他从多少方面教导他们，向他们显示的爱和给予的尊荣。天主为我们做了什么呢？祂为我们创造了可朽坏的世界，又为我们创造了不朽坏的。祂容许我们的先知受虐待，为我们使他们被掳，让他们落入火窑，忍受无数苦难。不仅如此，祂还为我们立了先知，也为我们立了宗徒。祂为我们舍弃了祂的独生子，祂为我们惩罚了魔鬼，祂使我们坐在祂右边，祂为我们受了辱骂。\BibleQuote{辱骂你者的辱骂}，经上说，\BibleQuote{落在我身上}（\BibleRef{咏 68:9}）。然而，在如此大恩之后，当我们退缩时，祂仍不放弃我们，反而再次恳求，并为我们激动别人为我们恳求，好施恩于我们。梅瑟也是如此。因为祂对他说：\BibleQuote{你且由我，让我消灭他们}（\BibleRef{出 32:10}），为要促使他为百姓祈求。如今祂也这样做。因此祂赐予了祈祷的恩赐。但祂这样做，并非祂自己需要恳求，而是免得我们因\OriginalTerm{简单地}{\GreekText{ἁ} \GreekText{πλὥς}}得救而变得冷漠。为此，祂常因达味和这个或那个人说祂与他们和好，以再次确立这一点，使和好合乎一切礼仪。然而，如果祂不是通过这位或那位先知，而是亲自告诉他们祂息怒了，那会显得更爱人。但祂不坚持那一点的缘由是，免得这和好的理由成为懒惰的借口。因此，祂也对耶肋米亚说：\BibleQuote{你不要为这人民祈祷，因为我不听你}（\BibleRef{耶 11:14}），并非要阻止他的祈祷（祂切望我们的得救），而是要恐吓他们；先知也看到这一点，并没有停止祈祷。为使你们看到祂这样说不是要使他放弃，而是要羞辱他们，请听经文：\BibleQuote{你没有看见他们所做的事吗？}（\BibleRef{则 8:6}）而当祂对城市说：\BibleQuote{虽然你用碱洗涤，多用肥皂（希腊文：草），但我面前你仍有污渍}（\BibleRef{耶 2:22}），祂这样说并非要使他们陷入绝望，而是要激发他们悔改。如同对尼尼微人，祂给出无条件的判决，不给予任何好的希望，反而更吓唬他们，引导他们悔改；祂在此也如此，既为激发他们，也为使先知更受尊敬，好使他们至少这样听他。然后，由于他们继续处于不可救药的疯狂中，即使其他人被掳走也不能使他们清醒，祂先劝他们留在那里。但既然他们没有遵守，反而逃到埃及，祂确实允许了他们，但要求他们不要同时背叛信仰和逃到埃及（\BibleRef{耶 44:8}）。但他们在这一点上也没有顺从，祂便派先知与他们同去，免得他们最终彻底沉沦（\BibleRef{耶 44:28}）。因为当他们不听从祂的召唤时，祂随后就跟随他们以管教他们，阻止他们更深地陷入恶习，如同一位充满慈爱的父亲对待一个无论怎样管教都同样任性的孩子，带着他到处走，跟随着他。这就是为什么祂不仅派耶肋米亚去埃及，也派厄则克耳去巴比伦，他们也没有拒绝去。因为当他们发现他们的主人极其爱百姓时，他们自己也继续如此行。就像一位正直的仆人，看到父亲为不听话的儿子悲伤哀叹，便也怜悯那儿子。他们为他们受了什么苦呢？他们被锯死，被驱逐，被辱骂，被石头砸死，忍受了无数苦难。之后他们又跑回他们那里。例如，撒慕尔没有停止为撒乌耳哀悼，尽管撒乌耳恶劣地侮辱他，造成了不可弥补的伤害（\BibleRef{撒上 15:35}）。但他仍然不记念这些事。耶肋米亚为犹太人民写了哀歌。当波斯将军给他自由，让他安全地、完全自由地住在任何他喜欢的地方时，他宁愿选择人民的苦难和他们在异地的艰难处境，也不愿住在家里（\BibleRef{耶 11:5}）。所以梅瑟离开了宫殿和其中的生活，急忙置身于他们的灾难中。达尼尔连续二十天不吃东西，以最严格的斋戒克苦自己，为要使他们与天主和好（\BibleRef{达 10:2}）。那三个青年在火窑中和如此猛烈的火焰中，也为他们献上祈祷。因为他们并非为自己悲伤，因他们已得救；但他们认为那时是最大的胆量说话的时候，因此他们为他们祈祷；因此他们也说：\BibleQuote{但愿我们以痛悔的心和谦逊的精神蒙悦纳}。\EditorialNote{三圣童之歌第16节}若苏厄也为此撕裂衣服（\BibleRef{苏 7:6}）。厄则克耳也为此哀号，当他看到他们被砍倒时（\BibleRef{则 9:8}）。耶肋米亚说：\BibleQuote{让我独自一人痛哭}（\BibleRef{依 22:4}）。在这之前，他不敢公开祈祷免除他们悲惨的处境，就寻求某个有限的时期，他说：\BibleQuote{上主啊，要到几时？}（\BibleRef{依 6:11}）因为一切圣徒的种族都充满慈爱。因此圣保禄说：\BibleQuote{为此，你们既是天主的选民，圣洁蒙爱的人，就要穿上怜悯的心肠、仁慈、谦卑}（\BibleRef{哥 3:12}）。你们看到这个词的严格恰当，以及他如何要我们常怀怜悯。因为他不仅说\Quote{显示怜悯}，而是说穿上，好如同我们的衣服常与我们同在，怜悯也如此。他不仅说怜悯，而是说\Quote{怜悯的心肠}，好使我们仿效亲属间的自然情感。\par

但我们所做的恰恰相反：如果有人来向我们乞讨一文钱，我们就辱骂他、呵斥他、称他为骗子。人哪，你不发抖、不脸红吗？为了块面包就称他为骗子？即使他是在行骗，他也值得怜悯，因为他被饥饿逼迫到不得不扮演这样的角色。这恰恰是对我们残酷的谴责。因为我们没有慷慨施舍的心肠，所以他们被迫使用许多伎俩，用欺骗来对付我们的冷酷，缓和我们的苛刻。如果他向你索取的是金银，那么你的怀疑还有些道理。但他是来向你要必需的食物，为什么你却不合时宜地显出自己的聪明，对他如此苛求，指责他懒惰和懈怠？如果我们必须这样说话，那么我们该责怪的也不是别人，而是我们自己。因此，当你到天主面前求赦免你的罪时，请记住这些话，你就会知道，你有更充分的理由让天主对你说这些话，而非那个穷人对你说的。然而天主从未对你说过这样的话：\Quote{退开吧，因为你是个骗子，总是来圣堂听我的法律，但在外面却看重黄金、贪欲（\OriginalTerm{贪欲}{\GreekText{ἐπιθυμίαν}}）、友谊，事实上任何东西都重于我的诫命。现在你装作谦卑，但祈祷一结束，你就变得傲慢、残酷、不仁。所以，你离开吧，再也不要到我这里来。}然而，我们配得上被说这样的话，甚至更严厉的话；但祂从未这样责备我们，反而长久忍耐，尽祂的本分完成一切，赐给我们超过我们所求的。既然如此，让我们记念这事，周济那些向我们乞讨之人的贫穷；即使他们欺骗我们，也不要过于苛求。因为我们自己所需要的就是这样的救恩：包含着宽恕、慈爱（\OriginalTerm{慈爱}{\GreekText{φιλανθρωπίας}}）和大量的仁慈。因为这是不可能的——绝对不可能——如果我们的状况被严格审查，我们永远不可能得救；反而必须受罚，完全被毁灭。因此，我们不要作苛刻的审判官，以免我们也受到严格的清算。因为我们的罪太大，无法找到任何借口。因此，让我们对那些犯了不可饶恕之罪的人多施怜悯，以便我们也能预先为自己积攒同样的怜悯。然而，无论我们如何宽宏大量，也永远无法付出如我们所需要的那份来自爱人的天主的爱。那么，当我们自己需要许多东西的时候，却对同僚苛求，竭尽全力反对自己，这岂非荒谬绝伦？因为你这样做，与其说是证明他不配你的慷慨，不如说是证明你自己不配天主对人的爱。因为那对同僚苛求的人，在天主面前必定会得到同样的对待。我们不要自相矛盾；即使他们是出于懒惰或任性而来，我们也要施舍。因为我们也犯了许多任性的罪，更确切地说，我们所有的罪都是任性的；但天主并未立即召我们受罚，反而给了我们一段补赎的时间，天天培育我们、管教我们、教导我们、供应我们一切所需，好让我们也能效法祂的仁慈。因此，让我们平息这种残酷，摒弃这种野蛮的精神，因为这样做受益的是我们自己，而非他人。因为我们给他们金钱、面包和衣服，却为自己预先积累了极大的荣耀，这荣耀是言语无法描述的。因为我们将重获不朽的身体，与基督一同受荣耀、一同为王。这事何等浩大，我们由此可以看见——确切地说，现在无法使我们清楚看见。但为了从我们目前的福分出发，至少从中获得一点模糊的预兆，我将尽力向你们说明我所谈论的。那么请告诉我：假如你年事已高，生活贫困，有人突然许诺使你返老还童、回到青春鼎盛，使你极为强壮、极其美丽，并赐给你全地球的王国一千年，一个处于最深和平的国度，那么有什么你不愿意做、不愿意忍受来取得这个应许呢？（四份手稿和Sav. Mar.反对。）请看，基督应许的不是这个，而是比这更大得多。因为老年与青年的差距无法与朽坏与不朽坏的差别相比，王国与贫穷的差距也无法与现在的荣耀和将来的荣耀相比；其差别如同梦境与现实。更确切地说，我尚未说到点子上，因为没有言语能够向你展示未来之事与现在之事之间的巨大差异。至于时间，根本没有可比较的余地。因为人用什么方式能将无限的生命与我们现在的情形相比较呢？至于和平，它远离任何现在的和平，如同和平不同于战争；至于不朽坏，它更美好，如同明亮的珍珠优于泥块。或者，无论你说多么伟大的事物，也无法将其呈现在你面前。因为即使我将那时我们身体的美丽比作阳光或最明亮的闪电，我也没有说出任何与那光辉相称的话。那么，为了这些事物，有什么金钱是不值得放弃的呢？有什么身体——更确切地说，有什么灵魂是不值得为了它而舍弃的呢？现在，如果有人带你进入王宫，当着众人的面给你机会与国王交谈，让你坐在他的桌前，共享他的膳食，你会称自己为最幸福的人。但当你将要升到天堂，站立在宇宙的君王本身面前，与天使争辉，享受那不可言喻的荣耀时，你却在犹豫是否该放弃金钱？其实，即使你必须舍弃生命本身，你也应该跳跃、欢欣，乘着喜乐的翅膀飞翔。但你们为了得到一个官职（\OriginalTerm{官职}{\GreekText{ἀρχήν}}），作为搜刮之地（因为我不能称这类东西为获利），就冒一切风险，甚至向别人借钱，如果需要，还会毫不犹豫地典当妻子和儿女。\par

但当天国摆在你们面前，那无人能取代的职位（\OriginalTerm{职位}{\GreekText{ἀρχῆς}}），且天主命你们收取的不是大地的一角，而是整个天，你们还在犹豫、抗拒、贪恋钱财，忘记我们所见的这部分天既如此美丽悦目，那更高的天和天上的天该是何等美妙？但因我们尚无法用肉眼看见，就请用思想上升，站在这天之上，仰望那天外之天，那无边的崇高，那令人敬畏的光明，那无数天使的队列，那无尽的总领天使的行列，以及其余无形体的大能者。然后再次把握我们从高处下来后所拥有的那幅图像（参柏拉图 \WorkTitle{理想国} 卷七，第516页），并为我们描绘一位地上君王的景象：他的士兵身披金甲，他的白骡成对，骄傲地饰以黄金，他的战车镶嵌宝石，他的坐垫洁白如雪（\OriginalTerm{坐垫}{\GreekText{στρωμνήν}}），战车周围飘动的饰片，丝织帷幔上织出的龙形，盾牌上凸起的金钉，从金钉到盾缘的饰带镶嵌众多宝石，马匹披着镀金的鞍具和黄金嚼子。但当我们看到君王时，这一切立刻被忽略：因为只有他吸引了我们的目光，他的紫袍、冠冕、宝座、扣环、鞋履，以及他全部辉煌的仪容。准确地将这一切聚集起来后，再把你的思绪从这些转移到上面的事，和那可怕的日子，就是基督来临的日子。那时你看到的不是成对的骡子，不是金战车，不是龙和盾牌，而是充满极大敬畏的事物，令无形体的大能者都惊奇。因为祂说：\InnerQuote{诸天的权势}将要\InnerQuote{震动。}（\BibleRef{玛 24:29}）那时整个天敞开，那些穹苍的门户自行打开，天主的独生子降下，不是带着二十或一百个护卫，而是成千上万的天使和总领天使、革鲁宾和色辣芬，以及其他大能者，一切事物都充满恐惧和战栗，大地自行裂开，从亚当出生直到那日所有存在过的人都从大地复活，全被提起（\BibleRef{得前 4:17}），祂亲自显现，带着如此伟大的荣耀，以致太阳、月亮和一切光都黯然失色，被那光辉所掩盖。什么言语能向我们描述那福乐、光明、荣耀？唉！我的灵魂。因为哀哭和极大的叹息临到我，当我思虑我们失落了何等美善，我们疏远了何等福乐。因为我们疏远了——我现在仍就自己的情况说话——除非我们做某些伟大的、令人惊奇的事。现在不要再对我提地狱，因为失去这荣耀比任何地狱更痛苦，疏远那份福乐比无数惩罚更糟。然而我们仍贪恋今世，不理会魔鬼的诡计，他用小事夺走我们的大事，给我们泥土以便夺走金子，或者更确切地说，夺走天堂，给我们影子以便使我们失去实体，在梦中把幻象放在我们面前——今世的财富正是如此——以便天一亮就显明我们是最贫穷的人。将这些牢记在心，尽管为时已晚，让我们逃离这诡计，转向未来的事物。因为我们不能说不知道今世是如此易变：每天的事物比喇叭更响亮地传入我们耳中：今世场景的空虚、可笑、可耻、危险和陷阱。那么，我们有什么辩护之词，竟如此热切地追求这些充满危险和耻辱的事，而忽略那永恒、使我们光荣辉煌的事，并将我们全部身心献给钱财的奴役？因为对这些事的奴役比任何束缚更糟。那些有幸获得自由的人知道这一点。为使你们也能感受到这美好的自由，请挣断绳索，跳出陷阱。不要让金子藏在你们家中，而要让比千万钱财更宝贵的施舍和爱人之心作为你们的宝藏。因为后者使我们向天主放胆无惧，前者却使我们深深蒙羞，并使魔鬼对我们猛攻（\OriginalTerm{猛攻}{\GreekText{σφοδρὸν} \GreekText{πνεῖν}}）。为何要武装你的敌人，使他更强？用你的右手武装起来对抗他，将你家中的一切辉煌转移到你的灵魂里，将你的全部财富储存在你的心智中，让天堂而不是箱子和房屋保存你的金子。让我们把全部财产围绕自己；因为我们比墙壁好得多，比铺地石更尊贵。为何忽略自己，将全部注意力浪费在那些我们离开后无法再触及、甚至常在此世也无法把握的事物上，而我们本可以拥有这样的财富，使我们在今世和来世都处于最安逸的境地？因为那把自己的田产、房屋和金子携带在灵魂上的人，无论出现在哪里，都带着这笔财富出现。有人问：这如何可能？可能，而且极其容易。如果你通过穷人的手将它们转移到天堂，你就完整地将它们转移到你自己的灵魂里。如果后来死亡降临，没有人能从你手中夺走它们，你将在来世也作为富人离去。塔彼达（\Person{塔彼达}）的宝藏就是这种。因此，宣告她富有的不是她的房屋、墙壁、石头或柱子，而是寡妇们穿着衣服的身体、她们的眼泪、逃亡的死亡和复归的生命。让我们也为自己制造这样的宝藏，为自己建造这样的房屋。这样，我们将有天主作为我们的同工，我们自己也将与祂同工。因为祂使贫穷者从无有进入存在，而你们则在她们进入生命和存在之后，通过照顾和扶助她们，在各地支撑天主的圣殿，防止她们因饥饿和其他痛苦而灭亡。就效用和荣耀而言，有什么能与这相比？如果你们尚未清楚明白祂托付你们救济贫穷时所赐予你们的巨大装饰，请思考这一点：如果祂赐予你们如此大的权能，使你们能够扶起正在坠落的天空，难道你们不认为这荣耀过于你们应得的吗？请看，祂现在使你们配得上更大的荣耀。因为在天主眼中，比诸天更宝贵的事物，祂托付你们去修复。在一切可见之物中，没有任何东西在天主眼中与人同等。因为祂为人创造了天、地、海，并更喜悦与人同住，胜过住在天上。然而我们虽知道这一点，却不关心也不顾及天主的圣殿；任由它们被忽视，却为自己建造华丽宽大的房屋。这就是我们缺乏一切美善的原因，比最贫穷的人更贫穷，因为我们为这些不能带走的房屋而骄傲，却忽略了那些我们可以随身移动的房屋。因为穷人死后身体必将复活；天主，这位托付此事的，将使他们显现，并赞美那些照顾他们的人，重视（\OriginalTerm{重视}{\GreekText{θαυμάσεται}}）他们，因为当穷人因饥饿、赤身和寒冷即将倒塌时，这些人尽力修复了他们。然而，即使这些赞美摆在我们面前，我们仍拖延，拒绝承担这尊贵的职责。基督竟无栖身之处，到处流浪、赤身、饥饿，而你却为自己建造郊外的房屋、浴池、平台和无数的房间，出于虚荣；而基督你连一间小屋的份额都不给，却为乌鸦和秃鹫装饰楼阁。有什么比这更疯狂？有什么比这更严重？这是极度的疯狂，甚至无法命名，无论用什么词。然而，只要我们愿意，仍能摆脱这种顽疾；不仅可能，而且容易；不仅容易，甚至比摆脱身体的痛苦更容易，因为医生更伟大。那么，让我们吸引祂，邀请祂帮助我们努力，并贡献我们的一份：即善意和努力。因为祂不会要求更多；只要遇到这一点，祂就会赐下祂的一份。让我们贡献我们的一份，好使我们在此世获得真正的健康，并达到未来的美善，藉着祂的恩宠和爱人之心等等。\par

\SourceRecord{Translated by J. Walker, J. Sheppard and H. Browne, and revised by George B. Stevens.；Nicene and Post-Nicene Fathers, First Series；Vol. 11.；Edited by Philip Schaff.；Buffalo, NY: Christian Literature Publishing Co.,；1889.；<http://www.newadvent.org/fathers/210214.htm>.}

% structure: p0001 source=h1 target=section reason=page-title

\section{《罗马书》第十五讲}

\BibleRef{罗 8:28}\par

\BibleQuote{我们知道，凡爱天主的，一切事都协助他们获得益处。}\par

这里，在我看来，他似乎是为了那些处于危险中的人而提出整个议题；或者，更确切地说，不仅为此，也为稍前所说的话。因为\Quote{现今时代的苦难，与将来要在我们身上显现的光荣相比，是不值得相提并论的}这句话，以及\Quote{整个受造物都在叹息}那句话，和\Quote{我们因希望而得救}的言论，以及\Quote{我们耐心等候}的表述，还有\Quote{我们不知道该怎样祈祷才合宜}这句话，都是对他们说的。因为他教导他们，不要只选择自己认为有益的事，而要选择圣神所建议的；因为许多看似对自己有益的事，有时甚至会带来很大的害处。例如，安静、免于危险、无忧无虑的生活，对他们似乎是有利的。这又如何令人惊讶呢，既然连蒙福的保禄自己也这样认为？但后来他认识到，与这一切相反的事才是有益的，一旦认识到这一点，他就满足了。所以那位三次求主免于危险的人，一旦听到主说：\BibleQuote{我的恩宠够你用，因为我的能力在软弱中才得以完全完备} \BibleRef{格后 12:8}，后来他竟以遭受迫害、侮辱和不可挽回的伤害为乐。因为他说：\BibleQuote{我以迫害、侮辱（英文译本作辱骂）、困窘为光荣。} \BibleRef{格后 12:10} 这就是他说\Quote{因为我们不知道该怎样祈祷才合宜}的原因。他劝告所有人将这些事交托给圣神。因为圣神非常顾念我们，这是天主的旨意。在用了各种方法鼓励他们之后，他转向我们今天所听到的，提出了一个足以挽回他们的理由。因为他说：\BibleQuote{我们知道，凡爱天主的，一切事都协助他们获得益处。} 现在，当他谈到\Quote{一切事}时，他甚至提到了那些看似痛苦的事。因为即使患难、贫穷、监禁、饥荒、死亡，或任何其他事情临到我们，天主都能将这一切变为相反的好事。因为这是祂不可言喻之大能的明证：祂使看似痛苦的事成为我们的光明，并使之转为有益。所以他并没有说，\Quote{凡爱天主的}不会遭遇苦难，而是说\Quote{一切事协助他们获得益处}，也就是说，祂利用苦难本身使被如此谋害的人得到炼净。这比阻止这些苦难的来临，或在它们到来时停止它们，要伟大得多。这正是祂在巴比伦的火窑中所做的。祂既没有阻止他们掉进去，也没有在那些圣人被投入后熄灭火焰，而是让火焰继续燃烧，却借这火焰使他们成为更令人惊奇的奇迹；同样，祂也藉着宗徒们不断行其他的类似奇迹。\BibleRef{谷 16:18} 因为如果学会哲学思考的人能利用自然事物达到相反的目的，甚至在贫困中显得比富人更安逸，在耻辱中发光，那么天主为爱祂的人所做的不止于此，而且远超这些。因为只需要一件事，就是对祂真诚的爱，其他一切都会随之而来。所以，正如看似有害的事对爱主的人有益，同样，看似有益的事也会伤害那些不爱主的人。例如，祂在教导中展示的奇迹和智慧只伤害了犹太人，正如道理的纯正一样；前者他们称祂为附魔的人 \BibleRef{若 8:48}，后者他们称祂为与天主平等的 \BibleRef{若 5:18}；并且因为奇迹（同上 11:47,53），他们甚至图谋杀害祂。但那位被钉死的强盗，当他被钉在十字架上，被辱骂，遭受无数的苦难时，不仅没有受损，反而因此获得了最大的益处。请看，对于爱天主的人，一切事如何协助他们获得益处。在提到了这个远超人性的伟大祝福之后——因为许多人甚至觉得这难以置信——他从过去的祝福中引出一个证明，说了\Quote{就是那些按祂旨意蒙召的人}。现在你要思考，他意思是，例如从蒙召来看我刚才所说的。那么当初祂为什么不召叫所有人？为什么不像召叫其他人那样立刻召叫保禄自己？难道拖延不是有害的吗？但事件仍然表明这是最好的。然而，他在这里提到\Quote{旨意}，是为了不把一切归于召叫；因为这样既会使希腊人也会使犹太人必定挑剔。因为如果仅凭召叫就足够，那么为什么不是所有人都得救呢？因此他说，成就救恩的不仅是召叫，还有被召者的意愿。因为召叫不是强加于他们的，也不是强迫的。所以，所有人都被召叫了，但不是所有人都服从了召叫。\par

\BibleQuote{因为祂所预知的人，也预定他们与祂圣子的肖像相同。} \BibleRef{罗 8:29}\par

看，何等崇高的荣誉！因为那独生子按本性是什么，他们藉恩宠也成了什么。然而他仍不满足于使他们与此相称的召叫，还加上了另一点：\Quote{祂成为长子。} 即使在这里他也没有停止，之后又提到另一点：\Quote{在众多弟兄中。} 所以他希望用一切方法阐明这种关系。现在你要把这些话都理解为是关于降生成人所说的。因为按天主性，祂是独生子。看，祂赐予了我们何等伟大的事！那么，不要怀疑未来。因为他甚至从其他根据上表明天主对我们的关怀，说事情从起初就这样预先安排了。因为人必须从事物中获得对它们的认识，但对天主而言，这些事早已决定，并且自古以来祂就对我们怀有善意 \OriginalTerm{祂对我们怀有善意}{\GreekText{πρὸς} \GreekText{ἡμἅς} \GreekText{διέκειτο}}，他说。\par

Ver. 30. \BibleQuote{此外，祂所预定的，也召叫了他们；祂所召叫的，也使他们成义。}\BibleRef{罗 8:30}\par

如今祂藉着圣洗的重生使他们成义。\BibleQuote{祂所成义的人，也使他们享受光荣}，藉着恩赐，藉着义子之名。\par

Ver. 31. \BibleQuote{那么，我们对这些事还有什么可说的？}\BibleRef{罗 8:31}\par

仿佛他这样说：那么，让我不再听见来自各方的危险和恶谋。因为即使有些人不相信未来的事，但他们对已经发生的善事也无话可说；例如，天主从起初对你们的友谊、成义、光荣。然而，这些事是祂藉着看似痛苦的方式赐给你们的。那些你们以为羞辱的事——十字架、鞭打、锁链——正是这些使整个世界得以矫正。正如祂藉着自己所受的苦难（尽管在人眼中面目可憎），甚至藉着这些，祂为全人类实现了自由和救恩；同样，祂也习惯于对待你们所忍受的那些事，将你们的痛苦转化为光荣和名声。\BibleQuote{如果天主偕同我们，谁能反对我们？}\par

也许有人会说：谁不反对我们呢？世界反对我们，君王和人民、亲戚和同胞都反对我们。然而这些反对我们的人，远不能阻挠我们，甚至在他们不情愿的情况下，也成了我们冠冕的缘由和无数祝福的提供者，因为天主的智慧将他们的阴谋转变为我们的救恩和光荣。看，实际上没有人反对我们！因为正是这一点使\Person{约伯}焕然一新，即魔鬼武装反对他的事实。因为魔鬼立刻发动朋友反对他、妻子反对他、创伤、仆人和无数其他阴谋。结果发现，总体上没有一样是反对他的。然而，这对他不算什么大事，尽管本身是伟大的，但更伟大的事实是，它们竟然都为他效力。这也发生在宗徒们身上，因为犹太人、外邦人、假兄弟、统治者、人民、饥荒、贫穷和无数事物反对他们，但没有什么反对他们。因为使他们最光辉、最引人注目、在天主和人眼中伟大的，正是这些。那么请思想，\Person{保禄}关于信友们以及那些真正（\OriginalTerm{精确地}{\GreekText{ἀκριβῶς}}）被钉十字架的人说了怎样的话，即使是头戴王冠的皇帝也无法成就。因为反对他的有大量武装的蛮族、入侵的敌人、阴谋的卫兵、众多时常反叛的臣民，以及成千上万的其他事物。但反对那谨守天主法律的信友，没有人、魔鬼或其他任何东西可以站立！因为如果你夺走他的钱财，你成了他奖赏的提供者。如果你说他坏话，藉着恶名他在天主眼中获得新的光彩。如果你使他挨饿，他的光荣和奖赏会更大。如果你（似乎是最严重的一击）置他于死地，你是在为他编织殉道的冠冕。那么，有什么能与这种生活方式等同呢？它无往不利，甚至那些看似图谋不轨的人对他的服务也不亚于恩人。因此他说：\BibleQuote{如果天主偕同我们，谁能反对我们？}接下来，他不满足于已经说的话，还在此处列出了祂对我们爱的最大标志，即祂始终反复思考的那一点：祂圣子的被杀害。因为祂的意思是，祂不仅使我们成义、得光荣、使我们符合那肖像，而且为了你们连祂的圣子也没有顾惜。因此祂接着说：\par

Ver. 32. \BibleQuote{祂没有顾惜自己的儿子，反而为我们众人把祂交出，岂不也把一切同祂一同赐给我们吗？}\BibleRef{罗 8:32}\par

在这里，他使用的语言是极度（\OriginalTerm{极度}{\GreekText{μεθ᾿} \GreekText{ὑπερβολῆς}}）热情洋溢的，以显示他的爱。那么，祂怎么会忽视我们呢？为了我们，祂连自己的儿子都没有顾惜，反而把祂交出。请思考，这是何等良善：不仅不顾惜自己的儿子，反而把祂交出，且为所有人交出，为那些无价值的、无情的、敌人和亵渎者。那么，祂岂不也把一切同祂一同赐给我们吗？他的意思大概如此：如果祂给了自己的儿子，不仅是给予，而是赐予死亡，那么既然你有了主，为何还怀疑其余的事？既然你有了主，为何还怀疑财物？因为那将更大的赐给了祂敌人的人，难道不也将较小的赐给祂的朋友吗？\par

Ver. 33. \BibleQuote{谁能控告天主所拣选的人呢？}\BibleRef{罗 8:33}\par

在这里，他反对那些说信德无益、不相信彻底改变的人（即在圣洗中，见第349页）。请看，他如何迅速堵住他们的嘴，凭藉拣选者的尊贵。他没有说：\Quote{谁能控告天主的}仆人呢？或\Quote{天主的信友？}而是说\BibleQuote{天主所拣选的人？}拣选是德行的标志。因为如果驯马师挑选了适合赛马的小马，没有人能找它们的毛病，找毛病的人会被人嘲笑；更何况当天主拣选灵魂时，那些控告他们的人是可笑的。\par

\BibleQuote{是天主使人成义。}\par

Ver. 34. \BibleQuote{谁能定他们的罪？}\BibleRef{罗 8:34}\par

祂没有说，是天主赦免了我们的罪，而是更伟大的：\BibleQuote{是天主使人成义。}因为当法官的判决宣告我们正义，而且是这样一位法官时，控告者有什么意义？因此，既不该害怕诱惑，因为天主为我们，并藉着祂的作为表明了；也不该害怕犹太人的琐碎小事，因为祂既拣选了我们，又使我们成义，而奇妙的是，祂也是藉着祂圣子的死亡这样做的。那么，既然天主给我们加冕，基督为我们而死，而且不仅死了，之后还为我们转求，谁能定我们的罪呢？\par

因为，他说：\Quote{基督}，\Quote{是那死了的，但更是从死者中复活的，现今在天主右边，也为我们转求。}\par

因为，虽然现在看见祂在自己的尊位中，祂并未停止关心我们，反而甚至\Quote{为我们转求}，并依然保持同样的爱。因为祂不满足于仅仅被处死。这在大部分情况下是极伟大爱情的表现，不仅做自己分内的事，还为此向另一位求情。因为这是他藉着转求所要表明的一切，采用一种更适合人且更俯就人情的说话方式，为的是彰显爱。因为除非我们以同样的理解来对待\Quote{祂没有顾惜}这话，否则会产生许多矛盾。为使你看见这正是他所指的方向，在首先说了祂\Quote{坐在右边}之后，他接着就说祂\Quote{为我们转求}，那时他已表明了尊荣与品级的平等，好由此看出转求不是低微的标志，而只是爱的标志。因为祂是生命本身\OriginalTerm{生命本身}{\GreekText{αὐτοζωή}}\BibleRef{咏 36:10}，是一切美善之物的泉源，并拥有与圣父同样的权能——既能复活死者，也能赐予生命，以及做其他一切祂所做之事——祂如何需要作为恳求者来帮助我们？\BibleRef{若 5:19}  祂凭自己的力量释放了那些被交付并定罪的人，甚至从那定案中释放了他们；使他们成为义人、成为儿子，引领他们达至最高的尊荣，成就了从未希望过的事：在完成这一切，并让我们的本性显于君王宝座之后，祂怎会需要成为一个恳求者来做更容易的事呢？\BibleRef{宗 7:55}  你看，每一条论证都表明，他之所以提到转求，没有别的原因，只是为了显示祂对我们之爱的温暖与活力；因为圣父也被描绘成恳求人们与祂和好。\Quote{我们作基督的使者，就好像天主藉着我们劝勉你们一般。}\BibleRef{格后 5:20}  尽管天主恳求，人们\Quote{以基督的名义作使者}到人中间，但我们并不因此理解任何有损于那尊位的事；我们只从一切告诉我们的中吸取一点，就是爱的炽烈。那么，让我们在这里也这样做吧。如果连圣神也\Quote{以无可言喻的叹息为我们转求}，基督又死了并为我们转求，圣父也\Quote{没有顾惜自己的儿子}为你，而拣选了你，成义了你，那为何还要害怕呢？或者，当如此伟大的爱眷顾着你，为何还要战栗呢？这样，在自从起初就向我们展示祂的伟大眷顾之后，他随后以大胆的笔法引出下面的话，并没有说：\Quote{你们也应当爱祂}，而是仿佛被这无法言喻的宰制我们的眷顾所激励，他说：\par

第35节。\BibleRef{罗 8:35} \BibleQuote{谁能使我们与基督的爱隔绝？}\par

他没有说\Quote{天主}，对他来说提到基督或天主的名无所谓。\Quote{是困苦吗？是窘迫吗？是迫害吗？是饥饿吗？是赤身吗？是危险吗？是刀剑吗？} 请看有福的保禄的判断。因为他没有提到我们每日被夺取的东西——贪爱金钱、虚荣欲望以及愤怒的束缚——而是列举了比这些更令人着迷、能对自然本身施加压力、并多次违背我们的意愿而撬开决心之严厉的事物，即困苦和窘迫。因为虽然所提及的事易说清，但每一个词中都包含成千上万条诱惑的线索。因当他提到困苦时，他就提到了监禁、锁链、诬蔑、流放和所有其他艰难，因此用一个词不吝啬地历遍了危险的大洋，实际上用一个词向我们展示了人类所遇到的一切邪恶。然而，他仍向它们挑战！因此他以提问的方式提出它们，仿佛无可争议：没有什么能动摇一个如此被爱、并享受如此大眷顾的人。然后，为了避免显得他忘乎所以，他也引用了很久以前就预言此事的先知，说：\par

第36节。\BibleRef{罗 8:36} \BibleQuote{我们为了你的缘故终日被杀，被视作待宰的羊。}\BibleRef{咏 44:22}\par

这就是说，我们被暴露在所有人面前，任凭他们恶待。然而面对如此众多且巨大的危险，以及这些最近的恐怖之事，我们斗争的目标给了我们充分的安慰——不，不仅是充分的，甚至更多。因为我们受苦不是为人类，也不是为今世的其他任何事物，而是为宇宙的君王（他说）。但这并不是唯一的花冠；他还用另一顶花冠环绕他们，那花冠多样且繁多。由于他们作为人，无法承受无数次的死亡，他表明这样奖赏并不减少。因为即使按本性注定只能死一次，天主却藉着选择允许我们每天受苦，只要我们愿意。由此可知，我们离开时会有多少天的花冠，或者更多。因为在一天之内，不仅可能死一次或两次，而是许多次。因为那始终准备赴死的人，不断领受圆满的赏报。这就是圣咏作者（\OriginalTerm{先知}{\GreekText{Προφήτης}}）所暗示的，当他说\Quote{终日}时。为此，宗徒也把他引到他们面前，以更加振奋他们。他的意思是：如果旧约时代那些以土地为赏报，以及其它与今世一同结束之事物的人，尚且如此轻视现世生命及其诱惑和危险，那么我们蒙受天堂、天国及其不可言喻的福乐之应许后，若仍如此懈怠，以至于连他们的标准都达不到，又怎能得到宽恕呢？他虽然没有明说，却将此留于听者的良心，仅以引用为满足。他也表明他们的身体成了祭品，并且我们不应因天主如此安排而困扰或不安。此外，他还以其他方式劝勉他们。为了防止有人说他只是在没有任何实际经验的情况下在这里空谈哲理，他补充说：\Quote{我们被看作待宰的羊群，}意指宗徒们每日的死亡。你看他的勇气和良善。他的意思是：正如羊群被宰杀时不抵抗，我们也是如此。但既然人灵因软弱，即使经历了如此大事，仍害怕众多诱惑，请看他又如何再次振奋听者，赐予他崇高而欢欣的精神，说道：\par

\BibleQuote{然而，靠着那爱我们的主，在这一切事上，我们大获全胜。} \BibleRef{罗 8:37}\par

因为真正奇妙的是，我们不仅得胜，而且正是藉着那些针对我们的阴谋而得以得胜。我们不仅是得胜者，更是\Quote{大获全胜}，也就是说，轻而易举，不费辛劳。因为无需经历实际的事，只须调整我们的心思，我们就能向敌人举起凯旋的旗帜。这理所应当，因为天主与我们一同奋斗。所以不要疑惑：即使我们被打，我们却胜过了打我们的人；即使被驱逐，我们却胜过了迫害我们的人；即使死亡，我们却使活着的人战斗。因为当你考虑到天主的权能和祂的爱时，没有任何事能阻止这些奇妙而奇异的事发生，且最有益的胜利应照在我们身上。因为他们不仅得胜，而且以奇妙的方式得胜，使人得知那些谋害他们的人所进行的战争不是针对人，而是针对那不可战胜的大能。请看当时犹太人处在他们中间，完全不知所措，说：\BibleQuote{我们拿这些人怎么办？} \BibleRef{宗 4:16} 这确实奇妙：虽然他们抓了他们，把他们置于法庭之下，监禁他们，鞭打他们，却仍不知所措、困惑不解，因为他们被那些本指望用以征服的事所征服。无论君王、民众、魔鬼的等级，还是魔鬼自己，都无力胜过他们，反而都大败亏输，发现他们一切针对他们的计谋都变成了他们的优势。因此他说：\Quote{我们大获全胜。} 因为这是一种新的胜利法则：人藉着敌人获胜，在任何情况下都不被征服，而是投身这些争战，仿佛结局掌握在自己手中。\par

\BibleQuote{因为我深信：无论是死亡，是生命，是天使，是掌权者，是异能，是现存的，是将来的，是高天，是深渊，或是任何别的受造物，都不能使我们与天主的爱隔绝，这爱是在我们的主基督耶稣内的。}\par

这里所提及的都是伟大的事物。但我们之所以不能深入理解，是因为我们没有如此伟大的爱。然而，尽管它们伟大，正如他想要表明它们与他被天主所爱的爱相比算不得什么；在提及那爱之后，他才放置自己的爱，免得他似乎在说自己伟大的事。而他所说的大致如此。他说，我提现世的事物和今生承受的邪恶做什么呢？因为即使有人告诉我将来的事和权能；诸如死亡和生命之类的事；权能之类的事，诸如天使和总领天使，以及所有更高的品级；这些与基督的爱相比，对我也是微不足道的。因为即使有人用那永死来威胁我，使我与基督分离，或者应许我不朽的生命，我也不会同意。何必提世上的君王和执政官呢？何必提这个或那个呢？因为如果你告诉我天使，或上面所有的权能者，或所有存在的事物，或所有将来的事物，对我来说，无论是地上的、天上的、地下的、还是超越天上的，与这爱慕相比，都是渺小的。然后，好像这些还不足以向他们展示他强烈的渴望，他另外赋予同样伟大的存在，并说：\BibleQuote{也没有任何其他受造物}\BibleRef{罗 8:39}。他所说的几乎就是：即使有其他受造物，像可见的与可理解的那么伟大，也没有一个能使我与那爱分离。他这样说，并非因为天使或其他权能者曾试图这样做，绝非如此，而是为了极其充分地表明他对基督的爱慕。因为他爱基督，并非因基督的事物，而是因基督本身而爱属于基督的事物，他只注视基督，他只畏惧一件事，就是失去对基督的爱。因为这事本身比地狱更可怕，而持守这爱比天国更可向往。\par

那么，我们如今该当如何呢？当发现他（保禄）连天上的事物也不如对基督的渴慕，而我们将污泥和尘土的事看得比基督更重时。他因渴望基督，甚至愿意堕入地狱、被逐出天国，若让他在两者间选择的话；但我们却连现世生命都不能超脱。我们与他那宽广的心胸相差如此之远，还配触摸他的鞋吗？他为了基督的缘故，连王国都看如无物；而我们却轻看基督自身，反倒看重祂的事物。但愿我们看重的是祂的事物。但现在连这也不是；天国摆在我们面前，我们却弃之不顾，终日追逐阴影和梦幻。然而，天主因祂对人的爱与无比的温和，做了如同一位慈父所做的事：当他的儿子不愿长久与他同住时，他便以另一种方式促成此事。因为我们没有正当的渴望之情，祂便不断将各种事物放在我们面前，好将我们拉住。即便如此，我们仍不与祂同在，反而不断跳向孩童的玩物。保禄却不是这样；他就像一位心地高尚的孩子，坦诚依恋父亲，只寻求父亲的同在，而不甚看重其他事物；或者说，他比孩子更胜一筹。因他不把父亲与父亲的事物同等看待；当他注视父亲时，便视那些事物为无物，宁愿受责打惩罚而与父亲同在，也不愿离开父亲而享受安逸。为此，我们所有人都该战栗——我们这些为了基督甚至不能超脱钱财的人，不如说我们这些连为了自己也不能超脱钱财的人。因为只有保禄真正为了基督的缘故忍受一切，不是为了天国或自己的荣誉，而是出于对基督的深情。至于我们，基督和基督的事物都不能将我们从今生的事物中拉开；我们却像蛇、蛆虫、猪，甚至像所有这些动物合在一起，在泥淖中拖曳。因为当如此众多伟大的榜样摆在我们面前时，我们若仍低头弯腰，甚至片刻不敢仰望天堂，我们比那些兽类好在哪里呢？然而，天主甚至舍弃了祂的圣子。你却连一块面包也不肯与祂分享——这位为你而被舍弃、为你而被宰杀的基督。圣父为了你，没有顾惜祂，而这圣子确实是祂的圣子；但祂忍饥挨饿时，你却不看顾祂，何况你只需花费祂自己的东西，且花费是为了你自己的益处！还有什么比这更严重的违法呢？祂为你被舍弃，为你被宰杀，为你忍饥挨饿；你应从祂所赐的中施舍，好使你自己获益，你竟不施舍！哪一块石头比这些人的感觉更迟钝？他们面对如此巨大的感召，仍执迷于这种魔鬼般的铁石心肠。因为祂不仅满足于死亡和十字架，还甘愿成为贫穷者、陌生者、乞讨者、赤身露体者、被投入监牢者、患病者，为的是至少能唤你回头。祂说：你若不愿因我为你受苦而报答我，至少因我的贫穷而怜悯我。你若不愿因我的贫穷而怜恤我，至少因我的疾病而感动，因我的监禁而软化。若这些仍不能使你慷慨，至少因要求的轻易而顺从。因为我所求的不是昂贵的礼物，而是面包、住处和安慰的话。但若连这之后你仍不软化，至少因我所应许的赏报——为了天国——而改进。难道你对这些也无动于衷？然而，至少因着本性，你看到我赤身露体时，当受感动；要记得我在十字架上为你赤身露体；若不然，也要记得我如今借着穷人赤身露体。我那时为你被捆绑，如今仍为你被捆绑，为要使你或因前者、或因后者而受到感动，生出一些怜悯。我为你禁食，如今又为你饥饿。我挂在十字架上时口渴，如今也借着穷人而口渴，为要借着前者和后者将你吸引到我身边，使你为自己的得救而慷慨施舍。因此，你欠我数不清的恩惠，我并不像向欠债者那样向你索取，反而像宠爱我者那样加冕你，并因这些小事赐你天国。因为我甚至没有说：结束我的贫穷，或给我财富，尽管我曾为你成为贫穷；然而我只求面包、衣服和对我饥饿的一点慰藉。若我被投入监牢，我不要求你解开我的捆绑、释放我；我只寻求一件事：请你探望我——这位为你被捆绑的——我便心满意足；仅为此，我将赐你天堂。然而，我曾将你从最紧的捆绑中释放；而对我，你若在监牢中探望我便足够了。因为我本可不需这一切就加冕你；但我愿成为你的债务人，好使那冠冕给你带来一些信心。这就是为何，我本可养活自己，却前来乞讨，站在你的门口，伸我的手；因为我渴望由你供养。我非常爱你，所以渴望在你的桌上就餐——这是爱一个人的表现。我以此为荣（见：\BibleRef{若 15:8}）。当整个世界作观众时，我将宣告你，在众人面前展示你是我的供养者。然而，当我们受人供养时，会感到羞愧并掩面；但祂极其爱我们，即使我们沉默，祂也会带着许多赞美述说我们所做的，并不羞于说：祂赤身时我们给祂穿上，祂饥饿时我们给祂吃的。为此，让我们将这些都放在心里，不要满足于对它们的空洞赞美，而要实际实行我所说的。因为这些掌声和欢呼有什么益处呢？我只向你们要求一件事：在实际行动中表现出来，即以行为服从。这是我的赞美，你们的收益，这比任何冠冕更使我光彩。当你们离开教会后，应当借着穷人之手为我、也为你们编织这冠冕；好使我们今生怀着美好的希望得到滋养，并在离世后获得那无数的美善。愿我们众人借着祂对人的恩宠与慈爱，都能达到这些美善，等等。\par

\SourceRecord{Translated by J. Walker, J. Sheppard and H. Browne, and revised by George B. Stevens.；Nicene and Post-Nicene Fathers, First Series；Vol. 11.；Edited by Philip Schaff.；Buffalo, NY: Christian Literature Publishing Co.,；1889.；<http://www.newadvent.org/fathers/210215.htm>.}

% structure: p0001 source=h1 target=section reason=page-title

\section{罗马书第十六讲}

\BibleRef{罗 9:1}\par

\BibleQuote{我在基督内说实话，并不说谎，我的良心藉着圣神也为我作证。}\par

昨天我向你们谈论保禄对基督的爱时，难道不显得说了些伟大而过分的话吗？那些话确实伟大，伟大到言语无法表达。然而，你们今天所听到的，比那些话更超越，如同那些话超越我们的一样。我本以为它们已无法被超越，但当我读到今天这段经文时，它显得比先前的一切都更加光荣。保禄自己也知道这一点，他的开场白便显示了这一点。因为当他正要谈论比那些更伟大的事，因而容易受到一般人的怀疑时，他首先就他将要说的事发出强烈的誓言；许多人在要说出一些不为一般人所相信、而他们自己心中却万分确定的话时，常有这样的习惯。因此他说：\Quote{我在基督内说实话，并不说谎，我的良心也为我作证，}\par

第2至3节。\BibleQuote{我甚忧愁，内心不断痛苦；因为我宁愿自己受诅咒，与基督隔绝。}\par

第2至3节。\BibleQuote{我甚忧愁，内心不断痛苦；因为我宁愿自己受诅咒，与基督隔绝。}\newline{}\newline{}\Quote{保禄啊，你说什么？从基督，你所爱的那一位，无论是王国还是地狱，是可见的或理智的事物，还是另一个同样伟大的世界，都不能使你与祂分离，现在你却要受诅咒，与祂分离吗？发生了什么事？你改变了，放弃了那份爱吗？不，他回答说，不要害怕。我倒使它更加炽烈了。那么，你怎么会愿意受诅咒，寻求分离，并且远离到如此地步，以至于在那之后不可能找到更远的距离呢？因为他会说：因为我极其爱祂。请问，是怎样的爱？因为这事似乎是个谜。或者，如果你愿意，让我们先了解什么是诅咒，然后我们将就这些点质问他，从而明白这不可言说且非凡的爱。那么，什么是诅咒？听他自己的话：\BibleQuote{若有人不爱我们的主耶稣基督，就让他受诅咒。}\BibleRef{格前 16:22} 这就是说，让他与众人分离，从众人中除去。因为正如奉献之物（\OriginalTerm{奉献之物}{\GreekText{ἀνάθημα}}的情况），它被奉献给天主，没有人敢用手触摸它，甚至不敢靠近它；同样，对于一个与教会分离的人，将他与众人隔绝，尽可能远地移除，他称呼他为\OriginalTerm{诅咒}{\GreekText{ἀνάθεμα}}，这是相反的意义，从而以极大的恐惧告诫所有人要远离他，避开他。因为那奉献之物，出于尊敬，没有人敢靠近。但与被隔绝的人，所有人都因完全相反的情感而远离他。因此，分离是相同的，二者同样被从人群中移除。然而，分离的方式并不相同，在这种情况下，它与此相反。因为对前者他们回避，因为是奉献给天主的；对后者，因为是疏离了天主，并与教会断绝了关系。这正是保禄说\Quote{我宁愿自己受诅咒，与基督隔绝}时的意思。他不仅说\Quote{我愿意}，而且用了一个更强的词，他甚至说：\Quote{我祈求}（\OriginalTerm{祈求}{\GreekText{ηὐχόμην}}）。但若他所说的在你（\OriginalTerm{较软弱的}{\GreekText{ἀσθενέστερον}}）的软弱中困扰你，请考虑实际情况：不仅他愿意被分离，而且他愿意分离的原因，然后你就会看到他的爱是多么伟大。因为他甚至行了割损礼\BibleRef{宗 16:3}，我们不注意所做的事，而是注意其意图和原因，因此我们更加钦佩他。他不仅给一个人行了割损礼，甚至剃了头并献了祭\BibleRef{宗 18:18}，但我们当然不因此断言他是犹太人，反而因此理由认为他完全脱离了犹太化，是基督的真正崇拜者。正如当你看见他行割损和献祭时，你并不因此谴责他犹太化，反而因此有最好的理由为他加冕，因为他完全疏离了犹太主义；同样，当你看见他渴望受诅咒时，不要因此困扰，反而当你知道他愿意如此的原因时，要给予他最热烈的赞扬。因为如果我们不仔细考察原因，我们会称厄里亚为杀人者，亚巴郎不仅是杀人者，而且是他儿子的谋杀者。我们也会控告丕乃哈斯和伯多禄犯有谋杀罪。不仅是在圣人的情况下，而且对于宇宙的天主，不遵循这一规则的人，会怀疑各种不当之事。现在，为了防止在所有类似情况下发生这样的事，让我们将原因、意图、时间以及所有有利于如此行事的事情汇集起来，并以这种方式研究行动。现在我们也必须为这位有福的灵魂这样做。那么原因是什么？就是那如此被爱的耶稣自己。但他没有说\Emph{为了}祂；因为他所说的是：我愿意为我的弟兄们从祂受诅咒。这源于他的谦卑。因为他不想引人注目，好像他在说什么伟大的事，并在这事上给基督帮忙。所以他称他们为\Quote{我的骨肉}，以隐藏他崇高的目标（\OriginalTerm{优越}{\GreekText{πλεονέκτημα}}）。因为要知道他完全为了基督而愿意这一切，请听接下来所说的。在提及骨肉之后，他继续道：}\par

第4至5节。\BibleQuote{他们是以色列人，义子的名分、光荣、盟约、法律的颁布、礼仪以及恩许，都属于他们；圣祖也是他们的，并且基督按血统说，也是从他们来的，祂是在万有之上，世世代代应受赞美的天主！亚孟。}\par

这是什么意思？有人会问。因为如果为了使别人相信，他甘愿成为\OriginalTerm{被诅咒的}{accursed}，那么他也理应在外邦人身上盼望同样的事。但如果他只愿在犹太人身上如此，就证明他不是为基督的缘故，而是为了自己与他们的关系。但事实上，如果他只为外邦人祈祷，这一点就不那么明显。然而既然只为犹太人，就清楚证明他这样热切完全是出于基督的光荣。而且我知道，我所说的话在你们看来似乎是悖论。但你们若不起哄，我即刻设法解释清楚。因为他所说的话并非直白说出；但因众人议论纷纷，指责天主：他们既被配称为义子，领受法律，比所有人更认识祂，享受如此大的光荣，比全世界更事奉祂，领受许诺，出自是祂朋友的祖先，而最伟大的是，成为基督本人的祖先（因为这句话的含义就是：\Quote{出自这些人，按\OriginalTerm{血肉}{flesh}说，基督是从他们来的}），如今却被驱逐和羞辱；代替他们的是从未认识祂的外邦人。既然他们说了这一切，亵渎了天主，保禄听见了，心如刀割，为天主的荣耀而焦虑，就希望自己成为被诅咒的，若可能的话，好使他们得救，亵渎得以止息，天主似乎没有欺骗那些蒙祂应许赐予恩惠者的后代。为使你们明白，他因恐怕天主的许诺落空而忧伤，这许诺对\Person{亚巴郎}说：\Quote{我要将这地赐给你和你的\OriginalTerm{后裔}{seed}}，他就发出这愿望，接着说道：\par

第6节。\Quote{并不是说天主的圣言失效了。}\BibleRef{罗 9:6}\par

为显示他之所以胆敢（\EditorialNote{马尔谷抄本和四个抄本作\Quote{希望}}）为\OriginalTerm{天主的圣言}{the word of God}承受这一切，即对\Person{亚巴郎}所作的\OriginalTerm{许诺}{promise}。因为正如\Person{梅瑟}似乎在为犹太人辩护，却一切为天主的光荣而行（因为他说：\Quote{恐怕他们说：因为祂不能救他们，就领他们出埃及，要在旷野杀害他们。}\BibleRef{申 9:28}；求祢止息祢的忿怒），同样，保禄也是如此。他的意思是：唯恐他们说天主的许诺落了空，祂使我们失望，没有兑现所应许的，这话没有实现，所以我情愿成为被诅咒的。这也就是为何他不提外邦人（因为天主未曾向他们作过许诺，他们也没有敬拜祂，因此也没有人因他们亵渎祂）；正是为了那些既领受许诺、又比他人与祂有更亲密关系的犹太人，他才表达了这个愿望。你们看见了么？如果他愿为外邦人成为\OriginalTerm{被诅咒的}{accursed}，就不会如此清楚地显出这是纯粹为基督的光荣。但由于他愿为犹太人成为被诅咒的，最明显地证明了他所求的唯一是为基督。为此他说：\par

\Quote{他们拥有\OriginalTerm{义子}{adoption}、\OriginalTerm{光荣}{glory}、\OriginalTerm{敬拜天主}{service of God}和\OriginalTerm{许诺}{promises}。}\par

因为法律，他说，那宣讲基督的，是从那里来的，所有与他们订立的盟约，以及基督自己是从他们而来的，承受恩许的祖先也都是从他们来的。然而，结果却相反，他们从一切福分中坠落了。因此，他说，我苦恼；如果可能的话，我宁愿与基督的同伴分离，成为局外人，不是离开对祂的爱（这在他绝不可能；因为他所做的一切都是出于爱），而是离开那一切享乐和光荣，我愿意接受那种命运，只要我的主人不被亵渎，不让祂听到有人说，这一切都是做戏；祂向一个人许诺，却给了另一个人。祂从一个民族出身，却拯救了另一个民族。祂曾向犹太人的祖先赐下恩许，却遗弃了他们的子孙，把那些从未认识祂的人放在了他们的福分中。他们辛苦遵行法律，诵读先知书，而那些昨天还从异教祭台和偶像那里来的人，却被高举在他们之上。这一切中有何远见呢？如今，他说，为了不让我的主人被这样议论，即使是不公正的议论，我也情愿失去天国和那不可言喻的光荣，并承受任何苦难，因为我认为，最大的安慰莫过于不再听见我所切慕的祂被如此亵渎。但如果你们仍然不接受这个解释，只需想想，许多父母常常为儿女承受如此之多，选择与他们分离，而宁愿看到他们受尊敬，认为他们的尊荣比与己同在更为宝贵。但既然我们如此缺乏这样的爱（培根, N. O. Aph. lib. 2, §7），我们甚至无法想象这里的意思。因为有些人完全不配听到保禄的名字，与他的热诚相距如此遥远，以至于以为他说的是暂时的死亡。我敢说，他们对保禄的无知，就如盲人对阳光的无知，甚至更加严重。因为那位天天死亡，把危险如雪暴般放在面前的人，然后说：\BibleQuote{谁能使我们与基督的爱隔绝呢？是困苦吗？是窘迫吗？是迫害吗？是饥饿吗？} 并且对自己所说的话仍不满足，在上升到诸天和天上的天之后，穿越天使、总领天使和所有更高的品级，同时把握现在和将来的事，有形和无形的事，痛苦和幸福的事，左边和右边的一切，毫无遗漏，却仍不满足，甚至构想出另一个不存在的受造物，在所有这一切之后，他怎能为了说出一件大事而提及暂时的死亡呢？不是这样，绝对不是！但这种观念就像蛆虫在粪堆中蠕动。因为如果他真的这样说，他怎么会希望自己与基督隔绝呢？因为那种死亡\BibleRef{斐 1:23}会更紧密地将他与基督的队伍联合，使他更多地享受那光荣。然而，有些人竟敢说出与此不同甚至更荒谬的话。他们说，他所希望的不是死亡，而是成为基督的宝藏，一个分别出来的事物。谁，即使是最卑贱懒惰的人，不愿这样呢？而这对他的同胞又有什么益处呢？那么，让我们抛开这些神话和琐事（因为对这些东西的回答不比回答玩耍中的孩子更有价值），让我们回到这些话本身，徜徉在这爱的海洋中，无所畏惧地向四面八方游去，默思那不可言说的爱火——或者，无论说什么，都无法说尽这个主题。因为没有如此广阔的海洋，没有如此炽烈的火焰。任何语言都无法恰如其分地表达它，只有真正获得它的人才知道它。现在，让我再把这些话带到你们面前。\par

\BibleQuote{因为我曾切愿我自己被诅咒。} 这个\Quote{我自己}是什么意思？意思是：我，作为万人的教师\BibleRef{格前 9:27}，积累了无数的善功，等待无数的冠冕，如此渴慕祂，以至于视祂的爱高于一切，我终日为祂燃烧，视万物\BibleRef{斐 3:8}为次要，相比对祂的爱。因为不仅被基督爱是他所关心的，而且极其爱基督也是他所关心的。而他最关心的是后者\OriginalTerm{他最为关注}{\GreekText{τούτου} \GreekText{μάλιστα} \GreekText{ἦν}}。因此，他只看重这一点，并对一切事泰然处之。因为他在一切环境中只持守一个目标：实现这卓越的爱。这也是他所愿望的。但既然事情不会如此发展，他也不会被诅咒，于是他接着尝试进行辩护，以反驳那些指控，从而否定众人所流传的话。在他公开辩护之前，他预先埋下一些种子。因为当他说：\BibleQuote{他们拥有义子的名分、光荣、法律、礼仪和恩许}，他只是在说天主确实愿意他们得救，并且祂通过先前的作为、通过基督从他们而出、通过祂向祖先所作的恩许，显明了这一点。但他们由于自身顽固的性情，把恩惠推开了。这也是他列出这些描述天主恩赐、而非赞扬他们的事物的原因。因为义子的名分来自祂的恩宠，光荣、恩许和法律也是如此。在考虑了这一切之后，反省天主连同祂的圣子对他们得救的热切渴望，他高声说：\BibleQuote{祂是永远应受赞美的。阿们。}\par

因此他亲自为众人向天主的独生子献上感恩。他说：即使别人亵渎又怎样？我们这些知道祂的奥迹、祂不可言喻的智慧，以及祂对我们伟大眷顾的人，深知祂不是应受亵渎，而是应受光荣。然而，他并不满足于自己意识到这一点，接着努力运用论据，并以更尖锐的方式对他们说话。他没有把矛头指向他们，而是先消除他们可能有的一个疑虑。为避免他显得在把他们当作敌人来对待，他在后面说：\BibleQuote{弟兄们，我心里所喜悦的，向天主所求的，是要以色列人得救。}\BibleRef{罗 10:1} 在这里，连同其他言论，他如此安排，以便不显得他说的话是出于对犹太人的敌意。因此，他甚至不避讳称他们为亲人和弟兄。因为即使他是为了基督的缘故才说这些话，他仍然\OriginalTerm{在吸引}{\GreekText{ἐπισπἅται}}他们的心归向自己，为他要说的话铺路，并消除因必须说的反对他们的话而引起的一切怀疑，然后他才最终进入大多数人所期待的主题。因为正如我已经说过的，许多人想知道为什么那些承受了应许的人反而没有得着，而那些从未听说过应许的人却比他们先得救。因此，为了澄清这个难题，他在提出异议之前先给出答案。为阻止有人说：\Quote{什么？你比天主更关心祂的光荣吗？难道天主需要你的帮助，免得祂的话落空吗？}针对这些，他说：\Quote{我说这话，并非因为天主的话落空了，而是为了显明我对基督的爱。因为事情既然这样发展，我们并不缺少为天主辩护的话，也不缺少证明祂的应许确实站立得住。}天主对亚巴郎说：\BibleQuote{我要将这地赐给你和你的后裔。}以及\BibleQuote{地上万民都要因你的后裔蒙受祝福。}\BibleRef{创 12:7} 那么，让我们看看，他说，这后裔是什么样的。因为并非所有从他而出的都是他的后裔。因此他说：\BibleQuote{因为从以色列生的，不都是以色列人（或出自以色列）。}\BibleRef{罗 9:6}\par

第7节：\BibleQuote{也不因为是亚巴郎的后裔，就都是他的子女。}\BibleRef{罗 9:7}\par

当你认识亚巴郎的后裔是什么样子时，你就会看到那应许是赐给他的后裔的，并且知道那话没有落空。请问，这后裔是怎样的呢？我的意思是，这不是我自己的话，而是旧约本身通过如下的话来解释自己：\BibleQuote{从依撒格生的，才称为你的后裔。}\BibleRef{创 21:12} \Quote{从依撒格生的}是什么意思？请解释。\par

第8节：\BibleQuote{这就是说，肉身所生的子女不是天主的子女，惟有那应许的子女才算是后裔。}\BibleRef{罗 9:8}\par

请看保禄的见识和深度。他在解释时，并没有说\Quote{肉身所生的子女不是}亚巴郎的子女，而是说\Quote{天主的子女}，从而把旧事与现今融合在一起，表明连依撒格也不仅仅是亚巴郎的儿子。他的意思大致是这样：凡像依撒格那样出生的，就是天主的儿子，也是亚巴郎的后裔。这就是为什么他说：\BibleQuote{从依撒格生的，才称为你的后裔。}\BibleRef{创 21:12} 好叫人知道，那些按依撒格方式出生的人，才是亚巴郎真正的子女。那么依撒格是怎样出生的呢？不是按照自然律，不是按照肉身的能力，而是按照应许的能力。那么\Quote{应许}的能力指的是什么呢？\par

第9节：\BibleQuote{到明年这时候我要来，撒辣必生一个儿子。}\BibleRef{罗 9:9}\par

这应许和天主的话塑造了依撒格，生了他。因为即使子宫和妇人的肚腹是工具，但生这孩子不是靠肚腹的能力，而是靠应许的德能。我们也是这样由天主的话所生。因为在水中之池，是天主的话产生并塑造我们。我们是因受洗归于父、子及圣神之名而生的。\BibleRef{若 3:3} 这种出生不是出于本性，而是出于天主的应许。因为正如祂先预言依撒格的出生，然后才成就它；同样，祂也早已在多个世纪以前通过众先知预言了我们的出生，后来才使之实现。你们知道祂把这出生说得多么伟大，并且祂应许了大事，却以极大的轻易成就了它！\BibleRef{欧 2:1} 但如果犹太人要说，\Quote{从依撒格生的，才称为你的后裔}这句话的意思是指那些由依撒格所生的人算为他的后裔，那么厄东人和所有那些人也应被称为他的儿子，因为他们的祖先厄撒乌也是他的儿子。但现在他们远未被称儿子，反而是最大的外人。可见，肉身所生的子女不是天主的子女，即使在自然本身，那从上面藉着洗礼而产生的出生也早已预表了。如果你对我提到子宫，我反过来对你提到水。但正如在这种情况下一切都是出于圣神，同样在另一种情况下一切都是出于应许。因为由于不孕和年老，子宫比任何水都要冷。因此，让我们准确认识我们自己的高贵身份，并活出与之相称的生命。因为这生命中没有任何肉身的或属地的因素；所以我们里面也不应有这些。因为生我们的，不是睡眠，也不是肉身的意愿\BibleRef{若 1:13}，也不是拥抱，也不是欲望的疯狂，而是\BibleQuote{天主对人的爱}成就了这一切。\BibleRef{铎 3:5} 正如在依撒格的情形中，年龄已无希望；同样，在我们身上，当罪恶的老年已笼罩我们时，依撒格突然在青春中迸发，我们都成了天主的子女和亚巴郎的后裔。\BibleRef{依 40:31}\par

第10节：\BibleQuote{不但如此，而且黎贝加也是这样：她怀了孕，就是从我们的祖宗依撒格一个人怀的。}\BibleRef{罗 9:10}\par

所讨论的主题十分重要。因此，他转而使用多个论证，并设法用各种方法解决这个难题。因为，如果他们在接受如此伟大的许诺之后被抛弃，这件事既奇怪又新鲜，那么，我们这些从未期望过此类事情的人竟然能进入他们的美物之中，就更奇怪了。这就像一位国王的儿子，他曾被许诺要继承他现有的权力，却被贬到声名狼藉的人的行列中，而一个被定罪的、背负无数罪恶的人，从监狱中被提出后，却得到了本当属于另一人的权力。因为他的意思是：你要说什么？儿子不配吗？好吧，但这个人也不配，而且更不配。因此，他要么应该与前者一同受罚，要么应该与他一同受尊荣。现在，类似的事情就发生在犹太人和外邦人身上，或者比这更奇怪得多。现在，所有人都配不上，他已在上面显示过了，他说：\BibleQuote{因为所有人都犯了罪，亏缺了天主的光荣。}\BibleRef{罗 3:23} 但新的事情是，当所有人都配不上时，外邦人却唯独得到了救恩。除此之外，他说，还有另一个难题有人可能提出：如果天主无意向他们实现许诺，为何还要许诺呢？因为那些不知道未来并且屡次受骗的人，也会向不配得的人许诺他们会得到赏赐。但那位事先知道未来如同知道现在一样，并且清楚知道他们会使自己不配得到许诺，因此不会得到任何指定事物——他为何要许诺呢？那么，保禄应对这一切的方式是什么？那就是指出他许诺的对象是哪一种以色列。因为这一点一旦被指出，就同时证明了所有的许诺都已实现。为了指明这一点，他说：\BibleQuote{凡属于以色列的，并不都是以色列人。}\BibleRef{罗 9:6} 正因如此，他没有使用雅各伯的名字，而是用了以色列，以色列是那位义人的德能、来自高天的恩赐以及见过天主的标志。\BibleRef{创 32:28} 然而，他说：\BibleQuote{众人都犯了罪，亏缺了天主的光荣。}\BibleRef{罗 3:23} 如果众人都犯了罪，为何有些人得救，有些人灭亡？这是因为并非所有人都愿意来到他面前，因为就他而言，众人都得救了，因为众人都被召叫。然而，他暂时还没有摆出这一点，而是从一个有利的位置，从其他例子中，通过提出另一个问题来应对它，并且像前一个案例那样，用另一个很大的困难来应对困难。因为当他讨论如何因着基督的成义而让其余的人享有那义时，他引用了亚当的例子，说：\BibleQuote{如果因一人的过犯，死亡就因那一人作了王，那么那些丰富地蒙受了恩宠和正义恩惠的人，更要藉着耶稣基督一人在生命中为王了。}\BibleRef{罗 5:17} 实际上，他并没有澄清亚当的例子，而是藉此澄清他的（或他自己的）例子，并表明那为他们而死的人按自己的意志掌管他们是更合理的。因为一人犯罪众人受罚，对大多数人来说似乎并不十分合理。但一人行义众人成义，则既更合理又更合乎天主。然而，他仍然没有解决他所提出的困难。因为那一点越是隐晦，犹太人就越被堵住口。而他立场的困难就转移到了另一点上，这一点也因此变得更清晰了（\EditorialNote{马尔谷福音抄本和四份手稿作\Quote{比那点}}）。同样，在这段经文中，他也是通过提出其他困难来应对所提出的问题，因为他是在与犹太人争辩。因此，他并不费心去解决他向我们提出的例子。因为他不需要为它们负责，就像在与犹太人争辩时一样。但他从这些例子中使自己的主题始终更加清晰。他的意思是：为什么你们感到惊讶，此时有些犹太人得救了，有些却未得救？为什么在古时，在先祖的时代，就可以看到这种事情发生？为什么依撒格单单被称为后裔，而他也是依市玛耳和其他几个人的父亲？\Quote{但他是女奴所生的。} 这与他的父亲有什么关系？我仍然不会吹毛求疵。让这个儿子因他的母亲而被放在一边吧。我们该如何谈论刻突辣所生的那些儿子？难道他们不是自由的，并且生于自由的母亲吗？为什么他们没有蒙受同样的偏爱如依撒格一样？我为什么说这些？因为黎贝加甚至是依撒格唯一的妻子，她生了两个孩子，都是为依撒格所生；然而这些出生的人，虽然有同样的父亲和母亲，并且是同一劳苦的果实，都出于一父一母，而且还是双胞胎，却没有享受同样的命运。然而在这里，你们没有母亲的奴隶身份来解释它，如同在依市玛耳的情况一样，你们也不能说一个出于这个子宫而另一个出于不同的子宫，如同刻突辣和撒辣的情况，因为在这种情况下，他们出生时有共同的时辰。这就是为什么保禄当时为了给出一个更清晰的例子，说这件事不仅发生在依撒格身上，\BibleQuote{而且黎贝加也从一人，即我们的先祖依撒格怀孕。}\BibleRef{罗 9:10}\par

第11至13节。\BibleQuote{双胞胎还没有出生，也没有行善或作恶，为要使天主按拣选所定的旨意得以成立，不是出于行为，而是出于召叫人的那一位，就有人对她说：长子要服务于幼子。正如经上所记载：\InnerQuote{我喜爱雅各伯，而憎恨了厄撒乌。}\BibleRef{拉 1:2-3}}\BibleRef{罗 9:11}\par

是什么原因导致一个被爱，另一个被恨？为什么一个服事，另一个被服事？是因为一个邪恶，另一个良善。然而在孩子们还未出生时，一个已被尊荣，另一个已被定罪。因为他们尚未出生时，天主曾说：\BibleQuote{年长的要服事年幼的}\BibleRef{创 25:23}。天主说这话的意图是什么？因为祂不像人一样等着从行为的结果看出谁是善的、谁不是，甚至在行为之前，祂就知道谁是邪恶的、谁不是。这件事在以色列子民身上也发生了，且以更奇妙的方式。他说，我何必提\OriginalTerm{厄撒乌}{Esau}和\OriginalTerm{雅各伯}{Jacob}呢？他们中一个邪恶，另一个良善？因为在以色列子民的事上，罪属于所有人，因为他们都崇拜了牛犊。然而，有些人蒙受了仁慈，有些人却没有。\par

第15节。\BibleQuote{因为我要对谁显示\OriginalTerm{仁慈}{mercy}，祂说，就对谁显示\OriginalTerm{仁慈}{mercy}；我要对谁显示\OriginalTerm{怜悯}{compassion}，就对谁显示\OriginalTerm{怜悯}{compassion}。}\BibleRef{罗 9:15}\BibleRef{出 33:19}\par

这件事也可以在那些受惩罚的人身上看到，因为对于受惩罚且付出了如此沉重代价的\OriginalTerm{法郎}{Pharaoh}，你会说什么呢？你说他心硬且悖逆。难道只有他如此，而没有其他人吗？那么他为什么受到如此严厉的惩罚？为什么甚至在\OriginalTerm{犹太人}{Jews}的事上，祂称那不是子民的为子民，或者，为什么不把所有人都视为配得同等的尊荣？因为（经上说）：\BibleQuote{如果他们像海沙那样多，但只有\OriginalTerm{遗民}{remnant}会得救。}\BibleRef{依 10:22} 为什么只能是一群\OriginalTerm{遗民}{remnant}？你看他让这个话题充满了多大的困难。而且非常恰当。因为当你有能力使你的对手陷入困惑时，不要立刻提出答案，因为如果他发现自己也对同样的无知负有责任，为什么要将不必要的危险揽到自己身上？为什么通过把一切揽到自己身上使他更加大胆？现在告诉我，你这\OriginalTerm{犹太人}{Jews}，你有这么多令人困惑的问题，却无法回答任何一个，你怎么能在\OriginalTerm{外邦人}{Gentiles}的蒙召之事上打扰我们？然而，我有一个很好的理由可以告诉你，为什么\OriginalTerm{外邦人}{Gentiles}成了义而你被抛弃。理由是什么？就是他们由于\OriginalTerm{信德}{faith}，你由于\OriginalTerm{法律的行为}{works of the Law}。正因为你的这种固执，你才在各方面（旁注和一些手抄本都）被遗弃了。因为他们不认识天主的\OriginalTerm{正义}{righteousness}，却企图建立自己的\OriginalTerm{正义}{righteousness}，而没有服从天主的\OriginalTerm{正义}{righteousness}。\BibleRef{罗 10:3} 这段经文的整个澄清，总结地说，就是由那位蒙福的人在此阐明。但为使这点更清楚，让我们逐一探究他所说的话；要知道，蒙福的\Person{保禄}所致力于的是，通过他所说的一切表明，只有天主知道谁是配得的，没有任何人知道，即使他自以为非常了解；但在他的判决中，存在各种偏差。因为那洞察人心的那位，只有祂确切地知道谁该得冠冕，谁该受惩罚和报应。因此，许多被人类视为善的人，祂定罪并惩罚；而那些被怀疑为恶的人，祂却在证明并非如此之后给予冠冕；这样，祂的判决不是依据我们这些奴隶的判断，而是依据祂自己敏锐而公正的决断，并且不等行为的结果来从中看出谁是恶的、谁不是。但为使我们不使主题更加模糊，让我们再次回到宗徒的原话。\par

第10节。\BibleQuote{不仅如此，而且当\OriginalTerm{黎贝加}{Rebecca}由一人怀了孕。}\BibleRef{罗 9:10}\par

\EditorialNote{他（保禄）暗示，}我本可以再提到\Person{刻突辣}的子孙，但我没有。我这样做，是为了从一个优势位置取得胜利，我所提出的是那些出于同一父亲、也是同一母亲的子孙。因为他们都是\Person{黎贝加}所生，来自\Person{依撒格}——他才是真生的、蒙拣选的、备受尊崇的儿子，关于他天主说过：\BibleQuote{在\Person{依撒格}内，你的后裔将被称呼}，他成了\BibleQuote{我们众人的父亲}；但如果他是我们父亲，那么他的儿子本应是我们的父亲；然而事实并非如此。你看，这事不仅发生在\Person{亚巴郎}身上，也发生在他自己的儿子身上，并且表明在一切情况下，信德和德行才是显著的，并赋予真正的亲属关系以特征。因为由此我们得知，儿女被称为他的儿女，不仅是由于出生方式，更是因为他们配得上父亲的德行。因为如果仅仅是由于出生方式，那么\Person{厄撒乌}也应该享有与\Person{雅各伯}相同的待遇。因为他也是从几乎已死的母腹中出生，他的母亲本是不生育的。然而，仅有这一点还不够，还需要有品行，这一事实为我们提供了不少实际的教训。他没有说一个是好人，另一个是坏人，因此前者得了尊荣；以免这种论据被人用来攻击他：\Quote{怎么，外邦人中的好人比受过割损的人中的好人更好吗？}因为即使事实如此，他还没有这样说，因为那会显得令人不快。而是他把一切全寄托于天主的预知，而这是没有人敢反驳的，即使他如何疯狂。\BibleQuote{因为孩子们还没有出生}，他说，\BibleQuote{她就得到指示：大的要服事小的。}他指出，出于肉体的高贵出身毫无益处，我们必须追求灵魂的德行，这在尚未有行为之前，天主就已经知道了。因为他说：\BibleQuote{孩子们还没有出生，也没有行善或作恶，为了使天主按照拣选的旨意得以坚定，她就得到指示：大的要服事小的；}这就是预知的标记，即他们从出生就被拣选了。为要使那按预知所作的拣选，明显是从天主而来，从第一天起，天主就立刻看到并宣告了谁是善的，谁不是。所以不要告诉我，你读了法律和先知，并服事了这么久（他意指）。因为那能试探灵魂的主，知道谁配得拯救。因此要顺服于拣选之不可测度性。因为只有祂知道如何正确地给予冠冕。例如，有多少人看起来比\Person{圣玛窦}更好——按当时可见的行为表现。但那知道未显明之事、并能考验心灵资质的主，知道那珍珠虽在污泥中，仍拣选了它；在越过其他人后，对这高贵美善感到喜悦，祂拣选了它，并加上来自祂自己的恩宠于那高贵的自由意志，使它成为可悦纳的。因为在那些可朽坏的艺术中，甚至在其他事物中，好的判断者不以无知者形成判断的理由来选用眼前的材料；而是从他们自己清楚知道的要点出发，他们常常贬低无知者所赞赏的，并决定贬低什么：相马师也常常如此对待马匹，宝石鉴赏家和其他手艺工人也是这样。那么，更爱人的天主、无限智慧、独有万物清晰知识的主，更不会容许人的猜测，而是凭祂自己准确无误的智慧，对众人作出判决。因此，祂拣选了税吏、盗贼和娼妓，却羞辱了司祭、长老和官长，把他们赶出去。这一点在殉道者身上也可以看到。因此，许多被完全抛弃的人，在考验的时期却得了冠冕。另一方面，一些被许多人视为伟大的人却绊倒失足。所以，不要质问造物主，不要说：为什么这个人得了冠冕，那个人受了惩罚？因为祂知道如何准确完成这些事。因此祂也说：\BibleQuote{我爱了\Person{雅各伯}，我恨了\Person{厄撒乌}。}你固然从结果知道这是公正的；但祂自己在结果之前就已经清楚知道了。因为天主所寻求的，不只是行为的展示，还有高贵的抉择和\OriginalTerm{顺从的性情}{\GreekText{γνώμην} \GreekText{εὐγνώμονα}}。因为这样的人，如果因一时意外而犯罪，必会迅速恢复自己。即使他偶然停留在恶习中，也不会被忽视，因为那知道万事的天主必会迅速把他拉出来。相反，在这种性情上败坏的，即使他似乎行了些善事，也会灭亡，因为他怀着恶意去做。因此，连\Person{达味}在犯了凶杀和奸淫之后，由于他是被意外所冲昏，并非惯行邪恶，便迅速洗涤了罪过。但是那法利塞人，虽然没有犯过这样的罪（\BibleRef{路 18:11}），甚至还有善行可夸，却因他所选择的坏精神而失掉了一切。\par

第14节。\BibleQuote{那么，我们说什么呢？难道天主有不义吗？绝对不可。}\BibleRef{罗 9:14}\par

因此，在我们和犹太人之间没有这种事。然后他继续讲另一件事，比这个更清楚。那是怎样的事呢？\par

第15节。因为他对\Person{梅瑟}说：\BibleQuote{我要恩待谁，就恩待谁；我要怜悯谁，就怜悯谁。}\BibleRef{罗 9:15}\par

这里他又通过将异议一分为二并加以应对，同时又提出另一个新的难题，来加强异议的力量。但为了使我所说的更清楚，有必要解释一下。他的意思是：\Person{天主}在分娩之前就说\BibleQuote{年长的要服侍年幼的}。那又怎样？\Quote{天主不义吗？}绝不是。现在请听下文所述。因为在那情况下，德行或恶行或许是决定性的因素。但在这里有一件所有\Person{犹太人}都参与的罪，即金牛犊的罪，仍然有一些人受了惩罚，一些人没有受惩罚。这就是为什么祂说：\BibleQuote{我要怜悯谁就怜悯谁，我要恩待谁就恩待谁。}\BibleRef{出 33:19} 他的意思是：梅瑟啊，要知道谁配得我对人的爱，这不是你的事，而是把这事交给我。但若\Person{梅瑟}无权知道，我们更无权知道。这就是为什么他不只是引用这段经文，还提醒我们这是对谁说的。他的意思是：他是在对\Person{梅瑟}说话，至少借着说话者的威严，使反对者知所收敛。在解决了所提出的难题之后，他将问题一分为二，又提出了另一个异议，如下：\par

第16至17节：\BibleQuote{如此看来，这不在于那愿意的，也不在于那奔跑的，而在于发仁慈的\Person{天主}。因为\WorkTitle{圣经}对\Person{法郎}说：\InnerQuote{我将你兴起来，正是为此目的，为要在你身上显示我的大能，并使我名传遍天下。}}\par

他的意思是：正如在那一种情况下，一些人得救，一些人受罚，这里也是如此。这个人正是为此目的而被保留的。然后他又再次提出异议。\par

第18至19节：\BibleQuote{因此，祂愿意怜悯谁，就怜悯谁；愿意使谁硬心，就使谁硬心。这样，你必要对我说：那么祂为什么还指责人呢？因为谁能抗拒祂的旨意呢？}\par

看他如何费尽心思从各方面使问题变得棘手。他没有立即给出答案，因为不立即给出是有益的，而是先止住辩驳者的口，如下所说：\par

第20节：\BibleQuote{然而，人啊！你是谁，竟敢向\Person{天主}抗辩？}\BibleRef{罗 9:20}\par

他这样做是为了抑制反对者不合时宜的追根究底和过分的好奇心，加以约束，并教他认识\Person{天主}是什么，人是什么，祂的预知是多么不可理解，远超我们的理性，以及在一切事上顺从祂是多么必要。所以，当他这样在听众心中预备了基础，并使他们的精神平静柔和之后，便极其恰当地引出答案，使他的话易于被接受。他并没有说回答这类问题是不可能的，而是说（五个抄本：不，而是说什么？说）提出这些问题本身就是僭妄。因为我们的本分是服从\Person{天主}的作为，即使不知道其中的缘由，也不应好奇。因此他说：\BibleQuote{你是谁，竟敢向\Person{天主}抗辩？}你看他把对方看得多么轻，如何压下他高傲的气焰！\BibleQuote{你是谁？}你与祂的能力有分吗？\EditorialNote{比较约伯传38章} 不，你是在审判\Person{天主}吗？比起祂，你连存在都不能！不是这样或那样的存在，而是完全没有！因为\BibleQuote{你是谁}这个说法比\BibleQuote{你什么都不是}更加藐视对方。他又以其他方式进一步表示质问中的愤怒，他说的不是\BibleQuote{你是谁，敢回答\Person{天主}？}而是\BibleQuote{敢向\Person{天主}抗辩}，即反驳和对抗。因为说事情应该这样、不应该那样，正是\BibleQuote{抗辩}的人所做的。看他如何惊吓他们，如何恐吓他们，如何使他们战栗，而不是质疑和好奇。这正是优秀教师所做的：他不处处迎合门徒的喜好，而是引导他们接受自己的心意，拔除荆棘，然后播下种子，而且并不立即回答所有向他提出的问题。\par

第20至21节：\BibleQuote{受造之物岂能对造它的说：\InnerQuote{你为什么这样造了我？}难道陶工没有用同一团泥，作成贵重的和卑贱的器皿的权柄吗？}\BibleRef{耶 18:1}\par

他这样说，并非要废除自由意志，而是要表明我们应当在何种程度上服从天主。因为就向天主质问而言，我们应当像泥土一样少有这种倾向。我们不仅应当戒绝反驳或质询，甚至应当戒绝谈论或思想这类事，而要成为那无生命的物质，它随陶匠的手转动，任他随意牵引。而他所用的比喻也只适用于这一点，即不是要阐述生活规则的任何宣告，而是要我们完全服从和保持静默。我们在所有情况下都应当注意这一点：不要完全照搬比喻，而要选取其中的益处，即它们被引入的目的，其余部分则不予理会。例如，当经上说：\BibleQuote{他蹲伏、躺卧如雄狮}\BibleRef{户 24:9}，我们只取其不可征服和可畏的一面，并非取兽性，也不是狮子所具有的任何其他特性。再者，当主说：\BibleQuote{我要像失仔的熊一样迎击他们}\BibleRef{欧 13:8}，我们只取其报复的一面。当经上说：\BibleQuote{我们的天主是吞噬的烈火}\BibleRef{希 12:29}，我们只取其惩罚时的毁灭力。同样，在这里我们也必须单独取出泥土、陶匠和器皿。当经上接着说：\BibleQuote{陶匠难道没有权柄，用同一团泥，造一个器皿作为贵重的，另一个作为卑贱的吗？}，不要以为这是保禄在讲论受造物的本质，或暗示意志的必然性，而是为了说明天主的宰制和不同安排的差异；因为如果我们不这样理解，就会产生各种矛盾。因为如果在此他是在论及意志，以及善与不善之人，那么天主便是这些人的制造者，而人将不负任何责任。如此，保禄也将显得自相矛盾，因为他始终最推崇自由选择。因此，他在这里别无他意，只是要劝说听者完全顺服天主，绝不可在任何事上向天主质问。因为他说，陶匠用同一团泥随己意制作，无人禁止；同样，当天主从同一人类中惩罚一些人，荣耀另一些人时，你不要对此好奇或干预，只需朝拜，并效法泥土。正如泥土随陶匠的手转动，你也当随那如此安排者的心意。因为祂所行没有一件是随意或偶然的，尽管你不懂祂智慧的秘密。然而，你允许陶匠用同一团泥制作各种器皿，并不责难；却对天主追究祂惩罚和荣耀的理由，不让祂知道谁配得荣耀、谁不配；但因为同一团泥具有同一本质，你就断言它们有着同样的倾向。这是何等荒唐！而且，器皿的尊贵与卑贱甚至不取决于陶匠，而取决于使用它们的人之手；同样，在这里也取决于自由选择。然而，如我先前所说，必须只取这个比喻的一个含义，即人不可违抗天主，而应顺服祂不可理解的智慧。因为例子应当大于主题和它们被引用来说明的事物，以便更好地吸引听者。如果它们不是更大，且没有远超其上，他就不能恰当攻击并羞辱反对者。然而，他以适当的优越性制止了他们不合时宜的固执。然后他引入了他的回答。那么，这回答是什么呢？\par

\BibleQuote{第22、23、24节。如果天主，愿意显示祂的义怒，并彰显祂的能力，以极大的宽容，容忍了那些注定灭亡的、可怒的器皿；并且，为要将祂荣耀的丰富，彰显在那蒙怜悯的、早预备得荣耀的器皿上，就是我们所蒙选的，不但从犹太人中，也从外邦人中。}\par

他的意思大致如下。法郎是一个义怒之器，即一个因自己的刚硬之心激起了天主义怒的人。因为在享受了极大的含忍之后，他并未变好，反而毫无改善。因此，他不仅称他为\BibleQuote{义怒之器}，也称他为\BibleQuote{注定灭亡的}。即，完全注定，但出于其自身。因为天主既没有遗漏任何可能使他回头的事，他也没有遗漏任何会毁灭他、使他无可宽恕的事。尽管如此，天主明知这一切，却\BibleQuote{以极大的含忍容忍了他}，愿意引导他悔改。因为如果祂不愿意这样，就不会如此含忍。但他不愿利用含忍来悔改，反而完全使自己成为义怒之器，天主便利用他来纠正其他人，通过对他的惩罚使他们变好，并以此显示祂的权能。因为天主不愿意以这种方式彰显祂的权能，而是以另一种方式，即通过祂的恩惠和仁慈，祂以上所有方式都表明了这一点。因为如果保禄不希望用这种方式显示能力（\BibleQuote{不是要我们显得合格}，他说，\BibleQuote{而是要你们做诚实的事} \BibleRef{格后13:7}），何况天主呢。但在祂显示了含忍，为引导他悔改之后，他却不悔改，祂仍容忍了他很长时间，为同时显示祂的良善和权能，即使那个人无意从这极大的含忍中获得任何益处。正如通过惩罚这个持续不改的人，祂显示了祂的权能，同样通过怜悯那些犯了许多罪但悔改了的人，祂显明了祂对人的爱。但此处并没有说对人的爱，而是说光荣，为显示这尤其是天主的荣耀，且为此祂最是恳切。但在说\BibleQuote{祂预先预备进入光荣的}时，他并非意指一切都是天主的工作。因为若是这样，就没什么能阻止所有人得救。但他再次显明祂的预知，并消除了犹太人与外邦人之间的区别。在这个话题上，他再次为他的陈述提供了一个不小的辩护。因为不仅是在犹太人中有一些人丧亡、一些人得救，在外邦人中也是如此。因此他没有说所有的外邦人，而是说\Quote{外邦人}，也没有说所有的犹太人，而是说\Quote{犹太人}。正如法郎因自己的无法无天成为义怒之器，这些人也因自己的服从之准备成为怜悯之器。因为虽然大部分是天主的工作，但他们自己也贡献了一点点。因此他也没有说善行之器或胆量之器(\OriginalTerm{胆量}{\GreekText{παρρησίας}})，而是说\BibleQuote{怜悯之器}，为显示一切都是出于天主。因为\BibleQuote{不在乎那愿意的，也不在乎那奔跑的}这句话，即使出现在反驳中，但若出自保禄之口，也不会造成困难。因为当他说\BibleQuote{不在乎那愿意的，也不在乎那奔跑的}时，他并未剥夺我们的自由意志，而是表明一切并非出于自己，因为需要来自上天的恩宠。因为我们有责任愿意和奔跑：但不要信赖自己的劳苦，而要信赖天主对人的爱。他在别处也表达了这一点。他说过：\BibleQuote{不是我，而是与我同在的恩宠} \BibleRef{格前15:10}。他说得好：\BibleQuote{祂预先预备进入光荣的}。因为既然他们责备这些人，说他们是因恩宠得救，并以为能让他们羞愧，他远远超过了这种暗示。因为如果这事甚至给天主带来了光荣，那么给那些通过他们天主受到光荣的人带来的荣耀就更多了。但请注意他的克制和不可言喻的智慧。因为当他本可以提出作为受罚例子的不是法郎，而是犹太人中的罪人，从而使他的论述更加清晰，并表明在同一祖先、同一罪恶的情况下，有些人丧亡，有些人蒙受了怜悯，并说服他们不要心存疑虑，即使有些外邦人得救，而犹太人丧亡；为了不使他的论述令人厌烦，他从外邦人那里引出惩罚的显明，以免被迫称他们为\BibleQuote{义怒之器}。但那些获得怜悯的人他从犹太人中引出。此外，他也充分为天主辩护，因为祂虽然明知这民族正使自己成为注定灭亡的器皿，但祂仍贡献了祂所有的一切：祂的忍耐、祂的含忍，而且不仅是含忍，而是\BibleQuote{极大的含忍}；然而他并不愿意只是针对犹太人陈述这些。那么，为什么有些人是义怒之器，有些人是怜悯之器？出于他们自己的自由抉择。然而，天主极其良善，对两者都显示同样的仁慈。因为不仅仅是得救的人蒙受了怜悯，法郎也蒙受了，就天主这方面而言。因为同样的含忍，他们和他都受益。如果他没有得救，那完全是出于他自己的意愿：因为就天主而言，祂为他所做的，与为得救之人所做的同样多。然后，在给问题一个基于事实的回答之后，为了使他的论述得到其他证词的支持，他引用了先知们从前所作的相同宣言。他说，欧瑟亚古时曾写下如下的话：\par

25节。\BibleQuote{我要叫那非我民的，成为我的人民；那非蒙爱怜的，成为蒙爱怜的。} \BibleRef{罗9:25}\par

这里，为阻止他们说：你是在用似是而非的推理欺骗我们，他引欧瑟亚为证，他呼喊说：\BibleQuote{我要叫那非我民的，成为我的人民；} \BibleRef{欧2:23} 那么，谁是非民？显然，是外邦人。谁是非蒙爱怜的？同样是他们。然而，他说，他们将立刻成为人民、蒙爱怜的和天主的子女。\par

26节。\BibleQuote{他们将要被称为，}他说，\BibleQuote{永生天主的子女。} \BibleRef{罗9:26}\par

但若他们坚持说这是指那些相信的犹太人说的，那么论点仍然成立。因为若那些在领受了如此多的恩惠后仍然心硬和疏离、失去了作为选民资格的犹太人，能发生如此巨大的改变，那么有什么能阻止那些并非疏离后归向祂、原本是外邦人的人被召叫，并且，只要他们服从，被视为堪得同样的福分呢？于是，处理完\Person{欧瑟亚}后，他并未止步于此，又接着引用了\Person{依撒意亚}，与他和谐一致地发言。\par

\Quote{因为\Person{依撒意亚}，}他说，\Quote{关于\Place{以色列}大声呼号着。}\BibleRef{罗 9:27}\par

这就是说，他大胆宣告，毫不掩饰。那么，当他们在先前以超过号角的响亮宣告了同样的事时，为何要指控我们呢？依撒意亚呼喊什么？\BibleQuote{虽然以色列子民的人数如海沙，但剩余的必得救。}\BibleRef{依 10:22}\par

你看见吗？他也并没有说所有人都要得救，而是说那堪当的人要得救。因为他认为，我并不在乎群众，他意思是说，一个如此广布的种族也不会使我困扰，我只拯救那些使自己堪当的人。他提到\Quote{海沙}也不是没有理由的，而是为了提醒他们那古老的许诺，他们已使自己不配继承这许诺。那么，当所有先知都表明并非所有人都要得救时，你为什么因许诺似乎落空而烦恼呢？接着他也提到了得救的方式。请看先知的准确性和宗徒的判断，他引用了多么贴切的证言！这不仅向我们表明得救的是部分人而非全部，而且还补充了他们得救的方式。那么，他们如何得救，天主又如何使他们堪得这恩惠呢？\par

\BibleQuote{祂要完成祂的工作，并迅速行义，}他说，\BibleQuote{因为主要在地上实行短暂的工作。}\BibleRef{罗 9:28}\BibleRef{依 10:23}\par

他所说的意思是这样的。无须绕圈子，无须麻烦和因法律工作而烦恼，因为救恩是藉着一条非常短的道路。因为信德就是如此，它以简短的话语包含救恩。\BibleQuote{因为如果你口里承认主耶稣，心里相信天主从死者中复活了祂，你必得救。}\BibleRef{罗 10:9}现在你明白这\Quote{主在地上要作简短的话}是什么意思了。更奇妙的是，这句简短的话不仅带来救恩，也带来义德。\par

\BibleQuote{正如\Person{依撒意亚}先前所说：若非万军的上主给我们留下后裔，我们早已像\Place{索多玛}，与\Place{哥摩拉}相似了。}\BibleRef{罗 9:29}\BibleRef{依 1:9}\par

这里他又显示了另一件事，即连这少数人也不是凭自己的资源得救的。因为他们也本会灭亡，遭遇索多玛的命运，即他们将不得不遭受彻底的毁灭（因为索多玛的人也被连根铲除，甚至连最小的残余也没有留下），他意思是说，他们也会像这些人一样，除非天主以极大的仁慈对待他们，并藉信德拯救了他们。这也在可见的被掳事件中发生了，大多数被掳走灭亡，只有少数人被救。\par

\BibleQuote{那么，我们该说什么呢？那未追求义德的外邦人，却得到了义德，即基于信德的义德。但以色列追求法律的义德，却没有达到法律的义德。}\BibleRef{罗 9:30-31}\par

这里终于有了最清楚的答案。因为他既已从事实（\Quote{因为凡以色列的并不都是以色列人}）以及族长\Person{雅各伯}和\Person{厄撒乌}，并从先知\Person{欧瑟亚}和\Person{依撒意亚}的事例中使用了证明，他在增添了困惑之后，进一步给出了最决定的答案。那么，讨论的问题有两个：一是外邦人得到了，二是他们在没有追求的情况下得到了，即没有为此费心。而在犹太人的情况下，也有两个同样的难点：一是以色列没有得到，二是虽然他们费了心，却没有得到。为此，他的用词也更加强烈。因为他没有说他们获得了，而是说他们\Quote{获得了义德}。因为特别新奇的是，那些追求的人没有得到，而那些没有追求的人却得到了。他似乎是在纵容他们，说\Quote{追求}。但随后他给予了重击。因为他有强有力的答案要给他们，所以他不怕使异议变得稍加尖锐。因此他既不谈论信德，也不谈论随之而来的义德，而是表明甚至在信德之前，在他们的立场上，他们就已败诉并被定罪。因为，你，犹太人，他说，你甚至连法律下的义德都没有找到。因为你触犯了它，成了诅咒的对象。但这些不是通过法律，而是通过另一条路来的人，找到了比这更大的义德，即信德的义德。他先前也说过：\Quote{因为如果\Person{亚巴郎}因行为成义，他就有可夸之处，但在天主面前却没有}\EditorialNote{罗马书第4章}：如此表明另一种义德比这更大。之前我说有两个难点，但现在它们甚至变成了三个问题：外邦人找到了义德，且在没有追求的情况下找到，并且找到了比法律更大的义德。相同的难点在犹太人的情况下又有相反的景象。以色列没有找到，尽管费了心却没有找到，甚至没有找到较小的。既然使听众陷入了困惑，他就开始给出一个简洁的答案，并告诉他所说的一切的缘由。那么，缘由是什么呢？\par

\BibleQuote{因为他们不凭信德，而凭法律的行为去追求。}\BibleRef{罗 9:32}\par

这是全段中最清晰的回答。如果他一开始就说出这话，就不容易被人接纳。但既然是在许多困惑、预备和论证之后，并用了无数准备步骤才写下来，他终于使它更易理解，也更容易被接受。因为他说，这就是他们灭亡的原因：\BibleQuote{因为他们不是凭信德，而是可以说是凭法律的行为}而想成义。他并没有说\BibleQuote{凭行为}，而是说\BibleQuote{可以说是凭法律的行为}，为表明他们连这种义德也没有。\par

\BibleQuote{因为他们正跌在那绊脚石上。}\par

33节. \BibleQuote{正如经上所载：看，我在熙雍安放了一块绊脚石，一块使人跌倒的磐石；凡相信祂的，必不蒙羞。}\BibleRef{罗 9:33}\par

你们又看到，勇敢来自信德，恩赐是普遍的；因为这话不仅是对犹太人说的，也是对整个人类说的。他要说，每一个人，无论是犹太人、希腊人、西古提人、色雷斯人，或其他任何人，只要相信，就享有极大勇敢的特恩。先知奇妙之处在于，他不仅预言了他们会相信，也预言了他们不会相信。因为跌倒就是不信。正如在前一段中，他指出了那些丧亡的人和得救的人，那里说：\BibleQuote{以色列子民的数目若如海沙，只有残余的将得救。并且，万军的上主若不给我们留下后裔，我们早就像索多玛了。}以及\Quote{祂所召的不只是犹太人，也是外邦人}；同样，这里也暗示有人会相信，有人会跌倒。但跌倒来自于不留心、贪恋其他事物。既然他们留心于法律，反而跌在这石头上，\Quote{一块绊脚石和使人跌倒的磐石}他如此称呼它，这是根据那些不信的人的特性与结局。\par

我所讲的话对你们够明白了吗？还是仍然不够清楚？我认为，对一直留心听的人，容易看明白。但若有人忘记了，可以私下见面，学习所讲的内容。正因如此，我在这解释部分讲得较长，以免中断上下文的连贯，从而破坏论述的清晰。为此，我也将在此结束讲道，不再像往常那样对你们讲一些更直接实践的道理，以免因讲得太多而在你们的记忆中造成新的模糊。现在是时候进入适当的结尾，以对万有之天主的光荣颂来结束讲道。让我们，说话的我与听讲的你们，一同停住，将光荣归于祂。因为国度、权柄、荣耀都属于祂，直到永远。阿们。\par

\SourceRecord{Translated by J. Walker, J. Sheppard and H. Browne, and revised by George B. Stevens.；Nicene and Post-Nicene Fathers, First Series；Vol. 11.；Edited by Philip Schaff.；Buffalo, NY: Christian Literature Publishing Co.,；1889.；<http://www.newadvent.org/fathers/210216.htm>.}

% structure: p0001 source=h1 target=section reason=page-title

\section{罗马书讲道 第十七篇}

\BibleRef{罗 10:1}\par

\Quote{弟兄们，我心中所渴望的和向天主为他们所求的，就是使他们得救。}\par

他现在又要更严厉地责备他们了。因此，他再次消除一切仇恨的猜疑，并事先努力纠正误解。他说道：\Quote{那么，不要在意言语或指责，而要观察我说这话并非出于敌对的心态。因为同一个人既渴望他们得救，不只是渴望，甚至还为此祈祷，却又恨他们、厌恶他们，这是不可能的。}因为在这里，他称自己的极度渴望和祈祷（\OriginalTerm{心愿}{\GreekText{εὐδοκίαν}}）为\Quote{心愿}。他所极力主张并祈求的，不仅是从惩罚中得到释放，更是使他们得救。不仅如此，从随后的话也表明他对他们的善意。因为从他所能及的，他尽力而为，并努力寻找哪怕一丝一毫的借口为他们开脱。但他无能为力，被事实的本质所压倒。\par

\BibleRef{罗 10:2} \Quote{我为他们作证}，他说，\Quote{他们对天主有热心，但不是按着智识。}\par

这不应该是宽恕而不是控告他们的理由吗？因为如果他们分离不是出于人的原因，而是出于热心，那他们就应当被怜悯而不是被惩罚。但请注意，他是如何巧妙地在言辞中偏袒他们，同时又显露出他们不合时宜的固执。\par

\BibleRef{罗 10:3} \Quote{因为他们不认识}，他说，\Quote{天主的义德。}\par

这话再次可能导致宽恕。但随后的话却导致更强烈的控告，并且如此消除了任何辩护。\par

\Quote{而且四处奔走}，他说，\Quote{想要建立自己的义德，却不顺服于天主的义德。}\par

他说这些话是为了表明，他们犯错是由于任性妄为和好权，而非由于无知，并且他们甚至连这出自法律行为的义德也没有建立。\BibleRef{玛 21:38} 因为说\Quote{四处奔走要建立}正表明这一点。实际上，他并没有明说（因为他没有说他们在两种义德上都欠缺），但他以一种非常明智且合宜的智慧暗示了这一点。因为如果他们仍然\Quote{四处奔走}要建立那义德，很明显他们还没有建立它。如果他们不曾顺服于这义德，那么他们在这义德上也是欠缺的。但他称它为\Quote{自己的义德}，或是因为法律已不再有效，或是因为它是劳苦的义德。而他称这出自信德的义德为天主的义德，因为它完全来自上天的恩宠，并且人在此情况下成义不是靠劳苦，而是靠天主的恩赐。但那些一直抵制圣神，并且执意要靠法律成义的人，没有归信。既然他们没有归信，也没有领受随之而来的义德，又无法靠法律成义，他们就被剥夺了一切资源。\par

\BibleRef{罗 10:4} \Quote{因为基督是法律的终末，使凡信的人都获得义德。}\par

请看保禄的智慧。因为他既谈到一种义德，又谈到另一种义德，唯恐那些信主的犹太人似乎拥有其中一种，却被排除在另一种之外，并受到无法无天的指控（因为就连这些人，由于刚刚信主，也并非没有同样的忧虑），又唯恐犹太人再次期望达到另一种，并说：\Quote{虽然我们现在尚未履行，但我们将来必定会履行。}看他采取了什么立场。他指出只有一种义德，它完全成就于此（即信德），而凡将此义德（即因信德而来的）归于己身的人，也成就了那另一种。但拒绝这一种的人，在那另一种上也同样欠缺。因为如果基督是\Quote{法律的终末}，那么没有基督的人，即使看似拥有那义德，实际上也没有。但拥有基督的人，即使没有完全履行法律，也领受了全部。因为医生的技艺的终末是健康。因此，能够使人痊愈的人，即使没有医生的技艺，也拥有一切；但不懂医治的人，即使看似是技艺的追随者，也欠缺一切：法律和信德的情况也是如此。拥有信德的人同样拥有了法律的终末，但没有信德的人则与两者都无缘。因为法律的目标是什么？是使人成义。但它没有能力，因为没有人履行它。因此，这就是法律的终末，它始终指向这一点，为此制定了其所有部分，包括节期、诫命、祭献以及其它一切，好使人成义。但基督通过信德使这个终末得到更圆满的实现。所以，他说，不要以为归信信德就是违犯法律。因为当你为法律而不信基督时，你才真正违犯法律。如果你相信他，那么你也履行了法律，甚至超出了它的要求。因为你领受了更大的义德。接下来，既然这是一个论断，他又从圣经中提出证据。\par

\BibleRef{罗 10:5} \Quote{因为梅瑟}，他说，\Quote{描述了出自法律的义德。}\par

他的意思是这样的：梅瑟向我们展示了出自法律的义德是什么样的，以及从哪里来的。那么它是什么样的，又包含什么？在于履行诫命。他（按：修订本作\Quote{那人}）说：\Quote{遵守这些事的人，必因此活着。} \BibleRef{肋 18:5} 除了完全履行法律，没有其它方法在法律中成为义人。但这对任何人来说都不可能，因此这种义德使他们失败了（\OriginalTerm{失败}{\GreekText{διαπέπτωκεν}}）。但请你告诉我们，保禄，还有另一种义德，即出自恩宠的义德。它又是什么，包含什么？请听他用什么话来清楚地描绘它。因为在他驳斥了那一种之后，接着就谈到这一种，并说：\par

第6、7、8、9节。\Quote{但出于信德的义德这样讲：\InnerQuote{你心里不要说：\InnerQuote{谁能升到天上去呢？}}（意思是把基督从天上带下来），或\InnerQuote{谁能下到深渊里呢？}（意思是把基督从死者中带上来）。而是说了什么？\InnerQuote{圣言离你很近，就在你口中，就在你心里}，即我们所宣讲的信德之言。因为你若口里承认主耶稣，心里相信天主使他从死者中复活了，你便得救。}\par

为防止犹太人这样说：\Quote{那些没有找到较小义德的人，如何找到较大的呢？}他给出一个无可辩驳的理由：这条路比那条更容易。因为那条要求履行一切（当你做了全部，你才得生存），但出于信德的义德并不这样说，而是说什么呢？\par

然后，为不使我们因显得轻易廉价而轻视它，注意他是如何扩展描述的。因为他没有立即说出刚才的话，而是说什么呢？\BibleQuote{但出于信德的义德是这样讲的：\InnerQuote{你心里不要说：谁能升到天上去？}（意思是把基督带下来）；或\InnerQuote{谁能下到深渊里？}（意思是把基督从死者中带上来）。}因为在行为表现的德行中，有一种懒惰与之对抗，放松我们的劳作，这需要一个非常警醒的灵魂才不屈服于此；同样，当要求人相信时，有些理性思考会扰乱并破坏大多数人的心智，这需要一个有些活力的灵魂才能彻底摆脱。这正是他为什么把同样的事摆在前面的原因。正如他对待亚巴郎那样，他这里也一样。在那处，他显示亚巴郎是因信德而成义，免得他好像因偶然获得如此大的冠冕，仿佛无足轻重，为了颂扬信德的本性，他说：\BibleQuote{他在绝望中仍怀着希望而相信了，为成为众多民族之父。他虽已年老，却不软弱于信德，反而视自己的身体已死，撒辣的胎也已死去；对于天主的恩许，他没有因不信而动摇，反而因信德而坚强，归光荣于天主，且坚信天主所应许的，他也有能力实现。}\BibleRef{罗 4:18}这样他显示了需要活力，一个高贵的灵魂，能容纳超出期望的事，不为表象所绊倒。他这里也同样做，显示了需要智慧的心，一个属灵（希腊：触天）而伟大的灵魂。他不仅说：\Quote{不要说，}而且\Quote{不要在心里说，}即连怀疑并与自己说\Quote{这怎能可能？}的念头也不要有。你看，这是信德的主要特征：抛下此世的一切推论，追求超乎自然的事，排除计算的虚弱，完全接受天主的大能。然而，犹太人不仅主张这一点，而且说不可能因信德而成义。但他自己却把已发生的事转到另一解释上，为了显示这事如此伟大，即使发生后仍需要信德，他才能有充分理由将冠冕给予这些事；并且他使用在旧约中出现的词句，总是努力完全摆脱喜好新奇和反对圣经的指控。因为他这里论信德所说的，梅瑟对他们论诫命时也说过，以此表明他们从天主手中得到极大的恩惠。他的意思是：不需要上到天上或跨越大海然后领受诫命，但如此伟大庄严的事，天主已使我们容易获得。那么，\Quote{圣言离你很近}是什么意思？就是容易。因为在你的心思和你的口中就是你的救恩。不需要长途跋涉，不需要航海，不需要翻山越岭而得救。你甚至不必跨过门槛，就可以坐在家中得救。因为\InnerQuote{在你们的口中，在你们的心中}就是救恩的源泉。然后从另一角度他也使信德之言显得容易，说：\Quote{天主使他从死者中复活了。}你只需思考行这事的配得者，你就不再看到任何困难。那么，他是主，从复活即可知道。他在书信开头就说了：\BibleQuote{按圣德的神，藉着从死者中复活，以大能显明是天主子。}\BibleRef{罗 1:4}但复活也是容易的，甚至对那些非常不信的人，也从行这事的权能中显明了。既然义德更大，而且轻省容易得到，那么放弃轻省的而去追求不可能的事，岂不是极端的争论吗？因为他们不能说他们避之惟恐不及。那么看，他是如何剥夺他们一切的借口。如果谁选择负担沉重且不能实行的事，而错过了轻省且能拯救他们、给予法律不能给予的，他们还有什么辩护词可说？这一切只能来自与天主作对的争论精神。因为法律是沉重的\OriginalTerm{沉重的}{\GreekText{ἐπαχθής}}，而恩宠是轻省的。法律，尽管他们极力争议，却无法拯救；恩宠却显出自己的义德，也给出了法律的义德。那么还有什么借口能拯救他们？因为他们对恩宠争论不休，却毫无意义地紧抓法律不放。然后，因为他做出了有力的断言，又用圣经来证实。\par

第11-13节。\Quote{因为圣经说：}他继续说：\BibleQuote{凡信他的人，必不受羞辱。因为犹太人与希腊人之间没有区别：因为同一的主对一切呼求他的人都是富足的。因为凡呼求主名的人，必得救。}\BibleRef{罗 10:11}\par

你看他如何引出见证，无论是对信德，还是对它的宣认。因为\Quote{凡信的人}这句话指出信德，而\Quote{凡呼求的}这句话表明宣认。然后，为宣扬恩宠的普遍性，并压低他们的夸耀，他将先前详细证明的，在这里简要地提醒他们，再次显示犹太人与未受割损者之间没有区别。他说：\Quote{因为}他说，\Quote{犹太人与希腊人之间没有区别。}他论及圣父时所说的话，在论证这一点时，他在这里关于圣子说了同样的话。因为正如他以前在肯定这一点时说过：\BibleQuote{祂只是犹太人的天主吗？不也是外邦人的吗？是的，也是外邦人的，因为天主只有一位。}\BibleRef{罗 3:29} ——同样，他在这里也说：\BibleQuote{因为同一主对众人都是富足的，厚待一切（并加给一切）。}\BibleRef{罗 3:22} 你看他如何将祂描述为极其渴望我们的救恩，因为祂甚至将此视为自己的财富；因此他们现在也不应当绝望，或以为只要他们悔改，就不可赦免。因为那将拯救我们视为自己财富的，不会停止富有。因为就连这一点也是财富，即恩赐倾注于众人。因为最令他苦恼的是，那些享有超越全世界特权的人，如今却因信德从这些宝座上被贬抑，丝毫未比他人优越，于是他不断引述先知们预言他们将与他们享有同等的尊荣。他说：\Quote{因为凡信祂的人，}他说，\Quote{必不致蒙羞}\BibleRef{依 28:16}；以及\BibleQuote{凡呼求主名的人，必得救}\BibleRef{岳 2:32}。而\Quote{凡}一词用在所有情况下，使他们无话可答。但没有什么比虚荣更糟糕的了。因为正是这一点，尤其这一点，导致了他们的毁灭。因此基督也对他们说：\BibleQuote{你们彼此接受光荣，却不寻求那唯独来自天主的荣耀，怎么能相信呢？}\BibleRef{若 5:44} 这一点，除了毁灭，还使人遭受许多嘲笑，并在另一世界的惩罚之前，就已在此世使他们陷入无数的灾祸。如果你们愿意清楚了解这一点，暂时撇开它使我们失去的天堂和它把我们投入的地狱，让我们就眼前的事调查整个问题。那么还有什么比这更浪费？还有什么更可耻、更令人反感？因为这种混乱是一种浪费，从那些毫无目的地花费于戏院、赛马及其他无关开支的人身上可以明显看出；从那些建造精美昂贵的房屋，并以无用的奢华风格布置一切的人身上也可以看出，这些我不必在此详述。但一个人染上这种病必然变得奢侈、浪费、贪婪、吝啬，任何人都能看到。为了给那野兽食物，他伸手攫取他人的财物。我为什么谈论财物呢？这火不仅吞噬金钱，也吞噬灵魂，不仅在此世造成死亡，也在来世造成死亡。因为虚荣是地狱之母，大大点燃了那火和毒虫。可以看到它甚至对死者也有权势。还有什么比这更糟？因为其他情欲被死亡终止，但这情欲即使在死后也展现其力量，并试图在死尸中显示其本性。因为当人们在临终时下令为他们建造华丽的纪念碑，这将耗尽他们所有财物，并预先为葬礼安排巨大的奢华，而在生前却侮辱前来乞求一分钱和一块面包的穷人，但死后却为虫子摆设丰盛的宴席，为什么还要寻找更过分的对疾病的奴役呢？从这祸患中也产生了不正当的爱情。因为有许多人，并非因为外表的美丽，也非因为与她同床的欲望，而是因为想夸耀\Quote{我征服了某某人}，甚至被拖入通奸。我为什么还要提及由此产生的其他祸患呢？因为我宁愿长期（三份手稿\OriginalTerm{连续地}{\GreekText{διηνεκώς}}）成为一万个野蛮人的奴隶，也不愿做一次虚荣的奴隶。因为就连他们也不会加给俘虏这样的命令，如同这恶习加于其崇拜者身上的那样。因为它说：你要做每个人的奴隶，无论他比你高贵还是卑微。藐视你的灵魂，忽视德行，嘲笑自由，牺牲你的救恩，如果你做任何善事，不要为悦乐天主而做，而要为了向众人炫耀，为这些事你甚至会失去你的冠冕。如果你施舍或禁食，承受痛苦，但要注意失去收获。还有什么比这些命令更残忍？因此嫉妒掌权，因此傲慢，因此贪婪——万恶之母。因为成群的家仆、身穿金色制服的奴仆、随从、谄媚者、银饰的战车，以及比这些更大的荒谬事物，并非为了任何快乐或需要，而仅仅是为了虚荣。是的，有人会说，这种痛苦是邪恶的，任何人都能看到；但如何完全避免它，这才是你应当告诉我们的。那么，首先，如果你说服自己这种混乱是有害的，你就为纠正它开了一个很好的头。因为当一个人生病时，他很快会派人去请医生，如果他首先意识到自己生病的话。但如果你寻求除此以外的另一条路以逃离此处，就不断仰望天主，并满足于来自祂的荣耀；如果你发现情欲在挑逗你，激动你向同伴讲述你的善行，那么接下来要想到，告诉之后你一无所获。熄灭这荒谬的欲望，并对你的灵魂说：看，你长久以来内心充满自己的善行想要告诉他人，却没有勇气将其保留给自己，反而向所有人泄露了。那么你从中得到了什么好处？一点也没有，只有极大的损失，并失去了所有辛劳积攒的成果。除此以外，还要考虑另一件事，即大多数人的意见是扭曲的，不仅扭曲，而且很快就消逝。\par

因为即使他们一时钦佩你，当那机会过去后，他们就会忘记一切，并夺去天主已赐给你的冠冕，且不能为你确保来自他们自己的冠冕。然而，如果这些是持久的，那么用那持久的交换这不持久的将是一件最可悲的事。但当这不持久也过去时，我们能为出卖那持久的而为了那不持久的，为失去如此祝福而为了博得少数人的好评，作何辩护呢？而且确实，即使那些赞美者众多，即使为此他们也值得可怜，而且越多人这样做，就越可怜。但若你对我所说感到惊讶，请听基督如此判决：\BibleQuote{祸哉，当众人都赞美你们时}\BibleRef{路 6:26}。确实如此。因为如果在每一种技艺中，你寻求其中的工匠（\OriginalTerm{工匠}{\GreekText{δημιουργους}}）作为评判者，你怎能将美德的检验托付给众人，而非最托付给那位比任何人更确切了解它、且最能鼓掌并加冕它的那一位呢？那么，让我们把这句铭刻在我们的墙壁、门户和脑海中，并不断对自己说：\Quote{祸哉，当众人都赞美我们时。}因为那些如此说话的人，后来也会诽谤你，说你是一个虚荣的人、爱荣誉的人、贪图他们好话的人。但天主并非如此。但当祂看见你渴慕来自祂的荣耀时，那时祂会最赞美你，并尊重（\OriginalTerm{尊重}{\GreekText{θαυμάσεται}}，多数抄本省略）你，并宣布你为得胜者。人却不然；但当他发现你是奴性而非自由的，常以空洞言语的虚假赞美满足你时，他就从你那里夺走了你真正的奖赏，使你比买来的奴隶更为卑贱。因为那些人得到他们的奴隶遵守他们的命令，但你甚至没有命令就使自己成为奴隶。因为你甚至不等听到他们的任何话，只要你知晓如何在何事上能取悦他们，即使没有他们的命令，你也做一切。那么，我们若取悦恶人，并在他们命令之前就服事他们，而却不肯听从天主，即使祂每天命令并劝告我们，我们岂不该受何等的地狱呢？然而，你若贪图荣耀和赞美，就要避开来自人的赞美，这样你才会得到荣耀。远离谄媚的言语，这样你才会从天主和人那里得到无数的赞美。因为没有人比那位轻视荣耀的人更受我们荣耀，或比那位认为被尊重和被赞美是无关紧要的人更受我们赞美和尊重。如果我们这样做，宇宙的天主更会如此。当祂荣耀你并赞美你时，有谁更加可称为有福呢？因为荣耀与耻辱之间的差别，没有比来自天上的荣耀与来自人的荣耀之间的差别更大。或者更确切地说，还有更大的差别，乃是无限的区别。因为如果这（属人的荣耀），即使不与其他放在一起，也是卑贱而不美的，当我们用另一者来审视它时，试想它将是何等卑贱！因为如同娼妓站在自己的位置，向任何人出卖自己，同样，那些做虚荣奴隶的人也是如此。或者更确切地说，这些人比那娼妓更卑贱。因为那种女人在许多情况下还轻视那些迷恋她们的人。但你对每个人出卖自己，不论是逃奴、盗贼还是扒手（因为由这些人及类似的人组成了那些为你鼓掌的剧场），而那些你作为个人视作无价值的人，当他们聚集成群时，你尊敬他们胜过你的救恩，并显明自己比他们中任何一个都更不值得尊敬。因为你怎么可能不是更不值得呢？当你需要别人的好话，并以为凭自己不够，除非你接受来自别人的荣耀？你难道没有觉察到，除了我所说的，还有：既然你是众人注目的对象，为人所熟知，如果你犯了错误，你会有数不尽的控告者；但如果你不为人知，你会处于安全中？是的，有人可能会说，但如果我做得对，我将有数不尽的崇拜者。可怕的是，不仅当你犯罪时，而且当你做对时，虚荣之乱仍对你造成伤害，在前一情形中颠覆了千万人，在后一情形中完全剥夺了你的赏报。因此，即使在世俗事务中，热衷于荣耀也是一件可悲的事，且充满各种耻辱。但当你在属灵事务中也有同样的处境时，你还有什么借口留下呢？既然你不愿意像你从仆人那里得到尊荣那样，给予天主同样多的尊荣？因为奴隶\BibleQuote{仰望主人的眼}\BibleRef{咏 122:2}，雇工仰望他的雇主（那要付他工资的人），门徒仰望他的师傅。但你却反其道而行。你离开了那雇佣你的天主，即你的主，却仰望你的同仆；而你知道天主在今生之后仍记念你的善行，但人只在当下记念。当你已有天堂的观众聚集时，你却在地上聚集观众。而摔跤手在哪里奋力拼搏，他就在哪里希望得尊重；但你，你的摔跤是在上面，却焦虑要在下面获得冠冕。还有什么比这样的疯狂更糟糕呢？但如果我们认为合适，也看看那些冠冕。一种是由傲慢形成，第二种是由嫉妒他人形成，第三种是由虚伪和谄媚形成，另一种又由财富形成，还有一种由奴性的奉承形成。如同孩子们在孩童游戏中彼此戴上草冠，且常背后嘲笑那被冠冕的人；同样，现在那些在你身上颂扬的人，也常私下嘲笑他们把草冠加在我们头上。但愿只是草！但现在那冠冕满载许多祸害，并毁坏我们所有的善行。因此，考虑到它的卑劣，逃避它所带来的损害。因为，你希望有多少人赞美你？一百人？或两倍、三倍、四倍多人？或者，如果你愿意，算做十倍、二十倍，并有两千、四千，甚至一万人为你欢呼。然而，这些人并不比众多从上空呱呱叫的乌鸦更好。或者，再考虑到天使的集会，这些人会比蠕虫更加卑劣，他们的好话也不比蜘蛛网、烟雾或梦境更坚实。\par

那么请听，深知这一切的\Person{保禄}，他非但不追求这些，甚至加以拒绝，他说：\BibleQuote{我绝不夸耀，只夸耀我们的主耶稣基督的十字架。}\BibleRef{迦 6:14} 那么你也要渴望这种光荣，免得激怒上主，因为这样做，你不仅侮辱自己，更是侮辱天主。因为即使你是一位画家，有一个门徒，他若省略向你展示他的绘画实践，却公开向任何愿意观赏的人展示他的画，你也不会安静地接受。但如果连对你的同仆这样做都是侮辱，更何况是对上主呢！但如果你愿意从其他方面学会轻视此事，你就该怀抱崇高的心胸，嘲笑外表，增加对真正光荣的热爱，充满属灵的精神，如同\Person{保禄}那样对你的灵魂说：\BibleQuote{你们不知道我们将来要审判天使吗？}\BibleRef{格前 6:3} 借此唤醒它，然后继续责备它说：\Quote{你这位审判天使的人，怎能容许自己被渣滓审判，与舞蹈者、模仿者、角斗士和御马者一同受称赞呢？} 因为这些人都追求这样的喝彩。但你要将你的翅膀高举在这喧嚣之上，效法那位旷野的居民——\Person{若翰}，学习他是如何超脱于众人之上，不回头看奉承者；当他看见全\Place{巴勒斯坦}的居民涌到他周围，惊叹诧异时，他并没有被这样的尊荣冲昏头脑，反而起来谴责他们，向这大群会众如同对一个人一样训话，责斥他们说：\BibleQuote{毒蛇的种类！}\BibleRef{玛 3:7} 然而他们正是为他而聚集，离开城市，为要见到那位圣者，但这些事丝毫没有使他软弱。因为他远在光荣之上，摆脱了一切虚荣。同样，\Person{斯德望}看见同一群人非但不尊敬他，反而对他疯狂，咬牙切齿，他高居他们的愤怒之上，说：\BibleQuote{你们这些执拗和心耳未受割损的人！}\BibleRef{宗 7:51} 同样，\Person{厄里亚}当那些军队、君王和全体百姓在场时，说：\BibleQuote{你们摇摆不定，模棱两可，要到几时呢？}\BibleRef{列上 18:21} 但我们却奉承所有人，讨好所有人，用这种卑躬屈膝的谄媚换取他们的尊敬。因此一切都颠倒了，为了这份恩宠，基督信仰的事业被出卖，一切因大众的看法而被忽视。那么让我们摒弃这种激情，这样我们就会对自由、港湾和宁静有正确的认识。因为虚荣的人总是像风暴中的人，颤抖、恐惧、服侍千万个主人。但摆脱这种束缚的人，就像在港湾中的人，享受纯洁的自由。那人却不是这样，他有多少熟人，就有多少主人，被迫成为所有人的奴隶。那么我们如何摆脱这种严酷的束缚呢？就是通过对另一种光荣——真正的光荣——产生爱慕。因为如同那些迷恋容貌的人，看到一个更英俊的人就会从第一个转移注意力；同样，那些追求来自人的光荣的人，如果天上的光荣照耀他们，就能将他们从此吸引开。那么让我们注视这天上的光荣，并完全认识它，这样，因着欣赏它的美丽，我们就能避开另一者的丑陋，并因持续享受它而获得许多喜乐。愿我们所有人都能藉着天主的恩宠和祂对人的慈爱达到这一点。等等。\par

\SourceRecord{Translated by J. Walker, J. Sheppard and H. Browne, and revised by George B. Stevens.；Nicene and Post-Nicene Fathers, First Series；Vol. 11.；Edited by Philip Schaff.；Buffalo, NY: Christian Literature Publishing Co.,；1889.；<http://www.newadvent.org/fathers/210217.htm>.}

% structure: p0001 source=h1 target=section reason=page-title

\section{罗马书第十八篇讲道}

\BibleRef{罗 10:14,15}\par

\Quote{那么，他们怎能呼求他们所未曾信的人呢？他们怎能信他们所未曾听见的人呢？没有宣讲者，他们怎能听见呢？除非被派遣，他们怎能宣讲呢？正如经上所载。}\par

这里他再次夺去他们的一切借口。因为他既说过：\Quote{我为他们作证，他们对天主有热心，但不是按着真知，}和\Quote{他们因为不认识天主的义，而没有服从它，}他接着指出，连这无知本身，他们在天主面前也是可罚的。他确实没有这样说，但他以提问的方式展开论述，从而更清楚地定了他们的罪，将整段经文编排成质问与回答。但请回顾更早之处。先知说：\Quote{凡呼求主的名的人，必获得救恩。}现在或许有人说：但他们怎能呼求他们所不信的人呢？然后他提出一个问题：他们为什么不信？接着又是一个异议：人当然可以说，他们怎能相信，既然他们没有听见？但他暗示他们确实听见了。然后又是一个异议：\Quote{没有宣讲者，他们怎能听见呢？}接着是回答：他们确实宣讲了，并且有许多人被派遣正是为此。如何证明这些人是被派遣的呢？接着他引述先知的话：\Quote{那传布和平福音、传报美善佳音者的脚，多么美丽啊！}\BibleRef{依 3:7}你看，他如何藉宣讲的方式指出宣讲者。因为这些人到处所传讲的，无非是那不可言喻的美善，以及天主与人类所缔造的和平。因此，不信的原因并非我们，而是\Person{依撒意亚}先知，他多年前就预言我们要被派遣、宣讲并说出我们所说的话。既然如此，得救来自呼求他，呼求来自相信，相信来自听闻，听闻来自宣讲，宣讲来自被派遣；而他们确实被派遣、宣讲了，先知也陪同他们，指出并宣告他们，说这些就是他们许多世纪前所预指的人，甚至连他们的脚也因宣讲的内容而受赞美；那么，显然不信只是他们自己的过错，因为天主的部分已经圆满完成。\par

第16至17节。\Quote{但他们并没有都服从福音。因为依撒意亚说：主啊，有谁相信了我们的报道呢？由此可见，信德出于听闻，听闻出于天主的话。}\BibleRef{依 53:1}\par

既然他们又用另一个异议来质问他：如果这些人是天主派遣的使者，那么所有人都应该听从他们。请看\Person{保禄}的智慧，看他如何表明正是这个造成困惑的事，反而消除了困惑和窘迫。他会说：犹太人啊，在如此大量而充分的证据和证明之后，是什么使你不安呢？是并非所有人都服从了福音吗？正是这件事，与其他证据一同，足以向你证实我所说的是真理，即有些人并不相信。因为连这也被先知预言了。还要注意他那难以言表的智慧：他如何表明确实超出了他们所期望或预料的回答。因为你的话是什么意思？他问道。是不是所有人都没有相信福音？是的！\Person{依撒意亚}从古时就预言了这一点。或者说，不仅如此，甚至远不止于此。因为你抱怨的正是为什么不是所有人都相信？但\Person{依撒意亚}走得更远。因为他说了什么？\Quote{主啊，有谁相信了我们的报道呢？}（\OriginalTerm{听闻}{\GreekText{ἀκοή}}）？他巧妙地抓住了这个引文，作为他所说话的证明：\Quote{因此，信德出于听闻}（\OriginalTerm{听闻}{\GreekText{ἀκοή}}）。他并非仅仅陈述这一点。但由于犹太人一直在寻求标记和复活的情景，对此极其渴慕，他说：然而先知并没有应许任何这类事，而是说我们应当藉着听闻而相信。因此他首先证实这一点，说：\Quote{因此，信德出于听闻。}然后，由于这话似乎有些卑微，请看他是如何提升它的。因为他说：我所说的不是普通的听闻，也不是需要听人的话而相信，而是指一种伟大的听闻。因为那听闻是\Quote{出于天主的话。}他们并非讲说自己的话，而是传述从天主所学来的。这比奇迹更高。因为无论天主说话还是行奇迹，我们都同样有义务相信和服从。因为事工和奇迹都出自他的话语。因为连天和一切也都是如此建立的。\BibleRef{咏 32:6}在表明我们应当相信总是讲说天主话语的先知，而不必寻求更多之后，他接着提出我所说的异议，说道：\par

第18节。\Quote{但我问：他们难道没有听见吗？}\BibleRef{罗 10:18}\par

他的意思是：如果宣讲者被派遣，并宣讲了他们奉命要讲的，而这些人没有听见，那又怎样？随后，他对这个异议给出了一个最完美的答复。\par

\Quote{诚然，他们的声音传遍了大地，他们的言语传到了地极。}\par

你说什么？他的意思是。他们没有听见吗？须知整个世界，甚至地极，都听见了。而你们，这些传报者居住如此之久、且出自你们土地的人，却没有听见？这怎么可能？倘若地极都听见了，你们更应当听见。接着又是另一个反驳。\par

\newline{}\BibleRef{罗 10:19}。\BibleQuote{但我问：以色列人不知道吗？}\par

因为即使他们听见了，他的意思是，却不知道所说的是什么，也不明白这些就是被派遣的人？他们因无知而应被宽恕吗？绝不然。因为依撒意亚已用这些话描述了他们的特征：\BibleQuote{那传报和平福音的人的脚步是多么美丽啊！}\newline{}\BibleRef{依 52:7}。而在他之前，立法者本人也已如此说。于是他继续讲。\par

\BibleQuote{首先梅瑟说：\BibleQuote{我要以那不成民族的外邦人，激怒你们；以一个愚昧的民族，惹起你们的怒气。}\newline{}}\newline{}\BibleRef{申 32:21}\par

所以，他们甚至本应从他就能分辨出传报者，不仅是从这些不信的事实，不仅是从他们传报和平的事实，不仅是从他们带来那些美好事物的喜讯的事实，也不仅是从圣言播撒在世界各地的事实，而且还从他们看见那些比自己卑微的人、即外邦人获得更大光荣的事实本身。因为那些他们从未听闻、他们的祖先也从未听闻的事，这些外邦人却突然\OriginalTerm{追求智慧}{\GreekText{ἑ} \GreekText{φιλοσόφουν}}。这是一种如此强烈的光荣标志，足以刺激他们，引发他们的嫉妒，并使他们想起梅瑟的预言，其中说：\BibleQuote{我要以那不成民族的外邦人，激怒你们。}因为不仅是光荣的伟大足以使他们陷入嫉妒，而且还有一个事实：一个如此微不足道、几乎不被视为民族的民族，却得以享受这些事物。\BibleQuote{因为我要以那不成民族的外邦人，激怒你们；以一个愚昧的民族，惹起你们的怒气。}因为还有什么比希腊人（外邦人，见第373、377页）更愚昧呢？又有什么比他们更微不足道呢？请看天主如何从古时就用各种方式预示了这些时代的迹象和清楚标记，以消除他们的盲目。因为这事并非发生在某个小角落，而是在陆地、海洋和地球的每一角落。他们看见那些从前为他们所鄙视的人，如今享受无数的祝福。因此，人们应当认为，这个民族就是梅瑟所说的：\BibleQuote{我要以那不成民族的外邦人，激怒你们；以一个愚昧的民族，惹起你们的怒气。}难道只有梅瑟这样说吗？不，因为依撒意亚在他之后也这样说。这就是为什么保禄说\Quote{首先梅瑟}，以表明将有第二位以更清晰、更明白的方式说同样的事。正如他上面说依撒意亚呼喊，这里也一样。\par

\newline{}\BibleRef{罗 10:20}。\BibleQuote{但依撒意亚放胆地说。}\par

他的意思是这样的：他强迫自己，渴望说话，不是含糊其辞，而是赤裸裸地摆在你们眼前，宁愿因直言而冒险（奥利金\Emph{在本段注释}），也不顾自己的安全，为的是不给你们的顽固留下任何借口。虽然预言的风格并非如此清楚地说出这些，但为了彻底堵住你们的嘴，他清晰而明确地预言了一切。一切！什么一切？就是你们被舍弃，以及他们被接纳，他如此说：\BibleQuote{我没有去寻找我的，我让他们找到了我；我没有求问我的，我向他们显现。}\newline{}\BibleRef{依 65:1}。那么，那些没有寻找的是谁？那些没有求问的是谁？显然不是犹太人，而是那些外邦人，他们至今不认识他。如同梅瑟用\Quote{不成民族}和\Quote{愚昧的民族}来给他们的特征作标记，同样，这里他也用同样的基础来指出他们，即他们的极端无知。这对犹太人是一个极大的责备：那些没有寻找他的反倒找到了他，而那些寻找他的却失去了他。\par

\newline{}\BibleRef{罗 10:21}。\BibleQuote{至于以色列，他说：\BibleQuote{我整天向那悖逆和抗辩的百姓伸出我的手。}\newline{}}\newline{}\BibleRef{依 65:2}\par

如今请注意，那使许多人成为疑问的难题，早已藏在先知的话中，并且连同明确的解答也一并显明了。这难题是什么呢？你们先前听见\Person{保罗}说：\Quote{那么，我们说什么呢？那未追求义的外邦人却得了义。但以色列人追求义的法律，反而未能达到义的法律。}（\BibleRef{罗 9:30}）\Person{依撒意亚}在这里也说了同样的话。因为他说：\Quote{我没有寻求我的，我让他们寻见了；我没有求问我的，我向他们显现了}（\BibleRef{罗 10:20}），这与\Quote{那未追求义的外邦人却得了义}是同一意思。然后，为了表明所发生的事不仅出于天主的恩宠，也出于那些归向祂之人的心性，正如其他人的被弃也出于那些悖逆者的争辩心，请听他接着说的话：\Quote{但关于以色列人，祂却说：我整日向一个悖逆和抗辩的民族伸出手。}（\BibleRef{罗 10:21}）这里以\Quote{整日}指整个前代时期。而\Quote{伸出手}意思乃是呼唤、吸引他们，并邀请他们。然后，为了表明全是他们自己的过错，祂说：\Quote{向一个悖逆和抗辩的民族。}你看，这是对他们多大的指控！因为即使祂邀请他们，他们也不听从，反而抗辩祂，而且他们看到祂如此行，不是一次、两次或三次，而是整个时期。但其他从未认识祂的人，却有力量将祂吸引到自己身边。不过，他不是说他们自己有能力这样做，而是为了消除外邦人中的高傲想象，并表明这一切都是祂的恩宠运作，祂说：\Quote{我显现了，我被寻见了。}或许有人会说：那么他们完全没有功劳吗？绝不是的，因为他们自己贡献了寻找那已经发现的，以及认识那已向他们显现的。然后，为防止这些人说：\Quote{但为什么你不也向我们显现呢？}祂又加上了比这更甚的事实：我不但显现了，还整日伸着手呼唤他们，显出一个慈父和爱子母亲的全部关怀。请看，他如何通过对先前提问的解答，最清晰地给我们指明了，是由于他们自己的心性而招致了毁灭，并且他们完全不配得到宽恕。因为他们虽然既听见又明白了所说的话，却仍不愿意归向祂。而且更甚的是，祂不仅使他们听见并明白了这些事，还加上了另一件更能激动他们、吸引他们的事，就是当他们悖逆和抗辩时，祂加上了这一点。那么这是什么呢？那就是激怒他们，使他们嫉妒。因为你们知道这情欲的支配力量，以及嫉妒天生具有何等大的能力，足以平息一切争辩，并唤醒那些懈怠的人。何必说人，就是在无理性的畜生和未成年的孩子身上，它所显示的力量也是如此之大呢？因为一个孩子常常在父亲呼唤时不愿服从，而是固执己见。但当他看到另一个孩子受到注意时，那么他即使没有被呼唤，也会跑到父亲的怀中；呼唤做不到的事，嫉妒的刺激却能做到。天主也这样做了。因为祂不仅呼唤、伸出了手，还通过把远不如他们的人带进（这最使人极其嫉妒的事）他们的好东西中，而且（这是更强烈的步骤，使嫉妒更盛）带进了比他们更好、更必要、他们连做梦也未曾想过的极大好东西中，以此来激起他们的嫉妒之情。然而他们仍不屈服。那么，他们如此顽固，还配得什么宽恕呢？一点也没有。但他自己并不说出这一点，而是让听众从他所陈述的结论中去领会，随后又凭着他惯常的智慧，通过接着说的内容来证实这一点。他在上面也这样做了，既在律法（\BibleRef{罗 7:7}）上，也在百姓的事上，引入了比实际情况更严重的异议；然后在回答中，为了推翻这个异议，他酌情让步，正如情况所允许的那样，以使他的话不令人反感。他现在也正是这样做的，写下了以下内容：\par

第十一章第一节。\Quote{那么，我说：难道天主弃绝了祂所预知的百姓吗？绝对不是。}（\BibleRef{罗 11:1}）\par

并且他引入了人们存疑时所使用的表达形式，仿佛是从前面所说的话中取得机会，在作出这一令人警觉的声明之后，通过否认它，使后续内容得以顺利地被人接受；他以前所有论点所努力证明的，在这里也加以证实。那么这是什么？即使只有少数人得救，但许诺仍然有效。这就是为什么他不仅仅是说\Quote{百姓}，而是说\Quote{祂所预知的百姓}。然后继续证明\Quote{百姓}没有被弃绝，说：\Quote{因为我也是一名以色列人，出自\Person{亚巴郎}的后裔，本雅明支派。}（\BibleRef{罗 11:1}）\par

他说：\Quote{我是教导者、宣讲者。}既然这与之前所说的似乎矛盾，即\BibleQuote{谁肯相信我们的报道？}以及\BibleQuote{我整天向悖逆、顶嘴的百姓伸出我的手}，还有\BibleQuote{我要借那不成子民的，惹动你们的妒恨}，他并不满足于反对之言，也不满足于说了\Quote{决不如此}，而是重新拾起话题加以证实，说：\BibleQuote{天主并没有弃绝自己的百姓。}但有人会说，这不是证实，而是断言。那么请看证实，首先是第一点，其次是接下来的一点。第一点是他本人就是那民族的一份子。如果天主正要将他们弃绝，便不会从他们中间拣选那位托付了所有宣讲、世务、一切奥秘及整个救恩计划的人。这便是一个证明。接下来的第二点是他说：\BibleQuote{他所预知的百姓}，即他清楚知道哪些人是相称的，会接受信德。\BibleRef{宗 2:41} 因为已有三千、五千甚至一万人相信了。因此，为了防止有人说：\Quote{那么你们就是百姓吗？因你被召叫，整个民族就被召叫了吗？}他继续往下说。\par

Ver. 2. \BibleQuote{天主并没有弃绝自己预先所知道的百姓。}\BibleRef{罗 11:2}\par

仿佛他说：\Quote{我有三千、五千或一万人。}那么怎么样？百姓就变成了三千、五千或一万人吗？那与天上星辰、海边沙粒一样多的后裔？你就是这样欺骗我们，愚弄我们吗？你把自己和跟随你的少数人当作整个百姓，并让我们抱有空洞希望，说应许已经实现，而所有人都灭亡了，救恩只降临到少数人身上？这全是虚张声势！我们无法容忍这样的诡辩！为避免他们这样说，请看他在下文如何答复，虽然没有提出异议，却在回答问题之前，先以古事为依据。那么答案是什么呢？\par

Ver. 2-5. 他说：\Quote{难道你们不知道，}他说，\BibleQuote{圣经关于厄里亚说了什么？他怎样（大多数抄本；萨伏伊抄本作\Quote{谁}）向天主控告以色列说：\InnerQuote{主啊，他们杀了你的先知，拆了你的祭台，只剩下我一个人，他们还要寻索我的命。}但天主回复他说的是什么？\InnerQuote{我为自己保留了七千人，是未曾向巴耳屈膝的。}同样，现在这时刻也按恩宠的拣选有了残余。}\BibleRef{罗 11:2-5}\par

他所说的大致是这样的。\Quote{天主并没有弃绝他的子民。}因为如果他这样做了，就不会接纳他们中的任何一个。但如果他确实接纳了一些，他就没有弃绝他们。但仍有人说，如果他未曾弃绝，他就会接纳所有的人。这并不成立，因为在\Person{厄里亚}的时代，得救的部分只剩下\Quote{七千人}，现在也可能有许多人相信。但你若不知道他们是谁，这并不奇怪，因为那位先知，如此伟大而善良的人，也不知道。但天主自己安排了事情，即使先知也不知道。但请考虑他的判断。现在，在证明他面前的事时，他巧妙地加重了对他们的指控。正是为此，他引用了整段经文，为了在他们面前展示他们的乖戾，并表明他们自古以来就是如此。因为如果他不想这样，而是全神贯注地证明民众只剩下少数，他就会说，即使在\Person{厄里亚}的时代，也留下了七千人。但现在他向它们读了更早的段落，似乎一直煞费苦心地表明，他们对待\Person{基督}和宗徒们所做之事并不奇怪，而是他们一贯的做法。为了防止他们说我们杀死\Person{基督}是因为他是个骗子，我们迫害宗徒是因为他们是冒牌货，他引用了经文说：\BibleQuote{主，他们杀了你的先知，拆毁了你的祭台。}\BibleRef{列上 19:14}然后，为了使他的言论不刺激他们，他为引用经文附加了另一个理由。因为他引用经文并非有意指责他们，而是好像要表明其他事情。他让他们在已经发生的事上毫无借口。因为请注意，这指控甚至从说话者本人那里是多么强烈。因为既不是\Person{保罗}，也不是\Person{伯多禄}，也不是\Person{雅各伯}，也不是\Person{若望}，而是他们极为敬重的一位，先知之首，天主的朋友，一个曾为他们如此热心以至于为他们忍受饥饿的人，一个直到今天仍未死去的人。这个人说了什么？\BibleQuote{主，他们杀了你的先知，拆毁了你的祭台；只剩下我一个人，他们还在寻索我的性命。}还有什么比这更残忍的暴行呢？当他们本应为已经犯下的罪行祈求宽恕时，他们却甚至想杀他。所有这些事都让他们完全无法得到宽恕。因为并非在饥荒盛行时，而是在季节适宜、他们的羞耻被除去、魔鬼（即假神）被羞辱、天主的大能得以彰显、君王也屈服于此时，他们犯下了这些狂妄之举，从谋杀到谋杀，除掉了他们的导师和那些会引导他们转向更好心思的人。那么他们对此能说什么呢？他们也是骗子吗？他们也是冒牌货吗？难道他们也不知道自己从哪里来吗？但他们困扰了你。是的，但他们也告诉了你们美好的事。但那些祭台呢？祭台肯定没有困扰你们吧？它们也激怒你们了吗？看，他们一直在证明自己的顽固和傲慢！这就是为什么在另一处，\Person{保罗}在给得撒洛尼人写信时说：\BibleQuote{你们也遭受了同族人的同样苦难，正如他们从犹太人那里遭受的一样，那些犹太人杀死了主和他们的先知，并且迫害了我们，不中悦天主，与众人为敌。}\BibleRef{得前 2:14}这也是他在这里所说的，他们既拆毁了祭台，又杀了先知。但是天主对他有什么答复呢？\BibleQuote{我为我自己保留了七千人，他们未曾向巴耳像屈膝。}\BibleRef{列上 19:18}这与当前的主题有什么关系呢？有人可能会说。它与当前的主题有很大关系。因为他在这里表明，即使承诺是对整个民族作出的，天主也使用那些值得的人来拯救。这也是他上面所指出的，当他说：\BibleQuote{虽然以色列子民的数目多如海沙，但得救的只是残余。}以及\BibleQuote{若不是万军的上主给我们留下苗裔，我们早已像索多玛一样了。}\BibleRef{罗 9:27}他也从这段经文中指出了这一点。所以他接着说：\BibleQuote{同样，在现今的时候，也按照恩宠的拣选，留有残余。}请注意，每一个词都保持自己的地位，同时显示出天主的恩宠和接受救恩者的顺服心态。因为通过说\Quote{拣选}，他表明了他们的得蒙悦纳；但通过说\Quote{恩宠}，他表明了天主的恩赐。\par

第6节。\BibleQuote{如果出于恩宠，就不再是由于行为了，否则恩宠就不成其为恩宠；但如果出于行为，就不再是由于恩宠了，否则行为就不成其为行为。}\BibleRef{罗 11:6}\par

他再次针对刚才所引经文中的犹太人的好辩性情，并藉此剥夺了他们的借口。他的意思是：你们甚至不能说，先知们确实呼唤了，天主邀请了，实况大声呼喊了，激起嫉妒之情足以将我们引向祂，但所吩咐的事太沉重，这就是我们为何不能亲近的原因，因为我们被要求展示行为，以及艰苦的善行。因为你们连这话也不能说。因为如果这样做只会使祂的恩宠黯然失色，天主怎么会向你们要求这个呢？他说这话，是出于意愿显明天主极欲他们得救。\BibleRef{申 5:29} 因为不仅他们的得救易于实现，并且向人显示祂的爱也是天主极大的光荣。那么，既然没有行为向你们要求，你们为何害怕亲近呢？当恩宠就在你们面前时，你们为何争辩争吵？又为何毫无意义地总是向我提起法律？因为你们不会藉着那法律得救，也会糟蹋这份恩赐。因为如果你们顽固地坚持靠法律得救，你们就废弃了天主的这份恩宠。然后，为了不让他们觉得奇怪，他先引用了那七千人；他说他们是因恩宠得救的。因为当他说：\Quote{同样，现在也有按照恩宠拣选的残余。}时，表明他们也是因恩宠得救的。不但如此，而且他也藉着说：\Quote{我为自己保留了。}因为这是那位显示自己是主要施予者的话语。如果因恩宠，那么有人会说，为什么我们所有人没有得救？因为你们不愿意。因为恩宠，虽然是恩宠，却拯救那些愿意的人，而非那些不愿接受、转离恩宠、坚持与它争战、反对它的人。请留意，他始终在证明这一点：\Quote{并不是说天主的话落了空，}藉此表明那些配得的人就是应许所临到的人，而且这些人，尽管为数不多，但仍然是天主的人民；实际上他在书信开头已经很有力地陈述了这一点，他说：\Quote{即便有些人不信，} \BibleRef{罗 3:3}，并且他并未止步于此，还继续说：\Quote{但愿天主是真实的，众人都是虚谎的。} \BibleRef{罗 3:4} 在这里他又以另一种方式证实了这一点，并显明恩宠的力量，以及总有一些人得救，另一些人丧亡。那么让我们感恩，因为我们属于那些得救的人，并且自己无法藉行为得救，而是藉天主的恩赐得救的。但在感恩时，我们不要只在言语上如此行，也要在行为和行动上。因为真正的感恩就在于：我们行事确保天主得到光荣，并逃避那些我们已从中被释放的事物。因为如果我们辱骂了君王，反受尊荣而非惩罚，随后又去重新辱骂祂，那么既然我们被发觉极度忘恩负义，就理当遭受极重的惩罚，远较先前为大。因为先前的傲慢并不如受尊荣和蒙眷顾之后所犯的傲慢那样显明我们的忘恩负义。那么让我们逃避那些我们已经从中被释放的事物，不要仅在口舌上感恩，免得也对我们说：\Quote{这民族用嘴唇尊敬我，他们的心却远离我。} \BibleRef{依 29:13} 因为当\Quote{诸天宣扬天主的荣耀} \BibleRef{咏 18:1}，而你，诸天为荣耀祂而为你受造，却行出那样的事，以至藉着你，那造你的天主受亵渎，这岂非不合宜吗？为此，不仅那亵渎者，连你自己也要受罚。因为诸天也不是藉着发出声音来光荣天主，而是藉着使别人看到它们就行动，然而它们仍被说成\Quote{宣扬天主的荣耀。}同样，那些提供可赞叹的生活的人，即使沉默不语，也光荣了天主，因为别人藉着他们光荣祂。因为祂因诸天所受的尊敬，远不如因无瑕的生活。因此，当我们与外邦人交谈时，我们引用（有四份抄本或读作\Quote{我们不要引用}）的不是他们眼前的诸天，而是那些人——他们虽曾处于比禽兽更糟的境地，他却说服他们成为天使的竞争者。而我们（一份抄本作\Quote{让我们}）藉着谈论这种转变来堵住他们的口。因为人远胜于诸天，他可能拥有比诸天之华美更明亮的灵魂。因为诸天虽然长久可见，却未能说服许多人。但\Person{保禄}在宣讲短时间之后，将整个世界吸引到自己身边。\EditorialNote{圣奥斯定论《圣咏》19:4} 因为他拥有一个不亚于诸天的灵魂，能够吸引所有人归向他。我们的灵魂连与地相比都不配，但他的灵魂却与诸天相当。地确实保持自己的界限和规律；但他的灵魂之高超越诸天，并与基督亲自交谈。\BibleRef{格后 10:15} 其美丽如此伟大，连天主都传扬它。因为星辰受造时，天使都惊叹。\BibleRef{约 38:7} 但祂在说\Quote{他是我特选的器皿} \BibleRef{宗 9:15}时，却惊叹于这个。而诸天常常被云彩遮蔽。但\Person{保禄}的灵魂从未被任何诱惑遮蔽，即使在风暴中，他也比那在正午时分的坚实天空（\OriginalTerm{坚实}{\GreekText{σταθερἅς}} \OriginalTerm{正午}{\GreekText{μεσημβρίας}}）更为清晰可见，并如云彩来临之前一样恒常闪耀。因为那在他内照耀的太阳，所发出的光线并未被众多诱惑遮蔽，反而照耀得更加灿烂。因此他说：\Quote{我的恩宠为你够了，因为我的能力在软弱中才成全。} \BibleRef{格后 12:9} 那么让我们努力效法他，这样，即使现在这个诸天对我们也将如同无物，如果我们愿意，太阳和整个世界也是如此。因为这些都是为我们而造，我们并非为它们而存在。让我们证明自己配得享有这些为我们所造之物。因为如果我们被发觉不配享用这些，又怎能配得国度呢？因为凡是如此生活以致亵渎天主的人，都不配看见太阳。那些亵渎祂的人，不配享受光荣祂的受造物：因为连一个侮辱父亲的儿子，也不配被忠仆服侍。因此，这些人将享受荣耀，而且是大荣耀；但我们却要遭受惩罚与报应。\par

那么，那为你们而受造的受造物，将要按\BibleQuote{天主儿女光荣的自由}被塑造，这是何等可怜的事（\BibleRef{罗 8:21}）；但我们这些成为天主儿女的人，却因我们的极大懈怠，被遣往毁灭和地狱；而为了我们，那受造物却将享受那伟大的欢庆？现在，为免此事发生，让我们中那些拥有纯洁灵魂的人保持其纯洁，甚至要使其光辉更加炽烈。而让我们中那些灵魂玷污的人不要绝望。因为（他说）\BibleQuote{如果你们的罪如朱红，我必使它们白如雪；如果它们如丹霞，我必使它们白如羊毛。}（\BibleRef{依 1:18}）但是，当天主应许时，不要怀疑，而要去做那些可以使你获得这些应许的事。你所做的可怕和残暴之事难道数不胜数吗？那又怎样？因为直到现在，你还没有进入那个无人再能忏悔的坟墓（\BibleRef{依 38:18}；\BibleRef{咏 6:5}）。至今，剧场（\OriginalTerm{剧场（或竞技场）}{\GreekText{θέατρον}}）还没有为你被拆解，你仍站在界线之内，甚至可以通过最后的挣扎挽回所有的失败。你还没有来到有财主的地方，以致你听到说：\BibleQuote{在你我之间有一深坑}（\BibleRef{路 16:26}）。新郎还没有临近，以致有人害怕分给你油。你仍然可以购买和积蓄。还没有人对你说：\BibleQuote{不行，恐怕不够我们和你们用的}（\BibleRef{玛 25:9}）；但有许多在售卖的人：赤身的、饥饿的、患病的、被囚禁的。给这些食物，给那些衣物，看望病人，油会比泉水更多。算账的日子还没有到。善用时间，酌情减少债务，对欠你\Quote{一百桶油}的人说：\BibleQuote{拿出你的账，写五十桶}（\BibleRef{路 16:6}）。对于金钱、言语，以及一切其他事物，都要照样去做，效法那个管家。并且把这劝告给你自己和你的亲属，因为你仍然还有说话的力量。你还没有到需要替他们叫别人的地步，但你有能力立刻给自己和他人提出建议（\BibleRef{路 16:28}）。但是，当你离开那里之后，这些事你一样也无法在必要时做了。而且理当如此。因为你去世之时已被定了这么长的期限，你既没有使自己获益，也没有使任何人获益，那么当你落在审判者手中时，你怎么能获得这恩宠呢？所以综合这一切，让我们紧紧抓住我们自己的救恩，不要失去今生的机会。因为甚至在我们最后一口气时，也有可能讨天主的喜悦。通过你的临终遗嘱获得认可，这并非不可能，尽管不像有生之年那样，但仍然是可能的。怎么做，以什么方式？如果你把祂列入你的继承人中，并且也给他（\OriginalTerm{也给他}{\GreekText{καὶ} \GreekText{αὐτῷ}}）一份你的全部财产。你在有生之年没有喂养过祂吗？至少在你离世后，当你不再是主人时，给祂一份你的财物。祂是爱人者，不会吝啬地对待你。当然，在有生之年喂养祂，是更大渴望的标志，因此回报更大。但如果你没有这样做，至少做次好的事。使祂与你的儿女同为继承人（参阅第384页），如果你对此拖延，请思考：祂的父曾使你与祂同为继承人，并克服你非人的精神。因为连那些使你有份于天堂、为你而被杀的，你甚至都不让祂与你儿女分享，你还有什么借口呢？然而祂所做的一切，都不是偿还债务，而是施予恩惠。但你在接受了如此大的恩惠之后，也成了债务人。虽然如此，祂加冕给你，仿佛是在接受恩惠，而不是要求偿还债务；并且，祂要收回的也是祂自己的东西。那么，交出你的金钱吧，它现在对你毫无用处，而且你也不再是主人；祂将赐给你一个永远有益于你的国度，并且也将赐予今世之物。因为如果祂成为你儿女的共同继承人，祂将减轻他们的孤儿之苦，消除对他们的阴谋，击退侮辱，堵住挑拨者的嘴。如果他们自己无法捍卫所继承的遗产，祂将亲自站起来，不让它们被破坏。但如果祂允许这样，那么祂会以更大的慷慨亲自补足遗嘱中所有安排，因为祂在遗嘱中只是被提及而已。那么，立祂为你的继承人吧。因为你正要去祂那里。祂将亲自作你的审判官，审判在这里所做的一切。但是有些人如此可怜和吝啬，即使他们没有儿女，也不敢这样做，反而赞成把他们的财物给予附从者、谄媚者、张三李四，而不给基督——那位对他们施予如此大恩惠者。有什么比这行为更无理的呢？因为如果有人要把这类人比作驴子，甚至石头，仍然不足以说明他们的无理和愚钝。也找不到一个比喻来向你展示他们的疯狂和昏聩。对于这些有生之年没有喂养祂，甚至临终时也不愿从那些他们已不再拥有的财物中给祂一点，反而怀有如此敌对和敌意的性情，连对自己无用的东西都不肯分给祂的人，他们能得到什么宽恕呢？难道你不知道有多少人甚至不配得到这样的结局，而是突然被夺走吗？但是天主却授权你向亲属发号施令，与他们谈论你的财产，并安排你家中一切。那么，在接受了祂这个恩惠之后，你却背信弃义地放弃了这益处，几乎与信仰中的祖先正面对立，你还能有什么辩护呢？因为他们甚至在生前就变卖一切，放在宗徒脚前。但你即使在临终，也不肯分给有需要的人任何一份。更好的部分，并给予人大胆的，是在生前救济贫穷。但如果你没有这样打算，至少在你临终时做一些高尚的事。因为这不是强烈的爱基督的表现，但仍然是爱的行为。因为如果你不愿与羔羊同在高处，那么即使只是跟在它们之后，也不是小事，因此不会与山羊一同站在左边。但如果你连这也不做，那么什么辩护能拯救你呢？当死亡的恐惧、金钱从此对你有无用处、为儿女留下平安、在那里为将来预备许多宽恕，都不能使你发慈悲之时？因此，我劝告，最好的做法是在你有生之年将财产的大半施舍给穷人。但如果有谁心胸狭窄，不忍心这样做，至少让他们因必要而变得慈悲。因为当你还活着，仿佛没有死亡时，你紧紧抓住你的财物。但现在你既然知道你要死了，至少现在放弃你的想法，像必死之人一样处理你的事务。或者，更像一个应当永远享受永生的人。因为如果我要说的话令你不悦，而且充满恐怖，但仍然必须说出来。与主清算你的奴仆。你正在给予你的奴仆自由吗？给予基督自由，使他免于饥饿、困苦、囚禁、赤身露体。你对这些话感到恐惧吗？那么当你连这一点都不做时，岂不更可怕？在这里，这话使你毛骨悚然。但当你去了那个世界，必须听到比这些更可怕的事，见到无法医治的刑罚，你会说什么呢？你会逃往谁那里避难？你会呼求谁与你联盟和援助？会是亚巴郎吗？他不会听你。或是那些童女？她们不会分给你油。那么你的父亲或祖父？这些人中，即使他再圣洁，也没有能力推翻那个判决。那么，权衡这一切，向那位唯一能涂抹你的债卷、熄灭那火焰的主祈祷恳求，并赢得祂的欢心，现在持续喂养祂、给祂衣穿：使你在此世可以带着良好的希望离世，并在那里享受永恒的福乐。愿我们众人借着恩宠与对人的爱得以达到，等等。\par

\SourceRecord{Translated by J. Walker, J. Sheppard and H. Browne, and revised by George B. Stevens.；Nicene and Post-Nicene Fathers, First Series；Vol. 11.；Edited by Philip Schaff.；Buffalo, NY: Christian Literature Publishing Co.,；1889.；<http://www.newadvent.org/fathers/210218.htm>.}

% structure: p0001 source=h1 target=section reason=page-title

\section{罗马书讲道集 第十九篇}

\BibleRef{罗 11:7}\par

\BibleQuote{\Emph{那么怎样？}\Emph{以色列人所寻求的，他们没有获得；而选民却获得了；其余的都成了盲目的。}}\par

他曾说过，天主并没有抛弃祂的子民；为了表明在何种意义上祂没有抛弃他们，他再次求助于先知。借着先知表明多数犹太人失丧之后，为了避免显得是在提出自己的指控，使他的言论冒犯人，并像是攻击敌人一样，他求助于\Person{达味}和\Person{依撒意亚}，说道：\par

\BibleRef{罗 11:8}. \BibleQuote{正如经上所记载：天主赐给他们沉睡的灵。} \BibleRef{依 29:10}\par

或者更好的说法是，我们应该回到他的论证的开头。然后，在提及\Person{厄里亚}时代的情况并表明恩宠是什么之后，他接着说：\BibleQuote{\Emph{那么怎样？以色列所寻求的并没有获得。}} 这既像一个控告者的话，也像一个质问者的话。他的意思是，\Person{犹太人}寻求义却不愿接受，是前后矛盾的。然后，为使他们无可推诿，他通过那些接受了的人指出他们的无情心态，如他所说：\BibleQuote{但选民获得了。} 他们就成为其他人的定罪。这就是\Person{基督}所说的：\BibleQuote{如果我藉\Person{贝耳则步}驱魔，你们的子弟是藉谁驱魔？因此他们将是你们的审判者。} \BibleRef{路 11:19} 为防止有人归咎于事情的性质而不归咎于他们自己的态度，他指出了那些获得了的人。因此他非常恰当地使用了这个词，以同时显示来自上天的恩宠和这些人的热忱。他说他们\Quote{获得了}（如同偶然获得，希腊文\OriginalTerm{获得了}{\GreekText{ἐπέτυχε}}），并不是否定自由意志，而是要显示恩惠的伟大，且大部分来自恩宠，虽然不是全部。因为我们也常说，\Quote{某人偶然获得了}（同一个词），\Quote{某人碰上了}，当所得很大时。因为大部分不是靠人的劳力，而是靠天主的恩赐所带来的。\par

\BibleQuote{其余的都成了盲目的。}\par

请看，他多么大胆地用自己的声音讲述其余人的遗弃。他先前已经提到过，但那是借着先知作为控告者。但从这里起，他亲自宣布了。然而即使在此，他也不满足于自己的宣告，而是再次引用了\Person{依撒意亚}先知。因为说了\Quote{成了盲目的}之后，他接着说：\Quote{正如经上所记载：天主赐给他们沉睡的灵。} 那么这盲目从何而来？他先前已经提到它的原因，并将其全归咎于他们自己，表明这是由于他们不合时宜的顽固。现在他也讲述了。因为当他说\BibleQuote{眼睛不得看见，耳朵不得听见}时，他正是在责备他们的好斗精神。因为他们有\Quote{眼睛看}奇迹，有\Quote{耳朵听}那奇妙的教导，却从未适当使用它们。至于\BibleQuote{祂赐给}，不要想象这里指主动作用，而只是许可。但\Quote{沉睡}（\OriginalTerm{沉睡}{\GreekText{κατάνυξις}}，原意刺透）是他给那种倾向于恶的习性起的名称，当这种习性不可救药且不可改变时。因为在另一处，\Person{达味}说：\BibleQuote{为叫我的荣耀歌颂祢，而不致沉睡。} \BibleRef{咏 29:12} 意即，我不会改变，不会转变。因为一个人若因虔诚的敬畏而昏昏沉睡，就不会轻易改变立场；同样，沉睡在邪恶中的人也不会轻易改变。因为这里的昏沉沉睡，无非是牢牢固定在一件事上。所以为了指出他们精神的不可救药和不变，他称之为\Quote{沉睡的灵}。然后，为了显示这种不信将受到最严厉的惩罚，他再次引入先知，预言了后来确实发生的事。\par

\BibleRef{罗 11:9}. \BibleQuote{让他们的筵席成为网罗、陷阱和绊脚石。} \BibleRef{咏 69:22}\par

也就是说，让他们的安慰和一切好东西都改变、消逝，让他们容易受到任何人的攻击。为了表明这是因罪受罚，他加上：\BibleQuote{并成为他们的报应。}\par

\BibleRef{罗 11:10}. \BibleQuote{让他们的眼睛变暗，以致看不见，并时常弯下他们的腰。}\par

这些事难道还需要解释吗？它们不是连最愚钝的人也明白吗？在我说这些话之前，事实本身已经为我们所讲的做了见证。因为什么时候他们曾如此易受攻击？什么时候他们曾如此易于得手？什么时候祂曾\Quote{弯曲了他们的脊背}？什么时候他们曾处于这样的奴役之下？而且，这些恐怖之事没有任何解脱的余地。先知也暗示了这一点。因为他不仅说\Quote{求祢弯曲他们的脊背}，而且说\Quote{永远弯曲他们}。但是，犹太人啊，如果你要就这事的结果争论，从先前的事也可以知道现在的处境。你下到埃及；两百年过去了，天主迅速将你从那奴役中解救出来，尽管你是不虔敬的，并以最有害的淫行行淫。你从埃及被解救出来，却崇拜牛犊，将你的儿子祭献给巴耳培敖尔，玷污了圣殿，追求各种恶行，甚至不认识本性本身。你用可咒骂的祭献充满了山、丛林、丘陵、泉源、河流、花园，你杀了先知，推翻了祭台，表现出各种邪恶和不虔敬的极端。尽管如此，在将你交给巴比伦人七十年之后，祂再次将你带回先前的自由，将圣殿、你的国家、你旧有的政体还给你，又有先知和圣神的恩赐。甚至在你被掳的时期，你也没有被遗弃，那里还有\Person{达尼尔}、\Person{厄则克耳}，在埃及有\Person{耶肋米亚}，在旷野有\Person{梅瑟}。之后你再次回到先前的恶行，成了一个狂欢\OriginalTerm{狂欢}{\GreekText{ἐξεβακχεύθης}} \BibleRef{加下 14:33}的人，并在不虔敬的\Person{安提约古}时代将你的生活方式\OriginalTerm{生活方式}{\GreekText{πολιτείαν}}改变为希腊式的\BibleRef{达 8:14}；\BibleRef{加上 4:54}。但即使那时，你们只被交给安提约古三年多一点，然后通过\Person{玛加伯}们，你们再次竖立了那些光辉的胜利。但现在完全相反，没有这样的事。这是最令人惊讶的，因为恶行已经停止，惩罚却增加了，而且没有任何改变的希望。因为不仅七十年过去了，也不是一百年，也不是两倍，而是三百多年，还多出不少，却找不到一丝这样的希望。尽管你们现在既不是偶像崇拜者，也没有做先前那些放肆的事。那么原因是什么？实体取代了预像，恩宠排除了法律。这就是先知从古时预言所说的：\Quote{永远弯曲他们的脊背}。请看预言的细致之处，它如何预言了他们的不信，也指出了他们的好辩，并显示了随之而来的审判，表明了惩罚的无尽。因为许多较为迟钝的人，由于不信将来要发生的事，想通过现在的事看到将来的事，从这个时候起，天主在两方面证明了自己的能力：使相信的外邦人高升于天之上，使不信的犹太人降到最荒凉的地位，并把他们交到永无止境的灾祸中。在严厉地指责他们既不信，又遭受和将要遭受的苦难之后，他再次缓和了他所说的话，接着写道：\par

第11节。\Quote{我且问：他们失足是要跌倒吗？断乎不是！}\BibleRef{罗 11:11}\par

当他表明他们应受无数灾祸之后，他便想出一个缓解之策。请考量\Person{保禄}的判断。他从先知那里引入了控告，但缓解之策却来自他自己。因为他说，他们犯了大罪，无人能否认。但我们来看，这失足是否不可治愈，是否完全阻碍了他们重新站起。其实并非如此。你看他如何再次攻击他们，而在期待缓解之下，他证明他们犯有自己承认的罪。但我们来看，他究竟为他们设计了怎样的缓解之策。这缓解之策是什么？他说：\Quote{等到外邦人的数目满了}，\Quote{那时全以色列都要得救}，这要在他第二次来临和世界末日的时候。但他并没有立刻这样说。因为他猛烈地攻击他们，将控告连于控告，带来一个又一个先知向他们大声呼喊：\Person{依撒意亚}、\Person{厄里亚}、\Person{达味}、\Person{梅瑟}、\Person{欧瑟亚}，不是一次两次，而是多次；唯恐这样一方面使他们陷入绝望，筑起一道墙阻挡他们接近信仰，另一方面又使那些信了的外邦人无端骄傲，他们若自高自大，也会在信仰上受损，他便进一步安慰他们说：\Quote{然而因他们的过犯，救恩临到了外邦人}。但我们不可按字面理解这里的话，而要领会说话者的精神与目的，以及他意在达成什么。这是我一向恳求你们的爱德要做的。因为若我们心中怀有这一点，拿起这里的话，就不会在任何部分遇到困难。因为他现在的忧虑是要除去那些从他的话中外邦人可能产生的傲慢。因为这样，他们学会了谦逊，就更有可能在信仰中保持不变；同样，那些犹太人若摆脱了绝望，也更可能乐意接受恩宠。因此，顾及他的这一目的，我们就这样来听这段经文所说的一切。那么他说什么？他哪里显明他们的失足不是无可挽救的，他们的被弃不是最终的？他从外邦人论证，这样说：\par

\Quote{因他们的过犯，救恩临到了外邦人，为要激起他们的妒愤。}\par

这番话不仅是他自己的，也是福音中比喻的含义。因为那为祂儿子摆设婚宴的，当客人不愿来时，就召了路上的人来。\BibleRef{玛 22:9}。那栽种葡萄园的，当园户杀了继承人时，就把葡萄园租给别人。\BibleRef{玛 21:38}。不用任何比喻，祂亲自说：\BibleQuote{我奉派，无非是到以色列家迷失的羊那里去。}\BibleRef{玛 15:24}。对于那客纳罕妇人，当她坚持时，祂还多说了一些：\BibleQuote{拿儿女的饼扔给小狗，是不对的。}\BibleRef{玛 15:26}。保罗对那些骚动的犹太人说：\BibleQuote{天主的道先讲给你们听，原是应当的；但你们既弃绝自己，断定自己不配得永生，我们就转向外邦人。}\BibleRef{宗 13:46}。整件事清楚表明，事物的自然顺序是先有他们进来，然后才有外邦人；但因他们的不信，顺序颠倒了；他们的不信和跌倒反倒使这些外邦人先被带进来。故此他说：\BibleQuote{因他们的过犯，救恩到了外邦人，为要激发他们嫉妒。}但如果他将事物发展的结果说得好像上智安排的主要意图，不要感到惊讶。因为他想安慰他们沮丧的心灵，他的意思大致是这样：耶稣先来到他们那里；他们却不接纳祂，尽管祂行了无数奇迹，反而把祂钉死了。因此祂将外邦人吸引过来，使得他们所享有的荣誉，因刺痛他们麻木的心，至少由于对他人的怨恨，诱使他们归附。他们应当先被接纳，然后才是我们。这正是为何他说：\BibleQuote{这福音就是天主的德能，为救一切相信的人，先为犹太人，后为希腊人。}\BibleRef{罗 1:16}。但既然他们已背离，我们这些后来的反而成了先来的。请看他由此为他们聚集了多大的荣誉。第一，他说，我们在他们不愿意时被召；第二，他说，我们被召的原因不仅是为使我们得救，也是为使他们因嫉妒我们的得救而变得更好。那么他说什么呢？难道是为了犹太人的缘故，我们才根本不会被召而获救吗？在我们之前不会，但按正常顺序会。因此，当祂对门徒说话时，祂并非仅仅说：\BibleQuote{到以色列家迷失的羊那里去}\BibleRef{玛 10:6}，而是说：\BibleQuote{宁可往迷失的羊那里去}，以示此后也必须到那些地方去。保罗又说：\BibleQuote{天主的道先讲给你们听，原是应当的}，而不是\BibleQuote{应当讲给你们听}\BibleRef{宗 13:46}，以示其次也当讲给我们。这一切言行，都是为使他们不能无耻地推诿说被忽略，并因此不信。因此，基督虽然事先知道一切，却仍然先来到他们那里。\par

\BibleQuote{现在，如果他们的跌倒成为世界的富足，他们的减少成为外邦人的富足，那么，他们的丰满岂不更加如此？}\BibleRef{罗 11:12}\par

在此他说话是为了取悦他们。因为即使他们跌倒了一千次，外邦人若不显示信德，也不会得救。同样，犹太人若不信从、不争辩，也不会灭亡。但如我所说，他现在是在安慰他们跌倒之后，给予他们更多的理由相信，如果他们改变，就必得救。因为如果当他们绊倒时，他说，那么多人得了救恩，当他们被驱逐时，那么多人被召，试想当他们回头时将如何。但他没有这样说：当他们回头时。他没有说\Quote{他们的回头}或\Quote{他们的改变}或\Quote{他们的善行}，而是\Quote{他们的丰满}，即当他们全都将进来的时刻。他说这话是为了显示，那时也是恩宠和天主的恩赐将做大部分，或几乎全部的工作。\par

\BibleQuote{因此，我对你们外邦人说：我既然是外邦人的宗徒，我就光荣我的职务，这样，或许可以激励我的同胞，拯救他们中的一些人。}\par

他再次极力使自己摆脱不当的猜疑。他似乎是在责备外邦人，压制他们的骄矜，但同时也温和地触动了犹太人。实际上，他四处周旋，试图掩饰和消解他们这极大的败落。但碍于事实的性质，他找不到办法做到这一点。因为根据他所说的，当那些远不及他们的人已领取了为他们预备的美物时，他们反而该受更大的谴责。因此，他于是从犹太人转到外邦人，并将关于外邦人的部分插入他的讲论中间，仿佛想表明他说这一切是为了教导他们变得明理。他说，我因两件事称赞你们：一是因为作为你们受委任的仆人，我不得不这样做；二是藉着你们我可以拯救他人。他没有说\Quote{我的弟兄，我的骨肉}，而是说\Quote{我的血肉}。接着，在指出他们好辩的性情时，他没有说\Quote{我或许能}说服，而是说\Quote{激励嫉妒并拯救}；这里又不是全部，而是\Quote{他们中的一些}。他们是何等刚硬！即使在责备中，他再次显明外邦人是受尊敬的，因为他们是他们得救的原因，且方式不同。因为外邦人藉着不信成了福分的提供者，而这些（信徒）却藉着信德成了犹太人的福分。因此外邦人的地位似乎既平等又优越。因为犹太人啊，你能说什么呢？说若我们没有被抛弃，他就不会这么快被召吗？外邦人也可以说：若我没有得救，你就不会受到激励而嫉妒。但若你要知道我们的优势何在，我是藉着信德拯救你，而你是藉着绊跌，在你之前为我们提供了通途。然后，他感到又刺中了他们的要害，于是重拾之前的论点，说：\par

\BibleRef{罗 11:15}\BibleQuote{如果他们被弃绝，就是世界的和好，那么他们被收纳，岂不是死而复生吗？}\par

然而这又定他们的罪，因为别人因他们的罪获益，他们自己却没有从别人的善行中得益。但若他断言那必然发生的事是他们的作为，不要惊讶：因为（正如我多次说过的）他这样做是为了谦逊这些人，劝勉那些人，才将他的讲话采用这种形式。因为正如我之前所说，即使犹太人被弃绝一千次，而外邦人不显明信德，他们也绝不会得救。但他站在软弱的一方，给予受困者援助。但请看他甚至在恩待他们时，如何只在言语上安慰他们。\BibleQuote{如果他们被弃绝，就是世界的和好}，（这对犹太人又算什么呢？）\BibleQuote{那么他们被收纳，岂不是死而复生吗？}然而，除非他们被收纳，这对他们也不是恩赐。他所说的要点是这样的：若在愤怒中天主赐予别人如此大的恩赐，那么当他与他们和好时，他还会赐予什么呢？但死人复活并不是因他们被收纳，同样，现今我们的得救也不是藉着他们。他们因自己的愚昧被弃绝，而我们因信德和自上而下的恩宠得救。但这一切对他们都无帮助，除非他们显出应有的信德。然而，他照常行事，继续另一个赞美，而这并不是真正的赞美，只是看起来像是，如此效法最明智的医师，他们给予病人尽可能多的安慰，正如疾病的性质所允许的。他说了什么呢？\par

\BibleRef{罗 11:16}\BibleQuote{如果初熟的果子是圣的，那团面也是圣的；如果根是圣的，那枝子也是圣的。}\par

他在此以初熟果子和根称呼亚巴郎、依撒格、雅各伯、先知、族长、以及旧约中所有显要人物；而枝子则指那些信了的人。既然事实是许多人没有信，请看他是如何再次（\OriginalTerm{暗中破坏}{\GreekText{ὑποτέμνεται}}，见345页）这一点，又说：\par

\BibleRef{罗 11:17}\BibleQuote{若有几根枝子被折下来。}\par

然而上文你说大多数人都灭亡了，只有少数得救。那么，你论及那些灭亡的人，为什么用了\Quote{有几根}这表示少数的词呢？他回答说，这并不是与我自己相矛盾，而是出于渴望讨好并挽回那些忧伤的人。请看在整个段落中，你都能发现他致力于这一目标，即安慰他们。你若否认，就会有许多矛盾。但请容我留意他的智慧，他如何在似乎为他们说话、为他们寻求安慰时，暗中打击他们，显明他们毫无借口，甚至从\Quote{根}、从\Quote{初熟果子}也是如此。因为要想想枝子的恶劣，当它们有甘甜的根时，却仍不效法它；而当面团未被初熟果子改变时，也是有过错的。\BibleQuote{若有几根枝子被折下来。}然而，大部分枝子被折下来了。但正如我所说，他愿意安慰他们。这就是为什么他不是以自己身份，而是以他们的身份引入所用的话语，即便如此也暗中打击他们，显明他们已不再是亚巴郎的亲属。\BibleRef{玛 3:9}因为他想要说的是，他们与亚巴郎毫无共同之处。\BibleRef{若 8:39}因为如果根是圣的，而这些不圣，那么它们就远离了根。然后在安慰犹太人的表象下，他又藉着指责击打外邦人。因为说了\BibleQuote{若有几根枝子被折下来}之后，他继续说：\par

\BibleQuote{而你本是野橄榄，却被接在枝上。}\par

因为外邦人越不受重视，犹太人看见他享受犹太人的福分就越发恼怒。而对那另一人来说，他所受的微不足道的鄙视，与这转变的荣耀相比算不得什么。且看他何等巧妙。他并没有说\Quote{你被}接\Quote{上}，而是说\Quote{你被接枝}，这再次刺痛了犹太人的心，表明外邦人正站在他自己的树上，而犹太人自己躺在地上。因此他并未在此止步，在说了接枝之后也没有停（虽然他已借此表明了全部），但他仍对外邦人的繁荣状态耿耿于怀，并以这些话夸大他的美名：\Quote{并且一同分享那橄榄树的根和肥汁。}他似乎确实将他视为一个附加物。但他表明他并未因此而变得逊色，反而拥有一切，即那从根长出的枝条所拥有的一切。以免你听见\Quote{你被接枝}这话，就以为他比自然枝条有所欠缺，请看他是如何借说\Quote{与他们一同分享橄榄树的根和肥汁}，使他与之平等：也就是说，你被置于同一高贵等级、同一本性中。然后在责备他时，说，\par

第十八节。\BibleRef{罗 11:18}\Quote{不可向枝条夸口。}他似乎确实在安慰犹太人，但却指出他的卑贱和极大的羞辱。这就是为什么他不说\Quote{不可夸口}，而是说\Quote{不可向}夸口——不要向他们夸口以致切断他们。因为你是被放置在他们的位置上，享受他们的福分。你看见他如何好像在责备一人，实际上却在尖锐地抨击另一人吗？\par

\Quote{但你若夸口，}他说，\Quote{不是你托着根，而是根托着你。}\par

这对被砍下的枝条有什么意义呢？毫无意义。因为如我所说，他似乎编织了一种微弱的安慰影子，却在瞄准外邦人的同时，给了他们致命的一击；因为他说\Quote{不可向他们夸口}和\Quote{你若夸口，不是你托着根}，他向犹太人表明所发生的事确实值得夸耀，即使不该夸耀，从而一面激励他、激发他信德，一面以辩护人的姿态打击他，指出他所受的惩罚，以及他人已拥有了他们自己的东西。\par

第十九节。\BibleRef{罗 11:19}\Quote{你必要说，}他接着说，\Quote{枝条被折下来，是为叫我接上。}\par

他又以反对的方式，借着反驳来确立与先前立场相反的观点，以表明他先前所说的话并非直接与主题相关，而是为把他们吸引过来。因为救恩临到外邦人不再是由于他们的跌倒，他们的跌倒也不是世界的富足。我们得救也不是因为他们跌倒了，反倒是相反。他表明对外邦人的眷顾是一个主要目的，尽管他似乎以另一种形式表达他的话。整段是一连串的反驳，在其中他洗脱了仇恨的嫌疑，使他的言语变得可以接受。\par

第二十节。\BibleRef{罗 11:20}\Quote{不错，}他称赞他们所说的，然后再次警告他们说：\Quote{因不信他们被折下来，你因信被接上。}\par

所以这里又一次赞扬，而对另一方则是控告。但他继续说\Quote{不要骄傲，但要畏惧}，从而再次压制他们的傲慢。因为这事不是本性的问题，而是信与不信的问题。他似乎再次勒住外邦人，却在教导犹太人不可固守自然的亲属关系。因此他接着说\Quote{不要骄傲}，他并没有说\Quote{但要谦卑}，而是说\Quote{畏惧}。因为骄傲滋生轻视和懈怠。然后，当他为了减少陈述的痛苦而进入他们灾难的全部悲哀时，他以下列方式作为对另一方的责备表达：\par

第二十一节。\BibleRef{罗 11:21}\Quote{因为天主若没有顾惜自然的枝条，}然后他并没有说\Quote{也不会顾惜你}，而是\Quote{小心，免得祂也不顾惜你。}因此他以修剪（\OriginalTerm{}{\GreekText{ὑποτεμνόμενος}}）的方式，削去他陈述中令人不快之处，将信徒描绘为在挣扎中，他同时将犹太人拉近自己，也使这些人谦卑。\par

第二十二节。\BibleRef{罗 11:22}\Quote{所以你要看见天主的恩慈和严厉：向跌倒的人是严厉，但向你是有恩慈的，只要你继续在他的恩慈里；不然你也要被砍下。}\par

他并不是说，看你所行的善，看你所劳的，而是说，\Quote{看天主的恩慈}待人，以示一切都来自上方的恩宠，并使我们战栗。因此，这夸口的理由应当使你恐惧：因为主（\OriginalTerm{}{\GreekText{δεσπότης}}）曾善待你，所以你当恐惧。因为祝福不会一成不变地停留在你身上，你若变得懈怠；同样，灾祸也不会一成不变地停留在他们身上，他们若能改变；\Quote{因为你，}他说，\Quote{除非你继续在信德中，也要被砍下。}\par

第二十三节。\BibleRef{罗 11:23}\Quote{他们若不长久在不信中，也要被接上。}\par

因为不是天主砍下了他们，而是他们自己折断了、跌倒了，他说\Quote{自己折断了}说得很好。因为祂从未如此（\EditorialNote{Sav. conj. 抄本校勘}\OriginalTerm{如此}{\GreekText{οὗτος}}）抛弃他们，尽管他们犯了如此多、如此频繁的罪。你看见人的自由选择是何等伟大，心智的功效何等巨大。因为这些事情中没有一件是不可变的，无论是你的善还是他的恶。你也看见他如何使沮丧者兴起，使自信者谦卑；你不可因听到严厉而气馁，也不可因听到良善而自恃。祂严厉地砍下你的原因，是使你渴望回来。祂向你显示良善的原因，是使你继续（他没有说信德，而是说）祂的良善，也就是，如果你做相称于天主对人的爱的事。因为需要比信德更多的东西。你看见他如何既不容这些低沉，也不容那些高傲，反而激起他们的嫉妒，通过他们给予犹太人能力重新占据这个人的位置，正如那人先占据了另一个人的位置。他通过犹太人和发生在他们身上的事使外邦人感到恐惧，免得他们因此自高。但他试图通过赐给希腊人的事鼓励犹太人。因为他说：你也会被砍下，如果你变得懈怠（因为犹太人已被砍下），而他会重新被接上，如果他恳切，因为你也曾被接上。但非常明智的是，他将一切话都指向外邦人，这是他惯常的做法，通过责备强者来纠正弱者。在本书信的末尾，当他谈论遵守食物时，他也这样做。然后，他将此基于已经发生的事，而非仅基于将来的事。这更可能说服他的听众。当他打算展开一连串无可反驳的推理时，他首先使用从天主的大能而来的论证。因为如果他们被砍下、被抛弃，而别人在他们的位置上占了先，即使现在也不要绝望。\par

他说：\Quote{天主有能力把他们重新接上}，因为祂行的是超乎期望的事。但如果你希望事情有序、推理连贯，你从自己身上就有一个远超所需的证据。\par

第24节。\BibleQuote{如果你从按本性是野的橄榄树上被砍下来，逆着本性被接在好橄榄树上，何况这些按本性是天然的枝条，岂不更要被接在自己的橄榄树上吗？}\BibleRef{罗 11:24}\par

既然如此，如果信德能做逆本性的事，那么更能做合乎本性的事。因为如果此人从按本性是他的父辈中被砍下来，逆着本性归到亚伯拉罕身上，那么你岂不更能恢复自己的吗？因为外邦人的厄运是合乎本性的（他按本性是野橄榄），而好事是逆本性的（他逆本性被接在亚伯拉罕上），但你的好运是合乎本性的。因为你并非像外邦人那样要接在别人的根上，而是要接在自己的根上，只要你愿意回来。那么，你还有什么可夸的呢？当外邦人能做到逆本性的事，你却不能做合乎本性的事，反而连这都放弃了吗？然后，既然他说了\Quote{逆着本性}和\Quote{被接上}，免得你以为犹太人占优势，他又纠正说，他也要被接上。他说：\Quote{何况这些按本性是天然的枝条，岂不更要被接在自己的橄榄树上吗？}又说：\Quote{天主有能力把他们接上。}在此之前他说：如果他们\Quote{仍然不信，他们也要被接上。}当你听见他不断说\Quote{合乎本性}和\Quote{逆本性}时，不要以为他指的是不可变的本性，而是用这些话告诉我们可能和连贯之事，以及不可能之事。因为善与恶不是本性的，而只在于习性和决心。也要看他多么柔和。因为他在说\Quote{你也会被砍下，如果你不信，这些也会被接上，如果他们不再不信}之后，便放下了严厉的一面，坚持更温和的声音，并以之结束，给犹太人摆出很大的希望，只要他们愿意不再如此。因此他接着说：\par

第25节。\BibleQuote{弟兄们，我不愿意你们不知道这奥秘，免得你们自以为聪明。}\BibleRef{罗 11:25}\par

这里所说的奥秘，是指那未知、不可言说、充满惊奇和出乎意料的事。正如他在另一处也说：\Quote{看，我告诉你们一个奥秘：我们不全都要睡眠，而是全都要改变。}\BibleRef{格前 15:51}那么，这奥秘是什么？\par

\Quote{以色列发生了部分的盲目。}这里他再次打击犹太人，同时似乎又压低外邦人。但他的意思大致与先前所说的一样：不信不是普遍的，而只是\Quote{部分的}。正如他说：\Quote{但如果有人使人忧愁，他并没有使我忧愁，只是部分地使我忧愁。}\BibleRef{格后 2:5}所以，他在这里也重复上面所说的话：\Quote{天主并没有弃绝祂预先所知道的子民。}\BibleRef{罗 11:2}又说：\Quote{那么，他们跌倒了，是要他们跌倒吗？绝对不可！}\BibleRef{罗 11:11}他在这里也是说：被拔起的不是整个民族，而是许多人已经相信，还有更多人将要相信。既然他许下了一件大事，他就引证先知来证明，如下所述。他引用这段经文不是为了证明盲目的发生（人人都能看到），而是为了证明他们将相信并得救，他请依撒意亚作证，他高声说：\par

第26节。\BibleQuote{拯救者必从熙雍出来，消除雅各伯的不敬。}\BibleRef{罗 11:26}\BibleRef{依 59:20}\par

然后，为了给出一个标记，将它固定在救恩的意义上，免得有人把它转向并归诸过往的时期，他说，\par

第27节。 \BibleRef{罗 11:27} \BibleQuote{因为这是我的盟约与他们所立的，就是当我除去他们的罪过的时候。}\par

不是在他们受割礼的时候，不是在他们献祭的时候，也不是在他们遵行法律其他行为的时候，而是在他们获得罪过赦免的时候。那么，如果这事是应许的，却从未在他们身上实现，他们也从未藉着圣洗获得罪过的赦免，它必定会发生。因此他接着说：\par

第29节。 \BibleRef{罗 11:29} \BibleQuote{因为天主的恩赐和召叫是不会撤回的。}\par

而且他所说的还不止这些来安慰他们，因为他利用了已经发生的事。而后来发生的事，他陈述为主要的意图，用这些话表达：\par

第28节。 \BibleRef{罗 11:28} \BibleQuote{就福音而论，他们为你们的缘故是仇敌；但就拣选而论，他们因祖宗的缘故是蒙爱的。}\par

这样，外邦人就不致自高自大，说： \Quote{我站住了，不要告诉我本会如何，只告诉我已经如何。} 他用这个考虑来使他谦卑下来，并说： \BibleRef{罗 11:28} \BibleQuote{就福音而论，他们为你们的缘故是仇敌。} 因为当你们被召叫时，他们变得更加诡辩。然而，天主至今也没有中止对你们的召叫，而是等待所有将要相信的外邦人都进来，然后他们也要进来。接着，他为他们行了另一种恩惠，说： \BibleRef{罗 11:28} \BibleQuote{就拣选而论，他们因祖宗的缘故是蒙爱的。} 这又是什么意思呢？因为在他们为仇敌的事上，惩罚是属于他们的；但在他们为蒙爱的事上，祖宗的功德对他们并无影响，如果他们不相信的话。然而，正如我所说，他不停地用言语安慰他们，为要争取他们。为此，作为对他先前断言的新证据，他说：\par

第30至32节。 \BibleRef{罗 11:30-32} \BibleQuote{因为正如你们从前不信天主，如今却因他们的不信而蒙受了怜悯；同样，他们现今也不信，为要藉着你们所得的怜悯，使他们也蒙受怜悯。因为天主把众人都圈在不信之中，为要怜悯众人。}\par

因为他会说： \Quote{你们从前不顺服，} 你们是外邦人，然后犹太人进来了。同样，当这些不顺服时，你们进来了。然而，他们不会永远灭亡。 \BibleRef{罗 11:32} \BibleQuote{因为天主把众人都圈在不信之中，} 就是说，祂说服了他们，显明了他们的不顺服；不是要他们停留在不顺服中，而是为要藉着一方的诡辩拯救另一方，这些藉着那些，那些藉着这些。现在想一想：你们不顺服，他们得救了。再者，他们不顺服，你们得救了。然而，你们得救并非像犹太人那样被重新抛弃，而是为了在你们存留时，藉着嫉妒把他们吸引过来。\par

第33节。 \BibleRef{罗 11:33} \BibleQuote{啊，天主的富饶、智慧与知识的深渊！祂的判断何其难测！}\par

在这里，他回溯了往昔，回顾了天主最初的安排，世界由此存续至今，并思考了祂为一切事情所做的特殊预备，他感到敬畏，大声呼喊，使听众确信他所说的一定会实现。因为若不确实，他不会呼喊并感到敬畏。那么，他知道这是一个深渊，但有多深，他却不知道。因为这是惊奇之人的话语，而非知晓全部之人。他赞叹并敬畏其美善，尽己所能，用两个强有力的词——\Quote{富饶}和\Quote{深渊}——来宣扬它，然后敬畏于祂既有意愿又有能力成就这一切，并藉着对立面产生对立面。 \BibleQuote{祂的判断何其难测！} 因为它们不仅无法理解，甚至连探寻也不可能。 \BibleQuote{祂的道路何其难寻！} 就是说，祂的安排不仅无法知道，甚至连探究也不可能。因为连我，他说，也没有完全找到，只找到一小部分，不是全部。因为惟有祂清楚地知道自己的事。因此他接着说：\par

第34至35节。 \BibleRef{罗 11:34-35} \BibleQuote{因为谁曾知道主的心意？或者谁曾作过祂的谋士？或者谁先给了祂，使祂偿还呢？}\par

他的意思大致是这样：虽然祂如此智慧，但祂的智慧并非来自他人，而是祂自己就是美善的源泉。而且，虽然祂成就了如此伟大的事，赐予了我们如此巨大的恩赐，但祂并非借用他人的东西来赐予，而是使之从自身涌出；也不欠任何人回报，因为祂未曾从任何人领受，而是始终率先施恩；这是富饶的主要标志：丰盈地溢出，却不需要援助。因此他接着说： \BibleRef{罗 11:36} \BibleQuote{因为万物都出于祂、藉着祂、归于祂。} 祂亲自设计，亲自创造，亲自持守（武加大译本作 \OriginalTerm{持守}{\GreekText{συγκρατεἵ}}，手抄本作 \OriginalTerm{形成}{\GreekText{συγκροτεἵ}}）。因为祂是富有的，不需要从他人领受。祂是智慧的，不需要谋士。我说谋士何用？除祂自己，这富有而智慧的一位，无人能知道祂的事。因为外邦人得到如此丰盛的供应，是富有的证明；使犹太人的后辈成为他们的老师，是智慧的证明。然后，他感到敬畏，也在话语中献上感恩： \Quote{愿荣耀归于祂，直到永远。阿们。}\par

因为当他讲述任何这类伟大而不可言喻的事情时，他总会以赞颂辞作为惊叹的结尾。关于圣子，他也是这样做的。因为在那段经文里，他也继续做了同样的事情，就像在这里一样。\BibleQuote{基督按肉身说，也是从他们来的；他是在万有之上，永远应受赞颂的天主。阿们。}\BibleRef{罗 9:5}\par

然后，我们也当效法祂，藉着儆醒的生活方式，在一切事上光荣天主；不要因祖先的德行而自负，要知道犹太人已成了鉴戒。因为基督徒的关系绝不是这样的，他们的亲属关系是属圣神的。因此，\Person{斯基泰人}成了\Person{亚巴郎}的儿子；而他的儿子反而比\Person{斯基泰人}更疏远。那么，我们不要因父亲的善行（大多数手抄本作 \Quote{他人} ）而自信；但若你有一位甚至令人惊叹的父母，也不要以为这就足以拯救你，或为你赢得尊荣和光荣，除非你与他有品格的相似关系。同样，若你有一位坏父亲，也不要以为你会因此而被定罪或蒙羞，只要你妥善安排自己的行为。因为有什么比外邦人更不光彩呢？然而他们在信德中不久就与圣徒成了亲属。又有什么比犹太人更亲近呢？然而他们因不信而成了外人。因为那种亲属关系是本性和必然的，在它之下我们都是亲属。因为我们都出自\Person{亚当}，没有人能比另一个人更亲近，无论是就\Person{亚当}还是\Person{诺厄}而言，或是就大地——众人共同的母亲而言。但那堪受尊荣的亲属关系，是那将我们与恶人区分开来的。因为并非所有人都能以这种方式成为亲属，只有那些具有同样品格的人。我们也不称那些与我们同劳苦的人为兄弟，而是称那些表现同样热忱的人。基督就是这样赐给人天主儿女的名号，同样地，也赐给人魔鬼的儿女、悖逆的儿女、地狱的儿女和丧亡的儿女。同样，\Person{弟茂德}因良善而成为\Person{保禄}的儿子，被称为 \BibleQuote{我自己的儿子} \BibleRef{弟前 1:2} ；但我们甚至不知道他姐姐的儿子的名字。然而那一个是与他有本性的亲属，这对他并无益处。而另一个虽在本性和地域上与他相距甚远（因他是\Place{吕斯特拉}人），却成了最亲近的亲属。那么，让我们也成为圣徒的儿子，或者说，让我们甚至成为天主的儿子。因为成为天主的儿子是可能的，请听祂所说： \BibleQuote{所以你们应当是成全的，如同你们的天父是成全的一样} \BibleRef{玛 5:48} 。这就是为什么我们在祈祷中称祂为父，不仅是为提醒我们恩宠，也是为提醒我们德行，使我们不做任何有损于这种关系的事。有人说，人如何能成为天主的儿子呢？——借着脱离一切情欲，并对那侮辱和亏待我们的人显示温良。因为你的父对待那亵渎祂的人也是如此。因此，虽然祂在不同的时候说不同的话，但祂从未说过你们可以像你们的父，除非当祂说： \BibleQuote{要为那迫害你们的人祈祷，要为那恨你们的人行善} \BibleRef{玛 5:44} ，那时祂才引入这作为赏报。因此，\Person{保禄}在说 \BibleQuote{你们该效法天主} \BibleRef{弗 5:1} 时，就是要他们在这一点上效法。因为我们需要一切善行，但尤其是对人的爱和温良，因为我们自己极度需要祂对人的爱。因为我们每天犯许多过犯。因此我们也需要显示许多慈悲。但多与少不是由所给之物的数量来衡量，而是由给予者的财力。因此，富人不可傲慢，穷人也不可因给得少而沮丧，因为后者往往比前者给得更多。我们不可因贫穷而苦恼，因为贫穷使我们的施舍更容易。因为收集了大量财物的人，会被傲慢以及对他所拥有的（或 \Quote{超乎所拥有的欲望} ）更大的执着所抓住。但只有一点财物的人，却免除了这两种主宰的情欲：因此他找到了更多行善的机会。因为这样的人会欣然进入监狱，探访病人，并给一杯凉水。但另一个人却不承担任何这类职分，因为他被财富所骄纵（ \OriginalTerm{被财富所骄纵}{\GreekText{φλεγμαίνων}} ）。那么，不要因贫穷而灰心。因为你的贫穷使你在天堂的贸易变得更容易。即使你一无所有，但有一颗同情的心，即使这一点也会为你存留赏报。因此，\Person{保禄}也命令我们 \BibleQuote{与哭泣的人一同哭泣} \BibleRef{罗 12:15} ，并劝勉我们要与囚犯如同被捆绑在一起的人 \BibleRef{希 13:3} 。因为这不只是给那些哭泣的人带来安慰，使他们有许多人同情，也是给那些处于其他痛苦环境中的人带来安慰。因为有些情况下，交谈与金钱同样有能力使沮丧的人复原。为此，天主劝勉我们赐给那求我们的人金钱，这不仅是为了救济他们的贫穷，更是为了教导我们同情邻人的不幸。为此，贪婪的人是可憎的，因为他不仅无视那些乞讨的人，而且因为他自己训练（ \OriginalTerm{训练}{\GreekText{ἀλείφεται}} ）成残忍和极不人道。因此，那为他们的缘故而轻视金钱的人，正是因此才成为可爱的，因为他对人慈悲和良善。基督在祝福慈悲的人时，不仅祝福和赞美那些施舍金钱的人，也祝福那些有如此意愿的人。那么，让我们如此倾向于慈悲，其他所有福气都会随之而来，因为一个有爱心和慈悲精神的人，若有金钱，就会施舍；若看见有人受苦，就会哭泣和哀叹；若遇见受冤屈的人，就会为他辩护；若看见有人受虐待，就会伸手相助。因为他有了那福气的宝库——一个有爱心和慈悲的灵魂，他将它变为一个满足所有弟兄需要的源泉，并享有那与天主（费尔德依据四份手抄本作 \OriginalTerm{与天主}{\GreekText{τᾥ} \GreekText{θεᾥ}} ）一起预备的一切赏报。那么，为达到这些，让我们首先按此塑造我们的灵魂。因为这样，我们在现世就会行无数的善举，并享受将来的荣冠。愿我们众人藉着恩宠和爱人之德，达到这一切，等等。\par

\SourceRecord{Translated by J. Walker, J. Sheppard and H. Browne, and revised by George B. Stevens.；Nicene and Post-Nicene Fathers, First Series；Vol. 11.；Edited by Philip Schaff.；Buffalo, NY: Christian Literature Publishing Co.,；1889.；<http://www.newadvent.org/fathers/210219.htm>.}

% structure: p0001 source=h1 target=section reason=page-title

\section{罗马书讲道集 第二十篇}

\BibleRef{罗 12:1}\par

\Quote{我因天主的仁慈，劝你们、弟兄们，将你们的身体献作生活的、圣洁的、中悦天主的祭献，这才是你们合理的敬拜。}\par

在广泛论述天主对人的爱，指出祂对我们不可言喻的关怀和无法测度的良善之后，他接着提出这一点，为要说服那些已蒙受恩惠的人，表现出与这恩赐相称的生活。尽管他是如此伟大而良善的人，他却不拒绝恳求他们，而且不是为他自己可能获得的任何快乐，而是为他们将得到的益处。为什么惊讶他不拒绝恳求，甚至将天主的仁慈摆在他们面前？因为，他的意思是，你们既是从这仁慈获得无数祝福，就当敬畏它们，被它们所感动而生怜悯。因为它们本身采取恳求者的姿态，要求你们不做出有愧于它们的行为。那么，他的意思是，我以你们借以得救的这一切恳求你们。就好像有人希望使一个曾受大恩的人表示敬意，就带着恩人亲自来作恳求者。你恳求什么？请听：\BibleQuote{你们要将自己的身体献上，当作生活的祭品，圣洁的，中悦天主的，这才是你们合理的敬礼。} 因为他提到祭品时，为防止有人认为他叫他们自杀，他立即加上（希腊语顺序）\BibleQuote{生活的}。然后为与犹太人的区别，他称它为\BibleQuote{圣洁的，中悦天主的，你们合理的敬礼}。因为他们的祭品是物质的，也不太蒙悦纳。因为祂说：\BibleQuote{谁向你们的手要求了这些？}\BibleRef{依 1:12} 而且在其他多处经文中祂明显弃绝它们。因为祂所期望的献祭不是这个，而是这个与其他。因此他说：\BibleQuote{颂赞的祭品将光荣我。} 又说：\BibleQuote{我要以诗歌赞美我天主的名，这要比献上公牛犊、有角有蹄的更中悦祂。}\BibleRef{咏 49:23} 另外一处祂拒绝说：\BibleQuote{难道我吃牛肉，或喝山羊的血吗？}\BibleRef{咏 49:13} 接着又说：\BibleQuote{你当向天主奉献颂赞的祭，并向上至高者还你的誓愿。}\BibleRef{咏 49:14} 同样，保禄也在这里吩咐我们\BibleQuote{将你们的身体献上，当作生活的祭品}。怎么可以说身体成为祭品呢？让眼睛不看邪恶的事，它就成了祭品；让你的舌头不说污秽的话，它就成了奉献；让你的手不做不法的事，它就成了全燔祭。或者这还不够，我们还必须有善工：让手行施舍，口祝福那些反对自己的人，耳朵常闲于圣经诵读。因为祭品不允许有任何不洁之物：祭品是其它行为的初果。那么让我们从手、足、口和所有其他肢体，向天主献上初果。这样的祭品是蒙悦纳的，而犹太人的祭品甚至是不洁的，因为经上说：\BibleQuote{他们的祭品对他们好像丧饼。}\BibleRef{欧 9:4} 我们的祭品却不是这样。那祭品使祭物死亡，这个却使祭物活着。因为当我们克治了我们的肢体，我们才能活。因为这祭品的法是新的，所以火的种类也是奇妙的。因为它不需要柴或物质垫底；我们的火自生，不焚毁牺牲，反而使它活起来。这就是天主自古寻求的祭献。因此先知说：\BibleQuote{天主所要的祭献就是破碎的精神。}\BibleRef{咏 50:17} 三个青年也献上这个，当他们说：\BibleQuote{在此时期，没有君王，没有先知，没有首领，没有全燔祭，没有地方在祢面前献祭，以求得祢的怜悯。但让我们怀着忏悔的心和谦卑的精神蒙祢悦纳。}\BibleRef{达 3:15} 还要观察他用词多么精确。因为他不是说，\OriginalTerm{作}{\GreekText{ποιήσατε}}\BibleRef{出 29:39}你们的身体为祭品，而是\BibleQuote{献上}（\OriginalTerm{献上}{\GreekText{παραστήσατε}}）它们，好像说，你们不再对它们有任何利益。你们已将它们交给别人。因为甚至那些提供（同一词）战马的人也不再对它们有兴趣。你也已献上你的肢体为与魔鬼作战和那可怕的阵列。不要把它们用于自私的用途。他也从这指出另一件事：如果一个人想要献上它们，就必须使它们被认可。因为我们不是向任何凡人献上，而是向天主，宇宙的君王；不仅是为战争，而且是为让君王亲自坐在上面。因为祂不拒绝甚至坐在我们的肢体上，反而极大渴望。任何与我们同为仆人的君王都不愿做的事，天使的主却选择做。既然它既是献上（即为君王使用）又是祭品，就要除去每一点瑕疵，因为如果有瑕疵，就不再是祭品。因为淫荡的眼睛不能献祭，贪婪的手不能献上，跛足的脚和去戏院的脚不能献上，纵欲的肚腹、燃起情欲的肚腹不能献上，充满愤怒和妓女之爱的心不能献上，吐出污秽之言的舌头不能献上。因此我们必须从各方面查看我们身体上的瑕疵。因为如果那些献旧祭的人被吩咐从各方面查看，并且不被允许献上\BibleQuote{有任何多余或缺少，或是疥癣的，或是癞疮的}\EditorialNote{参阅肋未纪 22:22-23}的动物，那么我们这些不是献无知牲畜而是献自己为祭的人，必须更加严格，在各方面纯洁，以便我们也能像保禄一样说：\BibleQuote{我现在已被奠祭，我离世的时期已经近了。}\BibleRef{弟后 4:6} 因为他比任何祭品更纯洁，所以他自称\BibleQuote{已被奠祭}。但这将实现，如果我们杀死老亚当，如果我们克治在地上的肢体，如果我们把世界钉在十字架上。这样我们不再需要刀、祭台或火，或者更确切地说，我们都需要这些，但不是人手所造的，而是这一切都从上面降临我们，火从上面，刀也从上面，我们的祭台将是广大天空。因为当厄里亚献上可见的祭品时，一团从上而来的火焰吞没了水、木柴和石头，更何况这将成就在你身上。如果你内心有任何松懈和世俗，仍怀着善意献祭，圣神的火将降临，既消除那世俗性，又成全（田野本：手抄本作\Quote{提升}）整个祭品。但什么是\BibleQuote{合理的（\OriginalTerm{合理的}{\GreekText{λογική}}）敬礼（\OriginalTerm{敬礼}{\GreekText{λατρεία}}）}？它意指属灵的事奉，符合基督的生活。正如在天主家中服务并执行礼仪的人，无论他是谁，都会收敛自己（\OriginalTerm{收敛}{\GreekText{συστέλλεται}}\BibleRef{则 44:19}），变得更庄严；所以我们一生都当以服务与事奉为念。如果你每天向他献上祭品（三份手抄本作\Quote{自己为祭品}），并成为你自己身体的司祭和灵魂美德的司祭，例如当你奉献节制、施舍、良善和忍耐时，这事就会实现。因为这样做你就献上了\BibleQuote{合理的敬礼}——即没有肉体、粗俗、可见之物。那么，藉着所赐的名号提升了听者，并表明每个人因行为而成为自己肉身的司祭之后，他也提到了我们实现这一切的方法。是什么方法呢？\par

第2节。\BibleQuote{你们不可与此世同化，反而应以更新的心思变化自己。}\BibleRef{罗 12:2}\par

因为这个世界的形态是卑微而无价值的，但是暂时的，既无崇高，也无恒久，也无正直，全是扭曲的。如果你愿意正直行走（或正确，\OriginalTerm{正直}{\GreekText{ὀρθά}}），就不要效法这现世生活的形态。因为在其中没有持久或稳固之物。这就是为什么他称之为一种\Emph{形态}（\OriginalTerm{形态}{\GreekText{σχῆμα}}）；并且在另一处经文中，\BibleQuote{这个世界的形态正在消逝。}\BibleRef{格前 7:31}因为它没有耐久性也没有固定性，其中所有都只是暂时的；因此他称之为这世代（或世界，\OriginalTerm{世代}{\GreekText{αἰών}}）。\par

\SourceRecord{Translated by J. Walker, J. Sheppard and H. Browne, and revised by George B. Stevens.；Nicene and Post-Nicene Fathers, First Series；Vol. 11.；Edited by Philip Schaff.；Buffalo, NY: Christian Literature Publishing Co.,；1889.；<http://www.newadvent.org/fathers/210220.htm>.}

% structure: p0001 source=h1 target=section reason=page-title

\section{《罗马书》讲道集第二十一篇}

\BibleRef{罗 12:4-5}\par

\BibleQuote{就如我们在一个身体上有许多肢体，但每个肢体都有不同的职务；同样，我们众人在基督内，也是一个身体，彼此之间，每个都是肢体。}\par

他又用对格林多人所讲的同一个比喻，并且是为了平息同样的情绪。因为这剂良药的力量，以及这个比喻对矫正高傲之病的效力是巨大的。他说：你为什么自我高抬？或者，为什么另一个人又完全鄙视自己？我们不是所有人，不论大小，都是一个身体吗？既然我们在总数上只是一个，而且彼此互为肢体，你为什么因高傲而分离自己？你为什么羞辱你的弟兄？因为他是你的肢体，你也是他的肢体。恰好在这点上，你们获得尊荣的机会是相等的。因为他陈述了两件事，可以压下他们高傲的精神：一是我们彼此互为肢体，不仅小的对大的，也有大的对小的；二是我们都是一个身体。或者不如说有三点，因为他显示这恩赐是出于恩宠。\Quote{因此，不要高傲。}因为是天主赐给你的；你没有取来，甚至没有找到它。因此，当他提到恩赐时，他并没有说一个人领受的更多，另一个人更少，而是说什么？\Emph{不同的}。因为他说的正是：\Quote{既有不同的恩赐}，不是更少和更大，而是\Quote{不同}。即使你没有担任同样的职务，身体仍然是同一个。而且从恩赐开始，他以善行结束；所以在提到预言、服务等之后，他以怜悯、勤勉和帮助作结。既然有些人可能有德行，却没有预言，他便显示这也是一种恩赐，而且比另一种大得多（如他在致格林多人书中所显示的），并且大得多，因为一种有赏报，另一种却没有报酬。因为整个事情都是出于恩赐和恩宠。因此他说：\par

\Quote{既有不同的恩赐，按那赐给我们的恩宠：或说预言，就该按信德的比例说预言。} \BibleRef{罗 12:6}\par

既然他已经充分安慰了他们，他也想使他们彼此竞争，并更认真地努力，通过显示是他们自己为领受多或少提供了基础。因为他说这确实是天主所赐的（如他所说：\Quote{按天主分给各人信德的尺度}；又说：\Quote{按那赐给我们的恩宠} \BibleRef{罗 12:3}），为的是制服高傲者。但他也说开端在于他们自己，以激励懒惰者。他在致格林多人书里也这样做，以激起这两种情绪。因为当他说：\Quote{你们该热切追求恩赐} \BibleRef{格前 12:31}，他显示他们自己是所赐之物差异的原因。但当他说：\Quote{这一切都是唯一且同一圣神所运作的，随他的意愿，个别地分给每人} \BibleRef{格前 12:11}，他是在证明领受者不应自高自大，用尽一切可行的方法平息他们的混乱。他在这里也这样做。再者，为了唤醒那些昏昏欲睡的人，他说：\Quote{或说预言，就该按信德的比例说预言。}因为虽然它是恩宠，却不是随意倾倒，而是按照领受者的容量调整其尺度，它流出的多少取决于它所发现带去的信德容器的大小。\par

\Quote{或服务，就该专务服务。} \BibleRef{罗 12:7}\par

这里他提到一个综合性的事物。因为连宗徒职务也称为服务，而且每一个属神的工作都是服务。这确实是一个特别职务的名称（即执事职），不过这里是一般意义。\Quote{或作教导的，就该专务教导。}看他何等不经意地将它们并列，小的在先，大的在后，再次给我们同样的教训：不要自高或骄傲。\par

\Quote{或作劝慰的，就该专务劝慰。} \BibleRef{罗 12:8}\par

这也是一种教导的方式。因为经上说：\Quote{若有任何劝勉的话，就向百姓讲说。}\BibleRef{宗 13:15}\newline{}然后，为表明若不以合宜的法则行善，则追求德行并无大益，他接着说：\Quote{分施　\OriginalTerm{分施}{\GreekText{μεταδιδοὺς}} 的，要本着纯朴去做。}因为仅仅分施还不够，还须慷慨大方，这正与纯朴之名相符。因为连那些童贞女也拥有油，但因不够充足，便被全然抛弃。\Quote{管理　\OriginalTerm{管理}{\GreekText{προἵστάμενος}} 的，要殷勤；}因为仅承担管理之责不够。\Quote{怜悯人的，要喜乐。}因为仅施怜悯不够，还须以宽宏和不吝啬的精神行善，更确切地说，不仅要不吝啬，还要以喜乐和欢欣的态度去做，因为不吝啬并不等于喜乐。这一点，他在致格林多人书函中也曾极力强调。为激励他们如此慷慨，他说：\Quote{少种的，也少收；多种的，也多收。}\BibleRef{格后 9:6}但为纠正他们的心态，他又补充道：\Quote{不要出于勉强，也不要出于不得已。}\BibleRef{格后 9:7}因为怜悯的雨露应当兼具不勉强与愉悦。你为何在施舍时自叹自怜？\OriginalTerm{亚里士多德《尼各马可伦理学》卷二第3章、卷四第1章}{Eth. N. ii. 3 and iv. 1}为何在施怜悯时忧伤，从而丧失善行的益处？因为若你忧伤，就不是在施怜悯，而是残忍无情。若你忧伤，怎能扶起那忧伤的人呢？即使你喜乐地施与，那人若能不疑有恶已是万幸。因为对人们而言，没有比受人施舍更觉耻辱的了；除非你以极其喜乐的容色消除他的疑虑，表明你是在领受而非施与，否则你反而会摧垮受施者，而非扶起他。故此他说：\Quote{怜悯人的，要喜乐。}因为那领受王国的人，岂会有忧伤的面容？那领受罪赦的人，岂会继续沮丧？不要计较金钱的付出，而要看付出所带来的增长。因为播种的人尚且因播种而喜乐，尽管收成未定；那耕种天国的人岂不更应如此？因为这样，即使你付出很少，也会成为很多；同样，无论你付出多少，若面带愁容，也会使你的许多变为很少。如此，那寡妇以两个小钱胜过了许多银子，因为她的心是慷慨的。\newline{}或许有人说：一个生活在极度贫穷中、倾其所有的人，怎能以乐意的心这样做呢？问问那寡妇，你便会知道途径，也会明白使人拮据的不是贫穷，而是人的心性，它同时造成拮据与富足。因为人在贫穷中也可以慷慨　\OriginalTerm{慷慨}{\GreekText{μεγαλόψυχον}}，在富足中也可以吝啬。因此，他在施与时要求纯朴，在怜悯中要求喜乐，在辅助　\OriginalTerm{辅助}{\GreekText{προστασίαν}} 中要求殷勤。因为他不仅希望我们用金钱去帮助有需要的人，也包括言语、行为、亲身以及其他一切方式。在提及辅助的主要形式（即在教导和劝勉方面，因为这是更必要的，能滋养灵魂）之后，他进而论及金钱及所有其他方式的辅助；然后为表明如何正确实行这些事，他引入了这些事的母亲——爱德。\par

第九节：\Quote{爱德不可虚假。}\BibleRef{罗 12:9}\par

你若拥有这爱德，就不会感到金钱的损失、身体的劳苦、言语的辛劳、服务的麻烦；你必勇敢地承担一切，无论是以身体、金钱、言语或任何其他方式帮助邻人。因此，他不仅要求施与，还要纯朴；不仅要求辅助，还要殷勤；不仅要求施舍，还要喜乐；同样，他也要求爱德不止于此，还要不虚假。因为这才是爱德的本色。人若拥有这爱德，其余一切都随之而来。因为施怜悯的人必喜乐地去做（他在施与自己），辅助的人必殷勤地去做（他在辅助自己），分施的人必慷慨地去做（他在赐予自己）。然而，既然也存在对恶事的爱——如放纵者的爱、为钱财为掠夺而同心的爱、以及为狂欢宴饮的爱——他借以下的话清除这一切：\Quote{恶要深恶痛绝（\OriginalTerm{深恶痛绝}{\GreekText{ἀποστυγοὕντες}}）。}他不是说戒避恶事，而是说憎恨它，不仅憎恨，且要极其憎恨。因为\OriginalTerm{加强前缀}{\GreekText{ἀπὸ}}在他那里常具有强调作用，如他说\Quote{热切期待、切望（完全）救赎}时。因为许多人虽不做恶事，却仍对恶事有欲望，故此他说\Quote{深恶痛绝}。因为他所要的是净化意念，使我们对恶事怀有极大的敌意、仇恨与争战。他的意思是：不要因我说\Quote{彼此相爱}，就以为我允许你们在恶事上也彼此合作；我所立的诫命恰恰相反。因为它要求你不仅远离恶行，甚至要远离对恶事的倾向；不仅如此，还要对那倾向有极度的厌恶和憎恨。他不仅满足于此，还引入德行的实践：\Quote{善要持守。}\par

他不仅谈到行动，也谈到心态。因为\Quote{貼緊}这命令就指明了这一点。当日天主将男人与女人结合时说：\BibleQuote{因為他要與妻子結合。}\BibleRef{创 2:24}。然后他提出我们应该彼此相爱的理由。\par

第10节。\BibleQuote{要以兄弟之爱彼此亲切相待。}\BibleRef{罗 12:10}\par

你们是弟兄，他说，且出自同一产痛。因此即便在这点上你们也应该彼此相爱。梅瑟对那些在埃及争吵的人说过：\BibleQuote{你们是弟兄，为什么彼此欺负？}\BibleRef{出 2:13}。那么当他谈到外人时，他说：\BibleQuote{若是可能，尽量与众人和睦相处。}\BibleRef{罗 12:18}。但当他谈到自己的人时，他说：\Quote{要以兄弟之爱彼此亲切相待。}因为在前边他要求避免争吵、仇恨和厌恶，但这里还要求爱，且不仅是爱，而是亲人之爱。因为一个人的爱不仅必须\Quote{没有伪善}，而且必须强烈、温暖、炽热。因此他说：彼此亲切相待，也就是要做朋友，而且是热心的朋友。不要等待被爱，而要主动去爱，并率先开始。因为这样你也会收获他的爱的报酬。既提出了应该彼此相爱的理由，他也告诉我们这种情感如何能持久不变。因此他接着说：\Quote{在尊敬上彼此谦让。}因为这是产生爱的方式，也是爱产生之后持久的方式。没有什么比努力在尊敬上胜过邻舍更能使人成为朋友。因为他先前所说的来自爱，而爱来自尊敬，正如尊敬也来自爱。那么，我们不但要尊敬，他还要求更多，他说：\par

第11节。\BibleQuote{不迟延的热诚。}\BibleRef{罗 12:11}\par

因为当我们尊敬并显出保护之热诚时，这也生出爱：没有什么比尊敬和关怀更使人蒙爱。因为爱是不够的，还必须加上这个；或者说，这个也来自爱，正如爱也从之获得温暖，二者相互巩固。因为有许多人在心里爱，却没有伸出援手。这就是他用各种方法建立爱的原因。我们如何能成为\Quote{不迟延的热诚}？\par

\Quote{在圣神内火热。}请看他在每一点上都追求更高的程度；因他不仅说\Quote{施舍}，还说\Quote{慷慨}；不仅说\Quote{治理}，还要\Quote{殷勤}；不仅说\Quote{怜悯}，还要\Quote{愉快}；不仅说\Quote{尊敬}，还要\Quote{彼此谦让}；不仅说\Quote{爱}，还要\Quote{没有伪善}；不仅说远离\Quote{邪恶}，还要\Quote{憎恨}；不仅说持守\Quote{良善}，还要\Quote{貼緊}；不仅说\Quote{爱}，还要\Quote{以兄弟之爱}；不仅说要热诚，还要\Quote{不迟延}；不仅要有\Quote{圣神}，还要\Quote{火热}，即要你温暖而警醒。因为如果你拥有上述之事，你必吸引圣神到你身边。如果圣神与你同在，祂也会使你为这些目的成为良善，而一切因着圣神和爱变得容易，你被两面加热。你们难道没有看见那背上载着火焰的公牛（汉尼拔，见李维《罗马史》22.16）如何无人能抵挡？同样，如果你取得这两种火焰，你也会使魔鬼难以承受。\BibleQuote{事奉主。}因为通过这些方式你都能事奉天主；因为你为你的弟兄所做的一切都归于你的主人，而祂会看作自己受了恩惠，从而按此报答你。请看他把行这些事之人的心灵提升到何等高度！然后，为显示圣神之火如何点燃，他说：\par

第12节。\BibleQuote{在希望中喜乐，在困苦中忍耐，恒切祈祷。}\BibleRef{罗 12:12}\par

因为这一切都是那火的燃料。当他要求金钱的支出、个人的辛劳、治理、热诚、教导以及其他劳苦的工作后，他又用爱、用圣神、借着希望来装备这斗士。因为没有什么比美好的希望更能使灵魂勇敢无畏。然后，甚至在所希望的善事之前，他又给予另一个报酬。既然希望是关于未来的事，他说：\Quote{在困苦中忍耐。}而且在未来之事以前，在今世生活中，你从困苦中获得一个大善\BibleRef{罗 5:4}，即成为坚忍和经过考验的人。此后他又提供另一个帮助，说：\Quote{恒切祈祷。}那么，既然爱使事情容易，圣神协助，希望减轻负担，困苦使你坚忍而能够高尚地承受一切，而且你又有另一个极大的武器，即\Quote{祈祷}及祈祷带来的辅助，那么他所命令的还有什么困难呢？显然没有。你看见他如何从各方面给斗士稳固的立足点，并显示这些命令完全容易。再请注意他如何为施舍辩护，或者说，不是绝对的施舍，而是向圣徒施舍。因为上面他说\Quote{怜悯的要愉快}时，他使我们慷慨对待众人。但这里他是为信徒说话。所以他接着说：\par

第13节。\BibleQuote{分担圣徒的\OriginalTerm{必需}{\GreekText{χρείαις}}（或作\OriginalTerm{记念}{\GreekText{μνείαις}}）。}\BibleRef{罗 12:13}\par

他没有说\Quote{赐予，}而是说\Quote{分担圣人的急需，}表明他们领受的比给予的多，这是一笔交易，因为它是共融。你带进金钱？他们带给你在天主面前的胆量。\Quote{致力于（希腊文：pursuing）款待。}他没有说\Quote{做}，而是说\Quote{致力于}它，以此教导我们不要等待那些请求的人，看他们何时会到我们这里来，而是要奔向他们，并致力于寻找他们。\par

罗特如此行，亚巴郎也如此行。因为他整日在此等待这美好的猎物，一看见便扑上去，跑去迎接他们，俯伏在地朝拜，并说：\BibleQuote{我主，若我在祢眼前蒙恩，求祢不要离开祢的仆人。}\BibleRef{创 18:3} 我们却不如此。若我们偶然看见一个陌生人或穷人，便皱起眉头，甚至不屑与之交谈。即使在上千次恳求后我们心软了，吩咐仆人给他们一点小钱，便自以为尽了本分。但他并不如此，反而摆出恳求者和仆人的姿态，尽管他并不知道要接待的是谁。而我们却明知所接待的是基督，竟仍不因此变得温和。他既恳求、哀求，又向他们下跪；我们却侮辱来到我们这里的人。他亲自与妻子一同做这一切，我们却连差遣仆人也做不到。但若你愿意看看他摆在客人面前的餐桌，那里你也会看到极大的慷慨，但这慷慨并非来自财富的充裕，而是来自乐意之心的丰盈。那时岂不是有许多富人？却无人如此行。以色列中有多少寡妇？却无人款待厄里亚。厄里叟的时代又有多少富人？但只有叔能女子收获了款待的果实；亚巴郎也是如此，除了他的宽宏和乐意之心外，尤其值得钦佩的是，他并不知道来者是谁，却仍如此行事。你也不可过于好奇：因为你是为基督而接待他。你若总是如此多虑，往往会错过一位可敬之人，从而失去从他应得的赏报。而且，接待一个不值得尊敬的人，也无可指责，反而可得赏报。因为\BibleQuote{因先知的名接待先知的，必得先知的赏报。}\BibleRef{玛 10:41} 但若因这不合时宜的多虑而错过一个值得钦佩的人，反而要受惩罚。因此，不要忙于过问别人的生活和行为。为了一块饼而计较一个人的全部生平，这是吝啬的极致。因为即使他是个杀人犯、强盗或别的什么，难道你便认为他不配得一块饼和几个小钱吗？你的主尚且使太阳照在他身上！你竟判断他连一日之食也不配吗？我再给你说另一情况。即使你确实知道他罪孽深重，你也没有理由剥夺他今日的饮食。因为你是那说\BibleQuote{你们不知道你们是什么精神}的那一位的仆人。\BibleRef{路 9:55} 你是那位医治了用石头打祂的人，甚或为这些人被钉十字架的那一位的仆人。你不要对我说他杀了别人，即使他正要杀你本人，你也不可在他饥饿时忽略他。因为你是那位渴望拯救那些钉死祂的人的门徒，祂在十字架上说：\BibleQuote{父啊，赦免他们，因为他们不知道他们所做的。}\BibleRef{路 23:34} 你是那位医治了打祂的人、在十字架上还加冕了曾轻视祂的人的那一位的仆人。有什么能与此相比？因为两个强盗起初都曾轻视祂。但祂仍为其中之一打开了乐园。祂哀哭那些将要杀害祂的人，看见叛徒时感到忧愁和困惑，不是因为祂将要被钉十字架，而是因为那人丧亡了。祂之所以忧愁，是因为预先知道那人的悬吊以及悬吊后的惩罚。虽然祂明知他的邪恶，仍容忍他到最后一刻，没有推开叛徒，反而亲吻了他。你的主亲吻并接纳那正要流祂宝血的人。而你竟认为穷人连一块饼都不配，也不尊重基督所订立的律法？祂借此表明，我们不仅不应远离穷人，甚至不应远离那些要置我们于死地的人。因此，不要对我说某人如何深深地伤害了我，只要想想基督在十字架近旁做了什么：祂想用亲吻来挽回那就要出卖祂的叛徒。请看这些话有多大的能力使他羞愧。因为祂说：\BibleQuote{犹达斯，你用亲吻来出卖人子吗？}\BibleRef{路 22:48} 有谁是他不能软化的？有谁听了这番话会不屈服？什么野兽？什么金刚石？然而那可怜的人却不为所动。因此，不要说某人杀了某人，所以我远离他。即使他正要拿剑刺你，要把手插入你的脖子，也要亲吻这只手！因为基督亲吻了那导致祂死亡的嘴！所以，你也不要仇恨，而要哀哭和怜悯那图谋害你的人。这样的人值得我们怜悯和流泪。因为我们是那甚至亲吻叛徒的仆人的仆人（我不停止反复强调这一点），祂对叛徒所说的话比亲吻更温柔。祂甚至没有说：\Quote{你这污秽恶毒的叛徒，你就是这样回报这么大的恩惠吗？}祂说的是什么话？\BibleQuote{犹达斯}，叫他的名字，更像是一个哀伤的人呼唤他、挽回他，而不是愤怒地对待他。祂没有说\Quote{你的老师、你的主、你的恩主}，而是说\Quote{人子}。因为即使祂既不是老师也不是主人，但祂对你如此温柔、如此真诚，以至于在背叛的时候还亲吻你，而这亲吻正是背叛的记号；难道你竟对祂扮演叛徒的角色吗？主啊，祢是有福的！祢赐给了我们何等谦卑、何等忍耐的榜样！祂对叛徒如此。对那些拿着棍棒刀剑来的人，不也是这样吗？有什么比祂对他们说的话更温柔？因为祂有能力一瞬间消灭他们所有人，却丝毫不做，反而以劝诫的方式（\OriginalTerm{劝诫地}{\GreekText{ἐντρεπτικῶς}}）对他们说：\BibleQuote{你们带着刀棒出来，如同对付强盗吗？}\BibleRef{玛 26:55} 祂使他们退后跌倒（\BibleRef{若 18:6}），他们仍麻木不仁，祂便主动交出自己，任由他们用锁链捆绑祂圣洁的手，而祂本有能力立刻摧毁万物、推翻他们。但连这之后，你仍凶狠地对待穷人吗？即使他犯了万种罪恶，匮乏和饥饿也足以软化任何刚硬的心灵。你却变得残忍，模仿狮子的暴怒。然而狮子从不吃死尸。你见他因苦难而如同尸体（\OriginalTerm{如同尸体}{\GreekText{τεταριχευμένον}}），却在他跌倒时扑向他，用侮辱撕碎他的身体，兴起风暴接风暴，使他在逃往避难所时触礁，造成比海难更悲惨的沉船。你怎敢向天主说\Quote{可怜我}，求祂赦免你的罪，而你却傲慢地对待一个没有犯罪的人，追问他饥饿和急需的原因，以你的残酷使所有野兽相形见绌？因为野兽是因肚腹的驱使才夺取所需食物。你却在没有任何推动或驱使的情况下，吞噬你的兄弟，用你的言语啃咬、撕裂他，即使不是用牙齿，也用比牙齿更锋利的言语。那么，当你用舌头染满人血之后，怎敢领受神圣的祭品（\OriginalTerm{祭品}{\GreekText{προσφορὰν}}）？怎敢用充满战争的嘴给予平安之吻？又怎能享用普通食物，当你积聚如此多的毒液时？你不解除贫穷，为何使它更加沉重？你不扶起跌倒的人，为何还要推倒他？你不消除沮丧，为何反而增加？你不施舍钱财，为何还要加上侮辱的话？难道你没有听见不喂养穷人要受什么惩罚？他们被判什么刑罚？因为祂说：\BibleQuote{离开我，到那为魔鬼和他的使者预备的永火里去。}\BibleRef{玛 25:41} 如果那些不喂养的人已被如此定罪，那么那些不仅不喂养，还加以侮辱的人，要受什么刑罚？他们将受什么惩罚？什么地狱？为避免我们自招如此巨大灾祸，趁着我们尚有能力，让我们也改正这恶习，给舌头加上辔头。让我们不仅不侮辱，反而用言语和行为邀请他们，好为自己积蓄许多慈悲，以获得所应许的福分。愿天主以其恩宠和爱人之心，使我们众人皆能获得，等等。\par

\SourceRecord{Translated by J. Walker, J. Sheppard and H. Browne, and revised by George B. Stevens.；Nicene and Post-Nicene Fathers, First Series；Vol. 11.；Edited by Philip Schaff.；Buffalo, NY: Christian Literature Publishing Co.,；1889.；<http://www.newadvent.org/fathers/210221.htm>.}

% structure: p0001 source=h1 target=section reason=page-title

\section{罗马人书释义 第22讲}

\BibleRef{罗 12:14}\par

\BibleQuote{祝福那迫害你们的人；只祝福，不可诅咒。}\par

在教导他们应当如何彼此相处，并将肢体紧密地结合为一个身体之后，他接着带领他们外出作战，并且自此使战斗变得更容易。因为一个人若不能妥善处理自己内部的事，就会在安排与外人相关的事上遇到更多困难；同样，若有人在自己人中间恰当操练，就更易在外人面前占优势。因此，保禄在继续前行的旅程中，在这一点之后放置另一点，并说：\BibleQuote{祝福那迫害你们的人。}他没有说：不要怀恨或报复，而是要求远胜于此的事。因为那是一个明智之人所能做到的，但这完全是天使的份。在说了\Quote{祝福}之后，他又接着说：\Quote{不可诅咒}，免得我们两者并行，而不仅仅做前者。因为那些迫害我们的人，乃是为我们提供赏报的人。但如果你头脑清醒，在那一赏报之外，还有另一赏报，是你自己可获得的。因为前者（迫害者）会因迫害而给你赏报，而你自己却因祝福他人而获得另一赏报，因为你表现了爱基督的一个非常伟大的标记。因为那诅咒自己迫害者的人，表明他并不十分乐意为了基督而受苦；反之，那祝福的人则显示他爱的高深。因此，不要辱骂他，好使你自己获得更大的赏报，并教导他：这事情是出于甘愿，而非出于勉强，是喜庆与节日，而非灾难或沮丧。因此，基督自己曾说过：\BibleQuote{当人们捏造一切坏话毁谤你们时，你们要欢喜踊跃。}\BibleRef{玛 5:11}因此，宗徒们不仅因受毁谤，而且因受鞭打而欢喜地从公议会回来。\BibleRef{宗 5:40}因为除了我所提到的之外，还有另一收获，且是不小的，即你届时会使你的对手羞愧，并以你的行动教导他们：你正在奔赴另一种生命；因为如果他看见你因受苦而欢欣鼓舞（\OriginalTerm{欢欣鼓舞}{\GreekText{πτερούμενον}}），他将从这些行动清楚看到，你拥有超越此世生活的更大希望。因此，如果你不如此行，反而哭泣哀叹，他怎能从那里学到你期待另一种生命呢？除此之外，你还会成就另一件事。即如果他看见你因所受的侮辱而不烦恼，反而祝福他，他将停止折磨你。看，由此产生多少益处：为你自己更大的赏报、更少的试探，他将停止迫害你，天主也将受光荣；而那位在错误中的人，你的忍耐将成为虔诚的教导。为此，他不仅命令我们以相反的态度回报那些侮辱我们的人，甚至包括那些迫害我们和虐待我们的人。现在他命令他们祝福，但随着他的进展，他还劝勉他们要以行动施惠于人。\par

\BibleQuote{与喜乐的人一同喜乐，与哭泣的人一同哭泣。}\BibleRef{罗 12:15}\par

既然可以祝福而不诅咒，却可能不是出于爱，他希望我们整个都被友谊的温暖所渗透。这就是为什么他接着说这些话，使我们不仅祝福，而且还要在他痛苦和受苦时同情他们，每当我们看见他们陷入困境时。是的，有人会说，但为何他吩咐他们加入悲伤者的悲伤？但为何他又命令另一件事，这并非如此重要呢？是啊，但后者更要求基督徒高尚的品性：与喜乐的人一同喜乐，比与哭泣的人一同哭泣更难。因为天然本性本身就完全实现了这一点：没有铁石心肠到不为遭难者哭泣的。但是，前者却需要一个非常高尚的灵魂，不仅能保持不嫉妒，甚至能对受尊敬的人感到高兴。这就是为什么他将它放在前面。因为没有什么比彼此分享喜乐和痛苦更牢固地维系爱。所以，不要因为你自己远离困境，就置身事外而不去同情。因为当你的邻人受虐待时，你应当将灾难当作自己的。那么，要分担他的眼泪，好减轻他的低落情绪。分担他的喜乐，好使喜乐深深扎根，并使爱紧固，这样做对你自己比对邻人更有益处；借你的哭泣，你使自己变得慈悲；借感受他的喜乐，你洗净了自己的嫉妒和吝啬。让我提醒你注意保禄的体贴。因为他不说：结束灾难，免得你常说（或\OriginalTerm{常常}{\GreekText{πολλάκις}}）这是不可能的：而是命令了更容易的事，即你力所能及的事。因为即使你不能消除祸患，也要贡献眼泪，这样你就拿走了一半的痛苦。如果你不能增加一个人的幸福，就贡献喜乐，这样你就为它做出很大的增加。因此，他劝勉我们的不仅仅是远离嫉妒，而是一件更大的事，即加入欢乐。因为这比不嫉妒大得多。\par

\BibleQuote{彼此同心。不要思想高傲的事，而要俯就卑微的人。}\BibleRef{罗 12:16}\par

此处他再次强调谦卑的心，这是他开始劝勉时所提出的主题。因为那时他们很可能因自己的城市（参见第343页）以及其他种种原因而充满高傲；因此他不停地\OriginalTerm{抽出}{\GreekText{ὑποσύρει}}（两份抄本作\OriginalTerm{聚拢}{\GreekText{ὐπορύττει}}）病态的东西，减低炎症。因为没有什么比虚荣更能使教会分裂。他所说的\Quote{你们要彼此同心}是什么意思呢？若有一个穷人进入你家，你要在态度上与他相似，不要因自己的财富而摆出任何不寻常的浮夸气派。在基督内并没有富人与穷人之分。那么，不要因他外表的衣着而羞于接纳他，而要因他内在的信德而接纳他。如果你见他忧愁，不要不屑于安慰他；如果你见他顺遂，也不要羞于分享他的快乐并与他同乐，反而要在他身上存有你在自己身上所存的同一心意。因为经上说：\Quote{你们要彼此同心。}举例来说，如果你自视为大人物，那么也要同样看别人。你怀疑他卑贱渺小吗？那么，你就对自己也下同样的判断，并抛开一切不平等。这如何做到呢？靠你抛开那鲁莽的性情。因此他接着说：\Quote{不要思想高傲的事，反要俯就卑微的人。}也就是说，你要降低自己到他们谦卑的境况，与他们交往，与他们同行，不仅在心里谦卑，还要帮助他们，向他们伸出援手，不是借他人之手，而是亲身，如同父亲照顾孩子，如同头照顾身体。正如他在另一处所说：\BibleQuote{与那些被捆绑的人一同被捆绑。}\BibleRef{希 13:3}但此处他所谓\Quote{卑微的人}不仅指那些谦卑的人，也指那些地位低下、容易被人轻视的人。\par

\Quote{不要自以为聪明。}这是说，不要以为你们可以只靠自己。因为圣经又在另一处说：\BibleQuote{祸哉，那些在自己眼中看为有智慧，在自己看为精明的人。}\BibleRef{依 5:22}借此，他再次暗暗地除去鲁莽，减少自负与膨胀。因为没有什么比认为自己可以只靠自己，更能使人得意扬扬，并觉得自己与众不同。因此，天主也使我们彼此需要；尽管你聪明，你仍需要别人；但如果你认为你不需要他，你就是最愚笨、最软弱的人。因为这样的人会摒弃一切援助，无论陷入什么错误，既无法得到纠正，也无法获得宽恕，并且会因自己的鲁莽而触怒天主，陷入许多错误。事实上，的确常常如此：一个有智慧的人未必能看清需要什么，而一个不太精明的人却可能碰巧想到合适的方法。这发生在梅瑟和他岳父之间，发生在撒乌耳和他的仆人之间，也发生在依撒格和黎贝加之间。那么，不要以为需要别人就是降低了你自己，因为这反而更提升你，使你更坚强、更光彩、更稳妥。\par

17节。\BibleQuote{不要以恶报恶。}\BibleRef{罗 12:17}\par

因为如果你责备一个图谋反对你的人，为什么你要使自己也受同样的指控呢？如果他做了错事，你怎么能不回避效仿他呢？请注意，他在这里不作区分，而是为众人定下同一律法。因为他不是说\Quote{不要以恶报恶}只针对信徒，而是针对\Quote{所有的人}，无论他是外邦人，是被玷污的人，或别的什么人。\Quote{在众人面前做诚实的事。}\par

18节。\BibleQuote{若是可能，总要尽力与众人和睦。}\BibleRef{罗 12:18}\par

这就是：\BibleQuote{你们的光当在人前照耀}\BibleRef{玛 5:16}，不是要我们为虚荣而活，而是不要给那些有心找碴的人留下把柄。因此，他在另一处说：\BibleQuote{不可给人绊脚石，无论是犹太人、外邦人，还是天主的教会。}\BibleRef{格前 10:32}随后他又恰当地限制了这句话的意思，说：\Quote{若是可能}。因为有的时候是不可能的，例如，当我们必须为宗教辩护，或为受冤屈的人争辩时。为什么别人不能普遍做到这一点时我们感到惊讶呢？即使在夫妻之间，他也破例了：\BibleQuote{倘若不信的一方要离去，就由他离去吧。}\BibleRef{格前 7:15}他的意思大致如下：你尽好自己的本分，不给任何人以争战或争斗的借口，无论是犹太人还是外邦人。但是，如果你看见宗教的事业在任何地方受到损害，不要贵和谐而轻真理，而要勇敢地坚持，甚至至死不渝。即使在那时，内心不可争战，情绪不可敌对，只与事争战。因为这就是\Quote{总要尽力与众人和睦}的涵义。但如果对方不愿和睦，你不要让风暴充满你的灵魂，而要像我之前所说的，在心态上做朋友（\OriginalTerm{朋友}{\GreekText{φίλος}}；几份抄本作\OriginalTerm{哲学家}{\GreekText{φιλόσοφος}}），在任何情况下都不放弃真理。\par

19节。\BibleQuote{亲爱的诸位，不要自己报复，宁可给义怒留地步，因为经上记载：\InnerQuote{报复是我的，我必报复，上主说。}}\BibleRef{罗 12:19}\par

给什么义怒？给天主的义怒。既然受伤的人最渴望看到的是自己得享报复的快感，他便充分地给予了这一点：意思是，如果你不为自己报复，天主将成为你的伸冤者。那么，就将你的冤情交托给天主吧。因为这就是\Quote{给义怒留地步}的涵义。然后，为了给予进一步的安慰，他又引用了经文，就这样更彻底地赢得人心之后，他要求他们表现出更大的\OriginalTerm{基督徒的英勇}{\GreekText{φιλοσόφιάν}}，并说：\par

20至21节。\Quote{如果你的仇敌饿了，就给他吃；渴了，就给他喝；因为你这样做，就是把炭火堆在他头上。不要被恶所胜，反要以善胜恶。}\par

他说：我为什么要告诉你们必须保持和平？因为我甚至坚持要你们行善。因为他说道：\\BibleQuote{你给他吃，给他喝。}因为他所给的命令既非常困难又重大，他就继续说：\\BibleQuote{因为你这样做，就是把炭火堆在他的头上。}他说这话，既是为了用恐惧使那人谦卑，也是为了用赏报的望德使另一人更乐意。因为受冤屈的人，当他软弱时，与其说关心自己的利益，不如说关心报复那伤害他的人。因为没有比看见敌人受罚更甜蜜的了。他所渴望的，那他就先给他；等他把毒气释放了，然后他又给出更高层次的劝告，说道：\\BibleQuote{不要被恶所胜。}因为他知道，如果敌人是畜生，他也不会在被喂养之后继续为敌。即使受伤害的人灵魂再小，在喂养他、给他喝的时候，他自己也不会再渴望惩罚他。因此，出于对行为结果的信心，他不仅威胁，甚至大大强调报复。因为他没有说\\Quote{你要报复}，而是说\\BibleQuote{你要把炭火堆在他的头上。}然后他进一步宣布他为胜利者，说道：\\BibleQuote{不要被恶所胜，反要以善胜恶。}他给出一种温和的暗示，即人不可怀着那种意图去做，因为仍然怀恨在心就是\\Quote{被恶所胜。}但他没有立刻说，因为他还觉得不合适。但当他使那人卸下愤怒后，他就接着说：\\BibleQuote{以善胜恶。}因为这将是胜利。因为斗士那时才是征服者，不是当他俯身挨打，而是当他退开，让对手把力气浪费在空中。这样，他不仅自己不被击中，还会耗尽对方全部力气。这在面对侮辱时也是如此。因为当你以侮辱还侮辱时，你遭受了更糟的——不是被人征服（如某抄本写作\\OriginalTerm{被征服}{\GreekText{νικηθεὶς}}，Sav. 写作\\OriginalTerm{被激动}{\GreekText{κινηθεὶς}}），而是被愤怒这奴性的情欲所征服，这远更可耻。但如果你沉默，你就会得胜，不战而树立凯旋碑，会有千万人给你加冕，并谴责虚假的诽谤。因为那回答的人，好像被刺痛了而还口。而那被刺痛的人，让人有理由怀疑他自知有愧于所说的事。但如果你一笑置之，你的笑就消除了对你的指控。如果你想对所说的有明确的证据，去问问敌人自己：他何时最烦恼？是你发怒回骂他的时候？还是你在他辱骂时嘲笑他的时候？他会告诉你后者更甚。因为他也不因未被回骂而那么高兴，而是因他的辱骂未能对你产生影响而烦恼。你难道没有见过发怒的人吗？他们如何不顾自己的伤口，猛烈冲撞，比野猪更甚地寻求伤害邻人，并且只专注于这一点，比防备自己受害更热衷于此事？因此，当你剥夺了他最渴望的东西，你就夺走了他的一切，因为你这样轻视他，显明他易被轻看，更像小孩而非成人；你确实获得了智者的声誉，而把他赋予有害野兽的特性。这也应是我们被击打时所做的；当我们想要还手时，应当避免再次击打。但是，你想给他致命一击吗？\\BibleQuote{把另一边面颊也转给他}\\BibleRef{玛 5:39}，你就要用无数伤口击打他。因为那些为你鼓掌、惊叹你的人，比那些用石头砸他的人更让他难受；在他们面前，他的良心会定他的罪，并给他最严厉的惩罚，于是他就会带着困惑的表情离去，仿佛受到了最严厉的对待。如果你追求的是众人的评价，你也会得到更大的份。一般来说，我们对受苦的人有一种同情；但当我们看到他们不仅不还手（几个抄本写作\\OriginalTerm{抵抗者}{\GreekText{ἀντιπίπτοντας}}），反而顺服时，我们不仅怜悯他们，甚至钦佩他们。\par

于是，我在此找到哀叹的理由：我们这些本可拥有现世之物——若我们如理当听从基督的律法——也能获得来世之物的人，竟因不留意所告知我们的事，反而沉溺于对此的无端哲学思辨，而两者俱失。因为祂为我们的益处，就这一切事赐给我们律法，并指示我们何事使我们得美名，何事令我们蒙羞。若这会使祂的门徒显得可笑，祂就不会命令此事。但正因这使人成为最卓越的人——即，被恶言相向时不说恶言；被恶待时不行恶——祂才命令此事。既然如此，那么被恶言相向时反说善言，凌辱我们时反加赞美，谋害我们时反行善事，更会使人卓越。因此祂才赐下这些律法。因为祂关怀自己的门徒，深知何事使人渺小或伟大。既然祂既关怀又深知，你为何与祂争论，并愿另走他路？因为以恶制胜是魔鬼的律法之一。因此，在为其举办的奥林匹克竞技中，所有竞争者都以这种方式获胜。但在基督的赛跑中，关于奖赏的规则并非如此；相反，规则是使挨打者受冠冕，而非打人者。因为祂的赛跑之特点，就是所有规定都相反；如此，不仅在胜利本身，更在胜利的方式上，奇迹更为伟大。如今，当那些在另一边是胜利标志的事，在这一边却显示为导致失败时，这便是天主的大能，这便是属天的赛程，这便是天使的剧场。我知道你们现在已完全被温暖，变得如蜡般柔软，但一旦离开这里，你们将全部吐尽。这就是我悲哀的原因：我们所说的事，并未在行动中显明，尽管我们本应因此获益最大。因为若我们让人看到我们的温和，我们将在任何人面前不可战胜；无论尊卑，无人能伤害我们。因为若有人辱骂你，他并未伤害你分毫，而是严重伤害了自己。再者，若他冤屈你，伤害将归于那行冤屈的人。你难道从未注意到，即使在法庭上，受冤屈的人受尊重，站立发言完全自由；而行冤屈的人则因羞愧和恐惧而低头？我何必再谈恶言\EditorialNote{Sav. conj. 和五个抄本，原文 \OriginalTerm{恶言}{\GreekText{κακηγορίαν}}}和冤屈？即使他磨刀向你，将右手染于你的性命\EditorialNote{原文 \OriginalTerm{喉咙}{\GreekText{εἰς} \GreekText{τὸν} \GreekText{λαιμόν}}，见第505页}，他未曾伤害你，却屠宰了自己。那第一个被兄弟之手如此杀死的人将为我作证。因为他去了一个无风浪的港湾，赢得了不朽的荣耀；而另一人却过了一种比死亡更糟的生活，呻吟、颤抖，身体上背负着所行之事的控诉。因此，我们不要追求后者，而要追求前者。因为受恶待的人并不拥有恒久居住的邪恶，因为他不是其制造者；而是他从别人那里接受它，并以忍耐将其化为善。但行恶的人，心中怀着邪恶的泉源。若瑟不是在狱中，而设计害他的娼妓却在华美宏大的房屋里？你愿意曾是哪一个？且让我不要先听将来的报应，而只审视已发生之事本身。因为如此，你必会认为若瑟的监狱远胜那娼妓所在的房屋。因为若你看到两者的灵魂，你会发觉若瑟的灵魂充满宽宏和勇气，而埃及妇人的灵魂则狭窄、羞耻、沮丧、混乱且极度消沉。然而她似乎得胜了；但这并非真正的胜利者。既知这一切，让我们使自己适合承受恶事，甚至为的是免于承受恶事，并得以获得来世的福分。愿天主恩赐我们众人皆能达到，藉着恩宠和对人的爱，等等。\par

\SourceRecord{Translated by J. Walker, J. Sheppard and H. Browne, and revised by George B. Stevens.；Nicene and Post-Nicene Fathers, First Series；Vol. 11.；Edited by Philip Schaff.；Buffalo, NY: Christian Literature Publishing Co.,；1889.；<http://www.newadvent.org/fathers/210222.htm>.}

% structure: p0001 source=h1 target=section reason=page-title

\section{罗马书释义 第二十三讲}

\BibleRef{罗 13:1}\par

\BibleQuote{人人都当服从在上的权柄。}\par

他在其他书信中也十分重视这一主题，将臣民置于统治者的权下，犹如家仆服从主人一般。他这样做，乃是为表明基督订立祂的诫命不是为了颠覆政体，而是为了使其更有序，并教导人不要挑起不必要且无益的纷争。因为我们为真理所受的迫害已足够，无需再增添多余而无用的试探。且看，他谈论这一主题的时机多么恰当。当他要求人怀有那伟大的英雄气概，使人既能与友相处亦能与敌周旋，使其既能造福顺境者又能服务逆境及匮乏之人，实际上对所有人都有利，并树立了堪当天使的言行，消除怒气，抑制冒失，在各方面使人心平稳之后，他便引入关于这些事情的劝勉。因为若连伤害我们的人我们尚且应以相反之举回报，那么服从施恩于我们的人更是我们的本分。但他将这点放在劝勉的末尾，此前并未涉及我所说的那些推理，而仅以负有义务的方式命令人这样做。为表明这些规条适用于所有人，甚至司铎、隐修士，不仅限于世俗者，他开头的计划便如是说：\Quote{人人都当服从在上的权柄}，即使你是宗徒、传福音者、先知或任何身份，因为这种服从并不颠覆信仰。他不单说\Quote{服从}，而是说\Quote{服从在上的}。这一法令对于我们首要的要求，以及适合信徒的推理，便是这一切都是天主的安排。\par

\Quote{因为没有什么权柄不是出于天主的。}他说：\Quote{没有权柄不是出于天主的。}你们怎么说？也许有人说：难道每一个统治者都是天主选立的吗？他答道：我并非此意。我现在谈论的不是个别统治者，而是权柄本身。因为应当有统治者，有人统治有人被治理，万事不可混乱如一，民众如波浪般左摇右摆——我所说，这是天主智慧的工程。因此他不说：\Quote{没有什么统治者不是出于天主的，}而是谈论权柄本身，说：\Quote{没有权柄不是出于天主的。存在的权柄都是天主规定的。}因此，当某位智者说：\Quote{是上主使男人与女人相配} \BibleRef{箴 19:14}，他的意思是天主设立了婚姻，并非指每一次男女的结合都是天主的作为。因为我们看到许多人因邪恶而结合，即使是依婚姻之律，这种结合我们不应归于天主。就如祂自己所说：\Quote{那起初造人的，是造男造女，并且说：\InnerQuote{为此，人要离开父母，与妻子连合，二人成为一体。}} \BibleRef{玛 19:4}那位智者想要说明的正是这一点。因为平等的尊荣常导致争斗，所以祂设立了多种统治与服从的形式：例如男人与女人、儿子与父亲、老年人年轻人、奴隶与自由人、统治者与被统治者、老师与门徒。在人类中如此，你们何必惊奇呢？因为甚至在身体里祂也这样做了。在这里，祂也没有使所有肢体具有同等的尊荣，而是造了一个较小一个较大，有些肢体令其统治，有些令其服从。在无理性的受造物中，也能看到同样的原则，如蜜蜂、鹤、野牛群中。甚至海洋本身也不缺乏这种美好的隶属关系；因为在那里，许多鱼群都是由一条鱼统率，如军队般被带领，离家远行。因为无政府，无论在哪里，都是邪恶与混乱的根源。在说明了权柄的来源之后，他接着说：\Quote{所以，谁反对权柄，就是反对天主的制度。}请看，他将话题引向了何处，如何令人畏惧，如何表明这是人应尽的义务。为防止信徒们说：你将我们这些将要享受天国的人置于统治者的权下，使我们显得卑贱可鄙，他便指出，这样做不是使人服从统治者，而是再次使人服从天主。因为那服从权柄的人，便是顺从天主。然而他并未明说——例如，听从权柄的人就是听从未天主——而是用相反的情况来威慑他们，并以更精确的方式说：那不听从的，是与制定这些法律的天主抗争。他总是不遗余力地表明，我们服从他们不是出于恩惠，而是出于义务。因为这样更能使不信的统治者皈依信仰，使信徒服从。因为当时常有传言（\OriginalTerm{戴尔都良}{Tertullian} \WorkTitle{护教学} 1, 31, 32）诽谤宗徒，称他们是叛乱和革命的策划者，其言行皆旨在颠覆既定的制度。当你表明我们共同的导师将这点吩咐给祂的所有人时，你便能立刻堵住那些诽谤我们为革命者之口，并能极其大胆地为真理的教义辩护。他说道：因此，不要为这种服从感到羞耻。因为天主制定了这条法律，并要严厉惩罚那些藐视它的人。你们若不服从，祂将施予极重的惩罚，而非普通之罚；没有任何事能使你们免除，你们既会遭受人最严厉的报复，也无人为你们辩护，且更会激怒天主。这一切他都在以下言语中暗示了：\par

\BibleQuote{抗拒的必自取刑罚。} 然后为显明恐惧之后的好处，他又用理由来说服他们：\par

第3节。\BibleQuote{因为长官对行善的，不是可怕的；对行恶的才是可怕的。} \BibleRef{罗 13:3}\par

因为当他给予了深重的打击、使他们跌倒之后，他又如同一位明智的医生，使用温和的治疗，敷上舒慰的药物，安慰他们说：为什么害怕？为什么战栗？难道他惩罚行善的人吗？或者他对度美德生活的人是可怕的吗？为此他又继续说道：\BibleQuote{你愿意不害怕掌权的吗？你行善，就可从他得称赞。} 你看他如何藉着显示甚至连从宝座上他也称赞他们，而使他与掌权者\OriginalTerm{成为朋友}{\GreekText{ᾥ} \GreekText{κείωσεν}}。你看他如何使忿怒变得无意义。\par

第4节。\BibleQuote{因为他是天主的仆役，是为你行善的。} \BibleRef{罗 13:4}\par

他说，他不仅不恐吓你，反而称赞你；不仅不阻碍你，反而协助你。那么，既然你得到他的称赞和帮助，为何不顺服他呢？因为他也以其他方式使美德更容易：惩罚恶人，益惠并尊敬善人，并与天主的旨意合作。因此他甚至给他\BibleQuote{仆役}之名。试想：我劝你谨守明智，而他通过法律说同样的话。我劝你不要贪婪掠夺，而他审理此类案件，所以他是我们的同工和助手，且被天主为此目的而派遣。因此有理由尊敬他：既因他被天主派遣，又因他为如此目的。\BibleQuote{但你若行恶，就该害怕。} 可见令人害怕的不是掌权者，而是我们自己的邪恶。\par

\BibleQuote{因为他不是无缘无故佩剑的。} 你看他如何为他装备了武器，并像士兵一样安置他守卫，以恐吓犯罪的人。\BibleQuote{他是天主的仆役，为执行忿怒，惩治行恶的人。} 现在，为免你再听到惩罚、报复和剑就退缩，他说这不过是执行天主的法律。因为即使他自己不知道，但天主\OriginalTerm{如此规定}{\GreekText{οὕτως} \GreekText{ἐτύπωσεν}}。那么，若他在惩罚或尊荣时都是仆役，为美德伸冤，驱除邪恶，按天主的旨意行事，你为什么对他吹毛求疵呢？他成就了如此多的善事，也为你的善事铺平了道路。因为有许多人起初是因畏惧天主而实践美德的。有一些较为迟钝的人，未来的事对他们的影响不如现今的事。那么，那藉着恐惧和赏报使大多数人的灵魂预先转向，使之更适合听道的言语，此人理当被称为\BibleQuote{天主的仆役}。\par

第5节。\BibleQuote{因此，你们必须服从，不但是因为忿怒，也是因为良心。} \BibleRef{罗 13:5}\par

\BibleQuote{不但是因为忿怒}是什么意思？意思是：不仅因你不顺服而抗拒天主，不仅你为自己招致来自天主和掌权者的大祸，还因为他在你最重要的事上是恩人，因为他为你带来和平与公民制度的福祉。因为通过这些权威，国家有无数的祝福；如果你废除它们，一切都会毁灭，城市、乡村、私人和公共建筑，以及其他一切都会坍塌，整个世界将颠倒，强者吞噬弱者。因此，即使某些忿怒不跟随人的不顺服，甚至基于此，你也应当顺服，以免显得对恩人缺乏良心和感情。\par

第6节。\BibleQuote{因此，你们也该纳税，因为他们是天主的仆役，专为这事劳碌。} \BibleRef{罗 12:6}\par

他不逐一列举统治者对国家的好处，如秩序与和平、军队和公共事务等，而是用一个例子就表明了全部。他说，你们自己付薪水给他们，就是证明你们受益于他。请注意圣保禄的智慧和判断。因为那看似繁重讨厌的税收制度，他竟转变成他们照顾人的证明。他说，我们为何向国王纳税？难道不是因为他供应我们吗？但除非我们首先知道从这种管理中获益，我们本不会付税。自古以来，众人同意统治者应由我们供养，因为他们忽略自己的事务，管理公共事务，并为此付出全部闲暇，从而我们的财产也得以保全。在说了外在的利益之后，他又回到先前的论证（因为这样更能吸引信徒归向他），再次表明这是天主的命令，并以此作为劝告的最终依据，说：\BibleQuote{他们是天主的仆役。} 然后为显示他们的辛劳和艰苦的生活，他接着说：\par

\BibleQuote{专为这事劳碌。}\par

因为这就是他们的生活，他们的事务，好使你们享受和平。因此，在另一封书信中，他不仅命令要顺服，还要\BibleQuote{为他们祈祷}。且为表明那里的利益也是共同的，他加上：\BibleQuote{使我们安居乐业，事事虔敬。} \BibleRef{弟前 2:1} 因为他们对现世生活的安定贡献不小：守卫、击退敌人、制止城市中的叛乱、消除人们之间的分歧。不要向我提某人滥用职权的事，而要看制度本身中的秩序，你就会看到最初制定这律法者的伟大智慧。\par

第七、八节：\BibleQuote{因此，当将一切应得的归于众人；应课的，归于课；应税的，归于税；应敬畏的，归于敬畏；应敬重的，归于敬重。不可欠人任何事物（或\Quote{你们不可欠人任何事物}），但彼此相爱。}\par

他仍继续同一思路，吩咐他们不仅要付金钱，也要付出敬畏和敬重。那么，当他上面说：\BibleQuote{你们不愿害怕那权力吗？你们行善吧；} 而这里却说：\Quote{付出敬畏？} 他这样说是出于极度的尊敬，并非出于先前所指的那种来自恶劣良心的恐惧。他说的不是 \Quote{给}，而是 \Quote{归还}（或 \Quote{偿还}，\OriginalTerm{归还}{\GreekText{ἀπόδοτε}}），然后又加上 \Quote{应得的}。因为你这样做并非施恩，而是本分之事。你若不做，必因顽固受罚。不要以为你在长官面前起立或脱帽是贬低自己、有损你哲学的尊严。因为他在那时，当统治者还是外邦人时，就制定这些律法，如今他们成了信徒，更当如此行。但若你说自己受托付了更大的特权，须知这不是你的时候。因为你是旅客和寄居者。终有一天你要显出比众人更光明。现在你的\BibleQuote{生命与基督一同藏在天主内；当基督显现时，你们也要与他一同在光荣中显现} \BibleRef{哥 3:3}。所以，不要在这个充满意外的人生中追求改变；即使你不得不在长官面前带着敬畏，也不要认为这有损你的高贵出身。因为天主愿意这样：那由祂所设位的统治者，应拥有自己的权力。当那自觉无恶的人在审判官面前战栗时，行恶者岂不更加惊恐？而你反而因此更受尊重。因为贬低自己并非来自尊敬，而是来自不尊敬。而且，即使长官是不信者，他也会因你而更尊敬你，并因这缘故光荣你的主。\BibleQuote{除了彼此相爱外，你们不可再欠人什么。} 他再次求助众善之母、前述诸事的导师，也是各样德行的孕育者，并说这也是一种债务，却不像课税或关税那样，而是持续不断的。因为他希望它永不还清，或者更确切地说，他要它一直被偿还，却从不完全偿还，而是永远欠着。因为这债的特性就是人不断给予、不断欠债。在说了应如何爱之后，他也指出了爱的益处，说道：\par

\BibleQuote{因为爱别人的，就满全了法律。}\par

并且，我恳求你，连这事也不要视作恩惠；因为这也是一笔债。借着属灵的关系，你应爱你的弟兄。不仅为此，还因为\BibleQuote{我们彼此互为肢体}。如果爱离开我们，整个身体就会碎裂。因此，要爱你的弟兄。因为若从与他的友谊中你获得如此大的益处，以致满全了全部法律，你因受惠于他而欠他爱。\par

第九节：\BibleQuote{因为，不可奸淫，不可杀人，不可偷盗，不可妄证，以及其他任何诫命，都包含在这一句话里：爱你的近人如你自己。}\BibleRef{罗 13:9}\par

他并不是仅仅说它被满全，而是说\Quote{它被简要概括}，也就是说，诫命的全部工作被简洁地用几句话完成了。因为德行的开端和终结就是爱。它以此为本根，以此为根基，以此为顶峰。那么，若它既是开端又是满全，还有什么能与之相比呢？但他寻求的不仅是爱，而是炽热的爱。因为他不仅说\BibleQuote{爱你的近人}，而且说\BibleQuote{如同你自己}。因此基督也曾说，\BibleQuote{全部法律和先知都系于}爱。在划分两种爱时，请看祂是如何尊崇这爱的！因为说了第一条诫命\BibleQuote{你当爱主，你的天主}之后，祂又加上第二条；而且祂没有停住，又加上\BibleQuote{与此相似：你当爱你的近人如你自己}。对人的爱或温柔，有什么能与之相比呢？当我们与祂相距无限远时，祂将我们对祂的爱与对人的爱相比，说\BibleQuote{与此相似}。因此，为了使二者的尺度近乎相同，对一条祂说\BibleQuote{全心、全灵}，而对近人的爱祂说\BibleQuote{如同你自己}。但保禄说，若这爱不存在，即使另一种爱对我们也无大益。正如我们喜爱某人时，会说：你若爱他，就是爱我；同样，祂为表明这一点，也说\BibleQuote{与此相似}；并对伯多禄说\BibleQuote{你若爱我，就牧放我的羊群} \BibleRef{若 21:16}。\par

第十节：\BibleQuote{爱不做任何损害近人的事，所以爱是法律的满全。}\BibleRef{罗 13:10}\par

请留意它如何兼备两种美德：远离恶事（因为他说，\Quote{不作恶事}），以及行善。他说，\Quote{因为爱德就是法律的满全（或法律的填满）}；它不只以简洁的形式教导我们伦理责任，更使这些责任的实行变得容易。因为祂所关心的不仅是让我们认识有益之事（这是法律所关心的），而且也是帮助我们实行这些事，给我们带来极大的助益；不是完成部分诫命，而是在我们身上达成全部德行的总和。让我们彼此相爱，因为这样我们也必爱那爱我们的天主。因为在人际关系中，你若爱某人所爱的人，那爱他的人会因此嫉妒。但在此，祂却认为你配分享祂的爱，而当你不分享时，祂就恨你。因为人的爱充满嫉妒和猜忌；但天主的爱毫无私情，所以祂也寻求人来分享祂的爱。因为祂说：你与我一同爱，然后我自己也要更加爱你。你看这热切爱者的话语！你若爱我之所爱，那么我就必看自己为你所深爱。因为祂切望我们的得救，这从古时就已经显明。请听当祂造人时所说的话：\BibleQuote{让我们照我们的肖像造人}（\BibleRef{创 1:26}）；又说：\BibleQuote{让我们给他造一个相称的助手。人单独不好。}（\BibleRef{创 1:26}）当他违命时，祂责备他，且留意祂多么温和；祂没有说：可怜的东西！你这个可怜虫！受了如此大的恩惠，你竟还是相信了魔鬼？离开你的恩主，去跟随恶神？祂说了什么？\BibleQuote{谁告诉你，你是赤身露体的？除非你吃了那棵树上的果子，就是我命令你不可吃的那树？}（\BibleRef{创 3:11}）好像一个父亲对那被命令不能碰刀、却违命受伤的孩子说：\Quote{你怎么受伤了？你之所以如此，是因为不听我的话。}你看，这些话更像是一位朋友而非主人，是一位被轻视却仍未离弃的朋友的话。让我们效法祂，当我们责备时，也要保持这种温和。因为祂也用同样的温柔责备那女人。或者更确切地说，祂所说的与其说是责备，不如说是劝诫和矫正，并预防未来。因此祂对蛇什么也没说。因为它是祸害的策划者，无法将罪责推给别人，所以祂严厉惩罚了它；即使如此，祂也没有就此止步，还使大地也分担咒诅。但若祂将他们逐出乐园，判他们劳作，即使为此我们也该最崇拜和敬畏祂。因为放纵导致懈怠，祂借着痛苦构筑抵御懈怠的堡垒，从而限制享乐，好使我们回心转意爱祂。该隐的情形又如何？他岂不也受到同样的温和对待吗？因为虽然被他侮辱，祂没有以侮辱还报，反而恳求（或安慰）他，说：\BibleQuote{你为什么垂头丧气？}（\BibleRef{创 4:6}）然而他所做的事毫无借口。这一点幼弟也显示了。但即便如此，祂也没有责备他；祂说了什么？\BibleQuote{你犯了罪，要安静；}\BibleQuote{不要再这样做。}\BibleQuote{他的归属必归于你，你要管辖他}（\BibleRef{创 4:7}），意思是指他的弟弟。\Quote{因为你怕，}祂的意思是，\Quote{恐怕因为这个祭献，我要剥夺你长子的特权，你放心，我把他完全交在你手中。只要你变得更好，爱那没有亏负你的人；因为我关心你们二人。最使我高兴的是，你们二人彼此不纷争。}如同一位慈爱的母亲，天主做一切事、安排一切事，都是为了阻止一个人与另一个人分离；但为让你通过例子更清楚明白我的意思，请你回想，恳请你回想，黎贝加在苦恼中，跑来跑去，当时长子与幼子为敌。因为她若爱雅各伯，她也不恨厄撒乌。因此她说：\BibleQuote{免得我一天内失去你们两个儿子。}（\BibleRef{创 27:45}）因此天主在那时也说：\BibleQuote{你犯了罪，要平安：他的归属必归于你}（\BibleRef{创 4:7}），这样预先制止了凶杀，旨在使他们二人和睦。但当他杀害了弟弟后，祂甚至那时也没有停止对他的关怀，再次用温和的话语回答那杀弟者，说：\BibleQuote{你的弟弟亚伯尔在哪里？}这样，现在他若愿意，就可以完全认罪。但他却为自己的前恶辩护，带着更大更可耻的无耻。但即使那时，天主也没有离开他，再次说出一个受伤和被轻视的爱者的话语，说：\BibleQuote{你弟弟的血由地上向我喊冤。}（\BibleRef{创 4:10}）祂又与凶手一起责备大地，将怒气转向它，说：\BibleQuote{地开了口由你手中接受你弟弟的血，你必受地咒骂}（同上），且像那些哀悼的人（\OriginalTerm{哀悼者}{\GreekText{ἀνακαλοὕντας}}）所做的那样，正如达味在撒乌耳倒下时所行的。他对着迎接他死去的群山致辞，说：\BibleQuote{基耳波亚山，愿雨露不降于你，因为在那里勇士的盾牌被丢弃了。}（\BibleRef{撒下 1:21}）同样，天主也如同唱一首孤独的挽歌（\OriginalTerm{挽歌}{\GreekText{μονῳδίαν}}），说：\BibleQuote{你弟弟的血由地上向我喊冤；从此你从地上受咒骂，地开了口从你手中接受你弟弟的血。}祂说这话，是为降卑他炽热的情绪，并说服他现在至少去爱那已离世的人。你熄灭了祂的生命，祂会说：为何你还不熄灭仇恨？但祂做了什么呢？祂爱这一个也爱那一个，因为祂创造了他们两个。那么怎样呢？祂会让凶手不受惩罚吗？不，那样他会变得更坏。那么祂会惩罚他吗？不，祂比父亲更温柔。请看她如何一面惩罚，一面也在惩罚中显出祂的爱。或者更确切地说，祂不算惩罚，只是纠正。因为祂没有杀死他，只用颤抖束缚他，使他能摆脱罪行，这样至少他能恢复对那人的天然温情，并最终与那已离世的人讲和；因为祂不愿他带着对已故之人的敌意转到另一个世界。这就是相爱之人的方式：当他们行善却得不到爱的回报时，他们会变得激烈和威胁，并非出于本意，而是被爱驱使这样做；这样至少能赢得那些蔑视他们的人。但这种爱是不得已的，然而即便如此，因着爱的激烈，它仍能安慰他们。所以惩罚本身就来自爱，因为若非因为被恨而感到痛苦，他们也不会选择惩罚。现在，注意保禄对格林多人所说的话。因为\BibleQuote{谁使我快乐呢？岂不是那因我而悲伤的人吗？}（\BibleRef{格后 2:2}）所以当他走到惩罚的极致时，他反而显出他的爱。同样，那埃及妇人因着强烈的爱，也强烈地惩罚了若瑟：她确实是为邪恶，因为那爱是不洁的；但天主是为善，因为那爱与爱者相称。这就是为什么祂不拒绝屈就于粗俗的词语，说出人间情欲的名称，并称自己为嫉妒的。\BibleQuote{因为我是嫉妒的天主}（\BibleRef{出 20:5}），祂说，好使你认识这爱的强烈。让我们如祂所愿爱祂：因为祂非常看重这个。若我们转离，祂不断邀请我们；若我们不肯悔改，祂出于爱而惩戒我们，并非出于惩罚我们的愿望。请看祂在厄则克耳中对那被爱却轻视祂的城说了什么。\BibleQuote{我要使你的情敌来攻击你，把你交在他们手中；他们要石击你，杀死你，这样我的嫉妒就离开你，我就安静，不再烦恼。}（\BibleRef{则 16:37}）一个热烈的爱人，当他被所爱者轻视、又再次热烈爱她时，还能说出比这更多的话吗？因为天主做一切事，为使我们爱祂，为此祂连自己的圣子也没有保留。但我们却顽固不化，野蛮无情。然而，让我们终于变得温柔，按当爱的方式爱天主，好使我们愉快地享有德行。因为若有人有爱妻，就觉察不到每天到来的烦恼，那怀着这神圣纯洁之爱而爱的人，请想他将享受何等大的喜乐！因为这就是，确实这就是天国；这是善事的享用、快乐、喜乐和幸福。或者，无论我说多少，我仍无法向你们充分描绘它，唯有亲身体验才能认识这美善。所以先知也说：\BibleQuote{你要因上主而喜乐}（\BibleRef{咏 36:4}），\BibleQuote{请你们品尝，并请观看，上主是何等的甘饴}（\BibleRef{咏 33:8}）。让我们受此劝服，沉浸于祂的爱中。因为这样，我们即使在今生也将看见祂的国度，并度天使般的生活；当我们在地上居住时，也将如天上居民一般幸福；在我们离世后，将成为在基督审判座前最光辉的受造物，并享受那不可言喻的荣耀。愿我们众人因我们的主耶稣基督的恩宠和慈爱，都能到达那里。因为光荣归于祂，直到永远。阿们。\par

\SourceRecord{Translated by J. Walker, J. Sheppard and H. Browne, and revised by George B. Stevens.；Nicene and Post-Nicene Fathers, First Series；Vol. 11.；Edited by Philip Schaff.；Buffalo, NY: Christian Literature Publishing Co.,；1889.；<http://www.newadvent.org/fathers/210223.htm>.}

% structure: p0001 source=h1 target=section reason=page-title

\section{\WorkTitle{论《罗马书》第二十四篇讲道}}

\BibleRef{罗 13:11}\par

\BibleQuote{而且，你们晓得这时期，现在正是该从睡梦中醒来的时刻。}\par

既然他已给了他们合宜的命令，他再次因临近的催促，推动他们行善。因为他意指审判的时刻就在门口。他同样写信给格林多人说：\BibleQuote{剩下的时间不多了。} \BibleRef{格前 7:29} 又写给希伯来人：\BibleQuote{因为还有一点点时候，那要来的就要来，决不会迟延。} \BibleRef{希 10:37} 但在那些情形下，他说这些话是为了鼓励受苦的人，安慰他们接连不断试探的辛劳；而在我们面前的这段经文，他是为了唤醒沉睡的人。这语词对我们都很有用。他说的是什么？\BibleQuote{现在正是该从睡梦中醒来的时刻？} 意思是：复活近了，可畏的审判近了，那烧如熔炉的日子近了。因此，我们必须摆脱懈怠，\BibleQuote{因为现在我们的救恩比我们初信时更近了。} 你看他如何把复活放在他们身旁。因为随着时间推移，他意指，我们今生季节正在消逝，来生季节愈发临近。如果你已预备好，并完成了他所命令的一切，那日子就是你的\OriginalTerm{救恩对你}{\GreekText{σωτηρία} \GreekText{σοι}}；如果不是，则不然。但此刻，他不是用可怕的理由劝诫他们，而是用慈爱的理由，如此也使他们摆脱对现世之物的依恋。既然在最初努力时，他们很可能最为热切，因为那时欲望正盛，但随着时间推移，他们的全部热忱可能消减至无；他说他们反而应当反其道而行之，不要随时间变得松弛，而应更加充满活力。因为君王越近，他们就越应做好准备；奖赏越近，他们越应警醒奋战，因为连赛跑者也如此：当接近终点、即将接受奖品时，他们会更加振奋。这就是为什么他说：\BibleQuote{现在我们的救恩比我们初信时更近了。}\par

\BibleRef{罗 13:12} \BibleQuote{夜已深了，白日已近。}\par

如果这事即将结束，而后者正临近，让我们从此做属于后者的事，而非前者。这是今生之事中的做法。当我们看到黑夜催促着黎明，听到燕子啁啾，我们相互唤醒邻人，尽管还是夜晚。但当夜真正离去时，我们彼此催促，说：现在已是白日！于是我们都开始做白日的工作，穿衣、摒弃梦境、彻底驱散睡意，好让白日发现我们已准备好，而不必在阳光升起时才开始起身伸腰。那么，我们在那种情形下的做法，也让我们在此做吧。让我们抛开幻想，摆脱今生的梦境，卸下深沉的睡眠，披上美德作衣裳。因为为指出这一切，他说：\par

\BibleQuote{让我们抛弃黑暗的行为，穿上光明的武器。}\par

是的，因为白日呼唤我们整装待发，投入战斗。但听到阵列和武器，不要害怕。因为就可见的盔甲而言，穿上它是沉重而可憎的任务。但这里却是令人向往、值得祈求的。因为这是光明的武器！因此它们会使你比阳光更明亮，发出耀眼的光芒，并保你安全：因为它们是武器，且使你闪耀：因为它们是光明的武器！那么，你是否必须战斗？是的，必须战斗，但不必苦恼辛劳。因为这其实并非战争，而是庄严的舞蹈和节庆，这就是武器的性质，这就是统帅的能力。正如新郎欢欣地从洞房走出，如此，被这些武器保护的人也如此。他既是战士，又是新郎。但当他说\BibleQuote{白日已近}时，他不仅容许它临近，反而把它置于我们身边。因为他说：\par

\BibleQuote{让我们行事得体，\EditorialNote{英译\Quote{honestly}在此意为\Quote{得体地}}如同在白日。} 因为此时已是白日。大多数人劝诫时极为强调的，他也用来引导他们，即得体的意识。因为他们非常看重民众的赞誉。他并非说\Quote{你们行走}，而是说\Quote{我们行走}，使劝诫免于刺耳，责备显得温和。\par

\BibleQuote{不可在宴乐与醉酒中。} 并非禁止饮酒，而是禁止无度；并非禁止享用酒，而是禁止\OriginalTerm{饮酒过度}{\GreekText{μετά} \GreekText{παροινίας}}。正如接下来他以同样尺度陈述的：\par

\BibleQuote{不在好色与放荡中；} 因为他此处并非禁止两性交往，而是禁止行淫。\BibleQuote{不在争吵与嫉妒中。} 他要熄灭的是致命的私欲，即情欲和愤怒。因此，他不仅除去它们本身，甚至除去其根源。因为没有什么像醉酒和久坐酒席那样点燃情欲、激起愤怒。因此，在先说\BibleQuote{不在宴乐与醉酒中}之后，他接着道\BibleQuote{不在好色与放荡中，不在争吵与嫉妒中。} 即使在此他也没有停步，而是在脱下我们这些邪恶的衣裳之后，听他又如何装饰我们，他说：\par

\BibleRef{罗 13:14} \BibleQuote{但要穿上主耶稣基督。}\par

他不再谈工作，而是激励他们追求更伟大的事。因为他谈论恶习时，提到了恶行；但论及德行时，他不谈工作，而谈武器，以表明德行使拥有它的人完全安全、完全光明。甚至他并不止步于此，而是将言论引向更伟大、远更令人敬畏的事：他将主本身赐给我们作为衣服，将君王本身赐给我们：因为穿上祂的人，就拥有了一切德行。但他说\Quote{穿上}，是吩咐我们从四面八方被祂环绕。正如在另一处他说：\BibleQuote{但如果基督在你们内。}\BibleRef{罗 8:10}又说：\BibleQuote{使基督住在内里。}\BibleRef{弗 3:16}因为祂愿我们的灵魂成为祂的居所，而祂自己被铺在我们周围作为衣服，好使祂从内从外成为我们的一切。因为祂是我们的圆满；因为祂就是\BibleQuote{那充满万有的圆满}\BibleRef{弗 1:23}；又是道路、丈夫、新郎——因为\BibleQuote{我曾把你们许配给一个丈夫，如同贞洁的童女}\BibleRef{格后 11:2}；又是根、饮料、食物、生命——因为他说：\BibleQuote{我活着，但不是我，而是基督在我内活着}\BibleRef{迦 2:20}；又是宗徒、大司祭、教师、父亲、兄弟、同承继者、坟墓与十字架的分担者——因为经上说：\BibleQuote{我们与祂同葬}，和\BibleQuote{在祂死亡的样式上同植}\BibleRef{罗 6:4}；又是恳求者——\BibleQuote{我们代基督作大使}\BibleRef{格后 5:20}；又是\Quote{在父那里的护慰者}——因为经上说：\BibleQuote{祂也为我们转求}\BibleRef{罗 8:34}；又是房屋和居住者——因为祂说：\BibleQuote{那住在我内，我也住在他内}\BibleRef{若 15:5}；又是朋友，因为\BibleQuote{你们是我的朋友}\BibleRef{若 15:14}；又是基础和基石。而我们是祂的肢体、祂的产业、建筑物、枝条和同工。当祂使我们以各种方式依附并适合祂时，有什么祂不愿意成为我们的呢？这是极爱之人的标志。所以，要受劝诫，从睡眠中醒来，穿上祂，并在穿上之后，将你的肉身交给祂的缰绳。因为这就是他在说以下话时所暗示的：\par

\Quote{不要为肉身安排，去满足它的私欲。}因为他并没有禁止饮酒，而是禁止酗酒；没有禁止结婚，而是禁止行邪淫；同样，他也不禁止为肉身安排，但禁止这样安排乃是为了\Quote{满足它的私欲}，例如超越必需。因为他确实吩咐要为肉身安排，请听他向弟茂德说的话：\BibleQuote{你该略微用点酒，为了你的胃病，并因你常害病。}\BibleRef{弟前 5:23}所以这里他也主张照顾肉身，但为了健康，而非邪淫。因为当你点燃火焰、使炉火旺盛时，这就不是为它安排了。但为了让你更清楚地明白什么是\Quote{为肉身安排}以\Quote{满足它的私欲}，并避而远之，只需回想那些醉汉、贪食者、以衣着为傲者、柔弱者、生活安逸松懈者，你就会明白所指。因为他们所做的一切，不是为健康，而是为行邪淫并点燃情欲。但愿你，既已穿上基督，就剪除这一切，只求一事：使你的肉身健康。要这样为它安排，而不再多行，却将你所有的努力用于灵性之事的照料。因为那时你就能从这睡眠中醒来，不被这许多私欲所重压。因为现世生活是一场睡眠，其中的事物与梦境无异。如同睡者常常说出和看见不健康的事物，我们也是如此，甚至更糟。因为若有人在梦中做了或说了可耻之事，当他醒来，也就脱离了羞耻，且不受惩罚。但在现实情形中却不然，羞耻与惩罚都是不朽的。再者，那些在梦中致富的人，天亮时就被证明是徒然富有。但在现实中，甚至在天亮之前，他们常被定罪，而在他们离世归入另一生命之前，那些梦已飞散。\par

那么，让我们摆脱这邪恶的睡眠，因为如果白昼发现我们在睡觉，一种不死之死将接踵而至；而在那白昼之前，我们就会暴露于世上一切仇敌的攻击之下，无论是人还是魔鬼；如果他们存心来毁灭我们，没有人能阻止他们。因为如果有很多人在守望，危险就不会这么大；然而，也许只有一两个人点着一盏灯，仿佛在深夜中守望，而人们都在睡觉；因此现在我们需要很多的警醒，很多的戒备，以防止我们陷入最不可挽救的祸患。现在难道不是似乎已是大白天吗？我们难道不是以为所有的人都醒着并清醒吗？然而仍然（也许你会对我所说的话微笑，但我还是要说出来）我们大家都像沉睡在深夜中打鼾的人。而且，如果真能看见一个无形体之物，我本可以向你们展示大多数人是如何打鼾，魔鬼是如何破墙而入，在我们躺卧时屠杀我们，偷走里面的财物，无所畏惧地做一切，如同在深沉的黑暗中。或者更确切地说，即使无法用肉眼看见，就让我们用言语勾勒出来吧，并想想有多少人被邪恶的欲望压垮，有多少人被放荡的剧痛压制，熄灭了圣神的光。因此，他们以这件事代替那件事，听见一个事物代替另一个，对这里告诉他们的一切都不加理会。或者，如果我这样说错了，而你们是醒着的，那就告诉我今天这里做了什么，如果你们不是像在梦里听了这些。我确实知道有些人能告诉我（我并非指所有人）；但你这个被所说的话所涵盖的人，你徒劳地来到这里，告诉我今天哪位先知、哪位宗徒向我们讲了话？讲了些什么？而你却没有能力告诉我。因为你在这里说了很多，就像在梦里一样，没有听见现实。这话我也要对妇女们说，因为她们中间有很多人睡着了。但愿这就是睡觉！因为睡着的人不说好也不说坏。但像你这样醒着的人，却说出许多甚至自招祸害的话来，谈论自己的利益，计算债务账目，回忆某些无耻的交易，在自己灵魂里密植荆棘，不让种子有丝毫前进。但要振作起来，把这些荆棘连根拔起，抖掉醉酒：因为这是睡眠的原因。但我所说的醉酒，不仅指酒的醉，也指世俗思想的醉，以及与之相随的酒的醉。（参第443页）我给这个建议不仅是对富人，也是对穷人，主要是那些合伙举办社交聚会的人。因为这并不是真正的放纵或放松，而是惩罚和报复。因为放纵不在于说脏话，而在于庄重地谈话，在于满足，而不是要撑破。但如果你认为这是享乐，那就到晚上把享乐给我看看！你不能！至今我对它导致的祸害还只字未提，现在只对你们讲了如此迅速凋零的享乐。因为聚会一散，所有寻欢作乐的东西就都飞走了。但当我提到呕吐、头痛、无数的不适和灵魂的俘虏时，你对这一切还有什么可说？我们难道因为穷，就可以举止不当吗？我说这些并不是禁止你们聚在一起，或在公共餐桌吃晚饭，而是为了防止你们举止不当，并希望放纵真正成为放纵，而不是惩罚、报复、醉酒或狂欢。让外邦人（\OriginalTerm{外邦人}{\GreekText{Ἕλληνες}}）看到基督徒最懂得如何放纵，并且有秩序地放纵。因为经上说：\BibleQuote{应当战兢而喜乐于主。} \BibleRef{咏 2:11} 那么，人怎能喜乐呢？就是通过唱圣歌、祈祷、用圣咏来取代那些低俗的歌曲。这样基督也会在我们的筵席上，用祝福充满整个盛宴，当你祈祷、唱属灵的歌曲、邀请穷人分享摆在你面前的、在宴会上设立很多秩序和节制的时候。这样你就会使聚会成为教会，用颂赞万有之主来代替不合时宜的喊叫和欢呼。不要告诉我另一种习俗已经盛行，而要纠正这种错误。因为经上说：\BibleQuote{因为你们或吃或喝}，经上说，\BibleQuote{或无论做什么，一切都要为光荣天主而做。} \BibleRef{格前 10:31} 因为从那种宴席上你们生出邪恶的欲望和不洁，妻子在你们中受轻视，娼妓却受尊敬。由此而来的是家庭倾覆和无数祸患，一切都颠倒混乱，你们离开了纯净的泉源，奔向泥沟。因为娼妓的身体是泥，我不向别人查问，只问你自己，你这在泥中打滚的人，难道你不为自己感到羞耻吗？难道犯罪之后你不觉得自己不洁吗？因此我恳求你们逃避奸淫，以及它的母亲——醉酒。为什么在不可能收割的地方撒种，或者即使你收割了，那果实也给你带来极大的耻辱？因为即使孩子出生，也立即玷污了你，而它本身也因你而出生为私生子、卑贱者，遭受了不公。即使你给它留下很多钱，无论是娼妓之子还是婢女之子，在家中不名誉，在城中不名誉，在法庭上不名誉；你也会不名誉，无论是生前还是死后。因为即使你离世了，你不端行为的纪念仍存留。那么为什么给所有这些带来耻辱？为什么在土地刻意毁坏果实的地方撒种？那里有许多堕胎的努力？那里有出生前的谋杀？因为你甚至不让娼妓仅仅做个娼妓，还使她成为杀人的凶手。你看见醉酒如何导致淫乱，淫乱导致奸淫，奸淫导致谋杀；或者甚至比谋杀更糟的事。因为我无法给它起名，因为它不是把已出生的取下，而是阻止其出生。那么，你为什么滥用天主的恩赐，与祂的法律争斗，追求如同祝福的诅咒，使生育之室成为谋杀之室，并武装原本为生育而赐予的女人去杀人？为了通过讨人喜欢、成为情人们渴慕的对象来赚取更多钱财，她甚至不犹豫做这事，从而在你头上堆积一大堆火。因为即使这胆大妄为的行为是她的，但导致它的却是你。由此也产生了拜偶像的事，因为很多人为了变得可接受，设计咒语、奠酒、爱情药水以及无数其他计谋。然而，在如此大的不端之后，在屠杀之后，在拜偶像之后，这件事在许多人看来仍属无关紧要之事，而且甚至对许多有妻子的人来说也是如此。因此祸害的混杂（\OriginalTerm{混杂}{\GreekText{φορυτὸς}}）就更大。因为邪术不仅施用于被出卖的子宮，也施用于受伤害的妻子；有无数的阴谋、对魔鬼的祈求、通灵术、每日的战争、无休止的打斗和家中珍藏的嫉妒。因此保禄在说了\BibleQuote{不在荒宴和放荡}之后，又接着说\BibleQuote{不在争竞和嫉妒}，因为他知道由此产生的战争、家庭的倾覆、对合法子女的损害以及无数其他祸患。那么，为了逃避这一切，让我们穿上基督，并常常与祂同在。因为这就是穿上祂的意思：从不离开祂，使祂在我们的圣化和节制中，在我们内永远可见。所以我们论及朋友时说，某人穿上了（\OriginalTerm{穿上}{\GreekText{ἐνεδύσατο}}）另一个人，意指他们极大的爱和不断的相守。因为穿上某物的人，似乎就是他穿上的那个东西。那么，让基督在我们每一部分中被看见。如何被看见？如果你行祂的作为。祂做了什么？祂说：\BibleQuote{人子没有枕头的地方。} \BibleRef{路 9:58} 你也要以此为目标。祂需要食物，祂吃的是大麦饼。祂需要旅行，但到处没有马或驮畜，祂步行甚至到疲惫。祂需要睡眠，祂\BibleQuote{在船头的枕上睡着了}（\OriginalTerm{船尾（此处指船首）}{\GreekText{πρύμνῃ}}，此处\OriginalTerm{船首}{\GreekText{πρῶρα}}）\BibleRef{谷 4:38}。需要坐下吃饭时，祂吩咐他们躺在草地上。祂的衣服是便宜的；祂常常独自留下，没有随从跟在后面。祂在十字架上所行的，在侮辱中所行的，以及总而言之祂所做的一切，你要学到心里（\OriginalTerm{学到心里}{\GreekText{καταμαθών}}）并效仿。这样你就穿上了基督，如果你\BibleQuote{不使肉体筹备满足其私欲}。因为这并没有真正的快乐，因为这些私欲又生出更强烈的其他私欲，你永远得不到满足，只会使你受很大的折磨。因为如同一个持续口渴的人，即使身边有上万泉水，也得不到好处，因为他无法熄灭这失调，同样，持续活在私欲中的人也是如此。但如果你守住必需之物，你就绝不会陷入这种恐惧，这一切都会消失，无论是醉酒还是放荡。那么就只吃够解除饥饿的量，只穿够遮蔽身体的量，不要好奇地用衣服装饰你的肉体，免得毁坏它。因为你会使它更娇弱，用如此多的柔软使之无力，损害它的健康。为了使它成为灵魂的合适载体，使舵手安全地坐在舵上，士兵轻松地挥舞武器，你必须使所有部分合适地组合在一起。因为使我们不受伤害的不是拥有很多，而是需求很少。因为一个人即使没有受冤屈也害怕；而另一个人即使受冤屈，也比那些没有受冤屈的人处境更好，并且正是为此他精神更佳。那么，我们追求的目标，并不是如何防止别人恶待我们，而是即使他愿意这样做，他也无能为力。除了守住在必需之物上，不贪求更多之外，没有其他来源可以得着。因为这样我们就能在此享受，并藉着恩宠和对人的爱等，达到将来的美好事物。\par

\SourceRecord{Translated by J. Walker, J. Sheppard and H. Browne, and revised by George B. Stevens.；Nicene and Post-Nicene Fathers, First Series；Vol. 11.；Edited by Philip Schaff.；Buffalo, NY: Christian Literature Publishing Co.,；1889.；<http://www.newadvent.org/fathers/210224.htm>.}

% structure: p0001 source=h1 target=section reason=page-title

\section{\WorkTitle{论罗马书第二十五篇讲道}}

罗马书第十四章1,2节\par

\BibleQuote{信德软弱的人，你们要接纳，但不要为辩论所疑惑。有人信什么都能吃；但那软弱的，只吃蔬菜。}\par

我知道这对大多数人而言是个难题。因此我必须首先说明整段经文的主题，以及保禄写这段话所要纠正的是什么。那么他要纠正什么呢？有许多信主的犹太人，他们因良心而持守法律，在信主之后仍然遵守关于食物的规条，因为他们还没有勇气完全脱离法律的束缚。为避免因只不吃猪肉而被注意，他们便戒除一切肉类，只吃蔬菜，这样他们的行为更像是守斋而非遵守法律。另一些人则更为精进（\OriginalTerm{精进}{\GreekText{τελει}\'\GreekText{οτεροι}}），毫不遵守这些规条，他们成为那些遵守者苦恼和冒犯的来源，因为他们责备、指责那些人，使他们灰心丧气。因此，蒙福的保禄担心：若为了一点小事而想要正确行事，反而可能颠覆全局；若想要使他们对于吃什么漠不关心，反而可能使他们易于放弃信仰；若因急于一次纠正所有事，在适当时机到来之前，反而在关键问题上造成损害——由于不断的责备，使他们偏离了对基督的共识（\OriginalTerm{共识}{\'\GreekText{ομολογ}\'\GreekText{ιας} \GreekText{ε}\'\GreekText{ις}}），结果两方面都未能纠正。请看他是如何以智慧运用伟大的判断，并同时顾及双方的利益。他既不敢对责备者说\BibleQuote{你们做错了}，以免显得肯定对方的守则；也不敢说\BibleQuote{你们做对了}，以免使他们更加严厉地指责。而是使他的责备恰到好处。表面上他责备那较强的，但实际上他通过向这些人的讲话，将所有要说的话都倾注于那较弱的一方。因为最不容易引起反感的纠正方式，就是当一个人对另一个人说话时，却打击了不同的对象，这样不会让被责备者动怒，并悄悄地引入纠正的良药。请看他是如何以判断来行事，时机又是多么恰当。因为说了\BibleQuote{不要为肉体安排，去放纵私欲}之后，他就进入对这些问题的讨论，以免显得是在为那些责备者、也就是主张什么都可以吃的人辩护。因为较弱的一方总是需要更多的考虑。因此他直接打击那强的，立刻说道：\BibleQuote{信德软弱的人}。你看到他立刻给了他一击。因为称他为\Quote{软弱}（\OriginalTerm{软弱}{\'\GreekText{ασθενο}\'\GreekText{υντα}}），就指明他并非健康（\OriginalTerm{健康}{\'\GreekText{αρρωστον}}）。然后他接着说\BibleQuote{你们要接纳}，再次表明他需要许多关注。这是极度虚弱的标志。\BibleQuote{但不要为辩论所疑惑}。看，他又打了第三下。因为他表明他的错误性质如此严重，以至于那些没有同样犯错、却仍以爱心接纳他、并热切想要医治他的人，也都感到疑惑。你看他表面上是在与这些人谈话，却暗中不动声色地责备了别人。然后他将两者并列，一方得到赞扬，另一方受到指责。因为他接着说：\BibleQuote{有人信什么都能吃}，称赞他的信德。\BibleQuote{但那软弱的，只吃蔬菜}，再次贬低后者，指出他的软弱。既然那一击是致命的（\OriginalTerm{致命的}{\GreekText{καιρ}\'\GreekText{ιαν}}，夸张用法），他又用这些话安慰他：\par

第3节。\BibleQuote{吃的人不可轻视不吃的人。}\BibleRef{罗 14:3}\par

他没有说\Quote{放过他}，也没有说\Quote{不要责备他}，更没有说\Quote{不要纠正他}，而是说不要辱骂他，不要\Quote{轻视}他，以表明他们的行为是完全可笑的。但他以别的话来说这事。\BibleQuote{不吃的人不可判断吃的人。}因为那更精进的人看不起这些持守者，视他们为信德小、虚假医治、假冒和仍犹太化的人；同样，这些持守者判断那吃得人，视他们为违法者或贪吃者。很可能这些持守者中也有许多外邦人。因此他接着说：\BibleQuote{因为天主已经收纳了他。}但在另一情况下他没有这样说。而被轻视的似乎是吃的人，因为被认为贪吃；而被判断的则是不吃的人，因为被认为信德小。但他使他们的位置互换，以表明他不仅不该被轻视，反而能够轻视别人。他是说\Quote{我定他的罪吗}？绝不是。这就是为什么他接着说：\BibleQuote{因为天主已经收纳了他。}那么你为什么对他讲法律，如同对犯法者一样？\BibleQuote{因为天主已经收纳了他}，意思是天主已在他身上显示了不可言喻的恩宠，免除了他所有指控；然后他又转向那强的。\par

第4节。\BibleQuote{你是谁，竟判断别人的仆人？}\BibleRef{罗 14:4}\par

由此可见，他们不但轻看，而且判断了。\Quote{他或站住或跌倒，自有他的主人。}请看，这里又有一击。这愤慨似乎针对强健的人，而（保禄）攻击了他。当他说：\Quote{他必被扶持起来}，这显示他仍在摇摆，需要如此大的关注，以致请天主来作为医治者，因为他说：\Quote{天主能使他站住。}而我们对完全绝望的事才这样说。然后，为了不使他绝望，他既称他为仆人，说：\Quote{你是谁，竟判断别人的仆人？}这里他又暗暗地攻击他。因为我叫你们不要判断他，不是因为他行了当得起不被判断的事，而是因为他是别人的仆人，即不是你们的，而是天主的。接着，为再次安慰他，他没有说：\Quote{跌倒}，而是说什么？\Quote{站住或跌倒。}无论是后者还是前者，都是主人的事，因为若他跌倒，损失归于主，若他站住，财富也归于主。如果我们不注意保禄的目的（他不愿在他们时机未成熟时受责备），这对基督徒应有的彼此关怀是极不相称的。但（如我常说的）我们必须审视说话时的用心、所谈论的主题以及他说这话所要达到的目的。他以此并非微小的动机使他们心存敬畏，因为他的意思是：如果天主（祂承受损失）至今什么也不做，你怎能不过分苛刻、不合时宜地抓住他\OriginalTerm{扼喉}{\GreekText{ἄγκων}}并烦扰他呢？\par

第5节。\BibleQuote{有人看这日比那日强，有人看日日都一样。}\BibleRef{罗 14:5}\par

这里他似乎是温和地暗示了关于禁食的事。因为很可能有些禁食的人常判断那些不禁食的人，或者在那些遵守中，很可能有些人在固定日期禁食，在固定日期不禁食。因此他也说：\BibleQuote{各人在自己心里要心意坚定。}\BibleRef{罗 14:5}这样，他借说这事是无所谓的，解除了那些遵守者（从遵守而来的）恐惧，也消除了那些攻击他们的人的争辩之性，表明总是为这些事找麻烦并不是非常可欲的\OriginalTerm{非常可欲的}{\GreekText{περισπούδαστον}}（或紧迫的）事。然而，这不是非常可欲的，不是出于这事本身的性质，而是因为所选的时候，因为他们是在信仰上的初学。因为当他写信给\BibleBook{哥罗森}{书}时，他极其恳切地禁止这事，说：\BibleQuote{你们要谨慎，恐怕有人用他的哲学和虚空的妄言，照着人的传统，照着世俗的小学，不照着基督，把你们掳去。}\BibleRef{哥 2:8}又说：\BibleQuote{所以，不拘在饮食上，你们不可让人论断你们。}\BibleRef{哥 2:16}并且：\BibleQuote{不可让人夺去你们的奖赏。}\BibleRef{哥 2:18}而当写信给\BibleBook{迦拉达}{书}时，他以极大的精确要求他们在这事上有基督徒的精神和成全。但在这里他没有使用这种激烈，因为信德是最近在他们中栽种的。因此，我们不可将\Quote{各人在自己心里要心意坚定}应用于一切主题。因为当他谈论道理时，请听他说的：\BibleQuote{若有人给你们传讲什么福音，与你们所领受的不同，他就当被咒诅。}\BibleRef{迦 1:9}即使是\Quote{一位天使}也一样。又说：\BibleQuote{我只怕你们的心或偏于邪，失去那向基督所存纯一清洁的心，就像蛇用诡诈诱惑了夏娃一样。}\BibleRef{格后 11:3}而在写给\BibleBook{斐理伯}{书}时，他说：\BibleQuote{应当防备狗，防备作恶的，防备妄自行割的。}\BibleRef{斐 3:2}但对罗马人，因为还不是纠正这类事的适当时候，他说：\Quote{各人在自己心里要心意坚定。}\BibleRef{罗 14:5}因为他刚才是在谈论禁食。正是为了消除其他人的虚荣和解除这些人的恐惧，他说了下面的话：\par

第6节。\BibleQuote{守日的人是为主守的；不守日的人是为主不守的。吃的人是为主吃的，因为他感谢天主；不吃的人是为主不吃的，也感谢天主。}\BibleRef{罗 14:6}\par

他仍保持同一主题。他所说的意思是：这事不关乎基本要理。因为必要的是，如果这人和那人都是为天主而行，必要的是（这些字在三份手稿中重复），如果两者都归结于感恩。因为实际上这人和那人都感谢天主。如果他们都感谢天主，那么区别就不大了。但请你们注意，他在这里如何暗中打击了犹太化者。因为如果必要的是这\Quote{感谢}，那么很明显，吃的人是\Quote{感谢}的，而不\Quote{吃的人}则不是。因为当他还持守法律时，如何能呢？正如他在\BibleBook{迦拉达}{书}中所说：\BibleQuote{你们这要靠法律成义的，是与基督隔绝，从恩宠中坠落了。}\BibleRef{迦 5:4}所以在这里他只暗示，并未完全展开。因为当时还不是这样做的时候。但此刻他容忍了（参看337页）：然而，借着下面的话，他进一步展开了。因为当他说：\par

第7、8节。\BibleQuote{因为我们中没有一人为自己活，也没有一人为自己死。因为我们若活，是为主而活；若死，是为主而死。}\BibleRef{罗 14:7-8} 藉此他也使同一道理更清楚。因为那为法律而活的人，怎能是为主而活呢？但这并不是他如此做的唯一原因，他也阻止那急于纠正他人的人，并说服他忍耐，表明天主不可能轻看他们，而会在适当的时候纠正他们。那么，\BibleQuote{我们中没有一人为自己活}的意思是什么？意思是：我们不自由，我们有一位主人，祂愿意我们活，不愿我们死，而且祂对这两者的关心远超过我们。因为藉着这里所说的，他表明祂对我们比我们自己更关心，比我们更看重我们的生命为财富，视我们的死亡为损失。因为我们若死，并不是只为自己死，也是为主人死。但这里的死，他指的是离开信仰的死。然而，这足以使我们相信祂眷顾我们，因为我们活着是为祂，死了也是为祂。但他不满足于此，继续深入。因为在说了\BibleQuote{所以我们或活或死，都是属于主}之后，并从那种死亡转向肉体的死亡，以免他的言语显得严厉，他给出了另一个非常重大的迹象，表明祂对我们的眷顾。那么，这迹象是怎样的呢？\par

第9节。\BibleQuote{基督为此死了，又复活了，且活了过来，为要作死人和活人的主。}\BibleRef{罗 14:9}\par

因此至少让我们说服你，祂顾念我们的救恩。因为如果祂没有对我们如此大的关怀，哪里需要这\OriginalTerm{经世}{\GreekText{οἰκονομία}}（或降生成人）呢？那么，祂既然如此焦急地要使我们成为祂的，甚至取了仆人的形体，并死亡，难道在我们成为祂的之后，祂会轻视我们吗？这不可能，确实不可能！祂也不愿白费如此大的苦心。\BibleQuote{为此他也死了，}\BibleRef{罗 14:9} 就如有人会说：这样的人不会忍心轻视自己的仆人。因为祂顾念自己的钱袋。\BibleRef{出 21:21} 因为我们对金钱的爱，远不及祂对我们救恩的爱。因此祂用来作我们担保的不是金钱，而是自己的血。为此，祂不会忍心放弃那些祂为之付出如此大代价的人。看他也如何显示祂的能力是无法言喻的。因为他说：\BibleQuote{为此他死了又复活了，为要作死人和活人的主。} 上面他说：\BibleQuote{因为我们或活或死，都是属于主的。} 看这多么广大的主权！看这无敌的力量！看这多么精确的眷顾！因为他的意思是：不要只对我谈活着的人，连去世的人祂也照顾。但如果祂照顾去世的人，显然祂也照顾活着的人。因为祂没有忽略任何一点来建立这主权，比人更有把握，尤其除了其他一切之外，为照顾我们。因为人付出金钱，因此紧紧抓住自己的奴隶。但祂自己付出了死亡，一个以如此大代价买来的人的救恩，以及祂以如此焦虑和困苦获得的主权，祂不会视为无价值。但他说这些是为了使犹太化者羞愧，并说服他记起这恩惠的伟大，以及他如何曾死而复活，从法律中没有得到任何益处，并且如果离开这位对他如此关怀的主，逃回法律，将是极其无情的末等行为。因此，在充分攻击他之后，他又缓和下来，说：\par

第10节。\BibleQuote{然而你为什么判断你的弟兄？又为什么轻看你的弟兄？}\BibleRef{罗 14:10}\par

因此他似乎将他们放在同一水平上，但从他所说的，他显示他们之间的区别是大的。首先，藉着\Quote{弟兄}的称呼，他消除了争论，然后又唤起他们想到那可怕的日子。因为在说了\BibleQuote{你为什么轻看你的弟兄？}之后，他接着说：\BibleQuote{因为我们众人都要站在基督的审判台前。}\BibleRef{罗 14:10}\par

他说这些时似乎又在责备那更进步的人，但他也使犹太化者的心混乱，不仅叫他敬畏那对他已完成的恩惠，也叫他害怕将来的惩罚。\BibleQuote{因为我们众人都要站在基督的审判台前。}\BibleRef{罗 14:10}\par

第11、12节。\BibleQuote{因为经上记载：我指着我的永生起誓，上主说，一切膝都要向我跪拜，一切口舌都要承认天主。这样，我们每人都要向天主交自己的账。}\BibleRef{罗 14:11-12}\par

看他在似乎责备对方时，又如何再次使他的心混乱。因为他暗示类似这样的事：这与你何干？难道你该为他受罚吗？但他没有明说，而是以更温和的方式暗示，说：\BibleQuote{因为我们众人都要站在基督的审判台前}\BibleRef{罗 14:10}，并说：\BibleQuote{这样，我们每人都要向天主交自己的账。}\BibleRef{罗 14:12} 并且他引先知作证，证明万有都服从祂，是的，连旧约中的人也绝对服从。因为他不仅说每人要敬拜，还说\Quote{要承认}，即每人要为自己的行为交账。所以你要焦虑，如同看见万有的主坐在审判台上，不要因为脱离恩宠、跑回法律而在教会中制造分裂。因为法律也是祂的。我为何这样说法律？连法律之下和法律之前的人也是祂的。不是法律要向你追账，而是基督，向你并向全人类。看他是如何释放我们脱离对法律的恐惧。然后，为了不让人以为他这样说只是为了恐吓他们一时，而是他自己原本计划中必然达到的，他再次回到同一主题，说：\par

第13节。\BibleQuote{所以我们不可再彼此判断，反倒要这样判断：不放置绊脚石或跌倒的缘由在他弟兄的路上。}\BibleRef{罗 14:13}\par

这并不适用于一方而非另一方；因此它完全可以适用于两者，即那因遵守食物规定而跌倒的进展者，以及那因所给严厉斥责而绊倒的未进展者。但是，我请求你们想一想，我们若是使人跌倒，将要受何等大的惩罚。因为若有一件事本是违反法律的，然而他们不合时宜地责备时，他禁止他们这样做，为使弟兄不致被绊倒而犯罪；当我们使人跌倒，却甚至没有任何要纠正的事时，我们将受到怎样的对待呢？因为若不拯救他人便是有罪（那埋藏元宝的人证明了这一点），那么也使他跌倒又会有什么后果呢？可是，你也许会说：他若因软弱而自己跌倒呢？这正是为什么你应当忍耐。因为他若是坚强的，就不需要如此多的关注。但如今，既然他是较软弱的一类，他就以此为理由需要相当大的照顾。那么，让我们把这给他，并在一切方面承担他的重担，因为我们不仅要为自己的罪交账，也要为那些我们使之跌倒的罪交账。因为若那交账本身已是难以通过的，再加上这些（使我们跌倒的罪），我们将如何得救呢？并且，我们不要以为，若我们能找到同谋者，那就会成为借口；因为这反会加重我们的惩罚。因为蛇所受的惩罚比女人更重，同样女人所受的惩罚也比男人更重；\BibleRef{弟前 2:14}并且依则贝耳所受的惩罚比夺取葡萄园的阿哈布更严厉；因为是她策划了整件事，并使君王犯罪。\BibleRef{列上 21:23}因此，当你成为他人毁灭的根源时，你将比那些被你败坏的人受更重的苦。因为犯罪本身不如引人同犯那样具有毁灭性。因此他说那些人\BibleQuote{不仅自己行恶，还喜欢别人行恶。}\BibleRef{罗 1:32}因此，当我们看见任何人犯罪时，我们不但不要推他们一把，反而要拉他们脱离罪恶的深坑，以免我们除了自己的毁灭外，还要为别人的毁灭受罚。并且我们要时常记念那可怕的审判台、火的河流、永不松解的铁链、无光的黑暗、切齿和咬人的毒虫。\Quote{啊，但天主是仁慈的！}这些难道只是空话吗？那个富翁岂非因轻视拉匝禄而受罚吗？那些愚拙的童女岂非被逐出婚宴吗？那些没有喂养祂的人岂非要进入\BibleQuote{给魔鬼预备的火}\BibleRef{玛 25:41}里去吗？那穿着污秽衣服的人岂非被\BibleQuote{捆起手脚}\BibleRef{玛 22:13}而灭亡吗？那要求偿还一百银币的人岂非被交给刑役吗？论到奸淫者所说的那句话岂非真实：\BibleQuote{他们的虫不死，他们的火不灭}\BibleRef{谷 9:43}？这些难道只是威胁吗？某人回答说：是的。请问你从哪里胆敢作这样的断言，尤其是你在凭自己的意见下判断的时候呢？我就能从祂所说的和所做的来证明相反的事。\BibleRef{若 5:22}因为若你不相信将来的惩罚，至少相信已经发生的惩罚。因为那些已经发生并成为现实的事，当然不是威胁和空话。那么，是谁淹没了全世界，造成了那可怕的毁灭，以及我们整个人类完全的覆灭呢？此后是谁把那些雷轰和闪电倾注在索多玛之地呢？是谁把整个埃及淹没在海中呢？是谁在旷野中烧灭了六十万人呢？是谁烧毁了阿彼兰的会众呢？是谁命令大地开口吞没了科辣黑和达堂的团伙呢？是谁在达味年间一举消灭了七万人呢？我是否还要提到那些个别受罚的人？加音，被交付于不断的惩罚；加尔米的儿子，与他全家一同被石头砸死；或者那在安息日捡柴而受同样惩罚的人；那四十个被猛兽吞吃的孩童，连他们的年龄也未能获得宽恕？若你愿意看到在恩宠时期之后同样的事，就请想想犹太人受了多大的苦，妇女们怎样吃自己的儿女，有的烤，有的用别的方式吃；他们如何被交付于无法补救的饥荒，以及多种而残酷的战争，他们以自己的灾祸之巨，使先前所有的灾难都黯然失色。因为正是基督向她们行了这些事，请听祂在比喻中以及明明白白地如此宣告。在比喻中，如祂说：\BibleQuote{至于那些不愿意我作他们君王的人，把他们拉到这里来，在我面前杀掉。}\BibleRef{路 19:27}还有葡萄园的比喻和婚宴的比喻。而明明白白地，如祂警告说他们将倒在刀锋下，被掳到万国中，地上将有\BibleQuote{民族因混乱而困苦，海洋波涛咆哮，人因恐惧而昏厥。}\BibleRef{路 21:24-26}\BibleQuote{那时必有大灾难，是从有世界以来直到如今没有过的，今后也决不会再有。}\BibleRef{玛 24:21}至于阿纳尼雅和撒斐辣因偷窃几块银钱而受的惩罚，你们都知道。你们没有看见每日的灾难吗？或者这些也没有发生吗？你们没有看见现在有人因饥荒而憔悴？有人患象皮病或身体残缺？有人常处在贫穷中？有人忍受无数不可补救的祸患吗？那么，有些人受罚，有些人不罚，这合理吗？因为若天主不是不义的（祂确实不是不义的），你如果犯罪，也必定要受惩罚。但若因为祂是仁慈的，祂就不惩罚，那么这些人也不该受惩罚。然而，因着你们的这些话，天主甚至在此世惩罚许多人，为叫你们既然不相信威胁的话，至少可以相信报应的事。\par

既然旧事不能如此令你们惧怕，祂便藉着每一世代发生的事，来纠正那些在每一世代中变得懈怠的人。或许有人会说：为什么祂不在这里惩罚所有人？——为给其他人留下悔改的余地。那么，为什么祂不在来世报复所有人？——免得许多人怀疑祂的眷顾。有多少强盗被抓获，又有多少未受惩罚就离开了此生？那么，天主的仁慈在哪里？现在轮到我来问你们。因为假设根本没有一个人受到报复，你们或许可以以此为托辞。但现在，有些人受了惩罚，有些人却没有，尽管他们是更坏的罪人，怎么能说同样的罪没有同样的惩罚是合理的呢？那些受惩罚的人岂不显得受了冤枉？那么，为什么不是所有人在此世都受惩罚？请听祂对这些事的辩护。因为当有些人被倒塌的塔楼压死时，祂对那些对此提出问题的人说：\Quote{你们以为那些人比众人更罪大恶极吗？我告诉你们：不是的；但如果你们不悔改，你们都要同样丧亡。}\BibleRef{路 13:4}——如此告诫我们：当他人受惩罚而自己虽犯了许多过犯却未受罚时，不可自恃。因为如果我们不改变行为，我们必定会受苦。又有人说：我们在此世犯罪时间短暂，怎么会在来世受无尽的惩罚？我倒要问：在此世，那些在瞬间犯下一件杀人罪的人，不是被判处矿坑中的恒常惩罚吗？\Quote{但这不是天主做的，}有人会说。那么，祂如何使那瘫痪了三十八年的人遭受如此大的惩罚？因为祂是因罪惩罚他，请听祂说的：\Quote{看，你已经痊愈了，不要再犯罪了。}\BibleRef{若 5:14}可是，有人说，他得到了解脱。但来世的情况并非如此。因为祂亲口说：\Quote{他们的虫不死，他们的火也不灭。}\BibleRef{谷 9:44}并且\Quote{这些人要进入永生，那些人要进入永罚。}\BibleRef{玛 25:46}既然生命是永恒的，惩罚也是永恒的。你们没看见祂如何严厉地威胁犹太人吗？那么，那些威胁的事实现了吗？还是对它们说的话只是空谈？\Quote{没有一块石头会留在另一块上。}\BibleRef{路 21:6}那么，石头留下了吗？但是，当祂说\Quote{那时要有大灾难，是从未有过的}时，\BibleRef{玛 24:21}难道没有发生吗？请读\Person{若瑟夫}的史书，你们即使只是听到他们因自己的行为所受的苦难，也会无法呼吸。我这样说，不是为使你们痛苦，而是为使你们稳妥，免得我过于迁就你们，反而为你们招致更严厉的惩罚。因为，请问，为什么你们认为犯罪不应受罚？祂岂不是预先告诉你们一切？岂不是警告过你们？不是来援助你们？为你们的救恩行了无数的事？祂不是给了你们重生的水，并赦免了你们先前所有的罪吗？在这次赦免和\OriginalTerm{圣洗}{laver}之后，祂不是还在你们犯罪时给了你们悔改的帮助吗？祂不是使此后获得罪赦的途径对你们变得容易了吗？那么，请听祂所吩咐的：\Quote{如果你们宽恕你们的近人，我也要宽恕你们}\BibleRef{玛 6:14}，祂说。这有什么艰辛呢？\Quote{你们应为孤儿伸冤，为寡妇辩护，然后来，我们一同辩论，}祂说，\Quote{如果你们的罪如朱红，我要使它们洁白如雪。}\BibleRef{依 1:17}这有什么劳苦呢？\Quote{陈述你们的罪，好使你们成义。}\BibleRef{依 43:26}这有什么艰辛呢？\Quote{以施舍赎你的罪过。}\BibleRef{达 4:24}这有什么烦难呢？税吏说：\Quote{天主，可怜我这个罪人吧！}，就\Quote{回家去，成了义人。}\BibleRef{路 18:13}模仿税吏有什么劳苦呢？即使这样，你们还是不肯相信有惩罚和报复吗？这样下去，你们甚至会否认魔鬼受惩罚。因为祂说：\Quote{离开我，到那为魔鬼和他的使者预备的烈火里去吧！}\BibleRef{玛 25:41}如果没有地狱，那么魔鬼也没有受惩罚。但如果他受惩罚，显然我们也会受惩罚。因为我们也违命了，即使方式不同。你怎么能不害怕说出这样大胆的话呢？因为当你说天主是仁慈的，祂不惩罚时，如果祂惩罚，祂在你身上就不再是仁慈的了。你看魔鬼把你引向什么样的言论？而那些占据山岭、做出如此多克己榜样的隐修士们，难道要空手而去，没有冠冕吗？因为如果恶人不受惩罚，也没有人得到报偿，也许有人会说，善人也不会得到冠冕。不，有人会说，因为与天主相称的是，只有天国，而没有地狱。那么，娼妓、奸夫和行无数恶事的人，要与显示节制和圣洁的人享受同样的好处吗？\Person{保禄}要与\Person{尼禄}站在一起，甚至魔鬼要与保禄站在一起？因为如果没有地狱，却仍有众人的复活，那么恶人也要获得同样的好处！谁会这样说？即使是疯狂的人中，有谁会这样说？或者，即使是魔鬼中，有谁会这样说？因为他们也承认有地狱。因此他们喊叫说：\Quote{你来到这里，在时期还没有到之前就折磨我们吗？}\BibleRef{玛 8:29}\par

那么，你怎能不畏惧战栗，连魔鬼都承认你所否认的事呢？或者，你如何看不见这邪恶教义的导师是谁？因为那欺骗了第一个人，并以更美好希望的借口，甚至将他们从已拥有的福乐中抛出者，正是他现在唆使人说出并幻想这些事。为此，他劝诱一些人怀疑没有地狱，好把他们推入地狱。反之，天主却以地狱相威胁，并准备了地狱，使你通过认识它而如此生活，以免落入地狱。然而，如果真有地狱，魔鬼还劝你这些事，那么魔鬼们又怎能承认它——若它不存在——而他们的目的和愿望正是要我们丝毫不怀疑这类事，好使我们因无所畏惧而更加懈怠，从而与他们一同坠入那火中？那么（有人会说），他们是如何承认的呢？那是因为他们无法忍受加于他们的强制。既然如此，让那些这样说话的人停止欺骗自己和他人吧，因为他们连这些话也将受到惩罚，因其诋毁（\OriginalTerm{诋毁}{\GreekText{διασύροντες}}）了那些可怕的事，并松懈了许多有心认真之人的毅力，甚至做得还不如那些野蛮人：他们虽然一无所知，但一听说那城要被毁灭，不但没有不信，反而呻吟，披上苦衣，惊慌失措，并不停地用尽一切方法，直到平息了愤怒。\BibleRef{纳 3:5} 但你，既经历了如此多的事实和教导，却轻视所告诉你的事吗？那么，你的命运将是相反的。因为，正如他们因惧怕话语而实际上未遭受报复，同样，你因轻视言语的威胁，将在实际上遭受惩罚。如果现在你被告知的事在你看来像寓言，但在那一日，当事实本身说服你时，它就不会像寓言了。你难道没有注意到祂在这世上所做的事吗？当祂遇到两个盗贼时，祂并不认为他们配得同样的境遇，而是将一人领入天国，将另一人送入了地狱？我何必提强盗和杀人犯呢？因为连宗徒，当他成了叛徒时，祂也没有宽恕，反而在看见他冲向绞索、悬挂、并从中间爆裂时（因为他确实\Quote{爆裂，他的全部肠子都流了出来}）\BibleRef{宗 1:18}，尽管祂预见了这一切，仍然让他照样受苦，以此从现在给你来世一切的凭据。所以，不要被魔鬼说服而欺骗自己。这些诡计是他的。因为如果连法官、主人、教师和野蛮人都尊重善人、惩罚恶人，那么在天主那里却相反，善人和非善人被视为配得同样的境遇，这有道理吗？他们何时才会放弃邪恶呢？因为那些现在期待惩罚、处于如此多的恐怖（来自法官和法律）之中，却仍不因此离开不义的人；当他们离开此生时，甚至要放下这种恐惧，不但不被投入地狱，反而获得天国——那时他们何时才会停止作恶呢？请问，这算是慈悲吗？为邪恶加添力量，为不义设立赏报，将清醒者与不受约束者、信徒与不信者、保禄与魔鬼视为同等报应？我还要胡扯到何时？因此，我劝你们脱离这种疯狂，成为自己的主人，说服你们的灵魂恐惧战栗，好使它们立即从将来的地狱中被拯救出来，并在现世清醒地度生后，藉着恩宠和爱人等达到未来的福乐。\par

\SourceRecord{Translated by J. Walker, J. Sheppard and H. Browne, and revised by George B. Stevens.；Nicene and Post-Nicene Fathers, First Series；Vol. 11.；Edited by Philip Schaff.；Buffalo, NY: Christian Literature Publishing Co.,；1889.；<http://www.newadvent.org/fathers/210225.htm>.}

% structure: p0001 source=h1 target=section reason=page-title

\section{论\WorkTitle{罗马书}讲道 第廿六篇}

\BibleRef{羅 14:14}\par

\BibleRef{羅 14:14} \BibleQuote{我知道，并且因主耶稣深信，没有一样东西本身是不洁的；但若有人以为某物是不洁的，那为他便是不洁的。}\par

在最初斥责那论断弟兄的人，并劝他放弃这种指责之后，他进一步解释教义部分，以冷静的口吻教导软弱者，在此事上也显示出极大的温良。因为他没有说他要受惩罚，也没有任何类似的话，只是解除他在此事上的恐惧，并且是为了让他更容易被他的话说服；他说：\BibleQuote{我知道，并且深信。} 然后，为了阻止那些不相信他的人（或\Quote{相信}，\OriginalTerm{不信的人}{\GreekText{τῶν} \GreekText{οὐ} \GreekText{πιστῶν}}）说：\Quote{你深信与我们何干？因为你所说的，不能当作可靠的证据，与如此伟大的法律和天上降下的神谕相抗衡。}他继续说道：\BibleQuote{在主内。} 意思是，因从他学习，因从他获得信心。那么，这个判断不是出于人的心思。你所深信和知道的是什么？告诉我们。\BibleQuote{没有一样东西本身是不洁的。} 他说，本性上无物是不洁的，但因人使用它时的心意而成为不洁。因此，它只对他自己成为不洁，而非对所有人。\BibleQuote{因为若有人以为某物是不洁的，那为他便是不洁的。} 那么，为何不纠正你的弟兄，使他认为它不是不洁的？为何不以全权使他脱离这种心念习惯和观念，使他永远不将它视为平常？他说，我的理由是，我害怕使他忧愁。因此他接着说：\par

第十五节。\BibleRef{羅 14:15} \BibleQuote{如果因你的食物而使你的弟兄忧愁，那么你行事就不是出于爱德。}\par

你看他现在对弟兄的关爱到了何种程度，显示他如此看重他，以至于为了避免使他忧愁，甚至不敢命令极其紧急的事，而是以屈就的态度，宁愿借爱德将他吸引到自己身边。因为即使他解除了他的恐惧，他也不生拉硬拽，而是让他自己做主。因为禁止人食用某些食物，并不像使之充满忧愁那样严重。你看他多么强调爱德。这是因为他知道爱德能够成就一切。基于此，他向对方提出了稍大的要求。因为他说，不仅不应该让他们使你忧愁，反而甚至，如果有必要，他们也不应犹豫迁就你。因此他接着说：\BibleQuote{不要因你的食物毁灭了他，基督曾为他而死。} 难道你为了以戒除食物为代价来换取你弟兄的救恩，都不愿看重他吗？然而，基督不仅没有拒绝成为奴仆，也没有拒绝为他而死；但你不屑于用甚至食物来拯救他吗？而且，尽管基督并没有赢得所有人，但他仍然为众人死了，尽了自己的本分。但你意识到没有，你因食物而在更重要的事上推翻了他，还为此争论？那为基督如此关爱的弟兄，你竟视为如此可鄙，而羞辱他所爱的人吗？他不仅为软弱者而死，甚至为仇敌而死。而你，难道连为了那软弱者也不肯放弃食物吗？然而，基督行了极伟大的事，而你在最小的事上也不愿付出。他是主，你是弟兄。这些话足以堵住他的口。因为这些话显示他心胸狭隘，在领受了天主极大的恩惠之后，连微小的事也不肯回报。\par

第十六、十七节。\BibleRef{羅 14:16} \BibleRef{羅 14:17} \BibleQuote{所以，不要让你们的善被人诋毁。因为天主的国不在于饮食。}\par

在这里，他说的\Quote{善}指的是他们的信德，或将来酬报的望德，或他们宗教状态的完美。因为他说，你不仅未能使你的弟兄受益，反而使教义本身、天主的恩宠和他的恩赐被人诋毁。现在，当你争斗、争吵、烦扰、在教会中制造分裂、责备你的弟兄、疏远他时，外人就会说你坏话。因此，善远非由此而来，恰恰相反。因为你的善是爱德、爱弟兄之情、合一、联系、平安生活和生活在温良中（\OriginalTerm{温良}{\GreekText{ἐπιεικείας}}）。他再次为了消除他的恐惧和另一方的争辩，说：\BibleQuote{因为天主的国不在于饮食。} 他意思是，难道凭这些我们就能蒙悦纳吗？正如他在另一处经文也说：\BibleQuote{我们若吃，并不因此更好；若不吃，也不因此更差。} 他不需要任何证明，只满足于陈述。他所说的是：你若吃，这就能引你进入天国吗？正因如此，为了讽刺他们在此事上自鸣得意，他不仅说\Quote{肉}，还说\Quote{酒}。那么，是哪些事物将我们带到天国呢？\BibleQuote{正义、平安和喜乐}，以及有德的生活、与弟兄们的平安（与此好争吵相对立）、来自同心合意的喜乐，而这样的责备却使这喜乐终止。但他这番话并非只对一方，而是对双方而言，此时正是对双方说话的适当时机。然后，既然他提到了平安和喜乐，但还有基于恶行的平安和喜乐，他便加上：\BibleQuote{在圣神内。} 因为那毁灭他弟兄的人，一下子就颠覆了平安，伤害了喜乐，比抢劫钱财的人更严重。更糟的是，另一位拯救了他，而你却伤害并毁灭他。既然吃和所谓的完美状态并不带来这些德行，反而带来了破坏它们的东西，那么，轻视小事以保证大事的稳固，岂不是正确的吗？既然这种责备在某种程度上是出于虚荣，他接着说：\par

第十八节。\BibleRef{羅 14:18} \BibleQuote{因为凡在这几件事上服侍基督的，必为天主所悦纳，也为众人所称许。}\par

因为他们不会因为你的完美状态而如此钦佩你，反而所有人都会因和平与友好而钦佩你。因为这是一件美事，每个人都将获益，而那种完美状态却连一个人也获益不了。\par

19节。\BibleQuote{因此，我们当追求和平的事，以及彼此建立德行的事。}\BibleRef{罗 14:19}\par

这对前一个人适用，使他变得和平。但另一点对后一个人也适用，使他不要毁坏他的弟兄。然而，他再次将两者都应用于双方，通过说\Quote{彼此}，表明没有和平就不容易建立德行。\par

20节。\BibleQuote{不可因食物毁坏天主的工程。}\BibleRef{罗 14:20}\par

他以此称呼弟兄的救恩，并大大增加了恐惧，表明他正在做与他所期望相反的事。因为他说，你非但没有像你所打算的那样建立，反而毁坏，而且那工程不是人的，而是天主的，且不是为了什么伟大的目标，而是为了一件琐事。因为他说，那是\Quote{因食物}。然后，免得如此多的纵容使软弱的弟兄在他的误解中更加坚定，他又变得教义化，如下所述：\par

\Quote{一切固然是洁净的，但若有人带着冒犯而吃，那对他便是恶的。}\par

也就是，带着不良的良心去做。所以如果你强迫他，他却吃了，那将毫无益处。因为不是吃使人不洁，而是人吃的意图。那么，如果你不纠正这一点，你就白费功夫，并且使情况更糟：因为认为某物不洁，还不如在认为它不洁时去尝它那样糟糕。在这里，你犯了两项错误：一是通过你的争吵增加他的偏见，二是让他尝了不洁之物。因此，只要你没有说服他，就不要强迫他。\par

21节。\BibleQuote{不吃肉，不喝酒，不涉足任何使你弟兄绊倒、或冒犯他、或使他软弱的事，这原是好的。}\BibleRef{罗 14:21}\par

他再次要求更大的取舍，即他们不仅不要强迫他，甚至要俯就他。因为他自己也常常这样做，例如当他行割礼时（\BibleRef{宗 16:3}），当他剃发时（\BibleRef{宗 18:18}），当他献上犹太人的祭献时（\BibleRef{宗 21:26}）。他并没有对那人说\Quote{这样做}，而是以格言的形式陈述，以免再次使那软弱的弟兄更加懈怠。他的话是什么？\Quote{不吃肉是好的。}我为什么提到肉呢？如果是酒，或任何其他类似的东西，引起了冒犯，就要戒绝。因为没有什么比你弟兄的救恩更重要。基督也向我们显示了这一点，因为他从天降下，为我们忍受了他所经历的一切。请观察他如何用\Quote{绊倒、或冒犯、或使他软弱}这些话来同样告诫那前一个人。他的意思是：不要告诉我他的软弱是没有理由的，但你有能力去纠正它。因为那弟兄的软弱有足够的理由得到帮助，而对你来说这不是损失，不是伪善的情况（\BibleRef{迦 2:13}），而是建立德行和权宜之计。因为如果你强迫他，他立刻被毁坏，并将谴责你，并更加坚持不吃。但如果你俯就他，他就会爱你，不会怀疑你是一位导师，之后你就能获得机会，不知不觉地将正确观点播种在他心中。但一旦他恨你，你就关闭了你的推理进入的门。因此，不要强迫他，反而要为他克制自己，不是认为它不洁而克制，而是因为他被冒犯，而他将会更爱你。保禄也这样建议，当他说\Quote{不吃肉是好的}时，不是因为肉不洁，而是因为弟兄被冒犯、软弱。\par

22节。\BibleQuote{你有信德吗？你就在自己内持守它。}\BibleRef{罗 14:22}\par

在这里，在我看来，他是在温和地警告那更进步的人，关于虚荣心。他说的是：你想向我表明你是完美的、充分装备的吗？不要向我表明，让你的良心足够。在这里，他所指的\Quote{信德}不是关于教义的，而是关于当前的主题。因为关于前者，经上说：\BibleQuote{口里宣认，便得救恩}\BibleRef{罗 10:10}；以及\BibleQuote{凡在人前否认我的，我也必否认他}\BibleRef{路 9:26}。因为前者若不宣认，就会毁了我们；而后者若不恰当地宣认，也会如此。\Quote{凡在自己所认可的事上不自责的，便是有福的。}他再次打击那软弱的，并给他（即那坚强的）足够的冠冕，在于他的良心。即使没有人看见，也就是说，你能在自己内快乐。因为在他说了\Quote{你就在自己内持守它}之后，免得他认为这是一个可鄙视的裁判所，他告诉他，这对你来说比整个世界更好。如果所有人都责备你，而你不自责，你的良心也不指控你，你就是有福的。但这一陈述他并没有应用于任何人。因为有许多人并不自责，却是大罪人：这些人是人间最痛苦的。但他仍然紧扣当前的主题。\par

23节。\BibleQuote{但那疑惑的，若吃了，就被定罪。}\BibleRef{罗 14:23}\par

他这样说，再次是为了劝勉他要宽容那软弱的。因为如果他疑惑着吃，并且定自己的罪，那有什么好处呢？我赞同那既吃又不疑惑的人。看他如何引导他不仅吃，而且带着好良心吃。然后他接着也提到了他被定罪的理由，继续这些话：\par

\BibleQuote{因为他吃不是出于信德。}不是因为它不洁，而是因为不是出于信德。因为他不相信它是洁净的，却在不洁的状态下吃了它。但借此他也向他们表明，他们强迫人、而不是说服人去触碰那些他们至今仍视为不洁之物，造成了多大的危害；为的是至少因这一点他们可以停止责备。\BibleQuote{因为凡不出于信德的，都是罪。}因为当一个人不确定、也不相信一件事是洁净时，他除了犯罪还能做什么呢？现在\Person{保禄}所说的这一切，都是针对当前讨论的事，并非泛指一切。并且请注意他多么小心谨慎，不冒犯任何人；他之前就说过：\BibleQuote{如果你因着食物使你的兄弟心乱，你便不是按爱德行事。}但如果一个人不该使他心乱，就更不该冒犯他。他又说：\BibleQuote{不可因食物而摧毁天主的事业。}因为如果推倒一座教堂是严重的罪行，那么对属神的圣殿这样做就更严重了。因为一个人比一座教堂更有尊严：因为\Person{基督}不是为墙壁而死，而是为这些圣殿。\par

那么，让我们在各方面留意自己的行为，不给任何人丝毫的把柄。因为今世的生命就如赛跑场，我们应在各方面有成千的眼目\\EditorialNote{希拉利注《圣咏》第一一九篇}，甚至不要幻想无知将是充分的借口。因为确实有，确实有这样一种情况：因无知而受罚，当这种无知是不可原谅的时候。因为犹太人也曾无知，但他们的无知并非情有可原。外邦人也无知，但他们无可推诿。\\BibleRef{罗 1:20}因为当你对那些无法认识的事无知时，你并不会因此受任何控告；但如果是对简易且可能的事无知，你将受到最严厉的惩罚。否则，如果我们不是过分怠惰，而是充分贡献自己的一份力量，那么天主也必在我们无知的事上向我们伸出援手。这正是保禄同样对斐理伯人所说的。\\BibleQuote{你们若在任何事上另有想法，天主也要将这启示给你们。}\\BibleRef{斐 3:15} 但是，当我们不愿意做那些自己力所能及的事时，我们在这一点上也不会得到祂的助佑。犹太人的情况也是如此。\\BibleQuote{为此，}祂说，\\BibleQuote{我用比喻对他们讲话，因为他们看却看不见。}\\BibleRef{玛 13:13} 他们看却看不见，这是什么意思呢？他们看见魔鬼被驱逐，却说祂附有魔鬼。他们看见死人复活，却不朝拜，反而企图杀害祂。但科尔乃略却不是这样的人。\\BibleRef{玛 12:24}因此，当他诚心尽全部本分时，天主也给他加添了所缺少的。那么不要说：为何天主忽略了某某人，他并非\\OriginalTerm{不拘泥形式的人}{\GreekText{ἄ} \GreekText{πλαστος}}，又是好人，虽然他是外邦人？首先，没有人能确切知道一个人是否不拘泥形式，唯有祂，\\BibleQuote{\\OriginalTerm{塑造}{\GreekText{πλάσαντι}}众心的那一位。}\\BibleRef{咏 33:15}此外还有一点要说的：\\OriginalTerm{或许}{\GreekText{πολλάκις}}这样的人既不深思也不热切。也许有人会说，他怎么能这样呢，因为他非常无知？\\OriginalTerm{不拘泥形式的人}{\GreekText{ἄ} \GreekText{πλαστος}}那么，我请你考虑这个纯朴而专一的人，在生活事务中观察他，你将看到他是最谨慎小心的榜样，以致如果他能在属灵的事情上也表现出同样的谨慎，就不会被忽视：因为真理的事实比太阳还要清晰。无论一个人走到哪里，他都能轻易抓住自己的救恩，只要他愿意留心，不把此事当作副业。难道这些事仅限于巴勒斯坦，或世界的小角落里吗？你们没有听先知说：\\BibleQuote{从最小的到最大的，都要认识我。}\\BibleRef{耶 31:34}难道你们没有看见事情本身在宣告真理吗？那么，他们看见真理的教义广泛传播，却不肯费心去学习，又怎能被原谅呢？有人问：你要求一个粗野的蛮族做到这一切吗？不，不仅是粗野的蛮族，而且包括任何比现今的人更野蛮的人。因为请告诉我，为什么在今世的事务中，他知道在受冤屈时如何答辩，在受暴力时如何抵抗，并做一切事，想一切办法，以避免自己的意愿被丝毫挫败；但在属灵的事情上，他却没有运用同样的判断力？当一个人崇拜石头，认为它是神时，他为其举行庆节，花费金钱，并表现出极大的敬畏，绝不因他的单纯而变得冷漠。但当他必须寻求真实的天主时，你却对我提起单纯和淳朴？这些事并非如此，绝非如此！因为这些抱怨只是出于怠惰。你认为哪些人最单纯粗野，是亚巴郎时代的人还是现今的人？\\BibleRef{苏 24:2}显然是前者。那么，在哪一个时代更容易发现宗教呢？显然是现在。因为现在天主的名号甚至被所有人传扬，先知们宣讲了，事情成就了，外邦人被说服了。\\BibleRef{创 32:29}但在那时，大多数人仍处于未受教化的状态，罪恶当道。没有法律来教导，没有先知，没有奇迹，没有教义，没有众多认识它的人，也没有任何类似的事物，一切事物都好像处于深沉的黑暗中，一个无月且狂暴的黑夜。然而，即使在那时，那位奇妙而高贵的人，虽然障碍如此巨大，仍然认识天主并实践德行，并引领许多人同样热忱；尽管他甚至连外邦人的智慧也没有。他怎么可能有呢？那时连文字都尚未发明。但他仍然贡献了自己的一份，天主也合作加添了他所缺少的。因为你甚至不能说亚巴郎从他的祖先接受了宗教，因为\\BibleRef{苏 24:2}他曾是偶像崇拜者。尽管如此，尽管他来自这样的祖先，未开化，生活在未开化的人中间，没有宗教导师，他却达到了对天主的认识，与所有他的后代相比——他们既有法律又有先知的优势——他的光辉是无可言喻的。这是为什么呢？这是因为在今世的事上他没有给自己过多的焦虑，而在属灵的事上全心投入。（\\WorkTitle{创世纪讲道集}第三十三篇等）默基瑟德又如何？他不也是出生于那个时代，如此光明以至于被称为天主的司祭吗？（\\WorkTitle{创世纪讲道集}第三十五、三十六篇）因为\\OriginalTerm{清醒的人}{\GreekText{νήφοντα}}绝不可能被忽视。不要让这些事困扰我们，但要知道，在每种情况下，权力在于心志，让我们关注自己的本分，以便成长得更好。不要追查天主的安排，或查问祂为何忽略这个人而召叫那个人。因为我们这样做，就像得罪了主人的仆人，却去窥探主人的家务。可怜而可悲的人！当你应当关心自己要交的账，以及如何与主人和好时，你却去查问那些你无需交代的事，而忽略了你必须交代的事。我对外邦人说什么呢？有人问。为什么，就是我一直说的。不要只看你对外邦人说什么，也要看如何改善自己。当他因审视你的生活而被冒犯时，你要考虑你将说什么。因为如果他因你的生活方式而被冒犯，你不必为他被召算账；但如果是你的生活方式伤害了他，你将面临最大的危险。当他看到你谈论天国却为今生之事焦虑，既害怕地狱又为今世的灾祸战栗时，你要记在心里。当他看到这一点并指责你，说：如果你热爱天国，为何不轻视今世之事？如果你期待可怕的审判，为何不藐视今世的恐怖？如果你盼望不朽，为何不轻视死亡？当他说这些话时，你要担心你将如何辩护。当他看到你因失去金钱而战栗，你期待天堂，却为一分钱而极度欢喜，并为了小钱再次出卖灵魂时，你要记在心里。因为正是这些，正是这些，使外邦人跌倒。因此，如果你关心他的得救，你要在这些方面辩护，但不是用言语，而是用行动。因为从来没有人因那个问题而亵渎天主，而是因人们的恶行，各处都有无数的亵渎。所以要纠正他。因为外邦人接着会问你：我如何知道天主的命令是可行的？因为你，有基督徒血统，并被这种美好宗教养育的人，却不做任何这类的事。你将对他说什么？你一定会说：我指给你看别人这样做：住在沙漠的隐修士。你承认自己是基督徒，却把他人指向别人，好像不能表现出基督徒的气魄，难道不羞愧吗？因为他也会直接说：我何必去山上，去寻找沙漠？如果居住在城中的人不能做门徒，那是对这种行为准则的严重谴责，即我们要离开城市，跑到沙漠去。但请给我看一个有妻子、儿女、家庭，却追求智慧的人。我们对此能说什么呢？我们岂不要低下头，感到羞愧吗？因为基督并没有给我们这样的命令；而是什么？\\BibleQuote{你们的光当在人前照耀，}\\BibleRef{玛 5:16}而不是在山上、沙漠、旷野和偏僻之处。我说这话，并非责备那些住在山上的人，而是哀叹那些住在城里的人，因为他们把德行从那里放逐了。因此我请求你们，让我们把他们那里的纪律也带到这里来，使城市真正成为城市。这将改善外邦人。这将使他免于无尽的绊脚石。因此，如果你想使他免于丑闻，并享受无数奖赏，就要整顿自己的生活，使其在各方面发光，\\BibleQuote{好使他们看见你们的善行，光荣你们在天上的父。}\\BibleRef{玛 5:16}这样，我们也必享受那不可言喻的伟大荣耀，愿天主赐予我们都能达到，藉着恩宠和爱人，等等。\par

\SourceRecord{Translated by J. Walker, J. Sheppard and H. Browne, and revised by George B. Stevens.；Nicene and Post-Nicene Fathers, First Series；Vol. 11.；Edited by Philip Schaff.；Buffalo, NY: Christian Literature Publishing Co.,；1889.；<http://www.newadvent.org/fathers/210226.htm>.}

% structure: p0001 source=h1 target=section reason=page-title

\section{罗马书讲道集 第27篇}

\BibleRef{罗 14:25-27}\par

\Quote{\Emph{如今愿那能照我的福音和耶稣基督的宣讲，按照奥秘的启示，这奥秘从世界创立以来一直隐藏，如今却彰显出来，并且（手抄本} \Emph{\OriginalTerm{\GreekText{τε}}{\GreekText{τε}} 系Sav. 省略）藉着先知们的经书，按照永生的天主的命令，向万民宣示，为信德的服从：归给独一智慧的天主，藉着我们的主耶稣基督，愿荣耀归于祂，直到永远。阿们。}}\par

保禄常以祈祷和赞颂结束劝言，因为他深知此事非同小可。出于慈爱与谨慎，他养成此习惯。因为一位忠于子女、忠于天主的教师，不仅以言语教导，更藉祈祷将天主的助佑带入教导。在此处他也如此行。但上下文如下：\Quote{愿那能坚固你们的天主，荣耀归于祂，直到永远。阿们。}因为他再次贴近那些软弱的弟兄，向他们发言。当他责备时，他使众人同受责备；如今祈祷时，他却为这些人代求。说过\Quote{坚固}后，他接着说明方式：\Quote{按照我的福音}；这显示出他们尚未坚稳，虽站立却动摇。为使他的话可信，他接着说：\Quote{和耶稣基督的宣讲}，即祂亲自宣讲的。若是祂宣讲的，则教义不是出于我们，而是祂的法则。随后，在讨论宣讲的性质时，他表明这恩赐大有裨益且极其荣耀；先以宣讲者的身份证明，再以宣讲的内容证明。因为这是喜讯。此外，因祂未曾向任何人揭示这些事，只向我们启示。这由以下话语暗示：\Quote{按照奥秘的启示}。这是最大友谊的标志：让我们分享奥秘，且无人先于我们。\Quote{这奥秘从世界创立以来一直隐藏，如今却彰显出来。}因为久已预定，如今才彰显。如何彰显？\Quote{藉着先知们的经书。}此处他再次释放软弱者的恐惧。因为你惧怕什么？是怕偏离法律吗？法律正期望如此，它自古已预言。但若你追究为何如今彰显，便是越分，探究天主的奥秘并要祂解释。对于这类事，我们不当好奇，而当满足。因此他亲自制止此种精神，加上：\Quote{按照永生的天主的命令，为信德的服从。}因为信德要求服从，而非好奇。当天主命令时，人当服从，而非好奇。接着他用另一论据鼓励他们，说：\Quote{向万民宣示。}即不独你一人，整个世界都持此信经，因为老师不是人，而是天主。故他又加：\Quote{藉着耶稣基督。}不但宣示，而且确认。两者都是祂的工作。据此，诵读方式应是：\Quote{愿那能藉着耶稣基督坚固你们的天主；}如我所说，他将两者归于祂；或者更正确地说，不仅两者，连那归于圣父（\OriginalTerm{归于}{\GreekText{τὴν} \GreekText{εἰς}}）的荣耀也归于祂。因此他说：\Quote{愿荣耀归于祂，直到永远，阿们。}由于这些奥秘的不可测度，他再次用赞颂。因为即使如今显现，仍无法以推理理解，只能藉信德认识，别无他法。他说得好：\Quote{独一智慧的天主。}因为只要你思考祂如何纳入外邦人，如何将他们与古时善行之人融合，如何拯救绝望者，如何将不配地上之人带上天国，如何使今世堕落者进入不死不变的生命，如何使被魔鬼践踏者与天使争辉，如何开启乐园，止息一切旧恶，且在一段时间内以简易之途成就这一切，你便会认识祂的智慧——当你看到天使、总领天使所不知的事，外邦人却通过耶稣骤然学到。（两份手抄本加：\Quote{那时你将知道祂的能力。}）故当赞美祂的智慧，归荣耀于祂！但你仍拘泥小事，居于阴影之下。这不像归荣耀者。因为对祂无信心、不信赖信德的人，不能为祂的伟业作证。但他亲自为他们献上荣耀，为引导他们也有同样热忱。当你听见他说\Quote{独一智慧的天主}时，不要以为这是贬抑圣子。因为这一切彰显智慧的事都是藉着基督而成，若没有祂无一是，那么显然祂在智慧上也是同等的。\BibleRef{若 1:3} 那么为何说\Quote{独一}？是为与一切受造物对比。献上赞颂后，他再次从祈祷转向劝言，对强者发言如下：\par

\BibleRef{罗 15:1} \Quote{我们刚强的人，应当}——是\Quote{应当}，不是\Quote{我们如此仁慈}。应当做什么？——\Quote{承担软弱者的软弱。}\par

请看祂如何藉赞美激起他们的注意，不仅称他们为强有力者，还将他们与祂自己并列。不仅如此，祂又藉此事的好处再次吸引他们，并因其并不沉重。因为祂说：你是强有力者，屈尊俯就对你毫无损害。但对他而言，若不被包容，则危险是最终的。祂没有说软弱者，却说\Quote{软弱者的软弱}，以此吸引他并使他倾向怜悯。正如在另一处祂也说：\BibleQuote{你们属神的人，要用温和的心扶助这样的人。}（\BibleRef{迦 6:1}）你变得强有力了吗？向使你如此的天主回报吧。而你若纠正病弱者的软弱，便是回报。因为我们也曾软弱，但藉恩宠变得强有力。我们不仅在此事上这样做，也应如此对待在其他方面软弱的人。例如，若有人易怒、傲慢或有任何此类缺失，就包容他。这如何做到呢？请听下文。祂在说\Quote{我们应当担当}之后，接着说：\Quote{且不可求自己的喜悦。}\par

第2节：\BibleQuote{我们各人务要叫邻舍喜悦，使他得益处，建立德行。}（\BibleRef{罗 15:2}）\par

祂所说的意思是：你是强有力者吗？让软弱者体验你的力量。让他认识你的力量，叫他喜悦。祂不仅说喜悦，还说\Quote{使他得益处}，且不单说得益处，以免那长进的人说：看，我在引导他走向益处！祂又加上\Quote{建立德行}。因此，你若富有或有权势，不要叫自己喜悦，而要叫穷乏者喜悦，因为这样你既享有真正的荣耀，又做了许多服务。因为来自世俗的荣耀转瞬即逝，而来自属神之事的荣耀是长存的，只要你为建立德行而行。因此祂要求所有人这样做。因为不是这个人或那个人要做，而是\Quote{你们各人}。然后，由于祂所命令的是一件大事，祂甚至叫他们放松自己的完美以纠正他人的软弱；祂再次引入基督，在下文中说：\par

第3节：\BibleQuote{因为连基督也没有求自己的喜悦。}（\BibleRef{罗 15:3}）\par

这是祂一贯的做法。因为当祂论及施舍时，便引出基督说：\BibleQuote{你们知道我们主耶稣基督的恩宠：祂本是富有的，却为你们成了贫穷。}（\BibleRef{格后 8:9}）当祂劝勉爱德时，便从基督引出劝说：\BibleQuote{正如基督爱了我们。}（\BibleRef{弗 5:25}）当祂给出关于忍受羞辱和危险的劝告时，便投靠基督说：\BibleQuote{祂因那摆在前面的喜乐，轻视了羞辱，忍受了十字架。}（\BibleRef{希 12:2}）同样在本段，祂也显示基督如何做了这事，以及先知如何早已预言。因此祂接着说：\par

\BibleQuote{辱骂祢的辱骂都落在我身上。}（\BibleRef{咏 68:9}）但\BibleQuote{祂不求自己的喜悦}有何含义？祂本可以不受辱骂，本可以不承受所承受的，若祂在意自己的事。但祂却不在意。因祂顾念我们的益处而忽略了祂自己的。为何祂不说\BibleQuote{祂空虚了自己？}（\BibleRef{斐 2:7}）因为祂不仅想指出祂成为人，还想指出祂受到恶待，且在许多人中名声不好，被视为软弱。因为经上说：\BibleQuote{祢若是天主子，就从十字架上下来吧。}（\BibleRef{玛 27:40}）以及\BibleQuote{祂救了别人，却不能救自己。}（\BibleRef{玛 27:42}）因此祂提及一件与当前主题相关的事，证明的远超过祂所承诺的；祂显示不仅是基督受辱骂，连圣父也受辱骂。\BibleQuote{因为辱骂祢的辱骂都落在我身上，}祂说。祂所说的近乎此：所发生的事并非新事或怪事。因为在旧约中那些惯于辱骂祂的人，也向祂的圣子发狂。但这些事被记录下来，是叫我们不效法他们。然后祂装备（希腊文：\OriginalTerm{膏抹}{\GreekText{ἀλείφω}}）他们以忍耐地经受试探。\par

第4节：\BibleQuote{从前所写的圣经都是为教训我们写的，叫我们因圣经所生的忍耐和安慰，可以得着盼望。}（\BibleRef{罗 15:4}）\par

这就是说，叫我们不跌倒（因为内有外有各种争战），藉着圣经的激励和安慰，我们得以忍耐，藉着活在忍耐中，我们持守盼望。因为这些事互相产生：忍耐产生盼望，盼望产生忍耐。两者皆由圣经而来。然后祂又以祈祷的形式继续讲论：\par

第5节：\BibleQuote{但愿赐忍耐和安慰的天主，使你们彼此同心，效法基督耶稣。}（\BibleRef{罗 15:5}）\par

因为祂既给了自己的劝告，又引用了基督的榜样，还加上了圣经的见证，以显示圣经本身也赐予忍耐。因此祂说：\Quote{但愿赐忍耐和安慰的天主，使你们彼此同心，效法基督耶稣。}因为这就是爱所要做的：对待他人如同对待自己。然后为显示祂所要求的不仅是普通的爱，祂加上\Quote{效法基督耶稣}。祂处处如此，因为还有其他种类的爱。他们同心合意有何益处呢？\par

第6节：\BibleQuote{叫你们同心合意地荣耀天主，我们主耶稣基督的父。}（\BibleRef{罗 15:6}）\par

祂不仅说用同一张嘴，还吩咐我们也要有同一意志。请看祂如何将全身结合为一，并以赞美天主作为结束，从而给予合一与和谐最大的动力。然后从这一点祂又回到先前的劝勉中，说：\par

第7节：\BibleQuote{所以你们要彼此接纳，如同基督接纳了你们，使荣耀归于天主。}（\BibleRef{罗 15:7}）\par

例子再次与先前相同，而收益不可言喻。因为这是一件特别光荣天主的事，即紧密结合。所以，即使违背你的意愿（Field\Quote{为祂的缘故而忧伤}，依据Savile，但违反手稿）并为祂的缘故，你与弟兄不和，要想到，平息你的愤怒就是光荣你的主；若不为弟兄的缘故，无论如何要为此与他和好：或者更确切地说，先为这缘故。因为基督也以各种方式坚持这一点，当向圣父说话时，他说：\BibleQuote{众人因此就会认出你派遣了我，如果他们合而为一。}\BibleRef{若 17:21}\par

让我们服从，并彼此结合。因为在此，他不再只激励软弱者，而是激励众人。若有人决意与你决裂，你也不可决裂。不要说那冷漠的话：\Quote{我爱那爱我的人；若我的右眼不爱我，我就把它挖出来。}因为这些是撒旦的话语，适合税吏和外邦人的小气精神。但你被召叫归属更高的国籍，且名字记在天上的册中，须遵守更高的律法。不可这样说；反而当他不愿爱你时，你更要显示更多的爱，好能把他吸引到你身边。因为他是肢体；当任何力量使肢体与身体分离时，我们尽一切努力使之重新结合，并给予更多关注。因为那时的赏报更大，当你把不愿爱的人吸引过来。若祂吩咐我们邀请那些无法回报我们的人赴宴，为使回报更大，那么在友谊方面我们更应如此。被爱又去爱的人，确实回报了你。但被爱却不去爱的人，使天主成为你的债主。此外，当他爱你时，你不需要多大辛劳；但当他恨你时，他正需要你的帮助。因此不要把辛劳的缘由变成懈怠的借口；不要说因为他生病了，所以我不照顾他（因为爱的冷淡确是一种病），你反而要温暖那已冷却的。但假设他仍不愿被温暖，那又怎样？继续尽你的本分。如果他变得更乖僻？这只不过为你赢得更大的回报，并显示你更是基督的追随者。因为如果彼此相爱是门徒的特征（祂说：\Quote{如果你们彼此相爱，众人因此就可认出你们是我的门徒。}\BibleRef{若 13:35}），那么想想爱那恨我们的人是多么伟大。因为你的主爱了那些恨祂的人，并召唤他们到祂面前；他们越软弱，祂越显关怀；祂大声说：\Quote{健康的人不需要医生，而是有病的人。}\BibleRef{玛 9:12}祂认为税吏和罪人堪与祂同席。犹太民族对祂的羞辱越大，祂对他们显示的光荣和关怀也越大，远超过他们应得的。你也要效法祂：因为这善工绝非小事；若无它，即使殉道者也不能大大中悦天主，正如保禄所说。那么不要说，我被人恨，所以我不去爱。正因如此，你更应去爱。此外，按常理，一个爱人的人不会被很快憎恨；即使人如禽兽，也会爱那爱他的人。因为祂说，外邦人和税吏也这样做。\BibleRef{玛 5:46}但如果人人都爱那爱他的人，又有谁不愛那虽被恨仍去爱的人呢？请显出这样的品行，不断说这句话：\Quote{你尽管恨我，我绝不会停止爱你。}这样，你将折服他的好争吵之心，驱散一切冷漠。因为这种混乱要么源于过度发热（\OriginalTerm{发热}{\GreekText{φλεγμονἥς}}），要么源于寒冷；而爱的力量常以其温暖矫正这两者。你们岂没见过那些沉溺于卑劣爱情的人，被那些娼妓殴打、唾弃、辱骂、百般虐待？那么，侮辱能中断这种爱情吗？绝不；反而更燃起它。然而，做这些事的，除了是娼妓，还是名声败坏、身份低微的人。但忍受这些的人，往往可以数算显赫的祖先和许多其他可夸口的贵族身份。然而，这些事都无法打断他们的联系，也不能使他们远离所爱的人。那么我们岂不羞愧，发现魔鬼和恶魔的爱竟有如此强大的力量（v. p. 520），而我们却不能为合乎天主的爱显示同等力量吗？你们不知道这是对抗魔鬼的极好武器吗？你们没有看见那邪恶的魔鬼站在一旁，把你们所恨的人拉向自己，想要夺走那肢体吗？而你们却跑开，放弃战斗的奖品？因为你们的弟兄被夹在你们中间，就是那奖品。如果你得胜，你将接受冠冕；如果你懈怠，你就空手而去。那么，不要再说出那撒旦的话：\Quote{若我的眼恨我，我不能看它。}没有什么比这话更可耻的了，然而一般人却把它当作高尚精神的标志。但这一切没有比这更卑鄙、更没头脑、更愚蠢的了。因此我实在悲痛：恶行竟被视为美德；轻视人、藐视人竟被看作光荣和尊严。这是魔鬼最大的圈套：给不义戴上好名声，使之难以抹去。因为我常听人夸耀自己不去接近那些厌恶他们的人。然而你们的主却以此为荣。人们多少次鄙视（\OriginalTerm{鄙视}{\GreekText{διέπτυσαν}}）祂？多少次厌恶祂？但祂不停跑向他们。那么不要说\Quote{我无法接近那些恨我的人}，而要说\Quote{我无法鄙视（\OriginalTerm{鄙视}{\GreekText{διαπτύσαι}}）那些鄙视我的人}。这是基督门徒的话，前者是魔鬼的话。后者使人光荣荣耀，前者使人羞耻可笑。正因如此，我们钦佩梅瑟，因为即使当天主说：\Quote{你且由我向他们发怒，消灭他们。}\BibleRef{出 32:10}他也不忍鄙视那些多次厌恶他的人，反而说：\Quote{倘若你赦免他们的罪，就赦免；否则，求你从你所写的书上涂抹我的名字。}\BibleRef{出 32:32}这是因为他是天主的朋友，并效法祂。我们不要因那些本应蒙羞的事而自夸。也不要说这些下流人、人中渣滓的话：\Quote{我知道如何唾弃（\OriginalTerm{唾弃}{\GreekText{καταπτύσαι}}）成千上万的人。}即便别人说这话，我们也要嘲笑并阻止他，因为他以本应羞愧的事为乐。请问，你竟鄙视一个信徒？基督在不信时都不鄙视他。为何说\Quote{不鄙视}？祂如此爱那卑贱丑陋之人，甚至为他而死。祂如此爱了这样的人，而此人现在已变得美丽可敬，你却鄙视他？如今他已成为基督的肢体，成为你们主的身体？你岂不知你在说什么？岂不知你冒险做什么？他有基督作头，作桌子，作衣服，作生命，作光明，作新郎，基督是他的一切；你竟敢说\Quote{我鄙视这家伙}？不仅如此，还连同成千上万的人？住口吧，人啊，停止你的疯狂；认识你的弟兄。要知道这是无理和狂妄的话，相反要说：即使他鄙视我一万次，我也绝不远离他。这样，你既赢得你的弟兄，也将生活于天主的荣耀中，并分享未来的福乐。愿天主恩赐我们众人皆能达到，并藉祂的恩宠和爱人之心。\par

\SourceRecord{Translated by J. Walker, J. Sheppard and H. Browne, and revised by George B. Stevens.；Nicene and Post-Nicene Fathers, First Series；Vol. 11.；Edited by Philip Schaff.；Buffalo, NY: Christian Literature Publishing Co.,；1889.；<http://www.newadvent.org/fathers/210227.htm>.}

% structure: p0001 source=h1 target=section reason=page-title

\section{\WorkTitle{罗马书释义 第二十八讲}}

\BibleRef{罗 15:8}\par

\BibleQuote{如今我说，耶稣基督为了天主的真理，成了割礼的仆役，以证实对祖先的恩许。}\par

他再次谈论基督对我们的关怀，仍然持守同一主题，并显示祂为我们做了何等伟大的事，以及祂如何\BibleQuote{没有寻求自己的喜悦}。\BibleRef{罗 15:3}除此以外，他还提出了另一论点，即外邦人欠天主的更多。如果欠得更多，他们就应当容忍犹太人中的软弱者。因为他曾严厉地责备这些人，恐怕这会使那些外邦人骄傲，他压制了他们的无理，指出他们所获得的恩赐是借着\Quote{对祖先的恩许}，而外邦人却只出于怜悯和天主对人的爱。这就是他说\Quote{并使外邦人因天主的仁慈而光荣天主}的原因。但为使所说的话更清楚，最好再听一遍这些话本身，好叫你们看到基督成为\Quote{割礼的仆役，为天主的真理，以证实对祖先的恩许}是什么意思。那么，所说的是什么？曾有一个恩许许诺给亚巴郎，说：\BibleQuote{我要将这块土地赐给你和你的后裔，并且万民都要因你的后裔蒙受祝福。}\BibleRef{创 12:7}但此后，亚巴郎的所有后裔都陷入了惩罚。因为法律因被违反而给他们带来义怒，并从此剥夺了他们那对祖先的恩许。因此，圣子来了，与圣父一同工作，好使那些恩许成为现实并结出果实。因为祂成全了整个法律（在其中祂也成全了割礼），并借着法律和十字架，使他们从违命的诅咒中解放出来，祂没有让这恩许落空。所以当祂称祂为\Quote{割礼的仆役}时，意思是：祂来了，成全了法律，受了割礼，生于亚巴郎的后裔，从而解除了诅咒，平息了天主的义怒，并且使那些将要领受恩许的人适合领受，因为他们已从疏离中被一次而永远地释放。为阻止这些被谴责的人说：\Quote{那么基督怎么受了割礼并遵守了整个法律？}祂将他们的论证转向相反的结论。因为祂不是为使法律继续，而是为终结它，将你们从法律的诅咒中解救出来，并使你们完全摆脱那法律的统治。正是由于你们违犯了法律，祂才成全了它，不是为叫你们成全它，而是为向你们证实对祖先的恩许，这恩许因法律而悬置，因为法律显出你们有过失，不配承受产业。因此，你们也是因恩宠得救，既然你们已被抛弃。那么，你们不要争吵，也不要在不合时宜的时候固执地抓住法律，因为除非基督为你们受了这么多苦，法律也会将你们逐出恩许。祂受苦，不是因为你们配得救恩，而是为叫天主是真实的。随后，为使外邦人不因此骄傲，他说：\par

第9节。\BibleQuote{并使外邦人因天主的仁慈而光荣天主。}\BibleRef{罗 15:9}\par

他的意思是这样的：犹太人即使不配，也会有恩许。但你们连这个也没有，而是仅因对人的爱而得救，即使退一步说，他们也不能从恩许中受益，除非基督来了。然而，为\OriginalTerm{调和}{\GreekText{κεράσῃ}}他们，不让他们起来反对软弱者，他提到了恩许。但关于这些人，他说他们只是因仁慈而得救。因此，他们更应该光荣天主。天主的光荣在于他们彼此融合、团结、一心赞美、承担软弱者、不忽略折断的肢体。然后他附加了证词，显示犹太人应当与外邦人融合；因此他说：\BibleQuote{正如经上所载：因此，我将在外邦人中向祢承认，主，我要歌颂祢的名。}\BibleRef{咏 17:46}\par

第10-12节。\BibleRef{罗 15:10} \BibleQuote{外邦人！你们要会同祂的子民一同喜乐！赞美主，所有外邦人！}\BibleRef{申 32:43}；\BibleQuote{愿万民都赞美祂！}\BibleRef{咏 116:1} \BibleQuote{必有一根从叶瑟的根茎发出，那起来统治外邦人的，外邦人将寄望于祂。}\BibleRef{依 11:1}\par

他给出所有这些引证，为显示我们应当合一，并光荣天主；同时也为谦卑犹太人，使他不在这些人面前自高，因为所有先知都召叫了这些人，也为劝服外邦人谦卑，向他显示他须为更大的恩宠负责。然后他再次以祈祷结束他的论证。\par

第13节。\BibleQuote{愿赐望德的天主，因着信德，以各种喜乐与平安充满你们，使你们因圣神的德能，富于望德。}\BibleRef{罗 15:13}\par

那就是，使你們擺脫那種彼此間的冷淡（\OriginalTerm{冷淡}{\GreekText{ἀθυμίας}}），並且永不因誘惑而沮喪。而這將要藉著你們在望德中充盈而實現。如今，這是一切美善之事的緣由，且出自聖神。但這並非單單來自聖神，而是以我們也貢獻自己的一份為條件。故此，他說：\Quote{因相信}。因為這便是你們充滿喜樂的途徑——若你們相信，若你們盼望。然而，他並未說\Quote{若你們盼望}，而是說\Quote{若你們在望德中充盈}，好叫你們不僅在患難中得到安慰，甚至因信德和望德的豐盈而滿懷喜樂。如此，你們也將吸引聖神來到你們中間。這樣，當祂來到，你們將藉著一切美善之事持守不渝。因為正如食物維持我們的生命，並藉此管理身體，同樣，若我們有善工，我們將擁有聖神；若我們擁有聖神，我們也將有善工。反之，若我們沒有善工，聖神便飛離。但若我們被聖神遺棄，我們在善工上也必跛行。因為當聖神離去，不潔者便來臨：這從\Person{撒烏耳}的事蹟可以明瞭。因為即使他不像對待\Person{撒烏耳}那樣使我們窒息，他仍藉邪惡的作為以另一種方式絞殺我們。因此，我們需要\Person{達味}的豎琴，好能用神聖的詩歌——既包含這些，也包含來自善行的詩歌——來撫慰我們的靈魂。因為若我們只做其中一樣，當我們聆聽撫慰之音時，卻藉行為與那撫慰者爭戰，如同他昔日所為（\BibleRef{撒上 19:10}），那麼這藥方甚至將成為我們的審判，而瘋狂將變得更為猛烈。因為在我們聆聽之前，邪魔害怕我們會聽見而痊癒。但當我們聽了之後，仍然與從前一樣，這恰恰消除了他的恐懼。因此，讓我們歌唱善行的聖詠，好能驅逐比邪魔更惡劣的罪。因為邪魔固然不會使我們喪失天國，在某些情況下甚至幫助清醒的人；但罪卻必然將我們逐出。因為這罪是我們情願接受的邪魔，一種自選的瘋狂。因此，無人憐憫它或寬恕它。那麼，讓我們對如此處境的靈魂吟唱撫慰之歌，既從其他聖經，也從蒙福的\Person{達味}。讓口舌歌唱，讓心靈受教。這已非小事。因為一旦我們教導舌頭歌唱，靈魂便羞於籌劃與所唱相反的事。我們不僅將獲得這益處，還將認識許多與我們相關的事。因為他向你們論及現在與將來、可見與不可見的受造界。若你願學習關於蒼天：它是恆常不變還是將被改變，他給你明確的答覆，說道：\BibleQuote{諸天要如衣服一樣陳舊，上主，你要如衣裳一般將它們捲起，它們就被更換。}（\BibleRef{詠 101:26}）若你願再聽它們的形狀，你將聽見：\BibleQuote{鋪張蒼天如同帳幔}（\OriginalTerm{帳幔}{\GreekText{δέρριν}}）（\BibleRef{詠 104:2}）。若你更想知道它們的背面，他將再次告訴你：\BibleQuote{祂以水遮蓋自己的樓閣。}（\BibleRef{詠 104:2}）而即使在此，他也不停歇，同樣向你論述寬與高，並向你展示它們度量相等。因為，他說：\BibleQuote{東離西有多遠，祂使我們的過犯離我們也有多遠。上主對敬畏祂者的慈愛，有如蒼天高於大地。}（\BibleRef{詠 103:11-12}）但若你願操心大地的根基，這點他也不向你隱藏，你將聽見他歌唱並說道：\BibleQuote{祂將大地奠基於海洋。}（\BibleRef{詠 23:2}）若你渴望知道地震的由來，他也要解除你此一困惑，說道：\BibleQuote{祂注視大地，大地就震動。}（\BibleRef{詠 103:32}）若你探究黑夜的用途，這點你也可以從他學習並知道。因為\BibleQuote{其中林中的百獸徜徉。}（\BibleRef{詠 103:20}）群山有何用途，他會告訴你：\BibleQuote{高山是為野羊。}（\BibleRef{詠 103:18}）為何有岩石？\BibleQuote{岩石是為石獾的避難所。}（\BibleRef{詠 103:18}）為何有不結果實的樹木？從他學習：因為\BibleQuote{那裡雀鳥築巢。}（\BibleRef{詠 103:17}）為何曠野中有泉源？聽！\BibleQuote{它們使天上的飛鳥棲息，使野獸存活。}（\BibleRef{詠 103:12}）為何有酒？並非為使你只飲（因為水足以勝任此需要），而是為使你喜樂：\BibleQuote{美酒使人心情愉快。}（\BibleRef{詠 103:15}）而知道這點，你將明白飲酒之容許限度。飛鳥和野獸從何得食？你將從他的話語聽到：\BibleQuote{這一切都在等待祢，按時給牠們食物。}（\BibleRef{詠 103:27}）若你問：牲畜有何用？他回答你，這些也是為你：\BibleQuote{祂使青草生長為牲畜，使蔬菜生長為人的服役。}（\BibleRef{詠 103:14}）月亮有何用？聽他說：\BibleQuote{祂造月亮為定節令。}（\BibleRef{詠 104:19}）一切可見與不可見之物都是被造的，這點他也清楚地教導我們，說：\BibleQuote{祂命令，萬有便造成；祂發言，萬有便創造。}（\BibleRef{詠 32:9}）死亡有終結，這點他也教導，當他說：\BibleQuote{天主必救我靈脫離陰府的手，當祂收納我時。}（\BibleRef{詠 43:15}）我們的身體從何造成？他也告訴我們：\BibleQuote{祂記念我們不過是塵土。}（\BibleRef{詠 102:14}）又，往何處去？\BibleQuote{歸回塵土。}（\BibleRef{詠 103:29}）這宇宙為何被造？是為你：\BibleQuote{祢以尊榮威嚴為冠冕，使他統治祢手的造化。}（\BibleRef{詠 8:5-6}）我們人類與天使有無相通？這點他也告訴我們，說：\BibleQuote{祢使他略遜於天使。}（\BibleRef{詠 8:6}）論天主的愛：\BibleQuote{如同父親憐憫自己的兒女，上主也憐憫敬畏祂的人。}（\BibleRef{詠 102:13}）論我們今生之後將臨之事及那安寧之境，他教導：\BibleQuote{我的靈魂，歸於安息。}（\BibleRef{詠 114:7}）為何蒼天如此廣大？他也要說。因為\BibleQuote{諸天宣講天主的榮耀。}（\BibleRef{詠 18:1}）為何造了晝夜——不僅為使他們光照和我們休息，也為教導我們：\BibleQuote{沒有言語，沒有詞句，也聽不到它們（即晝夜）的聲音。}（\BibleRef{詠 18:3}）海洋如何環繞大地，這點你也從此學得：\Quote{深淵如衣袍裹住大地。}因為希伯來文如此。\par

但藉着我所提到的这个例子，你将对其余的一切有所了解：关于基督、关于复活、关于来世、关于安息、关于惩罚、关于道德事务，一切涉及教义之事，你会发现这卷书充满了无数的祝福。如果你陷入诱惑，你将从此获得极大的安慰。如果你甚至陷入罪恶，你将发现这里储藏着无数的良药；或者如果陷入贫困或磨难，你将看到许多避风港。如果你正义，你将从此获得许多安全感；如果你是个罪人，也将得到许多慰藉。因为如果你正义却受虐待，你将听到他说：\BibleQuote{因祢的缘故，我们终日被杀，被视为待宰的群羊。} \BibleRef{咏43:22} \BibleQuote{这一切都临到了我们，但我们却没有忘记祢。} \BibleRef{咏43:17} 如果你的善行使你自高，你将听到他说：\BibleQuote{不要与祢的仆人进入审判，因为在祢面前凡活着的人没有一个能成义。} \BibleRef{咏142:2} 你就会立刻谦卑下来。如果你是个罪人，并对自己绝望了，你会不断听到他唱道：\BibleQuote{今天，如果你们听从祂的声音，就不要硬着心，像在激怒之时那样} \BibleRef{咏95:7-8}，你会迅速得到扶持。如果你头上甚至戴着冠冕，并且心高气傲，你将学到：\BibleQuote{君王不是因大军而得救，巨人也不是因力大而获救} \BibleRef{咏32:16}：你会发现自已能变得理性。如果你富有且有名望，你又会听到他唱道：\BibleQuote{祸哉，那些倚靠自己力量，并以财富众多而夸耀的人} \BibleRef{咏48:6} 以及\BibleQuote{至于人，他的岁月如草} \BibleRef{咏102:15}，\BibleQuote{他的荣耀不随他同去，也不在他以后} \BibleRef{咏48:17}：你就不会认为世上任何事物是伟大的。因为当比一切更光辉的事物，即荣耀和能力，都如此无价值时，世上还有什么值得重视呢？你是否沮丧？听他说：\BibleQuote{我的灵魂，你为何如此忧伤？为何如此烦乱？请仰望天主，因为我还要称谢祂。} \BibleRef{咏41:5} 或者你看到不配的人得荣耀？\BibleQuote{不要因作恶的人心怀不平。因为他们如草快要枯干，如青绿的嫩草很快就要凋谢。} \BibleRef{咏37:1-2} 你看到义人和罪人都受惩罚吗？要明白原因不同。因为他说：\BibleQuote{罪恶的祸患很多。} \BibleRef{咏31:10} 但论到义人，他不说祸患，而是说：\BibleQuote{义人的苦难很多，但主拯救他们脱离这一切。} \BibleRef{咏33:19} 又说：\BibleQuote{恶人的死亡是悲惨的。} \BibleRef{咏33:21} \BibleQuote{主圣人的逝世，在祂眼中是宝贵的。} \BibleRef{咏115:15} 你要不断思想这些事，并由此受教。因为这里的每一句话都包含着无法测度的意义之海。我们只是浅尝辄止，但如果你愿意对这些章节进行精确的研究，你会发现其中的丰富是巨大的。但就目前而言，即使通过我所给出的这些，也能够清除你身上的情欲。因为既然他禁止我们嫉妒、忧伤、不合时宜地沮丧，或认为财富、磨难、贫穷有什么了不起，或幻想生命本身有什么了不起，他就使你脱离了一切情欲。因此让我们为此感谢天主，并让我们手中常握这宝藏，\BibleQuote{藉着圣经的忍耐和安慰，我们怀有望德} \BibleRef{罗15:4}，并享受将来的福乐。愿天主赐我们众人都能获得，因我们的主耶稣基督的恩宠和慈爱。借着祂并同祂，等等。\par

\SourceRecord{Translated by J. Walker, J. Sheppard and H. Browne, and revised by George B. Stevens.；Nicene and Post-Nicene Fathers, First Series；Vol. 11.；Edited by Philip Schaff.；Buffalo, NY: Christian Literature Publishing Co.,；1889.；<http://www.newadvent.org/fathers/210228.htm>.}

% structure: p0001 source=h1 target=section reason=page-title

\section{论\WorkTitle{罗马人书}第二十九篇讲道}

\BibleRef{罗 15:14}\par

\BibleQuote{我本人也对你们怀有信心，我的弟兄们，深知你们也充满良善，满备一切知识，也能彼此劝诫。}（多数文本如此；\Person{圣金口若望}作\InnerQuote{他人}。）\par

他曾说过：\BibleQuote{我既是外邦人的宗徒，我光荣我的职务。}\BibleRef{罗 11:13} 他曾说过：\BibleQuote{小心，恐怕他也不饶恕你。}\BibleRef{罗 11:21} 他曾说过：\BibleQuote{不要自以为聪明}\BibleRef{罗 12:16}；又说：\BibleQuote{你为什么判断你的弟兄呢？}\BibleRef{罗 14:10} 以及：\BibleQuote{你是谁，竟判断别人的仆人？}\BibleRef{罗 14:4} 此外还有许多类似的话。既然他常常使自己的言辞有些严厉，现在他就以温和的话语来\OriginalTerm{安慰}{\GreekText{θεραπεύει}}他们。他在开头所作的，在结尾也作了。开头他说：\BibleQuote{我感谢我的天主，因你们的信德传遍了全世界。}\BibleRef{罗 1:8} 但在这里他说：\BibleQuote{我相信你们也充满良善，也能劝诫别人；}这比前者更进一层。他并没有说：我听说了，而是说：\Quote{我相信}，无需从别人那里听说。\Quote{我本人}，就是说，我这个责备你们、控告你们的人。\Quote{你们充满良善}，这是针对刚才的劝勉。好像他说：我给你们那个劝勉，要你们接纳、不忽视、不摧毁\Quote{天主的工程}，并不是因为你们残忍或仇恨弟兄，因为我知道\Quote{你们充满良善}。但在我看来，他在这里称他们的德行是完美的。他没有说你们有，而是\Quote{你们充满}。下文也用了同样的强调：\Quote{满备一切知识}。因为他们可能虽然富于爱心，却不知道如何恰当地对待所爱的人。这就是为何他加上\Quote{一切知识}。\Quote{也能劝诫别人}，不仅是学习，而且也能教导。\par

\BibleRef{罗 15:15}  \BibleQuote{然而，我更为大胆地写信给你们，在某种程度上。}\par

请看\Person{保禄}的谦卑，请看他的智慧：他在前一部分给了深深的切割，然后当他成功地实现了愿望时，他又如何大施温柔。因为即使没有他所说的话，仅仅承认自己曾大胆，就足以消解他们的激动。他在写给\BibleBook{希伯来}{人书}的信中也这样做，如此说：\BibleQuote{可是，亲爱的诸位，我们虽这样讲，但对你们，我们确信你们有更美的事，且关乎救恩。}\BibleRef{希 6:9} 在写给\BibleBook{格林多}{人书}的信中又说：\BibleQuote{我称赞你们，弟兄们，因为你们在一切事上记念我，并遵守我所传授给你们的规传。}\BibleRef{格前 11:2} 在写给\BibleBook{迦拉达}{人书}的信中他说：\BibleQuote{我在主内深信你们不会存别样的心。}\BibleRef{迦 5:10} 在书信的各部分，人们可以发现这一点常被遵守。但在这里甚至更甚。因为他们处于更高的地位，需要既用收敛剂也用宣泻剂来降低他们的高傲心态。因为他以不同的方式这样做。所以在这里他也说：\BibleQuote{我更为大胆地写信给你们}，即使如此他仍不满足，又加上：\Quote{在某种程度上}，即温和地；他在这里还不停止，而是说什么呢？\Quote{为提醒你们。}他没有说教导，也没有说简单的提醒（\OriginalTerm{提醒}{\GreekText{ἀναμιμνήσκων}}），而是用了一个词（\OriginalTerm{温婉提醒}{\GreekText{ἐπαναμιμνήσκων}}），意思是温和地提醒你们。请看结束与开头一致。因为正如在那段他说：\BibleQuote{你们的信德传遍了全世界}\BibleRef{罗 1:8}，同样在书信的结尾也说：\BibleQuote{你们的服从已传到各处}\BibleRef{罗 16:19}。又如在开头他写道：\BibleQuote{我很想见你们，要把一些属神的恩赐分给你们，为使你们得以坚固；也就是说，我愿与你们同得安慰}\BibleRef{罗 1:11-12}，所以在这里他也说：\Quote{为提醒你们。}无论在那里还是在这里，他都从师傅的座位上下来，以同等的弟兄和朋友的身份对他们说话。这更是师傅的本分，使他的讲论富有变化，对听者有益。请看，在他说了\BibleQuote{我更为大胆地写信给你们}、\Quote{在某种程度上}、\Quote{为提醒你们}之后，他仍不满足于这些，而是使自己的言辞更加谦卑，继续说：\par

\BibleQuote{借着天主赐给我的恩宠。}正如他在开头所说：\BibleQuote{我负有债务。}\BibleRef{罗 1:14} 好像他在说：这尊荣不是我为自己夺取的，也不是我先跳上前去的，而是天主命令了这事，而且这也是按恩宠，好像他没有因我配得而分别我为这职务。所以你们不要恼怒，因为不是我自己高抬自己，而是天主吩咐的。又如在那边他说：\BibleQuote{我在他圣子的福音上，以心灵所事奉的天主}\BibleRef{罗 1:9}，同样在这里，在说了\Quote{借着天主赐给我的恩宠}之后，他又加上：\par

\BibleRef{罗 15:16}  \BibleQuote{使我为外邦人作耶稣基督的仆役，司祭般地\OriginalTerm{服务}{\GreekText{ἱερουργοὕντα}}天主的福音。}\par

因为他充分地证明了自己的论点后，便将论调提升至更高层次，不再像开头那样仅仅谈论服务，而是谈论\OriginalTerm{服务与司祭职务}{\GreekText{λειτουργίαν} \GreekText{καί} \GreekText{ὶερουργίαν}}。因为对我来说，这宣讲与宣告就是一种司祭职务。这是我所献的祭。没有人会责备一位司铎，因为他渴望献上无瑕的祭品。他说这话是为了\OriginalTerm{提升}{\GreekText{πτερὥν}}他们的思想，向他们表明他们就是祭品，同时也是为自己在这件事上的职责辩护，因为他被委派担任这一职务。他说：我的刀就是福音，就是宣讲之言。原因不在于使我得荣耀，不在于使我显得显赫，而在于使那\Quote{外邦人的\OriginalTerm{奉献}{\GreekText{προσφορὰ}}当成为可悦纳的，藉着圣神得以圣化。}\par

也就是说，那些由我教导的灵魂可以蒙悦纳。因为天主引领我达到这一高度，并非为尊荣我，而是出于对你们的关心。那么他们如何成为可悦纳的呢？是藉着圣神。因为不仅需要信德，还需要属灵的生活方式，使我们能保持那一次而永远赐下的圣神。因为不是木柴和火，也不是祭台和刀，而是圣神在我们内成为一切。因此，我竭尽所能防止那火熄灭，正如我已受命去做的。那么你为何对那些不需要的人说话呢？他回答说：这正是我不教导你们，而是提醒你们的原因。如同司铎站在一旁扇旺火焰，我也如此，激发你们的准备心态。请注意，他没有说\Quote{你们的奉献当成为……}而是说\Quote{外邦人的奉献}。但当他提到外邦人时，他意指整个世界、大地和整个海洋，为抑制他们的高傲，使他们不会轻视这位\OriginalTerm{伸展}{\GreekText{τεινόμενον}}到地极的教师。正如他在开头所说的：\BibleQuote{如同在其他外邦人中一样，我也是希腊人和化外人、聪明人和愚拙人的负债者。} \BibleRef{罗 1:13}\par

17. \BibleQuote{所以，我藉着耶稣基督，在关乎天主的事上，有可夸耀的。} \BibleRef{罗 15:17}\par

由于他极其谦卑了自己，他再次提升了他的风格，这样做也是为了他们，免得他变得轻易被轻视。当他提升自己时，他记起了自己的适当气质，说道：\BibleQuote{所以我有所可夸耀的。} 他意指，我不是在自己内夸耀，不是在我们自己的热忱内，而是在天主的恩宠内。\par

18. \BibleQuote{我不敢提说，基督藉着我以外的事，有什么没有藉着言语和作为，藉着神迹和奇事，并天主圣神的能力，使外邦人顺服。} \BibleRef{罗 15:18}\par

他的意思是，没有人能说我的话只是夸口。因为我的这司铎职务，我所拥有的标记，以及任命的证据，比比皆是。不是长衣（\OriginalTerm{长衣}{\GreekText{ποδήρης}}）和铃铛，如同古人那样，也不是冠冕和头巾（\OriginalTerm{头巾}{\GreekText{κίδαρις}}），而是比这些更可敬畏的神迹和奇事。也不能说，我固然被托付了这职责，却没有执行。更确切地说，不是我执行的，而是基督。因此，我也是在祂内夸口，不是关于寻常的事，而是关于属神的事。这就是\Quote{关于天主的事}的意思。因为我已完成了我被派遣的目的，我的话不是夸口，神迹和外邦人的服从便证明了这一点。\BibleQuote{除了基督藉我作的那些事，我什么都不敢提，只提祂藉我的言语和作为，藉着神迹和奇事，并藉着天主圣神的能力，使外邦人服从。}请看，他是多么极力要证明一切都是天主的作为，没有什么是自己的。因为无论我说什么，或做什么，或行神迹，都是祂做的，都是圣神做的。他说这话也是为了显示圣神的尊威。请看，这些事比古时的事更奇妙、更可敬畏，即祭献、奉献和象征。因为当他说\Quote{藉着言语和作为，藉着大能的神迹和奇事}时，他指的是这个：教训、关于天国的体系（\OriginalTerm{体系}{\GreekText{φιλοσοφίαν}}）、行为和生活方式的表现、死人复活、魔鬼被驱逐、瞎子复明、瘸子跳跃，以及其他奇事，这一切都是圣神在我们内做的。然后，这些事的证明（因为这一切还只是断言）就是门徒的众多。因此他接着说：\BibleQuote{所以从耶路撒冷及其周围直到依里黎苛，我已把基督的福音传遍了。}那么，请数算城市、地方、民族、人民，不仅是罗马统治下的，还有蛮族统治下的。因为我不希望你只走遍腓尼基、叙利亚、基里基雅和卡帕多细雅，还要算上那些背后的地区，即撒辣森人、波斯人、亚美尼亚人以及其他野蛮民族的地方。这就是为什么他说\Quote{周围}，好叫你不仅走直路，还要在思想中跑遍整个地区，甚至亚洲的南部。如同他用一句话\Quote{藉着大能的神迹和奇事}密集地列举了神迹，他也用\Quote{周围}一词概括了无数的城市、民族、人民和地方。因为他远远脱离了所有的夸口。他说这话是为了他们，免得他们自高自大。他在开头说\Quote{我也可以在你们中间得些果实，如同在其他外邦人中一样。}但这里他说明了他司铎职位的必要性。因为他用更尖锐的语气说话，也借此更清楚地显示了他的权能。这就是为什么他在那里只说\Quote{如同在其他外邦人中一样}，但在这里他充分坚持这个话题，好从各方面削除骄傲。他不仅说传了福音，还说\BibleQuote{把基督的福音传遍了。}\par

\BibleRef{罗 15:20} \BibleQuote{然而，我这样竭力传福音，不是在基督被称名的地方。}\par

请看，这里又一个卓越之处：他不仅向这么多人传了福音并说服了他们，他甚至不去那些已成门徒的人那里。他如此远离将自身强加于他人门徒，或为荣耀而做这事，以至于他特意去教导那些没有听过的人。因为他不是说\Quote{在那些没有被说服的地方}，而是说\Quote{在基督连名都没有被提过的地方}，这更显著。他这样抱负的原因是什么？他说：\Quote{免得我建造在别人的根基上。}\par

他说这话是为了表明自己远离虚荣，并教导他们，他写信并非出于贪图荣耀或从他们得荣誉，而是履行他的职务，完成他的司铎职责，关爱他们的得救。但他称宗徒的根基为\Quote{别人的}，不是就人的品质或宣讲的性质而言，而是就赏报的问题而言。因为宣讲并不是别人的，但就归于别人的赏报而言，这赏报对这人来说是别人的。然后他显示一个预言也应验了，说：\par

\BibleRef{罗 15:21} \BibleQuote{正如经上所载：\InnerQuote{未曾听见过祂的，将要看见；未曾听过的，将要明白。}} \BibleRef{依 3:15}\par

你看，他跑向劳动更多、辛劳更大的地方。\par

\BibleRef{罗 15:22} \BibleQuote{因此，我也多次被拦阻，不能到你们那里去。}\par

\BibleQuote{我屡次定意要到你们那里去，但直到如今仍被拦阻。} \BibleRef{罗 1:13} 但这里他给出了被拦阻的原因，而且不是一次，而是两次，甚至多次。因为正如他在那里说\BibleQuote{我屡次定意要到你们那里去}，这里也说\BibleQuote{我多次（或常常，\OriginalTerm{多次}{\GreekText{τὰ} \GreekText{πολλὰ}}）被拦阻}到你们那里去。这证明了一种非常强烈的愿望，因为他尝试了那么多次。\par

\BibleRef{罗 15:23} \BibleQuote{但如今在这些地方再没有可传的地方了。}\par

请看，他如何表明他写信和将要来，都不是出于贪图从他们得荣耀。\BibleQuote{而且多年来我切切想望到你们那里去，}\par

\BibleRef{罗 15:24} \BibleQuote{当我往西班牙去的时候，我希望在途中经过你们那里，并得你们送我往那里去，只要先让我在你们那里稍得满足。}\par

因为他不想显得过于轻视他们，如果说\Quote{我没有什么事要做，因此我要到你们那里去}，他就再次触及爱的主题，说道：\Quote{我多年来有一个强烈的愿望，要到你们那里去。}我渴望前来的原因，并非因为我空闲，而是为了孕育那长久以来使我痛苦的心愿。然而，为了不让他们因此自高，请看他是如何贬低他们，说道：\Quote{每当我要前往西班牙时，希望途中能见到你们。}他说这话正是为了不让他们心高气傲。他想要表达爱意，同时又防止他们变得挑剔。因此他将这话紧接在前言之后，交替使用以确认二者。也为了不让他们说：\Quote{你只是顺路看我们吗？}他补充说，并且蒙你们送我上路（即：你们可以为证，我经过你们那里并非轻视你们，而是出于不得已。）然而，这仍然令人痛苦，他更细心地治愈它，说道：\Quote{我先稍得与你们相伴满足。}因为他说\Quote{在旅途中}，表明他并非贪求他们的好感。但他说\Quote{满足}，表明他渴望他们的爱，不仅是渴望，而且极其渴望；因此他没有说\Quote{满足}，而是说\Quote{稍得}。也就是说，任何时间长度都不能使我满足，或使我对你们的相伴产生厌倦。请看，他如何在匆忙中仍不离去，直到满足，以此表达他的爱。这是他深情的标志，他以如此温暖的方式使用言辞。因为他甚至不说\Quote{我将会见到}，而是说\Quote{将会满足}，效法父母的语言。起初他说：\BibleQuote{好叫我获得一些果实。}\BibleRef{罗 1:13}但这里说\Quote{满足}。这两者都像是一个吸引他人的人。因为前者是极大地称赞他们，如果他们的顺服能为他结出果实；而后者则是他自己友谊的真实证据。在写给格林多人的书信中，他如此说：\BibleQuote{你们可以送我上路，无论我去哪里。}\BibleRef{格前 16:6}这样，他在各方面向他的门徒表现出无与伦比的爱。因此，在他所有书信的开头，他从这一点出发，结尾也以此收束。正如溺爱的父亲对待独生且亲生的儿子，他同样爱所有信众。因此他说：\BibleQuote{谁软弱，我不软弱呢？谁跌倒，我不着急呢？}\BibleRef{格后 11:29}\par

因为教师首先应当具备这一点。因此基督也对伯多禄说：\BibleQuote{你若爱我，就喂养我的羊。}\BibleRef{若 21:16} 因为爱基督的人也爱他的羊群。梅瑟也是如此，当他显露出对犹太百姓的善意时，天主便立他管理他们。达味也是如此成为君王，因为他先被看出对他们怀有深情；他虽年轻，却如此为百姓忧伤，以致当杀死那蛮人时，甘愿为他们冒生命危险。但他若说\BibleQuote{为那杀死这培肋舍特人的人，要怎样对待？}\BibleRef{撒上 19:5}，他这样说并非为了要求赏报，而是希望别人信赖他，并将与他作战的使命交托给他。因此，当他得胜后来到君王面前，对这些事只字未提。撒慕尔也非常慈爱；因此他说：\BibleQuote{总之，我若停止为你们向主祈祷，愿天主阻止我犯罪。}\BibleRef{撒上 12:23} 保禄也是如此，确切地说，不是如此，而是以更大的程度，向他的所有\OriginalTerm{属民}{\GreekText{τὥν} \GreekText{ἀρχομένων}}燃烧着爱心。因此，他使门徒们对他怀有这样深厚的感情，以致他说：\BibleQuote{如果可能，你们会把眼睛挖出来给我。}\BibleRef{迦 4:15} 也是因此，天主尤其以此责备犹太人的教师，说：\BibleQuote{以色列的牧人哪，牧人岂是喂养自己？难道不是喂养羊群吗？}\BibleRef{厄 34:2} 但他们却反其道而行。因为他说：\Quote{你们喝牛奶，穿羊毛，宰杀肥羊，却不喂养羊群。} 基督在提出最合宜的牧者准则时说：\BibleQuote{善牧为羊舍命。}\BibleRef{若 10:11} 达味也是如此，不仅在多次其他场合，而且当那从上而来的可怕义怒降于全民时。因为当所有人都被击杀时，他说：\BibleQuote{我，牧人，犯了罪；我，牧人，行了恶；而这些羊群，他们做了什么？}\BibleRef{撒下 24:17} 同样，在选择那些惩罚时，他没有选择饥荒或逃避敌人，而是选择了天主降下的瘟疫，借此他希望使所有其他人都安全，而他自己比其余所有人更先被带走。但既然并非如此，他哀叹并说：\Quote{愿你的手落在我身上}；若这还不够，也\Quote{落在我父家}。他说：\Quote{因为我，}他说，\Quote{牧人，犯了罪。} 仿佛他曾说，即使他们也犯了罪，我才是应受惩罚的人，因为我没有纠正他们。但因为罪也是我的，我理应承受惩罚。为了加重罪行，他用了\Quote{牧人}这个名称。于是就这样平息了义怒，就这样使判决撤销！认罪的力量何等伟大！\Quote{因为义人先控告自己。} 一位善牧的关怀与同情何等巨大！因为他的肠子因他们的跌倒而绞痛，如同自己的儿女被杀。因此他恳求义怒降临到自己身上。在屠杀之初，他本已这样做，除非他看见它正在前进并预料它会临到自己。因此，当他看到这并未发生，而灾祸在他们中间肆虐时，他不再忍耐，而是比为他长子阿默农更悲伤。因为那时他并未求死，但现在他恳求自己先于他人倒下。统治者应当如此，当为别人的灾难而非自己的灾难悲伤。他在他儿子的情形上同样遭受了某种类似的事，使你们看到他爱儿子不如爱臣民，而那青年人不洁，且是父亲的\OriginalTerm{弑父者}{\GreekText{πατραλοίας}}，但他仍然说：\BibleQuote{我恨不得替你死！}\BibleRef{撒下 18:33} 你说什么，真福者，万民中最温良的人？你的儿子决意杀你，用无数灾祸围困你。当他被除掉，战利品举起时，你反倒祈求被杀？是的，他说，因为军队不是为我获胜，而是我比从前更猛烈地受攻击，我的肠子现在比从前更撕裂。然而，这些都为他们所负责的人着想，但有福的亚巴郎甚至为那些并未托付给他的人也极为关心，以至于投身于惊险之中。因为当他行事时，不仅为了他的侄子，也为了索多玛百姓，他没有停止追赶那些波斯人，直到全部释放了他们；然而他本可以在救出他之后离去，但他没有选择那样做。因为他同样关心所有人，并在他后来的行为中也显示了这一点。当时并非一群蛮族即将包围他们，而是天主的义怒要将他们的城市连根拔起，不再是武器、战斗和阵列的时候，而是祈求的时候；他为他们显示的热诚如此之大，仿佛他自己即将灭亡。因此，他一次、两次、三次，甚至多次来到天主面前，并以他的本性找到避难所（即借口），说：\BibleQuote{我是灰尘和灰烬。}\BibleRef{创 18:27}；既然他看到他们背叛自己，他便恳求他们因别人而获救。因此天主也说：\BibleQuote{我岂能瞒着亚巴郎我的仆人我要做的事呢？}\BibleRef{创 18:17}，好让我们学习义人是何等爱人。若不是天主先停止，他也不会停止恳求。\BibleRef{创 18:33} 他看似在为义人祈祷，但却是为所有人做一切。因为圣人们的灵魂非常温柔且爱人，既关心自己的人，也关心陌生人。甚至对无理性的受造物，他们也延伸他们的温柔。因此有一位智者说：\Quote{义人顾惜他牲畜的性命。} 但如果他这样对待牲畜，更何况对待人呢？但既然我提到了牲畜，让我们只考虑在卡帕多细雅地方的牧羊人，他们在守护无理性的受造物时所受的苦难的种类和程度。他们常常一连三天埋在雪下。而在利比亚的那些人据说也经历不少于这些的艰苦，他们在那个荒芜的旷野中整月徘徊，旷野充满最凶猛的\OriginalTerm{野兽}{\GreekText{θηρία}}（可能包括蛇）。如果为无理性之物有如此大的热诚，我们这些受托管理有理性灵魂的人，却仍在沉睡中打盹，还能提出什么辩护呢？因为安息、静处、不四处奔跑、不为这些羊献上无数次死亡，难道是正当的吗？或者你们不知道这羊群的高贵吗？你们的主不是为此承受了无尽的痛苦，然后倾流了祂的血吗？而你们却寻求安息？还有什么比这些牧人更坏的呢？你们没有觉察到，环绕这些羊的狼比这世界的狼更凶猛残暴吗？你们难道不自我思量，担负这职务的人应当有何等灵魂？如今，领导民众的人，即使只商议普通事务，也夜以继日地警醒。而我们这些为天国奋斗的人，竟然白天也睡觉。那么，谁将救我们脱离这些事的惩罚呢？因为即使身体被肢解，即使遭受万死，难道人不应该像赴宴一样奔赴吗？但愿不仅牧人，而且羊群也听到这话；好使他们使牧人更加积极，更加鼓励他们的善意：我并非指别的，而是指以完全服从和顺从来。保禄也这样命令他们，说：\BibleQuote{要服从你们的领袖，并服从他们；因为他们为你们的灵魂警醒，如同要交账的人。}\BibleRef{希 13:17} 当他说\Quote{警醒}时，他指的是成千上万的辛劳、忧虑和危险。因为那善牧，即基督所期望的那种牧人，在无数见证人面前搏斗。因为基督为他死了一次；但这人却为羊群死一万次，如果，就是说，他是他应当成为的那种牧人；因为这样的人可以天天死。\BibleRef{罗 8:36} 因此，你们既然知道这劳苦是什么，就与他们合作，以祈祷、热诚、敏捷、爱心，好使我们能以你们为荣，你们也以我们为荣。因为正是因此，祂把这使命托付给宗徒之长，他也爱祂超过其余的人；在首先问他是否爱祂之后，好使你们知道，这比其他事更被视为爱祂的证明。因为这需要一个坚强的灵魂。我这是就最好的牧人说的；不是就我和我们这时代的人，而是就任何可能像保禄、像伯多禄、像梅瑟那样的人。那么，让我们效法这些，无论是我们的统治者还是被统治者。因为被统治者也可以成为牧人，相对于他的家庭、他的朋友、他的仆人、他的妻子、他的儿女：如果我们这样安排我们的事务，我们就会获得一切美善。愿天主恩赐我们都能获得，藉着恩宠和爱人之心，等等。\par

\SourceRecord{Translated by J. Walker, J. Sheppard and H. Browne, and revised by George B. Stevens.；Nicene and Post-Nicene Fathers, First Series；Vol. 11.；Edited by Philip Schaff.；Buffalo, NY: Christian Literature Publishing Co.,；1889.；<http://www.newadvent.org/fathers/210229.htm>.}

% structure: p0001 source=h1 target=section reason=page-title

\section{罗马人书讲道辞第三十篇}

\BibleRef{罗 15:25-27}\par

\BibleQuote{但现在我要往耶路撒冷去，为圣人服务。因为马其顿和阿哈雅人乐意凑出一笔捐项，为耶路撒冷的贫穷圣人。他们确实乐意，而且他们是他们的债务人。}\par

既然他曾说过，我不再在这些地方有\BibleQuote{更多余地}，并且\BibleQuote{这些年我一直切愿到你们那里去}，但他仍有意延迟；免得人们以为他在戏弄他们，他也提到为什么还拖延的原因，并说：\BibleQuote{我现在往耶路撒冷去}，显然是在为延迟找借口。借此，他也成就了另一个目的，就是劝勉他们施舍，并使他们更加认真对待此事。因为如果他无意成就此事，说\BibleQuote{我现在往耶路撒冷去}就已足够。但现在他加上他旅程的原因。他说：\BibleQuote{因为我去，是为圣人服务}。他仔细思考这个话题，进入推理，说他们是\BibleQuote{债务人}，并且说：\BibleQuote{如果外邦人已分享他们的属灵事物，他们也有义务在属血肉的事物上为他们服务}，好让他们学会效法这些。因此，我们很有理由钦佩他的智慧，设计出这种劝告的方式。因为比起以劝诫的形式说，他们更可能以这种方式接受；因为如果为了激励他们而把格林多人和马其顿人摆在他们面前，他那时似乎就是在侮辱他们。事实上，这正是他如下激励他人的依据，说：\BibleQuote{此外，弟兄们，我们告诉你们天主赐给马其顿众教会的恩宠}。\BibleRef{格后 8:1} 他又用这些激励马其顿人。他说：\BibleQuote{因为你们的热情，激发了许多人}。\BibleRef{格后 9:2} 他也同样对迦拉达人这样做，如他说：\BibleQuote{我怎样吩咐了迦拉达众教会，你们也照样做}。\BibleRef{格前 16:1} 但对罗马人，他没有这样做，而是以更隐蔽的方式。他在讲道方面也这样做，如他说：\BibleQuote{什么？天主的话从你们出来吗？还是只临到你们身上？} \BibleRef{格前 14:36} 因为没有什么比争竞更有力量。因此他常使用它。因为在别处他也说：\BibleQuote{我也在众教会中这样规定} \BibleRef{格前 7:17}；又说：\BibleQuote{正如我在各处各教会中所教导的}。\BibleRef{格前 4:17} 他对哥罗森人说：\BibleQuote{福音在世界各处增长结果}。\BibleRef{哥 1:6} 他在这里对施舍也是如此。请看他措辞中何等尊严。因为他不说\Quote{我去带施舍}，而是说\Quote{去服务} \OriginalTerm{服务}{\GreekText{διακονὥν}}。但如果保禄服务，试想这是何等伟大的事，当世界的导师承担起搬运者，并且当他正要起程去罗马，又如此渴望见到他们，却仍更偏好此事。\Quote{因为马其顿和阿哈雅人乐意}，也就是说，符合他们的认可、他们的愿望。\Quote{某种捐献}。他不说施舍，而是说\Quote{共融} \OriginalTerm{共融}{\GreekText{κοινωνίαν}}。\Quote{某种}一词并非无意义，而是为了避免显得在责备他们。他不只说贫穷人，而是\Quote{贫穷的圣人}，这样使他的推荐有双重效果，既基于他们的美德，也基于他们的贫穷。即使如此他还不满足，还加上\Quote{他们是他们的债务人}。然后他表明他们如何是债务人。因为他说：\BibleQuote{如果外邦人已被分享他们的属灵事物，他们的债务也是要在属血肉的事物上为他们服务}。但他意思是：基督是为他们而来的。所有承诺都是给他们的，给犹太人的。基督出自他们。（因此也说：\BibleQuote{救恩出自犹太人。}）\BibleRef{若 4:22} 宗徒出自他们，先知出自他们，一切美善出自他们。因此，世界在这些事物上都成了分享者。如果，他说，你们在更大的事上成了分享者，而当宴席是为他们准备的，你们却被请来享用那铺陈的宴席 \BibleRef{玛 22:9}，按福音的比喻，你们也有债务与他们分享你们的属血肉的事物，并分给他们。但他不说分享，而是说\Quote{服务} \OriginalTerm{服务}{\GreekText{λειτουργἥσαι}}，从而把他们与仆役 \OriginalTerm{仆役}{\GreekText{διακόνων}} 以及向君王纳税的人同等看待。他不说\Quote{在你们属血肉的事物上}，如他说\Quote{在他们的属灵事物上}那样。因为属灵事物是他们的。但属血肉的事物不只属于他们，而是所有人的共同财产。因为他命令钱财应被视为属于所有人，而不仅仅是它的拥有者。\par

第28节。\BibleQuote{因此，当我办成这事，并为他们封印了这果子。} \BibleRef{罗 15:28}\par

即，当我好像把它存放在皇家宝库里，犹如在安全无盗贼和危险的地方。他不说施舍，而是再说\Quote{果子}，以表明施与者从中获利。\BibleQuote{我将从你们那里去西班牙。} 他再次提到西班牙，以显示他对他们的热情和热忱 \OriginalTerm{热忱}{\GreekText{ἀόκνον}}。\par

第29节。\BibleQuote{而且我确信，当我到你们那里去时，必带着基督福音的圆满祝福而来。} \BibleRef{罗 15:29}\par

\Quote{圆满的祝福}有何含义？或指施舍（钱财），或泛指善行。因为他常以\Quote{祝福}称呼施舍，正如他所说：\BibleQuote{作为祝福，而非作为贪婪}\BibleRef{格后 9:5}。这称呼在旧时也是惯例。但既然此处他加上了\Quote{福音的}，我们据此断言他不仅指钱财，也指一切事物。似乎他已说：我知道当我来时，我会发现你们具备一切善行的光荣与新鲜气，在福音中配得无数赞美。这是一种十分有效的劝诫方式，即通过赞颂预先引起他们的注意。因为当他以劝诫方式恳求他们时，这正是他采用的纠正方式。\par

第30节：\BibleQuote{如今我恳求你们，弟兄们，因主耶稣基督之名，并因圣神的爱。}\BibleRef{罗 15:30}\par

此处他再次提出基督与圣神，全然不提圣父。我说这话，为使你们发现他提及父与子、或仅提及父时，不至于轻视子或圣神。他并非说\Quote{圣神}，而是说\Quote{圣神的爱}。因为正如基督爱了世界，圣父也如此，圣神也同样。我们且听：你恳求我们什么？\Quote{与我一同在祈祷中为我在天主前奋斗}\par

第31节：\BibleQuote{使我得以脱离犹太地那些不信的人。}\BibleRef{罗 15:31}\par

一场巨大的挣扎正摆在他面前。这也是他呼吁他们祈祷的原因。他并未说\Quote{使我得以参与}，而是说\Quote{使我得以脱离}，正如基督所命：\BibleQuote{祈祷，免得你们进入诱惑}\BibleRef{玛 26:41}。他在此表明，某些恶狼、野兽般的仇敌将攻击他们。由此他也找到另一个理由，以证明他担当圣徒的职务是合理的，若不信者如此众多，以致他甚至祈祷脱离他们。因为身陷众多敌人中的人也可能因饥饿灭亡。因此，从别处向他（或\Quote{他的行程}）提供援助绝对必要。\Quote{并使我为耶路撒冷所作的服役，能被圣徒所悦纳。}\par

即：我的祭献得以被接纳，使接受者欣然领受所赐之物。看他如何再次抬高受施者的尊严。然后他请求如此伟大民族的祈祷，为使所送之物被接纳。由此他也指出另一点：施舍并不保证被接纳。因为当人勉强施舍、或以不义之财、或为虚荣时，其果实便消失了。\par

第32节：\BibleQuote{使我藉着天主的旨意，满心喜乐地来到你们那里。}\BibleRef{罗 15:32}\par

如他起初所说：\BibleQuote{或许如今，因天主旨意，终得顺利前行，来到你们那里}\BibleRef{罗 1:10}；此处他再次投靠同一旨意，并说这是我奋力前行并希望脱离他们的原因，为能尽快见到你们，且是愉悦地，不带任何来自彼处的沉重负担。\Quote{并与你们同得安息。}\par

看他再次如何显出谦逊。他并未说\Quote{我教导你们，警诫你们}，而是说\Quote{我与你们同得安息}。然而他正是那奋斗争战之人。那么他为何说\Quote{使我与你们同得安息（\OriginalTerm{同得安息}{\GreekText{συναναπαύσωμαι}}）}？这是为在这方面也使你们满足，使你们因分享他的冠冕而更愉快，并表明你们也一同奋斗劳苦。随后，按他惯常的作法，他在劝诫后加上祈祷，说道：\par

第33节：\BibleQuote{愿平安的天主与你们众人同在。阿们。}\BibleRef{罗 15:33}\par

第十六章第1节：\BibleQuote{我把我们的姊妹福依贝托付给你们，她是坚革里教会中的女执事。}\BibleRef{罗 16:1}\par

看他用了多少方式赋予她尊严。他既在众人之前提及她，又称她为姊妹。被称为保禄的姊妹并非小事。此外，他又提及她的职位，指明她是\Quote{女执事}。\par

第2节：\BibleQuote{请你们在主内，以合乎圣徒的样式接待她。}\BibleRef{罗 16:2}\par

即：为主的缘故，使她在你们中间得享尊荣。因为人若为主的缘故接待某人，即使接待的人身份不高，仍会用心接待。但若对方是圣徒，你想该当如何用心？故此他加上\Quote{合乎圣徒的样式}，即这样的人当如何被接待。因她有两重理由使你们关注她：既为主缘故被接待，又她自身是圣徒。并请你们\Quote{在她一切所需（即\InnerQuote{请求}，\OriginalTerm{需要}{\GreekText{χρήζῃ}}）的事上帮助她}。你们看他减轻了他们的负担多么微小。他并未说\Quote{请你们派遣}，而是\Quote{请你们贡献分内之事，向她伸出援手}；且是\Quote{在她一切所需的事上}。并非任何事，而是她向你们请求的事。但她只会在你们能力范围内的事上请求。随后，又是一句极大的赞美：\Quote{因为她曾帮助许多人，也帮助了我本人。}\par

请看他的判断。先是赞美，然后插入劝勉，再给予赞美，如此将这位蒙福女子的赞词置于她的需要两旁。因为那蒙受如此有利见证的祝福，且有能力援助那匡正全世界的人，这女子怎能不蒙福呢？这正是她善举的顶峰，因此他将其放在最后，正如他所说：\Quote{也帮助了我}。但\Quote{也帮助了我}这句话传达了什么？对那位世界宣报者，那位忍受了许多痛苦的人，那位足以援助万人（\OriginalTerm{足以援助万人}{\GreekText{μυρίοις} \GreekText{ἀρκοὕντος}}）的人。那么，我们男女众人都当效法这位圣妇以及随后的一位，连同她的丈夫。他们是谁呢？\par

第二节。他说：\BibleQuote{请问候普黎斯加和阿桂拉，他们在基督耶稣内是我的助手}。\BibleRef{罗 16:2}\par

关于这些人的卓越，\Person{圣路加}也作证。部分是在他说\Person{保禄}\BibleQuote{与他们同住，因为他们操业同为帐棚匠}（\BibleRef{宗 18:3}）；部分是在他指出那妇人接待\Person{阿颇罗}，并教导他天主的道（\BibleRef{宗 18:26}）。这些固然是大事，但保禄所说的更大。他说了什么呢？首先，他称他们为\BibleQuote{助手}，指明他们曾分享他极大的辛劳与危险。然后他说：\par

第四节。\BibleQuote{他们为我的性命置自己的颈项于险地}。\BibleRef{罗 16:4}\par

你看，他们乃是装备齐全的殉道者。因为在\Person{尼禄}时代，可能有成千的危险，那时他甚至命令所有犹太人离开罗马（\BibleRef{宗 8:2}）。\par

\Quote{不但我感谢他们，连外邦人的众教会也感谢他们}。\par

这里他暗示他们的好客与金钱援助，钦佩他们，因为他们不仅倾流了血，而且将自己的全部财产向众人开放。你看，这些是高尚的妇女，在德行道上丝毫没有受到性别的阻碍。这也正是可以预期的。\BibleQuote{因为在基督耶稣内，不再有男或女}（\BibleRef{迦 3:28}）。他对前者所说的话，也用在后者身上。因为对她，他也说过：\Quote{她曾帮助了许多人，也帮助了我}。同样，对这个妇人，他写道：\Quote{不但我感谢他们，连外邦人的众教会也感谢他们}。为了在此不显得是奉承，他还为这些妇女援引了许多其他见证。\par

第五节。\BibleQuote{同样，也问候在他们家里的教会}。\BibleRef{罗 16:5}\par

因为她如此可敬，甚至使他们的家成为教会，不仅使家中所有人都成为信徒，也因为他们向所有陌生人开放。因为他不惯于称任何房屋为教会，除非那里有深厚的虔诚和对天主的敬畏扎根其中。正因如此，他也曾对格林多人说：\BibleQuote{请问候阿桂拉和普黎斯加，以及在他们家里的教会}（\BibleRef{格前 16:19}）。而当他写到\Person{敖尼息摩}时：\Quote{保禄致费肋孟，并致亲爱的阿丕雅，以及在他们家里的教会}（\EditorialNote{费肋孟书1,2}）。因为一个人即使在婚姻状态中，也能成为值得仰望和高尚的人。你看，这些人都处于婚姻状态，却变得极其尊贵，而他们的职业远非尊贵，因为他们是\BibleQuote{帐棚匠}。然而，他们的德行遮盖了这一切，使他们比太阳更耀眼。他们的职业和他们的婚姻（\OriginalTerm{婚姻}{\GreekText{συζυγία}}，参\BibleRef{斐 4:3}）都没有伤害他们，反而显出了基督所要求的爱。\BibleQuote{人若为自己的朋友舍命，再没有比这更大的爱了}，祂说（\BibleRef{若 15:13}）。也成就了作门徒的证明，因为他们背起十字架跟随了祂。因为那些为保禄这样做的人，更会为基督显示他们的刚毅。\par

让富人和穷人都来听这一切。因为如果那些靠劳作谋生、管理作坊的人，竟能如此慷慨大方，以致有益于许多教会，那么那些富有却忽视穷人的人，还能期望得到什么宽恕呢？因为他们为了天主的旨意，连自己的血也不吝惜，而你们却吝惜微薄的钱财，多次连自己的灵魂也不顾惜。可是他们对老师尚且如此，对门徒却不这样吗？不，连这一点也不能说。因为\Quote{外邦人的教会}，他说，\Quote{感谢他们}。然而他们原是犹太人。但他们仍然有如此清晰的（\OriginalTerm{清晰}{\GreekText{εἰλικρινώς}}）信德，以致也心甘情愿地服事他们。妇女应当这样，不是以\BibleQuote{编发、黄金、或华服}（\BibleRef{弟前 2:9}）来装饰自己，而是以这些善行。因为试问，有哪位皇后曾像这位帐篷制造者的妻子那样引人注目、那样受人称颂呢？她为众人所传诵，不是十年二十年，而是直到基督再来，大家都宣扬她的名声，是出于那些比任何王冠更增添光彩的事。因为有什么比救助保禄、在自己危险中拯救世界的导师更伟大、或如此伟大的呢？试想：有多少皇后无人提及。但这位帐篷制造者的妻子却与帐篷制造者（或许意指圣保禄）一起到处被传扬；太阳所照的范围超不过这个世界，而这位妇女的荣耀所及却遍及各处。波斯人、西徐亚人、色雷斯人以及住在地极的人，都歌颂这位妇女的基督徒精神并祝福它。要获得这样的见证，你不愿付出多少财富、多少冠冕和紫袍呢？因为没有人能说，他们在危险中是这般的，对钱财慷慨，却忽视了宣讲。他因此称她们为\Quote{同工和助手}。而这个\BibleQuote{蒙选的器皿}（\BibleRef{宗 9:15}）并不羞于称一位妇女为他的助手，反而以此为荣。因为他所看重的不是性别（\OriginalTerm{本性}{\GreekText{φύσει}}），而是他所尊重的意愿。这样的装饰有什么可相比的呢？如今哪里还有泛滥的财富？哪里还有身体的装饰？哪里还有虚荣？要学习：妇女的装饰不是披在身上的，而是装饰灵魂的，这装饰永不脱去，不藏在箱子里，而是积蓄在天上。请看他们为宣讲所付出的劳动、殉道者的冠冕、钱财的慷慨、对保禄的爱、他们在基督里所找到的魅力（\OriginalTerm{魅力}{\GreekText{φίλτρον}}）。将你自己的处境、你对金钱的焦虑、你与娼妓（即衣着上）的竞争、你对草的效法与之相比，然后你就会看清他们是谁、你又是谁。或者不如不要比较，而要效法这位妇女，放下草的重担（\OriginalTerm{草}{\GreekText{χλόης}}）（因为你们昂贵的衣着就是如此），穿上从天上来的衣裳，学习普黎斯加是怎样变成这样的。他们是如何变成这样的呢？他们留保禄住了两年（\BibleRef{宗 19:10}），这两年对他们的灵魂有什么不能成就呢？那么你会说，我没有保禄，我该怎么办呢？如果你愿意，你可以比他们更真实地拥有他。因为即使是他们，看到保禄并非使他们成为这样的人，而是保禄的话语。所以，如果你愿意，你可以让保禄、伯多禄、若望以及整队先知和宗徒持续与你交往。因为取出这些蒙福之人的书卷，持续与他们的著作来往，它们就能使你成为帐篷制造者的妻子。我为何只说保禄呢？因为如果你愿意，你甚至可以拥有保禄的主自己。因为透过保禄的口，他甚至会与你交谈。还有另一种方式你也能接待这位，即当你们接待圣人、服事那些信从他的人时。因此即使他们离世后，你们仍有许多虔诚的纪念物。因为圣人的桌子、坐过的座位、躺过的床，甚至在圣人离去后，都能感动接待他的人。那么你想，那位叔纳米特女子进入厄里叟曾住过的楼房时，看见桌子、圣者睡过的床时，她岂不是深受感动吗？她必定从中感受到何等虔诚？若非如此，孩子死后她就不会把他放在那里（厄里叟的床上），如果她未从那里得到极大的益处。因为若在如此长的时间之后，当我们进入保禄住过的地方、他被捆绑的地方、他坐而论道的地方，我们的精神被提升，并从那些地方开始回到那记忆（按菲尔德的版本：武加大译本作\Quote{那一天的记忆}）；当情形尚更近时，那些曾虔诚接待他的人岂不更该有感触吗？既然知道这一切，让我们接待圣人，使房屋发光，脱离扼制的荆棘，使卧室成为避难所。让我们接待他们，洗他们的脚。你不比撒辣更好，更高贵，更富有，即使你是皇后。因为她在自己有三百一十八名家生仆役的时候——在那个时代有两名仆役就已是富有。我为何提三百一十八名仆役呢？她在自己的后代和应许中得着了整个世界，她有\BibleQuote{天主的朋友}（\BibleRef{依 41:8}）作丈夫，天主自己作保护者——这比任何王国都大。然而，虽然她处于如此显赫尊贵的地位，这位妇女却和面，做所有其他仆人的杂务，并在设宴时如同仆人般站在他们旁边。你并不比亚巴郎更高贵，他在得胜之后、在埃及王尊敬他之后、在驱逐波斯王之后、在树立光荣的旗帜之后，仍做仆人的事。不要看外表：住在你们家里的圣人可能是贫穷的、像乞丐、常穿破衣，但要记住那声音说：\BibleQuote{你们对我这些最小兄弟中的一个所做的，就是对我做的}（\BibleRef{玛 25:40}）。并且，\BibleQuote{不要轻视这些小子中的一个，因为他们的天使在天上常见我父的面}（\BibleRef{玛 18:10}）。因此要乐意接待他们，因为他们带来千万的祝福给你们，藉着平安的问候（\BibleRef{玛 10:12-13}）。在撒辣之后，也思想黎贝加，她打水给人喝，并招呼陌生人进来，践踏了一切傲慢。然而藉着这，她获得住宿接待的丰厚赏报！而你，如果你愿意，将获得比那些更大的。因为天主不仅会赐给你儿女的果实，而且赐给你天堂和那里的祝福，免于地狱和罪过的赦免。因为住宿接待的果实是巨大的（\BibleRef{路 11:41}）。同样，耶特洛虽是外邦人，却得以与那位以大能命令大海的人成为亲戚（\BibleRef{达 4:27}）。因为他的女儿们也将这尊贵的猎物拉入他的网中（\BibleRef{户 10:29}）。因此，将你们的思念放在这些事上，并反思那些妇女的刚毅英勇气概，践踏今世的奢华、衣着的装饰、昂贵的珠宝、香膏的涂抹。并抛弃那些轻浮、矫揉造作的态度和扭捏的行走，将一切注意力转向灵魂，并在你们心中燃起对天堂的渴望。因为一旦他的爱抓住你们，你们就会分辨出泥浆和尘土，并嘲笑现在被赞叹的事物。因为一个以属灵成就为装饰的妇女，甚至不可能追求这些可笑的东西。因此，你们要抛弃这些只有粗野男子的妻子、女演员和歌女才如此热衷的东西，披上智慧之爱、住宿接待、救助圣人、痛悔、不断祈祷。这些好过金衣，这些比宝石和项链更庄严，这些既使你们在众人中有好名声，又带你们在天主面前有丰厚赏报。这是教会的服饰，那是剧院的服饰。这配得上天堂，那配得上马和骡子；那甚至被裹在死尸上，这唯独在善灵中闪耀，基督居住在其中。因此，让我们获取这服饰，使我们也能到处受称赞，并中悦基督，祂是……等。阿们。\par

\SourceRecord{Translated by J. Walker, J. Sheppard and H. Browne, and revised by George B. Stevens.；Nicene and Post-Nicene Fathers, First Series；Vol. 11.；Edited by Philip Schaff.；Buffalo, NY: Christian Literature Publishing Co.,；1889.；<http://www.newadvent.org/fathers/210230.htm>.}

% structure: p0001 source=h1 target=section reason=page-title

\section{罗马书第三十一篇讲道}

\BibleRef{罗 16:5}\par

\BibleQuote{请问候我深爱的\Person{厄派内托}，他是\Place{阿哈雅}归向基督的初果。}\par

我认为，甚至许多表面上是极好的人，也匆匆掠过书信的这一部分，认为它是多余的，没有什么分量。而且我认为，他们在对待福音中的族谱时也是如此。因为那是一份名单，他们便以为不能从中获得什么大的益处。然而，铸金匠人甚至对小小的碎屑也小心谨慎；这些人却连如此大的金块也忽略过去。那么，为了使他们不致如此，我前面所说的已经足以使他们脱离懈怠。因为即便从这部分所得的利益也不可轻视，我们从前一次讲道时已经说明了，当时我们借着这些讲词振奋了你们的灵魂。因此，我们今天也要努力在同一处挖掘。因为即便从单纯的名字中，也能发现巨大的宝藏。例如，如果有人向你解释\Person{亚巴郎}为何如此称呼，\Person{撒辣}、\Person{以色列}、\Person{撒慕尔}又为何如此称呼，你就会从中发现许多真正值得研究的课题。而且从时间、地点，你也可以得到同样的好处。因为好人甚至从这些也能变得富有；但懒惰的人，即使从最明显的事物中也得不到益处。同样，\Person{亚当}这个名字本身教导我们大智慧，他儿子、妻子以及大多数其他人的名字也是如此。因为名字提醒我们许多情况。它们同时显示出天主的恩惠和妇女的感恩。因为当她们因天主的恩赐而怀孕时，是她们给孩子们取这些名字。但我们为何现在讨论名字，而如此重要的意义却被忽略，许多人甚至连圣书的名字都不知道？即便如此，我们也不应放弃对这类事情的关注。因为祂说：\BibleQuote{你应当把我的银子交给兑换银钱的人。}\BibleRef{玛 25:27}因此，即使没有人听从，我们也要尽自己的本分，证明圣经中没有多余、没有随意添加的内容。因为如果这些名字没有用处，它们就不会被添入书信，\Person{保禄}也不会写下他所写的内容。但有些人甚至如此卑微、空虚、不配天堂，以至于认为不仅名字，而且圣经中整卷书都没有用处，如\WorkTitle{肋未纪}、\WorkTitle{若苏厄书}以及其他更多。就这样，许多单纯的人一直拒绝旧约，并沿着这条由邪恶思想习惯所产生的道路前进，同样也删去了新约的许多部分。但对于这些人，他们如同醉酒、顺从肉身生活，我们不予重视。但如果有人是智慧的热爱者、属神娱乐的朋友，就让他知道，即使圣经中看似不重要的事物，也不是随意、无目的地放置在那里的，甚至旧律法也对我们大有裨益。因为经上说：\BibleQuote{这一切都是预像（英王钦定本作榜样），并是为教训我们而记录的。}\BibleRef{格前 10:11}因此，他也对\Person{弟茂德}说：\BibleQuote{要注重宣读、劝勉}\BibleRef{弟前 4:13}，以此敦促他阅读旧约书籍，尽管他是一位拥有如此伟大精神的人，甚至能驱逐魔鬼、复活死人。现在我们继续手头的话题。\BibleQuote{请问候我深爱的\Person{厄派内托}。}从这一点值得学习他是如何分配不同赞美给每个人的。因为这赞美并非微小，而是非常伟大，证明了他卓越的品质，使\Person{保禄}视他为可爱的；\Person{保禄}从不会出于偏袒、而非冷静判断去爱人。然后另一个赞扬来了：\BibleQuote{他是\Place{阿哈雅}的初果。}因为他的意思要么是他比任何人更早跃出，成为信徒（这也不是轻微的赞美），要么是他比其他人表现出更虔敬的行为。因此，在说了\BibleQuote{他是\Place{阿哈雅}的初果}之后，他并没有保持沉默，而是为了防止你怀疑这是世俗的荣耀，他补充道：\BibleQuote{归向基督。}如果在世俗事务中，首先的人显得伟大可敬，那么在这些事上就更当如此。由于他们很可能出身低微，他便谈论真正的出身高贵和卓越，并由此赐予他尊荣。并且他说，他\BibleQuote{是初果}，不仅是\Place{格林多}的，而且是整个民族的初果，因为他仿佛成了其余人的门和入口。对于这样的人，奖赏不小。因为这样的人也会从别人的成就中获得丰厚的回报，因为他的开始也为这些成就贡献了许多。\par

\BibleRef{罗 16:6}. \BibleQuote{问候\Person{玛利亚}，她为我们多受劳苦。}\par

这是怎么回事？一位妇女再次受到尊敬并被宣告得胜！我们男人再次蒙羞。或者更确切地说，我们不只蒙羞，反而获得了尊荣。因为我们蒙受尊荣，在于我们当中有这样的妇女，但我们蒙羞，在于我们男人被她们远远抛在后面。然而，如果我们知道她们的这些装饰从何而来，我们也会很快赶上她们。那么，她们的装饰来自何处？愿男女都听。并非来自手镯、项链，也不是来自他们的太监、婢女和锦绣衣裳，而是来自她们为真理所付出的辛劳。因为他说：\Quote{她们为我们付出了极大的辛劳}，即不单是为了她们自己，也不单是为了自身的进步（见第520页）（因为现今许多妇女通过禁食和睡硬地来这样做），而且还为了别人，如此承继了宗徒和传福音者所奔跑的赛程。那么，他说\BibleQuote{我不许女人施教？} \BibleRef{弟前 2:12} 是什么意思？他意思是阻止她公开站出来 \BibleRef{格前 14:35}，并阻止她登上讲坛，而不是阻止她施教的言语。因为如果情况如此，他怎么会对那有不信主之丈夫的妇女说：\BibleQuote{女人，你怎么知道你能不能救你的丈夫？} \BibleRef{格前 7:16} 或者，他怎么会允许她劝诫儿女，当他说：\BibleQuote{但她若持守信德、爱德和圣洁，持重自持，必因生育而得救} \BibleRef{弟前 2:15}？怎么\Person{普黎史拉}甚至教导了\Person{阿颇罗}？那么，他说这话并不是要切断私下有益的谈话，而是指在众人面前，以及属于教师在公众集会中应尽的职责；或者，又是在丈夫有信德且准备妥当，也能够教导她的情况下。当她更明智时，他不禁止她教导和改善他。并且他没有说\Quote{她教导了许多}，而是说\Quote{她付出了极大的辛劳}，因为除了教导（\OriginalTerm{言语}{\GreekText{τοὓ} \GreekText{λόγου}}）之外，她还从事其他服务，那些涉及危险、金钱、旅行方面的服务。因为当时的妇女比狮子还要英勇，为福音的缘故与宗徒分担辛劳。她们就这样与他们同行，并执行所有其他服务。甚至在基督的时代，有妇女跟随着祂，\BibleQuote{用自己的财物供养祂} \BibleRef{路 8:3}，并服侍这位导师。\par

第七节。\BibleQuote{请问候我的亲人安得洛尼科和犹尼亚。} \BibleRef{罗 16:7}\par

这也像是一个颂词。接下来的部分更是如此。这是什么样的颂词呢？\BibleQuote{我的囚友。} 因为这是最大的尊荣，是荣耀的宣告。保罗何时成为囚犯，竟称他们为\BibleQuote{我的囚友}？他确实未曾被囚，但他所遭受的比囚犯更甚，他不仅背井离乡，还与饥饿、不断的死亡以及成千上万的其他事物搏斗。因为囚犯的唯一不幸是，他与亲属分离，并且常常不再是自由人而成为奴隶。但在这里，可以提及多如雪片的诱惑，这个有福的人经受着被拖曳、被带走，受鞭打、被捆绑、被石头砸、遭遇海难，无数人谋害他。而被俘的人在押走后不再有敌人，反而受到俘获者的极大照顾。但这个人一直处于敌人之中，到处看见长矛、锋利的刀剑、阵列和战斗。既然他们很可能与他分担了许多危险，他就称他们为囚友。如同在另一处也写道：\BibleQuote{阿黎斯塔古，我的囚友。} \BibleRef{哥 4:10} 然后是另一种赞美。\BibleQuote{他们在宗徒中是有声望的。} 成为宗徒本身就是伟大的事。但甚至在宗徒中也闻名，想象这是多么大的赞美！然而，他们是因着他们的工作和成就而闻名。哦！这个女人的虔敬（\OriginalTerm{虔敬}{\GreekText{φιλοσοφία}}）多么伟大，她竟配得上称为宗徒！但在这里他也没有停止，又加上另一种赞美，说：\BibleQuote{他们也是在我之前已在基督内的。}\par

因为这也是一个非常大的赞美。但让我们注意这神圣的灵魂，它如何未被虚荣所玷污。因为在获得如此程度和种类的荣耀之后，他把别人放在自己前面，没有向我们隐瞒他落后于他们的事实，也不羞于承认这一点。既然他因此不能在这方面把别人放在前面，他就观察自己，即后来归信的人，从这里他找到了赞美他们的方式，说：\BibleQuote{他们是在我之前已在基督内的。} 当他不以承认此点为羞，竟然不害怕在众人面前展示他以前的生活，自称\BibleQuote{亵渎者、迫害者} \BibleRef{弟前 1:13}？\par

第八节。\BibleQuote{请问候我所爱的安普里亚。} \BibleRef{罗 16:8}\par

这里他通过对他的爱来赞美他。因为保罗的爱是为了天主，包含着无数的祝福。如果被国王喜爱是一件大事，那么被保罗喜爱是多么大的赞美！因为如果他没有获得极大的美德，怎能吸引他的爱？因为对于那些生活在罪恶和过犯中的人，他习惯（\OriginalTerm{习惯于}{\GreekText{οἶδε}}）不仅不喜爱他们，甚至诅咒他们。如他说：\BibleQuote{若有人不爱主耶稣，他就该受诅咒} \BibleRef{格前 16:22}；以及\BibleQuote{若有人给你们宣讲与你们所接受的不同福音，他就该受诅咒。} \BibleRef{迦 1:8}\par

第九节。\BibleQuote{请问候乌尔巴诺，他在主内是我的助手。} \BibleRef{罗 16:9}\par

这是比另一个更大的赞美。因为这个甚至包括了那个。\BibleQuote{还有斯塔辉，我所爱的。} 这又是同样的荣耀。\par

第十节。\BibleQuote{请问候阿培肋，他在基督里是经得起考验的。} \BibleRef{罗 16:10}\par

没有这样的赞美：无可指责，在天主的事上不给任何把柄。因为当他说\Quote{在基督里蒙悦纳}时，他就包括了整个德行的清单。那么他为何从不称某人为\Quote{我的主}或\Quote{我的师父}呢？因为这些颂词比那些更大。因为那些仅仅是地位的称号（\OriginalTerm{尊荣}{\GreekText{τιμής}}），但这些是德行的。他并非随意给予这种尊荣，也不是把低微德行的人与高尚伟大的人物一同称呼。就他致意而言，而且是把此人同那人一起，在同一封信中，他一视同仁地尊敬他们。但通过特别向每个人陈述赞美，他就给我们呈现了每个人特有的德行；这样既不会因尊敬一人而羞辱另一人而引起嫉妒，也不会因给予所有人同样的尊严（尽管他们不配得到同样的）而使他们产生懈怠和混乱。现在请看他是如何再次转向那些可敬佩的妇女的。因为在说了\Quote{要问候阿黎斯托步罗家中的人}之后，\par

第十一节。\BibleQuote{要问候我的同族黑落狄翁；问候纳尔基索家中的人。}\BibleRef{罗 16:11}\par

这些人很可能不如前面那些人那么配得，因此他甚至连他们的名字也未曾一一提及，而在给了他们合适的赞美（即忠诚）之后，（而这就是以下言辞的意思：）\par

\Quote{在主内的人。}\par

他又回到妇女们身上，说：\par

第十二节。\BibleQuote{要问候特黎斐纳和特黎佛撒，她们在主内劳苦。}\BibleRef{罗 16:12}\par

关于前一位妇女，他说她\Quote{曾在你们身上劳苦}；至于这些妇女，她们仍在劳苦。这是一个不小的赞美：她们应始终在工作，而且不仅要工作，还要劳苦。但他也称佩尔西斯为可爱的，以表明她比这些更大。\par

因为他说：\Quote{要问候我可爱的佩尔西斯。}\par

关于她巨大的劳苦，他也同样作证说：\Quote{她在主内劳苦甚多。}\par

他何等清楚如何按各人的功过称呼他们，如此既不让任何人失去应得的，甚至称赞最微小的优越，使这些人更加热切；也使其他人更加有德，通过他对这些人的赞美激励他们同样的热忱。\par

第十二节。\Quote{要问候鲁富斯，在主内蒙拣选的人，以及他的母亲，也是我的母亲。}\par

这里再次有毫无缺点的美善，因为儿子和母亲各具这样的品格，家中充满祝福，根与果实一致；因为他不会简单地说\Quote{他的母亲，也是我的母亲}，除非他在为这妇女的伟大德行作证。\par

第十四节。\BibleQuote{要问候阿斯因克利托、弗勒贡、赫尔玛斯、帕特洛巴、赫尔梅斯，以及与他们一起的弟兄们。}\BibleRef{罗 16:14}\par

这里不要关注他如何没有赞美就开始，而要关注他如何没有认为他们（尽管似乎远不如所有人）不配被他问候。甚至这也是一种不小的赞美：他甚至称他们为弟兄，正如后面那些人他称他们为圣人。因为他说：\par

第十五节。\BibleQuote{要问候非罗罗果、犹里、乃勒乌和他的姐妹、敖林帕，以及与他们一起的所有圣人。}\BibleRef{罗 16:15}\par

这是最大的尊严，无可言喻的尊荣高度。然后，为防止因他以不同方式致意不同的人（一些提名，一些不加区别，一些有更多赞美，一些较少）而引起嫉妒，他再次在爱德的平等中，在圣吻中，将他们混合起来，说：\par

第十六节。\BibleQuote{你们要以圣吻彼此问候。}\BibleRef{罗 16:16}\par

通过这个问候，要从他们中除去所有使他们混乱的争论，以及所有骄傲的小理由；这样，大的不会轻视小的，小的也不会嫉妒大的，而是当这个吻安抚并平等对待每个人时，傲慢和嫉妒就会被驱散。因此，他不仅命令他们以此方式问候，而且也将各教会的问候同样送给他们。因为他说，\Quote{问候你们}的，不是这个或那个人个人，而是你们众人共同，\par

\Quote{基督的众教会。}\par

你们看，我们从这些讲道中获得的益处并不小，我们本该匆匆略过的宝藏，若非在此书信的这一部分也仔细查考了它——即如我们能力所及——我们就会错过。所以，若发现任何有智慧且属神的人，他会潜入更深，发现更多的珍珠。但因为有些人常提出疑问：为何在这封书信中他问候了那么多人，而他在其他任何书信中都没有这样做？我们可以说，这是因为他还从未见过罗马人。然而有人可能说：\Quote{但他也没有见过哥罗森人，却并没有做任何类似的事。}但这些罗马人比其他城的人更尊贵，他们从其他城市来到那里，如同来到一个更安全、更皇家的城市。既然他们生活在异乡，需要许多安全的供应，而他们中有些是他的熟人，也有些为他提供了许多重要服务，他合理地通过书信称赞他们；因为\Person{保禄}的荣耀当时并不小，而是如此之大，以至于甚至他寄出书信，那些有幸收到他书信的人就获得了很大的保护。因为人们不仅尊敬他，甚至害怕他。若非如此，他就不会说那曾\BibleQuote{援助了许多人，也援助了我}\BibleRef{罗 16:2}。又说：\BibleQuote{我宁愿自己受诅咒。}\BibleRef{罗 9:3}又给\Person{费肋孟}写信说：\BibleQuote{像我\Person{保禄}这般年老，且是耶稣基督的囚犯。}\EditorialNote{斐理伯书 9}又给迦拉达人：\BibleQuote{看，我\Person{保禄}告诉你们。}\BibleRef{迦 5:2}又说：\BibleQuote{你们接待我如同耶稣基督。}\BibleRef{迦 4:14}又写信给格林多人说：\BibleQuote{有些人自高自大，以为我不会到你们那里去。}\BibleRef{格前 4:18}又说：\BibleQuote{我把这些事用了比喻转到我自己和\Person{阿颇罗}身上，好叫你们在我们身上学习，不可超过圣经所记载的看人。}\BibleRef{格前 4:6}从所有这些段落可以清楚看出，众人都对他有很高的评价。他愿意他们感到自在且受尊重，于是向他们每个人致意，尽力彰显他们的美德。因为他称这个为亲爱的，那个为同族，另一个既是亲爱的又是同族，另一个为同囚，另一个为同工，另一个为蒙悦纳的，另一个为蒙选的。至于妇女，他用她的头衔称呼一位，因为他没有含糊地称她为教会的仆人（因为若是如此，他也会给\Person{特黎费纳}和\Person{培尔息得}这个名称），而是称这位有女执事职分的，另一位为帮手和协助者，另一位为母亲，另一位则因她所受的劳苦；有些人他根据他们所属的家庭来称呼，有些人用弟兄的名称，有些人用圣徒的称号。有些人他单单因致意而尊荣他们，有些人则因直呼其名，有些人因称他们为初果，有些人因他们在时间上的优先，但最突出的是\Person{普黎史拉}和\Person{阿桂拉}（\OriginalTerm{普黎史拉和阿桂拉等人}{\GreekText{τοὺς} \GreekText{περὶ} \GreekText{Πρ}. \GreekText{κ}. \GreekText{᾿Α}}）。因为即使所有人都是信徒，但并非全都相同，而是功德各有不同。因此，为激励他们所有人更加努力，他没有隐藏任何人的赞美。因为当劳苦更多的人没有得到更大赏报时，许多人变得更懈怠。基于此，即使在王国（天国）里，光荣也不相等，门徒们也不全一样，而是那三位超乎其余。这三位之间又有很大差别。因为天主坚持这种精确的方法直到最后。因此，经上说：\BibleQuote{这星与那星的荣耀不同，}\BibleRef{格前 15:41}，它说。然而，所有人都是宗徒，都要坐在十二宝座上，都舍弃了财产，都与他相处；但祂仍只带了那三位。而且，对这三位，他说可能（\OriginalTerm{可能}{\GreekText{ἐγχωρεῖν}}）有人甚至更优越。他说：\BibleQuote{只是坐在我的右边和左边，不是我可以赐的，而是给谁预备的。}\BibleRef{谷 10:40}祂又将\Person{伯多禄}置于他们之前，说：\BibleQuote{你爱我胜过这些吗？}\BibleRef{若 21:15}\Person{若望}也更为蒙爱。\par

那么，那些看见这一切极大的精确，却仍不允许有地狱的人在哪里呢？因为如果所有义人不能享受同样的命运，如果他们比别人稍胜一筹——经上记载：\BibleQuote{这颗星与那颗星在光荣上有所不同}（\BibleRef{格前 15:41}）——那么罪人怎能与义人同等呢？这样的混乱连人也不会造成，何况天主呢！但如果你愿意，我要向你证明，即使在罪人的事上，根据现有的事实，也有这种区别和精确的公义审判。现在想想：亚当犯了罪，厄娃也犯了罪，两人都违命，但他们并非同样有罪，因此也没有受到同样的惩罚。因为差别如此之大，以致保禄说：\BibleQuote{亚当没有受骗，而是女人受了骗，陷于过犯。}然而欺骗是同一个。但天主的审查明察秋毫，指出了如此巨大的差别，以致保禄说了这话。再者，加音受了惩罚，但在他之后杀人的拉默客并没有遭受几乎同样大的惩罚。虽然这也是谋杀，那也是谋杀，而且后者更糟，因为他连榜样也没有使他变好。但由于前者既没有在劝诫后杀弟，也不需要原告，在天主质问他时也不退缩，反而在没有任何原告的情况下，自己为自己辩护，并更严厉地定自己的罪，他获得了宽恕。但后者因做了相反的事而受罚。请看天主如何精确地筛选事实。为此，祂用一种方式惩罚了洪水时期的人，用另一种方式惩罚了索多玛的人；又用不同的方式惩罚了以色列人，无论是巴比伦时期的，还是安提约古时代的：以此表明祂严格记录我们的行为。他们为奴七十年，那些为四百年，但另一些又吃自己的儿女，遭受无数其他更严重的灾祸，即便如此，无论是他们还是在索多玛被烧死的人，都没有得到释放。祂说：\BibleQuote{的确，为索多玛和哈摩辣地方，比那城更容易忍受。}（\BibleRef{玛 10:15}）因为如果祂不关心我们，无论我们犯罪或行善，也许说没有惩罚还有一些道理。但由于祂如此迫切地要我们不犯罪，并用这么多方法使我们持守正道，显然祂惩罚恶人，也加冠给行善的人。但请允许我请你们考虑多数人的不公。他们责怪天主，因为祂如此多次容忍，赦免了那么多不敬、不洁或暴力的人，而他们现在没有遭受惩罚。再者，如果祂威胁要在来世惩罚他们，他们就强烈而紧迫地指控。然而如果这是痛苦的，他们本应接受并钦佩另一面。但唉，愚蠢！这无理和愚蠢的灵！唉，这喜爱罪恶的灵魂，它凝视着邪恶！因为正是从这里产生所有这些意见。所以，如果说出这些事的人有心抓住美德，他们很快就会在地狱的问题上也感到满意，不再怀疑。有人说，这个地狱在哪里，在什么地方？因为有些神话家说它在约沙法特山谷，从而把关于某场过去战争的话引到地狱上。但圣经没有这样说。那么，请问，它会在什么地方？我想至少在这个世界之外。因为正如监狱和矿场远离王宫，地狱也将在这个世界之外。那么，我们不要寻求知道它在何处，而要寻求如何逃脱它。也不要因为天主没有在这里惩罚所有人，就因此不相信将来的事。因为祂是仁慈和忍耐的：这就是为什么祂威胁，而不立刻把我们投入其中。因为祂说：\BibleQuote{我不喜欢罪人丧亡。}（\BibleRef{则 18:32}）但如果没有罪人的死亡，这些话就是空谈。我确实知道，对你们来说没有什么比这些话更不愉快。但对我来说，没有什么比这更愉快。但愿我们在吃饭、晚宴、洗澡、任何地方都在谈论地狱。因为这样我们就不会感到今世邪恶的痛苦，也不会感到美好事物的快乐。因为你告诉我什么是邪恶？贫穷？疾病？被掳？身体残缺？为什么所有这些与那里的惩罚相比都是游戏，即使你提到那些一生受饥荒折磨的人；或那些从小残疾、乞讨的人，这与其他邪恶相比也是奢侈。让我们不断地谈论这些事。因为记住地狱可以防止我们落入地狱。你没有听到圣保禄说：\BibleQuote{他们要受永远丧亡之罚，离开主的面。}（\BibleRef{得后 1:9}）你没有听到尼禄的品性是什么？保禄甚至称他为\Quote{反基督的奥秘}？因为他说：\BibleQuote{不法的奥秘已经发动。}（\BibleRef{得后 2:7}）那么怎样？尼禄不受任何惩罚？反基督不受任何惩罚？魔鬼也不受惩罚？那么他将永远是反基督，魔鬼也如此。因为他们不会停止作恶，除非受惩罚。你说：\Quote{是，但有一个地狱是每个人都看见的。只有不信者才会掉进去。}请问，原因是什么？是因为信徒承认他们的主。但这有何相干？当他们的生活不洁时，他们会因此比不信者受到更严厉的惩罚。因为\BibleQuote{凡在法律之外犯了罪的，也必在法律之外丧亡；凡在法律之内犯了罪的，必按法律受审判。}（\BibleRef{罗 2:12}）并且，\BibleQuote{那知道主人的旨意，而不去实行的仆人，必受许多鞭打。}（\BibleRef{路 12:47}）但如果没有交账这回事，这一切只是随便说说，那么魔鬼也不会被报复。因为他也认识天主，而且远超过人，所有魔鬼都认识祂，并颤抖，承认祂是他们的审判者。那么，如果没有交账，也没有恶行，他们也会完全逃脱。这些事绝非如此，确实不是！亲爱的，不要自欺。因为如果没有地狱，宗徒们怎能审判以色列十二支派？保禄怎么说：\BibleQuote{你们不知道我们将审判天使吗？何况世俗的事呢？}（\BibleRef{格前 6:3}）基督怎么说：\BibleQuote{尼尼微人要起来判罚这一代}（\BibleRef{玛 12:41}）；并且\BibleQuote{在审判的日子，索多玛地方所受的惩罚要比那城更容易忍受？}（\BibleRef{玛 11:24}）那么，为什么拿不可嬉笑的事来嬉笑？为什么欺骗自己，在自己的理性上弄虚作假（\OriginalTerm{欺骗自己的理性}{\GreekText{παραλογίζῃ}}, om. \OriginalTerm{你的灵魂}{\GreekText{τὴν} \GreekText{ψυχήνσου}}）？为什么与天主对人的爱作对？因为正是借着这个爱，祂准备了地狱，并加以威胁，好叫我们因着恐惧而变好，以免被投入其中。因此，那废止谈论这些事的，无非是借着这欺骗把我们推进地狱，赶入其中。所以，不要放松那些为美德劳苦之人的手，也不要加增那些沉睡之人的怠惰。因为如果许多人被说服没有地狱，他们何时会离弃邪恶？何时会看见正义？我不是说罪人和义人之间，而是在罪人之间？因为为什么一个人在这里受罚，而另一个人没有受罚，尽管他犯了同样的罪，甚至更糟？因为如果没有地狱，你对那些以此反驳的人将无话可说。因此我的建议是，我们停止这种玩笑，堵住那些在这些问题上反对者的口。因为将会有对最小事情的精确审查，无论是罪还是善行，我们因不洁的目光、闲话、仅仅是辱骂的话受罚，因醉酒交账，同样，一杯凉水也会得到赏报，仅仅一声叹息。（\BibleRef{训 12:14}）因为经上说：\BibleQuote{要在那些叹息和哀号的人额上画一个记号。}（\BibleRef{则 9:4}）那么你怎么敢说，那用如此精确审查我们行为的主，只是用空话和轻率地威胁地狱？不要，我恳求你，不要用这些虚妄的希望摧毁你自己和那些被你说服的人！因为如果你不信我们的话，就问犹太人、外邦人和一切异端者。他们都会异口同声地回答，审判和报应一定存在。人还不够吗？问魔鬼自己，你会听见他们喊：\BibleQuote{你到这里来折磨我们，时候还没有到吗？}（\BibleRef{玛 8:29}）把这一切综合起来，说服你的灵魂不要无谓地玩笑，免得因经验而知道有地狱，而是借此清醒，从而能够逃脱那些折磨，获得将来的美好事物；愿我们所有人借着恩宠和对人的爱等，都能有分。\par

\SourceRecord{Translated by J. Walker, J. Sheppard and H. Browne, and revised by George B. Stevens.；Nicene and Post-Nicene Fathers, First Series；Vol. 11.；Edited by Philip Schaff.；Buffalo, NY: Christian Literature Publishing Co.,；1889.；<http://www.newadvent.org/fathers/210231.htm>.}

% structure: p0001 source=h1 target=section reason=page-title

\section{罗马书讲道集第三十二篇}

\BibleRef{罗 16:17,18}\par

\BibleQuote{弟兄们，我恳求你们，留意那些造成分裂和绊脚石、违反你们所学的道理的人，并要躲避他们。因为这样的人不是服事我们的主耶稣基督，而是服事自己的肚腹；并且以花言巧语欺骗了纯朴人的心。}\par

再次，是一个劝勉，以及劝勉之后的祈祷。因为在他告诉他们要\BibleQuote{留意那些造成分裂的人}，并且不要听从他们之后，他接着说：\BibleQuote{赐平安的天主快要将撒殚践踏在你们脚下；愿我们主的恩宠与你们同在。}并且注意他是多么温柔地劝勉他们：不是以顾问的身份，而是以仆人的身份，并带着极大的尊重。因为他称他们为弟兄，也同样恳求他们。因为他说：\BibleQuote{弟兄们，我恳求你们。}然后他也通过显示那些滥用他们之人的诡诈，使他们有所防备。因为好像他们不能立刻被识别出来，他说：\BibleQuote{我恳求你们留意，}也就是说，要特别仔细地注意、熟悉并彻底查究——请问是谁呢？就是\BibleQuote{那些造成分裂和绊脚石、违反你们所学的道理的人。}因为这正是教会的颠覆，就是分裂。这是魔鬼的武器，这使一切颠倒。因为只要身体合而为一，他就无法进入，但分裂带来绊脚石。分裂从何而来？从违反宗徒教训的意见而来。这种意见从何而来？从人成为肚腹和情欲的奴隶而来。因为\BibleQuote{这样的人，}他说，\BibleQuote{不是服事主，而是服事自己的肚腹。}所以，除非有人想出了违反宗徒教训的意见，就不会有绊脚石，不会有分裂。他在这里指出这一点，说：\BibleQuote{违反道理。}他没有说我们所教导的，而是说\BibleQuote{你们所学习的，}以此预先防范，表明他们相信了、听到了并接受了。对于那些以这种方式制造祸端的人，我们该做什么？他没有说开会和打架，而是说\BibleQuote{躲避他们。}因为如果他们是出于无知或错误而这样做，就应该纠正他们。但如果他们愿意犯罪，就离开他们。他在另一处也这样说。因为他说：\BibleQuote{你们要离开每一个行事不按规矩的弟兄}\BibleRef{得后3:6}；在对弟茂德谈到铜匠时，他也给予相同的劝告，说：\BibleQuote{你也要提防他。}\BibleRef{弟后4:15}然后，为了嘲弄那些胆敢做这些事的人，\OriginalTerm{嘲弄}{\GreekText{κωμῳδὥν}}，他也提到了他们策划分裂的原因。他说：\BibleQuote{因为这样的人不是服事我们的主基督，而是服事自己的肚腹。}他写信给斐理伯人时也说过：\BibleQuote{他们的天主是自己的肚腹。}\BibleRef{斐3:19}但在这里，他似乎暗示那些犹太人，他特别常常责备他们贪吃。因为他在写信给弟铎时也论到他们说：\BibleQuote{恶兽，懒惰的肚腹。}\BibleRef{铎1:12}基督也为此责备他们，他说：\BibleQuote{你们吞没寡妇的家产}\BibleRef{玛23:14}。先知们也责备他们这类事。因为他说：\BibleQuote{我的爱人，长得肥胖粗壮，就踢跳起来。}\BibleRef{申32:15}因此梅瑟也劝诫他们，说：\BibleQuote{当你吃了、喝了、饱了的时候，要记念上主你的天主。}\BibleRef{申6:11-12}在福音中，那些对基督说：\BibleQuote{你给我们显什么神迹？}\BibleRef{若6:30}的人，忽略其他一切，只记得玛纳。所以，他们到处都显得被这种情感所控制。那么，当你是基督的弟兄时，你怎能不以把肚腹的奴隶当作老师为耻呢？错误的根源是这样，但攻击的方式又是另一种混乱，即谄媚。因为他说：\BibleQuote{他们以花言巧语}\BibleQuote{欺骗了纯朴人的心。}因为他们的注意力只停留在言语上，但他们的意思并非如此，因为其中充满欺诈。他没有说他们欺骗了你们，而是说\BibleQuote{纯朴人的心。}即使这样他仍不满意，为了使这个说法不那么刺耳，他说：\par

\BibleRef{罗 16:19} \BibleQuote{因为你们的服从已经传于众人。}\par

他这样做，不是为了让他们可以无耻，而是先用称赞和众多证人来争取他们，以吸引他们的注意。因为不单是我作证，而是全世界。他没有说\BibleQuote{你们的理解}，而是说\BibleQuote{你们的服从}，即你们的顺服，这证明了你们极大的温良。\BibleQuote{所以我为你们欢喜。}这也是不小的称赞。然后，在称赞之后，是告诫。因为他怕在免除对他们的任何指责之后，他们会因为没有受到注意而更加懈怠；于是他用以下的话给他们另一个提示：\par

\BibleQuote{我愿你们在善事上聪明，在恶事上纯朴。}\par

你看，他如何又攻击他们，而且他们没有怀疑。因为这暗示他们中有些人容易被误导。\par

\BibleRef{罗 16:20} \BibleQuote{赐平安的天主快要将撒殚践踏在你们脚下。}\par

因为他提到那些\Quote{在你们中间制造分裂和绊脚石的人}后，也提到了\Quote{平安的天主}，好让他们对摆脱这些邪恶怀抱希望。因为那为此（即平安）喜乐的，必将终止那破坏平安的事。他并没有说\Quote{制伏}，而说\Quote{踏碎}\BibleRef{创 3:19}，这是更强烈的表达。而且不仅那些人，还有那在这事上作他们统帅的撒殚。他不仅说\Quote{踏碎}，还说\Quote{在你们脚下}，使他们亲自获胜，因这战利品而高贵。时间本身也成了安慰的理由，因为他加上\Quote{不久}。这既是祈祷，也是预言。\Quote{愿我们的主耶稣基督的恩宠与你们同在。}\par

这最强的武器！这不可攻破的城墙！这不摇撼的高塔！因为他提醒他们恩宠，好给他们更大的热忱。如果你们已从更严重的邪恶中被解救出来，且是单靠恩宠，那么现在你们已成为朋友，并同样贡献自己的份，就更会从较小的邪恶中被解救出来。你看他如何既不把祈祷置于行为之外，也不把行为置于祈祷之外。因为在他们因服从而得称赞后，他才祈祷，以表明我们两方面都需要：我们自己的份和天主的份，才能被妥善拯救。因为不仅从前，就是现在，即使我们是伟大的、受尊敬的，也需要来自他的恩宠。\par

21节。\Quote{我的同工提摩太问候你们。}\BibleRef{罗 16:21}\par

请注意惯常的赞美：\Quote{还有路求、雅松和索西帕特，我的亲族问候你们。}\par

路加也提到这位雅松，并展示了他的勇敢，当时说他们\Quote{拉}他\Quote{到城里的官长面前，喊叫着}等等。\BibleRef{宗 17:5}很可能其余的人也都是有名望的人。因为他不会仅仅提及亲族，除非他们也在虔诚上像他一样。\par

22节。\Quote{我特耳爵，代笔写这封信的，问候你们。}\BibleRef{罗 16:22}\par

这也不是小事——作保禄的代笔人。尽管如此，他说这话并非为自己颂扬，而是为了使他们因这服事而对他怀有热烈的爱。\par

23节。\Quote{该犹，我的主人（\OriginalTerm{主人}{\GreekText{ξένος}}），也是全教会的，问候你们。}\BibleRef{罗 16:23}\par

看他给他戴上了怎样的冠冕，为他见证如此伟大的款待，并把整个教会带进这人的家里！因为这里用的 \OriginalTerm{主人}{\GreekText{ξένον}} 一词，意思是主人，而不是客人。但当你听说他是保禄的主人时，不要只钦佩他的慷慨，也要钦佩他生活的严谨。因为除非他配得上保禄的卓越，保禄绝不会住在他那里；因为保禄竭力超越许多基督的命令，绝不会违背那条命令，即要非常谨慎地对待接待我们的人，并住在\Quote{相称}的人那里。\BibleRef{玛 10:11}\Quote{厄辣斯托，本城的财务官，问候你们，还有夸尔托兄弟。}他特意加上\Quote{本城的财务官}，因为正如他写给斐理伯人时说的\Quote{凯撒家里的人问候你们}\BibleRef{斐 4:22}，为表明福音已抓住了大人物，同样，这里他也提到职衔，目的是相同的，并表明对留意的人而言，财富、政务的忧虑或任何其他事物都不是阻碍。\par

24节。\Quote{愿我们的主耶稣基督的恩宠与你们众人同在。阿们。}\BibleRef{罗 16:24}\par

看，我们从头到尾都应以什么开始和结束？因为在这封信中，他以恩宠奠基，也以恩宠封顶，既为他们祈求一切好事的母亲，又提醒他们记起他的全部慈爱。因为这是一个慷慨教师的最好证明：不仅用言语，也用祈祷使学习者受益；为此缘故，有人曾说：\Quote{但我们专务祈祷和宣讲圣言的事。}\BibleRef{宗 6:4}\par

那么，既然保禄已经离去，谁还会为我们祈祷呢？是那些效法保禄的人。但让我们使自己堪当这样的\OriginalTerm{代祷}{\GreekText{συνηγορίας}}，这样我们不仅在此听到保禄的声音，而且将来离世后，也配得看见基督的这位斗士。或者更确切地说，如果我们在此听从他，将来必定会看见他，即使不是站在他身旁，也必定会看见他，在君王宝座附近闪闪发光。在那里，革鲁宾歌唱光荣，色辣芬飞翔，我们将看见保禄与伯多禄一起，作为圣人歌咏团的领唱和领袖，并享受他慷慨的爱。因为如果他在这里时如此爱人，以至于可以选择离世与基督同在，却选择了留在这里，那么他在那里必定会显示更热烈的爱。为此，我甚至爱罗马，尽管确实还有其他理由赞美它——它的伟大、古老、美丽、人口众多、权势、财富和战功。但我撇开这一切，并因这一点而认为它有福：他在世时曾写信给他们，如此爱他们，还在与我们同在时与他们交谈，并在那里结束了生命。因此，这座城市在这方面的名声，超过其他所有方面之和。如同一具巨大而强壮的身体，它拥有这两位圣人的身体作为两只闪烁的眼睛。天空在太阳放射光芒时，也不及罗马城将这两道光芒射向世界各处那样明亮。从那里，保禄将被提，从那里，伯多禄将被提。请想一想，并为之\OriginalTerm{战栗}{\GreekText{φρίξατε}}：当保禄突然从那安息之处与伯多禄一同起来，并被提升到空中与主相遇时，罗马将看到何等景象！\BibleRef{得前 4:17}罗马将向基督献上何等的玫瑰！\BibleRef{依 35:1}这座城市将拥有何等的两顶冠冕！她将束上何等的金链！她将拥有何等的泉源！因此，我钦佩这座城市，不是因其大量黄金，不是因其石柱，不是因其其他炫耀，而是因这些教会的支柱。\BibleRef{格前 15:38}但愿现在我能得以\OriginalTerm{环绕}{\GreekText{περιχυθἥναι}}保禄的身体，紧贴他的坟墓，并看见那身体的尘土——那身体\Quote{补充了}基督\Quote{所欠缺的}\BibleRef{哥 1:24}，那身体带着\OriginalTerm{烙印}{\GreekText{στίγματα}}\BibleRef{迦 6:17}，那身体将福音撒遍各处——是的，那身体的尘土，他藉此到处奔走！那身体的尘土，基督藉此说话，那光比任何闪电更明亮地照射，那声音对魔鬼比任何雷声更可畏！藉此他发出了那有福的声音，说：\Quote{为了我的弟兄，我宁愿自己成为诅咒}\BibleRef{罗 9:3}；藉此他在\Quote{君王面前说话，并不羞愧}\BibleRef{咏 118:46}；藉此我们认识保禄，也藉此认识保禄的主！对我们而言，雷声不如那声音对魔鬼那样可畏！因为如果他们对他的衣服颤抖\BibleRef{宗 19:12}，那么对他的声音就更甚。这声音使他们成为俘虏，这声音洁净了世界，这声音止住了疾病，驱除了邪恶，高举了真理，基督乘着它，它到处与基督同行；正如革鲁宾是那些力量，保禄的声音也是如此，因为正如基督坐于那些力量之上，他也坐在保禄的舌上。因为那舌头已堪当接受基督，只讲基督所喜悦的事，并像色辣芬一样飞向不可言喻的高处！因为有什么比那声音更高昂，它说：\Quote{我深信，无论是天使、掌权者、异能者，还是现在的事、将来的事，是高天、是深渊，还是任何别的受造物，都不能使我们与天主的爱隔绝，这爱是在我们的主基督耶稣内的。}\BibleRef{罗 8:38}你们觉得这话语似乎有何等的翅膀？何等的眼睛？\BibleRef{则 10:12}正因为如此，他说：\Quote{因为我们不是不知道他的阴谋。}\BibleRef{格后 2:11}正因为如此，魔鬼不仅听见他说话就逃跑，甚至看见他的衣服也逃跑。这就是那口腔，我切愿看见它的尘土，基督藉此说了伟大而隐密的事，且比他自己亲口说的更大（因为正如他行事，他也藉门徒说更大的事），圣神藉此向世界发出了那些奇妙的圣言！因为那口腔没有成就什么好事吗？它驱逐魔鬼，赦免罪恶，钳制暴君，封住哲学家的嘴，使世界归向天主，说服野蛮人学习智慧，改变了整个大地秩序。它也按照它的意愿安排天上的事\BibleRef{格前 5:3}，捆绑它所愿捆绑的，并在另一世界释放，\Quote{按着所赐给它的权柄}\BibleRef{格后 13:10}。不仅是那口腔，我也切愿看见那心灵的尘土，人若称它为世界之心、无数祝福的源泉、我们生命的开端和元素，也不算错。因为生命之神从它里面全部涌出，并通过基督的肢体分配，不是通过动脉发出，而是通过善行的自由选择。这心灵如此宽广，以至能容纳整个城市、民族和邦国。他说：\Quote{我们的心是宽大的。}\BibleRef{格后 6:11}然而，正是这使心宽大的爱，常常使这个本已宽广的心受窘迫和压抑。因为他说：\Quote{我心中有许多\OriginalTerm{患难}{\GreekText{θλίψεως}}和\OriginalTerm{焦虑}{\GreekText{συνοχἥς}}，写信给你们。}\BibleRef{格后 2:4}我渴望看见那心即使在分解之后，它为每个失丧的人燃烧，为那些流产的孩子再次受生产的痛苦\BibleRef{迦 4:19}，它看见天主（因为经上说：\Quote{心里洁净的人必看见天主}\BibleRef{玛 5:8}），它成为祭品（因为\Quote{天主所要的祭品，就是痛悔的心}\BibleRef{咏 50:17}），它比诸天更高，比世界更广，比太阳光更亮，比火更热，比金刚石更坚强，它涌出江河（因为经上说：\Quote{从他心中要流出活水的江河}\BibleRef{若 7:38}），它里面有泉源涌出，灌溉的不是大地表面，而是人的灵魂，从中不仅江河，而且泪泉日夜涌流，它活的新生命不是我们这种（因为他说：\Quote{我生活，不是我生活，而是基督在我内生活}\BibleRef{迦 2:20}，所以保禄的心是基督的心，是圣神的版，是恩宠的书）；它为别人的罪而战栗（因为他说：我恐怕\Quote{我白白地为你们劳苦了}\BibleRef{迦 4:11}；\Quote{恐怕蛇诱惑了厄娃}\BibleRef{格后 11:3}；\Quote{恐怕我来时，你们不是我所希望的那样}\BibleRef{格后 12:20}）；它既为自己害怕，也充满信心（因为他说：我恐怕\Quote{我传福音给别人，自己反而被弃绝了}\BibleRef{格前 9:27}；又说：\Quote{我深信无论是天使还是掌权者，都不能使我们分离}\BibleRef{罗 9:3}）；它配得爱基督，没有人像他那样爱基督；它轻视死亡和地狱，却被弟兄的眼泪所破碎（因为他说：\Quote{你们为什么哭泣，使我心碎呢？}\BibleRef{宗 21:13}）；它极其忍耐，却不能忍受与得撒洛尼人分开片刻\BibleRef{得前 2:17}！我切愿看见那双手的尘土，它们曾戴着锁链，藉着它们的覆手，圣神被赐予，藉着它们，神圣的著作得以写成（因为\Quote{请看我亲手写给你们的字是何等的大}\BibleRef{迦 6:11}；又说：\Quote{我保禄亲笔问安}\BibleRef{格前 16:21}），那双手，毒蛇看见它们就\Quote{掉进火里}\BibleRef{宗 28:5}。我切愿看见那双眼睛的尘土，它们曾荣耀地失明，又为世界的得救而复明；它们尚在肉身时，就配得看见基督；它们看见地上的事物，却似乎没有看见；它们看见那看不见的事物；它们不曾睡觉，在午夜警醒；它们不像眼睛那样被影响。我也切愿看见那双脚的尘土，它们跑遍世界而不疲倦；它们曾当监狱震动时被锁在木狗里；它们走过了可居住和不可居住的地方；它们行走了如此多的旅程。我何必细说各个部分？我切愿看见那坟墓，那里安放着正义的武器、光明的武器，那如今活着但在世时已死的肢体；在这些肢体中，基督活着，它们被钉在十字架上与世界分离，它们是基督的肢体，披戴着基督，是圣神的宫殿、神圣的建筑，\Quote{被圣神束缚}\BibleRef{宗 20:22}，被钉在敬畏天主上，带着基督的烙印。这身体是那座城的城墙，比所有塔楼和千万城垛更安全。与他一起的还有伯多禄的身体。因为保禄在世时尊敬伯多禄。他\Quote{上耶路撒冷去见伯多禄}\BibleRef{迦 1:18}，因此即使在离世后，恩宠也屈尊赐给他们同样的住所。我切愿看见那属灵的雄狮。因为如同一头\OriginalTerm{喷火}{\GreekText{πὕρ} \GreekText{ἀφιεὶς}}的狮子\BibleRef{歌 2:15}冲向狐狸群，他冲入魔鬼和哲学家的团伙，如同霹雳一般，猛扑进魔鬼的阵营。\BibleRef{路 13:32}因为他甚至没有来摆阵对抗魔鬼，因为魔鬼如此害怕和战栗，以至于只要看见他的影子、听见他的声音，便远远地逃跑。他这样把犯奸淫的人交给魔鬼，虽然远隔，却又把他从魔鬼手中夺回\BibleRef{格前 5:5}；同样，他也使其他的人这样受教训，\Quote{不再亵渎}\BibleRef{弟前 1:20}。试想他如何派遣自己的臣仆去对抗魔鬼，激励他们，装备他们。一时他对厄弗所人说：\Quote{我们战斗不是对抗血和肉，而是对抗掌权者和异能者。}\BibleRef{弗 6:12}然后他又把我们的奖品放在天上。他说：因为我们斗争的并不是地上的事物，而是天堂和天上的事物。对另一些人，他说：\Quote{你们不知道我们要审判天使吗？更何况今生的事呢？}\BibleRef{格前 6:3}因此，让我们将这些都放在心上，勇敢站立；因为保禄是一个人，与我们同享一样的人性，其他一切都与我们相同。但由于他向基督显示了如此伟大的爱，他升到诸天之上，与天使一同站立。同样，如果我们也能稍微振作，在我们内点燃那火焰，我们也能效法那位圣人。因为如果这是不可能的，他就不会大声呼喊说：\Quote{你们要效法我，如同我效法基督一样。}\BibleRef{格前 11:1}因此，让我们不仅钦佩他，不仅被他打动，还要效法他，这样我们将来离世时，也配得看见他，并分享那不可言喻的荣耀。愿天主因我们的主耶稣基督的恩宠和爱人，赐予我们大家都能达到那荣耀。愿光荣归于圣父、圣子及圣神，从现在直到永远。阿们。\par

\SourceRecord{Translated by J. Walker, J. Sheppard and H. Browne, and revised by George B. Stevens.；Nicene and Post-Nicene Fathers, First Series；Vol. 11.；Edited by Philip Schaff.；Buffalo, NY: Christian Literature Publishing Co.,；1889.；<http://www.newadvent.org/fathers/210232.htm>.}
